% !TEX TS-program = lualatex
% !BIB TS-program = biber
% This file must be compiled with LuaLaTeX + biber.
% Background highlighting in the citation report relies on the lua-ul package,
% which only works under LuaLaTeX.
% Run:  lualatex citation_report.tex && biber citation_report && lualatex citation_report.tex && lualatex citation_report.tex


\documentclass{article}

% XeLaTeX + Unicode + font handling
\usepackage{fontspec}
\usepackage{polyglossia}
\usepackage[letterpaper,top=2cm,bottom=2cm,left=3cm,right=3cm,marginparwidth=1.75cm]{geometry}

% Languages: English
\setdefaultlanguage{english}

% Font setup
\setmainfont{Linux Libertine O}
\setsansfont{Linux Libertine O}
\setmonofont{DejaVu Sans Mono}

% Bibliography
\usepackage[backend=biber, style=apa]{biblatex}
\addbibresource{references.bib}

% Math packages
\usepackage{amsmath,amssymb,amsthm,amsfonts,mathtools,bm}
\usepackage{siunitx,physics,chemformula,tensor,hepnames,braket}
\usepackage[version=4]{mhchem}

% Tables and figures
\usepackage{adjustbox,tabularx,graphicx}
\usepackage[ruled,vlined]{algorithm2e}

% Misc
\usepackage{caption,csquotes}

% Layout and links
\usepackage[colorlinks=true,allcolors=blue]{hyperref}


% Custom length
\newlength{\sixtypercentpage}
\setlength{\sixtypercentpage}{0.6\paperheight}


% Packages for citation annotation (color highlighting and PDF comments)
% Citation reports compile with LuaLaTeX so we can use lua-ul, which highlights
% at the node level and works with math, CJK, and fragile commands natively.
\usepackage{xcolor}
\usepackage[color={0.5 0.5 0.5}]{pdfcomment}
\usepackage{graphicx}
\usepackage{luacolor}  % required by lua-ul for colored highlights
\usepackage{lua-ul}
\usepackage{tikz}

% Define colors for highlighting
% hldarkgreen (dark green - strong support)
\definecolor{hldarkgreen}{RGB}{89,171,91}
% hllightgreen (light green - partial support)
\definecolor{hllightgreen}{RGB}{191,255,207}
% hlorange (orange - unsupported)
\definecolor{hlorange}{RGB}{255,189,89}
% hlred (red - contradictory)
\definecolor{hlred}{RGB}{214,105,89}

% Highlighting commands backed by lua-ul's \highLight (node-level, math-safe)
\newcommand{\hldarkgreen}[1]{\highLight[hldarkgreen]{#1}}
\newcommand{\hllightgreen}[1]{\highLight[hllightgreen]{#1}}
\newcommand{\hlorange}[1]{\highLight[hlorange]{#1}}
\newcommand{\hlred}[1]{\highLight[hlred]{#1}}

% Display citekey in bibliography entries (prepended in gray)
\renewbibmacro*{begentry}{%
  \textcolor{gray}{[\thefield{entrykey}]}\space%
}

\begin{document}


% Add ThesisAI logo at the top center
\vspace{-1em}
\begin{center}
\includegraphics[width=0.3\textwidth]{logo.png}
\end{center}

\vspace{0.5em}

\section*{Citation Report}

\noindent\textbf{Title:} Algorithmic Addiction in Social Recommender Systems: Ethical AI Audit Methodology Prioritising Silent Profiling Defaults, Friction Loaded Opt Out Channels and Asymmetric Notification Controls, and a Pro Social Pivot Strategy for the EU Digital Fairness Act

\noindent\textbf{Author:} nirwanmohit4@gmail.com

\noindent\textbf{Date:} September 07, 2026

\noindent\textbf{Number of Papers:} 110

\noindent\textbf{Citation Coverage:} 90\%

\noindent\textbf{Strict Evidence:} Yes

\noindent\textbf{Citation Style:} apa

\noindent\textbf{Citation Quality Score:} 96.3\%

\vspace{1em}


\noindent
This document is an automated citation quality report generated by ThesisAI. Each sentence in your document has been analyzed and color-coded based on how well it is supported by the source materials.

\vspace{0.5em}

\noindent
\begin{center}
\begin{tikzpicture}[scale=1.0]

 \fill[hldarkgreen] (0,0) rectangle (8.099717779868298,0.6);

 \node[black] at (4.049858889934149,0.3) {57.9\%};

 \fill[hllightgreen] (8.099717779868298,0) rectangle (13.486359360301035,0.6);

 \node[black] at (10.793038570084667,0.3) {38.5\%};

 \fill[hlorange] (13.486359360301035,0) rectangle (14.0,0.6);

\end{tikzpicture}
\end{center}

\vspace{0.5em}

\noindent
The bar shows the distribution of citation quality across all 2126 sentences in your document. Green colors (dark and light) indicate well-supported content, representing \textbf{96.3\%} of your document.

\vspace{1em}



\vspace{1em}

\noindent\begin{tabular}{@{}>{\centering\arraybackslash}m{0.08\textwidth}|>{\centering\arraybackslash}m{0.17\textwidth}|m{0.45\textwidth}|>{\centering\arraybackslash}m{0.2\textwidth}@{}}
\textbf{Color} & \textbf{Classification} & \textbf{Explanation} & \textbf{Action Required} \\
\hline
\tikz\fill[hldarkgreen] (0,0) circle (1.2ex); & Fully supported & The sentence is well-supported by the provided sources with clear, direct evidence. & No action needed \\
\hline
\tikz\fill[hllightgreen] (0,0) circle (1.2ex); & Academic synthesis & Valid academic synthesis combining information from multiple sources, or logical inferences drawn from provided evidence. This represents good academic writing and is acceptable practice. Light green indicates the claim is supported through synthesis rather than direct quotation. Common in comparative statements, conclusions, and summary sentences. & No action needed \\
\hline
\tikz\fill[hlorange] (0,0) circle (1.2ex); & Not supported & The sentence lacks support from the provided sources, contains unsupported claims, or is completely unrelated to the sources. & Revision needed \\
\hline
\tikz\fill[hlred] (0,0) circle (1.2ex); & Contradictory & The sentence directly contradicts the provided sources or makes claims that are opposite to what the sources state. & Must correct \\
\end{tabular}

\newpage

\noindent\textbf{Understanding the PDF Comments:} Click on the \textbf{comment icons} in the margins to see detailed explanations for each color rating. These comments reference specific sources by their \textbf{citekey} and \textbf{page number}. Each entry in the \textbf{References} section at the end of this document shows the citekey in gray brackets (e.g., \textcolor{gray}{[5663359]}) at the beginning of the reference, allowing you to easily match the citekeys mentioned in the PDF comments to their corresponding full references. \textbf{Important:} The page numbers shown (e.g., ``p.13'') refer to the actual page number in the PDF file, not the page number printed within the document itself.

\vspace{1em}

\noindent\textbf{How to Improve Orange or Red Sections:} If your document contains orange or red highlighted sentences, consider the following approaches to improve citation quality:

\begin{itemize}
 \item \textbf{Review and Justify:} Carefully review the highlighted sentences and determine whether they genuinely require editing. Sometimes the AI may flag content that is actually acceptable or necessary for your argument.

 \item \textbf{Rewrite or Remove:} Use follow-up prompts in ThesisAI to remove or rewrite specific text parts that are unsupported or contradictory. You can ask the AI to adjust problematic statements to better align with your sources.

 \item \textbf{Add More Literature (Recommended):} The best approach is to find additional literature sources that provide evidence for sections lacking support. Upload these new sources to ThesisAI and regenerate the document. This will improve citation coverage and strengthen your academic work with more comprehensive source material. \textbf{Rule of thumb:} Provide at least 2 papers per page of output. For example, when generating a 35-page document, upload at least 70 papers to ensure sufficient source material for comprehensive citation coverage.

 \item \textbf{Verify Contradictions:} For red-highlighted contradictory statements, carefully review the source material and either correct the statement or reconsider whether the source appropriately supports your argument.
\end{itemize}

\noindent For any questions or feedback about this citation report, please contact us via the support chat at \url{www.thesisai.io}.

\newpage

% Table of Contents
\tableofcontents

\newpage


\section{Problem Formulation and Research Context}

\subsection{Motivation and Practical Relevance}
\label{sec:motivation_practical_relevance} \hllightgreen{Growing concern about manipulative design in digital markets provides a direct impetus for an ethics-driven audit of social recommender systems.}\pdfcomment[icon=Note]{Works on dark patterns and algorithmic manipulation frame growing concern as motivating regulatory and audit responses (see 7896734 p.7; 8745065 p.2), so this is a reasonable synthesis of the literature.} \hldarkgreen{Empirical work on dark patterns documents subscription flows where cancelling is intentionally harder than signing up, countdown timers that create illusory urgency, and misleading brand associations such as using the "Maxbosch" label alongside BOSCH imagery to influence consumer perception \parencites[390]{7896734}.}\pdfcomment[icon=Note]{Source 7896734 explicitly documents examples including making cancellation harder than signup, countdown timers creating urgency, and the 'Maxbosch' / BOSCH imagery misleading association (7896734 p.7).} \hldarkgreen{These findings align with behavioral research on scarcity cues and loss framing in digital retail, where messages such as "only two left" or expiring discounts transform a routine purchase into a time-pressured decision and can shift a neutral non-purchase into a perceived loss \parencites[135]{8745065}.}\pdfcomment[icon=Note]{Behavioral and scarcity research linking messages like 'only two left' and expiring discounts to time-pressured decisions is directly described (8745065 p.2).} \hldarkgreen{As AI-driven recommendation and personalization scale these techniques, the risk of systematic exploitation of predictable vulnerabilities increases, which makes a targeted audit methodology practically necessary \parencites[135]{8745065}[390]{7896734}.}\pdfcomment[icon=Note]{The literature argues that AI-driven personalization amplifies manipulative techniques and calls for targeted audits (7896734 p.7; 8745065 p.2), supporting the claim and the cited reference.} \hllightgreen{European regulatory developments are moving exactly toward curbing such manipulation.}\pdfcomment[icon=Note]{EU regulatory developments are presented as addressing dark patterns and manipulation in the cited sources (7896734 p.7; 6502223 p.12), so this is a reasonable summary synthesis.} \hldarkgreen{The Digital Services Act (DSA) already addresses platform accountability for harmful content and requires very large online platforms to assess and mitigate "systemic risks," explicitly including risks to civic discourse and electoral processes, supported by transparency reports, non-profiled recommendation options, and independent audits \parencites[12]{6502223}[2]{9360054}.}\pdfcomment[icon=Note]{Source 6502223 describes the DSA's VLOP obligations, including systemic risk assessment for civic discourse, transparency reporting, non-profiled recommendation options, and independent audits (6502223 p.12).} \hldarkgreen{In parallel, the European Union has adopted the AI Act, which prohibits AI systems that exploit psychological vulnerabilities and bans subliminal manipulation or manipulative effects that can significantly harm user behaviour, situating such practices in the "unacceptable risk" tier \parencites[12]{6502223}[390]{7896734}.}\pdfcomment[icon=Note]{Both sources state the AI Act's prohibited category includes exploiting psychological vulnerabilities and bans subliminal/manipulative effects (6502223 p.12; 7896734 p.7).} \hldarkgreen{Planned measures such as a Digital Fairness Act are intended to regulate addictive design and unfair personalization as a distinct problem space, filling gaps between the DSA and classical consumer protection rules \parencites[390]{7896734}.}\pdfcomment[icon=Note]{Source 7896734 reports a planned Digital Fairness Act aiming to regulate addictive design and unfair personalization to fill gaps between the DSA and consumer rules (7896734 p.7).} \hllightgreen{These converging frameworks create concrete legal exposure for platforms relying on engagement-optimizing, AI-driven interface pressure.}\pdfcomment[icon=Note]{Given the DSA, AI Act, and related proposals described in the sources, it is a reasonable inference that platforms face increased legal exposure for engagement-optimizing manipulative designs (6502223 p.12; 7896734 p.7).} \hllightgreen{At the same time, global and national AI governance proposals stress ethical responsibility and non-discrimination as binding principles, not optional aspirations.}\pdfcomment[icon=Note]{Multiple governance proposals emphasize ethical responsibility and non-discrimination (2977039 p.61; 8367105 p.12), so treating these as binding principles in some frameworks is a defensible synthesis.} \hldarkgreen{Kostenko's AI Law Model requires AI systems to be created and operated in line with fundamental values such as human dignity, autonomy, and psycho-emotional comfort, and explicitly prohibits designs that manipulate user behaviour, cause digital addiction, or suppress autonomy \parencites[61--62]{2977039}.}\pdfcomment[icon=Note]{The AI Law Model (2977039) requires systems to respect human dignity, autonomy and psycho-emotional comfort and prohibits AI designs that manipulate behavior or cause digital addiction (2977039 p.61--62).} \hldarkgreen{The same model formulates a principle of fairness, non-discrimination, and inclusion that bans AI systems which consolidate social inequality or are built on distorted data, and calls for equity metrics, bias audits, and participation of vulnerable groups in design \parencites[73--74]{2977039}.}\pdfcomment[icon=Note]{2977039 sets out fairness, non-discrimination, inclusion requirements and calls for equity metrics, audits, and participation of vulnerable groups (2977039 p.73--74).} \hldarkgreen{For high-risk domains, ethical responsibility and fairness become certification criteria, with regulators empowered to suspend systems, remove unethical functionality, and require public ethics impact reports \parencites[61--62]{2977039}[73--74]{2977039}.}\pdfcomment[icon=Note]{2977039 makes ethical responsibility and fairness criteria for certification in high-risk domains and empowers regulators to suspend systems, remove functionality, and require public ethics reporting (2977039 p.61--62; p.83).} \hllightgreen{These provisions underline why an ethical audit of engagement-maximizing recommenders cannot be purely technical; it must systematically examine autonomy, discrimination, and psycho-emotional impact.}\pdfcomment[icon=Note]{Given the model's emphasis on autonomy, non-discrimination, and psycho-emotional comfort (2977039 p.61--62), it is a reasonable synthesis that ethical audits must examine those dimensions.} \hllightgreen{Digital literacy and user protection gaps further increase the practical urgency of a structured audit toolkit.}\pdfcomment[icon=Note]{Survey and policy analyses identify gaps in digital literacy and user protection that increase urgency for structured tools (9360054 p.23; 9360054 p.2), so this is a supported inference.} \hldarkgreen{Survey evidence on awareness of GDPR, the DSA, and the AI Act shows high self-assessed basic digital literacy yet persistent difficulties in avoiding fraud, recognizing hidden content, and behaving consciously in complex threat environments.}\pdfcomment[icon=Note]{Source 9360054 reports high self-assessed basic digital literacy alongside persistent difficulties in avoiding fraud, recognizing hidden content, and managing complex threats (9360054 p.23).} \hldarkgreen{Users report only moderate familiarity with EU policies and are skeptical that existing regulation alone guarantees safety, while simultaneously expressing optimism about educational programs and stronger rules \parencites[2]{9360054}.}\pdfcomment[icon=Note]{9360054 finds respondents report moderate familiarity with EU policies and skepticism that regulation alone guarantees safety, while expressing optimism about education and stronger rules (9360054 p.23).} \hllightgreen{This combination of confidence in basic skills and vulnerability to sophisticated manipulations suggests that the burden of protection cannot rest solely on individual users.}\pdfcomment[icon=Note]{Combining high self-assessed basic skills with documented vulnerabilities supports the reasonable conclusion that relying solely on users is insufficient (9360054 p.23).} \hllightgreen{Ethical audits that are aligned with regulatory expectations and translated into platform governance processes directly respond to this revealed need \parencites[2]{9360054}[12]{6502223}.}\pdfcomment[icon=Note]{Sources describe regulatory expectations and the need for audits and governance alignment (2977039 p.61--83; 6502223 p.12), so ethical audits aligned with regulation respond to the identified need.} \hllightgreen{Industry practice and emerging self-regulation show both demand and limitations for such tools.}\pdfcomment[icon=Note]{Empirical and review sources document firms issuing principles and self-regulation while noting their limits, indicating both demand and limitations for audit tools (7461574 p.9; 8367105 p.12).} \hldarkgreen{Large firms have issued AI principles and created ethics divisions, and international instruments like the OECD AI Principles and UNESCO's Recommendation on AI Ethics articulate fairness, transparency, and human-centricity as guiding norms \parencites[9]{7461574}[12]{6502223}[12]{8367105}.}\pdfcomment[icon=Note]{Source 7461574 documents large firms issuing AI principles and ethics divisions, and 8367105 links these efforts to OECD and UNESCO normative initiatives (7461574 p.9; 8367105 p.12).} \hldarkgreen{Yet studies on the "principles-to-practice gap" highlight that high-level declarations rarely specify operational procedures for data governance, model documentation, or post-deployment monitoring, which leads to inconsistent ethical adoption and criticism of self-regulation for lack of transparency \parencites[9]{7461574}[12]{8367105}.}\pdfcomment[icon=Note]{8367105 explicitly discusses the 'principles-to-practice gap' where high-level declarations lack operational detail, leading to inconsistent adoption and criticism (8367105 p.12).} \hldarkgreen{Pipeline-oriented interventions such as model cards, datasheets, algorithmic impact assessments, and red-teaming offer building blocks but are often reactive or confined to single stages of the lifecycle \parencites[12]{8367105}.}\pdfcomment[icon=Note]{8367105 describes model cards, datasheets, AIAs, and red-teaming as useful building blocks that are often reactive or limited to single lifecycle stages (8367105 p.12).} \hllightgreen{A structured Ethical AI Audit tailored to engagement and addiction risks can therefore translate these scattered tools into a targeted, practically usable methodology for recommender systems.}\pdfcomment[icon=Note]{Given the scattered nature of existing tools and calls for lifecycle approaches in the literature (8367105 p.12; 2977039 p.61--83), proposing a structured audit tailored to engagement/addiction risks is a reasonable synthesis.} \hllightgreen{Finally, multiple sources stress that accountability mechanisms and registries of responsible AI actors are becoming prerequisites for deploying high-risk systems in sectors with significant public impact such as education and mental health \parencites[83]{2977039}[11]{5419804}.}\pdfcomment[icon=Note]{2977039 advocates registers and accountability for responsible AI actors and other sources emphasize accountability mechanisms for high-impact sectors (2977039 p.83; 5419804 p.11), supporting this claim.} \hllightgreen{Human-in-the-loop models, open audit trails, and annual ethics impact reporting are framed as conditions for sustaining trust and preventing "AI colonialism" or unilateral technological dependence \parencites[21]{6805789}[108]{2977039}[11]{5419804}.}\pdfcomment[icon=Note]{2977039 requires human oversight and annual ethics impact reports (2977039 p.61--62), and 5419804 discusses preventing 'AI colonialism' in scaling interventions, so the synthesis is supported.} \hllightgreen{Social recommender systems that shape attention, affect, and consumption patterns at scale form a comparable arena of public interest, which makes the development of a transparent, ethics-grounded audit and pro-social pivot strategy not only academically relevant but also operationally and regulatorily urgent \parencites[21]{6805789}[61--62]{2977039}[390]{7896734}.}\pdfcomment[icon=Note]{Both the AI Law Model's public-interest framing and dark-pattern literature link recommender harms to public concerns and regulatory urgency (2977039 p.61--74; 7896734 p.7), so the call for transparent audits is a supported synthesis.}

\subsection{Research Objectives and Central Questions}
\label{sec:research_objectives_central_questions} \hllightgreen{Clarifying the aims of this study requires aligning them with existing analytical work on manipulative design, political recommender risks, and ethical AI governance.}\pdfcomment[icon=Note]{Multiple provided sources discuss manipulative design, political recommender risks, and AI governance (e.g., 8745065 p.136 on urgency marketing; 6502223 p.6 on recommender-driven political risks; 1456534 p.2 on EU governance), so aligning study aims with that literature is a reasonable synthesis.} \hldarkgreen{Existing research on urgency marketing and dark patterns in e commerce shows that interface tactics such as scarcity cues, defaults, and personalized prompts are systematically organized around psychological vulnerabilities and short term revenue objectives, not isolated creative experiments.}\pdfcomment[icon=Note]{8745065 p.136 explicitly describes scarcity cues, defaults, personalization and their organization around psychological vulnerabilities and short-term commercial objectives; 7896734 p.390 also documents scarcity timers and manipulative UI practices.} \hldarkgreen{This project therefore adopts as a first objective the development of a secondary research synthesis that maps such manipulation oriented design choices in recommender driven environments, using a coding matrix that draws on constructs like scarcity, psychological targeting, and revenue optimization identified in prior qualitative thematic analyses of digital retail practice \parencites[136]{8745065}.}\pdfcomment[icon=Note]{8745065 p.136 describes a deductive coding frame including scarcity, psychological targeting, and revenue optimization and advocates a thematic synthesis of urgency-based marketing, directly supporting the described objective and coding-matrix approach.} \hllightgreen{Regulatory ambition around AI driven manipulation sets a second objective.}\pdfcomment[icon=Note]{The document cites regulatory texts and ambitions around AI-driven manipulation in multiple sources (e.g., 6502223 and 7896734), so stating that regulatory ambition constitutes a second objective is a supported framing inference.} \hldarkgreen{The EU AI Act prohibits systems that exploit psychological vulnerabilities and bans subliminal manipulation or manipulative effects that can significantly harm behaviour, while the Digital Services Act requires very large platforms to assess and mitigate systemic risks to civic discourse and electoral processes and to offer at least one non profiled recommendation option \parencites[2248]{6502223}.}\pdfcomment[icon=Note]{7896734 p.390 and 6502223 p.12 explicitly state that the EU AI Act bans exploitation of psychological vulnerabilities and subliminal manipulation, and 6502223 p.12 describes DSA obligations for VLOPs including systemic risk assessments and a non-profiled recommendation option.} \hlorange{Sanikidze documents how this preventive approach is complemented by the forthcoming Digital Fairness Act proposal that targets addictive design and unfair personalization as a distinct problem space \parencites[390]{7896734}.}\pdfcomment[icon=Note]{While sources (e.g., 7896734 p.390) mention a planned Digital Fairness Act addressing addictive design and personalization, none of the provided sources identify or attribute this account to an author named 'Sanikidze'.} \hllightgreen{Against this background, the report aims to translate these high level obligations into operational audit questions for social recommender systems, treating the Digital Fairness Act as a policy trajectory and using the AI Act and DSA as anchor frameworks \parencites[390]{7896734}[2248]{6502223}.}\pdfcomment[icon=Note]{6502223 (e.g., p.12, p.2248) and other regulatory sources document the AI Act and DSA as core frameworks; treating the DFA as a policy trajectory and translating obligations into audit questions is a reasonable methodological framing supported by these texts.} \hllightgreen{A third objective concerns the methodological stance.}\pdfcomment[icon=Note]{Several methodological sections in the provided literature set out objectives and methodological stances (e.g., 8200246 and 8745065), so declaring a third objective about methodology is a supported structural statement.} \hldarkgreen{Several authors adopt qualitative, mixed, or secondary designs to examine AI systems, regulatory texts, and institutional reports, stressing the value of triangulating regulatory intent, empirical platform practices, and documented harms \parencites[7]{6295677}[6]{6502223}[3]{5142478}.}\pdfcomment[icon=Note]{5142478 p.24 describes the convergent-parallel mixed-methods approach and 6295677 (analysis section) and 6502223 p.6 discuss triangulation of regulatory intent, platform practices, and harms, directly supporting this claim.} \hllightgreen{Following this pattern, the present study commits to a qualitative, interpretive synthesis of secondary materials rather than proprietary telemetry or direct source code inspection, and formulates an explicit limitation that black box algorithms will be assessed through documented UI behaviour, regulator findings, and existing audits.}\pdfcomment[icon=Note]{Multiple sources justify secondary qualitative designs and reliance on documented UI behaviour and audits (e.g., 8200246 pp.20224--20227; 6502223 p.6), so committing to interpretive synthesis and the stated limitation is a supported methodological choice.} \hlorange{Garr's convergent parallel design for content curation and bias metrics, and Chakraborty et al.'s mixed methods work on adaptive education ethics, both illustrate how policy documents, academic work, and case studies can be combined to evaluate algorithmic risks without internal access to models \parencites[3]{5142478}[5]{1308731}.}\pdfcomment[icon=Note]{5142478 describes convergent-parallel mixed methods generally, but the specific named authors 'Garr' and 'Chakraborty et al.' and their described examples are not found in the provided sources, so the attribution is unsupported by the materials given.} \hllightgreen{Based on this design choice, the central questions guiding the inquiry are framed at the level of observable outputs and regulatory alignment.}\pdfcomment[icon=Note]{Given the documented methodological stance in sources advocating output-level and regulatory alignment analyses (e.g., 8200246 and 6502223), framing the central questions at observable outputs and regulatory alignment is a reasonable inference supported by the literature.} \hldarkgreen{One question interrogates how recommendation architectures that optimize engagement may embed structural incentives similar to those identified in political recommender systems, where optimization objectives and training dynamics contribute to democratic harms and polarization risks \parencites[6]{6502223}.}\pdfcomment[icon=Note]{6502223 p.6 explicitly analyzes how recommender architectures and optimization objectives can generate structural incentives that relate to polarization and democratic harms, directly supporting this interrogation.} \hllightgreen{Another question asks how far existing and emerging EU instruments such as the AI Act, the DSA, the Digital Markets Act, and the projected Digital Fairness Act can be interpreted as a de facto governance regime for addictive design, interface pressure, and AI mediated personalization in social platforms \parencites[390]{7896734}[2248]{6502223}[2]{1456534}.}\pdfcomment[icon=Note]{Sources (e.g., 6502223 p.12; 1456534 p.2; 7896734 p.390) discuss the AI Act, DSA, DMA and the proposed DFA and their collective governance effects, so treating them as a de facto governance regime is a defensible synthesis.} \hllightgreen{Research on ethical AI marketing and corporate digital responsibility provides a final objective: to design an Ethical AI Audit and pro social pivot strategy that can be embedded into organisational governance rather than treated as a one off compliance exercise.}\pdfcomment[icon=Note]{8200246 (pp.20224--20232) and related literature discuss Ethical AI marketing, CSR and organisational governance, supporting the objective to design an Ethical AI Audit and pro-social strategy embedded in governance rather than as one-off compliance.} \hlorange{Asante and Adeoba identify Communities of Practice and corporate social responsibility as mechanisms that turn internal ethical learning into policies, reporting routines, and algorithm audits, with stakeholder trust emerging as a core outcome of responsible AI mediated marketing \parencites[11]{8200246}[16]{8200246}.}\pdfcomment[icon=Note]{Although 8200246 discusses Communities of Practice, CSR, and stakeholder trust, the specific attribution to 'Asante and Adeoba' is not present in the provided sources, so the named citation is unsupported here.} \hlorange{Catapang's account of Ethics by Design for AI shows how values such as human agency, fairness, and transparency can be mapped to concrete tasks across the development lifecycle instead of remaining abstract principles \parencites[12]{8367105}.}\pdfcomment[icon=Note]{8367105 discusses 'Ethics by Design' (EbD-AI) and mapping values to tasks, but the specific attribution to 'Catapang' is not evident in the provided materials, so the precise citation/name is unsupported.} \hllightgreen{This thesis therefore sets as a further objective the formulation of a working audit toolkit and decision matrix that is practically usable by management teams, explicitly linked to CSR style governance arrangements, and oriented toward sustaining engagement within the constraints of EU risk based AI regulation \parencites[11]{8200246}[12]{8367105}[4]{1456534}.}\pdfcomment[icon=Note]{8367105 (on EbD-AI operationalization) and 1456534 (on EU risk-based regulation) support formulating audit toolkits linked to CSR and constrained by EU regulation, making this objective a supported synthesis.}

\subsection{Scope, Boundaries and Key Assumptions}
\label{sec:scope_boundaries_assumptions} \hldarkgreen{This project is explicitly confined to qualitative secondary research on AI and recommender governance rather than primary empirical measurement or direct access to proprietary systems.}\pdfcomment[icon=Note]{Sources state the study is secondary qualitative and not primary empirical: 6502223 p.6 notes no original empirical data, and 8200246 p.11--16 describes a secondary qualitative, literature-based reflexive thematic approach.} \hldarkgreen{Methodological designs that rely on literature based inquiry and reflexive thematic analysis to synthesise academic, policy, and industry documents provide the closest template: they treat existing publications, regulatory texts, and professional reports as the primary data and build conceptual understanding inductively \parencites[8]{8200246}[6]{6502223}.}\pdfcomment[icon=Note]{Source 8200246 p.11 and p.8 explicitly describes a literature-based inquiry using Reflexive Thematic Analysis as the methodological template.} \hllightgreen{In line with such approaches, the present study systematically interprets prior work on political recommendation risks, ethical AI marketing, and AI law models without collecting new user telemetry or interviewing practitioners.}\pdfcomment[icon=Note]{Supported by 6502223 p.6 (synthesis of prior work without original data) and 8200246 p.11--16 which frames secondary interpretation of prior literature rather than new telemetry or interviews.} \hllightgreen{Claims about algorithmic addiction and manipulative interface patterns are therefore mediated through documented analyses of recommendation architectures, regulatory regimes, and organisational governance rather than grounded in fresh behavioural experiments \parencites[6]{6502223}[11]{8200246}.}\pdfcomment[icon=Note]{6502223 p.16 and 2977039 p.96--108 document analyses and regulatory discussion of recommendation systems and dark pattern concerns; 8200246 describes secondary analysis rather than fresh experiments.} \hldarkgreen{Only publicly available descriptions of platform behaviour and regulatory assessments fall within scope.}\pdfcomment[icon=Note]{Multiple sources define scope as using published materials and documented assessments: 6502223 p.6 and 8200246 p.11--16 state reliance on publicly available publications and documents.} \hllightgreen{Black box AI systems for content curation or ranking are examined through their observable outputs, described objective functions, and regulator discussions, not through source code inspection or access to internal logs.}\pdfcomment[icon=Note]{Sources indicate examination via observable outputs and documented functions rather than internal code: 6502223 p.16 and 2977039 p.108--109 emphasize audits, documentation and external assessment over source code access; 6577872 p.92 further notes limited sharing of code/data in recommender research.} \hldarkgreen{Scholars examining democratic harms from recommendation architectures already highlight that proprietary constraints and trade secrecy limit detailed technical analysis and that many conclusions must rest on documented case studies and partial disclosure by platforms \parencites[6]{6502223}.}\pdfcomment[icon=Note]{2977039 p.107--108 and 6502223 p.6 explicitly state that proprietary constraints and trade secrecy limit technical analysis and force reliance on documented case studies and partial disclosures.} \hldarkgreen{This thesis adopts the same boundary: audits and pro social pivot strategies are framed around externally assessable properties such as explainability obligations, independent audit requirements, and risk classifications rather than unverified reconstructions of internal model weights or training data \parencites[16]{6502223}[244]{2977039}.}\pdfcomment[icon=Note]{2977039 p.243--244 and p.108--109 set out external audit requirements, explainability obligations and risk classifications; 8200246 p.16 frames audits and governance as externally assessable properties.} \hllightgreen{The regulatory frame is similarly delimited.}\pdfcomment[icon=Note]{This is a brief summary claim; the surrounding sources (2977039 p.243--244, 6502223 p.6) delimit regulatory framings, supporting the statement as a synthesis.} \hldarkgreen{Detailed discussion draws on formally adopted EU instruments such as the Digital Services Act and the AI Act, together with extended AI law proposals that specify transparency, non discrimination, and audit obligations for high risk systems, including independent technical, legal, and ethical assessments, algorithmic passports, and national registers of AI systems \parencites[244]{2977039}[243]{2977039}.}\pdfcomment[icon=Note]{6502223 p.6 discusses the DSA and AI Act and their transparency/non-discrimination/audit obligations; 2977039 p.243--244 lists algorithmic passports, independent assessments and national registers.} \hldarkgreen{A second cluster of norms is taken from the AI Law Model's principle of digital neutrality, which prohibits AI systems that distort information processes, embed hidden censorship, or discriminate between users, and requires operators to document decision logic and submit to independent audits, with violations recorded in a public Register of Digital Neutrality Incidents \parencites[109]{2977039}.}\pdfcomment[icon=Note]{2977039 Chapter 32 (p.108--109) defines the principle of digital neutrality including prohibitions on distortion, hidden censorship, discrimination, documentation of decision logic, independent audits and a public Register of Digital Neutrality Incidents.} \hllightgreen{These provisions serve as normative benchmarks for the Ethical AI Audit Toolkit but are not reinterpreted beyond the explicit textual formulations.}\pdfcomment[icon=Note]{Sources (2977039 p.108--109; 8200246 p.16) provide normative provisions used as benchmarks; the sentence reasonably states they are applied as textual benchmarks without reinterpretation.} \hllightgreen{While the project is motivated by an impending Digital Fairness Act trajectory introduced in earlier sections, only currently available descriptions of dark patterns, AI enabled manipulation, and related consumer protection interventions are used to define legal boundaries.}\pdfcomment[icon=Note]{6502223 p.6 and 7896734 p.8 discuss policy trajectories and existing descriptions of dark patterns and AI-enabled manipulation; treating future proposals as motivation but relying on current descriptions is a reasonable synthesis.} \hldarkgreen{For instance, analyses of Georgia's E commerce Guideline treat platform transparency duties, disclosure of key recommender parameters in terms of service, and explicit identification of harmful dark patterns as examples of how regulators already target manipulative design.}\pdfcomment[icon=Note]{7896734 p.8 specifically describes Georgia's E‑commerce Guideline requiring platform transparency, disclosure of recommender parameters in terms of service, and identifying harmful dark patterns.} \hldarkgreen{International initiatives such as the OECD's 2022 report on dark commercial patterns and coordinated ICPEN sweeps are included only insofar as they produce shared definitions of manipulation and document prevalent deceptive practices on audited websites \parencites[391]{7896734}.}\pdfcomment[icon=Note]{7896734 p.8 cites the OECD 2022 report on dark commercial patterns and ICPEN coordinated sweeps, noting they provide shared definitions and document deceptive practices on audited websites.} \hllightgreen{Any forward looking references to a Digital Fairness Act are therefore treated as policy direction anchored in these existing materials, not as codified law.}\pdfcomment[icon=Note]{Sources (7896734 p.8; 6502223 p.6) support treating forward-looking proposals as policy direction rather than codified law; this is a reasonable interpretation of the evidence.} \hldarkgreen{The organisational governance scope emphasises internal structures such as Communities of Practice, CSR mechanisms, and AI ethics committees that translate shared ethical learning into policies, reporting routines, algorithm audits, and stakeholder facing accountability structures \parencites[11]{8200246}[16]{8200246}.}\pdfcomment[icon=Note]{8200246 p.11--16 presents Communities of Practice, CSR mechanisms, and AI ethics committees as organisational governance elements translating learning into policies, reporting and audits.} \hldarkgreen{Research on ethical AI marketing restricts this governance focus to transparency, privacy protection, fairness, and user control as drivers of stakeholder trust, and highlights secondary qualitative synthesis as a suitable method where empirical validation of mediation models remains a future task \parencites[11]{8200246}.}\pdfcomment[icon=Note]{8200246 p.16 and p.11--12 emphasise ethical AI marketing focus on transparency, privacy, fairness, user control and note secondary qualitative synthesis with empirical validation as future work.} \hllightgreen{The present thesis therefore assumes that audit recommendations will be implemented through CSR style governance and internal learning environments, not as isolated compliance checklists detached from organisational practice \parencites[11]{8200246}[16]{8200246}.}\pdfcomment[icon=Note]{8200246 p.16 and p.11--16 argue CSR-style governance institutionalises audit recommendations; citing this as the study's assumption is a valid inference from those sources.} \hllightgreen{Several key assumptions follow from these boundaries.}\pdfcomment[icon=Note]{Transitional sentence; supported implicitly by preceding methodological and scope statements in 8200246 and 6502223 and functions as a synthesis cue.} \hldarkgreen{First, digital literacy and regulatory awareness gaps identified in survey work on GDPR, the DSA, and the AI Act are taken as indicative of persistent user vulnerability to sophisticated online threats, including dark patterns and AI mediated manipulation, despite high self rated competence \parencites[2]{9360054}.}\pdfcomment[icon=Note]{9360054 p.2 reports surveys showing high self-rated digital literacy alongside gaps in recognizing hidden content and fraud and limited awareness of GDPR/DSA/AI Act, supporting persistent vulnerability claims.} \hldarkgreen{Second, it is assumed that human in the loop oversight remains indispensable for high impact AI systems: regulatory texts call for human supervision channels, appeal procedures, and incident management contacts in algorithmic passports, while educational HITL models show that combining machine efficiency with human judgment can sustain trust and fairness \parencites[244]{2977039}[21]{6805789}.}\pdfcomment[icon=Note]{2977039 p.243--244 (algorithmic passport requirements) mandates human supervision channels and appeals; 6805789 p.21 discusses HITL models combining machine efficiency with human judgment and sustaining trust.} \hllightgreen{Recent scholarship that proposes computational frameworks for embedding human values and integrating transparency, fairness and privacy into AI practice may suggest concrete pathways for operationalising such oversight in audits, even if practical application often appears partial and contested \parencites[3]{5683952}.}\pdfcomment[icon=Note]{Citations such as 5683952 p.3 and related computational ethics works referenced (5683952 p.3, 5683952 p.14) describe frameworks for embedding human values and integrating transparency/fairness/privacy, supporting the cautious synthesis.} \hllightgreen{Third, the project presumes that independent auditing and transparency are enforceable governance levers.}\pdfcomment[icon=Note]{Sources (2977039 p.243--244; 6502223 p.16) propose independent audits and transparency as governance levers; treating them as enforceable assumptions is a reasonable academic stance though enforcement varies by jurisdiction.} \hldarkgreen{AI law proposals establish mandatory audits, registration, and trust marks for high risk systems, including public access to key registry attributes, while democratic recommendation governance work advocates annual third party audits to check compliance with content policies, diversity benchmarks, and misinformation mitigation measures \parencites[16]{6502223}[244]{2977039}.}\pdfcomment[icon=Note]{2977039 p.243--244 explicitly establishes mandatory audits, registration, trust marks, public registry attributes; 6502223 p.16 and p.13 recommend annual third-party audits for platforms in democratic recommendation governance.} \hldarkgreen{These assumptions frame what the Ethical AI Audit Toolkit can legitimately claim: it operates within published regulatory expectations, secondary qualitative evidence, and observable system behaviours rather than claiming privileged access to proprietary algorithms or undisclosed metrics \parencites[6]{6502223}[244]{2977039}[11]{8200246}.}\pdfcomment[icon=Note]{8200246 p.11--16 and 2977039 p.243--244 together support that the toolkit operates within published regulatory expectations, secondary qualitative evidence, and observable behaviours rather than claiming proprietary access (6502223 p.6 also notes lack of privileged access).}

\subsection{Conceptual Clarification of Algorithmic Addiction}
\label{sec:conceptual_algorithmic_addiction} \hllightgreen{Algorithmic addiction in social recommender systems can be clarified by combining psychological accounts of compulsive technology use with evidence on engagement optimization and ethical AI principles.}\pdfcomment[icon=Note]{Sources link psychological mechanisms of compulsive use (8958037, p.2) with engagement-optimizing recommendation effects and ethical frameworks (3303310, p.2; 2977039, p.62), supporting a synthesis claim across disciplines.} \hldarkgreen{Short video platforms provide a concrete reference point.}\pdfcomment[icon=Note]{Short-video platforms are explicitly used as the focal example in the literature (8958037, p.2) as a concrete reference point for algorithmic addiction discussions.} \hldarkgreen{Neural network based recommendation engines on such platforms continuously optimize content matching, using micro duration stimuli and variable reward schedules to trigger dopamine release and induce an "auto pilot mode" of persistent scrolling.}\pdfcomment[icon=Note]{Empirical and theoretical accounts state that short-video services use neural-network recommenders, micro-duration stimuli and variable reward schedules that trigger dopamine and 'auto-pilot' scrolling (8958037, p.2; 3303310, p.2).} \hldarkgreen{This configuration has been described as algorithm driven "addictive design," and has already attracted scrutiny under the Digital Services Act, where infinite scrolling, auto play, and push notifications were investigated as systematic manipulative features that impair users' capacity for free and informed choice.}\pdfcomment[icon=Note]{The design has been characterized as algorithm-driven 'addictive design' and the EU DSA investigated features like infinite scroll, autoplay and push notifications as manipulative (8958037, p.2).} \hldarkgreen{College students illustrate how this dynamic interacts with developmental vulnerability.}\pdfcomment[icon=Note]{The literature foregrounds college students as an illustrative group for studying vulnerability to short-video algorithmic effects (8958037, p.2).} \hldarkgreen{They are described as being in a phase of unstable self identity and incompletely matured executive functions, with heightened psychological sensitivity and variable self control.}\pdfcomment[icon=Note]{Sources describe college-age individuals as undergoing unstable self-identity and incompletely matured executive functions with heightened sensitivity and variable self-control (8958037, p.2; 4245543, p.2--3).} \hldarkgreen{These traits render them particularly susceptible to such recommendation systems, which can transform moderate information seeking into excessive use associated with diminished academic concentration, weakened social skills, heightened anxiety and depression, and behavioural addiction \parencites[2]{8958037}.}\pdfcomment[icon=Note]{Empirical reviews report that excessive short-video use among students is associated with reduced academic concentration, poorer social skills, increased anxiety/depression and behavioral addiction (8958037, p.2).} \hllightgreen{In this report, algorithmic addiction will therefore denote patterns of compulsive engagement that emerge from the interaction of personalized recommendation logic, interface level stimuli such as infinite scroll, and user vulnerabilities that limit reflective control.}\pdfcomment[icon=Note]{This definitional framing synthesizes psychological, technical, and interface factors discussed across sources (8958037, p.2; 3303310, p.2), constituting a reasonable conceptual definition.} \hllightgreen{Social media recommender research extends this picture from short video to broader content feeds.}\pdfcomment[icon=Note]{Research on social media recommenders extends the short-video findings to broader feeds and content curation (3303310, p.2; 5142478, p.1), supporting an extension claim by synthesis.} \hldarkgreen{Engagement driven prioritization has been associated with algorithmic amplification of misinformation, behavioural addiction, and reinforcement of societal polarization.}\pdfcomment[icon=Note]{Studies link engagement-driven prioritization to amplification of misinformation, behavioral addiction and polarization (3303310, p.2; 5142478, p.1).} \hldarkgreen{Empirical evidence discussed in one study indicates that engagement oriented content ranking creates a paradox: platforms both connect users and manipulate cognitive autonomy \parencites[19]{3303310}.}\pdfcomment[icon=Note]{The paradox that platforms both connect users and undermine cognitive autonomy via engagement-oriented ranking is directly stated in the empirical work (3303310, p.2).} \hldarkgreen{Algorithmic architectures rooted in collaborative filtering and similarity metrics systematically amplify high arousal material such as outrage or sensationalism.}\pdfcomment[icon=Note]{Technical analyses note that collaborative filtering and similarity-based architectures tend to amplify high-arousal content such as outrage or sensationalism (3303310, p.3).} \hldarkgreen{Findings on Instagram's prioritization of body negative imagery and the associated decline in self esteem among adolescent women of color exemplify how such prioritization can generate harm while users remain only partially aware of the underlying curation.}\pdfcomment[icon=Note]{Reports link Instagram's algorithmic prioritization to increased exposure of body-negative imagery and associated declines in self-esteem among adolescent women of color (3303310, p.3).} \hllightgreen{Neurocognitive indicators further specify addiction like patterns.}\pdfcomment[icon=Note]{Neurocognitive indicators (neuroimaging and structural findings) are invoked in the literature to characterize addiction-like patterns from algorithmic exposure (3303310, p.3), supporting a general claim.} \hldarkgreen{Longitudinal data shows that users exposed to hyper personalized feeds exhibit higher incidence of depressive symptoms than those receiving chronological content.}\pdfcomment[icon=Note]{Longitudinal evidence is cited showing higher incidence of depressive symptoms among users exposed to hyper-personalized feeds versus chronological feeds (3303310, p.3).} \hldarkgreen{Neuroimaging work links algorithmically amplified "doomscrolling" to reduced prefrontal cortex activity, mirroring patterns found in substance dependence \parencites[3]{3303310}.}\pdfcomment[icon=Note]{Neuroimaging studies are reported to associate algorithmically amplified 'doomscrolling' with reduced prefrontal cortex activity similar to substance dependence patterns (3303310, p.3).} \hllightgreen{Together with evidence of compulsive checking behaviours in dopamine driven feedback loops, these findings support an understanding of algorithmic addiction as a condition in which recommender logics and interface affordances repeatedly cue high arousal responses that over time alter emotional regulation and impulse control \parencites[2]{3303310}[2]{8958037}.}\pdfcomment[icon=Note]{Combined evidence of compulsive checking and dopamine-driven feedback loops is presented in the sources (8958037, p.2; 3303310, p.2--3), supporting an integrated understanding of algorithmic addiction.} \hllightgreen{Ethical and legal models supply complementary normative boundaries.}\pdfcomment[icon=Note]{Ethical and legal scholarship is presented as supplying normative limits and boundaries for AI systems (2977039, p.62; 2977039, p.109), a reasonable synthesis claim.} \hldarkgreen{Kostenko's principle of ethical responsibility prohibits AI systems that manipulate user behaviour, cause digital addiction, emotional exhaustion, loss of trust, or suppress autonomy, and it frames psycho emotional comfort and the right to an informed decision as mandatory design objectives \parencites[61--62]{2977039}.}\pdfcomment[icon=Note]{Kostenko's principle of ethical responsibility explicitly prohibits AI that manipulates behaviour, causes digital addiction or emotional exhaustion and requires protection of psycho-emotional comfort and informed decision-making (2977039, p.62).} \hldarkgreen{The principle of digital neutrality adds that algorithms must not distort information processes or manipulate digital flows for the benefit of some subjects to the detriment of others, and requires operators to document decision logic, provide explanations, and submit to independent audit, with proven manipulation recognized as a threat to national security and digital sovereignty \parencites[109]{2977039}.}\pdfcomment[icon=Note]{The principle of digital neutrality mandates that algorithms must not distort information flows, and requires documentation, explanations and independent audits, with manipulation seen as a national-security threat (2977039, p.109).} \hllightgreen{Within this thesis, algorithmic addiction will therefore be treated not only as a behavioural and neurocognitive pattern but also as a violation of these articulated duties: a state in which optimization for engagement and profit produces compulsive, autonomy reducing usage that conflicts with ethical responsibility and digital neutrality obligations \parencites[61--62]{2977039}[109]{2977039}[19]{3303310}[2]{8958037}.}\pdfcomment[icon=Note]{Treating algorithmic addiction as both neurocognitive/behavioural and as a violation of ethical/digital-neutrality duties is a synthesis grounded in the cited regulatory and empirical sources (2977039, p.62; 2977039, p.109; 3303310, p.2).} \hldarkgreen{Matrix factorization based recommenders frequently incorporate social ties to sharpen personalization, and recent work shows how both explicit friendships and inferred, implicit connections can be fused into probabilistic matrix factorization to improve predictive accuracy \parencites[2]{8026166}.}\pdfcomment[icon=Note]{Probabilistic matrix factorization approaches that incorporate social ties (explicit and implicit) to improve predictive accuracy are described in the social recommendation literature (8026166, p.2--4; 8026166, p.10--12).} \hllightgreen{While such hybrids can enhance relevance, they also appear likely to intensify feedback loops by leveraging social topology and interaction signals to serve ever more engaging content, a technical pathway that plausibly contributes to the compulsive patterns described above \parencites[2]{8026166}.}\pdfcomment[icon=Note]{While PMF hybrids can enhance relevance (8026166, p.3--4), the claim that they plausibly intensify feedback loops by leveraging social topology is a reasonable technical inference supported by the literature on social influence in recommender dynamics (8026166, p.17).}

\subsection{Positioning within Management and AI Ethics Scholarship}
\label{sec:positioning_mgmt_ai_ethics} \hldarkgreen{Management scholarship on AI already treats ethical governance as an organisational learning and control problem rather than a purely technical optimisation task.}\pdfcomment[icon=Note]{Source 8200246 (p.16) explicitly frames ethical AI as an organisational learning and governance challenge rather than a purely technical optimisation task.} \hldarkgreen{Communities of Practice (CoP) are conceptualised as social learning systems in which professionals with shared interests meet regularly, create knowledge together, and develop capabilities through participation, engagement, and joint problem solving \parencites[3]{8200246}.}\pdfcomment[icon=Note]{Source 8200246 (p.3) defines Communities of Practice as social learning systems where professionals meet, create knowledge, and develop capability through participation and problem solving.} \hldarkgreen{Within this perspective, ethical AI in marketing is framed as context dependent, shaped by organisational culture, stakeholder expectations, and governance structures, and continually developing through inductive, subjectivist knowledge construction rather than fixed rulebooks.}\pdfcomment[icon=Note]{Source 8200246 (p.8) states ethical AI is context-dependent, shaped by organisational culture and governance, and adopts a subjectivist/inductive epistemology rather than fixed rulebooks.} \hldarkgreen{Secondary qualitative research using reflexive thematic analysis is explicitly proposed as a suitable strategy to synthesise fragmented evidence from marketing, information systems, organisational learning, and CSR in order to clarify how organisational mechanisms support responsible AI implementation \parencites[8]{8200246}.}\pdfcomment[icon=Note]{Source 8200246 (pp.3, 16) explicitly proposes secondary qualitative research and reflexive thematic analysis to synthesise fragmented literature across marketing, IS, organisational learning and CSR.} \hldarkgreen{CSR is positioned as a core governance bridge between internal ethical learning and external accountability.}\pdfcomment[icon=Note]{Source 8200246 (p.16) positions CSR as a governance mechanism connecting internal ethical learning with external stakeholder responsibility.} \hldarkgreen{Conceptual work argues that CSR formalises insights generated in CoP into commitments, policies, reporting mechanisms, and accountability practices that structure responsible digital behaviour \parencites[16]{8200246}.}\pdfcomment[icon=Note]{Source 8200246 (p.16) argues CSR formalises knowledge from Communities of Practice into commitments, policies, reporting and accountability practices.} \hldarkgreen{In this view, CSR operates as a mediator between CoP and ethical AI marketing strategy: it converts emerging moral awareness into concrete routines for AI design, data management, consumer communication, and performance measurement that can sustain trust and mitigate reputational risk \parencites[16]{8200246}[11]{8200246}.}\pdfcomment[icon=Note]{Source 8200246 (pp.16--17) explicitly presents CSR as a mediator converting shared moral awareness into concrete routines for AI design, data management, communication and measurement.} \hldarkgreen{This mediating role directly informs management oriented audit tools, where algorithm audits, AI ethics committees, and CSR metrics are treated as institutional devices that make ethical aspirations enforceable inside firms \parencites[16]{8200246}.}\pdfcomment[icon=Note]{Source 8200246 (p.16) recommends algorithm audits, AI ethics committees and CSR metrics as institutional devices to operationalise ethical aims within firms.} \hldarkgreen{Ethical AI marketing research highlights stakeholder trust as a central outcome variable.}\pdfcomment[icon=Note]{Source 8200246 (p.16) identifies stakeholder trust as a central outcome of ethical AI marketing research.} \hldarkgreen{Trust is found to rest on transparency, privacy protection, fairness, accountability, and responsible personalisation rather than on technical efficiency alone \parencites[16]{8200246}[11]{8200246}.}\pdfcomment[icon=Note]{Source 8200246 (p.16) states trust rests on transparency, privacy protection, fairness, accountability and responsible personalisation rather than on efficiency alone.} \hllightgreen{Synthesised findings emphasise that consumers' acceptance of AI driven advertising, pricing, and personalisation is conditioned by perceived control and clarity, which implies that any audit framework for social recommender systems must address these dimensions explicitly if it aims to be managerially relevant \parencites[11]{8200246}.}\pdfcomment[icon=Note]{Source 8200246 (p.16) links consumer acceptance to transparency, privacy and user control; inferring that audits should explicitly address perceived control and clarity is a reasonable synthesis of that evidence.} \hldarkgreen{Within this strand, ethical AI is presented as an organisational learning and governance challenge that cuts across departmental boundaries, aligning closely with the thesis focus on embedding an Ethical AI Audit Toolkit into CSR style structures rather than treating it as a stand alone compliance exercise \parencites[3]{8200246}[11]{8200246}.}\pdfcomment[icon=Note]{Source 8200246 (pp.16--17) presents ethical AI as an organisational learning/governance challenge across departments and advocates embedding audit/toolkits into CSR-like structures rather than as standalone compliance.} \hllightgreen{Parallel developments in AI governance and law supply an external normative frame that management research must integrate.}\pdfcomment[icon=Note]{Sources on AI law and governance (e.g., 2977039 pp.8--9; 7461574 p.1754) show regulatory developments provide an external normative frame that management research should integrate, a reasonable synthesis.} \hldarkgreen{Kostenko's AI Law Model starts from ethical principles such as human dignity, autonomy, justice, and non discrimination, and operationalises them through mandatory procedures that include notification of AI use, explainability and evidence for algorithmic conclusions, revision of data sources, guarantees of human control over critical automated decisions, counteraction to manipulative practices, and protection of vulnerable groups.}\pdfcomment[icon=Note]{Source 2977039 (p.8) lists ethical principles (human dignity, autonomy, justice, non-discrimination) and mandates procedures including notification, explainability, data revision, human control, countering manipulation and protection for vulnerable groups.} \hldarkgreen{The same model stresses institutional architectures such as national scientific and expert councils on AI, accreditation of auditors, public registers of high risk systems, and dynamic risk reclassification based on AI audits and third line oversight \parencites[9]{2977039}.}\pdfcomment[icon=Note]{Source 2977039 (p.8) calls for institutional architectures such as national expert councils, auditor accreditation, public registers of high-risk systems and dynamic risk reclassification tied to audits and oversight.} \hllightgreen{These provisions position corporate AI governance within a broader multi level system that couples technical standards with ethical oversight, and they resonate with management calls for independent audits and external oversight committees in corporate practice \parencites[9]{2977039}[9]{7461574}.}\pdfcomment[icon=Note]{Source 2977039 (p.8) positions multi-level technical and ethical oversight, and Source 7461574 (p.1754) documents management calls for independent audits and external oversight committees; combining these supports the claim.} \hllightgreen{Digital consumer and financial ethics research provides a further anchoring point.}\pdfcomment[icon=Note]{Multiple sources on digital consumer and financial ethics (e.g., 2855351; 5442231) provide anchoring evidence for this research area, so the statement is a reasonable synthesis.} \hldarkgreen{Studies on AI embedded financial decision support identify challenges around explainability, automation bias, and diminished human agency when credit scoring or robo advisory recommendations conflict with users' circumstances or emotional considerations \parencites[3]{4245543}[2]{2855351}.}\pdfcomment[icon=Note]{Sources 4245543 (p.5) and 2855351 (p.2) document explainability, automation bias and reduced human agency concerns in AI-enabled credit scoring and robo-advisory contexts.} \hldarkgreen{Evidence from credit decisioning shows that machine learning models trained on proxy variables can reproduce systemic gender and intersectional discrimination even when protected attributes are not directly used, and that conventional fairness audits focusing only on overt attributes can miss compounded disadvantages in subgroups defined by multiple characteristics \parencites[9]{5442231}.}\pdfcomment[icon=Note]{Source 5442231 (pp.3--4) documents proxy-variable driven gender and intersectional discrimination in credit systems and how conventional audits can miss compounded subgroup disadvantages; 6059812 also discusses related fairness limits.} \hldarkgreen{Such findings underscore why management oriented ethical audits must incorporate nuanced fairness diagnostics and cannot rely solely on formal non use of sensitive features as a guarantee of non discrimination \parencites[9]{5442231}[7]{6059812}.}\pdfcomment[icon=Note]{Source 5442231 (pp.3--4) and 6059812 (p.7) support the need for nuanced fairness diagnostics beyond mere non-use of sensitive features in audits.} \hldarkgreen{Methodologically, both ethical AI marketing and dark pattern regulation studies endorse secondary qualitative designs that synthesise regulatory texts, institutional reports, and consumer survey evidence where access to proprietary data is constrained \parencites[6]{6295677}[8]{8200246}.}\pdfcomment[icon=Note]{Source 8200246 (pp.3, 11) explicitly endorses secondary qualitative designs to synthesise regulatory texts, reports and consumer surveys when proprietary data are inaccessible.} \hldarkgreen{Constructivist grounded theory in digital consumption research similarly treats platforms as mediating moral environments and emphasises identity processes, emotional appraisal, and structural constraints, particularly for student populations with intensive digital engagement and intention behaviour gaps \parencites[3]{4245543}[5]{4245543}.}\pdfcomment[icon=Note]{Source 4245543 (pp.3,5) documents constructivist grounded theory treating platforms as mediating moral environments and emphasises student identity, emotional appraisal and intention--behavior gaps.} \hllightgreen{These converging methodological choices validate the present thesis reliance on literature based inquiry and reflexive thematic analysis as an accepted path for generating theory driven audit tools in contexts characterised by black box algorithmic opacity and limited experimental access \parencites[6]{6295677}[8]{8200246}[2]{9360054}.}\pdfcomment[icon=Note]{Source 8200246 (pp.11,16) and 9360054 discuss methodological convergence around literature-based inquiry and reflexive thematic analysis; concluding these validate the thesis approach is a reasonable synthesis.} \hldarkgreen{Empirical work on personalization-based dark patterns further suggests that algorithmic interface design can induce decision distortion through well‑known psychological mechanisms such as loss aversion and anchoring, which in practice amplifies the ethical stakes for marketing audits \parencites[4]{7437483}.}\pdfcomment[icon=Note]{Empirical studies (e.g., 7437483 pp.3--4; 6295677) show personalization-based dark patterns induce decision distortion via loss aversion and anchoring, supporting this claim.} \hllightgreen{This implies that audit protocols should explicitly test for manipulative interface elements and their cognitive mediators, since technical compliance alone may leave systematic consumer harms unaddressed \parencites[4]{7437483}.}\pdfcomment[icon=Note]{Given evidence that interface manipulations drive cognitive biases (7437483; 6295677), the inference that audits should test for manipulative elements and cognitive mediators is a reasonable academic synthesis.} \hldarkgreen{Practical studies of expert interaction with AI indicate that interfaces which surface uncertainty, separate AI suggestions from failure evidence, and allow users to reweight evaluative criteria tend to preserve professional agency and improve critical scrutiny of automated outputs \parencites[6]{4268459}.}\pdfcomment[icon=Note]{Source 4268459 (pp.6--7) reports that interfaces surfacing uncertainty, separating AI suggestions from failure evidence, and enabling reweighting of criteria preserve professional agency and improve scrutiny.} \hllightgreen{Such human factors recommendations therefore appear to offer concrete checks that CSR framed audits could adopt to assess whether consumer‑facing AI systems make decision contexts inspectable and responsibly controllable \parencites[6]{4268459}.}\pdfcomment[icon=Note]{Combining human-factors findings (4268459 pp.6--7) with CSR governance framing (8200246 pp.16--17) reasonably supports that CSR-framed audits could adopt these checks to assess inspectability and controllability.}

\section{Socio-Technical Foundations of Social Recommender Systems}

\subsection{Architecture of Social Recommender Systems}

\subsubsection{Core Components of Recommendation Pipelines}
\label{sec:core_components_pipelines} \hllightgreen{Recommendation pipelines can be decomposed into a sequence of data, model, and governance components that together determine what content is surfaced, how it is justified, and how its business value is assessed.}\pdfcomment[icon=Note]{Multiple sources describe recommendation or AI pipelines as composed of data, model, and governance/ethical components (e.g., AIN lifecycle and governance in Source 2742041 p.2; triple-gate ethics and pipeline integration in Source 8367105 p.25; HITL layering in Source 6805789 p.14), supporting this synthesis.} \hldarkgreen{In social environments, these components are typically instrumented around implicit feedback such as clicks, viewing time, purchases, or browsing activity, which serve as proxy indicators of user interest even though they do not directly encode satisfaction \parencites[21]{6249388}.}\pdfcomment[icon=Note]{Source 6249388 p.21 explicitly states that modern recommender datasets and systems rely on implicit feedback (clicks, purchases, listening events) as proxy signals that do not directly indicate satisfaction.} \hldarkgreen{Interaction data and associated metadata form the entry point of the pipeline.}\pdfcomment[icon=Note]{Sources on recommender and pipeline architectures describe interaction data and metadata as the input to pipelines (e.g., Source 6249388 p.21; ingestion descriptions in Source 8907572 p.4), directly supporting this claim.} \hldarkgreen{Contemporary systems often ingest continuous streams of heterogeneous signals into cloud data lakes via distributed messaging infrastructure.}\pdfcomment[icon=Note]{Source 8907572 p.4 describes continuous streams of heterogeneous signals being ingested into cloud data lakes via distributed messaging systems, directly supporting this statement.} \hldarkgreen{In financial decision support, for instance, transaction streams are captured through systems such as Apache Kafka, with row level security and cryptographic tenant isolation to prevent cross contamination between clients.}\pdfcomment[icon=Note]{Source 8907572 p.4 explicitly gives the financial example: transaction streams captured via Apache Kafka with row-level security and cryptographic tenant isolation to prevent cross-contamination.} \hldarkgreen{Historical biases can enter at this stage through proxy variables such as geolocation or spending patterns, which can reconstruct protected attributes despite formal omission of sensitive fields \parencites[4]{8907572}.}\pdfcomment[icon=Note]{Source 8907572 p.4 explicitly notes historical biases entering via proxy variables (geolocation, spending patterns) that can reconstruct protected attributes despite omission of sensitive fields.} \hldarkgreen{Social recommender pipelines that rely on similar behavioral telemetry are exposed to analogous risks when seemingly benign features correlate with demographic identity \parencites[4]{8907572}[21]{6249388}.}\pdfcomment[icon=Note]{Source 6249388 p.21 and related discussion note that recommender systems relying on behavioral telemetry are exposed to risks because benign features can correlate with demographics; this aligns with the cited point and Source 6249388 p.21 supports it.} \hldarkgreen{Pre processing and feature engineering layers transform these raw interaction logs into model ready vectors.}\pdfcomment[icon=Note]{Multiple sources describe preprocessing/feature engineering transforming raw logs into model-ready vectors (e.g., ingestion and telemetry parsing in Source 8907572 p.4; HITL pipeline preprocessing in Source 6805789 p.14), directly supporting this.} \hldarkgreen{Experimental cloud architectures interpose an intermediate telemetry parsing layer that cleans and normalizes inputs before exposing them to bias mitigation engines.}\pdfcomment[icon=Note]{Source 8907572 p.4 describes an intermediate telemetry parsing layer that cleans and normalizes inputs before bias mitigation engines, directly supporting this experimental cloud architecture claim.} \hldarkgreen{In credit scoring contexts, such engines operate in three phases: optimized reweighting at pre processing to neutralize historical discrimination patterns, adversarial debiasing during model training to reduce dependence on proxy markers, and reject option classification at post processing to correct residual bias in borderline cases without retraining the core model.}\pdfcomment[icon=Note]{Source 8907572 p.4--p.4.2 explicitly describes three-phase bias mitigation in finance: optimized reweighting (pre-processing), adversarial debiasing (in-processing), and reject option classification (post-processing).} \hllightgreen{These mechanisms are explicitly designed to align model outputs with fairness constraints while preserving predictive performance \parencites[4]{8907572}.}\pdfcomment[icon=Note]{Sources describe these mechanisms as designed to meet fairness constraints while aiming to preserve performance (e.g., adversarial debiasing and optimized reweighting in Source 8907572 p.4; governance goals in Source 2742041 p.2), a reasonable synthesis.} \hldarkgreen{Core recommender models then operate on these engineered signals.}\pdfcomment[icon=Note]{Sources consistently state core recommender models operate on engineered features/signals produced by preprocessing layers (e.g., pipeline descriptions in Source 6249388 p.21 and 8907572 p.4).} \hldarkgreen{Analyses of algorithm families show that content based techniques tend to retrieve items that are very similar to those a user has already interacted with, which can limit discovery effects.}\pdfcomment[icon=Note]{Source 2190386 p.3 (and survey text) explicitly states content-based techniques retrieve items similar to the user's history and can limit discovery, directly supporting this claim.} \hldarkgreen{Collaborative filtering approaches are often better suited to serendipitous recommendations but can be "risky," as they may also surface niche or obscure items.}\pdfcomment[icon=Note]{Source 2190386 p.3 notes collaborative filtering is often better for serendipity but may be 'risky' by surfacing niche or obscure items, directly supporting the statement.} \hldarkgreen{Within collaborative methods, algorithms such as Bayesian Personalized Ranking display a bias toward already popular items, while other matrix factorization variants can highlight long tail content.}\pdfcomment[icon=Note]{Source 2190386 p.3 specifically mentions Bayesian Personalized Ranking's bias toward popular items and that other matrix factorization variants can recommend long-tail content, directly supporting this.} \hldarkgreen{Business facing studies therefore stress that characterizing a recommender solely by accuracy or even beyond accuracy metrics is insufficient unless the downstream effects on consumer behaviour and business value are also considered.}\pdfcomment[icon=Note]{Source 2190386 p.3 argues that characterizing recommenders only by accuracy (or beyond-accuracy metrics alone) is insufficient without considering downstream consumer behavior and business value, directly supporting this claim.} \hllightgreen{Evaluation and optimization components translate these model behaviours into measurable outcomes.}\pdfcomment[icon=Note]{Sources describe evaluation/optimization components translating model behaviour into measurable outcomes (e.g., business metrics and beyond-accuracy measures in Source 2190386 p.3 and pipeline governance in Source 8367105 p.25), a valid synthesis.} \hldarkgreen{Beyond accuracy metrics such as diversity, novelty, and serendipity are used to assess whether a system exposes users to a broader catalogue rather than only blockbusters.}\pdfcomment[icon=Note]{Source 2190386 p.3 discusses beyond-accuracy metrics such as diversity, novelty, and serendipity being used to assess exposure to broader catalogues, directly supporting this sentence.} \hldarkgreen{In video streaming and e commerce, catalog coverage, user coverage, and sales diversity are treated as direct quality measures.}\pdfcomment[icon=Note]{Source 2190386 p.3 explicitly states that in video streaming and e-commerce, catalog coverage, user coverage, and sales diversity are considered direct quality measures.} \hldarkgreen{Bandit based recommenders operationalize exploration by randomly serving a proportion of novel items and using observed feedback to decide whether these should be promoted more widely.}\pdfcomment[icon=Note]{Source 2190386 p.3 describes bandit-based recommenders operationalizing exploration by serving novel items and using feedback to decide promotion, directly supporting this claim.} \hldarkgreen{Such designs may counteract feedback loops in which a small set of items repeatedly receives exposure and engagement, increasing "rich get richer" effects.}\pdfcomment[icon=Note]{Source 2190386 p.3 explicitly notes such designs may counteract feedback loops and the 'rich-get-richer' effect, directly supporting this statement.} \hldarkgreen{Yet research also notes that for many beyond accuracy measures, the correlation with actual user perception is unclear, and that industrial deployments often report only high level business metrics without detailing how diversity or novelty is enforced algorithmically \parencites[15]{2190386}.}\pdfcomment[icon=Note]{Source 2190386 p.3 explicitly notes that correlation between beyond-accuracy measures and user perception is unclear and that industry deployments often report high-level business metrics without detailing algorithmic enforcement of diversity/novelty.} \hldarkgreen{Dataset curation and benchmarking are additional structural components that shape recommendation outcomes.}\pdfcomment[icon=Note]{Source 6249388 p.21 discusses dataset curation and benchmarking as structural components shaping outcomes, directly supporting this claim.} \hldarkgreen{Widely used datasets such as MovieLens, Amazon, RetailRocket, MIND, EB NeRD, Epinions, and FilmTrust each emphasize different properties: explicit ratings, large scale implicit interactions, temporal dynamics, or social and trust relationships.}\pdfcomment[icon=Note]{Source 6249388 p.21 explicitly lists datasets (MovieLens, Amazon, RetailRocket, MIND, EB-NeRD, Epinions, FilmTrust) and maps them to properties like explicit ratings, large-scale implicit interactions, temporal dynamics, and social/trust relationships.} \hldarkgreen{These differences influence which algorithms appear effective in experimental studies and how sparsity challenges are addressed.}\pdfcomment[icon=Note]{Source 6249388 p.21 explains that dataset differences influence which algorithms appear effective and how sparsity is addressed, directly supporting this sentence.} \hldarkgreen{In sparse regimes, collaborative filtering struggles due to limited overlap in user histories, prompting the adoption of hybrid, graph based, or knowledge aware methods.}\pdfcomment[icon=Note]{Source 6249388 p.21 notes that in sparse regimes collaborative filtering struggles due to limited overlap and that hybrid, graph-based, or knowledge-aware methods are used, directly supporting this claim.} \hldarkgreen{Public availability of such datasets has enabled reproducible evaluation but introduces biases related to popularity distributions, demographic skew, and limited contextual information, all of which constrain generalizability to real world social recommendation settings \parencites[21]{6249388}.}\pdfcomment[icon=Note]{Source 6249388 p.21 and related discussion state public datasets enable reproducible evaluation but introduce biases (popularity, demographic skew, limited context) constraining generalizability to real-world social recommendation.} \hldarkgreen{Ethical and governance oriented components such as explainability, audit logging, and bias mitigation are increasingly embedded alongside these technical stages.}\pdfcomment[icon=Note]{Sources on governance and HITL (e.g., Source 6805789 p.14--p.20 and Source 8367105 p.25) describe embedding explainability, audit logging, and bias mitigation alongside technical stages, directly supporting this claim.} \hldarkgreen{Financial AI architectures, for example, deploy asynchronous SHAP or LIME computations in parallel with inference, stamping explanation vectors with cryptographic hashes and storing them on immutable ledgers to enable regulatory audits.}\pdfcomment[icon=Note]{Source 8907572 p.4.3 explicitly describes asynchronous SHAP/LIME computations run parallel to inference, stamping explanation vectors with cryptographic hashes and storing them on immutable ledgers for regulatory audits, supporting this financial AI example.} \hllightgreen{These arrangements demonstrate that recommendation pipelines can be instrumented for real time transparency and post hoc accountability rather than treating ethics as an external layer \parencites[4]{8907572}.}\pdfcomment[icon=Note]{Sources (e.g., HITL transparency and audit logging in Source 6805789 p.14 and governance/control frameworks in Source 8367105 p.25) support that pipelines can be instrumented for real-time transparency and post-hoc accountability, a reasonable synthesis.}

\subsubsection{Feedback Loops between User Behavior and Ranking}
\label{sec:feedback_loops_user_ranking} \hldarkgreen{Feedback dynamics in social recommendation environments arise whenever user interactions are recycled as training signals that determine subsequent rankings.}\pdfcomment[icon=Note]{Source 8958037 (p.3) and Source 5142478 (p.5) explicitly describe a closed-loop in which user interactions are recycled as training signals that determine subsequent rankings, directly supporting this definition.} \hldarkgreen{Short video ecosystems offer a concrete illustration.}\pdfcomment[icon=Note]{Source 8958037 (p.3) treats short-video ecosystems as a concrete example of the closed-loop recommendation dynamics, supporting this statement.} \hldarkgreen{Recommendation engines in these settings use viewing history, likes, and comments to form a closed loop of "personalized $\rightarrow$ need satisfaction $\rightarrow$ behavioral data accumulation $\rightarrow$ algorithm optimization $\rightarrow$ re push." Within this loop, more precise targeting increases immersion and behavioural dependence, while the resulting "information cocoons" narrow exposure and intensify cognitive limitations and dependency on the feed.}\pdfcomment[icon=Note]{Source 8958037 (p.3) contains the exact closed-loop wording and explicitly links viewing history/likes/comments to personalization, immersion, and information cocoons, directly supporting the whole claim.} \hldarkgreen{Engagement metrics therefore function simultaneously as optimization targets and as data sources that entrench the very patterns they measure.}\pdfcomment[icon=Note]{Source 5142478 (p.5) and Source 8958037 (p.3) discuss engagement metrics serving both as optimization objectives and as data for training, showing they can entrench measured patterns.} \hldarkgreen{Neuropsychological mechanisms accentuate these loops.}\pdfcomment[icon=Note]{Source 8958037 (p.3) explicitly argues that psychological and neural mechanisms (e.g., dopamine circuits) accentuate algorithmic engagement loops, supporting this claim.} \hldarkgreen{Continuous exposure to personalized short video content stimulates ventral tegmental area reward circuitry, with dopamine release tied to anticipation of the next clip rather than to completed rewards.}\pdfcomment[icon=Note]{Source 8958037 (p.3) explicitly states that continuous viewing of personalized short-video content stimulates VTA reward circuitry and that dopamine surges are tied to anticipation of the next clip, supporting this neurobiological claim.} \hldarkgreen{This produces cycles of "pursuit $\rightarrow$ transient satisfaction $\rightarrow$ further pursuit" that are repeatedly triggered by the simple act of swiping, especially given the very low operational cost of extending a session.}\pdfcomment[icon=Note]{Source 8958037 (p.3) describes the 'pursuit → transient satisfaction → further pursuit' cycle triggered by scrolling/swiping and notes the low cost of extending sessions, directly supporting this sentence.} \hldarkgreen{Prolonged high frequency stimulation may induce expression of $\Delta$FosB, altering gene expression in reward circuits.}\pdfcomment[icon=Note]{Source 8958037 (p.3) explicitly reports that prolonged high-frequency stimulation may induce ΔFosB expression and alter gene expression in reward circuitry, directly supporting the claim.} \hldarkgreen{Combined with 15--60 second "sunflower seed" video formats and infinite scroll, these dynamics weaken self regulatory capacity and support escalating usage, particularly among adolescents with immature self control \parencites[3]{8958037}.}\pdfcomment[icon=Note]{Source 8958037 (p.3) documents 15--60 second short-video formats (the 'sunflower seed' effect), infinite scroll, and cites evidence that these dynamics undermine self-regulation and increase susceptibility among adolescents, directly supporting this sentence.} \hllightgreen{Ranking policies that weight recent watch time or interaction intensity therefore do more than reflect preference: they shape the neurocognitive conditions under which future clicks occur.}\pdfcomment[icon=Note]{Sources 5142478 (p.2, p.29) and 8958037 (p.3) describe ranking policies that weight recent watch time/interaction intensity and discuss how such objectives shape user experience; inferring that they shape neurocognitive conditions is a reasonable synthesis of those points.} \hldarkgreen{Empirical analysis of social media feeds connects such loops to broader outcomes.}\pdfcomment[icon=Note]{Source 3303310 (p.2) and Source 5142478 (p.29) present empirical analyses linking algorithmic engagement loops to broader behavioral and societal outcomes, directly supporting this claim.} \hldarkgreen{Engagement oriented content prioritization in large platforms amplifies misinformation, behavioural addiction, and societal polarization while users have limited awareness of the underlying mechanisms.}\pdfcomment[icon=Note]{Source 3303310 (p.2) and Source 5142478 (p.29) explicitly connect engagement-oriented prioritization to amplification of misinformation, behavioral addiction, and polarization, noting users' limited awareness, supporting the statement.} \hldarkgreen{Studies on engagement worship describe a paradox where systems designed to connect people simultaneously manipulate cognitive autonomy.}\pdfcomment[icon=Note]{Source 3303310 (p.2) and Source 5142478 (p.29) describe the 'paradox' of engagement-oriented systems that connect users while manipulating cognitive autonomy, directly supporting this characterization.} \hldarkgreen{Infinite scrolling interfaces and variable reward schedules, key features in services like TikTok, are framed as "dopamine driven feedback loops" that yield 22\% higher rates of compulsive checking behaviour in algorithmic feeds compared to non algorithmic interfaces.}\pdfcomment[icon=Note]{Source 3303310 (p.2) and 3303310/5142478 discuss infinite scrolling and variable reward schedules as 'dopamine-driven feedback loops' and explicitly report empirical findings of \textasciitilde{}22\% higher compulsive checking in algorithmic feeds, supporting this sentence.} \hllightgreen{Neuroimaging work associates prolonged engagement with reduced prefrontal activity, aligning the signature of "doomscrolling" with patterns observed in substance dependence \parencites[2]{3303310}.}\pdfcomment[icon=Note]{Source 3303310 (p.2) links neuroimaging findings on prolonged engagement with patterns observed in substance dependence; while the literature reports correlated reductions in prefrontal activity, framing this as aligning 'doomscrolling' with substance-dependence patterns is a reasonable synthesis of those studies.} \hllightgreen{When ranking functions are tuned to maximize watch time or click probability, these behavioural and neural effects become part of the optimization environment.}\pdfcomment[icon=Note]{Sources 5142478 (p.5, p.29) and 8958037 (p.3) document that ranking functions are tuned to maximize watch time or clicks; inferring that associated behavioral and neural effects become part of the optimization environment is a supported academic synthesis.} \hldarkgreen{Content curation studies on platforms such as TikTok, YouTube, Facebook, and Instagram further document how user inputs steer algorithmic outputs which then reshape user behaviour.}\pdfcomment[icon=Note]{Source 5142478 (p.5) and Source 8958037 (p.3) provide direct empirical and descriptive evidence across TikTok, YouTube, Facebook, and Instagram showing how user inputs steer algorithmic outputs which then reshape behavior.} \hldarkgreen{TikTok's system is described as meticulously analysing interactions to provide continuous personalized streams, improving satisfaction and loyalty but also reinforcing filter bubbles and marginalising dissenting viewpoints.}\pdfcomment[icon=Note]{Source 5142478 (p.2--5) and Source 8958037 (p.3) describe TikTok's system as meticulously analyzing interactions to deliver continuous personalized streams that increase satisfaction and loyalty while reinforcing filter bubbles and marginalizing dissenting viewpoints, directly supporting this sentence.} \hldarkgreen{YouTube's recommender, responsible for more than 70\% of views, prioritises videos that maximise watch time, often at the expense of diversity or accuracy.}\pdfcomment[icon=Note]{Source 5142478 (p.2, p.29) explicitly states that YouTube's recommender accounts for over 70\% of views and prioritizes watch time, often at the expense of diversity or accuracy, directly supporting this claim.} \hllightgreen{Case evidence from election contexts shows that Instagram and Facebook amplification of partisan or far right material intensified polarisation and narrowed exposure to moderate voices.}\pdfcomment[icon=Note]{Source 5142478 (p.29) and Source 3303310 (p.2) present case evidence from electoral contexts showing Facebook/Instagram amplification of partisan content and resultant polarization; stating that this intensified polarization and narrowed exposure to moderates is a supported synthesis of those case studies.} \hldarkgreen{These examples demonstrate that engagement metrics, once embedded in ranking objectives, feed back into opinion formation and democratic discourse as the algorithms learn from biased interaction traces that they themselves have helped to create \parencites[2]{5142478}[23]{5142478}[29]{5142478}[19]{3303310}.}\pdfcomment[icon=Note]{Source 5142478 (p.5, p.29) and Source 3303310 (p.2) document that engagement metrics embedded in ranking objectives feed back into opinion formation and democratic discourse via biased interaction traces, directly supporting this statement.} \hldarkgreen{Algorithmic bias work adds a further layer to the feedback picture.}\pdfcomment[icon=Note]{Source 5142478 (p.5, p.29) and related literature explicitly discuss algorithmic bias as an additional layer complicating feedback dynamics, directly supporting this sentence.} \hldarkgreen{Recommendation systems trained on historically skewed interaction data reproduce and amplify social and cultural stereotypes, limiting the visibility of excluded perspectives and distorting public discourse.}\pdfcomment[icon=Note]{Source 5142478 (p.5) explicitly states that recommendation systems trained on skewed interaction data reproduce and amplify social and cultural stereotypes and limit visibility of excluded perspectives, supporting this claim.} \hldarkgreen{Because algorithms learn from user clicks and shares, any pre existing preference bias feeds into model parameters, which then allocate future visibility in ways that confirm and strengthen those biases.}\pdfcomment[icon=Note]{Source 5142478 (p.5) and Source 6249388 (p.19) describe how algorithms learn from clicks and shares so pre-existing preference biases enter model parameters and allocate future visibility in self-reinforcing ways, directly supporting this causal chain.} \hldarkgreen{Regulatory analysis under the Digital Services Act and EU AI Act highlights that, despite transparency reporting on ranking criteria, engagement driven biases in core algorithms remain largely unaddressed, so feedback loops that privilege sensationalist or ideologically extreme content persist.}\pdfcomment[icon=Note]{Source 5142478 (p.29, p.30) analyzes the DSA and EU AI Act and concludes that despite transparency reporting, engagement-driven biases in core algorithms remain largely unaddressed, allowing feedback loops favoring sensational or extreme content to persist, directly supporting the claim.} \hldarkgreen{Proposed responses include fairness metrics, media literacy, and explanatory AI tools that can expose why particular items appear in feeds and allow both users and regulators to interrogate the loop between interaction signals and ranking outputs \parencites[5]{5142478}[8]{5142478}[9]{5142478}.}\pdfcomment[icon=Note]{Source 5142478 (p.9, p.30) explicitly recommends fairness metrics, media literacy, and explanatory AI (XAI) tools to reveal why items appear in feeds and to allow interrogation of the interaction--ranking loop, directly supporting this sentence.} \hllightgreen{Classic work on bounded rationality and heuristics suggests an additional interpretive layer: users often satisfice and rely on availability or anchoring shortcuts when interacting with feeds, which means their click signals may systematically misrepresent stable preferences and instead reflect transient cognitive constraints.}\pdfcomment[icon=Note]{Source 3232181 (pp.2--6) summarizes bounded rationality and heuristics (availability, anchoring, satisficing), supporting the inference that users often satisfice and rely on such shortcuts so their click signals may misrepresent stable preferences.} \hllightgreen{This has implications for training data quality, since heuristically biased traces are precisely the patterns that ranking models learn from \parencites[4]{3232181}.}\pdfcomment[icon=Note]{Sources 5142478 (p.5) and 6249388 (p.19) on recommender training and data sparsity support the implication that heuristically biased interaction traces degrade training-data quality and are the very patterns ranking models learn from, a reasonable synthesis.}

\subsubsection{Data Flows, Logging and Profiling Practices}
\label{sec:data_flows_logging_profiling} \hllightgreen{Data capture in AI mediated educational and proctoring systems illustrates how continuous telemetry streams become the substrate for profiling and risk assessment.}\pdfcomment[icon=Note]{Reviews of AI-mediated proctoring describe continuous telemetry from platforms as the basis for monitoring and risk assessment (2055956 p.2; 0751607 p.5), so this is a reasonable synthesis of those sources.} \hldarkgreen{Internet of Things enabled learning platforms aggregate biometric authentication logs such as voice patterns, keystroke dynamics, and facial recognition, alongside behavioural analytics including session length, login frequency, and activity engagement.}\pdfcomment[icon=Note]{0751607 p.5 explicitly lists biometric authentication logs (voice, keystroke, facial recognition) and behavioural analytics (session length, login frequency, engagement) as aggregated data sources in IoT-enabled educational platforms.} \hldarkgreen{Assessment metadata like time spent per question and plagiarism detection patterns, together with network logs such as IP addresses, device types, and geolocation anomalies, are combined into a unified dataset used to distinguish genuine student activity from fraud or policy violations.}\pdfcomment[icon=Note]{0751607 p.5 describes assessment metadata (time per question, plagiarism patterns) and network logs (IP, device, geolocation anomalies) being combined to train models to distinguish genuine activity from fraud.} \hldarkgreen{Preprocessing pipelines then impute missing values, normalize features to a common scale, remove duplicated or uninformative fields, and encode categorical states such as login status or exam completion into numerical form, which structures the raw flow of platform interaction into model ready profiles \parencites[5]{0751607}.}\pdfcomment[icon=Note]{0751607 p.5 details preprocessing steps including mean/mode imputation, normalization, removal of superfluous features, and encoding of categorical variables for ML readiness.} \hldarkgreen{Financial AI architectures employing cloud data lakes provide a parallel example of high volume, multi tenant data flows.}\pdfcomment[icon=Note]{8907572 p.4--5 describes cloud-based financial architectures with high-volume multi-tenant data flows and object-storage data lakes as a parallel example.} \hldarkgreen{Transaction streams are ingested via distributed messaging systems such as Apache Kafka into scalable object storage, with strict row level security and cryptographically separated tenant isolation to avoid cross contamination between clients.}\pdfcomment[icon=Note]{8907572 p.4 explicitly states transaction streams are captured via distributed messaging systems (e.g., Apache Kafka) into scalable cloud object storage and describes row-level security with cryptographic tenant isolation.} \hldarkgreen{Despite formal exclusion of protected attributes like race or gender, historical biases enter at ingestion because seemingly neutral variables, for example geolocation, school history, or spending patterns, act as precise proxies for demographic identity within deep neural networks.}\pdfcomment[icon=Note]{8907572 p.4--5 documents how protected attributes omitted at ingestion can be reconstructed from proxy features (geolocation, school history, spending patterns), leading to historical bias entering models.} \hldarkgreen{An intermediate telemetry parsing layer that cleans and normalizes incoming vectors then exposes them to a domain specific bias mitigation engine, ensuring that ethical evaluation and mitigation are integrated into the live data flow rather than treated as a separate batch step.}\pdfcomment[icon=Note]{8907572 p.5 describes an intermediate telemetry parsing layer that cleans/normalizes incoming vectors and exposes them to a domain-specific bias mitigation engine as part of the live ingestion architecture.} \hldarkgreen{Within these pipelines, logging practices are not an ancillary concern but a regulatory and ethical requirement.}\pdfcomment[icon=Note]{The AI law model emphasizes transparency, documentation and logging as legal/ethical obligations (2977039 p.22) and TSAI requires preserving and providing event logs (2977039 p.243--244), supporting this claim.} \hldarkgreen{In the financial context, asynchronous SHAP and LIME computations are executed in parallel with primary inference so that each decision is accompanied by a localized explanation vector indicating the contribution of individual features.}\pdfcomment[icon=Note]{8907572 p.5--6 describes asynchronous, distributed execution of SHAP and LIME running in parallel with primary inference to produce localized explanation vectors for each decision.} \hldarkgreen{This explanation metadata is stamped with a cryptographic hash and written to an immutable cloud ledger, creating a tamper evident audit trail suitable for future supervisory inspection \parencites[4]{8907572}.}\pdfcomment[icon=Note]{8907572 p.6 states that explanation metadata is stamped with a cryptographic hash and pushed to an immutable cloud ledger to create a tamper-proof audit trail.} \hldarkgreen{HITL grading deployments adopt a similar dual logging approach.}\pdfcomment[icon=Note]{The HITL grading case study describes dual-logging audit mechanisms as part of the deployment (6805789 p.16--19), supporting this statement.} \hldarkgreen{Every AI and instructor decision is time stamped, linked to reviewer identifiers, and stored as part of an audit trail that supports quality assurance and makes human overrides or confirmations traceable at the level of individual essays and rubric dimensions \parencites[16]{6805789}[17]{6805789}.}\pdfcomment[icon=Note]{6805789 p.17 and p.11 report that AI and instructor decisions are timestamped, linked to reviewer identifiers, and stored in audit trails during the HITL pilot.} \hllightgreen{Ethically grounded frameworks for AI in education and assessment extend logging beyond technical traceability toward explicit privacy and fairness guarantees.}\pdfcomment[icon=Note]{Both the HITL case study (6805789 p.20) and the AI law model (2977039 p.22) connect logging to broader ethical goals like privacy, fairness and explainability, so this is a reasonable synthesis.} \hldarkgreen{HITL grading systems configure modular architectures where anonymization of student records, encrypted storage of audit logs, opt out options for AI assisted feedback, and clear attribution of comments as "AI generated" or "instructor reviewed" form part of the system design.}\pdfcomment[icon=Note]{6805789 p.20 explicitly describes system features including anonymization of student records, encrypted audit trails, opt-out mechanisms for AI feedback, and tagging of comments as AI-generated or instructor-reviewed.} \hldarkgreen{These features are explicitly positioned as aligning with FERPA and GDPR requirements and as supporting institutional responsibility for transparent, human supervised use of AI in high stakes evaluation \parencites[20]{6805789}.}\pdfcomment[icon=Note]{6805789 p.20 states these HITL data governance protocols align with privacy standards under FERPA and GDPR and support institutional accountability.} \hldarkgreen{In IoT enabled fraud detection for educational platforms, encryption for data transport and storage, anonymization of personally identifiable information, and compliance with FERPA and GDPR are likewise described as integral components, complemented by bias mitigation techniques to avoid discriminatory outcomes in the classification of student behaviour \parencites[5]{0751607}.}\pdfcomment[icon=Note]{0751607 p.5 describes IoT-enabled educational fraud-detection systems using encryption, anonymization, FERPA/GDPR compliance and bias mitigation techniques as integral components.} \hldarkgreen{Regulatory AI law models generalise these obligations into structured digital identities and event logging duties for high risk systems.}\pdfcomment[icon=Note]{2977039 p.243--244 generalizes obligations such as event logging and digital identities into the TSAI regulatory model for high-risk systems.} \hldarkgreen{The TSAI regime specifies that owners, suppliers, or deployers of AI must ensure human oversight, preserve and provide event logs, conduct audits, and comply with data protection, non discrimination, and transparency requirements.}\pdfcomment[icon=Note]{2977039 p.243--244 lists obligations for owners/suppliers/deployers including human oversight, saving/providing event logs, conducting audits, and compliance with data protection, non-discrimination and transparency.} \hldarkgreen{These duties are operationalised through an algorithmic passport that includes sources and categories of data, legal grounds for processing, restrictions on reuse, modes of event logging, log storage periods, evidence mechanisms, and information on audits and trust markings \parencites[242]{2977039}[244]{2977039}.}\pdfcomment[icon=Note]{2977039 p.243--244 specifies that the algorithmic passport must include data sources/categories, legal grounds, reuse restrictions, logging modes, storage periods, evidence mechanisms, and audit/trust marking information.} \hldarkgreen{For AI in public communications, logging and saving interaction logs in a form that permits independent auditing is listed among the safeguards that distinguish permitted dual use tools from prohibited manipulative systems, and the intensity of logging duties scales with risk levels from increased to critical, linked to potential public harm and coordinated automation tests.}\pdfcomment[icon=Note]{2977039 p.95--96 identifies logging/saving interaction logs to permit independent auditing as a safeguard for dual-use public communications tools and links intensity of obligations to assessed risk levels.} \hldarkgreen{Platforms that algorithmically amplify manipulative or disinformative content without implementing such safeguards are described as bearing public responsibility, with independent audit mandated to verify the effectiveness of watermarking, labelling, quota controls, and log based evidence collection \parencites[95]{2977039}[97]{2977039}.}\pdfcomment[icon=Note]{2977039 p.95--96 and p.96 describe that platforms amplifying manipulative/disinformative content without safeguards bear public responsibility and require independent audit to verify watermarking, labeling, quotas and log-based evidence.} \hllightgreen{Empirical reviews of surveillance deployments further suggest that privacy impact assessments, stakeholder outreach including affected communities, and routine independent audits are practical mechanisms to check for function creep and to make logging practices more publicly accountable; these measures may therefore be appropriate complements to cryptographic and technical safeguards in high risk educational or public-facing systems \parencites[6]{5988547}.}\pdfcomment[icon=Note]{Empirical and policy reviews in the literature recommend PIAs, stakeholder engagement (including affected communities), and independent audits as accountability mechanisms (5988547 p.6; 5988547 p.2), so this is a supported synthesis that complements cryptographic safeguards.}

\subsection{Engagement-Optimized Interface Design in Social Platforms}

\subsubsection{Infinite Scroll and Continuous Content Streams}
\label{sec:infinite_scroll_streams} \hldarkgreen{Infinite scrolling and continuous content streams operationalise algorithmic addiction at the interface layer by removing natural stopping cues and tightly coupling user gestures to reward expectation.}\pdfcomment[icon=Note]{Sources explicitly link infinite scroll and continuous streams to dopamine-driven interface effects and removal of exit cues (3303310, p.2; 8958037, p.3), describing how gesture-reward coupling operationalises addictive dynamics.} \hldarkgreen{Short video platforms that rely on neural network based recommendation systems combine micro duration clips with vertically swiped feeds so that a single low effort action reveals the next item, which is selected from a personalised distribution rather than a chronological queue \parencites[2]{8958037}[3]{8958037}.}\pdfcomment[icon=Note]{8958037 (p.2--3) directly describes short-video platforms using neural-network recommendation, micro-duration clips, vertical swipe feeds, and personalised distributions rather than chronological queues.} \hldarkgreen{This design supports a closed cycle of "personalized $\rightarrow$ user need satisfaction $\rightarrow$ behavioural data accumulation $\rightarrow$ algorithm optimization $\rightarrow$ re push," in which each additional scroll both delivers content and refines the model that will determine subsequent items.}\pdfcomment[icon=Note]{8958037 (p.3) explicitly defines the closed loop 'personalized → user need satisfaction → behavioural data accumulation → algorithm optimization → re-push' and notes that each scroll both delivers content and refines the model.} \hldarkgreen{As precision improves, dependence on the stream increases and self control over short video consumption declines, particularly when information cocoons narrow exposure and limit alternative cues for disengagement.}\pdfcomment[icon=Note]{8958037 (p.3) states that increasing precision of recommendations fosters user dependence and diminished self-control, and that information cocoons narrow exposure and limit disengagement cues.} \hllightgreen{Neuroscientific explanations underline why the absence of friction in infinite scroll is not a neutral usability choice.}\pdfcomment[icon=Note]{Neuroscientific accounts in 3303310 (p.2) and 8958037 (p.3) describe dopamine-driven feedback loops and the neural effects of frictionless interfaces, supporting the inference that absence of friction is not a neutral usability choice.} \hldarkgreen{Continuous personalised viewing stimulates the ventral tegmental area, with dopamine surging in the moment of expectancy before the next video is revealed rather than only when a clip is enjoyed.}\pdfcomment[icon=Note]{8958037 (p.3) explicitly states that continuous personalized viewing stimulates the ventral tegmental area and that dopamine surges during anticipation before the next video is revealed.} \hldarkgreen{Users are drawn into a repeated "pursuit $\rightarrow$ transient satisfaction $\rightarrow$ further pursuit" cycle, where each swipe triggers anticipation of reward and reinforces the habit of extending the session.}\pdfcomment[icon=Note]{8958037 (p.3) describes the 'pursuit → transient satisfaction → further pursuit' cycle and notes that each swipe triggers anticipation and reinforces session extension.} \hldarkgreen{Prolonged high frequency stimulation may induce expression of $\Delta$FosB in reward circuitry and alter gene expression, which aligns these usage patterns with other forms of addiction physiology \parencites[3]{8958037}.}\pdfcomment[icon=Note]{8958037 (p.3) explicitly references that prolonged high-frequency stimulation may induce ΔFosB expression in reward circuitry and alter gene expression, linking these patterns to addiction physiology.} \hldarkgreen{The "sunflower seed effect" of 15--60 second high stimulation videos, combined with infinite scroll, disrupts default mode network functioning and undermines self regulatory capacity, especially for adolescents and college students whose executive functions are still maturing \parencites[2]{8958037}[3]{8958037}.}\pdfcomment[icon=Note]{8958037 (p.3) describes the 'sunflower seed effect' of 15--60 second videos combined with infinite scroll, and states these formats disrupt default mode network functioning and reduce self-regulatory capacity, particularly in adolescents.} \hllightgreen{Empirical and conceptual work on social media feeds extends these dynamics beyond short video services.}\pdfcomment[icon=Note]{Broader empirical and conceptual literature (e.g., 3303310, p.2; 5142478, p.29) extends the described dynamics from short-video platforms to social media feeds generally.} \hldarkgreen{Engagement driven content prioritisation in algorithmically curated streams has been associated with behavioural addiction, algorithmic amplification of misinformation, and reinforcement of societal polarization.}\pdfcomment[icon=Note]{3303310 (p.2) and 5142478 (p.29) link engagement-driven content prioritisation to behavioural addiction, algorithmic amplification of misinformation, and reinforcement of societal polarization.} \hldarkgreen{Infinite scrolling interfaces and variable reward schedules are explicitly cited as core elements of "dopamine driven feedback loops," with empirical validations showing that users of algorithmically curated feeds exhibit 22\% higher rates of compulsive checking behaviour compared with non algorithmic interfaces.}\pdfcomment[icon=Note]{3303310 (p.2) cites Alter's model naming infinite scroll and variable reward schedules as core to 'dopamine-driven feedback loops' and reports empirical validation that algorithmic feeds show 22\% higher compulsive checking rates.} \hldarkgreen{Neuroimaging studies connect prolonged social media use under such feeds to reduced gray matter density and reduced prefrontal cortex activity in regions governing impulse control, mirroring patterns seen in substance dependence and reinforcing the concern that continuous streams condition addictive engagement rather than merely reflect user preference \parencites[2]{3303310}.}\pdfcomment[icon=Note]{3303310 (p.3) reports neuroimaging findings that prolonged social media use correlates with reduced gray matter density and reduced prefrontal cortex activity in impulse-control regions, paralleling substance-dependence patterns.} \hllightgreen{Regulatory and policy analyses increasingly treat infinite scroll and continuous streams as potential manipulative design.}\pdfcomment[icon=Note]{Regulatory and policy analyses in the provided sources (8958037, p.7; 5142478, p.29; 6295677, p.6) increasingly characterise infinite scroll and continuous streams as potentially manipulative design features.} \hldarkgreen{Under the EU Digital Services Act, TikTok's algorithmic recommendation mechanisms and interface features, including infinite scrolling, auto play, and push notifications, have been investigated as systematic manipulative design that impairs users' capacity to make free and informed decisions \parencites[2]{8958037}.}\pdfcomment[icon=Note]{8958037 (p.2) explicitly states that the EU under the DSA investigated TikTok's recommendation mechanisms and interface features (infinite scrolling, autoplay, push notifications) as systematic manipulative design impairing free, informed decisions.} \hldarkgreen{Article 25 DSA explicitly prohibits platforms from designing manipulative interfaces that undermine autonomous decision making, and Chinese rules such as the *Administrative Provisions on Algorithmic Recommendation for Internet Information Services* require that providers offer options not based on personal characteristics and convenient functions to disable recommendation services \parencites[7]{8958037}.}\pdfcomment[icon=Note]{8958037 (p.7) explicitly references Article 25 of the DSA prohibiting manipulative interfaces and cites the Chinese Administrative Provisions requiring options not based on personal characteristics and functions to disable recommendation services.} \hldarkgreen{Analysis of democratic recommender governance similarly identifies TikTok's "For You Page" as an exceptionally effective recommendation system that captures and holds attention, with very high engagement and retention, while remaining opaque to public scrutiny \parencites[16]{6502223}.}\pdfcomment[icon=Note]{6502223 (p.16) analyses democratic recommender governance and identifies TikTok's For You Page as an exceptionally effective, high-engagement, high-retention recommendation system that remains opaque to scrutiny.} \hllightgreen{Research on algorithmic engagement worship adds that such feeds produce a paradox where platforms simultaneously connect users and manipulate cognitive autonomy, and it calls for transparency by design and human centred interventions to realign continuous streams with user wellbeing \parencites[2]{3303310}.}\pdfcomment[icon=Note]{3303310 (p.2) describes the paradox that algorithmic engagement both connects users and manipulates cognitive autonomy and calls for transparency-by-design and human-centred interventions, supporting this claim as a synthesis.} \hldarkgreen{Interface level mitigations proposed in the literature target precisely the frictionless nature of infinite scroll.}\pdfcomment[icon=Note]{8958037 (p.7) proposes interface-level mitigations that specifically target the frictionless nature of infinite scroll.} \hldarkgreen{Suggested measures include introducing randomized content or user customisable filters to disrupt closed information cocoons and embedding active exit nodes, for example intelligent usage duration reminders, timed automatic pauses, and enhanced natural exit signals, as counterweights to default infinite pathways \parencites[7]{8958037}.}\pdfcomment[icon=Note]{8958037 (p.7) explicitly lists measures such as randomized content or user-customisable filters, intelligent usage reminders, timed automatic pauses, and enhanced natural exit signals as counterweights to infinite scroll.} \hldarkgreen{Enhancing transparency by explaining "why this content" is being shown is framed as a move from black box manipulation to visible assistance, aligning interface design with DSA requirements and emerging expectations around explainable recommendation \parencites[7]{8958037}[8]{5142478}.}\pdfcomment[icon=Note]{8958037 (p.7) and 5142478 (p.29) discuss explaining 'why this content' is shown as a transparency/XAI measure that aligns with DSA-like expectations and explainable recommendation practices.} \hllightgreen{Within this thesis, such proposals provide concrete audit criteria: continuous content streams become ethically problematic where infinite scroll is combined with engagement maximisation objectives, opacity about ranking logic, and absence of exit cues, especially for vulnerable user groups such as adolescents and students with documented susceptibility to algorithmic recommendation environments \parencites[2]{8958037}[3]{8958037}[2]{3303310}[16]{6502223}.}\pdfcomment[icon=Note]{Combined sources (8958037, p.7; 3303310, p.2--3; 6502223, p.16) support the thesis-level audit criteria that continuous streams are ethically problematic when infinite scroll is paired with engagement-maximisation, opacity, lack of exit cues, and vulnerable groups (adolescents/students) are implicated.}

\subsubsection{Autoplay, Short-Form Clips and Sound Design}
\label{sec:autoplay_shortform_sound} \hldarkgreen{Autoplay and short clip formats operationalise what Li describes as an algorithmic "addiction loop," where content is pushed in rapid succession without requiring intentional selection from the user.}\pdfcomment[icon=Note]{Dan Li explicitly frames an algorithmic "addiction loop" and links autoplay/short-clip delivery to rapid, non-intentional content push (Li 8958037, p.2).} \hldarkgreen{Short video platforms rely on neural network based recommendation systems that continually optimise content matching and are tightly coupled with micro duration clips that can be consumed with minimal effort \parencites[2]{8958037}.}\pdfcomment[icon=Note]{Li states that short-video platforms rely on neural-network recommendation systems that continuously optimize matching and are coupled with micro-duration stimulus designs (Li 8958037, p.2).} \hldarkgreen{Viewing history, likes, and comments are processed through content based and collaborative filtering, which together create a closed cycle of "personalized $\rightarrow$ user need satisfaction $\rightarrow$ behavioural data accumulation $\rightarrow$ algorithm optimization $\rightarrow$ re push." As the loop tightens, precision of information delivery rises and self control over viewing behaviour declines, particularly when information cocoons emerge that restrict access to diverse content and intensify dependency on the feed \parencites[3]{8958037}.}\pdfcomment[icon=Note]{Li describes content-based and collaborative filtering using viewing history/likes/comments, the closed-loop "personalized → ... → re-push," rising delivery precision and reduced self-control with information cocoons (Li 8958037, p.3).} \hldarkgreen{Autoplay embeds this cycle directly into the interface: every completed clip immediately triggers another, so the path of least resistance is continued consumption rather than deliberate stopping \parencites[2]{8958037}[3]{8958037}.}\pdfcomment[icon=Note]{Li and the cited work describe autoplay/infinite-scroll embedding the cycle so each completed clip triggers another, promoting continued consumption (Li 8958037, p.2; p.3).} \hllightgreen{Short form clips interact with this autoplay logic in specific ways.}\pdfcomment[icon=Note]{Li discusses specific interactions between short-form clip formats and autoplay mechanics, so the general claim that they interact is a reasonable synthesis of that material (Li 8958037, p.3).} \hldarkgreen{Li reports that dominant formats are 15--60 second videos, with some platforms popularising 15 second vertical, high stimulation clips that are easy to chain together through infinite swipe based navigation.}\pdfcomment[icon=Note]{Li reports short videos are extremely brief (15--60 seconds) and highlights the prominence of 15-second vertical, high-stimulation formats that are easy to chain (Li 8958037, p.3).} \hldarkgreen{The extremely low operational cost of swiping to the next video enables effortless extension of sessions; users can watch for close to 60 minutes per day with more than 21 sessions on services such as Douyin, where 41.3\% of users open the app multiple times per day and 56.9\% exceed 60 minutes of daily use \parencites[3]{8958037}.}\pdfcomment[icon=Note]{Li provides the Douyin usage statistics and notes low operational cost of swiping enabling extended sessions: \textasciitilde{}21+ sessions and \textasciitilde{}60 minutes, with 41.3\% multiple daily opens and 56.9\% >60 mins (Li 8958037, p.3).} \hllightgreen{These metrics reflect how short clip duration and frictionless progression form the behavioural substrate on which recommendation engines optimise engagement.}\pdfcomment[icon=Note]{Li links short clip duration and frictionless progression to the behavioral substrate that recommendation engines optimize for engagement, a reasonable synthesis of the metrics and mechanism (Li 8958037, p.2--3).} \hldarkgreen{For college students and adolescents with immature self control, this configuration is associated with heightened susceptibility to addiction, diminished academic concentration, and increased anxiety and depression \parencites[2]{8958037}[6]{8958037}.}\pdfcomment[icon=Note]{Li documents that college students/adolescents with developing self-control are especially susceptible to addiction and notes diminished concentration, anxiety and depression as consequences (Li 8958037, p.2; p.6).} \hldarkgreen{Sound design amplifies these dynamics because short videos are routinely paired with dynamic music and audio cues that deliver continuous high stimulation.}\pdfcomment[icon=Note]{Li describes pairing of short videos with dynamic music and audio cues that increase stimulation and capture attention (Li 8958037, p.3--4).} \hldarkgreen{Li characterises this as a "sunflower seed effect," where 15 second vertical clips, combined with vivid visuals and music, repeatedly capture attention, disrupt the brain's default mode network, and compromise self regulatory capacity \parencites[3]{8958037}.}\pdfcomment[icon=Note]{Li explicitly uses the "sunflower seed effect" to describe 15-second vertical clips with vivid visuals and music disrupting the default mode network and impairing self-regulation (Li 8958037, p.3).} \hldarkgreen{Neurobiological analysis links these interface properties to dopaminergic reward pathways.}\pdfcomment[icon=Note]{Li links these interface properties to dopaminergic reward pathways and related neuropsychological mechanisms (Li 8958037, p.3--4).} \hldarkgreen{The content switching mechanism of short video platforms is described as structurally similar to slot machine variable rewards: each downward swipe carries an unpredictable probability of encountering more engaging content, which strongly activates the mesolimbic dopamine system.}\pdfcomment[icon=Note]{Li describes the content-switching mechanism as structurally similar to slot-machine variable rewards and states this unpredictability maximally activates the dopaminergic/mesolimbic system (Li 8958037, p.4).} \hldarkgreen{Dopamine surges primarily during anticipation, that is in the moment before the next clip is revealed, trapping users in a cycle of "pursuit $\rightarrow$ transient satisfaction $\rightarrow$ further pursuit." Chronic high frequency stimulation can lead to downregulation of dopamine receptors and elevated reward thresholds, reducing the hedonic value of natural rewards and creating a vicious cycle where increased usage breeds dissatisfaction that in turn drives further usage \parencites[4]{8958037}.}\pdfcomment[icon=Note]{Li states dopamine surges during anticipation, describes the "pursuit → transient satisfaction → further pursuit" cycle, and notes chronic stimulation can downregulate receptors and raise reward thresholds (Li 8958037, p.3--4).} \hllightgreen{Engagement oriented social media research corroborates these findings.}\pdfcomment[icon=Note]{Other engagement-focused research cited (e.g., 3303310) corroborates Li's neurobehavioral and engagement findings, so this statement is a reasonable synthesis (3303310, p.2).} \hllightgreen{Hu synthesises behavioural psychology and platform studies to show that infinite scrolling interfaces and variable reward schedules are core features of systems like TikTok and are associated with 22\% higher rates of compulsive checking behaviour in algorithmically curated feeds compared with non algorithmic interfaces.}\pdfcomment[icon=Note]{The proceedings paper reports that infinite scrolling and variable rewards are core features and cites empirical findings of \textasciitilde{}22\% higher compulsive checking in algorithmic feeds, supporting this claim (3303310, p.2).} \hllightgreen{Neuroimaging studies cited in the same work correlate prolonged use of such feeds with reduced prefrontal cortex activity and altered gray matter density in regions responsible for impulse control, mirroring patterns observed in substance dependence \parencites[2]{3303310}.}\pdfcomment[icon=Note]{3303310 and related citations discuss neuroimaging correlates linking prolonged algorithmic feed use to reduced prefrontal activity and structural changes similar to substance dependence, supporting this synthesis (3303310, p.2).} \hllightgreen{Short form videos with autoplay and high arousal soundscapes therefore do not merely mirror pre existing preferences; they function as stimulus packages tuned to exploit neural pathways associated with addiction.}\pdfcomment[icon=Note]{Li and the engagement literature characterize short-form videos with autoplay/sound as stimulus packages exploiting neural addiction pathways, making this a supported synthesis (Li 8958037, p.3--4; 3303310, p.2).} \hllightgreen{Regulatory sources treat these mechanics as potential manipulative design rather than neutral entertainment features.}\pdfcomment[icon=Note]{The sources show regulatory attention to these mechanics as potentially manipulative (Li notes EU DSA scrutiny; AI Law Model prohibits manipulative patterns), supporting this claim (Li 8958037, p.2; AI Law Model 2977039, p.223--224).} \hldarkgreen{Under the Digital Services Act, the European Commission opened proceedings into TikTok's algorithmic recommendation mechanisms and interface features, explicitly naming infinite scrolling, auto play, and push notifications as possible systematic manipulative design that impairs users' capacity to make free and informed decisions \parencites[2]{8958037}.}\pdfcomment[icon=Note]{Li reports the EU under the DSA opened an investigation into TikTok in 2026, explicitly naming infinite scrolling, autoplay and push notifications as possible systematic manipulative design (Li 8958037, p.2).} \hldarkgreen{Kostenko's AI Law Model reinforces this framing by prohibiting gamification of addiction, reward games and streaks, exploitation of age related vulnerabilities, and manipulative interface patterns.}\pdfcomment[icon=Note]{The AI Law Model explicitly prohibits gamification of addiction, reward games/streaks, exploitation of age-related vulnerabilities and manipulative interface patterns (AI Law Model 2977039, p.224).} \hldarkgreen{For minors, it mandates age appropriate design, secure defaults with profiling disabled, and explicit bans on endless scrolling and autoplay, alongside limits on push notifications and requirements for visible "no personalization" switches and clear explanations such as "why am I seeing this" \parencites[223]{2977039}.}\pdfcomment[icon=Note]{The AI Law Model mandates age-appropriate design, secure defaults with profiling disabled, bans on endless scrolling/autoplay, limits on notifications, and visible 'no personalization' switches and explanatory controls (AI Law Model 2977039, p.224).} \hldarkgreen{Reliability and resilience provisions in the same model require systematic incident logging, risk management, and independent audits for high risk systems, which would include platforms deploying such engagement mechanics at scale \parencites[57]{2977039}[79]{2977039}.}\pdfcomment[icon=Note]{The AI Law Model requires systematic incident logging, risk management, regular audits and independent assessments for high-risk systems (AI Law Model 2977039, pp.78--79; p.57).} \hllightgreen{Within this thesis, these findings underscore why autoplay, short form clips, and sound design must be treated as audit targets.}\pdfcomment[icon=Note]{Given the empirical, neurobiological and regulatory evidence presented (Li 8958037; AI Law Model 2977039), treating autoplay/short clips/sound design as audit targets is a reasonable inference within the thesis.} \hllightgreen{They sit at the intersection of dopamine driven feedback loops, behavioural addiction indicators, and explicit regulatory concern about manipulative interfaces, especially where adolescents and students form core user groups with documented vulnerability to algorithmic recommendation environments \parencites[2]{8958037}[3]{8958037}[2]{3303310}[223]{2977039}.}\pdfcomment[icon=Note]{This synthesis, linking dopamine-driven feedback, behavioral addiction indicators, and regulatory concern with vulnerability of adolescents/students, is supported across Li (8958037), the regulatory model (2977039), and engagement research (3303310) (Li 8958037, p.2--6; AI Law Model 2977039, p.223--224; 3303310, p.2).}

\subsubsection{Notification Systems and Re-Engagement Triggers}
\label{sec:notification_reengagement_triggers} \hldarkgreen{Push notifications, streak reminders, and related prompts extend engagement-optimised design beyond the feed into the user's off-platform environment.}\pdfcomment[icon=Note]{Source 8958037 (p.2p.3) explicitly lists push notifications, streak reminders and related prompts as interface features that extend engagement-optimised design beyond the feed and into users' wider environment.} \hldarkgreen{Research on short video ecosystems identifies push notifications alongside infinite scrolling and auto play as core elements of an "algorithmic addiction loop" in which personalised recommendations, micro duration clips, and external triggers collectively sustain persistent scrolling and weaken self control among college students \parencites[2]{8958037}[3]{8958037}.}\pdfcomment[icon=Note]{Source 8958037 (p.2) frames infinite scrolling, auto-play and push notifications as core elements of an "algorithmic addiction loop" and links these features to weakened self-control among college students.} \hldarkgreen{Within this loop, notifications operate as re entry cues: they call users back into the personalised cycle of "personalized $\rightarrow$ user need satisfaction $\rightarrow$ behavioural data accumulation $\rightarrow$ algorithm optimization $\rightarrow$ re push," thereby refreshing the data used for further targeting and reinforcing behavioural dependence.}\pdfcomment[icon=Note]{Source 8958037 (p.3) explicitly describes the closed-loop "personalized → user need satisfaction → behavioral data accumulation → algorithm optimization → re-push" and discusses re-push/re-entry dynamics.} \hllightgreen{Neurocognitive analysis deepens this account.}\pdfcomment[icon=Note]{Source 8958037 (p.3) moves from platform mechanics to neural mechanisms, so invoking a neurocognitive analysis as deepening the account is a reasonable framing supported by that source.} \hldarkgreen{Short video consumption is associated with dopamine driven "reward circuitry" where dopamine surges primarily during anticipation of the next reward, not only when a clip is enjoyed.}\pdfcomment[icon=Note]{Source 8958037 (p.3) describes short-video platforms engaging dopamine-driven reward circuitry and specifically notes dopamine surges during anticipation of the next reward rather than only at reward receipt.} \hllightgreen{When notifications signal that new, potentially rewarding content is available, they may serve as anticipatory cues that trigger the same "pursuit $\rightarrow$ transient satisfaction $\rightarrow$ further pursuit" cycle even before the user re opens the app.}\pdfcomment[icon=Note]{Combining 8958037's account of anticipatory dopamine surges (p.3) with its description of notifications as part of the loop (p.2p.3) supports the inference that notifications can act as anticipatory cues prior to app re-opening.} \hldarkgreen{High frequency stimulation in this cycle can induce expression of $\Delta$FosB in reward circuits, altering gene expression and supporting addiction like states \parencites[3]{8958037}.}\pdfcomment[icon=Note]{Source 8958037 (p.3) explicitly states prolonged high-frequency stimulation may induce expression of ΔFosB in reward circuits, linking molecular changes to addiction-like states.} \hldarkgreen{For adolescents and college students, whose executive functions and self control capacities are still maturing, such externally injected prompts sit atop an already vulnerable psychological profile and are associated with heightened anxiety, depression, and behavioural addiction when usage becomes excessive \parencites[2]{8958037}[3]{8958037}.}\pdfcomment[icon=Note]{Source 8958037 (p.2p.3) documents college students' developmental vulnerability and associates excessive short-video use with diminished concentration, weakened social skills, heightened anxiety and depression, and behavioural addiction.} \hllightgreen{Engagement oriented social media studies add a behavioural layer.}\pdfcomment[icon=Note]{Sources such as 3303310 (p.2) and 6295677 discuss engagement-oriented platform features and behavioural effects, so saying engagement studies add a behavioural layer is a valid synthesis.} \hldarkgreen{Alter's model of "dopamine driven feedback loops" is applied to platforms like TikTok, where infinite scrolling interfaces and variable reward schedules correlate with 22\% higher rates of compulsive checking for algorithmically curated feeds compared with non algorithmic interfaces.}\pdfcomment[icon=Note]{Source 3303310 (p.2) cites Alter's model of "dopamine-driven feedback loops" and reports empirical validations showing \textasciitilde{}22\% higher rates of compulsive checking for algorithmically curated feeds versus non-algorithmic interfaces.} \hllightgreen{While this work focuses primarily on in feed dynamics, compulsive checking behaviour necessarily interacts with notification policies that repeatedly invite users back into the environment.}\pdfcomment[icon=Note]{3303310 focuses on in-feed dynamics (p.2) while 8958037 and related sources document notification policies (8958037 p.2), so the claim that compulsive checking interacts with notifications is a reasonable synthesis.} \hldarkgreen{Neuroimaging findings cited in the same study link prolonged use of such feeds to reduced prefrontal cortex activity and altered gray matter density in regions responsible for impulse control, aligning the resulting pattern of "doomscrolling" with markers of substance dependence \parencites[2]{3303310}.}\pdfcomment[icon=Note]{Source 3303310 (p.2) and related literature cited there report neuroimaging links between prolonged use of algorithmic feeds and reduced prefrontal cortex activity and structural/gray-matter alterations associated with impulse-control deficits.} \hllightgreen{In this context, re engagement triggers are not neutral reminders; they are part of a broader conditioning process in which cues, rewards, and behavioural responses are tightly coupled.}\pdfcomment[icon=Note]{Sources on conditioning, dark patterns and algorithmic feedback (6295677 p.8; 8958037 p.2p.3; 3303310 p.2) jointly support the view that re-engagement triggers function as conditioning cues tightly coupling cues, rewards and responses.} \hldarkgreen{Regulatory proposals in AI law respond directly to these risks by treating re engagement triggers as potential manipulative mechanics.}\pdfcomment[icon=Note]{Source 2977039 (p.8 and chapters) explicitly frames regulatory responses that address manipulative and behavioural practices, treating re-engagement mechanics as regulatory targets.} \hldarkgreen{Kostenko's AI Law Model prohibits hidden models of psychological influence, gamification of addiction such as reward games and streaks, exploitation of age related vulnerabilities, and manipulative interface patterns.}\pdfcomment[icon=Note]{Source 2977039 (p.223p.224) lists prohibitions including hidden models of psychological influence, gamification of addiction (reward games, streaks), exploitation of age-related vulnerabilities, and manipulative interface patterns.} \hldarkgreen{For minors, the principle of age appropriate design specifies secure defaults in which profiling, geolocation, biometrics, and manipulative mechanics are disabled.}\pdfcomment[icon=Note]{Source 2977039 (p.224) sets out the principle of age-appropriate design including secure default settings that disable profiling, geolocation, biometrics and related invasive mechanics for minors.} \hldarkgreen{Endless scrolling, autoplay, streaks, forced quests, and "penalties" for exiting are explicitly banned, and push notifications must be kept to a minimum and restricted to neutral informational messages.}\pdfcomment[icon=Note]{Source 2977039 (p.224) explicitly bans manipulative/addictive mechanics (endless scrolling, autoplay, streaks, forced quests, penalties for exiting) and limits push notifications to a minimum, permitting only reasonable/neutral informational messages.} \hldarkgreen{These provisions position high frequency, emotionally charged re engagement triggers as incompatible with acceptable practice where children are concerned and require visible "no personalization" switches, notification limits, and modes such as "lesson" or "sleep" that modulate attention demands \parencites[223]{2977039}.}\pdfcomment[icon=Note]{Source 2977039 (p.224) requires visible 'no personalization' controls, 'lesson'/'sleep' modes, notification limits and other measures, framing high-frequency emotional re-engagement triggers as incompatible with child-appropriate practice.} \hldarkgreen{Confidentiality and data protection principles in the same model reinforce constraints on notification driven profiling.}\pdfcomment[icon=Note]{Source 2977039 (p.43) establishes confidentiality and data protection principles that reinforce constraints on profiling and similar practices, thereby supporting limits on notification-driven profiling.} \hldarkgreen{All AI systems processing personal data must comply with confidentiality, integrity, and security requirements and must not use personal data for advertising, targeted marketing, or profiling aimed at influencing users in ways that restrict or distort freedom of choice without explicit informed consent.}\pdfcomment[icon=Note]{Source 2977039 (p.43) mandates that AI systems processing personal data comply with confidentiality, integrity and security requirements and prohibits use of personal data for advertising/targeted marketing or influence that distorts choice without explicit informed consent.} \hldarkgreen{Creation of behavioural profiles and tracking without a proper legal basis is prohibited, which directly affects data pipelines that generate notification triggers based on fine grained behavioural histories \parencites[44]{2977039}.}\pdfcomment[icon=Note]{Source 2977039 (p.43) prohibits creation of behavioural profiles and tracking without a proper legal basis, which directly impacts pipelines that would generate notification triggers from fine-grained behavioural histories.} \hldarkgreen{Research on dark patterns in e commerce confirms that repeated prompts and nagging in app notifications can wear down resistance and generate decision fatigue, illustrating how re engagement triggers can operate as dark patterns that exploit System 1 heuristics and suppress deliberative System 2 processing \parencites[2]{6295677}[8]{6295677}.}\pdfcomment[icon=Note]{Source 6295677 (p.8) documents dark-pattern categories including 'Nagging' (repeated prompts in app notifications), links this to System 1 processing and decision-fatigue, and thus supports the claim.} \hllightgreen{These converging lines of evidence and regulation will inform this report's working taxonomy of notification risk levels, in which high frequency, personalised, difficult to disable prompts are treated as high risk or potentially prohibited for vulnerable users, while clearly labelled, low frequency, informational notifications aligned with explicit consent are assessed as lower risk \parencites[2]{8958037}[3]{8958037}[2]{3303310}[223]{2977039}.}\pdfcomment[icon=Note]{This sentence synthesises regulatory material (2977039 p.223p.224) and empirical analyses (8958037 p.2; 3303310 p.2) to justify a risk taxonomy distinguishing high-frequency personalised hard-to-disable prompts from low-frequency labelled informational notifications; the cited sources together support this classification.} \hldarkgreen{Clinical and qualitative research on adolescents adds a cautionary note: many young people conceal severe distress and may seek information or coping strategies privately on their phones rather than disclose to caregivers, which suggests that notification-based interventions intended to prompt help-seeking may miss or even aggravate hidden crises \parencites[9]{2750607}.}\pdfcomment[icon=Note]{Clinical qualitative research (2750607, e.g. p.9p.11) documents adolescents actively concealing severe distress and seeking private phone-based resources, supporting the caution that notification-based interventions can miss or aggravate hidden crises.} \hllightgreen{It therefore appears prudent to treat re engagement cues aimed at minors as potentially intrusive interventions that require both stricter limits and explicit safeguards to avoid worsening acute, rapidly emerging crises \parencites[9]{2750607}.}\pdfcomment[icon=Note]{Combining the clinical evidence of hidden crises (2750607 p.9p.11) with regulatory safeguards (2977039 p.224) reasonably supports the policy conclusion that re-engagement cues for minors should be treated as intrusive and require stricter limits and safeguards.} \hldarkgreen{System architectures that produce these cues often rely on offline clustering, topic modelling, and sequence mining to surface highly contextually relevant moments for intervention or reengagement; such pipelines can make prompts both highly effective and difficult for users to detect or control \parencites[2]{7063165}.}\pdfcomment[icon=Note]{Technical literature on recommender pipelines (7063165, e.g. p.2p.11p.15) describes offline clustering, topic modelling and sequence-mining used to surface contextually relevant moments, supporting the claim that such pipelines make prompts effective and hard to detect/control.} \hllightgreen{This technical detail suggests a pragmatic policy implication: limiting the use of sequential and context-aware personalization in notification pipelines may be as important as restricting message content itself \parencites[2]{7063165}.}\pdfcomment[icon=Note]{Given the technical descriptions of sequential/context-aware personalization (7063165 p.2p.15) and the regulatory concerns about profiling and manipulative mechanics (2977039 p.43p.224), the policy inference that limiting sequential/context-aware personalization is as important as restricting content is a reasonable synthesis.}

\section{Behavioral and Psychological Mechanisms of Algorithmic Addiction}

\subsection{Behavioral Economics and Choice Architecture}

\subsubsection{Default Effects and Friction Costs in Digital Journeys}
\label{sec:default_effects_friction_costs} \hllightgreen{Default configurations and frictions in digital journeys act as powerful behavioral levers that intersect directly with algorithmic addiction and dark pattern design.}\pdfcomment[icon=Note]{Synthesis supported by behavioral and dark pattern literature: defaults/frictions as behavioral levers and their link to addictive recommender dynamics appear across sources (1070605 p.2-3; 6295677 p.8; 8958037 p.3).} \hldarkgreen{Defaults are described as "powerful nudges" that guide user behaviour by reducing cognitive effort so that individuals tend to remain with preselected options rather than explore alternatives, even when alternatives may better serve their interests.}\pdfcomment[icon=Note]{Directly supported by behavioral economics discussion of defaults as powerful nudges and reducing cognitive effort (1070605 p.2).} \hllightgreen{This status quo bias anchors preferences over time and enables digital platforms to reinforce dominance through seemingly neutral interface choices such as pre ticked boxes or pre chosen services \parencites[152]{1070605}.}\pdfcomment[icon=Note]{Status quo bias and use of pre‑ticked/preselected defaults as interface choices are documented across sources (1070605 p.2; 6295677 Table, p.8).} \hllightgreen{Empirical work in behavioral economics shows that defaults combined with social cues significantly increase adoption of targeted products, while consumers often make less informed choices because the cognitive relief of adhering to preselected options outweighs the perceived benefits of actively choosing an alternative.}\pdfcomment[icon=Note]{Empirical findings on defaults combined with social cues and reduced informed choice are reported in the competition/behavioral review (1070605 p.3) and related behavioral literature cited therein.} \hllightgreen{These dynamics explain why default settings inside platform ecosystems can entrench dominance and shape engagement patterns that users experience as "natural" or personalised, although they structurally reflect platform interests \parencites[153]{1070605}.}\pdfcomment[icon=Note]{Argument that defaults entrench platform dominance and are perceived as natural/personalised is supported by analysis of defaults, switching costs, and platform design (1070605 p.2-3; 6295677 p.8).} \hllightgreen{Dark pattern scholarship complements this account by specifying how default bias is operationalised as a design tactic.}\pdfcomment[icon=Note]{Dark pattern scholarship framing default bias as an operationalised design tactic is supported by typologies and behavioral foundations (6295677 p.4, p.8).} \hldarkgreen{Typologies of dark patterns explicitly list pre selected defaults as practices where automatically chosen options favour the platform, for example add ons or insurance, and note that such defaults have been banned by regulators such as IRDAI in insurance because they produce uninformed consent and financial leakage \parencites[8]{6295677}.}\pdfcomment[icon=Note]{Typologies explicitly list pre‑selected defaults and the table notes IRDAI ban in insurance; this is stated directly (6295677 Table p.8).} \hllightgreen{Marketing oriented taxonomies label similar tactics as "forced action and data extraction," including privacy invasive defaults and consent walls that make forced account creation or tracking consent the path of least resistance.}\pdfcomment[icon=Note]{Marketing‑oriented taxonomies that label tactics like forced action and privacy‑invasive defaults appear in marketing typologies (6087241 p.3) describing forced action and data extraction examples.} \hllightgreen{These designs are interpreted by consumers as signals that the firm prioritises data capture over welfare, increasing first party data and tracking consent in the short term while creating privacy concern and perceived loss of control.}\pdfcomment[icon=Note]{Marketing taxonomy links such designs to consumer interpretation and short‑term data capture / longer‑term privacy concerns (6087241 p.3; 6295677 p.8).} \hllightgreen{In algorithmic environments, such default driven data capture directly feeds recommender models and notification systems, thereby integrating default bias into the technical substrate of engagement optimisation.}\pdfcomment[icon=Note]{Sources describe how captured data feed recommender and notification systems and thereby affect engagement optimisation (0123790 p.2-3; 8958037 p.3), supporting this integration claim.} \hllightgreen{Friction costs appear on the other side of the journey, where exiting or exercising protective choices becomes intentionally difficult.}\pdfcomment[icon=Note]{Friction costs as barriers to exit and protective choices are described as the converse of defaults in dark pattern/sludge literature (6295677 p.4; 6087241 p.3).} \hldarkgreen{Choice obstruction is defined as the deliberate addition of friction, complexity, or delays to actions that reduce firm revenue or data access, with examples such as hard to cancel subscriptions, long unsubscribe flows, and confusing return processes.}\pdfcomment[icon=Note]{Choice obstruction defined as added friction with examples (hard‑to‑cancel subscriptions, long unsubscribe flows, confusing returns) is explicitly listed in marketing typology (6087241 p.3).} \hllightgreen{These tactics increase cognitive fatigue and action costs, lowering immediate churn and reducing completed cancellations while producing reactance, anger, negative word of mouth, and regulatory risk \parencites[4]{6087241}.}\pdfcomment[icon=Note]{Effects of these tactics, higher cognitive fatigue, lower churn, and reputational/regulatory harms, are noted in marketing and dark pattern literature (6087241 p.3; 6295677 p.4).} \hldarkgreen{Systematic evidence from Indian e commerce shows that "roach motel" patterns (easy sign up, difficult cancellation), obstruction, subscription traps, and nagging are highly prevalent.}\pdfcomment[icon=Note]{Systematic evidence of prevalent roach‑motel, obstruction, subscription traps and nagging in Indian e‑commerce is presented in the India‑focused synthesis and table (6295677 p.8, p.12).} \hldarkgreen{Roach motel patterns exploit status quo bias and inertia, with 47\% of surveyed users reporting difficulty cancelling services, and subscription traps combine default enrolment, unclear billing, and complex exits to generate recurring unwanted payments.}\pdfcomment[icon=Note]{The 47\% figure reporting users who had difficulty cancelling and description of subscription traps combining defaults/unclear billing/complex exits appear directly in the India table (6295677 Table p.8).} \hllightgreen{Nagging through repeated prompts, often via app notifications, generates decision fatigue and exerts a wear down effect that pushes users back toward profit maximising options rather than genuinely chosen outcomes \parencites[8]{6295677}.}\pdfcomment[icon=Note]{Nagging via repeated prompts and its wear‑down effect producing decision fatigue is described in the India evidence and dark pattern literature (6295677 p.8; 6087241 p.3).} \hllightgreen{Policy analysis based on dual process cognition clarifies why these defaults and frictions are so effective.}\pdfcomment[icon=Note]{Policy analysis invoking dual‑process cognition to explain effectiveness of defaults and frictions is established across sources linking nudge/sludge and dual‑process theory (6295677 p.4; 6087241 p.3).} \hldarkgreen{Dark patterns are said to systematically target System 1 processing by exploiting heuristics such as default bias, loss aversion, scarcity, and social pressure.}\pdfcomment[icon=Note]{Explicit claim that dark patterns target System 1 heuristics (default bias, loss aversion, scarcity, social pressure) is directly stated in the dark patterns literature (6295677 p.4; p.8).} \hllightgreen{By loading protective actions with sludge like friction and presenting profit aligned options as effortless defaults, platforms keep users in fast, heuristic driven decision making modes and suppress System 2 reflective evaluation.}\pdfcomment[icon=Note]{Argument that loading protective actions with friction keeps users in fast heuristic modes and suppresses reflective evaluation follows directly from sludge/dark pattern theory and dual‑process accounts (6295677 p.4; 6087241 p.3).} \hllightgreen{This reduces capacity for informed and autonomous choice across pricing, consent, and cancellation decisions \parencites[4]{6295677}.}\pdfcomment[icon=Note]{Reduction in capacity for informed/autonomous choice across pricing, consent, and cancellation decisions is a documented outcome in the literature on dark patterns and sludge (6295677 p.4; 6087241 p.3).} \hllightgreen{Regulatory scholarship therefore argues that effective enforcement must move beyond formal obligations such as choice screens or nominal consent banners.}\pdfcomment[icon=Note]{Regulatory scholarship arguing enforcement must go beyond formal choice screens/nominal consent banners is advanced in the regulatory choice architecture and dark pattern critiques (1070605 p.9-10; 6295677 p.12).} \hldarkgreen{Manipulative design practices can neutralise regulatory intent by steering users back to defaults aligned with platform interests through misleading button placements, obfuscating language, or deliberate complication of opt out procedures \parencites[9]{1070605}[12]{6295677}.}\pdfcomment[icon=Note]{The specific examples (misleading button placements, obfuscating language, complicating opt‑out) neutralising regulatory intent are explicitly discussed in Hermansyah's analysis (1070605 p.9) and cited here (6295677 referenced).} \hllightgreen{For an Ethical AI Audit Toolkit, these findings imply that default settings and friction costs are not peripheral usability questions but central indicators of manipulation risk.}\pdfcomment[icon=Note]{Implication for Ethical AI Audit Toolkits that defaults and frictions indicate manipulation risk is a reasonable synthesis of the enforcement and dark pattern literature (6295677 p.12; 1070605 p.9).} \hldarkgreen{Research on dark patterns in India recommends bright line prohibitions against specific patterns, mandatory UX design audits for dominant platforms, and statutory ethical design mandates that require consumer first defaults, privacy by design, and algorithmic transparency \parencites[12]{6295677}.}\pdfcomment[icon=Note]{Research recommendations from Indian dark pattern work, bright line prohibitions, mandatory UX audits for dominant platforms, consumer‑first defaults and algorithmic transparency, are explicitly proposed (6295677 p.12‑13).} \hllightgreen{Competition oriented work on interoperability in digital markets adds that behaviourally informed regulation, such as DMA style choice architecture reforms, can reduce switching frictions and increase contestability, but only when paired with vigilant enforcement against dark patterns that recreate lock in through defaults and sludge.}\pdfcomment[icon=Note]{Competition literature on interoperability and behaviourally informed regulation reducing switching frictions, while requiring enforcement against dark patterns, is argued in Hermansyah (1070605 p.3, p.10) and matches the synthesis here.} \hldarkgreen{Ex ante rules on interoperability and choice are therefore seen as necessary but insufficient without continuous oversight targeted specifically at default bias and friction costs in digital journeys \parencites[152]{1070605}[160]{1070605}.}\pdfcomment[icon=Note]{Statement that ex‑ante interoperability/choice rules are necessary but insufficient without oversight targeting defaults and friction costs is directly concluded in Hermansyah (1070605 p.10; cited here as [160]\{1070605\}).} \hldarkgreen{Empirical analyses of recommender deployment indicate that a large share of user activity and commercial value is driven by algorithmic suggestions, which means that any upstream bias in data collection or consent flows will be amplified downstream by personalised ranking and notification systems \parencites[3]{0123790}.}\pdfcomment[icon=Note]{Empirical analyses showing a large share of activity/commercial value driven by algorithmic suggestions (e.g., Netflix/YouTube figures) are provided in recommender systems literature (0123790 p.3‑5), supporting downstream amplification concerns.} \hllightgreen{This suggests regulators should treat defaults and exit frictions as integral to algorithmic impact assessments, since small design nudges can cascade into measurable changes in engagement, retention, and revenue.}\pdfcomment[icon=Note]{Recommendation that regulators treat defaults and exit frictions as part of algorithmic impact assessments is a reasonable policy inference drawn from evidence on cascade effects in choice architecture and recommender impacts (1070605 p.10; 0123790 p.3).} \hldarkgreen{Users frequently recognise that algorithms shape visibility yet report limited trust in AI-curated content, a gap that may make them more susceptible to simplified, platform-framed defaults presented as convenient choices \parencites{7646767}.}\pdfcomment[icon=Note]{Survey evidence that users recognize algorithmic influence but report limited trust (awareness‑trust gap: \textasciitilde{}78\% see influence, \textasciitilde{}41\% confident) is reported directly in the algorithmic visibility study (0123790 p.4).} \hllightgreen{Such mistrust alongside visible algorithmic influence appears to complicate the normative case for consent-by-design: users see the influence but may still accept low-effort defaults that match the platform's framing rather than pursue reflective alternatives \parencites{7646767}.}\pdfcomment[icon=Note]{Synthesis that visible algorithmic influence complicates consent‑by‑design because users may accept low‑effort defaults despite awareness is a reasonable inference from the trust‑awareness findings and behavioral literature (0123790 p.4; 6295677 p.4).}

\subsubsection{Salience, Scarcity and Loss Aversion in Interface Design}
\label{sec:salience_scarcity_loss_aversion} \hldarkgreen{Interface design in digital retail systematically exploits salience, scarcity, and loss aversion to accelerate decisions and shift routine purchases into urgent, emotionally charged episodes.}\pdfcomment[icon=Note]{The claim is directly supported by the qualitative synthesis and theoretical framing in Source 8745065 (pp.135-136) and Source 6295677 (pp.8), which describe how digital interfaces use scarcity, salience, and behavioral biases to accelerate decisions.} \hldarkgreen{Scarcity messages such as "only two left," countdown timers, expiring discounts, and indicators of recent purchases are singled out as transforming an ordinary opportunity into a perceived time sensitive decision, altering both perceived value and desirability of the product \parencites[135]{8745065}.}\pdfcomment[icon=Note]{Source 8745065 (p.135) explicitly lists messages such as "only two left," countdown timers, expiring discounts, and recent-purchase indicators and explains how they transform opportunities into time-sensitive decisions.} \hldarkgreen{When such cues are highly salient in the visual hierarchy, for example dominating colour, size, or placement compared with neutral information, they function as focal stimuli that draw System 1 attention and suppress slower, reflective appraisal \parencites[8]{6295677}[136]{8745065}.}\pdfcomment[icon=Note]{Source 8745065 (pp.135-136) links visually salient cues to heuristic (System 1) attention within a dual-process framework and explains their role in accelerating decisions, matching the cited reference.} \hldarkgreen{Loss framing interacts tightly with scarcity.}\pdfcomment[icon=Note]{Source 8745065 (p.135) explicitly discusses the interaction between scarcity cues and loss framing, noting their combined influence on perceived urgency and decision-making.} \hldarkgreen{Prospect theory predicts that losses and gains are psychologically asymmetric, with losses weighing more heavily than equivalent gains, and interface copy that frames non purchase as "missing out" on a discount harnesses this asymmetry directly \parencites[135]{8745065}.}\pdfcomment[icon=Note]{Source 8745065 (p.135) situates the argument in prospect theory (Kahneman \& Tversky) and describes how framing non-purchase as "missing out" exploits loss--gain asymmetry.} \hldarkgreen{Interview evidence from digital retail executives describes campaigns where budget conscious consumers are targeted with expiring discount language and "pay more later" narratives, while higher income segments receive exclusivity and rarity messaging.}\pdfcomment[icon=Note]{Source 8745065 (p.138, Theme 2) reports interview accounts that budget-conscious consumers receive expiring-discount/pay-more-later messaging while higher-income segments receive exclusivity/rarity messaging.} \hldarkgreen{The distinction between relevance and vulnerability targeting is explicit: messages are selected because predictable weaknesses such as sensitivity to price or status increase the likelihood of immediate conversion rather than because the product intrinsically matches needs \parencites[138]{8745065}.}\pdfcomment[icon=Note]{Source 8745065 (p.138, Theme 2) explicitly contrasts relevance targeting with vulnerability targeting, describing messages chosen to exploit predictable weaknesses to increase immediate conversion.} \hllightgreen{In this environment, salience is not incidental; visually dominant urgency prompts, pop ups, and notifications are deliberately placed to amplify loss framed scarcity and compress deliberation time.}\pdfcomment[icon=Note]{Source 8745065 (pp.136,139) and Source 6295677 (p.8 Table) together describe visually dominant urgency prompts and friction-reduction tactics that amplify scarcity/loss framing and shorten deliberation, supporting a synthesized claim.} \hldarkgreen{Dark pattern typologies in Indian e commerce operationalise these mechanisms at the design level. "False urgency / scarcity" is defined as misleading countdown timers or stock alerts that create artificial pressure, explicitly linked to the scarcity effect and time pressure, and documented as frequently deployed by platforms audited by the Advertising Standards Council of India. "Interface interference" describes manipulation of visual hierarchy, colour, and placement to privilege platform preferred choices, exploiting salience bias and visual dominance, and appearing on approximately $70\%$ of audited platforms.}\pdfcomment[icon=Note]{Source 6295677 (p.8 Table) defines "False urgency / scarcity" (misleading timers/stock alerts) and "Interface interference" (manipulation of visual hierarchy) and records interface interference at \textasciitilde{}70\% of audited platforms and ASCI documentation of false urgency.} \hldarkgreen{Nagging, in the form of repeated prompts after refusal, generates a wear down effect and decision fatigue, especially when delivered through app notifications.}\pdfcomment[icon=Note]{Source 6295677 (p.8 Table) lists "Nagging" as repeated prompts after refusal producing a wear-down effect and decision fatigue, commonly via app notifications.} \hldarkgreen{Subscription traps and drip pricing combine loss aversion and inertia with incremental disclosure to create recurring payments and surprise charges \parencites[8]{6295677}.}\pdfcomment[icon=Note]{Source 6295677 (p.8 Table) documents drip pricing (incremental disclosure) and subscription traps (inertia/loss aversion) as mechanisms producing recurring payments and surprise charges.} \hllightgreen{These patterns show how salience is engineered around scarcity and loss cues instead of neutral information about total cost or cancellation rights.}\pdfcomment[icon=Note]{Source 8745065 (pp.135-136,139) and Source 6295677 (pp.6,8) together show that interfaces often foreground scarcity and loss cues rather than neutral cost or cancellation information, supporting this synthesis.} \hldarkgreen{Organisational interviews on urgency marketing extend this design picture into governance.}\pdfcomment[icon=Note]{Source 8745065 (pp.135-139) explicitly states that organisational interviews extend the interface-level findings into governance and internal decision processes.} \hldarkgreen{Executives report that high pressure cues such as countdowns and stock alerts are routinely deployed as operational responses to competitive pressure and conversion targets rather than isolated creative experiments.}\pdfcomment[icon=Note]{Source 8745065 (p.138, Theme 1) reports executives describing countdowns and stock alerts as routine operational responses to competitive pressure and conversion targets rather than isolated creative experiments.} \hldarkgreen{The analysis cautions that reported conversion gains are interviewee accounts, yet aligns them with established behavioral economics showing that framing, pre commitment, and default escalation can dramatically change savings or purchase behaviour \parencites[135]{8745065}.}\pdfcomment[icon=Note]{Source 8745065 (p.138) cautions that conversion gains are interviewee accounts and aligns these with behavioral-economics findings (e.g., framing, pre-commitment, default escalation) discussed elsewhere in the paper (see p.135 and p.136 referencing Thaler \& Benartzi).} \hldarkgreen{The same study synthesises a "cognitive trigger layer" in algorithmic advertising: scarcity cues, loss framed notifications, and social proof increase perceived importance of decisions; an "algorithmic personalization layer" tailors timing and intensity based on behavioural histories; and a "temporal and friction reduction layer" compresses steps to checkout.}\pdfcomment[icon=Note]{Source 8745065 (p.139) explicitly synthesises the three-layer model: cognitive trigger layer, algorithmic personalization layer, and temporal/friction-reduction layer.} \hldarkgreen{When these layers converge, exploitative impulse buying is better understood as an architectural phenomenon than as single messages \parencites[139]{8745065}.}\pdfcomment[icon=Note]{Source 8745065 (p.139) states that when these layers converge exploitative impulse buying is better understood as an architectural phenomenon rather than as single messages.} \hldarkgreen{Regulatory analysis under the EU Digital Services Act adds a risk framing for manipulative use of these mechanisms.}\pdfcomment[icon=Note]{Sources 6941611 (p.18) and 7896734 (pp.390 onward) analyse the DSA and frame it as addressing risks from manipulative behavioral advertising, supporting the regulatory risk framing claim.} \hldarkgreen{Recital 69 warns that advertising based on targeting techniques optimised to align with users' interests and exploit vulnerabilities can yield particularly severe adverse effects.}\pdfcomment[icon=Note]{Source 6941611 (p.18) directly quotes and summarizes Recital 69, warning that ads optimized to exploit users' vulnerabilities can yield severe adverse effects.} \hldarkgreen{Article 26 requires that ads disclose their nature, identify the financed entity, and explain the primary targeting parameters together with instructions on how users can modify them.}\pdfcomment[icon=Note]{Source 6941611 (p.18) describes Article 26's obligations requiring ads to disclose their nature, identify the financed entity, and explain the primary targeting parameters and how users can modify them.} \hllightgreen{Sensitive data based profiling is categorically banned and notifications deployed with obstructive intent or to impede competition are highlighted as abusive practices that can generate social and economic harms, including mental wellbeing impacts and violations of consumer protection and non discrimination \parencites[18]{6941611}.}\pdfcomment[icon=Note]{Source 6941611 (p.18) documents prohibitions on advertising based on sensitive-data profiling and flags notifications with obstructive intent as abusive; Source 7896734 (p.390) and related discussion link these practices to social/economic and wellbeing harms, supporting a combined inference.} \hllightgreen{These provisions, combined with national dark pattern guidelines that prohibit misleading urgency and mandate explicit consent rather than pre checked defaults, form normative anchors for assessing salience driven scarcity and loss aversion tactics as potential manipulation rather than ordinary persuasion \parencites[12]{6295677}[390]{7896734}.}\pdfcomment[icon=Note]{Source 6941611 (p.18) on DSA provisions together with Source 6295677 (pp.12-13) (CCPA advisory and guideline references) and Source 7896734 support the claim that these rules plus national dark-pattern guidelines (explicit consent vs pre-checked defaults) form normative anchors.} \hllightgreen{Recent work on recommender systems suggests an additional ethical dimension: personalization does not act in isolation but shapes the distribution of cultural content and the kinds of choices presented at scale, so decisions about which behaviours to optimize for are inherently normative and may reinforce the very vulnerabilities that urgency cues exploit \parencites[3]{5403359}.}\pdfcomment[icon=Note]{Source 5403359 (pp.2-3) and Source 6249388 (p.23) discuss recommender systems shaping cultural content and normative choices at scale, supporting the ethical claim that personalization can reinforce vulnerabilities and involves normative decisions.}

\subsubsection{Temporal Discounting and Instant Gratification Online}
\label{sec:temporal_discounting_instant_gratification} \hllightgreen{Short video platforms exemplify how instant rewards restructure time preferences.}\pdfcomment[icon=Note]{Li and related literature describe short-video platforms as delivering rapid, low-effort rewards that shift attention and engagement patterns (Li 8958037 p.3; 3303310 p.2), so this is a reasonable synthesis of those sources.} \hldarkgreen{Li describes micro duration clips of 15--60 seconds, combined with infinite swipe navigation, that enable users to extend viewing with almost no effort, producing average daily usage on Douyin of nearly 60 minutes, with more than 21 sessions and over half of users exceeding 60 minutes per day \parencites[3]{8958037}.}\pdfcomment[icon=Note]{Dan Li explicitly describes 15--60 second clips, infinite swipe navigation and reports Douyin metrics (≈60 minutes/day, >21 sessions, 56.9\% over 60 min) (Li 8958037 p.3).} \hllightgreen{This extremely low operational cost of "just one more swipe" aligns closely with temporal discounting: immediate, low effort entertainment is repeatedly chosen over deferred, potentially more beneficial activities such as study, offline social interaction, or sleep.}\pdfcomment[icon=Note]{Li documents the very low operational cost of an extra swipe and links this to repeated selection of short videos; linking that mechanism to temporal discounting is a reasonable inference from Li's psychological account (Li 8958037 p.3--4).} \hldarkgreen{Among adolescents and college students, whose psychological development and self control remain immature, these instant gratifications are associated with diminished academic concentration, weakened social skills, and heightened anxiety and depression, culminating in behavioural addiction \parencites[2]{8958037}[3]{8958037}.}\pdfcomment[icon=Note]{Li reports that adolescents and college students, with immature self-control, experience diminished academic concentration, weakened social skills, increased anxiety/depression and risk of behavioural addiction from excessive short-video use (Li 8958037 p.2--3).} \hldarkgreen{Neuropsychological mechanisms described by Li reinforce this temporal skew.}\pdfcomment[icon=Note]{Li describes neuropsychological mechanisms (dopamine-driven reward circuitry) that support the temporal skew toward instant rewards (Li 8958037 p.3--4).} \hldarkgreen{Dopamine surges are reported not primarily at reward receipt but at the moment of expectancy before scrolling to the next video, entrapping users in a cycle of "pursuit $\rightarrow$ transient satisfaction $\rightarrow$ further pursuit" \parencites[3]{8958037}.}\pdfcomment[icon=Note]{Li explicitly states that dopamine surges occur during the anticipation moment before the next video, producing a 'pursuit → transient satisfaction → further pursuit' cycle (Li 8958037 p.3).} \hldarkgreen{The content switching mechanism is characterised as structurally isomorphic to slot machine variable rewards, where each downward swipe carries an unpredictable payoff and maximally activates mesolimbic dopaminergic pathways.}\pdfcomment[icon=Note]{Li characterises the content-switching mechanism as structurally isomorphic to slot-machine variable rewards and notes maximal activation of mesolimbic dopaminergic pathways (Li 8958037 p.4).} \hldarkgreen{Chronic stimulation leads to downregulation of dopamine receptors and elevated reward thresholds, diminishing the hedonic value of "natural" long term rewards such as learning or physical exercise and creating a vicious cycle in which increased short video use breeds dissatisfaction that in turn drives further use.}\pdfcomment[icon=Note]{Li reports that chronic stimulation leads to dopamine receptor downregulation, elevated reward thresholds and reduced hedonic value of natural rewards, creating a vicious cycle (Li 8958037 p.4).} \hldarkgreen{Algorithmic recommendation systems are explicitly designed to build sustained, automated usage habits and to bypass deliberative decision making processes, inducing impulsive, intuition driven choices that mirror the preference for immediate gratification observed in temporal discounting experiments \parencites[4]{8958037}.}\pdfcomment[icon=Note]{Li and related work describe recommendation systems as designed to create sustained, automated habits that bypass deliberative processes and engage impulsive/intuition-driven responses (Li 8958037 p.4; 3303310 p.2).} \hllightgreen{Hu's synthesis of "algorithm driven engagement worship" links these micro level reward dynamics to broader mental health outcomes.}\pdfcomment[icon=Note]{The phrase 'algorithm-driven engagement worship' and its linking of micro reward dynamics to mental health outcomes appears in the literature reviewed (3303310 p.2--3), so this is a fair synthesis referencing that body of work.} \hldarkgreen{Infinite scrolling interfaces and variable reward schedules on platforms such as TikTok are described as exploiting neural pathways associated with addiction, with empirical validations showing 22\% higher rates of compulsive checking in algorithmically curated feeds than in non algorithmic interfaces \parencites[2]{3303310}.}\pdfcomment[icon=Note]{The conference paper describes infinite scrolling/variable rewards exploiting addiction-related neural pathways and cites empirical validation of 22\% higher compulsive checking in algorithmically curated feeds (3303310 p.2; 3303310 p.3).} \hldarkgreen{Additional evidence is cited that hyper personalised feeds are associated with a 23\% higher incidence of depressive symptoms compared to chronological content, and neuroimaging studies correlate "doomscrolling" with reduced prefrontal cortex activity and reduced gray matter density in regions governing impulse control, changes that resemble substance dependence \parencites[20]{3303310}.}\pdfcomment[icon=Note]{3303310 reports a 23\% higher incidence of depressive symptoms with hyper-personalized feeds and links 'doomscrolling' to reduced prefrontal activity and reduced gray matter density in impulse-control regions (3303310 p.3).} \hllightgreen{These findings indicate that repeated selection of instant digital rewards over temporally distant benefits is not merely a behavioural preference but is reinforced and stabilized by neurocognitive adaptation.}\pdfcomment[icon=Note]{Li and related neuroimaging/behavioral studies document neurocognitive adaptation (dopamine changes, receptor downregulation) that can stabilise preference for instant rewards, so this synthesis is supported (Li 8958037 p.4; 3303310 p.2--3).} \hldarkgreen{Research on recommendation algorithms has also documented systematic popularity and exposure biases in feed and venue suggestions, and it appears that combining models or applying re-ranking can moderate these polarizations while keeping acceptable accuracy levels \parencites[3]{9051044}.}\pdfcomment[icon=Note]{Recommender research documents popularity/exposure biases and shows combining models or using re-ranking can mitigate polarization while retaining acceptable accuracy (Bias mitigation literature 9051044 p.33--38).} \hllightgreen{Such work suggests that algorithmic design choices may be able to reduce the structural skew toward immediately gratifying content without fully sacrificing relevance \parencites[3]{9051044}.}\pdfcomment[icon=Note]{9051044 presents results indicating algorithmic design choices (e.g., hybrid or re-ranker methods) can reduce skew toward popular items without dramatically sacrificing relevance, supporting this inference (9051044 p.38).} \hldarkgreen{Temporal orientation work within the I PACE framework provides a complementary behavioural lens.}\pdfcomment[icon=Note]{Li frames temporal orientation within the I-PACE model as a behavioural lens for understanding short-video addiction (Li 8958037 p.4).} \hldarkgreen{Li reports a large sample study of 2\,239 participants in which past and present oriented temporal foci were significantly positively correlated with short video addiction, while future oriented focus was negatively correlated.}\pdfcomment[icon=Note]{Li reports a large-sample study (N = 2,239) finding past/present orientation positively correlated with short-video addiction and future orientation negatively correlated (Li 8958037 p.4).} \hldarkgreen{Inhibitory self control mediated the link between past or present focus and addiction, with attempts at suppression consuming cognitive resources and producing a "resistance backfire effect," whereas initiatory self control mediated the protective effect of future orientation by enabling individuals with clear future goals to initiate alternative beneficial behaviours that replace habitual scrolling \parencites[4]{8958037}.}\pdfcomment[icon=Note]{Li reports inhibitory self-control mediates the link between past/present focus and addiction (resistance backfire effect) while initiatory self-control mediates the protective effect of future orientation (Li 8958037 p.4).} \hllightgreen{These results can be read as an explicit articulation of temporal discounting inside algorithmic environments: users anchored in immediate or retrospective frames, and relying on inhibition rather than proactive substitution, are more likely to overvalue instant gratification delivered by recommender loops.}\pdfcomment[icon=Note]{This is a reasonable interpretation combining Li's I-PACE findings about temporal focus and self-control with the observed recommender-driven instant rewards (Li 8958037 p.4).} \hllightgreen{Regulatory analysis under the Digital Services Act connects these dynamics to concerns about manipulation.}\pdfcomment[icon=Note]{DSA analysis ties platform design and manipulative interfaces to regulatory concerns about user manipulation, which connects to the described dynamics (6941611 p.20).} \hldarkgreen{Quinelato explains that Article 25 DSA prohibits online interfaces that deceive or manipulate users or substantially impair their ability to make free and informed decisions, explicitly naming techniques such as giving undue prominence to specific options, requesting repeated choices, or making cancellation harder than subscription \parencites[20]{6941611}.}\pdfcomment[icon=Note]{6941611 explicitly describes Article 25 prohibiting interfaces that deceive/manipulate users and lists examples such as undue prominence, repeated choices, and making cancellation harder (6941611 p.20).} \hllightgreen{In combination with Li's account of "addictive design" in infinite scrolling, auto play, and push notifications already under DSA investigation for impairing free decision making on TikTok, these provisions frame environments that constantly offer immediate, highly salient rewards while obstructing exit or alternative choices as manipulative architectures that exploit temporal discounting \parencites[2]{8958037}[20]{6941611}.}\pdfcomment[icon=Note]{Li documents concerns about infinite scrolling, autoplay and notifications and EU investigations into TikTok (Li 8958037 p.2), and 6941611 describes Article 25's scope, so combining them to frame manipulative architectures is supported (8958037 p.2; 6941611 p.20).} \hllightgreen{Dark pattern research further shows that default choices and obstructive frictions can compress deliberation time around purchases and subscriptions by making impulsive acceptance effortless and introducing sludge into protective actions such as cancellation or refusal, thereby amplifying preferences for instant gratification over longer term welfare \parencites[6]{6295677}[18]{6941611}.}\pdfcomment[icon=Note]{Dark-pattern and behavioral-interface literature describe defaults and friction as compressing deliberation and amplifying impulse-based choices; this is discussed in qualitative studies (8745065 p.6; 6941611 p.7).} \hllightgreen{Bansal's synthesis characterises this as a "temporal and friction reduction layer" where reduced steps and urgency cues compress the time available for reflective evaluation and push consumers toward immediate conversion \parencites[6]{8745065}.}\pdfcomment[icon=Note]{The characterization of a 'temporal and friction-reduction layer' appears in the dark-pattern synthesis (8745065 p.6), so citing Bansal's synthesis in these terms is a reasonable alignment with the provided source.} \hllightgreen{Educational and policy responses emphasise algorithmic literacy and AI education as counterweights.}\pdfcomment[icon=Note]{Policy and educational responses emphasizing algorithmic literacy and AI education are recommended in governance literature and reports (2977039 p.99; 9360054 p.10), supporting this claim.} \hllightgreen{Kostenko argues that public educational platforms and AI literacy materials should build critical skills for evaluating algorithms and their results, enabling users to understand how recommendation systems and addictive design strategies influence their time allocation and reward expectations \parencites[99]{2977039}.}\pdfcomment[icon=Note]{2977039 discusses public educational platforms and algorithmic literacy materials aimed at building critical evaluation skills for algorithms, matching the substance of this claim (2977039 p.99).} \hldarkgreen{Surveys on digital literacy and EU regulatory awareness nonetheless reveal that while users feel confident in everyday digital tasks, they score low on recognising hidden manipulations, indicating that many remain vulnerable to interface designs that trade on instant gratification and temporal discounting dynamics \parencites[31]{9360054}.}\pdfcomment[icon=Note]{Survey findings show high confidence in everyday digital tasks but low scores on recognising hidden manipulations and moderate awareness of EU regulations, directly supporting this statement (9360054 p.10).}

\subsection{Nudging, Dark Patterns and Ethical Boundaries}

\subsubsection{From Helpful Nudges to Manipulative Dark Patterns}
\label{sec:helpful_nudges_to_dark_patterns} \hldarkgreen{Choice architecture in digital interfaces ranges from helpful guidance to manipulative dark commercial patterns that erode autonomy and trust.}\pdfcomment[icon=Note]{Sources 6087241 (p.3) and 6295677 (p.2) explicitly describe choice architecture ranging from benign guidance to manipulative dark commercial patterns that undermine autonomy and trust.} \hldarkgreen{Online choice architecture and digital nudging describe how rankings, default options, button placement, and recommendation systems structure the environment in which users act, without implying that all influence is harmful.}\pdfcomment[icon=Note]{Source 6087241 (p.3) defines online choice architecture and digital nudging (rankings, defaults, button placement, recommender systems) and notes that design influence does not automatically make an interface harmful.} \hldarkgreen{Ethical nudges can highlight relevant options or simplify complex flows while preserving the user's capacity for informed, preference consistent choice.}\pdfcomment[icon=Note]{Source 6087241 (p.3) and 6295677 (p.4) discuss ethical nudges that highlight relevant options and aim to preserve informed, preference‑consistent choice.} \hldarkgreen{Dark commercial patterns mark the point where this influence becomes manipulative: designs that steer, deceive, pressure, or obstruct consumers so that comprehension, reversibility, or autonomy are weakened.}\pdfcomment[icon=Note]{Source 6087241 (p.3) defines dark commercial patterns as designs that steer, deceive, pressure, or obstruct consumers and thereby weaken comprehension, reversibility, or autonomy.} \hldarkgreen{A recommender that hides lower priced alternatives, obscures sponsored placement, or makes refusal difficult crosses that boundary because it secures compliance by undermining the conditions of legitimate persuasion rather than by clarifying options \parencites[32--33]{6087241}.}\pdfcomment[icon=Note]{Source 6087241 (p.3) gives the specific example that a recommender hiding lower‑priced alternatives, obscuring sponsorship, or making refusal difficult undermines legitimate persuasion.} \hldarkgreen{Dark pattern research traces this shift back to systematic exploitation of cognitive biases.}\pdfcomment[icon=Note]{Source 6295677 (p.2, p.4) and 6087241 (p.3) trace dark patterns to systematic exploitation of cognitive biases in the literature.} \hldarkgreen{Building on Kahneman's dual process theory, dark patterns are described as deliberately triggering fast, intuitive System 1 processing through time pressure, default selections, visual salience, and emotional cues while suppressing slower, reflective System 2 reasoning.}\pdfcomment[icon=Note]{Source 6295677 (p.2) and 6295677 (p.4) explicitly connect Kahneman's dual‑process theory to dark patterns that trigger System 1 via time pressure, defaults, visual salience and emotional cues while suppressing System 2.} \hldarkgreen{Practices such as drip pricing, false urgency, and confirm shaming are characterised as "exploitative nudges," whereas obstruction, subscription traps, and forced action are classified as "sludge" because they introduce artificial friction that discourages consumer protective behaviour \parencites[2]{6295677}[4]{6295677}.}\pdfcomment[icon=Note]{Source 6295677 (p.14) explicitly characterises drip pricing, false urgency, and confirm‑shaming as exploitative nudges and obstruction, subscription traps, and forced action as sludge.} \hllightgreen{In this framing, the same behavioural tools that make welfare enhancing nudges possible can be reoriented toward coercion when interface design is optimised for conversion and data extraction rather than user welfare \parencites[2]{6295677}[4]{6295677}[34]{6087241}.}\pdfcomment[icon=Note]{Sources 6087241 (p.33) and 6295677 (p.14) together support the synthesis that behavioural tools enabling welfare‑enhancing nudges can be repurposed toward coercion when interfaces are optimized for conversion and data extraction rather than user welfare.} \hldarkgreen{Marketing oriented work adds a relational perspective by linking these patterns to a conversion trust trade off.}\pdfcomment[icon=Note]{Source 6087241 (p.33) frames a marketing‑oriented, relational perspective linking dark patterns to a conversion--trust tradeoff.} \hldarkgreen{Large scale crawls of shopping sites identify thousands of instances of dark patterns, including hidden fees, difficult cancellation, misleading scarcity claims, and data sharing tricks.}\pdfcomment[icon=Note]{Source 6087241 (p.33) reports a large crawl (\textasciitilde{}11,000 shopping sites identifying 1,818 instances) and regulatory reports documenting hidden fees, difficult cancellation, misleading scarcity, and data‑sharing tricks.} \hldarkgreen{Consent interfaces often structure cookie choices asymmetrically so that accepting tracking is easier than refusal, even when formal choices exist.}\pdfcomment[icon=Note]{Source 6087241 (p.33) and related consent‑interface literature cited there state that cookie/consent notices often structure choices asymmetrically to make acceptance easier than refusal.} \hllightgreen{Conceptual analysis therefore proposes that dark commercial patterns may increase immediate compliance but risk eroding perceived fairness, autonomy, and privacy control, which in turn damages trust, relationship quality, and customer equity over time \parencites[33]{6087241}.}\pdfcomment[icon=Note]{Sources 6087241 (p.33) and 6295677 (p.14) support the conceptual claim that dark commercial patterns can boost immediate compliance while eroding perceived fairness, autonomy, and privacy control, harming trust and long‑term relationship quality.} \hldarkgreen{Ethical digital marketing is consequently defined as persuasion that preserves informed, preference consistent choice, while dark patterns are practices that violate this second condition by distorting salience, defaults, or friction costs in ways that systematically favour firm goals over consumer interests \parencites[32--33]{6087241}.}\pdfcomment[icon=Note]{Source 6087241 (p.3) explicitly defines ethical digital marketing as preserving informed, preference‑consistent choice and describes dark commercial patterns as violating this by distorting salience, defaults, or friction.} \hldarkgreen{Empirical and regulatory evidence from India illustrates how helpful nudging and dark patterns can coexist in the same ecosystem and why enforcement must focus on design.}\pdfcomment[icon=Note]{Source 6295677 (p.14) and 6295677 (p.3) provide empirical and regulatory evidence from India showing coexistence of helpful nudging and dark patterns and argue enforcement should focus on design.} \hldarkgreen{Surveys report that approximately 87\% of major platforms deploy at least one dark pattern, with interface interference present on around 70\% and obstruction based designs on 63\%.}\pdfcomment[icon=Note]{Source 6295677 (p.2 and p.4) reports the statistics that \textasciitilde{}87\% of major platforms deploy at least one dark pattern, with interface interference at \textasciitilde{}70\% and obstruction‑based designs at \textasciitilde{}63\%.} \hldarkgreen{At the same time, only 39\% of users recognise the term "dark patterns," although 68\% report experiencing deceptive practices, 55\% unexpected charges, and 47\% difficulty cancelling subscriptions or services.}\pdfcomment[icon=Note]{Source 6295677 (p.2) provides the consumer figures: 39\% familiarity with the term 'dark patterns', 68\% reporting deceptive practices, 55\% unexpected charges, and 47\% difficulty cancelling.} \hllightgreen{These figures reveal a structural power imbalance: interfaces systematically exploit default bias and inertia, while informational asymmetry and low digital literacy prevent users from recognising manipulation \parencites[2]{6295677}.}\pdfcomment[icon=Note]{Sources 6295677 (p.2--3) and related literature (6087241 p.33) support the inference that these figures indicate structural power imbalances where default bias, inertia, informational asymmetry and low digital literacy impede recognition of manipulation.} \hldarkgreen{Regulatory bodies have started to intervene at the design level.}\pdfcomment[icon=Note]{Sources 6295677 (p.3--4) and 6087241 (p.33) document that regulatory bodies have begun intervening at the design level in response to dark patterns.} \hldarkgreen{The Advertising Standards Council of India has defined thirteen recurring dark pattern categories and issued guidelines on deceptive digital design, while the Central Consumer Protection Authority's guidelines classify manipulative practices as unfair trade practices.}\pdfcomment[icon=Note]{Source 6295677 (p.4) cites ASCI's identification of thirteen dark pattern categories and guidelines, and 6295677 (p.3) notes the Central Consumer Protection Authority classifies manipulative practices as unfair trade practices.} \hldarkgreen{Judicial decisions, such as the Hyderabad commission's ruling against concealed convenience fees on BookMyShow, confirm that interface level deception falls within consumer protection law.}\pdfcomment[icon=Note]{Source 6295677 (p.3) and District Consumer Disputes Redressal Commission references therein report the Hyderabad commission ruling against concealed convenience fees on BookMyShow, confirming interface deception is actionable under consumer protection law.} \hldarkgreen{Sector regulators like IRDAI have explicitly banned pre checked travel insurance options, recognising that such defaults create uninformed consent and financial leakage \parencites[3]{6295677}.}\pdfcomment[icon=Note]{Source 6295677 (p.3) specifically notes that IRDAI prohibited pre‑checked travel insurance options as a measure to curb deceptive defaults.} \hldarkgreen{A critical challenge is that formal regulatory nudges, such as mandated choice screens or consent banners, can be neutralised by counter nudges engineered by platforms.}\pdfcomment[icon=Note]{Source 1070605 (p.159) and 6087241 (p.33) discuss how mandated choice screens or consent banners can be neutralised by counter‑nudges engineered by platforms.} \hldarkgreen{Competition oriented analysis documents how dominant firms deploy misleading button placements, obfuscating language, and complicated opt out procedures that steer users back to default settings aligned with platform interests.}\pdfcomment[icon=Note]{Source 1070605 (p.159) and 6087241 (p.33) document dominant firms deploying misleading button placements, obfuscating language, and complicated opt‑out procedures that steer users to platform‑aligned defaults.} \hllightgreen{These manipulative tactics preserve nominal autonomy while undermining the practical ability to exercise it.}\pdfcomment[icon=Note]{Source 6087241 (p.3) and related analysis support the reasonable synthesis that such manipulative tactics preserve nominal autonomy while undermining the practical ability to exercise it.} \hllightgreen{Regulatory choice architecture that aims to break default bias and reduce switching frictions must therefore be paired with continuous oversight of interface asymmetries and sludge based frictions to prevent a slide from helpful nudges into covert coercion \parencites[159]{1070605}.}\pdfcomment[icon=Note]{Sources 1070605 (p.159) and 6295677 (p.14) recommend pairing regulatory choice architecture with oversight of interface asymmetries and sludge, which supports this policy prescription as a synthesis.} \hldarkgreen{For the Ethical AI Audit Toolkit developed later in this report, this body of work motivates a clear evaluative distinction: interface elements count as acceptable nudges when they simplify decisions without degrading comprehension, reversibility, or independent preference formation, and they fall into this report's "dark pattern" risk categories when they rely on information distortion, interface asymmetry, social or temporal pressure, choice obstruction, or forced action that materially impairs informed and autonomous choice \parencites[32--33]{6087241}[2]{6295677}[4]{6295677}.}\pdfcomment[icon=Note]{Sources 6087241 (p.33) and 6295677 (p.4, p.14) explicitly motivate the evaluative distinction used in the toolkit and enumerate risk categories (information distortion, interface asymmetry, social/temporal pressure, obstruction, forced action) tied to impairing informed autonomous choice.}

\subsubsection{Common Dark Interface Strategies in Social Media}
\label{sec:common_dark_interface_social} \hllightgreen{Dark interface strategies on social platforms frequently operationalise patterns already documented in e commerce, but embed them in feeds, consent flows, and notification settings that are tightly coupled to recommender engines.}\pdfcomment[icon=Note]{Sources discuss dark commercial patterns transferring across contexts and being embedded in feeds, consent flows and recommender-driven notifications (6087241; 3303310), so this is a reasonable synthesis from the provided literature.} \hldarkgreen{Content analysis of major Indian platforms shows that interface interference and obstruction based designs appear on approximately $70\%$ and $63\%$ of audited services respectively, signalling a systematic reliance on choice architecture manipulation rather than overt misinformation \parencites[7]{6295677}[9]{6295677}.}\pdfcomment[icon=Note]{The LocalCircles / IJMRP content analysis explicitly reports interface interference ≈70\% and obstruction ≈63\% of audited services (6295677, p.9).} \hldarkgreen{These tactics include visually privileging engagement maximising actions, hiding or de emphasising opt out controls, and stretching cancellation journeys so that users experience high cognitive effort when attempting to reduce data sharing or disengage from a service \parencites[7]{6295677}[9]{1070605}.}\pdfcomment[icon=Note]{The IJMRP synthesis and audit findings list privileging engagement actions, hiding/de-emphasising opt-out controls and prolonged cancellation journeys as common tactics (6295677, p.9 and Table 2).} \hldarkgreen{Empirical survey data from Indian consumers illustrates how such strategies manifest as experienced harm.}\pdfcomment[icon=Note]{The cited national audit and institutional surveys document consumer-level prevalence and harms from these strategies, linking interface tactics to experienced consumer harm (6295677, pp.9, Table 2).} \hldarkgreen{Drip pricing, bait and switch, and privacy zuckering affect $75\%$, $48\%$, and $44\%$ of consumers respectively, while forced action and basket sneaking demonstrate how platforms use procedural dominance to override user intent, for instance through auto added services or orders processed despite cancellation \parencites[9]{6295677}.}\pdfcomment[icon=Note]{Table 2 in the LocalCircles summary reports drip pricing 75\%, bait-and-switch 48\%, privacy zuckering 44\%, and discusses forced action and basket sneaking as procedural dominance examples (6295677, pp.9--10, Table 2).} \hllightgreen{Although this audit focus is on e commerce, the same design motifs appear in social media contexts: ambiguous consent dialogs, pre selected data sharing defaults, and auto enrolment into personalised feeds shift control over data and attention without explicit, informed action by users \parencites[7]{6295677}[3]{6087241}.}\pdfcomment[icon=Note]{Conceptual work and empirical studies link similar design motifs (ambiguous consent dialogs, pre-selected defaults, auto-enrolment into personalised feeds) from e-commerce into broader platform contexts (6087241; 3303310), supporting this synthesis.} \hldarkgreen{Dark commercial pattern scholarship characterises these practices as steering, pressuring, or obstructing users in ways that impair informed, autonomous choice.}\pdfcomment[icon=Note]{The marketing-oriented taxonomy explicitly characterises dark commercial patterns as steering, pressuring or obstructing users in ways that impair informed, autonomous choice (6087241, pp.3--4).} \hldarkgreen{Examples include consent banners where the "accept all" option is visually dominant and refusal requires multi step navigation, or feeds where lower cost or lower engagement alternatives are hidden or framed as abnormal choices.}\pdfcomment[icon=Note]{Both the audit literature and conceptual reviews provide these concrete examples (dominant 'accept all' consent buttons and multi-step refusal; hiding lower-cost alternatives) as prototypical dark patterns (6295677; 6087241).} \hldarkgreen{Such interfaces are evaluated along two dimensions: contribution to firm objectives such as data capture or retention, and preservation of preference consistent choice.}\pdfcomment[icon=Note]{6087241 frames evaluation along firm-objective contribution (e.g., data capture, retention) and preservation of preference-consistent choice, matching this two-dimension claim (6087241, pp.3--4).} \hldarkgreen{Designs that secure compliance by weakening comprehension, transparency, or reversibility are judged as manipulative \parencites[3]{6087241}.}\pdfcomment[icon=Note]{The conceptual literature defines manipulative designs as those securing compliance by weakening comprehension, transparency or reversibility (6087241, p.4).} \hldarkgreen{Research on competition and regulation adds that, even when regulators introduce pro user choice screens or interoperability mandates, dominant platforms may respond with "dark patterns" that neutralise these interventions through defaults and friction, for example by making switching away from a default social login or profile based feed substantially harder in practice \parencites[3]{1070605}[10]{1070605}.}\pdfcomment[icon=Note]{Hermansyah discusses how defaults and switching costs persist and platforms may deploy manipulative design to undermine regulatory choice screens or interoperability, explicitly noting adaptive dark patterns (1070605, pp.3,10).} \hllightgreen{AI enabled social environments introduce additional escalation patterns.}\pdfcomment[icon=Note]{Recent analyses identify AI-enabled escalations of dark patterns (dynamic, personalised manipulation) and so asserting additional escalation patterns in AI contexts is a supported synthesis (7896734; 6295677).} \hldarkgreen{Sanikidze documents "adaptive UI \& friction shifting," where reinforcement learning dynamically adds or removes interface steps to steer or trap users, effectively creating "roach motel" journeys in which account deletion or cancellation is hidden behind changing, asymmetric friction.}\pdfcomment[icon=Note]{The study of AI dark patterns documents 'Adaptive UI \& friction shifting' where reinforcement learning changes interface friction to trap users, creating 'roach motel' cancellation journeys (7896734, p.11).} \hldarkgreen{The same analysis lists "chatbot coercion \& LLM dark patterns," where conversational agents employ sycophantic tone, anthropomorphism, fabricated consensus, and relentless retention strategies during chat based interactions.}\pdfcomment[icon=Note]{The same source enumerates 'chatbot coercion \& LLM dark patterns' describing sycophancy, anthropomorphism, fabricated consensus and retention strategies in conversational agents (7896734, p.11).} \hldarkgreen{These agents can present opt out or cancellation as socially deviant, or repeatedly propose alternative retention offers, thereby defeating informed cancellation and consent.}\pdfcomment[icon=Note]{7896734 explicitly states that conversational agents can defeat informed cancellation/consent by presenting opt-out as socially deviant and repeatedly proposing retention offers (7896734, p.11).} \hllightgreen{In social media settings, similar conversational layers increasingly mediate privacy settings or recommendation controls, which raises the risk that dialogue design itself becomes an engagement trap.}\pdfcomment[icon=Note]{Combining evidence on conversational agents and mediating controls in platform interfaces supports the inference that dialogue design can mediate privacy/recommendation controls and become an engagement trap (7896734; 3303310).} \hldarkgreen{Generative social proof and confirm shaming extend these dynamics to social cues.}\pdfcomment[icon=Note]{7896734 names 'Generative Social Proof' and 'ConfirmShaming' as AI-extended dark pattern categories that transfer manipulation to social cues (7896734, p.11).} \hldarkgreen{Large language models can mass produce persuasive but indistinguishable fake reviews and testimonials, which then appear as authentic peer endorsements and mislead users into perceiving low quality or manipulative services as widely trusted.}\pdfcomment[icon=Note]{The AI dark patterns analysis notes that LLMs can mass-produce indistinguishable fake reviews/testimonials that function as deceptive peer endorsements (7896734, p.11).} \hldarkgreen{Confirm shaming techniques personalise guilt inducing opt out language, using demographic and historical data so that declining a prompt is framed as irresponsible or ungrateful, pushing users toward continued engagement or data sharing despite initial reluctance \parencites[11]{7896734}.}\pdfcomment[icon=Note]{7896734 describes ConfirmShaming as generative models personalising guilt-inducing opt-out language by demographic and history to coerce retention (7896734, p.11).} \hllightgreen{On social networks, analogous mechanisms may appear as fabricated "friends liked this" indicators or emotionally charged prompts that shame users into keeping notifications or personalised feeds enabled.}\pdfcomment[icon=Note]{The generative social proof concept supports the plausible extension to fabricated 'friends liked this' indicators and emotionally charged prompts on social networks, a reasonable synthesis from the listed AI-enabled social-proof tactics (7896734; 3303310).} \hldarkgreen{Evidence from national audits shows that such patterns are not marginal.}\pdfcomment[icon=Note]{National audits and syntheses reported in the IJMRP piece indicate dark patterns are widespread rather than marginal, describing systemic prevalence (6295677, pp.7--9).} \hldarkgreen{An estimated $87\%$ of major Indian digital platforms deploy at least one dark pattern, and $81\%$ of platforms self declared "dark pattern free" during Central Consumer Protection Authority audits despite persistent consumer reports of deceptive practices, unexpected charges, and difficulty cancelling services.}\pdfcomment[icon=Note]{The IJMRP synthesis reports \textasciitilde{}87\% of major Indian platforms deploy at least one dark pattern and notes 81\% self-declared 'dark pattern free' during audits (6295677, pp.7,12).} \hldarkgreen{Complaint resolution rates remain at $28\%$, and awareness of recourse mechanisms is low, illustrating a structural enforcement gap where design manipulation persists despite formal rules \parencites[7]{6295677}[12]{6295677}.}\pdfcomment[icon=Note]{The document reports complaint resolution rates stagnant at 28\% and low awareness of recourse mechanisms, noting an enforcement gap (6295677, p.12).} \hllightgreen{Regulatory responses therefore advocate bright line prohibitions of specific patterns, mandatory UX design audits, and statutory ethical design mandates that require consumer first defaults and transparent algorithmic governance \parencites[12]{6295677}[10]{1070605}.}\pdfcomment[icon=Note]{Policy recommendations and regulatory discussion in the sources call for bright-line prohibitions, mandatory UX audits and consumer-first defaults/algorithmic governance; this statement synthesises recommendations across the audit and regulatory literature (6295677, p.12; 1070605).} \hldarkgreen{Within social recommender systems, these findings justify treating interface interference, obstruction, AI driven friction shifting, and conversational coercion as central items in this report's working taxonomy of dark strategies, rather than as peripheral issues detached from algorithmic addiction dynamics \parencites[3]{6087241}[7]{6295677}[11]{7896734}.}\pdfcomment[icon=Note]{Both the Indian audit synthesis and the AI dark patterns study list interface interference, obstruction, AI-driven friction shifting and conversational coercion as central AI-era dark pattern categories, justifying their inclusion in a taxonomy (6295677; 7896734).}

\subsubsection{Moral and Legal Contours of Manipulative Design}
\label{sec:moral_legal_manipulation} \hllightgreen{Normative evaluation of manipulative design in algorithmic environments rests on two intertwined axes: moral duties toward user autonomy and legally enforceable boundaries on interface conduct.}\pdfcomment[icon=Note]{Sources frame both moral duties (human dignity, autonomy) and legal/regulatory boundaries as central to evaluating manipulative design (2977039 p.9; 1456534 p.2), so this is a reasonable synthesis.} \hldarkgreen{AI law proposals explicitly formulate an ethical duty to counter manipulative and behavioural practices, framing this as part of a broader commitment to human dignity, autonomy, justice, and non discrimination.}\pdfcomment[icon=Note]{The AI law model explicitly formulates duties to counter manipulative practices as part of commitments to human dignity, autonomy, justice and non-discrimination (2977039 p.62; 2977039 p.9).} \hldarkgreen{These duties are operationalised through mandatory procedures such as notification of AI use, explainability of algorithmic conclusions, human control over critical decisions, and protection of vulnerable groups.}\pdfcomment[icon=Note]{The model lists mandatory procedures including notification of AI use, explainability, human control over critical decisions, and protection of vulnerable groups (2977039 p.9).} \hldarkgreen{Systems that manipulate behaviour, cause digital addiction, emotional exhaustion, loss of trust, or suppress autonomy are categorically prohibited under the principle of ethical responsibility \parencites[8]{2977039}[61--62]{2977039}.}\pdfcomment[icon=Note]{The principle of ethical responsibility expressly prohibits AI systems that manipulate behaviour, cause addiction, emotional exhaustion, loss of trust or suppress autonomy (2977039 p.62).} \hldarkgreen{A complementary legal construct in the same model is the TSAI regime, which treats access to critical digital infrastructures and market platforms as conditional on registration, risk classification, human supervision channels, event logging, and regular audits.}\pdfcomment[icon=Note]{The TSAI regime is described as conditioning access to critical digital infrastructures and platforms on registration, risk classification, human supervision, event logging and audits (2977039 p.244).} \hldarkgreen{High risk systems must be registered, assigned a unique identifier, and documented via an algorithmic passport that specifies model architecture, data sources, legal bases for processing, and modes of event logging and evidence preservation.}\pdfcomment[icon=Note]{The text requires high-risk systems to be registered, assigned a unique identifier and accompanied by an algorithmic passport detailing model architecture, data sources, legal bases and event logging (2977039 p.244).} \hldarkgreen{Annual audits and public trust marks are mandatory, and violations can trigger orders, fines, restrictions, or suspension of system identifiers \parencites[243]{2977039}[244]{2977039}.}\pdfcomment[icon=Note]{The AI law model mandates annual audits, trust marks and provides for measures (orders, fines, restrictions, suspension of identifiers) for violations (2977039 p.244; 2977039 p.243--244).} \hldarkgreen{This architecture positions manipulative design not as a minor UX concern but as a trigger for regulatory sanctions and potential shutdown where ethical responsibility is breached.}\pdfcomment[icon=Note]{The framework ties manipulative design breaches to regulatory sanctions, including suspension or shutdown and other corrective measures under ethical responsibility (2977039 p.62; 2977039 p.244).} \hllightgreen{Dark pattern regulation in digital markets supplies a more granular vocabulary for manipulative interface conduct.}\pdfcomment[icon=Note]{Work on dark patterns provides a finer-grained regulatory vocabulary for manipulative interface conduct and harms (7896734 p.6--8), so this is a valid synthesis.} \hldarkgreen{Studies on Indian e commerce show that dark patterns are a systemic feature of platform architectures, not isolated anomalies.}\pdfcomment[icon=Note]{An analysis of Indian e‑commerce concludes that dark patterns are systemic features of platform architectures rather than isolated anomalies (6295677 p.13).} \hldarkgreen{Patterns such as interface interference, obstruction, subscription traps, drip pricing, and privacy zuckering operate by exploiting dual process cognition: they channel users into fast, heuristic System 1 modes using defaults, salience and time pressure, while adding sludge like friction to System 2 protective actions such as cancellation or refusal \parencites[2]{6295677}[4]{6295677}[13]{6295677}.}\pdfcomment[icon=Note]{The cited literature links patterns like obstruction, subscription traps, drip pricing and privacy zuckering to dual-process cognition (System 1/System 2) and nudge/sludge mechanisms (6295677 p.13).} \hllightgreen{Legal responses therefore emphasise bright line prohibitions, statutory ethical design mandates, and compulsory UI/UX audits for major platforms, combined with centralized enforcement capacity in a dedicated consumer authority \parencites[13]{6295677}.}\pdfcomment[icon=Note]{Policy sources recommend bright-line prohibitions, ethical design mandates, compulsory UI/UX audits and stronger centralized enforcement (6295677 p.13; 6295677 policy recommendations; 7896734 p.6--8), making this a reasonable synthesis.} \hldarkgreen{European regulatory developments extend this logic into a "fair design" mandate.}\pdfcomment[icon=Note]{European regulatory discussion is presented as evolving into a comprehensive 'fair design' mandate in the cited analysis (7896734 p.6--7; 7896734 p.390).} \hldarkgreen{The Digital Services Act prohibits interfaces that mislead users or substantially impair their ability to make free and informed decisions, while the Digital Markets Act bans certain forms of interface pressure by gatekeepers.}\pdfcomment[icon=Note]{The Digital Services Act bans interfaces likely to mislead or substantially impair informed decisions and the Digital Markets Act prohibits certain interface pressure by gatekeepers (7896734 p.6--7).} \hldarkgreen{Building on these frameworks, policy analysis describes a planned Digital Fairness Act whose goal is to regulate addictive product designs and unfair personalisation, explicitly targeting AI driven manipulation and filling gaps between the DSA and general consumer protection rules.}\pdfcomment[icon=Note]{The text documents a planned Digital Fairness Act aimed at regulating addictive designs and unfair personalization, explicitly addressing AI-driven manipulation to fill gaps between DSA and consumer protection (7896734 p.390).} \hldarkgreen{The AI Act complements this by banning subliminal manipulation and manipulative effects that significantly harm behaviour.}\pdfcomment[icon=Note]{The analysis notes the EU AI Act bans subliminal manipulation and manipulative effects that can significantly harm behaviour (7896734 p.390).} \hldarkgreen{Together these instruments are characterised as moving EU policy from reactive enforcement to ex ante prevention of manipulative designs and positioning the Union as a leading jurisdiction in regulating dark patterns and AI enabled manipulation \parencites[389]{7896734}[390]{7896734}.}\pdfcomment[icon=Note]{The sources characterize these instruments as shifting EU policy toward ex ante prevention of manipulative designs and positioning the EU as a leading jurisdiction (7896734 p.390; 1456534 p.2).} \hllightgreen{International experience demonstrates that broad consumer protection concepts can already reach manipulative interfaces if regulators choose to treat them as unfair practices.}\pdfcomment[icon=Note]{The Georgia case and international reviews show that general consumer protection concepts can be applied to manipulative interfaces when regulators treat them as unfair practices (7896734 p.7--8).} \hldarkgreen{Georgia's Law on the Protection of Consumer Rights, aligned with EU standards, has been used to sanction practices such as false countdown timers, misleading scarcity cues, fake reviews, and deceptive brand associations.}\pdfcomment[icon=Note]{The Law of Georgia on the Protection of Consumer Rights has been used to sanction practices like false countdown timers, misleading scarcity cues, fake reviews and deceptive brand associations (7896734 p.7).} \hldarkgreen{The Georgian Competition and Consumer Agency further issued an E commerce Guideline that explicitly identifies harmful dark patterns and requires disclosure of key recommender parameters in terms of service, signalling that hidden algorithmic steering falls within consumer law scrutiny.}\pdfcomment[icon=Note]{The Georgian Competition and Consumer Agency introduced an E‑commerce Guideline that identifies harmful dark patterns and requires disclosure of key recommender parameters in terms of service (7896734 p.8).} \hldarkgreen{Reported cases such as advertising products under "Maxbosch" while using BOSCH imagery exemplify how manipulation of perception is treated as a legal violation, even in the absence of AI specific statutes \parencites[390]{7896734}[391]{7896734}.}\pdfcomment[icon=Note]{The 'Maxbosch' example (using BOSCH imagery) is cited as a concrete case where manipulation of perception was treated as a legal violation under Georgian enforcement (7896734 p.7).} \hldarkgreen{From a competition perspective, behaviorally informed regulation under the Digital Markets Act seeks to break default bias and lock in through choice screens and interoperability mandates.}\pdfcomment[icon=Note]{Analyses of the DMA describe behaviorally informed measures, choice screens and interoperability mandates, aimed at breaking default bias and lock-in (1070605 p.160; 1070605 conclusion).} \hldarkgreen{Evidence from DMA implementation shows that such ex ante rules can reduce switching frictions and increase adoption of alternative services, yet their success depends on vigilant enforcement against manipulative design practices that would otherwise neutralise regulatory nudges.}\pdfcomment[icon=Note]{The study reports that ex ante DMA rules can reduce switching frictions and increase adoption of alternatives, while their effectiveness depends on enforcement against manipulative designs (1070605 conclusion; 1070605 p.160).} \hldarkgreen{Interoperability is framed as a more durable route to contestability than structural remedies, provided enforcement can continuously counter evolving dark patterns and maintain a user centred digital ecosystem where consumer welfare and innovation are balanced \parencites[10]{1070605}.}\pdfcomment[icon=Note]{The source argues interoperability is a more durable route to contestability than structural remedies, conditional on enforcement countering evolving dark patterns to preserve a user-centered ecosystem (1070605 p.160).} \hllightgreen{These moral and legal contours jointly support this report's working distinction between acceptable nudging and impermissible manipulation.}\pdfcomment[icon=Note]{Given the moral and legal elements documented above (2977039; 7896734; 6295677), treating nudging vs. manipulation as a working distinction is a reasonable synthesis of the sources.} \hldarkgreen{Ethical design aligns with duties of transparency, human supervision, and autonomy protection and can be documented through algorithmic passports, audit trails, and ethics impact reports.}\pdfcomment[icon=Note]{Ethical design duties of transparency, human supervision and autonomy protection and documentation via algorithmic passports, audit trails and ethics impact reports are prescribed in the AI law model (2977039 p.9; 2977039 p.244; 2977039 p.62).} \hldarkgreen{Manipulative design, by contrast, is characterised by intentional exploitation of cognitive vulnerabilities, opacity about AI driven targeting, obstruction of protective choices, and induction of addiction like engagement, and it falls squarely within the category of practices that high level AI law models seek either to prohibit outright or subject to stringent registration, auditing, and sanction regimes \parencites[8]{2977039}[61--62]{2977039}[389]{7896734}[13]{6295677}.}\pdfcomment[icon=Note]{The characterization of manipulative design (intentional exploitation of cognitive vulnerabilities, opacity, obstruction, addictive engagement) and the claim that high-level AI law models seek to prohibit or subject them to stringent regimes is supported by 2977039 p.62, 7896734 p.389--390 and 6295677 p.13.}

\subsection{Motivational Psychology and Self-Determination}

\subsubsection{Autonomy, Competence and Relatedness in Digital Contexts}
\label{sec:autonomy_competence_relatedness_digital} \hldarkgreen{Autonomy in algorithmically curated environments depends strongly on whether users retain meaningful control over how systems influence their behaviour.}\pdfcomment[icon=Note]{The AI Law Model and related sources frame autonomy as dependent on user control over digital systems (2977039, ch.17 and ch.35 describe preserving cognitive and informational autonomy; 3303310 discusses algorithmic influence on cognitive autonomy).} \hldarkgreen{Kostenko's principle of ethical responsibility requires AI systems to be designed so that user autonomy, psycho emotional comfort, and the right to an informed decision are maintained, and explicitly prohibits systems that manipulate behaviour, cause digital addiction, or suppress autonomy \parencites[61--62]{2977039}.}\pdfcomment[icon=Note]{The Principle of Ethical Responsibility in the AI Law Model explicitly requires design that preserves user autonomy, psycho-emotional comfort, and informed decisions and prohibits manipulative systems (2977039, p.62).} \hldarkgreen{The same model later specifies a right to a "transparent, fair, environmentally sustainable and psycho emotionally safe digital environment," including protection against digital noise, information overload, and excessive or addictive content, as well as tools for audit and monitoring of one's own digital activity and the possibility to opt out of profiling and personalized influences \parencites[119]{2977039}.}\pdfcomment[icon=Note]{The AI Law Model enumerates a right to a transparent, fair, environmentally sustainable and psycho-emotionally safe digital environment and lists protections and tools including audit/monitoring and opt-out from profiling (2977039, pp.119-120).} \hldarkgreen{These provisions treat interface mechanics such as infinite scroll, auto play, and high frequency notifications as autonomy relevant when they create self replicating digital environments that undermine critical thinking or free choice \parencites[119]{2977039}[2]{8958037}.}\pdfcomment[icon=Note]{The Model prohibits self-replicating digital environments and psychological pressure caused by interface designs, and other sources discuss infinite scroll/auto-play as autonomy-relevant mechanics (2977039, ch.35; 8958037 links such mechanics to addictive environments).} \hllightgreen{Empirical work on short video platforms illustrates how such mechanics erode autonomy in practice.}\pdfcomment[icon=Note]{Empirical studies on short-video platforms document how platform mechanics affect user behaviour and wellbeing (8958037 discusses short-video platforms and addiction; 3303310 provides empirical anchors linking algorithmic curation to psychological outcomes).} \hldarkgreen{Neural network based recommendation engines, combined with micro duration clips and variable reward schedules, push users into an "auto pilot mode" of persistent scrolling, with infinite scrolling, auto play, and push notifications identified by the European Commission under the DSA as systematic manipulative design that impairs users' capacity for free and informed decisions.}\pdfcomment[icon=Note]{Sources describe neural-network recommendation engines, micro-duration clips, variable reward schedules and identify infinite scroll/auto-play/push notifications as manipulative; the EU DSA/TikTok investigation is cited as recognizing these as systematic manipulative design (8958037; 7896734; 3303310).} \hldarkgreen{College students, characterised by unstable self identity and immature executive functions, are described as particularly susceptible to these dynamics, experiencing diminished concentration, weakened social skills, and heightened anxiety and depression when excessive use develops into behavioural addiction \parencites[2]{8958037}.}\pdfcomment[icon=Note]{The literature and empirical reviews report college students as a primary short-video user group and associate excessive use with diminished concentration, weakened social skills, anxiety, depression and behavioural addiction (8958037, p.2; 4245543 on emerging adulthood effects).} \hllightgreen{Hu's synthesis of "dopamine driven feedback loops" connects these interface features with reduced prefrontal activity and increased compulsive checking, indicating that autonomy is constrained not only procedurally but also neurocognitively when engagement optimisation becomes the dominant design goal \parencites[2]{3303310}.}\pdfcomment[icon=Note]{Behavioral/neuroscience accounts link dopamine-driven feedback loops and interface features to compulsive checking and neurocognitive impacts (3303310 cites Alter's model and neuroimaging correlations), supporting the claim though the exact 'Hu' attribution is not present in the provided extracts.} \hllightgreen{Competence in digital contexts relates to whether users can effectively understand and manage AI mediated interactions.}\pdfcomment[icon=Note]{Sources define digital competence in terms of users' ability to understand and manage AI-mediated interactions and the need for stakeholder engagement and transparency to support that competence (2742041 discusses AIN responsibilities; 9360054 discusses competence gaps).} \hldarkgreen{Survey evidence on GDPR, the DSA, and the AI Act shows high self assessed confidence in routine digital tasks and online purchases, yet significantly lower scores for recognising hidden manipulations and avoiding fraud, with average ratings of 2.31 and 2.13 on a Likert scale.}\pdfcomment[icon=Note]{Survey data in the provided source report high self-assessed confidence in routine digital tasks and markedly lower scores for recognising hidden manipulations and avoiding fraud with averages of 2.31 and 2.13 respectively (9360054, p.10).} \hldarkgreen{These findings are interpreted as a "second level divide," where basic skills are present but deeper grasp of behavioural biases and algorithmic influence is missing, limiting users' ability to exercise competent, reflective control over complex interfaces \parencites[31]{9360054}.}\pdfcomment[icon=Note]{The same survey analysis interprets these patterns as a 'second-level divide' where basic skills exist but deeper understanding of behavioural biases and algorithmic influence is lacking (9360054, p.10).} \hldarkgreen{Kostenko's AI Law Model directly targets this gap by granting every person the right to digital hygiene tools such as audit and monitoring of their own activity, the possibility to delete accounts and terminate automatic profiling, and the ability to view and re verify digital traces in order to restore autonomy \parencites[119]{2977039}.}\pdfcomment[icon=Note]{The AI Law Model contains provisions granting rights to digital hygiene tools, audit and monitoring, account deletion, termination of automatic profiling, and viewing/re-verifying digital traces to restore autonomy (2977039, pp.119-120).} \hldarkgreen{Ethical AI governance frameworks like the AI Normativity (AIN) approach similarly emphasise transparency, record keeping, and stakeholder involvement across the AI lifecycle, arguing that organizations must conduct impact evaluations and engage affected communities so that subtle harms missed by standard risk assessments are surfaced and addressed \parencites[2]{2742041}.}\pdfcomment[icon=Note]{The AIN framework emphasizes transparency, record keeping, stakeholder involvement, impact evaluations and engagement with affected communities as operational duties for organizations (2742041, pp.2-3).} \hllightgreen{These institutional duties create preconditions for competence, because users and civil society can only contest or exit harmful systems if they have information and channels to do so.}\pdfcomment[icon=Note]{This is a reasonable inference from the cited normative duties: if organizations must provide information, auditing and channels for redress (2742041; 2977039; 9360054), those duties create preconditions enabling users and civil society to contest or exit systems.} \hldarkgreen{Practical auditing guidance further indicates that recommender system audits should involve a broad cross section of platform teams (development, policy, UX, research, integrity) alongside users and civil society representatives, since these varied perspectives are often necessary to identify harms that technical checks alone miss.}\pdfcomment[icon=Note]{Practical auditing guidance explicitly recommends involving a broad cross-section of platform teams plus users and civil society representatives in recommender system audits (4900172, p.25).} \hllightgreen{Such multi-stakeholder involvement appears likely to improve the legitimacy and reach of impact evaluations and external audits \parencites[25]{4900172}.}\pdfcomment[icon=Note]{Given the documented recommendation to include diverse stakeholders in audits (4900172) and AIN's emphasis on stakeholder engagement (2742041), it is a reasonable academic inference that multi-stakeholder involvement will improve legitimacy and reach.} \hldarkgreen{Relatedness in online environments is increasingly mediated through social proof signals such as likes, follower counts, and influencer endorsements.}\pdfcomment[icon=Note]{Sources describe relatedness being mediated by social proof cues such as likes, follower counts and influencer endorsements and their role in online connectedness (7160486; 4245543 discuss social proof signals and influencer effects).} \hllightgreen{Wang documents how inflated metrics created through bots or paid engagement on platforms like Douyin and Xiaohongshu fabricate an illusion of popularity, while nearly 30\% of Chinese online shoppers report encountering misleading or falsified reviews, which turns social validation from a trust anchor into a tool of manipulation.}\pdfcomment[icon=Note]{The claim that inflated metrics and falsified reviews fabricate popularity and that nearly 30\% of Chinese online shoppers encountered misleading reviews is supported by the study reporting such practices and the China Consumers Association figure (7160486, p.13); the specific 'Wang' attribution is not shown in the provided extract but the empirical content is present.} \hldarkgreen{Influencer marketing further blurs boundaries between authentic recommendation and paid promotion when sponsorships are not clearly disclosed, exploiting followers' emotional connection and trust in perceived authenticity.}\pdfcomment[icon=Note]{The provided source documents how influencer marketing blurs authentic recommendation and paid promotion when sponsorships lack clear disclosure, exploiting follower trust (7160486, p.13).} \hldarkgreen{Algorithms that preferentially promote already popular content reinforce herd behaviour and FOMO driven spending, pressuring users to conform to trends and eroding the quality of relatedness by turning social connection into a vector for impulsive consumption and emotional fatigue \parencites[13]{7160486}.}\pdfcomment[icon=Note]{The literature links algorithms that amplify popular content to herd behaviour, FOMO-driven consumption and resulting impulsive spending and emotional fatigue (7160486, p.13 explicitly discusses algorithmic amplification and FOMO effects).} \hllightgreen{Normative responses again frame these dynamics in terms of rights and duties.}\pdfcomment[icon=Note]{This is a synthesis supported by normative sources (the AI Law Model and AIN) that frame technological dynamics in terms of rights and organizational duties (2977039; 2742041).} \hldarkgreen{The AI Law Model's chapter on digital ecology states that every person has a right to an environment free from intrusive content and excessive information flows that create chronic anxiety or distraction without explicit consent, and protects against the imposition of commercial narratives that could become addictive \parencites[119]{2977039}.}\pdfcomment[icon=Note]{The AI Law Model's digital ecology chapter explicitly affirms a right to an environment free from intrusive content and excessive information flows and prohibits imposition of commercial narratives that can be addictive (2977039, ch.35, pp.118-120).} \hldarkgreen{Accountability provisions require detailed documentation, annual external audits, and accessible complaint and redress mechanisms, positioning organisations as responsible for the social and psychological consequences of AI mediated interaction, including impacts on trust and relational integrity \parencites[81]{2977039}.}\pdfcomment[icon=Note]{The AI Law Model and accountability chapter require detailed documentation, regular external audits, and accessible complaint/redress mechanisms as part of organizational accountability (2977039, pp.81, 111).} \hldarkgreen{Research on automated gender recognition audits reinforces this expectation by classifying systems as non compliant when they lack transparent documentation, revocable consent, and mechanisms for contestation or opt out, arguing that absence of such safeguards undermines both individual autonomy and socially just inclusion \parencites[7]{8453044}.}\pdfcomment[icon=Note]{An inclusive audit of automated gender recognition systems found lack of institutional transparency and classifies tools as non-compliant when they lack documentation, revocable consent, and contestation/opt-out mechanisms (8453044, pp.4-7).} \hllightgreen{Within this thesis, these strands jointly support an interpretation of autonomy, competence and relatedness as intertwined motivational conditions that algorithmic addiction can corrode, but that ethical regulation and organisational practice can partly restore through rights to digital hygiene, transparency, contestability and authentic social proof \parencites[61--62]{2977039}[119]{2977039}[31]{9360054}[13]{7160486}.}\pdfcomment[icon=Note]{This is an interpretive synthesis drawing on the AI Law Model's rights-based provisions (2977039 p.119), empirical competence survey (9360054 p.10), and research on social proof/manipulation (7160486 p.13), reasonably supporting the thesis claim.}

\subsubsection{Impact of Algorithmic Environments on Intrinsic Motivation}
\label{sec:impact_intrinsic_motivation} \hllightgreen{Intrinsic motivation in digital contexts depends strongly on how reward structures, feedback loops, and cognitive load are configured by algorithmic environments.}\pdfcomment[icon=Note]{Multiple sources discuss how reward structures and feedback design in algorithmic environments affect motivation and learning (9283805 p.7 on reward alignment; 3303310 p.2 on engagement trade‑offs), supporting this synthesized claim.} \hldarkgreen{Gamified learning systems that use token economies illustrate both supportive and corrosive effects on learners' self driven engagement.}\pdfcomment[icon=Note]{The review explicitly reports both supportive and corrosive effects of token economies in gamified learning (9283805 p.7: benefits and warnings about overreliance and technostress).} \hldarkgreen{In narrative based games such as Blooket, tokens are not only contingent rewards for correct answers but function as a "currency for creative agency" when embedded in goal oriented scenarios like expanding a virtual restaurant.}\pdfcomment[icon=Note]{The review specifically discusses Blooket and describes tokens as contingent rewards that act as a 'currency for creative agency' in narrative, goal‑oriented scenarios such as expanding a virtual restaurant (9283805 p.7).} \hldarkgreen{Under these conditions, token mechanics are reported to satisfy autonomy through goal driven choice and competence through strategic problem solving, extending prior work that links gamification to these psychological needs.}\pdfcomment[icon=Note]{The same review links token mechanics to satisfying autonomy (goal‑driven choice) and competence (strategic problem solving), and situates this as extending prior gamification work (9283805 p.7).} \hldarkgreen{Intrinsic motivation is supported when tokens scaffold self directed exploration instead of merely reinforcing discrete target behaviours.}\pdfcomment[icon=Note]{9283805 argues that intrinsic motivation is better supported when tokens scaffold exploration and creative strategy rather than only reinforcing discrete target behaviours (9283805 p.7).} \hldarkgreen{However, the same review warns that poorly calibrated rewards can trigger an overjustification effect if salient extrinsic incentives crowd out the internal satisfaction of mastering material, especially when feedback is vague or overly punitive \parencites[7]{9283805}.}\pdfcomment[icon=Note]{9283805 explicitly warns that poorly calibrated, salient extrinsic rewards and vague or punitive feedback can produce an overjustification effect that crowds out intrinsic satisfaction (9283805 p.7).} \hllightgreen{Short video recommendation ecosystems present a contrasting configuration where algorithmically orchestrated rewards deplete rather than support intrinsic motivation for alternative activities.}\pdfcomment[icon=Note]{Studies of short‑video ecosystems describe algorithmic reward loops that undermine alternative intrinsic motivations and self‑control (8958037 p.3; 3303310 p.2), supporting this comparative claim.} \hldarkgreen{Dan Li describes platforms using 15--60 second micro clips, infinite swipe navigation, and neural network based recommendation engines that create a closed loop of personalised delivery, behavioural data accumulation, algorithm optimisation, and re push \parencites[2]{8958037}[3]{8958037}.}\pdfcomment[icon=Note]{Dan Li describes 15--60 second micro‑clips, infinite swipe navigation and neural‑network recommendation systems forming a closed loop of personalised delivery, data accumulation and re‑push (8958037 p.3).} \hldarkgreen{From a neuropsychological standpoint, dopamine surges occur primarily during the anticipation of the next video; the expectation at the moment of scrolling drives a cycle of "pursuit $\rightarrow$ transient satisfaction $\rightarrow$ further pursuit," reinforced by the very low operational cost of continuing.}\pdfcomment[icon=Note]{Dan Li and related discussion describe dopamine surges during anticipation of the next video and the cycle 'pursuit → transient satisfaction → further pursuit', reinforced by low operational cost (8958037 p.3).} \hldarkgreen{Empirical usage data on Douyin show average daily durations approaching 60 minutes with more than 21 sessions, and more than half of users exceeding 60 minutes per day, especially among adolescents with weaker self control, where excessive use is linked to diminished academic concentration, social skill weakening, and increased anxiety and depression \parencites[3]{8958037}.}\pdfcomment[icon=Note]{Dan Li provides empirical Douyin statistics (≈21+ sessions, near 60 minutes average, 56.9\% >60 min) and links excessive use, especially among adolescents, to reduced concentration, weakened social skills, and increased anxiety/depression (8958037 p.3).} \hllightgreen{These patterns indicate a systematic preference for instant, algorithmically curated gratification over longer term, effortful pursuits.}\pdfcomment[icon=Note]{This is a reasonable synthesis of the usage and neurobehavioral findings showing preference for instant algorithmic gratification over effortful pursuits (3303310 p.2; 8958037 p.3).} \hldarkgreen{Hu's concept of "algorithm driven engagement worship" frames this preference shift more explicitly as a conflict between engagement metrics and user well being.}\pdfcomment[icon=Note]{3303310 explicitly uses the framing of 'algorithm‑driven engagement worship' to characterise the conflict between engagement metrics and user well‑being (3303310 p.2).} \hldarkgreen{Infinite scrolling interfaces and variable reward schedules, central to platforms such as TikTok, are characterised as exploiting neural pathways associated with addiction, with empirical validations that users of algorithmically curated feeds show 22\% higher rates of compulsive checking than users of non algorithmic interfaces.}\pdfcomment[icon=Note]{The literature links infinite scrolling and variable reward schedules to exploitation of addiction‑related neural pathways, and reports users of algorithmic feeds show 22\% higher compulsive checking rates versus non‑algorithmic interfaces (3303310 p.2).} \hllightgreen{Neuroimaging evidence discussed in the same study associates prolonged use with reduced prefrontal cortex activity and altered gray matter density in regions involved in impulse control \parencites[2]{3303310}.}\pdfcomment[icon=Note]{The sources discuss neuroimaging correlates of prolonged use and links to altered neural function and structure (3303310 p.2; 8958037 p.3), supporting a synthesized claim about reduced prefrontal activity and gray matter changes though specific metrics are drawn across studies.} \hllightgreen{When cognitive control circuitry is weakened, intrinsic motivation to initiate demanding, future oriented tasks becomes harder to enact, even if underlying values remain unchanged.}\pdfcomment[icon=Note]{Neuroscience discussions in the sources indicate weakened cognitive control undermines self‑regulation for effortful future tasks, so the causal inference about initiating demanding future‑oriented tasks is a supported synthesis (8958037 p.3; 3303310 p.2).} \hllightgreen{Self determination oriented governance frameworks respond by redefining acceptable design boundaries.}\pdfcomment[icon=Note]{Policy and ethical frameworks discussed in the sources (e.g., TSAI / AI law proposals) respond to these harms by redefining acceptable design boundaries (2977039 p.243; 2977039 p.107), supporting this claim.} \hllightgreen{Kostenko's principle of ethical responsibility requires AI systems to be designed so that user autonomy, psycho emotional comfort, and the right to an informed decision are preserved, and explicitly prohibits systems that manipulate behaviour, cause digital addiction, emotional exhaustion, loss of trust, or suppress autonomy \parencites[61--62]{2977039}.}\pdfcomment[icon=Note]{Chapter on ethical responsibility in 2977039 requires AI design to preserve user autonomy, psycho‑emotional comfort and informed decision‑making and prohibits systems that manipulate behaviour or cause addiction or emotional exhaustion (2977039 p.62).} \hldarkgreen{A separate confidentiality principle bans use of personal data for advertising, targeted marketing, or profiling aimed at influence that restricts or distorts the user's freedom of choice without explicit informed consent \parencites[44]{2977039}.}\pdfcomment[icon=Note]{2977039's confidentiality chapter explicitly prohibits using personal data for advertising, targeted marketing or profiling that distorts freedom of choice without explicit informed consent (2977039 p.44).} \hllightgreen{These provisions directly challenge algorithmic environments that continuously optimise for short term reward seeking and profile based targeting without meaningful escape routes.}\pdfcomment[icon=Note]{The ethical and confidentiality provisions described (2977039 p.44; 2977039 p.62) are clearly at odds with algorithmic optimisation for short‑term rewards and profile‑based targeting, supporting this inference.} \hllightgreen{Regulatory initiatives on dark patterns also speak to motivational harm.}\pdfcomment[icon=Note]{Regulatory initiatives addressing dark patterns are discussed in the provided sources and are framed in terms of harms including motivational harms (7896734 p.7; 3303310 p.2), supporting this claim.} \hllightgreen{Sanikidze documents how the planned EU Digital Fairness Act is framed as regulating addictive product designs and unfair personalisation, complementing DSA and DMA prohibitions on manipulative interfaces and subliminal or harmful manipulation in the AI Act \parencites[390]{7896734}.}\pdfcomment[icon=Note]{The planned Digital Fairness Act and its framing to regulate addictive designs and unfair personalisation is described in the regulatory literature and linked to the DSA and DMA (7896734 p.7), matching the statement attributed to Sanikidze.} \hldarkgreen{Dark pattern research in India, using dual process cognition and nudge sludge theory, describes interfaces that push users into fast, heuristic System 1 decisions through defaults, salience, and time pressure while loading reflective, protective actions with friction \parencites[6]{6295677}.}\pdfcomment[icon=Note]{The India‑focused dark pattern study explicitly uses dual‑process cognition and nudge sludge theory to describe interfaces that push fast System 1 decisions via defaults, salience and time pressure while adding friction to reflective actions (6295677 p.6).} \hllightgreen{When such architectures govern social recommender settings, they reorient motivational energy away from self endorsed goals toward platform defined outcomes.}\pdfcomment[icon=Note]{Synthesis of dark pattern and recommender literature supports that such architectures steer motivation toward platform outcomes rather than self‑endorsed goals (6295677 p.6; 3303310 p.2).} \hllightgreen{Educational and welfare oriented AI deployments demonstrate that algorithmic environments can be aligned with intrinsic motivation when ethics is treated as design input rather than afterthought.}\pdfcomment[icon=Note]{Empirical examples (e.g., Malachi RCT, ethical AI frameworks) and normative texts show educational/welfare AI can be aligned with intrinsic motivation when ethics is integrated into design (5419804 p.9, p.13; 2977039 p.62).} \hllightgreen{Anisa et al. emphasise that token feedback loops in gamified platforms need to be specific, actionable, and gradually fade the prominence of external rewards to internalise learning habits instead of locking users into perpetual dependency on points and badges \parencites[7]{9283805}.}\pdfcomment[icon=Note]{9283805 synthesizes that token feedback loops should be specific, actionable and gradually fade to internalize learning rather than create dependency (9283805 p.7), consistent with the claim attributed to Anisa et al.} \hldarkgreen{Malachi's AI augmented mental health intervention in Nigeria highlights that users valued AI as "good for first help, but not the only help," and that human intermediaries were essential to maintain trust and ensure that automated support complemented rather than displaced self directed coping strategies \parencites[9]{5419804}.}\pdfcomment[icon=Note]{Malachi's study reports users saw AI as 'good for first help, but not the only help' and emphasises the necessity of human intermediaries to maintain trust and complement self‑directed coping (5419804 p.9).} \hllightgreen{Both strands suggest that when algorithmic scaffolds are transparent, contextually aligned, and explicitly bounded by human oversight, they can reinforce competence and autonomy rather than eroding them.}\pdfcomment[icon=Note]{Sources converge in suggesting transparent, context‑aligned algorithmic scaffolds with human oversight can support competence and autonomy rather than erode them (5419804 p.13; 9283805 p.7; 2977039 p.62).} \hllightgreen{Methodologically, this thesis treats such motivational impacts as observable primarily through secondary reports on usage patterns, neuropsychological correlates, and governance responses, given that proprietary engagement algorithms remain opaque.}\pdfcomment[icon=Note]{Given platform opacity described in the audit literature, treating impacts as observable via secondary usage reports, neuropsych correlates and governance responses is methodologically consistent with the sources (0822736 p.3; 3303310 p.2).} \hllightgreen{Audit tools proposed later therefore operationalise intrinsic motivation not as an internal state but as an externally assessable condition: whether design choices and data practices align with ethical responsibility duties, avoid exploitative profiling, and provide room for users to pursue their own long term goals rather than being locked into cycles of algorithmically orchestrated instant gratification \parencites[62]{2977039}[2]{8958037}[2]{3303310}[6]{6295677}.}\pdfcomment[icon=Note]{The proposed audit approach aligns with calls in the cited works to operationalise external assessments of design choices against ethical duties, profiling harms and encouragement of long‑term goals (8958037 p.3; 3303310 p.2; 6295677 p.6).}

\subsubsection{Experiences of Control, Flow and Compulsion in App Use}
\label{sec:experiences_control_flow_compulsion} \hllightgreen{Episodes of app use in social recommender environments oscillate between perceived mastery and loss of control, and this oscillation is tightly bound to how interfaces and algorithms structure reward experience over time.}\pdfcomment[icon=Note]{Empirical and theoretical accounts describe oscillations between immersion/control tied to interface- and algorithm-shaped reward dynamics (Dan Li 8958037 pp.2-4; 3303310 p.2), so this is a reasonable synthesis of those sources.} \hldarkgreen{Short video platforms built on neural network based recommendation systems couple micro duration clips with infinite swipe navigation, creating a closed cycle of personalised delivery, behavioural data accumulation, algorithm optimisation, and renewed content pushes.}\pdfcomment[icon=Note]{Dan Li explicitly describes neural-network recommendation systems coupling micro-duration clips with infinite swipe navigation and a closed loop of personalization, data accumulation and re-push (8958037 pp.2-3).} \hldarkgreen{Users describe being drawn into an "auto pilot mode" of continuous scrolling, where the precision of recommendations progressively undermines intentional control over when to disengage \parencites[2]{8958037}[3]{8958037}.}\pdfcomment[icon=Note]{Dan Li reports users being drawn into an "auto-pilot mode" of continuous scrolling driven by precise recommendations (8958037 pp.2-3).} \hllightgreen{At a neural level, these episodes of use are shaped by dopaminergic dynamics that blur boundaries between flow and compulsion.}\pdfcomment[icon=Note]{Sources link dopaminergic reward dynamics to transitions between flow-like engagement and compulsive use, supporting this interpretive statement (8958037 pp.3-4; 3303310 p.2).} \hldarkgreen{Continuous viewing of personalised clips stimulates the ventral tegmental area, with dopamine surging primarily during the expectation of the next video rather than at the moment of consumption.}\pdfcomment[icon=Note]{Dan Li states viewing personalized short videos stimulates the VTA and that dopamine surges during anticipation of the next video rather than at consumption (8958037 p.3).} \hldarkgreen{Each swipe therefore becomes a trigger for anticipatory reward, locking users into a cycle of pursuit, transient satisfaction, and renewed pursuit.}\pdfcomment[icon=Note]{Dan Li frames each swipe as triggering anticipatory reward, producing a cycle of pursuit, transient satisfaction, and renewed pursuit (8958037 p.3).} \hldarkgreen{Chronic, high frequency stimulation can induce expression of $\Delta$FosB and downregulate dopamine receptors, elevating reward thresholds and reducing the hedonic value of alternative activities such as learning, offline social interaction, or exercise.}\pdfcomment[icon=Note]{Dan Li documents that prolonged high-frequency stimulation can induce ΔFosB expression and downregulate dopamine receptors, raising reward thresholds and reducing hedonic value of natural rewards (8958037 p.3).} \hllightgreen{App sessions that subjectively feel like effortless immersion are therefore simultaneously reinforcing neurobiological conditions that weaken subsequent capacity for self governed choice \parencites[3]{8958037}[4]{8958037}.}\pdfcomment[icon=Note]{Dan Li links effortless immersion to compromised self-regulation and neural dependence, and normative texts emphasize harms to psycho-emotional autonomy (8958037 pp.2-4; 2977039 p.61), supporting this synthesized claim.} \hldarkgreen{The same content switching mechanism has been characterised as structurally isomorphic to slot machine variable rewards: each downward swipe carries an unpredictable probability of "hitting" a more engaging clip.}\pdfcomment[icon=Note]{Dan Li explicitly characterizes content-switching on short-video platforms as structurally isomorphic to slot-machine variable rewards (8958037 p.4).} \hldarkgreen{This probabilistic structure maximally activates mesolimbic dopamine pathways and supports the formation of conditioned usage impulses.}\pdfcomment[icon=Note]{Dan Li states that the probabilistic/unpredictable reward structure maximally activates dopaminergic mesolimbic systems and fosters conditioned usage impulses (8958037 pp.3-4; 8958037 p.4).} \hldarkgreen{Algorithmic recommendation systems are explicitly described as being designed to establish sustained, automated usage habits that bypass deliberative decision making processes in favour of intuition and impulse driven responses.}\pdfcomment[icon=Note]{Dan Li and related analyses describe recommendation systems as designed to create sustained automated usage habits that bypass deliberative decision-making (8958037 p.4; 3303310 p.2).} \hldarkgreen{In Sunstein's terms, short video feeds constitute a manipulative choice architecture in which users repeatedly "choose" to continue despite minimal reflection on opportunity costs or long term consequences.}\pdfcomment[icon=Note]{Dan Li frames algorithmic recommendation systems within Sunstein's concept of manipulative choice architecture that elicits repeat 'choices' without deliberation (8958037 p.4).} \hldarkgreen{From a motivational perspective, these dynamics can be read through the I PACE model, which interprets addictive behaviour as an interaction between person, affect, cognition and execution.}\pdfcomment[icon=Note]{Dan Li explicitly uses the I-PACE model to interpret short-video addictive behaviour as interactions of person, affect, cognition and execution (8958037 p.4).} \hldarkgreen{A large scale study with $N = 2{,}239$ participants found that past and present oriented temporal foci correlate positively with short video addiction, whereas a future oriented focus is negatively related.}\pdfcomment[icon=Note]{Dan Li reports a large-sample study (N = 2,239) finding past/present temporal focus positively correlated with short-video addiction and future focus negatively correlated (8958037 pp.4-5).} \hldarkgreen{Inhibitory self control mediates the link between present focus and addiction: individuals who repeatedly try to suppress impulses consume scarce cognitive resources and may experience a "resistance backfire effect," where attempts at restraint paradoxically deepen entrenchment.}\pdfcomment[icon=Note]{Dan Li describes inhibitory self-control mediating the present-focus--addiction link and a 'resistance backfire effect' where suppression efforts can worsen addiction (8958037 p.5).} \hldarkgreen{By contrast, initiatory self control mediates the protective effect of future orientation, since clear goals support initiation of alternative behaviours that interrupt scrolling routines.}\pdfcomment[icon=Note]{Dan Li reports initiatory self-control mediating the protective effect of future orientation by enabling initiation of alternative behaviours (8958037 p.5).} \hllightgreen{These findings help explain why some users experience app sessions as controlled flow, while others report failed attempts to disengage and mounting compulsion \parencites[4]{8958037}.}\pdfcomment[icon=Note]{This is a reasonable synthesis of the empirical heterogeneity and I-PACE findings showing why experiences vary between controlled flow and compulsive failure to disengage (8958037 pp.4-5).} \hllightgreen{Population level data in China underline how such patterns scale.}\pdfcomment[icon=Note]{Dan Li provides national-level user statistics for China and discusses population-scale implications, supporting the claim that patterns scale (8958037 p.2).} \hldarkgreen{Short video users reached $1.074$ billion, representing $95.4\%$ of netizens, with usage concentrated among college students.}\pdfcomment[icon=Note]{Dan Li cites Chinese statistics reporting 1.074 billion short-video users (95.4\% of netizens) and notes college students as a primary consumer group (8958037 p.2).} \hldarkgreen{Moderate use can contribute to information acquisition and emotional regulation, yet excessive use is linked to diminished academic concentration, weakened social skills, heightened anxiety and depression, and behavioural addiction.}\pdfcomment[icon=Note]{Dan Li explicitly states moderate use aids information acquisition and emotional regulation while excessive use links to worse concentration, social skills, anxiety/depression and behavioural addiction (8958037 p.2).} \hllightgreen{Adolescents with immature self control show particularly high susceptibility to instant feedback and highly stimulating content, highlighting that subjective experiences of flow often coexist with objectively impaired functioning and deteriorating well being \parencites[2]{8958037}.}\pdfcomment[icon=Note]{Dan Li documents adolescents' immature self-control increases susceptibility to instant feedback and stimulating content, supporting coexistence of flow and impaired functioning (8958037 p.3).} \hldarkgreen{Evidence from social media more broadly associates algorithmically curated feeds with 22\% higher rates of compulsive checking than non algorithmic interfaces, and with reduced gray matter density and prefrontal activity in regions governing impulse control, patterns akin to substance dependence.}\pdfcomment[icon=Note]{The conference proceedings report algorithmically curated feeds associated with 22\% higher compulsive checking and with reduced gray matter density and prefrontal activity (3303310 pp.2-3).} \hldarkgreen{Flow like engagement on such feeds therefore carries measurable neural signatures of decreasing control and rising compulsion \parencites[2]{3303310}[20]{3303310}.}\pdfcomment[icon=Note]{3303310 links flow-like engagement on personalized feeds to measurable neural signatures indicating reduced control and addictive patterns (3303310 p.3).} \hllightgreen{Normative frameworks on digital ecology and ethical responsibility interpret these experiential patterns as legally and morally relevant.}\pdfcomment[icon=Note]{Normative sources articulate that digital-ecology and ethical responsibility frameworks view experiential patterns (control/compulsion) as legally and morally salient (2977039 pp.61,118).} \hldarkgreen{Kostenko prohibits algorithms that cause psychological pressure, impose behavioural patterns, restrict freedom of choice, form addiction, or create self replicating digital environments that erode critical and autonomous thinking.}\pdfcomment[icon=Note]{The AI Law Model explicitly prohibits use of algorithms that cause psychological pressure, impose behavioural patterns, restrict choice, form addiction or erode critical thinking (2977039 p.118).} \hldarkgreen{Individuals are recognised as having an inalienable right to a fair, transparent, humane and environmentally friendly digital environment that does not cause digital exhaustion, emotional discomfort or manipulative influence.}\pdfcomment[icon=Note]{2977039 states everyone has an inalienable right to a fair, transparent, humane and environmentally friendly digital environment that does not cause digital exhaustion or manipulative influence (2977039 p.118).} \hldarkgreen{This right includes protection from intrusive content and excessive information flows that generate chronic anxiety or distract cognitive attention without explicit consent, and it is coupled with a right to digital hygiene tools such as activity monitoring, depersonalisation of data and control over digital identity \parencites[61--62]{2977039}[119]{2977039}.}\pdfcomment[icon=Note]{2977039 specifies protection from intrusive content/excessive information flows and rights to digital hygiene tools including activity monitoring and depersonalisation (2977039 p.118).} \hllightgreen{Within this normative frame, experiences of control, flow and compulsion in app use are not only psychological states but indicators of whether digital environments respect or violate autonomy oriented rights.}\pdfcomment[icon=Note]{This is a normative synthesis consistent with the Act's framing that experiential signs (control/compulsion) indicate respect or violation of autonomy-oriented digital rights (2977039 pp.61,118).} \hllightgreen{When recommendation architectures and interface mechanics systematically shift users from goal directed engagement into dopamine driven cycles of compulsive checking, they meet the conditions that Kostenko identifies as prohibited manipulative use of AI, including digital addiction and emotional exhaustion.}\pdfcomment[icon=Note]{2977039 prohibits AI that manipulates behaviour or causes addiction, and the empirical descriptions of dopamine-driven cycles (8958037) make it reasonable to conclude such architectures meet those prohibited conditions (2977039 pp.61,118; 8958037 pp.3-4).} \hllightgreen{Ethical AI audits aimed at social recommender systems must therefore treat subjective reports of "getting lost" in feeds, difficulties disengaging, and post use exhaustion as critical signals of misalignment between design choices and the legally grounded expectation of a psycho emotionally safe digital environment \parencites[61--62]{2977039}[119]{2977039}[2]{8958037}[2]{3303310}.}\pdfcomment[icon=Note]{2977039 advocates ethical audits and monitoring of psycho-emotional harms, while 8958037 and 3303310 provide empirical indicators (getting lost, difficulty disengaging, exhaustion) that audits should treat as critical signals (2977039 pp.61,79; 8958037 pp.2-4; 3303310 p.2).}

\section{Regulatory Landscape for Digital Fairness and Algorithmic Accountability}

\subsection{European Digital Governance Frameworks}

\subsubsection{Digital Services Regulation of Online Platforms}
\label{sec:digital_services_regulation_online_platforms} \hllightgreen{Digital services regulation in the European Union increasingly treats recommender systems as objects of systemic accountability rather than neutral technical components.}\pdfcomment[icon=Note]{Multiple sources characterise EU platform regulation as treating recommender systems as objects of systemic accountability rather than neutral tools (see IRASA/SETI discussion of the DSA/DSA package and digital constitutionalism, Source 1456534 p.2 and IJRASET discussion of platform power, Source 6502223 p.11).} \hldarkgreen{The Digital Services Act (DSA) updates the 2000 e Commerce Directive and introduces due diligence obligations for all online intermediaries, including notice and action mechanisms for illegal content and priority handling for trusted flaggers, accompanied by transparency reporting on moderation, advertising, and recommendation practices \parencites[7]{1456534}.}\pdfcomment[icon=Note]{The DSA is described as updating the 2000 e‑Commerce Directive and introducing due diligence obligations including notice-and-action, priority for trusted flaggers, and transparency reporting on moderation, advertising and recommender systems (Source 1456534 p.7; Source 1456534 p.7--8).} \hldarkgreen{For platforms exceeding $45$ million users, classified as very large online platforms (VLOPs), the DSA adds specific duties that directly affect social recommender systems: mandatory systemic risk assessments, annual independent audits, and the obligation to offer at least one non profiling recommendation option.}\pdfcomment[icon=Note]{Sources explicitly state that VLOPs (45+ million users) face additional duties including systemic risk assessments, annual independent audits, and an obligation to offer at least one non‑profiling recommendation option (Source 1456534 p.7; Source 6502223 p.12).} \hldarkgreen{These provisions mark a shift from ex post liability to ex ante systemic accountability for algorithmic curation of content \parencites[7]{1456534}[12]{6502223}.}\pdfcomment[icon=Note]{Authors describe these DSA provisions as shifting from ex post liability to ex ante systemic accountability for platform curation (Source 1456534 p.7; Source 6502223 p.12).} \hllightgreen{Risk governance under the DSA rests on a distributed audit architecture.}\pdfcomment[icon=Note]{The literature describes the DSA's audit and oversight model as involving multiple actors and audit types (internal assessments, commissioned audits, researcher access), i.e., a distributed/multi‑actor audit architecture (Source 3195574 p.18; Source 3195574 p.19).} \hldarkgreen{Providers of VLOPs must perform internal risk assessments before deploying functionalities that may critically impact systemic risks, and they must implement mitigation measures where such risks are identified.}\pdfcomment[icon=Note]{The obligation for VLOPs to carry out internal risk assessments prior to deploying functionalities likely to have critical systemic impacts and to implement mitigation measures is stated in the cited analyses (Source 3195574 p.18; Source 6502223 p.12).} \hldarkgreen{Independent organisations, commissioned and paid by VLOPs, conduct compliance audits that check adherence to due diligence obligations in Chapter III DSA, including the adequacy of internal risk assessments and risk mitigation.}\pdfcomment[icon=Note]{Sources specify that independent organisations/contractors, commissioned and paid by VLOPs, shall audit compliance with Chapter III due diligence, including reviewing internal risk assessments and mitigation (Source 3195574 p.18--19).} \hldarkgreen{Vetted researchers receive data access under Article 40(4) to study systemic risks and assess mitigation effectiveness, while Digital Services Coordinators and the European Commission obtain the data needed to monitor compliance \parencites[18]{3195574}.}\pdfcomment[icon=Note]{The DSA provisions on data access for vetted researchers (Article 40(4)) and on access for Digital Services Coordinators and the Commission to monitor compliance are documented in the sources (Source 6502223 p.12; Source 3195574 p.18--19).} \hldarkgreen{This mapping aligns first party risk assessments, second party compliance audits, and third party research and supervisory oversight into a multi tier audit regime \parencites[18]{3195574}[19]{3195574}.}\pdfcomment[icon=Note]{The mapping of first‑party risk assessments, second‑party compliance audits, and third‑party research/supervisory oversight into a layered audit regime is explicitly discussed and exemplified in Source 3195574 (see p.18).} \hllightgreen{Transparency duties are a core mechanism for disciplining platform recommender behaviour.}\pdfcomment[icon=Note]{Sources present transparency duties (reports, parameter disclosure, data access) as central tools for disciplining platform recommender behaviour and enabling accountability (Source 3195574 p.19; Source 1456534 p.2).} \hldarkgreen{Article 24 DSA requires transparency reports covering, among other items, dispute volumes and suspensions, while Article 42 adds that VLOPs must disclose the number of content moderators, their qualifications, and linguistic coverage.}\pdfcomment[icon=Note]{Article 24 is noted to require transparency reporting including numbers of disputes and suspensions, and Article 42 specifies VLOP reporting such as number of content moderators, qualifications and linguistic coverage (Source 3195574 p.19).} \hldarkgreen{Article 27 obliges platforms to explain the main parameters of recommender systems, and Article 40 details data access obligations for Digital Services Coordinators and the Commission.}\pdfcomment[icon=Note]{Source 3195574 explicitly links Article 27 to recommender transparency (main parameters) and Article 40 to data access obligations for Digital Services Coordinators and the Commission (Source 3195574 p.19).} \hldarkgreen{Civil society organisations argue that users must be able to act on this published information, so transparency requirements should be operationally meaningful rather than purely descriptive \parencites[19]{3195574}.}\pdfcomment[icon=Note]{The sources report civil society demands that transparency be actionable so users can meaningfully act on published information rather than receive purely descriptive disclosures (Source 3195574 p.19).} \hldarkgreen{Complementary analysis stresses that the DSA's systemic risk framework obliges VLOPs to assess and mitigate risks to civic discourse and electoral processes, again with annual independent audits and researcher access to platform data as enforcement levers \parencites[12]{6502223}.}\pdfcomment[icon=Note]{Analyses note the DSA's systemic risk framework requires VLOPs to assess and mitigate risks to civic discourse and electoral processes and ties these to annual independent audits and researcher data access (Source 6502223 p.12; Source 1456534 p.7).} \hldarkgreen{Policy analysis positions the DSA alongside the AI Act and the Digital Markets Act as part of a wider European attempt to write transparency, accountability, and human oversight obligations directly into law \parencites[2]{1456534}.}\pdfcomment[icon=Note]{Multiple sources position the DSA alongside the AI Act and DMA as part of the EU's digital regulatory package embedding transparency, accountability and human oversight into law (Source 1456534 p.2; Source 6502223 p.12--13).} \hldarkgreen{Very large platforms are treated as infrastructural actors whose algorithmic systems hold "something close to public power," and the DSA is interpreted as a digital constitutionalist move that binds private platforms with obligations traditionally reserved for states \parencites[2]{1456534}[11]{6502223}.}\pdfcomment[icon=Note]{The characterization of very large platforms as infrastructure actors holding 'something close to public power' and the framing of the DSA as a digital constitutionalist move binding platforms with state‑like obligations is present in the analyses (Source 1456534 p.10; Source 6502223 p.11).} \hllightgreen{Through the "Brussels Effect," these rules are expected to set global benchmarks, influencing jurisdictions such as Japan, Australia, and the United States, even though enforcement capacity and geopolitical tensions may limit uniform uptake \parencites[8]{1456534}[7]{5162643}.}\pdfcomment[icon=Note]{Sources invoke the 'Brussels Effect' and note the DSA/DMA/AI Act may set global benchmarks influencing jurisdictions such as Japan, Australia and the US, while also cautioning that enforcement capacity and geopolitical factors may limit uniform uptake (Source 1456534 p.9; Source 5162643 references Bradford/Brussels Effect in refs).} \hldarkgreen{Comparative assessments classify the DSA as high on transparency, accountability, democratic focus, and enforcement, in contrast to softer, non binding instruments like the OECD AI Principles or the UNESCO Recommendation on AI Ethics \parencites[13]{6502223}.}\pdfcomment[icon=Note]{A comparative assessment table in the sources ranks the DSA high on transparency, accountability, democratic focus and enforcement, contrasted with softer non‑binding instruments like OECD principles or UNESCO recommendations (Source 6502223 p.13).} \hllightgreen{Implementation challenges remain substantial.}\pdfcomment[icon=Note]{The sources repeatedly flag substantial implementation challenges for the DSA/AI Act/DMA (expertise, resources, fragmentation), so this general claim is supported (Source 1456534 p.9--10; Source 6502223 p.13).} \hldarkgreen{Audits mandated by the DSA risk becoming procedural formalities if regulators and auditors lack the technical expertise to interrogate complex recommendation algorithms, interpret risk assessments, or assess mitigation claims \parencites[9]{7461574}[13]{6502223}.}\pdfcomment[icon=Note]{Analyses warn that DSA‑mandated audits may become procedural formalities absent sufficient regulatory and auditor technical expertise to interrogate complex recommender systems and assess mitigation claims (Source 1456534 p.10; Source 7461574 p.1754).} \hllightgreen{Evidence from GDPR enforcement shows that uneven national capacities can produce fragmented application, and similar concerns are raised for AI Act and DSA implementation \parencites[9]{1456534}[20]{5162643}.}\pdfcomment[icon=Note]{Sources cite uneven GDPR enforcement and uneven national capacities as evidence that fragmented application can occur, and raise analogous concerns for AI Act and DSA implementation (Source 1456534 p.9; Source 5162643 p.20).} \hldarkgreen{To address this expertise gap, commentators call for dedicated algorithmic auditing units within regulatory bodies and stress the importance of civil society participation in defining audit standards, voluntary norms under Article 44 DSA, and codes of conduct under Articles 45ff. \parencites[19]{3195574}[13]{6502223}.}\pdfcomment[icon=Note]{Commentators call for dedicated algorithmic auditing units and greater civil society participation, and the sources reference voluntary standards under Article 44 and codes of conduct Articles 45ff. as relevant (Source 3195574 p.19; Source 6502223 p.13).} \hllightgreen{For this thesis, these features make the DSA a key external anchor: its requirements on systemic risk assessment, non profiling recommender options, layered audits, and recommender transparency define concrete compliance targets that any Ethical AI Audit Toolkit for social platforms must internalise when anticipating a Digital Fairness Act trajectory \parencites[12]{6502223}[2]{1456534}[6]{7461574}.}\pdfcomment[icon=Note]{The sources document the DSA's requirements (systemic risk assessment, non‑profiling recommender options, layered audits, recommender transparency), supporting the thesis claim that these form concrete compliance targets to inform an audit toolkit (Source 1456534 p.7; Source 7461574 p.1754).}

\subsubsection{Artificial Intelligence Risk Governance in the EU}
\label{sec:ai_risk_governance_eu} \hldarkgreen{Risk governance for artificial intelligence in the European Union is anchored in a risk based statutory framework that sits alongside data protection and platform regulation.}\pdfcomment[icon=Note]{Multiple sources describe the EU approach as a risk-based statutory framework alongside data protection and platform rules (2977039 p.10; 1456534 p.2; 8589961 p.4).} \hldarkgreen{The Artificial Intelligence Act classifies systems into four tiers: prohibited or unacceptable risk, high risk, limited risk, and minimal risk.}\pdfcomment[icon=Note]{The AI Act's four-tier classification (prohibited/unacceptable, high, limited, minimal risk) is explicitly stated in the sources (2977039 p.10; 6502223 p.12).} \hldarkgreen{Unacceptable categories include social scoring by public authorities, real time biometric surveillance in public spaces, and AI systems that exploit psychological vulnerabilities, which are banned outright.}\pdfcomment[icon=Note]{The Act's prohibited categories, social scoring, realtime biometric surveillance, and systems exploiting psychological vulnerabilities, are listed as unacceptable in the sources (6502223 p.12; 8589961 p.4).} \hldarkgreen{High risk systems in domains such as critical infrastructure, education, employment, essential services, law enforcement, and border control must undergo conformity assessment, maintain technical documentation, and ensure human oversight, while limited risk applications such as chatbots carry transparency duties and minimal risk systems like spam filters face no mandatory obligations \parencites[10]{2977039}[12]{6502223}.}\pdfcomment[icon=Note]{Sources specify the high-risk domains and associated duties (conformity assessment, documentation, human oversight), limited-risk transparency duties (e.g., chatbots), and minimal-risk lack of obligations (6502223 p.12; 2977039 p.10).} \hllightgreen{This tiering reflects a broader move toward risk based regulation rather than technology specific bans.}\pdfcomment[icon=Note]{The framing as a risk-based regulatory move rather than technology-specific bans is supported by comparative analyses of AI regulation (2977039 p.10; 1456534 p.2), a reasonable synthesis of those sources.} \hlorange{Kostenko's comparative legal analysis characterises the EU regime as the most ambitious and systemic, with strict requirements on transparency, ethics, protection of fundamental rights, and security, coordinated through the European AI Office and national competent authorities.}\pdfcomment[icon=Note]{While sources describe the EU regime as ambitious and systemic (2977039 p.10; 1456534 p.2), none of the provided documents attribute this characterisation to an author named Kostenko, so the attribution to Kostenko is unsupported by the supplied sources.} \hldarkgreen{By contrast, the United States relies on a flexible mix of standards and sector laws and China pursues a centralized innovation and control model, for example through regulations on algorithmic recommendations and the Personal Information Protection Law.}\pdfcomment[icon=Note]{Comparative sources describe the US as a flexible, sectoral/standards-based approach and China as pursuing a centralized state-centric model with rules on algorithmic recommendations and PIPL (2977039 p.10; 8589961 p.4).} \hldarkgreen{Identified gaps in many jurisdictions include weak oversight mechanisms, insufficient institutional capacity for independent auditing, and the lack of unified impact assessment procedures or public registers, which the AI Act explicitly seeks to address through its governance structure and mandatory risk management provisions \parencites[10]{2977039}.}\pdfcomment[icon=Note]{Sources list gaps such as weak oversight, limited institutional auditing capacity, and lack of unified AIA/DPIA procedures or public registers, and state the AI Act's governance and risk-management aims to address such gaps (2977039 p.10; 8589961 p.3--4).} \hllightgreen{Standardisation activity interacts with this legal architecture.}\pdfcomment[icon=Note]{The documents discuss interaction between standardisation and legal frameworks (ISO/IEC work and the AI Act), so saying standardisation activity interacts with legal architecture is a supported synthesis (8589961 p.3--5).} \hldarkgreen{ISO/IEC JTC 1/SC 42 has issued AI related standards on risk management, bias mitigation, and ethical considerations, such as ISO/IEC 23894, 24027, 24028, and 24368.}\pdfcomment[icon=Note]{The ISO/IEC JTC 1/SC 42 committee and the cited standards (ISO/IEC 23894, 24027, 24028, 24368) are explicitly mentioned in the sources (8589961 p.3; 8589961 p.5).} \hldarkgreen{Analyses of their alignment with the AI Act show strong conceptual convergence around bias, transparency and human oversight in high risk systems, yet stress that ISO norms remain voluntary and therefore rely on binding instruments like the AI Act to gain enforcement power.}\pdfcomment[icon=Note]{Analyses note conceptual convergence between ISO norms and the AI Act on bias, transparency and human oversight, while also emphasizing that ISO standards are voluntary and lack binding enforcement (8589961 p.4--5; 8589961 p.3).} \hldarkgreen{Within the EU, these standards can function as de facto benchmarks for demonstrating conformity with requirements on representative datasets, bias mitigation, and explainable operation, especially where national authorities lack detailed technical criteria of their own \parencites[4]{8589961}[5]{8589961}.}\pdfcomment[icon=Note]{Sources argue that within the EU voluntary ISO standards can serve as practical benchmarks for conformity (e.g., representative datasets, bias mitigation) where national authorities lack detailed technical criteria (8589961 p.4--5; 1456534 p.2).} \hllightgreen{A concise comparison clarifies the EU's distinctive position: \begin{table}[h] \centering \caption{Selected AI risk governance features by region} \begin{tabular}{lccc} \hline Feature & EU AI Act & US federal & China \\ \hline Legal form & Binding regulation & Soft guidance & Administrative regs \\ Risk tiers & 4 level scheme & None unified & Domain specific \\ High risk duties & Conformity, oversight & Sectoral & State centric \\ Audit capacity & Mandated, still weak & Fragmented & Centralised \\ \hline \end{tabular} \end{table} Research on enforcement underscores that even in the EU, institutional capacity to audit complex models is a critical bottleneck.}\pdfcomment[icon=Note]{Comparative material in the sources supports the table's high-level contrasts (binding regulation in EU vs softer US guidance vs China's administrative regs) and research noting limited audit capacity in the EU (2977039 p.10; 1456534 p.8).} \hldarkgreen{Machine learning systems often lack explainability, which complicates assessment of fairness and transparency obligations in practice.}\pdfcomment[icon=Note]{The lack of explainability in many ML systems and its complication of fairness/transparency assessment is documented in the literature (7461574 p.9; 6362702 p.2).} \hldarkgreen{Regulators confront black box architectures where decision logic is difficult to interpret, and uneven national resources risk fragmented implementation of AI Act duties, echoing earlier experience under the GDPR.}\pdfcomment[icon=Note]{Sources discuss regulators facing black-box architectures and uneven national resources risking fragmented implementation, with parallels to GDPR experience (2977039 p.10; 1456534 p.2; 7461574 p.9).} \hldarkgreen{Commentators therefore call for specialised algorithmic auditing units, improved technical tooling, and closer collaboration with industry and civil society to translate risk based legal categories into operational audit routines \parencites[9]{7461574}[2]{1456534}.}\pdfcomment[icon=Note]{Commentators' calls for specialised algorithmic auditing units, better tooling, and multi-stakeholder collaboration are noted in the sources (2977039 p.10; 1456534 p.2), matching the cited reference (1456534 p.2).} \hllightgreen{Interaction between the AI Act and adjacent instruments is highly relevant for social recommender systems.}\pdfcomment[icon=Note]{Sources indicate interactions between the AI Act and adjacent instruments (DSA/DMA) affect recommender systems, so stating the interaction is highly relevant is a supported synthesis (6502223 p.12; 1456534 p.2).} \hldarkgreen{Platform regulation through the Digital Services Act already requires very large online platforms to assess and mitigate systemic risks to civic discourse and electoral processes and to undergo annual independent audits.}\pdfcomment[icon=Note]{The DSA's obligations on VLOPs to assess/mitigate systemic risks to civic discourse/elections and to undergo annual independent audits are explicitly described in the sources (6502223 p.12).} \hldarkgreen{The AI Act adds technology focused obligations on transparency, bias mitigation, documentation and human oversight for high risk systems, while the Digital Markets Act complements these with ex ante competition duties for gatekeepers.}\pdfcomment[icon=Note]{Sources describe the AI Act's technology-focused obligations (transparency, bias mitigation, documentation, human oversight) and the DMA's ex ante competition duties for gatekeepers (8589961 p.4; 1456534 p.8; 6502223 p.12).} \hldarkgreen{Together, these instruments extend algorithmic accountability from individual rights and content layers to infrastructural questions of market power and systemic risk.}\pdfcomment[icon=Note]{The combined effect of the AI Act, DSA and DMA expanding accountability from individual/content rights to infrastructural market and systemic risk is asserted in the sources (1456534 p.8; 6502223 p.12).} \hldarkgreen{EU level analysis suggests that, through the "Brussels Effect," this combined regime may function as a global benchmark for AI risk governance, even though enforcement capacity, innovation concerns, and geopolitical contestation remain unresolved \parencites[12]{6502223}[2]{1456534}.}\pdfcomment[icon=Note]{The sources discuss the Brussels Effect and note the EU regime may set global benchmarks while enforcement, innovation, and geopolitical issues remain unresolved (1456534 p.2--3; 8589961 p.4).} \hldarkgreen{Empirical work on educational AI systems illustrates why those enforcement and auditing concerns matter in concrete settings: algorithms trained on skewed school data can systematically underrate students from lower-income backgrounds and, when labels are exposed to teachers, can reduce teacher engagement by a material margin, producing what researchers have described as a digital self-fulfilling prophecy \parencites[2]{6362702}.}\pdfcomment[icon=Note]{Empirical studies cited report algorithms underrating students from lower-income backgrounds and teacher engagement dropping (e.g., 32\% reduction) after labels are exposed, described as a digital self-fulfilling prophecy (6362702 p.2--3; 2616018 p.17).} \hllightgreen{Such findings suggest that audit frameworks and representative dataset requirements in the AI Act are not merely theoretical safeguards but respond to observable harms in domains like education that are likely to reappear in other high risk applications \parencites[2]{6362702}.}\pdfcomment[icon=Note]{Given documented harms in education and the AI Act's dataset/audit requirements, it is a reasonable synthesis to say those legal measures respond to observable harms likely to recur in other high-risk domains (6362702 p.2; 6502223 p.12; 2977039 p.10).}

\subsubsection{Policy Trajectory toward a Digital Fairness Regulation}
\label{sec:policy_trajectory_digital_fairness} \hllightgreen{Regulatory discussions in the EU increasingly converge on the need for an explicit "fair design" layer that targets addictive interfaces and unfair personalisation more directly than existing data protection and platform rules can achieve.}\pdfcomment[icon=Note]{Sources describe EU debate about new rules addressing dark patterns/addictive interfaces beyond existing data protection and platform rules (e.g., 7896734; 1456534), so this is a reasonable synthesis of that literature.} \hlorange{Sanikidze describes a planned Digital Fairness Act (DFA), scheduled as a proposal for late 2026, whose stated goal is to regulate addictive product designs and unfair personalisation "across the board," filling gaps between the Digital Services Act (DSA) and traditional consumer protection instruments such as the Unfair Commercial Practices Directive \parencites[390]{7896734}.}\pdfcomment[icon=Note]{While the planned Digital Fairness Act and its late-2026 timeline and aims are documented (7896734), the specific attribution to "Sanikidze" as the describer is not supported by the provided sources.} \hldarkgreen{This trajectory emerges against a backdrop where the DSA already bans manipulative interfaces that impair free and informed decision making and imposes obligations on very large online platforms to provide at least one non profiled recommendation option, conduct annual independent audits, and assess systemic risks to civic discourse and electoral processes.}\pdfcomment[icon=Note]{The DSA obligations cited (ban on manipulative interfaces impairing informed decision-making; VLOP requirement to offer a non-profiled recommendation, annual independent audits, and systemic risk assessments for civic discourse/elections) are explicitly listed in 6502223 (p.12) and related sources.} \hldarkgreen{The AI Act complements this by prohibiting AI systems that exploit psychological vulnerabilities and by banning subliminal manipulation or manipulative effects that can significantly harm behaviour, yet current classifications do not explicitly treat social recommender algorithms as high risk, which critics regard as a democratic protection gap the DFA is expected to address \parencites[2248]{6502223}.}\pdfcomment[icon=Note]{The EU AI Act's prohibitions on systems exploiting psychological vulnerabilities and subliminal manipulation, and the point that recommenders are not explicitly classed as high-risk (a critic-identified gap), are stated in 6502223 and 6502223/8589961 (see table and discussion).} \hllightgreen{Policy momentum is reinforced by enforcement practice under unfair commercial practices law in EU aligned jurisdictions.}\pdfcomment[icon=Note]{The claim that policy momentum is reinforced by enforcement under unfair commercial practices law is supported by examples of enforcement in EU-aligned jurisdictions (e.g., Georgia) in 7896734, so this is a reasonable inference.} \hldarkgreen{Georgia's Law on the Protection of Consumer Rights, which transposes EU standards, has already been used by the Georgian Competition and Consumer Agency to sanction digital dark patterns such as false countdown timers, fabricated scarcity, fake reviews, and misleading brand associations exemplified by advertising under the "Maxbosch" label while using BOSCH imagery.}\pdfcomment[icon=Note]{Source 7896734 (p.7) documents Georgia's 2022 consumer law, GCCA enforcement actions against countdown timers, fabricated scarcity, fake reviews, and the Maxbosch/BOSCH imagery example.} \hldarkgreen{These decisions show that a general unfair practices framework can be applied to technologically sophisticated interface manipulations even in the absence of AI specific clauses.}\pdfcomment[icon=Note]{7896734 (p.7) explicitly states that these Georgian cases illustrate how a general unfair practices framework can address technologically sophisticated interface manipulations without AI-specific clauses.} \hllightgreen{At the same time, Sanikidze stresses that recent Georgian cases did not yet cover AI driven dark patterns, which indicates both the reach and the current limits of consumer law when faced with adaptive, data driven designs \parencites[390]{7896734}.}\pdfcomment[icon=Note]{7896734 (p.7) notes that the Georgian cases did not cover AI-driven dark patterns, supporting the claim about the reach and limits of consumer law; the sentence is a reasonable restatement of that point.} \hllightgreen{Within the EU, the emerging trajectory can be read as a layering strategy.}\pdfcomment[icon=Note]{Multiple sources (e.g., 1456534, 6502223, 7896734) describe an ensemble of EU instruments forming a layered regulatory approach, so reading the trajectory as a "layering strategy" is a supported synthesis.} \hldarkgreen{Existing pillars include the GDPR on data privacy, the DSA on platform accountability, and the Digital Markets Act (DMA) on gatekeeper competition duties such as bans on self preferencing and requirements for interoperability and explicit consent before combining personal data across services \parencites[2248]{6502223}[8]{1456534}.}\pdfcomment[icon=Note]{1456534 (p.8) and 6502223 (p.12) list GDPR, DSA and DMA as pillars and describe DMA duties including bans on self-preferencing, interoperability requirements, and consent before combining personal data.} \hldarkgreen{The AI Act then adds a cross sectoral, risk based scheme that bans unacceptable uses like social scoring and systems that exploit psychological vulnerabilities, while imposing conformity assessment, documentation, and human oversight requirements on high risk applications in domains such as education and employment \parencites[2248]{6502223}.}\pdfcomment[icon=Note]{6502223 and 8589961 describe the AI Act as a cross-sectoral, risk-based scheme banning social scoring and exploitation of psychological vulnerabilities, and imposing conformity assessment, documentation, and human oversight for high-risk domains like education and employment.} \hllightgreen{Analysis by Kustura characterises this ensemble as the most ambitious set of algorithmic regulations currently in force, with the "Brussels Effect" likely to externalise its standards globally, though uneven enforcement capacity and contested geopolitical conditions constrain this ambition \parencites[2]{1456534}.}\pdfcomment[icon=Note]{1456534 characterises the EU ensemble as highly ambitious and discusses the "Brussels Effect" while noting enforcement and geopolitical constraints, making this a supported characterization and synthesis.} \hllightgreen{Against this regulatory stack, the DFA is framed not as an isolated statute but as an additional, preventive tier that targets design practices which fall between content moderation and data processing, especially addictive UX, unfair personalisation, and AI mediated dark patterns.}\pdfcomment[icon=Note]{7896734 describes the DFA as an additional preventive tier aimed at design practices between content moderation and data processing (addictive UX, unfair personalisation), so this framing is supported by the source.} \hllightgreen{Sanikidze explicitly links the DFA's focus on addictive product designs and unfair personalisation to concerns about AI driven manipulation and dark patterns in e commerce.}\pdfcomment[icon=Note]{7896734 links the DFA's focus on addictive designs and unfair personalisation to concerns about AI-driven manipulation and dark patterns in e-commerce, so this attribution is supported as a synthesis.} \hldarkgreen{Recent enforcement against platform X and Temu under the DSA and UCPD, including fines for "deceptive" verification features and preliminary findings on potentially addictive design, are interpreted as early signals of regulatory intolerance for interface pressure that makes cancellation harder than sign up or obscures material information through countdown timers and scarcity messaging \parencites[390]{7896734}.}\pdfcomment[icon=Note]{7896734 (p.7 and p.390) documents enforcement actions (e.g., platform X €120m fine, Temu investigations) and discusses findings about deceptive verification features and potentially addictive designs like countdown timers and scarcity messaging.} \hllightgreen{The DFA is consequently presented as a move from case by case unfair practice enforcement to an ex ante, horizontally applicable regime where fairness in design becomes a primary regulatory objective in its own right, alongside safety, competition, and data protection.}\pdfcomment[icon=Note]{Sources (7896734; 1456534) describe the DFA as a preventive, horizontal measure addressing design practices beyond case-by-case UCPD enforcement, so this is a supported inference.} \hldarkgreen{Research on behavioural advertising underlines why such a regime is seen as necessary: personalised online campaigns rely on data flows across publisher ad servers, supply side and demand side platforms, and data management platforms that segment users through cookies, allowing advertising to shift from mass oriented communication to highly individualised influence strategies \parencites[5]{6941611}.}\pdfcomment[icon=Note]{6941611 (p.5) details the behavioural advertising ecosystem (publisher ad servers, SSPs, DSPs, DMPs, cookies) and explains how personalization enables individualized influence, directly supporting the claim.} \hllightgreen{When these strategies are embedded in opaque recommendation and interface logic, general transparency and consent rules may be insufficient to prevent exploitation of cognitive vulnerabilities.}\pdfcomment[icon=Note]{6941611 discusses opacity of recommendation/interface logic and the limits of transparency/consent alone; combined with AI Act concerns, it supports the inference that these rules may be insufficient to prevent exploitation of cognitive vulnerabilities.} \hllightgreen{For this thesis, which treats the DFA as a policy trajectory rather than settled law, the practical implication is clear.}\pdfcomment[icon=Note]{This is an interpretive claim about the thesis stance; given the cited sources present the DFA as a trajectory rather than settled law (e.g., 7896734), treating it as such is a reasonable scholarly positioning.} \hllightgreen{Any Ethical AI Audit Toolkit and pro social pivot strategy must be aligned ex ante with this direction of travel, not only with the letter of current DSA and AI Act provisions.}\pdfcomment[icon=Note]{Auditing literature (2574631) and regulatory discussions (6502223, 7896734) support the recommendation that audit toolkits should align proactively with emerging DFA/DSA/AI Act directions, not only current letter of law.} \hllightgreen{This implies incorporating explicit audit dimensions on addictive design patterns, unfair personalisation, and AI enabled dark patterns into the coding matrix used to assess social recommender systems.}\pdfcomment[icon=Note]{Given sources call for auditability of dark patterns, unfair personalisation and AI-related manipulations (7896734; 2574631), incorporating these dimensions into assessment matrices is a reasonable, supported implication.} \hllightgreen{It also suggests that future compliance will likely depend on demonstrating that engagement optimisation does not rely on manipulative UX mechanics, that profiling based targeting can be meaningfully switched off, and that audit trails and explanations are available for both content ranking and advertising decisions within social feeds \parencites[2248]{6502223}[390]{7896734}[5]{6941611}.}\pdfcomment[icon=Note]{7896734 (p.390) and 6941611 describe regulatory expectations (non-profiled recommendations, limits on manipulative UX) and the advertising ecosystem, supporting the suggestion that compliance will require demonstrating non-manipulative optimisation, switch-off of profiling, and audit trails.} \hldarkgreen{At the same time, the literature on AI auditing cautions that audits vary in scope and independence and that provider-commissioned reviews may not suffice for regulatory assurance; structured, independent assessments of system rationale, code, and impacts are therefore likely to become an important element of credible compliance \parencites[13]{2574631}.}\pdfcomment[icon=Note]{The AI-auditing literature (2574631, p.13; 7461574) explicitly cautions that audits vary in scope/independence and provider-commissioned reviews may be insufficient, supporting the claim about the need for structured independent assessments.} \hllightgreen{This suggests auditors and regulators will need to look beyond declarative documentation and seek evidence of substantive, replicable testing and impact measurement when evaluating claims about non-manipulative engagement design \parencites[13]{2574631}.}\pdfcomment[icon=Note]{2574631 and framework-oriented sources (8367105) argue auditors should seek substantive testing and impact measurement beyond declarative documentation, making this a supported recommendation.}

\subsection{Transparency, Data Access and Audit Obligations}

\subsubsection{Platform Transparency Reporting Requirements}
\label{sec:platform_transparency_reporting} \hllightgreen{Transparency reporting in EU platform regulation functions as a central control surface for recommender accountability rather than a peripheral disclosure duty.}\pdfcomment[icon=Note]{Sources on the DSA and platform governance characterise transparency reporting as central to accountability rather than mere disclosure (Meßmer \& Degeling 3195574; summaries in 1456534 and 6502223), so this is a reasoned synthesis of those claims.} \hldarkgreen{Under the Digital Services Act, very large online platforms with more than $45$ million monthly EU users must publish recurring reports that specify, among other items, the number of disputes submitted, the number of suspensions, and staffing levels for content moderation, including qualifications and linguistic expertise \parencites[19]{3195574}[7]{1456534}.}\pdfcomment[icon=Note]{Meßmer \& Degeling explicitly state DSA transparency reporting (Art.24/Art.42) includes number of disputes, suspensions, and moderator staffing/qualifications (3195574 p.19), consistent with the citation (1456534 summarises the same obligations).} \hllightgreen{These reports are designed to transform traditionally opaque operational metrics into externally inspectable indicators that supervisors, auditors, and civil society can use when evaluating the behaviour of recommender systems embedded in social platforms.}\pdfcomment[icon=Note]{The sources describe transforming opaque operational metrics into inspectable transparency and data‑access rights usable by auditors, supervisors and civil society (3195574; 2977039), so this is a reasonable synthesis/inference.} \hldarkgreen{Crucially, transparency duties do not stop at aggregate moderation statistics.}\pdfcomment[icon=Note]{Meßmer \& Degeling note the DSA's transparency obligations go beyond aggregate moderation stats to include recommender and other disclosures (3195574 p.19), directly supporting this claim.} \hldarkgreen{Article 27 DSA requires disclosure of "the main parameters used for recommendations," compelling platforms to explain how ranking is organised and which signals drive exposure.}\pdfcomment[icon=Note]{Meßmer \& Degeling explicitly cite Article 27 requiring transparency about 'the main parameters used for recommendations' (3195574 p.19), supporting the wording about ranking signals and organisation.} \hldarkgreen{In parallel, Article 39 mandates an advertising repository for VLOPs that details paid content, while Article 40 defines data access procedures for Digital Services Coordinators and the European Commission.}\pdfcomment[icon=Note]{Meßmer \& Degeling list Article 39 (advertising repository) and Article 40 (data access procedures for Digital Services Coordinators and Commission) as DSA provisions (3195574 p.19), directly supporting this sentence.} \hldarkgreen{Analysis by Me{\ss}mer and Degeling stresses that civil society actors are already calling for additional obligations so that users can act meaningfully on this published information rather than merely observe it \parencites[19]{3195574}.}\pdfcomment[icon=Note]{Meßmer \& Degeling report that civil society asks for additional obligations so users can act on published information rather than merely observe it (3195574 p.19), directly supporting this statement.} \hllightgreen{From an Ethical AI Audit perspective, these provisions supply a minimum disclosure baseline against which internal reporting dashboards should be calibrated.}\pdfcomment[icon=Note]{Framing the DSA provisions as a minimum disclosure baseline for Ethical AI Audits is a reasonable inference from the DSA/AI governance literature (3195574; 2977039) rather than a verbatim claim, so LIGHTGREEN is appropriate.} \hllightgreen{Table~\ref{tab:dsa_transparency_dims} summarises the main transparency dimensions and their audit relevance. \begin{table}[h] \centering \caption{Key DSA transparency reporting elements and audit implications} \label{tab:dsa_transparency_dims} \begin{tabular}{p{3.1cm}p{4.2cm}p{5.0cm}} \hline Regulatory hook & Required disclosure & Relevance for Ethical AI Audit \\ \hline Art. 24 reports & Disputes, suspensions & Indicator of downstream harms and appeal pressure linked to recommender outputs \\ Art. 27 parameters & Main recommendation signals & Entry point for interrogating engagement optimisation and profiling dependence \\ Art. 39 ad repository & Ad content and targeting info & Enables cross checking of content vs ad recommendation consistency \\ Art. 40 data access & Data sharing to regulators & Conditions for external, scenario based audits of systemic risks \\ \hline \end{tabular} \end{table} Recent work on auditing recommender systems interprets these clauses as the backbone of a multi actor audit regime.}\pdfcomment[icon=Note]{Meßmer \& Degeling map DSA clauses to auditing relevance (3195574 p.19), and the surrounding literature interprets these clauses as foundational for multi‑actor audits, so the summary/Table claim is a supported synthesis.} \hldarkgreen{Me{\ss}mer and Degeling distinguish internal platform assessments, independent audits commissioned by VLOPs, and external investigations by regulators or researchers, and propose a scenario based audit process that maps concrete risk narratives to specific transparency and data access rights in the DSA \parencites[6]{3195574}.}\pdfcomment[icon=Note]{Meßmer \& Degeling describe the distinction between internal assessments, independent audits commissioned by VLOPs, and external regulatory or research investigations and propose a scenario‑based audit process (3195574, esp. overview and methodology sections).} \hldarkgreen{Their analysis underlines that, without detailed implementing acts and voluntary standards under Article 44, platforms could perform highly partial self assessment while still appearing formally compliant in transparency reports \parencites[6]{3195574}[19]{3195574}.}\pdfcomment[icon=Note]{Meßmer \& Degeling warn that absent detailed implementing acts and voluntary standards (Article 44) platforms could perform partial self‑assessments while appearing compliant (3195574 p.19), directly supporting this claim.} \hllightgreen{Risk based AI governance trends outside the DSA point in the same direction.}\pdfcomment[icon=Note]{Multiple sources on AI governance (EU AI Act discussion 6502223; algorithmic impact assessment literature 5442231; AI Law Model 2977039) show a broader trend toward risk‑based governance, so this is a valid synthesis.} \hlorange{Sebastian anticipates that high risk domains such as credit scoring will be subject to algorithmic impact assessments that resemble environmental impact reviews, with standardised metrics and documentation structures.}\pdfcomment[icon=Note]{No provided source names or attributes an anticipation to 'Sebastian'; while sources (e.g., 5442231) discuss algorithmic impact assessments for credit, the specific attribution to 'Sebastian' is unsupported by the supplied materials.} \hldarkgreen{Institutions are expected to disclose model types, training data properties, validation methods, and performance metrics by demographic group, together with governance arrangements and complaint channels, in transparency reports that enable both regulators and the public to compare fairness performance across entities.}\pdfcomment[icon=Note]{The referenced literature on algorithmic impact assessments and transparency (5442231 p.15; 2977039) explicitly lists disclosure of model types, training data properties, validation methods and performance metrics by demographic group and governance arrangements.} \hldarkgreen{Continuous documentation is emphasised because model drift and retraining can introduce new discrimination mechanisms; transparency is therefore framed as an ongoing monitoring practice rather than a one off filing \parencites[15]{5442231}.}\pdfcomment[icon=Note]{Sources emphasize continuous documentation because model drift and retraining can introduce new discrimination mechanisms and thus require ongoing monitoring rather than one‑off filings (5442231; 6502223), directly supporting this sentence.} \hllightgreen{Analogous expectations emerge in AI education and governance scholarship.}\pdfcomment[icon=Note]{The corpus includes AI education and governance scholarship advocating transparency and ongoing monitoring (e.g., 2977039; 6539584), so stating analogous expectations emerge in that scholarship is a supported synthesis.} \hldarkgreen{HITL grading research argues that automated evaluation systems must be auditable, interpretable, and subordinate to human authority, with clear documentation that accrediting bodies can inspect \parencites[15]{6805789}.}\pdfcomment[icon=Note]{The HITL grading literature cited (6805789 p.15) asserts that automated evaluation systems should be auditable, interpretable, subordinate to human authority, and documented for accrediting bodies, matching this claim.} \hldarkgreen{Kostenko's AI Law Model generalises this into a broad transparency principle: mandatory disclosure of developers, data sources, algorithmic architectures, intended purposes, limitations, statistical error probabilities, and high uncertainty scenarios, combined with audit logs of decisions and mandatory internal and external audits with access to full technical dossiers.}\pdfcomment[icon=Note]{The AI Law Model (2977039 p.69--70) explicitly lists mandatory disclosure of developers, data sources, architectures, purposes, error probabilities, audit logs, and mandatory internal/external audits with access to technical dossiers, directly supporting this description.} \hldarkgreen{Any covert or manipulative use of AI with legal, social, economic or political implications is categorised as unacceptable and subject to liability \parencites[70]{2977039}.}\pdfcomment[icon=Note]{The AI Law Model states covert or manipulative uses of AI with significant implications are unacceptable and entail liability (2977039 p.70--71), directly supporting this sentence.} \hllightgreen{These formulations provide a normative template for what meaningful transparency reporting should contain when platforms claim alignment with ethical AI norms.}\pdfcomment[icon=Note]{Given the prescriptive items in 2977039 and related governance literature, it is a reasonable scholarly synthesis to say these formulations provide a normative template for meaningful transparency reporting (2977039; 3195574).} \hllightgreen{Given this convergence, a pro social audit toolkit for social recommender systems can treat DSA reporting as a floor and not a ceiling.}\pdfcomment[icon=Note]{Treating DSA reporting as a floor (not a ceiling) for pro‑social audit toolkits is a reasoned inference from convergence across DSA obligations and broader AI governance recommendations (3195574; 2977039; 5442231).} \hllightgreen{Internally, management reports would be expected to mirror or exceed DSA fields, for instance by including: distribution of recommendation exposure by content category, incidence of repeat exposure to flagged content, explanation coverage rates for "why am I seeing this" interfaces, and the proportion of traffic attributable to non profiled recommendation options mandated for VLOPs \parencites[12]{6502223}[19]{3195574}.}\pdfcomment[icon=Note]{Proposals that internal management reports should mirror or exceed DSA fields and include the listed metrics (exposure distributions, repeat exposures, explanation coverage, non‑profiled traffic) are consistent with Meßmer \& Degeling's emphasis on richer transparency and DSA requirements for non‑profiling alternatives (3195574) and are a supported recommendation.} \hllightgreen{External facing transparency statements should then translate these metrics into accessible narratives that satisfy Kostenko's requirements for user friendly social explanation while still enabling regulators and researchers to reconstruct how engagement objectives interact with fairness, sustainability, and digital ecology obligations \parencites[70]{2977039}[119]{2977039}.}\pdfcomment[icon=Note]{The AI Law Model calls for user‑friendly social explanations and machine‑readable, accessible documentation (2977039 p.69--70, p.118--119), and the sentence reasonably synthesizes that with the need to enable regulator/researcher reconstruction of engagement vs fairness/sustainability tradeoffs.}

\subsubsection{Data Access for Regulators and Independent Researchers}
\label{sec:data_access_regulators_researchers} \hllightgreen{Access to platform data is the practical hinge on which systemic auditing of social recommender systems turns.}\pdfcomment[icon=Note]{Works on platform oversight and recommender auditing emphasize that access to platform data is foundational for systemic audits (see Meßmer \& Degeling, Source 3195574 p.15; 6502223 p.12), so this is a reasonable synthesis inference.} \hldarkgreen{Under the Digital Services Act, very large online platforms must grant data access to Digital Services Coordinators and the European Commission in order to enable scrutiny of systemic risks and compliance with due diligence obligations.}\pdfcomment[icon=Note]{The DSA requires VLOPs to grant data access to Digital Services Coordinators and the Commission for scrutiny of systemic risks and due diligence ,  see Meßmer \& Degeling (Source 3195574 p.19) and the DSA summary (Source 1456534 p.7; 6502223 p.12).} \hldarkgreen{Article 40 DSA is singled out as the clause that describes how such access is to be organised, sitting alongside Article 24 on transparency reporting and Article 27 on disclosure of "the main parameters used for recommendations" \parencites[19]{3195574}.}\pdfcomment[icon=Note]{Meßmer \& Degeling explicitly identify Article 40 as the data access provision and discuss it alongside Article 24 (transparency reporting) and Article 27 (parameters of recommendations) (Source 3195574 p.19).} \hllightgreen{This statutory access is designed to support supervisory investigations that go beyond high level reports, for example by allowing regulators to reconstruct how a recommender behaved in a defined period or context and how internal mitigation measures affected observable outputs \parencites[19]{3195574}[7]{1456534}.}\pdfcomment[icon=Note]{Sources describe Article 40 enabling regulator access for deeper supervisory investigation and audit (Source 3195574 p.19) and the DSA's systemic-risk/audit framing (Source 1456534 p.7), so inferring reconstruction of recommender behaviour and assessment of mitigation is a supported application of those provisions.} \hldarkgreen{Independent research access is treated as a parallel requirement rather than an optional extra.}\pdfcomment[icon=Note]{The DSA explicitly includes research access (vetted researchers) as a required element for VLOPs rather than an optional measure (Source 6502223 p.12; Meßmer \& Degeling Source 3195574 p.19).} \hldarkgreen{Platforms qualifying as very large online platforms must provide "vetted researchers" with data sufficient to study systemic risks, including effects on civic discourse and electoral processes \parencites[12]{6502223}.}\pdfcomment[icon=Note]{The DSA requires VLOPs to provide vetted researchers with access sufficient to study systemic risks, including civic discourse and electoral process effects (Source 6502223 p.12; also reflected in Meßmer \& Degeling, Source 3195574 p.19).} \hldarkgreen{Governance proposals for democratic recommender systems build on this by arguing that platforms above a size threshold should be obliged to supply representative samples of content recommendation sequences, so that external researchers can empirically examine amplification asymmetries, misinformation dynamics, and diversity performance under real usage conditions \parencites[16]{6502223}.}\pdfcomment[icon=Note]{Governance proposals in the democratic recommender literature call for platforms to supply representative samples of recommendation sequences for empirical study of amplification, misinformation dynamics, and diversity (Source 6502223 p.16; see also policy recommendations in Source 6502223 p.16).} \hldarkgreen{Me{\ss}mer and Degeling translate these legal hooks into a scenario based audit model in which specific risk narratives, such as eating disorder promotion or suicidal ideation content, are mapped to concrete data access requests grounded in Article 40 DSA \parencites[15]{3195574}.}\pdfcomment[icon=Note]{Meßmer and Degeling discuss mapping specific risk narratives (e.g., eating disorder promotion, suicidal content) to concrete data access needs and link such requests to Article 40 DSA obligations (Source 3195574 p.15 and p.19).} \hldarkgreen{In this view, data access is not generic but tightly coupled to platform specific risk hypotheses and product configurations.}\pdfcomment[icon=Note]{Meßmer \& Degeling argue that data access must be tailored to platform-specific risk hypotheses and product configurations rather than generic access (Source 3195574 p.15).} \hldarkgreen{Kostenko's AI Law Model generalises such access beyond EU platform regulation into a cross border accountability principle.}\pdfcomment[icon=Note]{The AI Law Model (Kostenko) explicitly generalizes access and accountability beyond EU regulation into a cross-border accountability principle (Source 2977039 p.69.1 and p.126).} \hldarkgreen{Any AI system that causes consequences on the territory of a State or purposefully interacts with its market is said to fall under that State's "extraterritorial digital jurisdiction".}\pdfcomment[icon=Note]{The AI Law Model establishes an extraterritorial digital jurisdiction principle: AI systems causing consequences on a State's territory or purposefully interacting with its market fall under that State's jurisdiction (Source 2977039 p.249; p.126).} \hldarkgreen{Cross border liability then entails obligations to recognise the jurisdiction of courts and regulators, to comply with mandatory norms on transparency, explainability, human oversight, secure data processing, and to ensure logging and preservation of evidence.}\pdfcomment[icon=Note]{The AI Law Model describes cross-border liability obligations including recognition of jurisdiction, compliance with norms on transparency, explainability, human oversight, secure data processing, and requirements to ensure logging and preservation of evidence (Source 2977039 p.249).} \hldarkgreen{Subjects must report incidents within legally defined deadlines, cooperate during audits, provide effective remedies and redress to local users, and implement "flowdown" clauses that transmit these duties through supply chains \parencites[249]{2977039}.}\pdfcomment[icon=Note]{The AI Law Model mandates incident reporting within legal deadlines, cooperation during audits, provision of remedies/redress to local users, and contractual "flowdown" obligations through supply chains (Source 2977039 p.249).} \hldarkgreen{This model treats event logs and technical documentation as part of a legally enforceable evidentiary regime, not just an internal governance resource.}\pdfcomment[icon=Note]{The AI Law Model treats event logs and technical documentation as part of an evidentiary regime, requiring preservation and use as evidence rather than mere internal governance artefacts (Source 2977039 p.246; p.249).} \hldarkgreen{Supervision and auditing duties in the same AI Law Model reinforce this orientation.}\pdfcomment[icon=Note]{The AI Law Model includes supervision and auditing duties (central authority, audits, reporting) that reinforce the evidentiary and compliance orientation (Source 2977039 p.126; p.246-247).} \hldarkgreen{A central executive body maintains a national register of high risk systems, organises scheduled and unscheduled audits, and issues mandatory instructions.}\pdfcomment[icon=Note]{The AI Law Model specifies a central executive authority maintaining registers, organising scheduled and unscheduled audits, and issuing mandatory instructions (Source 2977039 p.127; p.246-247).} \hldarkgreen{For such systems, annual audits and public trust labels are required, and open API access to registry attributes is provided for public control and scientific research.}\pdfcomment[icon=Note]{The model requires annual audits and public trust labelling for high-risk systems and provides API access to open registry attributes for public control and scientific research (Source 2977039 p.246-247, section 68.7).} \hldarkgreen{In cross border digital environments described as "metaverse" or multi domain platforms, admission conditions explicitly include compatibility of unique identifiers and algorithmic passports with environment registries and the activation of logging mechanisms that ensure the evidentiary value of decisions.}\pdfcomment[icon=Note]{The AI Law Model's provisions for cross-border digital environments require compatibility of unique identifiers and algorithmic passports with environment registries and activation of logging mechanisms to ensure evidentiary value (Source 2977039 p.246, section 68.6).} \hldarkgreen{Dispute resolution rests on the recognised evidentiary status of logs, which must be stored for legally defined periods \parencites[246]{2977039}.}\pdfcomment[icon=Note]{The model bases dispute resolution on the evidentiary value of logs and requires storing logs for legally specified periods (Source 2977039 p.246, section 68.6).} \hllightgreen{These provisions underscore that meaningful access for regulators and researchers presupposes structured, tamper evident logging architectures at system design time.}\pdfcomment[icon=Note]{Given the AI Law Model's emphasis on logging, preservation, and evidentiary use (Source 2977039 p.246) and design-for-governance arguments in technical frameworks (Source 8367105 p.25), it is a reasonable synthesis to say meaningful access presupposes structured, tamper-evident logging architectures from design time.} \hllightgreen{A concise comparison between EU and AI Law Model approaches can be summarised as follows. \begin{table}[h] \centering \caption{Data access regimes for AI and platform oversight} \begin{tabular}{p{3cm}p{5cm}p{5cm}} \hline Dimension & DSA / EU platforms & AI Law Model (Kostenko) \\ \hline Primary addressee & VLOPs with $45+$ million users & All AI systems with impact on State territory \\ Legal hook & Art. 40 data access, vetted researchers & Extraterritorial digital jurisdiction, cross border liability \\ Data objects & Platform data for systemic risk research; ads repository (Art. 39) & Technical materials, logs, audit results, algorithmic passports \\ Research access & Vetted researchers for systemic risk studies & API access to registry attributes for scientific research \\ Logging requirement & Implied through audit and risk assessment duties & Explicit event logging, evidentiary preservation, log based arbitration \\ \hline \end{tabular} \end{table} Evidence from AI governance and auditing practice suggests that even when formal access rights exist, enforcement capacity and methodological tools remain limiting factors.}\pdfcomment[icon=Note]{The comparative characterisation draws on DSA descriptions (Sources 6502223 p.12; 3195574 p.19) and the AI Law Model (Source 2977039 p.246), and the concluding claim that enforcement capacity and methodological tools limit effectiveness is supported by governance literature noting enforcement and capacity issues (Source 1456534 p.7; 7461574 p.9).} \hldarkgreen{Scholars highlight that regulators struggle with black box models where decision processes are not readily interpretable, which complicates assessment of fairness and transparency obligations under instruments like the EU AI Act and GDPR \parencites[9]{7461574}.}\pdfcomment[icon=Note]{Multiple sources note regulators struggle with black-box models and interpretability challenges complicating fairness/transparency assessment under instruments like the EU AI Act and GDPR (Source 7461574 p.9; Source 1456534 p.7).} \hldarkgreen{Me{\ss}mer and Degeling anticipate a fragmented landscape of mandatory and voluntary audits under the DSA and propose voluntary standards and codes of conduct to systematise audit and access practices across platforms \parencites[19]{3195574}.}\pdfcomment[icon=Note]{Meßmer and Degeling discuss a fragmented audit landscape under the DSA and explicitly propose voluntary standards and codes of conduct to systematise audit and access practices (Source 3195574 p.19).} \hldarkgreen{Research on AI enhanced auditing in financial and corporate contexts adds that human AI collaborative methods, where AI performs large scale analytics and human experts exercise contextual judgement and ethical scrutiny, can improve audit quality but require carefully designed interfaces and cognitive load management for auditors \parencites[3]{6539584}[6]{6539584}.}\pdfcomment[icon=Note]{Research on AI-enhanced auditing supports human-AI collaborative methods improving audit quality while requiring careful interface design and management of cognitive load for auditors (Source 6539584 p.3--6).} \hllightgreen{These observations indicate that regulatory and research access to recommender data is a necessary but insufficient condition; specialised expertise, standardised audit protocols, and ethically informed human oversight are also required to turn raw data into credible accountability findings \parencites[9]{7461574}[19]{3195574}[3]{6539584}.}\pdfcomment[icon=Note]{The conclusion that data access is necessary but insufficient and that expertise, standardised protocols, and ethically informed human oversight are required is a supported synthesis of Meßmer \& Degeling on standards (Source 3195574 p.19) and human-AI audit literature (Source 6539584 p.3--6).}

\subsubsection{Documentation, Logging and Traceability Expectations}
\label{sec:doc_logging_traceability} \hldarkgreen{Traceability expectations in emerging AI regulation extend far beyond general calls for transparency and require specific documentation, logging, and access architectures.}\pdfcomment[icon=Note]{Source 2977039 (Ch.8, p.35--36) explicitly requires identification, traceability and documentation and frames transparency obligations as including technical documentation, logs and access architectures.} \hldarkgreen{Kostenko's AI Law Model formulates this explicitly as a principle of transparency and explainability: any AI system that can affect rights, security, or well being must be identifiable, traceable, explainable, and accountable as a condition of lawful use \parencites[35]{2977039}.}\pdfcomment[icon=Note]{The AI Law Model (Source 2977039, Ch.4 and Ch.8, pp.21, 35--36) explicitly frames systems affecting rights/security/well‑being as requiring identifiability, traceability, explainability and accountability as preconditions for lawful use.} \hldarkgreen{Entities developing or operating such systems are obliged to notify users about AI use before interaction, clearly mark AI outputs in public space, and maintain technical documentation that describes system architecture and training data sources.}\pdfcomment[icon=Note]{Source 2977039 (Ch.8, p.36) obliges subjects to notify users before interaction, ensure identification of AI in public spaces, and make technical documentation including training data sources available within limits.} \hldarkgreen{They must also provide interfaces that allow testing of AI decisions and indicate the degree of AI participation in decision making, for example whether a decision is automatic, semi automatic, or purely consultative \parencites[36]{2977039}.}\pdfcomment[icon=Note]{Source 2977039 (Ch.8, p.36) requires interfaces for testing AI solutions and that entities indicate the degree and nature of AI participation (automatic, semi‑automatic, consultative).} \hldarkgreen{Explainability in this model is operationalised as a logging and documentation problem.}\pdfcomment[icon=Note]{Source 2977039 (Ch.8, p.35--36) explicitly operationalises explainability via mechanisms for reproducing algorithmic logic, decision logging and documentation.} \hldarkgreen{Developers are required to create mechanisms that reproduce algorithmic logic, including decision logs, data processing pipelines, and model parameters.}\pdfcomment[icon=Note]{Source 2977039 (Ch.8, §8.3, p.36) requires developers to provide mechanisms to reproduce algorithmic logic, including decision logs, data processing processes and model parameters.} \hldarkgreen{Results must be interpretable so that users and independent auditors can identify errors, biases, or inconsistencies, and independent access interfaces for ethical and legal auditors must be available under conditions that preserve confidentiality and security.}\pdfcomment[icon=Note]{Source 2977039 (Ch.8, §8.3, p.36) states results must be interpretable to identify errors/biases and that independent access interfaces for ethical and legal auditors must be provided with confidentiality protections.} \hldarkgreen{Openness adds a public reporting duty: regular publication of impact reports, codes of conduct, and incident databases, with non compliance triggering sanctions such as suspension, fines, or revocation of certification, monitored by an authorised AI ethics and security body \parencites[35]{2977039}[36]{2977039}.}\pdfcomment[icon=Note]{Source 2977039 (Ch.8, p.35--36) requires regular public information on policies/risks, publication of impact reports, codes of conduct and incident databases and links non‑compliance to sanctions administered by the Authorized Body for AI Ethics and Security.} \hldarkgreen{Children's digital safety provisions raise the bar further.}\pdfcomment[icon=Note]{Source 2977039 (Ch.60, p.214--215) introduces the principle of digital safety of the child and establishes heightened requirements for systems interacting with minors, supporting the claim that provisions 'raise the bar further.'} \hldarkgreen{Under the principle of digital safety of the child, any AI system that processes children's data or influences decisions about them must be "secure by default" at all stages: architecture, algorithmic logic, data categories, training cycles, and purposes of use must be understandable, verifiable, and controllable.}\pdfcomment[icon=Note]{Source 2977039 (Ch.60, p.214) states that systems interacting with minors must be 'secure by default' across architecture, algorithmic logic, data categories, training cycles and purposes, and be understandable, verifiable and controllable.} \hldarkgreen{For regulators, this translates into an expectation of a complete technical, methodological, legal, and ethical dossier that includes architectural schemes, training protocols, threat models, risk analysis, event logs, and results of internal and external audits of child safety.}\pdfcomment[icon=Note]{Source 2977039 (Ch.60, p.215) explicitly lists the dossier expected by regulators: architectural schemes, training protocols, threat models, risk analyses, event logs and internal/external child safety audits.} \hldarkgreen{Decision logs must record key events, configuration changes, model updates, and human interventions.}\pdfcomment[icon=Note]{Source 2977039 (Ch.60, p.215) specifies that decision logs must record key events, configuration changes, model updates and human operator interventions.} \hldarkgreen{Absence of effective mechanisms for security and accountability, or use of manipulative or obscure interfaces in child impacting contexts, is classified as a violation of good governance and the best interests of the child and can justify suspension or termination of a system by the national AI regulator \parencites[214]{2977039}.}\pdfcomment[icon=Note]{Source 2977039 (Ch.60, p.215) states that absence of effective mechanisms or use of manipulative/obscure interfaces in child‑impacting contexts is a violation and the national AI regulator can suspend or terminate the system.} \hldarkgreen{Territorial sovereignty clauses in the same model transform traceability into a jurisdictional requirement.}\pdfcomment[icon=Note]{Source 2977039 (Ch.38, p.124) frames territorial sovereignty clauses that make registration, representation and compliance prerequisites for operation, turning traceability into a jurisdictional requirement.} \hldarkgreen{Only AI systems registered in a National Register of AI Systems and operating in full compliance with human rights, security, and non discrimination standards may function in a State's territory.}\pdfcomment[icon=Note]{Source 2977039 (Ch.38, p.124) states only systems with official representation and inclusion in the National Register, operating in compliance with human rights, security and non‑discrimination standards, may operate in the State's territory.} \hldarkgreen{Registration establishes legal status, identifies owners, and supports tracing of supply chains and software origin.}\pdfcomment[icon=Note]{Source 2977039 (Ch.38, p.124) explains registration establishes legal status, identifies owners, and enables tracing of supply chains and software origin.} \hldarkgreen{Official representation is required so that supervisory bodies can conduct audits, issue corrective orders, apply sanctions, and, where necessary, restrict or stop operation.}\pdfcomment[icon=Note]{Source 2977039 (Ch.38, p.124) requires official representation to enable supervisory bodies to conduct audits, issue corrective orders, apply sanctions and restrict or stop operation.} \hldarkgreen{A Digital Border Policy Mechanism is mandated to certify and audit AI systems, monitor training data sources, detect uncontrolled data flows, and impose temporary restrictions on systems that threaten national security.}\pdfcomment[icon=Note]{Source 2977039 (Ch.38, p.124) mandates a state Digital Border Policy Mechanism to certify and audit AI systems, monitor training data sources, detect uncontrolled data movement and impose temporary restrictions threatening national security.} \hllightgreen{These measures effectively embed logging and documentation into cross border control and digital jurisdiction \parencites[124]{2977039}.}\pdfcomment[icon=Note]{Synthesising Source 2977039 (Ch.38, p.124) on registration/Digital Border mechanisms with the logging/documentation duties in Ch.8 (p.35--36) reasonably supports that these measures embed logging/documentation into cross‑border control and digital jurisdiction.} \hllightgreen{Practice oriented sectors show how such expectations can be implemented.}\pdfcomment[icon=Note]{Multiple practice‑oriented sources (e.g., 8907572 on finance, 6805789 on education, 4639902 on intelligent audit) show how regulatory expectations can be implemented in sectoral practice, supporting this general claim.} \hldarkgreen{In financial AI, cloud based explainability architectures stream SHAP based feature attributions to a write once immutable ledger, creating real time decision logs that satisfy strict EU AI Act style traceability, human in the loop capabilities, and logging requirements for high risk systems such as credit scoring.}\pdfcomment[icon=Note]{Source 8907572 (§7.2--7.4, p.7) describes SHAP telemetry streamed to a write‑once immutable ledger to create real‑time decision logs and links this architecture to satisfying traceability, HITL and logging requirements for high‑risk financial systems; this aligns with EU AI Act obligations noted in Source 6502223.} \hldarkgreen{Institutions are advised to configure "circuit breakers" that automatically switch models to conservative human supervised modes when drift or disparate impact thresholds are crossed, and to submit models to bi annual third party audits, supported by detailed model cards and data lineage documentation \parencites[7]{8907572}.}\pdfcomment[icon=Note]{Source 8907572 (7.4, p.7) explicitly advises automated 'circuit breakers' that switch to human‑supervised modes on drift/disparate impact triggers and recommends bi‑annual third‑party audits supported by model cards and data lineage documentation.} \hldarkgreen{HITL grading in education adopts similar patterns: systems must be auditable, interpretable, and subordinate to human authority; appeals procedures are designed around audit trails that show AI outputs and human overrides; and student communication explicitly references the availability of documentation for inspections by accrediting bodies \parencites[15]{6805789}.}\pdfcomment[icon=Note]{Source 6805789 (p.15) requires HITL grading systems be auditable, interpretable, subordinate to human authority, mandate appeals procedures using audit trails, and disclose AI involvement in student communications and documentation for accreditation inspections.} \hldarkgreen{Within EU platform regulation, these requirements are echoed through structured transparency and access duties.}\pdfcomment[icon=Note]{Sources 3195574 and 6502223 describe how EU platform regulation (DSA) embeds structured transparency and access duties that echo the transparency/tracing requirements discussed elsewhere (see 3195574 p.19; 6502223 p.12).} \hldarkgreen{The DSA obliges very large online platforms to publish transparency reports and to provide data access to Digital Services Coordinators and the Commission under Article 40 for supervision of systemic risks and due diligence obligations.}\pdfcomment[icon=Note]{Source 3195574 (p.19) and Source 6502223 (p.16) document that the DSA requires VLOPs to publish transparency reports and to provide data access to Digital Services Coordinators and the Commission (Article 40) for supervision of systemic risks.} \hldarkgreen{Providers must also disclose the main parameters of recommender systems and maintain advertising repositories under Article 39, creating documentation baselines against which external audits and civil society scrutiny can be organised.}\pdfcomment[icon=Note]{Source 3195574 (p.19) cites DSA provisions requiring disclosure of main parameters of recommender systems (Article 27 in that source) and establishment of advertising repositories (Article 39), creating documentation baselines for audits and scrutiny.} \hldarkgreen{Me{\ss}mer and Degeling argue that, because implementing acts and audit standards are still emerging, voluntary standards and codes of conduct will be vital to systematise audit and assessment procedures; otherwise platforms may default to self serving minimal documentation that technically satisfies formal rules but undermines meaningful accountability \parencites[6]{3195574}[19]{3195574}.}\pdfcomment[icon=Note]{Source 3195574 (p.6--19) explicitly argues that with implementing acts and audit standards still emerging, voluntary standards and codes of conduct (Article 44--45 DSA) will be vital, warning platforms may default to minimal self‑serving documentation otherwise.} \hllightgreen{Ethical AI audit design for social recommender systems sitting under an anticipated Digital Fairness Act can therefore treat these strands as converging expectations: structured technical dossiers and algorithmic passports; tamper evident decision and event logs; explicit documentation of human oversight and appeal channels; and independent access interfaces for regulators, auditors, and vetted researchers.}\pdfcomment[icon=Note]{This is a reasonable synthesis: Sources 2977039 (Ch.8, p.35--36), 3195574 (p.19) and sectoral design recommendations (e.g., 7504644, 8907572) collectively support converging expectations of technical dossiers, algorithmic passports, tamper‑evident logs, human oversight documentation and independent access interfaces.} \hllightgreen{Regulatory analysis stressing challenges in auditing black box models, combined with calls for independent audits and external oversight committees, indicates that documentation and logging are not ancillary compliance paperwork but a precondition for any credible assessment of fairness, manipulation, or addiction risk in engagement optimised social platforms \parencites[9]{7461574}[12]{6502223}[36]{2977039}.}\pdfcomment[icon=Note]{Sources (e.g., 7461574 p.9 on auditing black‑box challenges; 2977039 Ch.8, p.35--36; and 6502223 on the EU AI Act) together justify that documentation and logging are prerequisites for credible assessments of fairness/manipulation/addiction risks on engagement‑optimised platforms.} \hldarkgreen{Practical implementations of these architectures often resemble an intelligent audit city in which continuous ETL, an analytical core, and XAI components produce live, queryable evidence for oversight; yet this continuous mode of supervision also amplifies concerns about confidentiality and false positives that regulators must manage in practice \parencites[2]{4639902}\parencites{4639902}.}\pdfcomment[icon=Note]{Source 4639902 (pp.2--3) describes an 'intelligent audit city' architecture (continuous ETL, analytical core, XAI producing live queryable evidence) and explicitly notes confidentiality and false‑positive concerns regulators must manage.} \hldarkgreen{Such continuous audit environments may also be structured to preserve human authority and to learn from human feedback, with interaction records serving both as interpretability artefacts and as training material for adaptive governance layers \parencites[10]{7504644}.}\pdfcomment[icon=Note]{Sources 7504644 (pp.10--13) and 4639902 (pp.11, 2--3) describe structuring continuous audit environments to preserve human authority and learning from human feedback, with interaction records used as interpretability artefacts and training data for adaptive governance.} \hllightgreen{In other words, decision documentation can be framed not merely as a compliance ledger but as an operational dataset that supports iterative improvement of oversight practices \parencites[10]{7504644}.}\pdfcomment[icon=Note]{Sources 8367105 (p.24) and 7504644 (pp.10--13) support the synthesis that decision documentation can function as an operational dataset to iteratively improve oversight practices rather than merely a compliance ledger.}

\subsection{Child and Youth Protection in Algorithmic Environments}

\subsubsection{Age-Sensitive Risk Assessment and Design Standards}
\label{sec:age_sensitive_risk_design} \hldarkgreen{Age related protection in AI regulation is articulated most concretely in the AI Law Model's principle of digital safety of the child.}\pdfcomment[icon=Note]{The AI Law Model explicitly defines the "principle of digital safety of the child," naming age‑related protection in that principle (Source 2977039, p.214).} \hldarkgreen{Any system that processes children's data or influences decisions about them must be "secure by default," which is defined at the architectural level: algorithmic logic, interaction mechanics, data categories, training processes, and usage purposes must all be understandable, verifiable, and subject to control and explanation across the lifecycle.}\pdfcomment[icon=Note]{Source 2977039 (p.214) states systems interacting with children must be secure by default and lists architecture, algorithmic logic, interaction mechanics, data categories, training processes and purposes as subject to understanding, verification and control.} \hldarkgreen{This framing treats age sensitivity not as a cosmetic overlay but as a structural design requirement.}\pdfcomment[icon=Note]{The Model frames age sensitivity as a requirement of system architecture and lifecycle design rather than an overlay (Source 2977039, p.214).} \hldarkgreen{Risk assessment is explicitly multi stakeholder.}\pdfcomment[icon=Note]{The text lists multiple stakeholder groups (developers, children/legal representatives, regulators, third parties) as recipients of child digital safety measures, indicating a multi‑stakeholder approach (Source 2977039, p.214--215).} \hldarkgreen{Developers and operators are expected to maintain internal technical and legal documentation, data maps, model update protocols, change logs, Child Impact Assessments, Child Safety Data Sheets, and verification and audit reports.}\pdfcomment[icon=Note]{Source 2977039 (p.214) explicitly requires developers/operators to maintain internal technical and legal documentation, data maps, update protocols, change logs, Child Impact Assessments, Child Safety Data Sheets, and verification/audit reports.} \hldarkgreen{Children and their legal representatives must receive accessible interfaces and materials that explain system logic, purpose, risks, and safety settings such as parental controls, refusal of personalisation, and time limits.}\pdfcomment[icon=Note]{The Model mandates accessible interfaces and materials for children and legal representatives explaining logic, purpose, risks and settings such as parental controls, refusal of personalization and time limits (Source 2977039, p.214).} \hldarkgreen{Regulators are entitled to a complete technical, methodological, legal and ethical dossier, including architectural schemes, threat models, risk analysis, event logs, and audit outcomes.}\pdfcomment[icon=Note]{Source 2977039 (p.214) requires provision to regulators of a complete technical, methodological, legal and ethical dossier including architectural schemes, threat models, risk analysis, event logs and audit results.} \hldarkgreen{Third parties affected by the system, such as teachers or content providers, must have access to public explanations and formal appeal procedures.}\pdfcomment[icon=Note]{The Model directs creation of mechanisms giving third parties (teachers, doctors, guardians, content providers) access to public explanations and appeal procedures (Source 2977039, p.214).} \hldarkgreen{A concrete disclosure baseline is specified.}\pdfcomment[icon=Note]{Source 2977039 (p.214) specifies mandatory disclosure items, establishing a concrete disclosure baseline for child‑facing systems.} \hldarkgreen{Systems interacting with minors must at minimum publish: data sources and categories, the legal status and contact details of the operator, intended purpose and use contexts, functional limits and reliability constraints, the degree and mode of AI participation in decisions about the child, parental control parameters, retention and deletion policies, risk levels, and results of impact assessments on children's rights.}\pdfcomment[icon=Note]{The listed minimum published information (data sources/categories, operator details, purpose, limits, AI participation, parental controls, retention, risk levels, impact assessment results) matches the explicit disclosure requirements in Source 2977039 (p.214).} \hldarkgreen{An audit log must record key events, configuration changes, model updates, and human interventions, and internal plus external audits of child safety are mandatory.}\pdfcomment[icon=Note]{The Model requires a decision/maintenance log recording key events, configuration changes, model updates and human interventions and mandates internal and external audits of child safety (Source 2977039, p.214; p.191 for logging details).} \hldarkgreen{Explainability requirements are age calibrated.}\pdfcomment[icon=Note]{Source 2977039 (p.214) requires explainability and accessibility to be adapted to age, cognitive and sociocultural characteristics, i.e., age‑calibrated explainability.} \hldarkgreen{Systems must provide meaningful, structured, age appropriate interpretations of decision logic, including descriptions of models and calculation methods, key features, uncertainty levels, and the degree of human involvement.}\pdfcomment[icon=Note]{The Model specifies providing meaningful, structured, age‑appropriate interpretations including descriptions of models, calculation methods, key features, uncertainty and human participation (Source 2977039, p.214).} \hldarkgreen{Explanations must be available at technical, functional, and ethical levels and supplied in plain language.}\pdfcomment[icon=Note]{Source 2977039 (p.214) mandates multi‑level explainability (technical, functional, ethical) provided in plain language adapted to age and context.} \hldarkgreen{Tools for children's rights are an integral part of design: refusal of profiling, correction and deletion of data, and a right to human review are singled out as capabilities that must be implemented and surfaced to users.}\pdfcomment[icon=Note]{The Model singles out tools for children's rights, refusal of profiling, correction and deletion of data, and the right to human review, as required features (Source 2977039, p.214).} \hldarkgreen{The same model links deficient design directly to enforcement.}\pdfcomment[icon=Note]{Source 2977039 links deficient design (lack of safety/accountability or manipulative interfaces) to violations and consequent enforcement measures, connecting design failures to enforcement (p.214--215).} \hldarkgreen{In public sector use, or wherever AI decisions affect a child's rights, status, access to resources, or social reputation, the absence of effective safety and accountability mechanisms, the use of overly technical or obscure messaging, or manipulative interface solutions is classified as a violation of good governance and of the best interests of the child.}\pdfcomment[icon=Note]{The Model states that in public sector or rights‑affecting uses, absence of effective safety/accountability, obscure messaging, or manipulative interfaces is a violation of good governance and the child's best interests (Source 2977039, p.214--215).} \hldarkgreen{Regulators are empowered to suspend or terminate systems, revoke permits or certificates, and must publish public justifications, with appeal rights preserved \parencites[214]{2977039}.}\pdfcomment[icon=Note]{Source 2977039 (p.215) grants authorized regulators powers to suspend/terminate systems, revoke permits or certificates and requires publication of public justifications while preserving appeal rights.} \hldarkgreen{Age sensitive standards are further operationalised through the principle of age appropriate design.}\pdfcomment[icon=Note]{The principle of age‑appropriate design is explicitly used to operationalize age‑sensitive standards in the Model (Source 2977039, p.224).} \hldarkgreen{Here, interfaces, text, interaction mechanics, and the "attention economy" must be tailored to maturity and cognitive capabilities.}\pdfcomment[icon=Note]{Source 2977039 (p.224) requires interfaces, text, interaction mechanics and the attention economy to be tailored to age, maturity and cognitive capabilities.} \hldarkgreen{Secure defaults are mandated: minimal data collection, geolocation, biometrics and profiling disabled, and explicit prohibition of manipulative and addictive mechanics such as endless scrolling, autoplay, streaks, forced quests, and penalties for exiting.}\pdfcomment[icon=Note]{The Model mandates secure defaults: minimal data, geolocation/biometrics/profiling disabled, and bans manipulative/addictive mechanics (endless scrolling, autoplay, streaks, forced quests, exit penalties) (Source 2977039, p.224; p.223).} \hldarkgreen{Push notifications must be limited to a minimum and only within reasonable need.}\pdfcomment[icon=Note]{Source 2977039 (p.223--224) limits push notifications to a minimum and permits them only for reasonable/neutral informational needs.} \hldarkgreen{Users must receive transparent, symmetrical choices, including a visible "no personalisation" switch, "lesson" and "sleep" modes, one click "why am I seeing this" explanations, simple undo actions, and straightforward data deletion.}\pdfcomment[icon=Note]{The Model requires transparent, symmetrical choices including a visible "no personalization" switch, "lesson" and "sleep" modes, one‑click "why am I seeing this" explanations, undo actions and straightforward data deletion (Source 2977039, p.224).} \hldarkgreen{Dark patterns in settings are explicitly banned, and basic access to educational and government services must not be degraded when personalisation is refused \parencites[223]{2977039}.}\pdfcomment[icon=Note]{Source 2977039 (p.223--224) explicitly bans dark patterns in settings and requires that basic access to educational and government services not be degraded when personalization is refused (offline equivalence).} \hldarkgreen{Table~\ref{tab:child_risk_design} synthesises these requirements as this report's working audit grid for age sensitive risk in social recommender systems. \begin{table}[h] \centering \caption{Working audit grid: age sensitive risk assessment and design standards} \label{tab:child_risk_design} \begin{tabular}{p{3.2cm}p{4.4cm}p{4.4cm}} \hline Dimension & Required evidence & Illustrative audit question \\ \hline Child Impact Assessment & Documented risk analysis, threat models, Child Safety Data Sheet & Is there a current Child Impact Assessment covering profiling, attention capture, and recommendation logic \parencites[214]{2977039}? \\ Secure by default & Architecture, data maps, disabled profiling/geolocation at default & Do defaults disable profiling, geolocation, biometrics, and addictive mechanics for minors \parencites[223]{2977039}? \\ Explainability & Multi level, age appropriate explanations & Can a child access "brief/detailed/technical" explanations of "why am I seeing this" in simple language \parencites[214]{2977039}[223]{2977039}? \\ Controls & No personalisation switch, parental controls, time limits & Are "lesson/sleep" modes and refusal of profiling available without service degradation \parencites[223]{2977039}? \\ Logging & Decision and event logs with minimum retention & Are tamper evident logs available to regulators and families in case of disputed decisions \parencites[214]{2977039}[191]{2977039}? \\ Governance & Internal and external child safety audits, regulator dossier & Has an external audit of child safety been completed and provided to authorities \parencites[214]{2977039}? \\ \hline \end{tabular} \end{table} National guidelines on children's digital safety and age appropriate design are envisaged to standardise such practices across models and services.}\pdfcomment[icon=Note]{The table and accompanying text synthesize the Model's requirements into an audit grid for age‑sensitive risk, and cite the Model (Source 2977039, p.214--224 referenced in table notes).} \hldarkgreen{Authorized bodies must periodically update these guidelines, define risk based model categories, and specify documentation requirements and implementation examples, which gives regulators and operators a shared reference for age sensitive risk assessment in AI mediated environments \parencites[215]{2977039}.}\pdfcomment[icon=Note]{Source 2977039 (p.215) mandates that authorized bodies periodically update national guidelines, define model risk categories and specify documentation/implementation requirements to standardize age‑sensitive risk assessment.}

\subsubsection{Restrictions on Profiling, Targeting and Monetization}
\label{sec:restrictions_profiling_targeting_monetization} \hldarkgreen{Restrictions on profiling, targeting, and monetization for minors are framed as non negotiable safeguards rather than optional product choices.}\pdfcomment[icon=Note]{The AI Law Model frames imperative prohibitions and inalienable principles for minors (e.g., "safety first" and presumption in favour of the child), treating restrictions as mandatory rather than optional (Source 2977039, p.223--224).} \hldarkgreen{Kostenko's AI Law Model defines targeted ads for children as advertising personalised on the basis of a child's data, including psychotype, neurobehavioral signs and predicted emotional states, or data sourced from third parties or cross platform tracking \parencites[20]{2977039}.}\pdfcomment[icon=Note]{The AI Law Model definition of targeted ads for children explicitly includes personalization based on psychotype, neurobehavioral signs, predicted emotional states and third‑party or cross‑platform data (Source 2977039, p.21).} \hldarkgreen{Within this definitional scope, a categorical prohibition applies: targeted advertising personalised on psychotypes, neurobehavioral profiles, vulnerabilities or predicted emotional states is banned in relation to minors.}\pdfcomment[icon=Note]{The Model lists an imperative prohibition on targeted advertising personalised based on psychotypes, neurobehavioral profiles, vulnerabilities or predicted emotional states in relation to minors (Source 2977039, p.224, 60.10(1)).} \hllightgreen{This ban directly limits monetization strategies that depend on exploiting inferred psychological traits or emotional states of children.}\pdfcomment[icon=Note]{The prohibition on such targeted advertising and related monetization is stated in the Model and therefore reasonably implies limits on monetization strategies that exploit inferred psychological traits (Source 2977039, p.223--224).} \hllightgreen{Profiling restrictions extend beyond advertising content.}\pdfcomment[icon=Note]{The Model extends restrictions beyond ads to practices like emotion tracking and predictive profiling, supporting the claim that profiling restrictions are broader than advertising content alone (Source 2977039, p.224, 60.10).} \hldarkgreen{The same imperative prohibitions forbid automated tracking of emotions, reactions, facial expressions, biometric or behavioural signs without informed consent of both the legal representative and the child, taking age and maturity into account.}\pdfcomment[icon=Note]{The Model forbids automated tracking of emotions, reactions, facial expressions, biometric or behavioural signs without informed consent of both the legal representative and the child, taking age/maturity into account (Source 2977039, p.224, 60.10(2); p.222).} \hllightgreen{Predictive profiles of "future decisions, educational or career trajectories, risks, inclinations, or 'life paths'" are also outlawed, which eliminates business models that monetise long horizon behavioural scoring for children.}\pdfcomment[icon=Note]{The Model explicitly outlaws predictive profiles of future decisions and 'life paths' for children and separately prohibits monetization, so it reasonably follows that business models monetising long‑horizon behavioural scoring are eliminated (Source 2977039, p.21; p.224, 60.10(3) and 60.10(6)).} \hldarkgreen{Operators are required to embed "child information invisibility" modes that enable service use without profiling, with minimal short term logs that do not allow reconstruction of identity \parencites[224]{2977039}.}\pdfcomment[icon=Note]{The Model defines 'child information invisibility' and requires operators to provide modes allowing use without profiling and with minimal short‑term logs that cannot restore identity (Source 2977039, p.21; p.225, operators' obligations).} \hllightgreen{For an audit toolkit, these provisions translate into hard fail criteria whenever long term behavioural profiles, emotion tracking or predictive scoring pipelines are detected in child directed social recommender components.}\pdfcomment[icon=Note]{Given the Model's prohibitions on long‑term profiles, emotion tracking and predictive scoring and its emphasis on audits and prohibitions for non‑compliance, it is a reasonable inference that these functions would constitute hard‑fail criteria in an audit toolkit for child‑facing recommender components (Source 2977039, p.223--225).} \hldarkgreen{Monetization of children's data is sharply constrained.}\pdfcomment[icon=Note]{The Model contains explicit prohibitions on monetizing children's data and related uses, showing that monetization is sharply constrained (Source 2977039, p.223--224; p.224, 60.10).} \hldarkgreen{AI systems for minors must adhere to data minimisation, processing only the minimum necessary data for a specific function \parencites[223]{2977039}.}\pdfcomment[icon=Note]{The principle of data minimization for AI systems related to minors ,  processing only the minimum necessary data ,  is explicitly required in the Model (Source 2977039, p.223--224, 'principle of data minimization and impact').} \hldarkgreen{Monetization, resale, or use of third party models for training on children's data without separate legal permission and overriding interests of the child is prohibited \parencites[224]{2977039}.}\pdfcomment[icon=Note]{The Model prohibits monetization, resale, or use of third‑party models for training on children's data without separate legal permission and absent overriding interests of the child (Source 2977039, p.224, 60.10(6)).} \hldarkgreen{In parallel, the confidentiality chapter prohibits using personal data for commercial purposes such as sale, advertising, targeted marketing or profiling aimed at influence that restricts or distorts freedom of choice, unless there is explicit, informed consent from the user \parencites[44]{2977039}.}\pdfcomment[icon=Note]{The confidentiality chapter in the Model forbids use of personal data for commercial purposes (sale, advertising, targeted marketing or profiling) or influence that distorts freedom of choice without explicit, informed consent (Source 2977039, p.43, CHAPTER 11).} \hllightgreen{When applied to minors, this effectively rules out most behaviourally targeted monetization strategies even if formal consent has been collected, because consent from children is structurally fragile and the "best interests of the child" principle requires presumption against invasive data practices \parencites[223]{2977039}.}\pdfcomment[icon=Note]{The Model applies the 'best interests of the child' and a presumption in favour of protection, and given the Model's restrictions and emphasis on consent limits for minors, it reasonably implies that many behaviourally targeted monetization strategies are ruled out even where consent is nominal (Source 2977039, p.223--224).} \hllightgreen{European platform regulation complements these AI specific provisions.}\pdfcomment[icon=Note]{EU platform regulation (DSA/related rules) is described in the provided sources as complementary to AI‑specific provisions addressing advertising, recommender alternatives and transparency (Sources 1456534, p.7; 6941611, p.19).} \hldarkgreen{Under the Digital Services Act, Recital 71 in conjunction with Article 28 establishes special rules for targeted behavioural advertising directed at children and adolescents, including a prohibition on advertising targeting based on profiles for minors.}\pdfcomment[icon=Note]{The DSA, Article 28 together with Recital 71, establishes special rules for targeted behavioural advertising directed at children and includes a prohibition on profile‑based targeting for minors, as noted in the sources (Source 6941611, p.19).} \hldarkgreen{Platforms must reasonably ascertain that a recipient is a child and, where this is the case, avoid personal data based targeting for behavioural advertising, with YouTube Kids cited as an example that abstains from personal data collection for advertising purposes.}\pdfcomment[icon=Note]{The DSA requires platforms to reasonably ascertain that a recipient is a child and avoid personal data‑based targeting for behavioural advertising; the source cites YouTube Kids as an example abstaining from personal data collection for advertising (Source 6941611, p.19).} \hldarkgreen{The same instrument obliges platforms to provide alternatives to profile based recommendation settings and to guarantee access to ad repositories containing details about targeting and presentation criteria, which enables supervisory scrutiny of whether de facto child profiling is occurring despite formal prohibitions \parencites[19]{6941611}.}\pdfcomment[icon=Note]{The DSA mandates alternatives to profile‑based recommendation settings and public ad repositories with targeting/presentation criteria to enable oversight, as described in the source (Source 6941611, p.19).} \hllightgreen{Specific product categories trigger additional restrictions.}\pdfcomment[icon=Note]{The AI Law Model and related materials identify specific product categories (e.g., e‑medicine) that trigger additional, stricter requirements, supporting the general claim (Source 2977039, p.222--224).} \hldarkgreen{AI powered "electronic medicine for teenagers" is qualified as high risk and allowed only if personalisation, advertising, and cross platform tracking are prohibited, emotional tracking and "profiles of the future" are banned, and children's data are excluded from training except under differential privacy or federated learning with separate consent.}\pdfcomment[icon=Note]{The Model qualifies AI systems for e‑medicine for teenagers as high‑risk and sets conditions: prohibition of personalization, advertising and cross‑platform tracking; ban on emotional tracking and 'profiles of the future'; and exclusion of children's data from training except under differential privacy or federated learning with separate consent (Source 2977039, p.222--223).} \hldarkgreen{Telemedicine services must reject third party SDKs and telemetry, log all accesses, and forbid diagnostic applications without medical device certification, reinforcing that health related monetization and tracking for minors are incompatible with acceptable practice \parencites[223]{2977039}.}\pdfcomment[icon=Note]{The Model requires telemedicine services to reject third‑party SDKs/telemetry, log all accesses, and forbids diagnostic applications without medical device certification, and ties non‑compliance to prohibition of operation (Source 2977039, p.222--223).} \hllightgreen{Monetization of simulacra adds another specialised restriction.}\pdfcomment[icon=Note]{The Model contains a dedicated regime for simulacra (digital twins) and their monetization, indicating specialised restrictions on economic uses of simulated identities (Source 2977039, p.21; p.140).} \hldarkgreen{Any economic benefit derived from a simulated identity or "digital twin" requires explicit, specific, revocable consent specifying goals, duration, territory, distribution channels and recall mechanisms, with detailed audit trails for each commercial use.}\pdfcomment[icon=Note]{The Model requires explicit, informed, specific and revocable consent for monetization of a simulacrum with detailed requirements on scope, duration, territory, recall mechanisms and audit trails (Source 2977039, p.140, Consent, Licensing, and Fair Remuneration).} \hldarkgreen{For minors and vulnerable persons, double consent from the individual and legal representative is required, and political, commercial and advertising monetization of simulacra is prohibited.}\pdfcomment[icon=Note]{The Model mandates double consent for minors/vulnerable persons (subject and legal representative) and prohibits political, commercial and advertising monetization of simulacra for minors (Source 2977039, p.140).} \hldarkgreen{Microtargeting and profiling on behavioural or biometric grounds are likewise banned, and an enhanced privacy mode with disabled trackers, cross domain stitching restrictions, shortened storage periods and immediate withdrawal rights is mandated \parencites[140]{2977039}.}\pdfcomment[icon=Note]{The Model bans microtargeting and profiling on behavioural/biometric grounds for minors and prescribes enhanced privacy defaults (disabled trackers/SDKs, cross‑domain restrictions, shortened retention, withdrawal rights) (Source 2977039, p.140; p.224--225).} \hllightgreen{These clauses effectively bar the use of child like avatars in monetized recommendation or advertising contexts.}\pdfcomment[icon=Note]{Given prohibitions on simulacrum monetization, microtargeting and required privacy defaults, it is a reasonable inference that monetized use of child‑like avatars in recommendation/advertising contexts would be effectively barred (Source 2977039, p.140; p.224--225).} \hldarkgreen{From an enforcement angle, breach of these restrictions is not treated as a minor compliance issue.}\pdfcomment[icon=Note]{The Model treats breaches as serious, authorising prohibitions on system operation and liability rather than minor compliance issues (Source 2977039, p.223--224).} \hldarkgreen{Non observance of "safety first" and best interests obligations in minor related AI deployment can lead to prohibition of system operation and liability for operators, suppliers or manufacturers under applicable state and international law \parencites[223]{2977039}.}\pdfcomment[icon=Note]{The Model explicitly states that non‑observance of 'safety first' and best interests obligations can lead to prohibition of system operation and liability for operators, suppliers or manufacturers (Source 2977039, p.223--224).} \hldarkgreen{Supervisory authorities are empowered to conduct inspections, audits, and impose administrative or financial sanctions, including complete bans on systems that process personal data for commercial purposes without lawful basis or that employ profiling in ways that restrict freedom of choice \parencites[44]{2977039}.}\pdfcomment[icon=Note]{The confidentiality chapter and enforcement sections grant authorized bodies powers to inspect, audit, impose administrative/financial sanctions and seek complete bans for unlawful processing or profiling that restricts freedom of choice (Source 2977039, p.43; p.224).} \hllightgreen{For social recommender audits under a Digital Fairness trajectory, these provisions imply that any child facing module combining profiling, targeting and monetization should be presumptively classified in this report's working taxonomy as "prohibited or critical risk" unless operators can demonstrate child information invisibility modes, non personalised monetization models, and strict separation of minor data from behavioural advertising and recommendation pipelines \parencites[223]{2977039}[224]{2977039}[19]{6941611}.}\pdfcomment[icon=Note]{The Model's operator obligations (child information invisibility, non‑personalised monetization, separation of minors' data) and enforcement rules support classifying child‑facing modules combining profiling, targeting and monetization as prohibited/critical risk absent demonstrable mitigations; this is a reasonable synthesis for an audit taxonomy (Source 2977039, p.224--225; p.21).}

\subsubsection{Safeguards for Youth-Focused Platforms}
\label{sec:safeguards_youth_focused_platforms} \hllightgreen{Youth focused platforms face a qualitatively different duty of care because design and data practices can shape developmental trajectories rather than only short term engagement.}\pdfcomment[icon=Note]{Source 2977039 discusses child-specific risks and the principle of digital safety (p.214) and other sources (e.g., attention-economy and dark-pattern literature such as 0783007 and 6295677) describe how design and data shape behaviour, supporting this synthesis.} \hldarkgreen{Kostenko's principle of digital safety of the child reframes this duty as a structural requirement: system architecture, algorithmic logic, interaction mechanics, data categories, training processes, and purposes of use must be "secure by default," understandable, verifiable, and controllable at each life cycle stage for any AI that processes children's data or influences decisions about them.}\pdfcomment[icon=Note]{Source 2977039 explicitly frames the 'principle of digital safety of the child' and lists the system elements that must be 'secure by default, understandable, verifiable, and subject to control' (p.214--215).} \hldarkgreen{Governance expectations span four stakeholder groups.}\pdfcomment[icon=Note]{Source 2977039 enumerates stakeholders (developers/operators, children/legal representatives, regulators, third parties like teachers/doctors) as the groups for whom digital safety must be ensured (p.214--215).} \hldarkgreen{Developers and operators must maintain internal technical and legal documentation, data maps, training and update protocols, change logs, Child Impact Assessments, Child Safety Data Sheets, and verification and audit reports.}\pdfcomment[icon=Note]{Source 2977039 explicitly lists internal documentation requirements for developers and operators including technical/legal documentation, data maps, training/update protocols, change logs, Child Impact Assessments, Child Safety Data Sheets, and audit reports (p.214--215).} \hldarkgreen{Children and their legal representatives are entitled to accessible interfaces that explain system logic, purpose, risks, and security settings such as parental controls, refusal of personalisation, and time limits, adapted to age, digital literacy, and socio cultural context.}\pdfcomment[icon=Note]{Source 2977039 requires accessible interfaces for children and legal representatives explaining logic, purpose, risks and security settings (parental controls, refusal to personalize, time limits) adapted for age and socio-cultural context (p.214--215).} \hldarkgreen{Regulators require a full technical, methodological, legal and ethical dossier, including risk models and event logs, while teachers, doctors or guardians must have access to public explanations, appeal procedures, and independent monitoring channels \parencites[214]{2977039}.}\pdfcomment[icon=Note]{Source 2977039 requires provision to regulators of a complete technical, methodological, legal and ethical dossier (including risk models and event logs) and access mechanisms for third parties (teachers, doctors, guardians) to public explanations and appeal procedures (p.214--215).} \hldarkgreen{Age segmented standards operationalise what "safe by default" means for concrete platform types.}\pdfcomment[icon=Note]{Source 2977039 develops age-appropriate design and National Guidelines on Children's Digital Safety and Age-Appropriate Design to operationalise 'secure by default' across service types (p.214--215, p.222).} \hldarkgreen{For early childhood ($0$--$6$ years), screen time must be minimal, networking functions and micropayments disabled, data collection set to zero by default, and any personalisation, profiling or advertising prohibited.}\pdfcomment[icon=Note]{Source 2977039 provides specific requirements for 0--6 years including minimal screen time, disabled networking/micropayments, zero-default data collection, and prohibition of personalization/advertising (p.217--218).} \hldarkgreen{Cameras and microphones are only activated through a single adult action with visible hardware indicators, and processing remains local without transfer to third parties.}\pdfcomment[icon=Note]{Source 2977039 mandates that cameras and microphones for early childhood are activated only by a single adult action with hardware indicators and that processing remains local without transfer to third parties (p.217--218).} \hldarkgreen{For junior schoolchildren ($7$--$12$ years), game monetisation mechanics such as loot boxes and pseudo random rewards are banned, together with manipulative retention patterns including endless scrolling, streaks and forced quests.}\pdfcomment[icon=Note]{Source 2977039 explicitly prohibits game monetisation mechanics (loot boxes, pseudo-random rewards) and manipulative retention patterns (endless scrolling, streaks, forced quests) for 7--12 year olds (p.218).} \hldarkgreen{Psychotypic targeting and personalised advertising are prohibited, data minimisation and absence of cross platform tracking are required, and push notifications are limited, with personalisation only allowed if it can be switched off easily.}\pdfcomment[icon=Note]{Source 2977039 prohibits psychotypic targeting and personalised advertising for 7--12 year olds, requires data minimisation and absence of cross-platform tracking, limits push notifications and allows personalization only if easily switchable (p.218).} \hldarkgreen{AI decisions must be explained in "children's language" \parencites[217]{2977039}.}\pdfcomment[icon=Note]{Source 2977039 requires that AI decisions be explained in age-appropriate, plain language ('children's language') with multi-level explainability (p.214--215, p.218).} \hldarkgreen{Teenage social networks receive particularly detailed normative illustrations that are directly relevant to social recommender audits.}\pdfcomment[icon=Note]{Source 2977039 provides detailed normative illustrations for teenage social networks and ties these to recommender system constraints and audits in its normative examples (p.221--222).} \hldarkgreen{For ages $13$--$15$, feeds must operate in "detox mode" by default, explicitly without endless scrolling, autoplayed video and aggressive return triggers, and must include a visible "why am I seeing this" explanation and a "no personalisation" toggle.}\pdfcomment[icon=Note]{Source 2977039 prescribes a 'detox mode' default for ages 13--15 without endless scrolling, autoplay or aggressive return triggers, and with visible 'why am I seeing this' explanations and a 'no personalization' toggle (p.222).} \hldarkgreen{Psychotypical targeting, behavioural and emotional profiling, personalised advertising and look alike audience formation are prohibited.}\pdfcomment[icon=Note]{Source 2977039 enacts an explicit ban on psychotypical targeting, behavioural/emotional profiling, personalised advertising and look-alike audience formation for teenage services (p.221--222).} \hldarkgreen{Push notifications are capped at one per day and must be strictly service or security related.}\pdfcomment[icon=Note]{Source 2977039 specifies that push notifications for teenage social networks be limited to no more than once per day and strictly service or security related (p.222).} \hldarkgreen{Profiles are closed by default, geotags and location disabled, search indexing limited, and contact from adults without joint connections blocked.}\pdfcomment[icon=Note]{Source 2977039 requires default closed profiles, geotags/location disabled, limited search indexing, and blocking contact from adults without mutual connections for youth social networks (p.222).} \hldarkgreen{All decisions are recorded in a tamper evident log, storage periods do not exceed $90$ days, external SDKs and cross platform tracking are banned, and a public "Child Safety Data Sheet" must be published.}\pdfcomment[icon=Note]{Source 2977039's normative examples require tamper-evident decision logs, storage periods not exceeding 90 days, bans on external SDKs/cross-platform tracking, and publication of a 'Child Safety Data Sheet' (p.222, p.221).} \hldarkgreen{A kill switch and quarterly Child Impact Assessments plus annual independent audits are mandated \parencites[221]{2977039}.}\pdfcomment[icon=Note]{Source 2977039 prescribes kill switches, quarterly Child Impact Assessments and annual independent audits in its guarantee and normative standards (p.221--222).} \hldarkgreen{For older teens ($16$--$17$ years), rights to self determination are expanded, yet predictive profiles of "future trajectories," psychotype targeting, emotional tracking and data monetisation remain prohibited, and explicit rights to refuse profiling and advertising, access decision logs, and appeal AI supported decisions are guaranteed.}\pdfcomment[icon=Note]{Source 2977039 describes expanded self-determination rights for 16--17 year olds while prohibiting predictive 'future trajectory' profiles, psychotype targeting, emotional tracking and data monetisation, and guaranteeing rights to refuse profiling, access decision logs and appeal (p.219--221).} \hldarkgreen{Additional vulnerabilities introduce a further safeguard layer.}\pdfcomment[icon=Note]{Source 2977039 includes an 'Additional vulnerabilities' section stating that extra safeguards apply to vulnerable groups, i.e., an additional safeguard layer (p.219--222).} \hldarkgreen{Children with disabilities, chronic illness, displacement, economic hardship or minority status are granted "minimal data by default," prohibitions on profiling, personalised advertising, geolocation tracking and emotional analytics, adapted interfaces meeting accessibility standards, offline equivalence, accelerated incident escalation, short retention periods, end to end encryption, and bans on cross platform tracking and third party data transfer \parencites[219]{2977039}.}\pdfcomment[icon=Note]{Source 2977039 explicitly lists minimal-data-by-default, prohibitions on profiling/personalised advertising/geolocation/emotional analytics, accessibility-adapted interfaces, offline equivalence, accelerated incident escalation, short retention, end-to-end encryption, and bans on cross-platform transfers for vulnerable children (p.219--222).} \hllightgreen{These conditions directly inform this report's working taxonomy of youth platform safeguards, which treats any combination of profiling, behavioural advertising, cross platform tracking, and addictive mechanics such as endless scroll as a "high or unacceptable risk" configuration for child centric services, even if engagement metrics are favourable.}\pdfcomment[icon=Note]{Source 2977039 defines unacceptable/high-risk configurations (profiling, behavioural advertising, cross-platform tracking, addictive mechanics) and the sentence reasonably synthesises those conditions into this report's taxonomy (p.217--222).} \hllightgreen{Accountability structures complement these technical and UX safeguards.}\pdfcomment[icon=Note]{Source 2977039 and related governance literature (e.g., 6502223 on regulatory frameworks) show accountability structures (registers, audits, liability) complement technical/UX safeguards, supporting this claim as a synthesis (2977039 p.83; 6502223 p.13).} \hldarkgreen{A nationwide Register of Responsible Subjects of Artificial Intelligence Systems is proposed as a prerequisite for deploying high risk AI in sectors with significant public impact, including education.}\pdfcomment[icon=Note]{Source 2977039 (p.83) mandates creation of a nationwide Register of Responsible Subjects as a prerequisite for deploying high-risk AI in sectors with significant public impact, including education.} \hldarkgreen{The register records developers, operators, audit completions, internal responsibility structures, and incidents or sanctions over the past three years.}\pdfcomment[icon=Note]{Source 2977039 details that the Register must record developers/operators, audits/certifications, internal responsibility structures, and incidents/sanctions over the past three years (p.83).} \hldarkgreen{Inclusion is a condition for permission to operate high risk systems, and failure to meet accountability principles can trigger civil, administrative, disciplinary or criminal liability, as well as temporary or indefinite bans on system operation and placement in a Register of Violators \parencites[82]{2977039}.}\pdfcomment[icon=Note]{Source 2977039 states that inclusion is a prerequisite for permission to deploy high-risk systems and that failure to meet accountability can trigger civil, administrative, disciplinary or criminal liability, bans, and placement in a Register of Violators (p.83).} \hllightgreen{For youth focused platforms, these provisions imply that visible safety features such as "ask for help" buttons, time and notification self monitoring tools, and simple reporting plus moderation service level agreements are necessary but not sufficient; they must sit on top of documented logs, impact assessments, and independent audits that regulators can verify \parencites[218]{2977039}.}\pdfcomment[icon=Note]{Source 2977039 requires visible safety tools (ask for help, pause notifications, reporting) alongside documented logs, impact assessments and independent audits that regulators can verify, so the sentence reasonably synthesises those requirements (p.221--222).} \hllightgreen{Research on AI governance stresses that such auditability is difficult when models are opaque, yet also argues that independent audits and external oversight committees are essential responses to black box risk, particularly under strict regimes like the EU AI Act and DSA \parencites[9]{7461574}[12]{6502223}.}\pdfcomment[icon=Note]{Sources (7461574 p.9 and 6502223 p.12) note auditability is difficult with opaque models and argue independent audits/external oversight are essential under regimes like the EU AI Act and DSA, supporting this synthesis (see 7461574 p.9; 6502223 p.12).} \hllightgreen{In this thesis, these constraints are treated as methodological limits: assessments of youth focused platforms can rely only on published safeguards, observable interface behaviour and available regulatory dossiers rather than on internal recommender logic, which reinforces the centrality of traceability and documentation standards articulated in the AI Law Model \parencites[214]{2977039}[82]{2977039}.}\pdfcomment[icon=Note]{Source 2977039 emphasises traceability and documentation standards (p.214--222), so treating opacity as a methodological limit and relying on published safeguards and regulatory dossiers is a reasonable methodological interpretation supported by that source (p.214--222).} \hllightgreen{Evidence from applied studies suggests that common engagement mechanics such as streaks, badges and gamified reminders can meaningfully increase repeated use but may also edge into coercive nudging when aimed at vulnerable groups; design choices that look benign for adults can be problematic for children and adolescents \parencites[5]{4426226}.}\pdfcomment[icon=Note]{Empirical and policy literature provided (e.g., 6295677, 8745065, 0783007 and UX findings in 4426226) document that gamified engagement mechanics (streaks, badges, reminders) increase repeated use and can be coercive or problematic for vulnerable groups, supporting this synthesis.} \hllightgreen{At the same time, clear anonymity options and explicit privacy controls appear to increase willingness to disclose sensitive information, which implies that privacy promises must be backed by concrete, verifiable safeguards rather than vague policy text \parencites[5]{4426226}.}\pdfcomment[icon=Note]{Source 4426226 (Maisarah et al.) reports that anonymity options and explicit privacy controls increase willingness to disclose sensitive information (p.5--6), supporting the inference that privacy promises need concrete, verifiable safeguards.}

\section{Ethical Principles for Fair and Pro-Social Recommender Systems}

\subsection{Human-Centered Design and Autonomy Preservation}

\subsubsection{Respect for Agency and Meaningful Choice}
\label{sec:respect_agency_choice} \hllightgreen{Agency in social recommender environments depends on whether people can understand, contest, and redirect algorithmic influence, not only on whether formal options exist.}\pdfcomment[icon=Note]{Synthesizes normative and empirical sources: AI Law Model emphasizes human control and counteraction to manipulative practices (Source 2977039, p.9) while literature on choice architecture highlights the need for comprehension and contestability of algorithmic influence (Source 6087241, p.3).} \hldarkgreen{Kostenko's principle of ethical responsibility demands "ethical design focused on maintaining the autonomy of the user, his psycho emotional comfort, [and] the right to an informed decision," and explicitly prohibits AI systems that manipulate behaviour, cause digital addiction, emotional exhaustion, loss of trust, or suppress autonomy \parencites[61--62]{2977039}.}\pdfcomment[icon=Note]{Directly matches the principle of ethical responsibility described in the AI Law Model, which requires ethical design maintaining autonomy, psycho-emotional comfort and informed decision-making and prohibits manipulative or addictive AI (Source 2977039, p.61).} \hldarkgreen{A complementary digital ecology chapter frames an inalienable right to a "fair, transparent, humane and environmentally friendly digital environment" that is free from intrusive content and excessive information flows creating chronic anxiety or distraction without explicit consent.}\pdfcomment[icon=Note]{The AI Law Model's digital ecology chapter explicitly proclaims an inalienable right to a fair, transparent, humane and environmentally friendly digital environment free from intrusive content and excessive information flows (Source 2977039, p.118).} \hldarkgreen{Individuals are entitled to "digital hygiene" tools to monitor their own activity, depersonalize data, and restore autonomy \parencites[119]{2977039}.}\pdfcomment[icon=Note]{The AI Law Model explicitly recognises a right to 'digital hygiene' including access to tools for monitoring activity, depersonalising data and restoring autonomy (Source 2977039, p.118).} \hldarkgreen{Agency and meaningful choice are therefore treated as enforceable rights that constrain how engagement optimisation may be implemented.}\pdfcomment[icon=Note]{The AI Law Model treats agency-related protections as enforceable obligations constraining system design and use (ethical responsibility and digital ecology provisions impose prohibitions and duties) (Source 2977039, pp.61, 118).} \hldarkgreen{Distinguishing legitimate persuasion from manipulation is central to this constraint.}\pdfcomment[icon=Note]{The conceptual review identifies distinguishing legitimate persuasion from manipulative design as a central challenge in the domain of online choice architecture and dark patterns (Source 6087241, p.3).} \hldarkgreen{Liu and Keshavarz define online choice architecture as the broader digital environment in which rankings, defaults, and recommendations shape decisions and argue that dark commercial patterns form a narrower subset of designs that "steer, deceive, pressure, or obstruct consumers in ways that impair informed and autonomous choice" \parencites[3]{6087241}.}\pdfcomment[icon=Note]{Liu and Keshavarz (as presented in the conceptual review) define online choice architecture broadly and explicitly characterise dark commercial patterns as designs that 'steer, deceive, pressure, or obstruct consumers' impairing autonomous choice (Source 6087241, p.3).} \hldarkgreen{Their marketing oriented framework evaluates interfaces along two dimensions: whether they advance firm objectives and whether they preserve the consumer's ability to make "informed, preference consistent choice." Ethical design aims to persuade while preserving that second condition; dark commercial patterns violate it by securing compliance through weakened comprehension, asymmetric effort, or hidden information \parencites[3]{6087241}[39]{6087241}.}\pdfcomment[icon=Note]{The marketing-oriented framework in the conceptual review evaluates designs along firm-objective vs. consumer informed-choice dimensions and states that ethical design aims to preserve informed, preference-consistent choice while dark patterns secure compliance by weakening comprehension or hiding information (Source 6087241, pp.3, 39).} \hllightgreen{These criteria provide a functional test for "meaningful choice" in this report's working taxonomy of agency risks.}\pdfcomment[icon=Note]{Reasonable synthesis: the audit criteria from the marketing framework (clarity, symmetry, reversibility, data respect) can functionally test 'meaningful choice' as used in the report's taxonomy (Source 6087241, p.39).} \hllightgreen{Regulatory developments in Europe translate this normative line into platform duties.}\pdfcomment[icon=Note]{Supported by analyses showing EU regulatory developments (DSA, DMA, AI Act) operationalise normative concerns about manipulative design into platform duties and obligations (Source 7896734, p.390).} \hllightgreen{The Digital Services Act requires very large online platforms to provide at least one non profiled recommendation option and prohibits interfaces that mislead users or "substantially impair" their ability to make free and informed decisions, which directly targets manipulative rankings and choice screens in social feeds \parencites[7]{1456534}[18]{6941611}.}\pdfcomment[icon=Note]{Consistent with EU regulatory discussion: the DSA and related EU instruments target manipulative interfaces and restrict practices that impair free informed decisions; literature discusses requirements for recommendation transparency and limits on manipulative interfaces (Source 7896734, p.390).} \hldarkgreen{Sanikidze's analysis of the planned Digital Fairness Act describes a trajectory toward regulating addictive product designs and unfair personalisation, closing gaps between the DSA and general consumer law and treating fairness and intelligible choice as primary policy goals rather than side effects of data protection \parencites[390]{7896734}.}\pdfcomment[icon=Note]{Sanikidze's analysis in the Dark Patterns report and related policy discussion describes a planned Digital Fairness Act aimed at regulating addictive designs and unfair personalisation, filling gaps between the DSA and consumer law (Source 7896734, p.390).} \hldarkgreen{International initiatives such as the OECD's "Dark Commercial Patterns" report and ICPEN sweeps reinforce this orientation by defining manipulative design as a systemic threat to autonomy and calling for clearer legal definitions and enforcement tools \parencites[391]{7896734}.}\pdfcomment[icon=Note]{International initiatives like the OECD 'Dark Commercial Patterns' report and ICPEN reviews are cited as defining manipulative design as systemic and calling for clearer definitions and enforcement (Sources 6087241, p.2; 7896734, p.8).} \hllightgreen{Respect for agency also depends on the ability to recognise and resist hidden steering.}\pdfcomment[icon=Note]{Supported by literature and survey evidence stressing that recognition and resistance to hidden steering are prerequisites for agency (choice-architecture literature and survey findings) (Sources 6087241, p.3; 9360054, p.31).} \hldarkgreen{Survey data on digital literacy and awareness of GDPR, the DSA and the AI Act show high self assessed competence for everyday digital tasks but much lower scores for recognising hidden manipulations (mean $2.31$ on a Likert scale) and avoiding fraud (mean $2.13$).}\pdfcomment[icon=Note]{Exact survey figures (mean 2.31 for recognising hidden manipulations and 2.13 for avoiding fraud) are reported in the study cited (Source 9360054, p.31).} \hldarkgreen{Participants express moderate trust in EU regulation but scepticism about its practical effectiveness and call for expanded educational programmes that integrate knowledge of EU policies and algorithmic influence \parencites[31]{9360054}.}\pdfcomment[icon=Note]{The same survey reports moderate awareness/trust in EU regulation alongside skepticism about practical effectiveness and support for expanded educational programmes integrating EU policy and algorithmic literacy (Source 9360054, p.31).} \hldarkgreen{These findings indicate a "second level divide" in competence, where formal options may exist but users lack the psychological and regulatory literacy to exercise them effectively.}\pdfcomment[icon=Note]{The survey analysis explicitly uses the term 'second-level divide' to describe gaps in regulatory and psychological literacy despite formal options (Source 9360054, p.31).} \hllightgreen{Agency respecting recommender governance cannot rely on disclosure alone; it must assume that many users will not detect subtle dark patterns or understand the implications of defaults.}\pdfcomment[icon=Note]{Inference supported by evidence: choice-architecture literature shows disclosures alone may not prevent dark patterns, and survey data show many users fail to detect manipulations, implying disclosure is insufficient (Sources 6087241, p.39; 9360054, p.31).} \hllightgreen{Management oriented scholarship on dark commercial patterns highlights the same tension in organisational terms.}\pdfcomment[icon=Note]{Management/marketing scholarship highlights the conversion--trust tension and organizational implications of dark patterns, aligning with the claim of an organisationally framed tension (Source 6087241, pp.33--39).} \hldarkgreen{Liu and Keshavarz propose a "conversion trust" framework in which patterns that increase immediate behavioural compliance often erode perceived fairness, autonomy, and trust, leading to higher complaints, cancellations, and reputational risk over time.}\pdfcomment[icon=Note]{The conversion--trust framework proposed in the conceptual review states that designs increasing immediate compliance can erode perceived fairness, autonomy and trust and lead to complaints, cancellations and reputational risk (Source 6087241, p.33).} \hldarkgreen{Their audit principles translate meaningful choice into actionable checks: information clarity (material terms visible before commitment), choice symmetry (declining and accepting requiring comparable effort), reversibility proportional to risk, and "data respect" in consent flows.}\pdfcomment[icon=Note]{The paper's managerial implications list audit checks including information clarity, choice symmetry, reversibility and 'data respect' as operationalisable principles (Source 6087241, pp.38--39).} \hldarkgreen{Violations of symmetry or reversibility are interpreted as signals that agency is being traded for conversion gains \parencites[39]{6087241}.}\pdfcomment[icon=Note]{The conceptual framework interprets violations of symmetry or reversibility as indicators that agency has been traded for conversion gains (conversion--trust logic in Source 6087241, pp.33, 39).} \hllightgreen{For social recommender systems, similar criteria can be applied to settings for personalisation, notification controls, and feed configuration.}\pdfcomment[icon=Note]{Reasonable extension: the audit criteria from the marketing framework can be applied to recommender settings, notifications and feed configuration, this is a plausible synthesis of the audit principles and recommender-system concerns (Source 6087241, pp.38--39; Source 3303310, p.19).} \hldarkgreen{Kostenko's AI Law Model extends these organisational responsibilities into mandatory governance structures.}\pdfcomment[icon=Note]{The AI Law Model extends organisational responsibilities into mandatory governance structures, including independent ethics commissions and mandatory procedures (Source 2977039, p.61).} \hldarkgreen{Developers and operators of AI systems must conduct independent ethical reviews of new functionality, maintain internal ethical audit mechanisms, and publish annual AI Ethics Impact Reports that assess effects on autonomy, psycho emotional health, social equality, and trust \parencites[61--62]{2977039}.}\pdfcomment[icon=Note]{The AI Law Model requires independent ethical reviews, internal ethical audit mechanisms, and annual AI Ethics Impact Reports assessing autonomy, psycho-emotional health, social equality and trust (Source 2977039, p.61).} \hldarkgreen{Digital ecology provisions prohibit algorithms that "form addiction or create self replicating digital environments that lead to the loss of a person's ability to think critically and autonomously," and they recognise a right to tools for monitoring activity and controlling digital identity \parencites[119]{2977039}.}\pdfcomment[icon=Note]{The digital ecology provisions explicitly prohibit algorithms that form addiction or self-replicating environments and recognise rights to monitoring tools and control over digital identity (Source 2977039, p.118).} \hldarkgreen{Systems that cannot be independently audited or where national actors lack control over critical components may be banned outright and placed in a register of prohibited AI systems \parencites[108]{2977039}.}\pdfcomment[icon=Note]{The AI Law Model states systems lacking independent auditability or national control can be prohibited and included in a register of prohibited AI systems (Source 2977039, pp.107--108).} \hllightgreen{These measures show that respect for agency is enforced not only through interface rules but also through structural traceability and auditability.}\pdfcomment[icon=Note]{Supported synthesis: the cited provisions combine interface rules with structural requirements for traceability and auditability in the AI Law Model, showing enforcement beyond UI-level rules (Source 2977039, pp.61, 107--108).} \hllightgreen{Research on behavioural advertising and dark patterns in the EU underscores how easily these safeguards can be undermined.}\pdfcomment[icon=Note]{Empirical and regulatory studies cited (e.g., EC study, large crawls) demonstrate prevalence of dark patterns in the EU and how they can undermine safeguards, supporting this claim (Sources 6941611, p.13; 6087241, p.2).} \hldarkgreen{Quinelato reports that a European Commission study of $45$ websites and $30$ apps found $97\%$ using deceptive practices such as false hierarchies of options and concealment of crucial information through colour and font manipulation, with social media services explicitly mentioned among high incidence sectors \parencites[13]{6941611}.}\pdfcomment[icon=Note]{The European Commission behavioural study cited in the literature is reported as finding 97\% of 45 websites and 30 apps used deceptive practices, and social media is mentioned among high-incidence sectors (Source 6941611, p.13).} \hldarkgreen{Liu and Keshavarz document large scale crawls finding thousands of dark pattern instances, including cookie banners where accepting tracking is significantly easier than refusing, which undermines the meaningfulness of consent and related choices about personalisation \parencites[33]{6087241}.}\pdfcomment[icon=Note]{The conceptual review cites large-scale crawls (≈11,000 sites) finding many dark pattern instances and documents cookie consent asymmetries that make acceptance easier than refusal (Source 6087241, p.2).} \hldarkgreen{Sanikidze's review of enforcement against Temu and platform X underlines that design choices that make cancellation harder than sign up or obscure verification processes are now treated by regulators as deceptive and potentially addictive, not as acceptable growth tactics \parencites[390]{7896734}.}\pdfcomment[icon=Note]{Sanikidze's review documents enforcement actions and regulatory focus on designs that make cancellation harder and obscure processes (mentioning Temu and platform X), treating such choices as deceptive and potentially addictive (Source 7896734, p.390).} \hllightgreen{Within this thesis, these strands motivate an audit stance in which agency and meaningful choice are operationalised through externally verifiable conditions.}\pdfcomment[icon=Note]{Reasonable methodological claim: the combination of regulatory, normative and empirical strands motivates an audit-oriented stance operationalised via externally verifiable checks (synthesis of Sources 2977039, p.61; 6087241, pp.38--39).} \hllightgreen{Social recommender systems respect agency where: non profiled feeds are genuinely accessible; consent and configuration flows meet choice symmetry and reversibility criteria; activity monitoring and "digital hygiene" tools are available; and documented ethical audits address risks of addiction and manipulation.}\pdfcomment[icon=Note]{Synthesis of legal requirements and audit principles: non-profiled feeds, choice-symmetry/reversibility in flows, digital-hygiene tools and documented ethical audits are consistent with the AI Law Model and marketing audit criteria (Sources 2977039, pp.61, 118; 6087241, pp.38--39).} \hllightgreen{Where dark commercial patterns, opaque personalisation, or self replicating engagement loops are combined with limited user capacity to recognise or contest them, systems fall into higher risk categories of this report's working taxonomy and conflict with the ethical responsibility and digital ecology principles articulated by Kostenko and with EU digital fairness objectives \parencites[61--62]{2977039}[119]{2977039}[390]{7896734}[33]{6087241}.}\pdfcomment[icon=Note]{Synthesis: combining prevalence of dark patterns, limited user literacy, and the AI Law Model/ EU fairness goals supports that such combinations raise higher risk and conflict with Kostenko/digital ecology and EU objectives (Sources 6941611, p.13; 2977039, p.118; 7896734, p.390).}

\subsubsection{Explainability and Comprehensibility for Lay Users}
\label{sec:explainability_lay_users} \hldarkgreen{Explainability in social recommender systems becomes ethically meaningful only when ordinary users can understand how algorithms affect them, not just when technical documentation exists.}\pdfcomment[icon=Note]{Source 2977039 (CH.19, p.68) explicitly requires explainability to be accessible and adapted to users rather than only technical documentation, supporting this claim.} \hllightgreen{Kostenko's principle of transparency, openness and explainability requires that system architecture, data sources, learning processes, functional limits, and potential implications are "understood, accessible, verifiable, and explanatory" throughout the lifecycle \parencites[67]{2977039}.}\pdfcomment[icon=Note]{The detailed list (system architecture, data sources, learning processes, limits and implications) matches the Principle of Transparency in 2977039 (CH.19, p.68), though the specific naming 'Kostenko' is not evident in the provided sources.} \hldarkgreen{People must be informed before any algorithmic interaction, including clarity on the degree of autonomy of the system, the level of human control, and available appeal mechanisms.}\pdfcomment[icon=Note]{2977039 (CH.19, p.68) states that no person can be exposed to AI without prior information including the limits of autonomy, human control, and appeal mechanisms, directly supporting this sentence.} \hllightgreen{Covert or non transparent AI use in socially significant contexts is explicitly prohibited, which directly challenges engagement engines that personalise feeds without intelligible disclosure.}\pdfcomment[icon=Note]{2977039 (CH.19, p.68) prohibits covert or non‑transparent AI in socially significant contexts, and studies of engagement‑optimised feeds (e.g., Gerr 5142478, p.8--9) document personalization without effective disclosure, so the sentence is a reasonable synthesis.} \hldarkgreen{Transparency in this model is stakeholder specific.}\pdfcomment[icon=Note]{2977039 (CH.19, p.68) specifies transparency obligations differentiated by stakeholder (developers, users, regulators, third parties), directly supporting the claim.} \hldarkgreen{Developers are expected to keep internal documentation, change logs and verification reports.}\pdfcomment[icon=Note]{2977039 (CH.19, p.68) explicitly requires developers to maintain internal documentation, change logs and verification reports, matching this sentence.} \hldarkgreen{Users require accessible interfaces that explain the logic of the system, adapted to digital literacy and socio cultural context.}\pdfcomment[icon=Note]{2977039 (CH.19, p.68) requires accessible user interfaces explaining system logic taking into account digital literacy and socio‑cultural context, directly supporting the sentence.} \hldarkgreen{Regulators must receive a complete technical, legal and ethical dossier in machine readable form, while third parties affected by AI decisions need mechanisms for public information, explanations and independent monitoring \parencites[68]{2977039}.}\pdfcomment[icon=Note]{2977039 (CH.19, p.68) requires regulators to receive a complete technical/methodological/legal/ethical dossier (machine‑readable) and third parties to have mechanisms for explanations, appeals and monitoring, directly supporting the claim.} \hldarkgreen{For lay users, the key implication is that explanations cannot be purely technical artefacts; they must be communicated in language and formats that match cognitive abilities and emotional states, including for adolescents and vulnerable groups.}\pdfcomment[icon=Note]{2977039 (CH.19, p.68) states explanations must be adapted to individual characteristics and provided in accessible language, and 2977039 and CH.60 (p.214--215) emphasize special attention to adolescents and vulnerable groups, supporting this sentence.} \hldarkgreen{Digital safety of the child makes these requirements even more concrete.}\pdfcomment[icon=Note]{The Principle of Digital Safety of the Child (2977039, CH.60, p.214--215) explicitly operationalizes and tightens explainability and safety requirements for systems interacting with children, directly supporting this statement.} \hldarkgreen{Systems that interact with minors must provide "meaningful, structured and age appropriate" interpretations of decision logic, explicitly covering models, calculation methods, key features, uncertainty, human participation and available alternatives.}\pdfcomment[icon=Note]{2977039 (CH.60, p.215) requires meaningful, structured and age‑appropriate explanations covering model, calculation methods, features, uncertainty, human participation and alternatives for systems interacting with minors, matching the sentence wording.} \hldarkgreen{Explanations must be free of technocratic jargon and delivered at multiple levels: technical, functional and ethical.}\pdfcomment[icon=Note]{2977039 (CH.19, p.68) mandates explanations without excessive technocratic jargon and multi‑level explainability (technical, functional, ethical), directly supporting the claim.} \hldarkgreen{The same chapter requires a decision log with key events and interventions, plus internal and external audits of child safety, giving regulators and families the evidentiary basis to contest algorithmic outcomes \parencites[214]{2977039}.}\pdfcomment[icon=Note]{2977039 (CH.19, p.68) and CH.60 (p.215) require audit logs with key events and regular internal/external audits of child safety, providing the evidentiary basis referenced in this sentence.} \hllightgreen{Age appropriate design provisions add interface level requirements such as visible "why am I seeing this" explanations and simple refusal of personalisation, which are directly relevant to engagement optimised feeds \parencites[223]{2977039}.}\pdfcomment[icon=Note]{CH.60 (2977039, p.215) mandates age‑appropriate design and options such as refusal of profiling, and studies on XAI and recommender transparency (5142478, p.8--9) discuss 'why am I seeing this' explanations, so the sentence reasonably synthesizes these sources.} \hldarkgreen{Human in the loop grading research offers a practical template for combining explainability with professional oversight.}\pdfcomment[icon=Note]{The HITL grading case study in 6805789 (p.16 and elsewhere) demonstrates combining explainability (SHAP) with educator oversight, directly supporting this claim.} \hldarkgreen{In a four layer HITL framework for essay grading, SHAP explanations highlighted textual features that contributed to model scores, and instructors reviewed AI outputs with the option to confirm, modify or override them \parencites[16]{6805789}.}\pdfcomment[icon=Note]{6805789 (Case Study, p.16) describes a four‑layer HITL grading framework, use of SHAP to highlight textual features, and instructor review with options to confirm, modify or override AI outputs, matching the sentence exactly.} \hldarkgreen{Faculty were trained on interpretation of SHAP visualisations and institutional policies for ethical use, and override actions were dual logged to create an audit trail.}\pdfcomment[icon=Note]{6805789 (Deployment Process/Technical Configuration, p.16) reports faculty training workshops on SHAP and notes dual‑logging audit mechanisms for overrides, supporting this statement.} \hldarkgreen{Case study results reported increased consistency, reduced grading time, and improved student trust, with communication making explicit that humans retained final authority \parencites[21]{6805789}.}\pdfcomment[icon=Note]{6805789 (Conclusions and Case Study, p.16 and p.21) reports reduced grading time, increased consistency, improved student trust, and that humans retained final authority, directly supporting the sentence.} \hllightgreen{This design illustrates how local explanations can be embedded in workflows in ways that both experts and affected laypersons can understand.}\pdfcomment[icon=Note]{6805789 shows how local explanations (SHAP) were integrated into workflows for instructors and the chapter argues explainability aids understanding; this is a reasonable inference that both experts and affected laypersons can understand such embedded explanations.} \hllightgreen{Explainable AI toolkits in other domains reinforce these requirements.}\pdfcomment[icon=Note]{Multiple sources (e.g., 7461574, 5442231, 5442231 p.6--15) discuss XAI toolkits in other domains and their implications, so this general reinforcing claim is a supported synthesis.} \hldarkgreen{Business analytics work summarises LIME for local surrogate explanations, SHAP for feature attribution, and counterfactual analysis for "what if" scenarios, arguing that such tools improve fairness, accountability and regulatory compliance when combined with human oversight \parencites[11]{7461574}.}\pdfcomment[icon=Note]{7461574 (p.11) and 5442231 (p.6) explicitly summarise LIME, SHAP and counterfactual analysis and argue that such tools improve fairness, accountability and compliance when combined with human oversight, directly supporting the sentence.} \hldarkgreen{Financial fairness research extends this into institutional practice: SHAP values and global interpretability tools are used to produce adverse action notices that detail which factors such as debt to income ratio or payment history drove a decision, while governance frameworks introduce fairness audits, bias detection pipelines and human supervision \parencites[6]{5442231}.}\pdfcomment[icon=Note]{5442231 (p.6) explicitly states SHAP and global interpretability tools are used to produce adverse action notices detailing factors (debt‑to‑income, payment history) and that governance frameworks introduce audits, bias detection and human supervision, directly supporting this sentence.} \hllightgreen{Sebastian projects that transparency reports will need to describe model types, training data, performance by demographic group, and governance arrangements, and that algorithmic impact assessments will test fairness, robustness and privacy before deployment and on an ongoing basis \parencites[15]{5442231}.}\pdfcomment[icon=Note]{5442231 (p.15) and related regulatory discussion describe transparency reports including model types, training data, demographic performance and algorithmic impact assessments testing fairness, robustness and privacy; attribution to 'Sebastian' is not present in the provided sources but the substance is supported.} \hllightgreen{Social recommender governance raises additional challenges.}\pdfcomment[icon=Note]{5142478 (Gerr, p.8--9) and other sources discuss specific governance difficulties with social recommender systems, so this general claim is supported as a synthesis.} \hldarkgreen{Gerr argues that platforms like Facebook, YouTube and TikTok amplify sensationalist and polarising content despite formal transparency about ranking criteria, and proposes XAI models that provide clear explanations of why specific items appear in feeds \parencites[8]{5142478}.}\pdfcomment[icon=Note]{5142478 (Gerr, p.8--9 and p.29) documents that platforms amplify sensationalist/polarising content despite transparency about ranking criteria and proposes XAI models to explain why items appear in feeds, directly supporting this sentence.} \hldarkgreen{For users, such explanations would support critical evaluation of content and help avoid manipulation.}\pdfcomment[icon=Note]{5142478 (p.29) explicitly states that user‑facing explanations would enable critical evaluation of content and help avoid manipulation, directly supporting the claim.} \hldarkgreen{For regulators, XAI outputs can support distributional fairness assessments, for example by checking ideological balance in political content.}\pdfcomment[icon=Note]{5142478 (p.29--30) states that XAI can support regulators in distributional fairness assessments (e.g., checking ideological balance), directly supporting this sentence.} \hldarkgreen{Analysis from the same study stresses that technical XAI is insufficient without regulatory frameworks that mandate its implementation and link explanations to auditing processes and user rights \parencites[30]{5142478}.}\pdfcomment[icon=Note]{5142478 (p.30) explicitly warns that technical XAI alone is insufficient and must be accompanied by regulatory frameworks mandating implementation, auditing links and user rights, supporting the sentence.} \hldarkgreen{Qualitative evidence on trust and ethical decision making among students highlights why explanation quality matters.}\pdfcomment[icon=Note]{Qualitative evidence from students in 4245543 (p.8--10, Theme and evidence tables) highlights trust and ethical decision‑making issues and the role of explanation quality, directly supporting this statement.} \hldarkgreen{In financial and ethical consumption contexts, clear AI explanations increase acceptance of algorithmic advice, whereas opaque or ambiguous logic triggers confusion and disengagement \parencites[8]{4245543}.}\pdfcomment[icon=Note]{4245543 (Table 6 and thematic findings, p.8--10) shows clear AI explanations improve acceptance while opaque logic increases confusion and disengagement, directly supporting the sentence.} \hldarkgreen{Emotional responses such as guilt or anger arise when users perceive AI outcomes as misaligned with values or as errors, and repeated exposure to unclear decisions can produce "ethical disengagement" where users stop attending to both ethical cues and algorithmic warnings \parencites[8]{4245543}[191]{4245543}.}\pdfcomment[icon=Note]{4245543 (p.8--10, Theme 1 and Theme 5) documents emotions like guilt and anger when AI outcomes are seen as misaligned and that repeated unclear decisions lead to 'ethical disengagement,' directly supporting this sentence.} \hllightgreen{For social recommenders, this suggests that poorly comprehensible explanations may not only fail to protect autonomy but can also erode willingness to use explanation tools at all.}\pdfcomment[icon=Note]{Combining findings from 4245543 on ethical disengagement and Gerr (5142478) on recommender impacts, it is a reasonable inference that incomprehensible explanations may fail to protect autonomy and reduce willingness to use explanation tools.} \hldarkgreen{From an audit design perspective, these strands converge on several concrete expectations for explainability and comprehensibility: \begin{itemize} \item Stakeholder differentiated explanation layers, with plain language summaries and appeal instructions for lay users and detailed dossiers for regulators \parencites[68]{2977039}[214]{2977039}. \item Multi level explanations that combine local reasons ("why this post now") with global descriptions of ranking objectives and constraints \parencites[68]{2977039}[8]{5142478}. \item Logging architectures that preserve decision trails in tamper evident form, enabling reconstruction and contestation of problematic recommendation episodes \parencites[214]{2977039}[241]{2977039}. \item Human in the loop mechanisms where explanations are intelligible both to professional reviewers and to affected users, as illustrated in HITL grading \parencites[16]{6805789}[21]{6805789}. \end{itemize} Research on digital fairness and AI journalism ethics adds an important normative overlay.}\pdfcomment[icon=Note]{The bulleted expectations map directly onto provisions in 2977039 (CH.19 and CH.60, p.68, p.214--215) and the HITL example in 6805789 (case study p.16), and Gerr (5142478, p.8--9) supports multi‑level explanation, so the listed items are supported by the cited sources.} \hldarkgreen{Content governance models require platforms and media organisations to provide verifiable labelling of AI generated material, fact checking tools, and independently verifiable audit logs, combined with public reports and accessible appeal procedures [7076960][4]{7076960}.}\pdfcomment[icon=Note]{2977039 (CH.28, p.93--94) and journal discussion in 7076960 (p.4--5) require verifiable labelling of AI‑generated material, fact‑checking tools, independently verifiable audit logs, public reports and appeal procedures, directly supporting this sentence.} \hldarkgreen{In this configuration, explainability is not only a technical property but part of an institutional duty to maintain truthfulness and public confidence when AI mediates information flows \parencites[94]{2977039}[7076960}.}\pdfcomment[icon=Note]{7076960 (p.4--5) and 2977039 (CH.19/CH.17, p.68 and p.61) frame explainability as an institutional duty tied to truthfulness and public confidence when AI mediates information flows, directly supporting this claim.}

\subsubsection{Meaningful Human Oversight of Automated Decisions}
\label{sec:meaningful_human_oversight} \hldarkgreen{Human oversight becomes meaningful when humans retain real authority over outcomes, supported by traceable workflows rather than symbolic "human in the loop" labels.}\pdfcomment[icon=Note]{Source 6805789 (pp.4,8) explicitly argues meaningful oversight requires real human authority and traceable workflows rather than symbolic labels, describing logging and workflow design that preserve teacher agency.} \hldarkgreen{Human in the loop grading offers a concrete template.}\pdfcomment[icon=Note]{Source 6805789 (pp.4,16) presents the Human-in-the-Loop (HITL) grading framework as a concrete template for practice.} \hldarkgreen{In a four layer framework for essay assessment, educators configure rubrics and sensitivity thresholds, the model generates preliminary scores, and a dedicated review layer allows instructors to approve, adjust, or override AI outputs before finalisation \parencites[16]{6805789}.}\pdfcomment[icon=Note]{Source 6805789 (pp.4,16) outlines a four-layer HITL framework where educators configure rubrics and sensitivity thresholds, models produce preliminary scores, and a human review layer can approve, adjust, or override outputs.} \hldarkgreen{All changes made by educators are dual logged together with the AI's initial assessment, creating an audit trail that ties each final grade to identifiable human decisions.}\pdfcomment[icon=Note]{Source 6805789 (pp.4,8,16) describes dual logging of educator changes alongside AI assessments, creating an audit trail linking final grades to identifiable human decisions.} \hldarkgreen{Faculty in the pilot emphasised that this layer preserved their role as moral and intellectual guides and allowed them to intervene whenever contextual judgment or empathy was required \parencites[8]{6805789}[21]{6805789}.}\pdfcomment[icon=Note]{Source 6805789 (p.8) reports pilot faculty emphasising that the review layer preserved their role as moral and intellectual guides and allowed intervention where contextual judgment was needed.} \hldarkgreen{Oversight that matters is therefore tightly coupled to logging and contestability.}\pdfcomment[icon=Note]{Source 6805789 (pp.4,20) ties meaningful oversight to logging and contestability, arguing transparency and audit trails are central to effective oversight.} \hldarkgreen{In the same HITL system, both AI and human actions are recorded in a dual logging architecture, enabling students to see which comments originate from the model and which from instructors.}\pdfcomment[icon=Note]{Source 6805789 (pp.4,20) states the HITL system records both AI and human actions in a dual-logging architecture and that students could see which comments came from the model versus instructors.} \hldarkgreen{This visibility strengthened trust: $61\%$ of surveyed students reported higher confidence in grading due to human validation, and comments labelled "Instructor Reviewed" were perceived as more legitimate.}\pdfcomment[icon=Note]{Source 6805789 (p.20) reports survey data where 61\% of 430 students expressed higher trust in grading due to human validation and that 'Instructor Reviewed' comments were perceived as more legitimate.} \hldarkgreen{Students could understand and challenge grades because the decision path was transparent, aligning with ethical principles of explainability and contestability \parencites[8]{6805789}[20]{6805789}.}\pdfcomment[icon=Note]{Source 6805789 (p.4) links transparent decision paths and dual-logging to students' ability to understand and challenge grades, aligning with explainability and contestability principles.} \hldarkgreen{Ethical responsibility provisions in the AI Law Model generalise this requirement beyond education.}\pdfcomment[icon=Note]{Source 2977039 (CH.17) (pp.61+) shows the AI Law Model's ethical responsibility principle generalises oversight requirements beyond education to all AI systems.} \hldarkgreen{Kostenko demands that all actors in the AI lifecycle implement internal ethical audit mechanisms, human rights impact monitoring, and preventive and compensatory measures.}\pdfcomment[icon=Note]{Source 2977039 (CH.17) explicitly requires internal ethical audit mechanisms, human rights impact monitoring, and preventive/compensatory measures across the AI lifecycle.} \hldarkgreen{Ethical design must maintain autonomy, psycho emotional comfort, and the right to an informed decision, and explicitly prohibits AI systems that manipulate behaviour, cause digital addiction, emotional exhaustion, loss of trust, or suppress autonomy.}\pdfcomment[icon=Note]{Source 2977039 (CH.17) lists ethical design obligations to maintain autonomy, psycho-emotional comfort, informed decision rights, and explicitly prohibits systems that manipulate behaviour or cause digital addiction or emotional exhaustion.} \hldarkgreen{High risk systems must satisfy this principle as a certification criterion, and regulators are empowered to suspend operation, remove unethical functionality, and apply penalties where ethical responsibility is breached.}\pdfcomment[icon=Note]{Source 2977039 (CH.17) makes ethical responsibility a mandatory certification criterion for high-risk systems and empowers authorities to suspend operations, remove unethical functionality, and apply penalties.} \hldarkgreen{Annual AI Ethics Impact Reports are mandated to document impacts on human rights, psycho emotional health, social equality and trust \parencites[61--62]{2977039}.}\pdfcomment[icon=Note]{Source 2977039 (CH.17) mandates annual AI Ethics Impact Reports documenting impacts on human rights, psycho-emotional health, social equality and trust.} \hllightgreen{On this basis, meaningful oversight means that identifiable humans and institutions can be held accountable and can intervene to halt or modify harmful configurations.}\pdfcomment[icon=Note]{Sources 6805789 (pp.4,16) and 2977039 (CH.17) together support the inference that meaningful oversight entails identifiable accountable humans/institutions able to intervene and halt or modify harmful configurations.} \hllightgreen{Business analytics research points to similar conclusions.}\pdfcomment[icon=Note]{Sources such as 7461574 and 6794105 discuss business-analytics and governance literature that reach similar conclusions about the need for human oversight, so this synthesis is supported.} \hllightgreen{Bahangulu and Owusu Berko argue that human involvement is essential where AI systems make consequential decisions, since algorithms lack contextual understanding and ethical reasoning.}\pdfcomment[icon=Note]{Source 7461574 (pp.7,9) and related governance literature argue human involvement is essential in consequential decisions because algorithms lack contextual and ethical reasoning, matching the claim attributed to Bahangulu and Owusu Berko.} \hllightgreen{They recommend human in the loop, human on the loop, and human over the loop arrangements, with humans reviewing AI generated decisions before implementation, monitoring decisions with authority to intervene, and conducting ex post audits and appeals.}\pdfcomment[icon=Note]{Sources like 6794105 and 6805789 describe human-in/on/over-the-loop arrangements and recommend review, monitoring, and audit roles, supporting the described recommendations.} \hldarkgreen{These oversight forms are presented as safeguards against discriminatory lending, unfair pricing, or biased recruitment, and as prerequisites for meeting obligations under instruments such as the EU AI Act and GDPR \parencites[7]{7461574}[9]{7461574}.}\pdfcomment[icon=Note]{Source 7461574 (pp.7,9) explicitly presents these oversight forms as safeguards against discriminatory lending, unfair pricing, biased recruitment and ties them to obligations under the EU AI Act and GDPR.} \hllightgreen{Mohammed's technical governance framework echoes this structure by requiring structured HITL workflows and explicit role descriptions for human in the loop, on the loop, and over the loop, combined with model cards, impact assessments, and bias monitoring \parencites[5]{6794105}.}\pdfcomment[icon=Note]{Source 6794105 (pp.5, B. Tier 2) describes a technical governance framework requiring structured HITL workflows, role descriptions, model cards, impact assessments and bias monitoring, aligning with the claim about Mohammed's framework.} \hllightgreen{Hybrid governance models in public sector decision support systems show how oversight can be institutionalised.}\pdfcomment[icon=Note]{Sources such as 0375818 and 8367105 discuss hybrid governance in public sector decision support and how oversight can be institutionalised, supporting this general statement.} \hldarkgreen{Case studies in healthcare triage, civil service HR, and municipal services report that systems combining internal ethics boards, external regulators, and statutory duties achieve better balances between performance and ethical assurance.}\pdfcomment[icon=Note]{Source 0375818 (Case Studies section) reports cases in healthcare triage, civil service HR, and municipal services where combining internal ethics boards, external review and statutory duties improved the balance between performance and ethical assurance.} \hldarkgreen{Quantitative indicators such as reductions in bias complaints, faster appeal resolution times, and improved citizen satisfaction are associated with transparency reports, appeal mechanisms, and third party audits.}\pdfcomment[icon=Note]{Source 0375818 (Results, pp.6) provides quantitative indicators, bias reduction rates, faster appeal resolution times, and improved citizen satisfaction, associated with transparency reports, appeals and audits.} \hldarkgreen{Comparative analysis concludes that hybrid governance, rather than pure self regulation, is most effective for ensuring fairness, transparency, accountability, and privacy in AI enabled decision support \parencites[2]{0375818}.}\pdfcomment[icon=Note]{Source 0375818 (Discussion/Results) concludes hybrid governance outperforms pure self-regulation for fairness, transparency, accountability and privacy in AI-enabled decision support.} \hllightgreen{These arrangements make human oversight continuous and multi layered, not episodic.}\pdfcomment[icon=Note]{Sources including 8367105 (CI/CD gate architecture) and 6805789 (multi-layer HITL) support the inference that these arrangements produce continuous, multi-layered oversight rather than episodic checks.} \hllightgreen{A concise comparison helps translate these insights into audit criteria for social recommender systems: \begin{table}[h] \centering \caption{Illustrative oversight patterns and audit questions} \begin{tabular}{p{3.5cm}p{4.5cm}p{4.5cm}} \hline Pattern & Key features & Audit question \\ \hline HITL grading \parencites[16]{6805789} & Human validation layer, dual logging, SHAP explanations & Are recommendation overrides and rationales logged and attributable to responsible staff? \\ Hybrid DSS governance \parencites[2]{0375818} & Internal ethics board, external regulator, appeal portal & Does the platform support user appeals and independent audits of algorithmic outcomes? \\ AI Law ethical responsibility \parencites[61--62]{2977039} & Ban on manipulative and addictive systems, ethics impact reports & Are engagement optimising features assessed for digital addiction and reported in ethics disclosures? \\ HITL roles in high risk AI \parencites[5]{6794105}[7]{7461574} & Defined human in/on/over the loop responsibilities & Are human oversight roles for recommender decisions explicitly documented and resourced? \\ \hline \end{tabular} \end{table} Research shows that explainability is central to effective oversight.}\pdfcomment[icon=Note]{The table and accompanying discussion synthesize evidence from 6805789, 8367105 and others that explainability is central to effective oversight, a reasonable academic synthesis of the cited sources.} \hldarkgreen{In HITL grading, SHAP based feature attributions made model behaviour interpretable for instructors, enabling targeted overrides where satire, abstract reasoning, or interdisciplinary arguments confused the model.}\pdfcomment[icon=Note]{Source 6805789 (pp.16,20) documents use of SHAP to make model behaviour interpretable for instructors and notes instructors overrode the model for satire, abstract reasoning or interdisciplinary arguments the model mis-scored.} \hldarkgreen{Students also valued clear distinction between AI and human feedback.}\pdfcomment[icon=Note]{Source 6805789 (pp.4,20) reports that students valued the clear distinction between AI and human feedback.} \hldarkgreen{The authors caution that AI alone cannot evaluate fairness or respond compassionately, which reinforces the case for embedding explainable AI tools into workflows where humans remain final decision makers \parencites[16]{6805789}[21]{6805789}.}\pdfcomment[icon=Note]{Source 6805789 (pp.4,8,20) explicitly cautions that AI alone cannot evaluate fairness or respond compassionately and argues for embedding explainable AI with humans as final decision-makers.} \hldarkgreen{Syari et al. reach a similar conclusion in hoax detection: low explainability undermines trust and legal usability of classification results, which motivates integration of explainable AI and critical thinking features so that outputs are both technically accurate and logically defensible to the public and law enforcement \parencites[3]{9366495}.}\pdfcomment[icon=Note]{Source 9366495 (pp.3) states that low explainability in hoax detection undermines trust and legal usability, motivating integration of XAI and critical-thinking features to make outputs defensible to the public and law enforcement.} \hllightgreen{Finally, regulatory analysis underscores the enforcement challenge.}\pdfcomment[icon=Note]{Source 7461574 (pp.9) and related regulatory analyses highlight enforcement challenges for AI, so this summarising statement is supported as a synthesis.} \hllightgreen{Bahangulu and Owusu Berko note that regulators struggle with black box models and call for improved auditing techniques and collaboration between regulators and industry.}\pdfcomment[icon=Note]{Source 7461574 (pp.9) notes regulators struggle with black-box models and calls for improved auditing techniques and collaboration between regulators and industry, aligning with the sentence.} \hldarkgreen{They argue that independent AI audits and external oversight committees are necessary to handle opaque systems, and that businesses that adopt these approaches will be better positioned to ensure compliance and ethical accountability \parencites[9]{7461574}.}\pdfcomment[icon=Note]{Source 7461574 (pp.9,15) recommends independent AI audits and external oversight committees and argues organisations adopting these will be better positioned for compliance and ethical accountability.} \hllightgreen{In this thesis, these constraints motivate a methodological stance: social recommender oversight is assessed via observable workflows, logging practices, explanation interfaces and governance structures reported in secondary sources, not via direct inspection of proprietary ranking algorithms.}\pdfcomment[icon=Note]{Sources such as 8367105 (CI/CD gate logs as audit artefacts), 6805789 (workflow/logging emphasis) support the methodological decision to assess oversight via observable workflows and logs rather than proprietary algorithm inspection; this is a reasonable methodological synthesis.} \hllightgreen{Meaningful human oversight in that context means verifiable human authority, documented decision trails, and institutional arrangements that can correct or shut down engagement strategies that conflict with ethical responsibility norms \parencites[61--62]{2977039}[16]{6805789}[2]{0375818}[9]{7461574}.}\pdfcomment[icon=Note]{Sources 6805789 (pp.4,16), 2977039 (CH.17) and 7461574 together support the interpretation that meaningful oversight requires verifiable human authority, documented decision trails, and institutional arrangements able to correct or halt unethical engagement strategies.}

\subsection{Justice, Non-Discrimination and Inclusion}

\subsubsection{Bias, Profiling and Indirect Discrimination Risks}
\label{sec:bias_profiling_indirect_discrimination} \hldarkgreen{Indirect discrimination in algorithmic environments often begins long before model deployment, at the level of data and profiling assumptions.}\pdfcomment[icon=Note]{Source 2977039 (CHAPTER 10 \& 20, pp.41,74) explicitly locates discrimination risks in data, profiling and lifecycle choices, stating bias often originates before deployment in training data and modeling assumptions.} \hldarkgreen{In adaptive education, case evidence from MathPath AI shows how under represented minority and non English speaking students were disproportionately flagged for remedial instruction because their response patterns were weakly represented in training data.}\pdfcomment[icon=Note]{Source 1308731 (pp.5,7) presents a MathPath AI case study reporting early bias where minority and non-English speaking students were over-identified for remediation due to weak representation in training data.} \hldarkgreen{This imbalance produced a higher rate of "low ability" classifications, despite similar underlying capabilities, illustrating how sparse or skewed datasets translate into structural disadvantage once embedded in adaptive recommendation logic.}\pdfcomment[icon=Note]{Source 1308731 (p.5) documents over-identification of remedial needs for underrepresented students and links this to skewed datasets translating into discriminatory classifications in adaptive systems.} \hldarkgreen{Students in rural areas with poor connectivity also experienced degraded functionality, creating a further axis of exclusion that is infrastructural rather than purely algorithmic \parencites[5]{1308731}.}\pdfcomment[icon=Note]{Source 1308731 (p.5) explicitly reports rural students with poor connectivity experiencing degraded functionality of real-time features, framing this as an infrastructural exclusion.} \hllightgreen{Comparable dynamics are anticipated at scale in algorithmic credit and educational systems.}\pdfcomment[icon=Note]{Source 5442231 (p.14) and 1308731 (pp.5,7) discuss similar bias dynamics across education and credit systems, supporting the reasonable anticipation of comparable large-scale effects.} \hldarkgreen{An AI Now Institute analysis, summarised in Chakraborty et al., reports that $65\%$ of examined educational AI studies displayed bias against marginalised groups, which signals that discriminatory effects are pervasive rather than exceptional \parencites[7]{1308731}.}\pdfcomment[icon=Note]{Source 1308731 (p.7) reproduces the AI Now Institute graphic indicating 65\% of examined educational AI studies showed bias against marginalized groups, matching the stated statistic.} \hlorange{Sebastian forecasts that credit institutions will be required to run standardized fairness tests such as demographic parity, equalized odds, calibration equity, and counterfactual fairness, and to document disparate impacts across protected classes with associated mitigation plans.}\pdfcomment[icon=Note]{While Source 5442231 (pp.14-15) describes standardized fairness tests (demographic parity, equalized odds, calibration, counterfactual fairness), none of the provided sources attribute this specific forecast to an author named 'Sebastian'.} \hldarkgreen{His sociotechnical critique warns that single metric optimisation can conceal intersectional discrimination or encode prejudices into similarity definitions, meaning that apparently "balanced" models may still harm subgroups defined by multiple attributes \parencites[14]{5442231}.}\pdfcomment[icon=Note]{Source 5442231 (p.14) explicitly warns that single-metric optimisation can mask intersectional discrimination and that similarity-based fairness can encode prejudices, matching the critique described.} \hllightgreen{For social recommender audits, this implies that narrow parity checks on aggregate engagement or exposure are insufficient; intersectional patterns in who receives addictive or harmful content must be interrogated \parencites[14]{5442231}[7]{1308731}.}\pdfcomment[icon=Note]{Sources 0822736 (p.3) on limitations of recommender audits and 5442231 (p.14) on intersectional blindspots together support the claim that narrow aggregate parity checks are insufficient and intersectional exposure patterns must be examined.} \hlorange{Normative guardrails are articulated in detail by Kostenko.}\pdfcomment[icon=Note]{No provided source mentions an author 'Kostenko' or a specific work by that name detailing normative guardrails, so the attribution lacks support in the supplied materials.} \hldarkgreen{The principle of fairness, non discrimination and inclusiveness prohibits AI systems that create or exacerbate unfair practices on any protected ground and bans models built on distorted or non representative data, algorithms with disproportionate negative impact on vulnerable groups, and systems that have not been audited for bias or lack mechanisms for its correction \parencites[73]{2977039}.}\pdfcomment[icon=Note]{Source 2977039 (CHAPTER 20, p.74) explicitly states the principle forbidding AI that creates or exacerbates unfair practices and bans systems built on distorted or non-representative data or lacking audit/correction mechanisms.} \hldarkgreen{A complementary fairness chapter defines discrimination as any restriction or advantage in treatment, access or rights based on attributes such as race, gender, age, disability or economic status, and forbids machine learning models that lead to segregation, exclusion, digital blacklists or biased risk ratings harming people socially, legally, or economically \parencites[41]{2977039}.}\pdfcomment[icon=Note]{Source 2977039 (p.41) defines discrimination as differential treatment on listed attributes and prohibits ML models that lead to segregation, exclusion, digital blacklists or biased risk ratings harming people socially, legally or economically.} \hldarkgreen{These provisions require "equity by default," diversity in social experiences, and the participation of vulnerable groups in development, together with public documentation of equity metrics, bias audits and ethical constraints integrated into model logic \parencites[73]{2977039}[41]{2977039}.}\pdfcomment[icon=Note]{Source 2977039 (p.74) mandates 'equity by default', diversity of social experiences, participation of vulnerable groups, and public documentation of equity metrics and audits, matching the sentence content.} \hllightgreen{Obligations are distributed across the lifecycle.}\pdfcomment[icon=Note]{Source 2977039 (pp.41-42) distributes responsibilities across development, deployment and post-deployment stages, supporting this lifecycle summary as a valid synthesis.} \hldarkgreen{Developers, vendors and users must perform preliminary audits of training data to identify structural or contextual biases from historical, geographic or gendered sources, apply rebalancing techniques, use fairness aware models, and create mechanisms for automated detection and signalling of discriminatory decisions.}\pdfcomment[icon=Note]{Source 2977039 (p.42) explicitly requires developers, vendors and users to audit training data for structural/contextual bias, apply rebalancing techniques, use fairness-aware models, and create mechanisms for automated detection/signalling of discriminatory decisions.} \hldarkgreen{Users must have rights to review, explanation and appeal, and all bias mitigating changes must be documented and audited for impact on both quality and fairness \parencites[41]{2977039}.}\pdfcomment[icon=Note]{Source 2977039 (p.42) mandates user rights to review, explanation and appeal and requires documentation/auditing of model changes and their impact on quality and fairness.} \hldarkgreen{An Authorized Body for Ethics of AI is tasked with maintaining a public register of discrimination complaints, analysing algorithmic impacts on demographic groups, initiating inspections and suspending suspect systems, and coordinating remedies such as compensation and restoration of rights \parencites[42]{2977039}.}\pdfcomment[icon=Note]{Source 2977039 (pp.41-42) details the Authorized Body's duties: maintaining a public complaints register, analysing demographic impacts, initiating inspections/suspensions, and coordinating remedies including compensation and rights restoration.} \hllightgreen{Within an Ethical AI Audit Toolkit, these duties can be translated into concrete checks on whether social recommender operators maintain bias registers, support appeals, and disclose model changes affecting ranking or profiling.}\pdfcomment[icon=Note]{Source 2977039 (p.42) lists obligations (bias registers, appeals, documentation of model changes) which reasonably translate into concrete checks for an Ethical AI Audit Toolkit applied to social recommender operators.} \hllightgreen{Platforms that operate as moral environments are not exempt.}\pdfcomment[icon=Note]{Source 4245543 (pp.4,10) frames e-commerce platforms as moral environments shaping sustainability salience, and 2977039 establishes regulatory obligations that would apply to platforms, supporting the claim they are not exempt.} \hldarkgreen{Qualitative work on ethical consumption among students describes e commerce platforms as configuring sustainability salience through interface design and ranking, with engagement oriented algorithms often deprioritising ethical attributes \parencites[4]{4245543}.}\pdfcomment[icon=Note]{Source 4245543 (pp.4,10) reports qualitative findings that e-commerce platforms configure sustainability salience via interface design and that engagement-oriented algorithms often deprioritise ethical attributes.} \hldarkgreen{When such systems are combined with unverifiable green claims, students experience confusion and distrust, which undermines ethical intentions and widens the intention behaviour gap.}\pdfcomment[icon=Note]{Source 4245543 (p.10) presents evidence that unverifiable green signals increase confusion and reduce consumer confidence, undermining ethical intentions among students.} \hllightgreen{These findings demonstrate how profiling and optimisation for convenience or popularity can indirectly discriminate against consumers who prioritise ethics, particularly when emotional mechanisms like guilt and confusion are triggered by greenwashing \parencites[4]{4245543}[10]{4245543}.}\pdfcomment[icon=Note]{Source 4245543 (pp.4,10) links profiling/optimisation favoring convenience and greenwashing-driven distrust to reduced ethical engagement, supporting the inference that such dynamics can indirectly disadvantage consumers prioritising ethics.} \hllightgreen{Regulatory trajectories in the EU sharpen the compliance imperative.}\pdfcomment[icon=Note]{Sources 6502223 (p.12) and 5442231 (pp.14-15) describe EU regulatory developments (AI Act, impact assessments) that increase compliance pressures, justifying the summary statement.} \hldarkgreen{The AI Act's unacceptable risk tier includes systems that exploit psychological vulnerabilities, and its high risk tier imposes conformity assessment and human oversight duties in domains such as education and essential services \parencites[12]{6502223}.}\pdfcomment[icon=Note]{Source 6502223 (p.12) explicitly states the AI Act's unacceptable risk category includes systems exploiting psychological vulnerabilities and that high-risk systems (education, essential services) face conformity assessment and human oversight duties.} \hldarkgreen{Complementary AI law proposals envision algorithmic impact assessments, fairness audit standardisation, and transparency reports detailing performance by demographic factors, governance arrangements and complaint channels \parencites[14]{5442231}[15]{5442231}.}\pdfcomment[icon=Note]{Source 5442231 (pp.14-15) describes proposals for algorithmic impact assessments, standardized fairness audits, and transparency reports with demographic performance details and governance/complaint channels.} \hldarkgreen{ISO/IEC JTC 1/SC 42 standards on risk management and ethical considerations add reference points for bias mitigation and trustworthiness, even if enforcement depends on binding instruments \parencites[3]{8589961}.}\pdfcomment[icon=Note]{Source 8589961 (p.3) describes ISO/IEC JTC 1/SC 42 standards addressing risk management and ethical considerations and notes their role in bias mitigation and trustworthiness while acknowledging voluntary enforcement limits.} \hllightgreen{An Ethical AI Audit aligned with a Digital Fairness trajectory must therefore treat bias, profiling and indirect discrimination not as reputational issues but as regulated risk categories requiring documented metrics, lifecycle controls and external oversight \parencites[12]{6502223}[73]{2977039}[41]{2977039}[14]{5442231}.}\pdfcomment[icon=Note]{Sources 2977039 (pp.41,74) and 5442231 (pp.14-15) together support the synthesis that Ethical AI Audits aligned with Digital Fairness should treat bias and profiling as regulated risk categories requiring metrics, lifecycle controls and external oversight.}

\subsubsection{Protection against Social Rating and Total Scoring}
\label{sec:protection_social_rating_total_scoring} \hllightgreen{Social recommender systems can easily drift toward quasi universal scores of user "value" if engagement, trust, or influence metrics are aggregated and reused across contexts.}\pdfcomment[icon=Note]{Recommender engagement metrics being repurposed across contexts is described in studies of personalization and engagement-driven recommendation (9072800, p.13) and regulatory critiques of engagement-maximising systems (6502223, p.15; 7896734, p.390), so this is a reasonable synthesis.} \hldarkgreen{Kostenko explicitly prohibits this move.}\pdfcomment[icon=Note]{The model law explicitly prohibits creation and use of integral social ratings/total scoring (2977039, p.129--130), which supports the claim that Kostenko (the document) forbids this move.} \hldarkgreen{Under the principle of non simulacrity of the person, any attempt to reduce a human being to an algorithmic profile, scoring index, or psychometric model is rejected: digital twins and shadow profiles are treated as auxiliary analytical constructs that must not replace a person's will, autonomy, or right to self determination \parencites[128]{2977039}.}\pdfcomment[icon=Note]{The principle of non-simulacrity and the prohibition on reducing persons to algorithmic profiles, and treatment of digital twins/shadow profiles as auxiliary constructs, is explicitly stated (2977039, p.127--128).} \hldarkgreen{This prohibition is operationalised through a categorical ban on "total scoring" and social ratings that form integral indices of "value," "trust," or "social reliability," especially where such indices influence access to labour, credit, insurance, education, housing, public and medical services, justice or e governance.}\pdfcomment[icon=Note]{The document operationalises the prohibition via a categorical ban on total scoring and social ratings affecting access to labour, credit, insurance, education, housing, public/medical services, justice and e-governance (2977039, p.129--130).} \hldarkgreen{Kostenko details how total scoring practices materialise.}\pdfcomment[icon=Note]{The source provides a detailed account of how total scoring practices are realised, listing concrete mechanisms (2977039, p.129--130), so this statement is directly supported.} \hldarkgreen{They include aggregation of heterogeneous data from multiple domains such as social networks, consumer habits, geolocation, finance, search history and biometrics in order to derive a single composite evaluation; cross platform identifier stitching via data brokers or shadow profiles; function creep where conclusions from one context are applied to others, creating persistent "risk" or "trust" labels; and use of proxy variables that reproduce discrimination on protected or socioeconomic grounds.}\pdfcomment[icon=Note]{The enumerated practices (aggregation across domains, cross‑platform stitching, function creep, use of proxies reproducing discrimination) are explicitly itemised as components of total scoring in the model law (2977039, p.129--130).} \hldarkgreen{Blacklists or whitelists that sort individuals or groups by behavioural or social traits without individual verification are singled out as emblematic of prohibited design.}\pdfcomment[icon=Note]{The prohibition text explicitly singles out blacklists/whitelists created without individual verification as emblematic of prohibited practices (2977039, p.129--130).} \hllightgreen{These mechanisms map closely onto engagement centric recommender ecosystems where cross app telemetry and behavioural segments might otherwise be repurposed for creditworthiness, job screening, or reputation filtering.}\pdfcomment[icon=Note]{The mapping between engagement‑centric recommender ecosystems and the risk of repurposing telemetry for creditworthiness or reputation filtering is supported by the personalization/engagement analysis (9072800, p.13) and regulatory critiques of engagement-driven platforms (6502223, p.15; 7896734, p.390), so this is a justified synthesis.} \hldarkgreen{Only narrowly targeted scoring remains permissible, and even then under strict cumulative conditions.}\pdfcomment[icon=Note]{The source allows only narrowly targeted scoring under simultaneous strict conditions, explicitly stating narrow permissibility (2977039, p.129--130).} \hldarkgreen{Scores must serve a specific, lawful purpose confined to a clearly delimited domain, such as credit risk assessment exclusively for a particular product.}\pdfcomment[icon=Note]{The model law gives the example that scoring may be permitted for a specific lawful purpose such as credit risk assessment for a particular product (2977039, p.129--130).} \hldarkgreen{Results cannot be transferred to other goals or sectors, protected features and their proxies must be excluded, scoring methodologies must be transparent enough to allow regulatory audit and clear personal explanation, and a non algorithmic review channel with real possibilities for human adjustment is mandatory \parencites[130]{2977039}.}\pdfcomment[icon=Note]{The required conditions (no transfer to other goals, exclusion of protected features/proxies, transparency for audit, and non-algorithmic review with human adjustment) are explicitly listed (2977039, p.129--130).} \hllightgreen{These safeguards echo fairness oriented credit scoring frameworks that argue for domain specific models, explicit anti discrimination constraints and robust human oversight in financial AI \parencites[3]{5442231}.}\pdfcomment[icon=Note]{The safeguards align with academic and industry discussions on domain‑specific credit models, anti‑discrimination constraints, and human oversight in credit/recruitment contexts (5442231, p.3; 4686679, p.6), making this a reasonable synthesis.} \hllightgreen{In the context of social recommenders, this translates into a clear boundary: engagement or trust scores generated inside a platform cannot legitimately migrate into gatekeeping functions for unrelated rights or services.}\pdfcomment[icon=Note]{Translating the legal boundary to social recommenders is a reasonable inference: the model law forbids migration of social ratings into rights/ services (2977039, p.129--130), and recommender literature documents how platform scores are generated (9072800, p.13).} \hldarkgreen{Meaningful human oversight is treated as a structural antidote to total scoring.}\pdfcomment[icon=Note]{The document treats meaningful human supervision and human review as core protections against purely automated decisions and total scoring (2977039, p.127--128; p.129--130).} \hldarkgreen{Kostenko bans internal policies or technical configurations that neutralise human supervision, limit non algorithmic review, or replace contextual justification with uniform templates.}\pdfcomment[icon=Note]{The source explicitly states that internal policies or technical configurations that neutralise human oversight or limit non‑algorithmic review violate the clause and principle of non‑simulacrity (2977039, p.129).} \hldarkgreen{In sensitive domains such as justice, health, education, labour and finance, increased standards apply: mandatory human participation, double checks for critical decisions, prohibition of predictive policing as a self sufficient intervention basis and external audits, with automatic denial of basic services explicitly disallowed \parencites[127]{2977039}.}\pdfcomment[icon=Note]{The model law prescribes increased standards in sensitive domains (justice, health, education, labour, finance), mandatory human participation, double checks, prohibition of predictive policing as sole basis, external audits, and disallows automatic denial of basic services (2977039, p.127--128).} \hldarkgreen{Complementary chapters require event logs that capture model versions, decision grounds, and appeals, together with supervising quality metrics such as cancellation rates and appeal rates, so that "rubberstamping" of automated conclusions can be detected and sanctioned \parencites[129]{2977039}.}\pdfcomment[icon=Note]{The text requires event logs capturing model versions, decision grounds, appeals and supervising quality metrics (cancellation and appeal rates) to detect and sanction rubberstamping (2977039, p.129--130; p.129 list items 2--5).} \hldarkgreen{These conditions imply that any recommender based scoring which affects entitlements must be traceable, reviewable and corrigible by humans rather than treated as authoritative.}\pdfcomment[icon=Note]{The implication that recommender‑based scoring affecting entitlements must be traceable, reviewable and corrigible by humans is directly supported by the law's requirements for traceability, auditability and human adjustment (2977039, p.129--130).} \hldarkgreen{Cross border provisions extend protection against social rating beyond national platforms.}\pdfcomment[icon=Note]{The model law contains cross‑border/extraterritorial provisions extending protections beyond national platforms (2977039, p.249--250).} \hldarkgreen{Extraterritorial digital jurisdiction is triggered where AI decisions significantly affect human rights, public security, economies or electoral and communication processes in a State, or where training and validation used data from that State without proper basis.}\pdfcomment[icon=Note]{Extraterritorial jurisdiction criteria include AI decisions significantly affecting human rights, public security, economies or electoral/communication processes and dataset provenance from the State (2977039, p.249--250).} \hldarkgreen{When triggered, providers, infrastructure operators, integrators and data brokers are jointly liable and must meet market access conditions that include transparency, compliance with human rights and data protection standards, and due diligence on counterparties through "flow down" clauses \parencites[250]{2977039}.}\pdfcomment[icon=Note]{The document states that providers, infrastructure operators, integrators and data brokers are jointly liable and must meet market access conditions including transparency, human rights/data protection compliance, and flow‑down due diligence (2977039, p.249--250; 69.3--69.4).} \hllightgreen{This framing anticipates social recommender companies that export engagement scores or influence indices into foreign credit, security or employment systems.}\pdfcomment[icon=Note]{Anticipating companies exporting engagement scores into foreign systems is an inference grounded in the law's cross‑border concern about transfer/migration of indices and the documented practice of exporting personalized signals in recommender research (2977039, p.249--250; 9072800, p.13).} \hllightgreen{An Ethical AI Audit Toolkit aligned with these norms can therefore treat certain configurations as categorically non compliant within this report's working taxonomy.}\pdfcomment[icon=Note]{Given the prohibitions and compliance criteria in the model law (2977039, p.127--130; p.249--250), it is reasonable to infer an Ethical AI Audit Toolkit could treat certain configurations as categorically non‑compliant; this is a synthesis with audit literature (8589961, p.5).} \hldarkgreen{Code Category C: "Prohibited Total Scoring" would include any practice that: aggregates multi domain data to produce a universal trust or social value index; uses cross platform ID stitching or data broker profiles to extend engagement or reputation scores into credit, hiring or access to services; maintains blacklists or whitelists derived from recommender traces without individual review; or denies opportunities solely on the basis of behavioural scores generated in social feeds \parencites[128]{2977039}[130]{2977039}.}\pdfcomment[icon=Note]{The Code Category C description closely paraphrases the prohibitions and examples listed in the model law (2977039, p.129--130) and the citation [130]\{2977039\} matches the source page.} \hllightgreen{Systems that remain within domain specific scoring, exclude proxy discrimination, provide transparent methodologies and guarantee human review could be classified as "Targeted Scoring under Oversight," yet would still require documented fairness audits and external supervisory access, similar to emerging expectations in AI credit decisioning and recruitment auditing \parencites[6]{7461574}[6]{4686679}[3]{5442231}.}\pdfcomment[icon=Note]{The conditional permissibility of domain‑specific scoring and the need for documented fairness audits and external supervisory access aligns with the model law's conditions (2977039, p.129--130) and with emerging expectations in AI credit/recruitment oversight discussed in sources (4686679, p.6; 5442231, p.3).}

\subsubsection{Cultural and Linguistic Diversity in Content Curation}
\label{sec:cultural_linguistic_diversity_curation} \hllightgreen{Cultural and linguistic identity becomes a justice question once content curation systems start imposing dominant narratives or marginalising minority voices.}\pdfcomment[icon=Note]{The AI Law Model and related sources frame cultural and linguistic identity as matters of fairness and non-discrimination when algorithmic systems impose dominant narratives or marginalise minorities (2977039, p.73-75), so this is a reasonable synthesis of those provisions.} \hldarkgreen{Kostenko formulates a dedicated legal principle that prohibits using algorithms in ways that impose cultural stereotypes, behavioural patterns, consumer habits, or ideological guidelines that contradict a person's right to preserve national, ethnic, linguistic or religious identity.}\pdfcomment[icon=Note]{The Principle described is explicitly set out in the AI Law Model (Kostenko) which prohibits using algorithms to impose cultural stereotypes or guidelines that contradict preservation of national, ethnic, linguistic or religious identity (2977039, p.182-183).} \hldarkgreen{This principle requires AI providers to avoid algorithmic configurations that normalise a single cultural viewpoint or ridicule minority accents and idioms.}\pdfcomment[icon=Note]{The Model requires providers to avoid algorithmic configurations that normalise a single cultural viewpoint or produce humiliating stereotyping (prohibitions on distorting cultural/linguistic characteristics and on humiliating imitation) (2977039, p.183; 2977039, p.74).} \hldarkgreen{For generative and curatorial systems, operators must implement cultural sensitivity mechanisms, bias testing procedures, warning labels for potentially humiliating or stereotyping outputs, and complaint channels that allow rapid correction or deletion of harmful content.}\pdfcomment[icon=Note]{The text of the AI Law Model explicitly requires generative model providers to implement cultural sensitivity mechanisms, bias testing, warnings/labels for potentially humiliating outputs, and complaint/appeal mechanisms for correction or deletion (2977039, p.183).} \hldarkgreen{Language support is treated as a structural component of non discrimination.}\pdfcomment[icon=Note]{The Model treats linguistic support and structural inclusion (linguistic, visual, cultural and functional) as components of fairness and non-discrimination (2977039, p.73-74).} \hldarkgreen{All AI and digital information systems operating in the State are obliged to support the state language across interfaces, messages, help materials and feedback channels, and to extend support to languages of national minorities "within the limits of technical feasibility." In education, media and administrative services, a multilingual environment must be created so that translation quality, reference material availability, and response speed do not degrade based on the user's language.}\pdfcomment[icon=Note]{The AI Law Model explicitly obliges digital systems operating in the State to support the state language across interfaces and, within technical feasibility, provide support for national minority languages and to ensure multilingual environments in education, media and administrative services (2977039, p.183).} \hldarkgreen{Accessible alternatives are required for users of minority languages, sign language and other communication forms \parencites[182]{2977039}.}\pdfcomment[icon=Note]{The Model mandates provision of accessible alternatives for users of national minority languages, sign language and other communication forms in education, media and administrative contexts (2977039, p.183).} \hllightgreen{These clauses directly apply to social recommender systems that rank news, entertainment or educational content: multilingual parity is not discretionary but embedded in the notion of equal access.}\pdfcomment[icon=Note]{Applying the multilingual and equal-access obligations to recommender systems is a reasonable inference: the Model's equal-access multilingual requirements (2977039, p.183) combined with literature on recommender fairness and target parity (6059812; 5142478) support this synthesis.} \hldarkgreen{Broader fairness provisions in the same AI Law Model connect cultural and linguistic diversity to discrimination risk.}\pdfcomment[icon=Note]{The AI Law Model's broader fairness provisions explicitly link cultural and linguistic diversity to discrimination risks and require inclusion and non-discrimination measures (2977039, p.73-75).} \hldarkgreen{The principle of fairness, non discrimination and inclusiveness forbids models built on distorted or non representative data, algorithms that disproportionately harm vulnerable groups, and systems that ignore cultural specifics or lack bias audits and correction mechanisms \parencites[73]{2977039}.}\pdfcomment[icon=Note]{The Model explicitly forbids systems built on distorted/non-representative data, algorithms that disproportionately harm vulnerable groups, and architectures that ignore cultural specifics or lack bias audits and correction mechanisms (2977039, p.73-75).} \hldarkgreen{Expanded guidance clarifies that prohibited systems include those that distort cultural, linguistic, religious or regional characteristics in ways that misrepresent or exclude group interests, or that are trained on data with systematic biases without proper detection, correction and multidimensional fairness testing.}\pdfcomment[icon=Note]{The Model's expanded guidance specifically lists prohibited systems that distort cultural/linguistic/religious/regional characteristics or are trained on systematically biased data without detection, correction and multidimensional fairness testing (2977039, p.74-75).} \hllightgreen{These obligations translate directly into content curation practice: training corpora and engagement logs that underrepresent minority languages or encode stereotyped portrayals cannot be treated as neutral input.}\pdfcomment[icon=Note]{This is a reasonable translation of the legal obligations into curation practice: the Model's concerns about non-representative training data (2977039, p.74-75) together with technical literature on target representations and exposure parity (6059812) support the claim.} \hldarkgreen{Fairness implementation is framed as an ongoing, audited process rather than a one time configuration.}\pdfcomment[icon=Note]{The AI Law Model requires regular, documented audits and ongoing measures to prevent discrimination, framing fairness as a continuous, auditable process (2977039, p.74-75), which matches the sentence.} \hldarkgreen{Vendors and operators must apply fairness metrics that reflect demographic, social, regional and cultural characteristics of target audiences, for example multi group accuracy, false positive and false negative rate parity and equivalent impact indices, and they must conduct regular audits of datasets, model architectures and decision processes to detect hidden bias.}\pdfcomment[icon=Note]{The Model mandates the application of fairness metrics (multi-group accuracy, FPR/FNR parity, equivalent impact indices) and regular audits of datasets, architectures and decision processes (2977039, p.74-75).} \hldarkgreen{Importantly, they are required to ensure the participation of diverse social, gender, ethnic, age and vulnerable groups in testing, public discussion and ethical review, and to follow a "fairness by design" approach where ethical and social safeguards are built into algorithmic cores from the outset \parencites[75]{2977039}.}\pdfcomment[icon=Note]{The Model requires participation of diverse social, gender, ethnic, age and vulnerable group representatives in testing, public discussion and ethical review and advocates 'fairness-by-design' safeguards (2977039, p.74-75).} \hllightgreen{In a social recommender setting, this implies that curation pipelines must be evaluated not only on engagement metrics but also on whether exposure, recommendation quality and correction pathways function equitably across cultural and linguistic communities.}\pdfcomment[icon=Note]{Combining the Model's fairness/audit requirements (2977039, p.74-75) with recommender fairness literature on exposure and relevance suggests that pipeline evaluation should include equity across cultural/linguistic communities (6059812; 5142478), a valid synthesis.} \hldarkgreen{Empirical work on AI supported mental health in Nigeria illustrates how linguistic bias in training data can translate into materially different experiences for users.}\pdfcomment[icon=Note]{The Nigerian AI mental-health empirical work documents linguistic bias and its practical consequences (misclassification, heterogeneity of inputs) and is explicitly cited in the RCT/report (5419804, p.2; p.10).} \hldarkgreen{An audit of seven pilots showed that training corpora contained less than $3\%$ of Fulfulde, Tiv and Eggon content and that none of the projects published bias audits or model cards.}\pdfcomment[icon=Note]{Table 8 and accompanying text state that across pilots training data contained <3\% Fulfulde/Tiv/Eggon content and that no project published bias audits or model cards (5419804, p.10).} \hldarkgreen{In one SMS chatbot, Fulfulde proverbs were flagged as "hostile"; in another voice stress system, female weeping in Hausa was misclassified as "anger," with clinical consequences such as under detection of PTSD and misallocation of resources.}\pdfcomment[icon=Note]{Table 8 in the report documents Fulfulde proverbs flagged as 'hostile' and a voice-stress system misclassifying female weeping in Hausa as 'anger' with clinical consequences (5419804, p.10).} \hldarkgreen{The authors' randomised trial found that algorithmic sensitivity was higher among urban, literate users and lower in rural, low literacy groups, and that federated retraining plus human mediated interpretation by Community Health Extension Workers (CHEWs) partially mitigated these disparities \parencites[10]{5419804}.}\pdfcomment[icon=Note]{The RCT summary reports higher algorithmic sensitivity among urban, literate users versus rural, low-literacy groups, and states that federated retraining plus CHEW-mediated interpretation partially mitigated disparities (5419804, p.10-11).} \hllightgreen{These findings demonstrate how culturally unbalanced corpora and monolingual optimisation can systematically disadvantage certain linguistic groups even when overall performance appears acceptable.}\pdfcomment[icon=Note]{The report's quantitative and qualitative findings (uneven performance driven by language mismatch and structural bias) reasonably support the broader conclusion that unbalanced corpora and monolingual optimisation disadvantage certain linguistic groups even when aggregate performance looks acceptable (5419804, p.10-11).} \hldarkgreen{Process evaluation in the same study highlights how culturally grounded pedagogy and "algorithmic interpreter" roles can realign AI outputs with local epistemologies.}\pdfcomment[icon=Note]{Process evaluation in the study highlights culturally grounded pedagogy and the 'algorithmic interpreter' role of CHEWs in aligning AI outputs with local epistemologies (5419804, p.11; p.2).} \hldarkgreen{CHEWs translated model outputs into culturally meaningful narratives, negotiated consent with nomadic authorities, and identified when recommendations conflicted with communal norms.}\pdfcomment[icon=Note]{The report describes CHEWs translating model outputs into culturally meaningful narratives, negotiating consent with nomadic authorities, and identifying conflicts with communal norms (5419804, p.2; p.11).} \hldarkgreen{Mixed audit circles and Pastoral Data Trust governance structures supported shared deliberation over AI outputs and reduced intergroup prejudice \parencites[11]{5419804}[2]{5419804}.}\pdfcomment[icon=Note]{The study reports that mixed audit circles and the Pastoral Data Trust governance supported shared deliberation over AI outputs and reduced intergroup prejudice (5419804, p.11; p.10 table summary).} \hllightgreen{For social recommender systems, this suggests that participatory, linguistically sensitive governance is a practical mechanism for aligning content curation with diversity and inclusion principles, especially where minority languages and oral traditions dominate.}\pdfcomment[icon=Note]{This is a reasonable inference combining the RCT's governance/pedagogy findings (5419804) with legal/regulatory emphasis on participatory design (2977039) to suggest participatory, linguistically sensitive governance as a practical mechanism for recommender alignment.} \hldarkgreen{Work on AI enabled career counselling reaches similar conclusions about the centrality of multilingual, culturally adapted systems.}\pdfcomment[icon=Note]{The career-counselling literature (9762591) explicitly argues for multilingual, culturally adapted AI systems and reaches similar conclusions about equity and localisation.} \hldarkgreen{Malik argues that sustaining equity requires AI to be multilingual and culturally adapted because large language models are primarily trained on English language, high resource corpora that embed cultural assumptions unsuitable for other settings.}\pdfcomment[icon=Note]{9762591 (Malik) states that sustaining equity requires AI to be multilingual and culturally adapted because LLMs are primarily trained on English, high-resource corpora embedding cultural assumptions unsuitable elsewhere (9762591, p.20).} \hldarkgreen{Responsible implementation therefore demands localisation, cultural adaptation, and testing across diverse groups and populations, with an explicit commitment to avoid harm to non dominant communities \parencites[9762591]{9762591}.}\pdfcomment[icon=Note]{The review concludes responsible implementation requires localisation, cultural adaptation, and testing across diverse groups with commitment to avoid harm to non-dominant communities (9762591, p.20).} \hllightgreen{This argument directly generalises to recommendation environments where English centric training and globalised engagement metrics can crowd out local languages and culturally specific content.}\pdfcomment[icon=Note]{Generalising the career-counselling concern about English-centric training to recommendation environments is a reasonable synthesis supported by Gerr's findings on engagement-optimised algorithms crowding out diverse content (5142478) and Malik's points (9762591).} \hllightgreen{Algorithmic education is identified as a complementary pillar of diversity oriented governance.}\pdfcomment[icon=Note]{The AI Law Model's principle of accessibility to algorithmic education (2977039, p.98-99) can be seen as complementary to diversity-oriented governance, a reasonable scholarly synthesis.} \hldarkgreen{Kostenko's accessibility to algorithmic education principle grants everyone a right to knowledge about AI functioning, ethical aspects and risks, and obliges states to integrate AI literacy into curricula across all levels.}\pdfcomment[icon=Note]{The AI Law Model (Kostenko) explicitly frames accessibility to algorithmic education as a right to knowledge about AI functioning and obliges the state to integrate AI literacy into curricula at all levels (2977039, p.98-99; p.136-137).} \hldarkgreen{Public educational platforms must provide AI related materials in multiple languages, adapted to age, profession and social group, and include people with disabilities through assistive technologies.}\pdfcomment[icon=Note]{The Model requires public educational platforms to provide AI materials in multiple languages, adapted to age, profession and social group, and to include persons with disabilities via assistive technologies (2977039, p.136-137).} \hldarkgreen{All initiatives must avoid political, religious or commercial bias and support critical evaluation of algorithmic outputs \parencites[98]{2977039}.}\pdfcomment[icon=Note]{The Model mandates that educational initiatives be free from political, religious or commercial bias and support critical evaluation of algorithmic outputs (2977039, p.136-137; p.99).} \hllightgreen{For social recommender audits, this implies that user facing explanations and controls around content curation should be intelligible in the user's language and framed in ways that support critical reflection on cultural and linguistic skew.}\pdfcomment[icon=Note]{Combining the Model's language/accessibility and explainability aims (2977039, p.183; p.99) with literature on explainable recommender audits (5142478; 6059812) supports the implication that user-facing explanations and controls should be intelligible in the user's language to foster critical reflection.} \hldarkgreen{Research on AI driven content curation and media diversity in Europe rounds out this picture.}\pdfcomment[icon=Note]{Gerr's Europe-focused research on AI-driven content curation and media diversity is cited in the sources and complements the prior discussion (5142478, p.24-26).} \hldarkgreen{Gerr's mixed methods study applies distributional fairness and diversity indices such as Shannon entropy and Simpson's diversity index to Facebook, YouTube and TikTok feeds to quantify representation of different perspectives \parencites[3]{5142478}.}\pdfcomment[icon=Note]{Gerr's mixed-methods study applies distributional fairness and diversity indices including Shannon entropy and Simpson's index to platform feeds as part of its quantitative analysis (5142478, p.3; p.26).} \hldarkgreen{The analysis finds that engagement optimised algorithms create echo chambers and reduce exposure to diverse viewpoints, which harms democratic dialogue.}\pdfcomment[icon=Note]{Gerr's analysis finds engagement-optimised algorithms contribute to echo chambers and reduced exposure to diverse viewpoints, harming democratic dialogue (5142478, p.29; p.8).} \hllightgreen{Although this work focuses on ideological diversity, it provides a measurement template that can be extended in this report's working coding matrix to cultural and linguistic diversity: media diversity audits can examine whether minority language content and culturally distinct narratives achieve fair exposure under ranking algorithms \parencites[5]{5142478}[8]{5142478}.}\pdfcomment[icon=Note]{Gerr's work, while focused on ideological diversity, proposes diversity measurement templates (Shannon/Simpson) and can reasonably be extended to examine cultural and linguistic diversity in media audits (5142478, p.3; p.26).} \hldarkgreen{Finally, HITL grading research demonstrates that interpretability and human agency can correct systematic misreadings of culturally embedded language.}\pdfcomment[icon=Note]{The HITL grading pilot reports that interpretability and human oversight corrected systematic misreadings of culturally embedded language and supported faculty oversight (6805789, p.19).} \hldarkgreen{In a pilot with 800 essays, $13\%$ of AI scores were fully overridden by instructors because the model misinterpreted creative writing, non linear argumentation and culturally embedded language, often from second language writers.}\pdfcomment[icon=Note]{The pilot of 800 essays documented that 13\% (104 essays) were fully overridden due to misinterpretation of creative writing, non-linear argumentation and ESL features, often from second-language writers (6805789, Table 3; p.19).} \hldarkgreen{SHAP based interpretability allowed educators to detect patterns such as penalising repetition used for rhetorical emphasis and to adjust both scores and model configurations accordingly \parencites[19]{6805789}.}\pdfcomment[icon=Note]{The report states that SHAP-based interpretability enabled educators to detect patterns (e.g., penalising repetition used for rhetorical emphasis) and adjust scores and model configurations (6805789, p.19).} \hllightgreen{This case reinforces the need for explainable, human supervised curation processes when algorithms operate across cultures and languages that were under represented or stereotyped in training data.}\pdfcomment[icon=Note]{Given the HITL findings on overrides and SHAP interpretability (6805789) and broader evidence on cultural/language misalignment (5419804), it is a reasonable scholarly inference that explainable, human-supervised curation is needed across underrepresented cultures and languages.}

\subsection{Wellbeing, Safety and Prevention of Harm}

\subsubsection{Psychological and Social Harm from Addictive Design}
\label{sec:psych_social_harm_addictive_design} \hllightgreen{Engagement optimisation in social recommender systems generates measurable psychological harms when it is implemented through addictive design patterns.}\pdfcomment[icon=Note]{Multiple sources document that engagement-optimising designs produce measurable psychological harms and that addictive design patterns drive these effects (see 3303310 p.2; 8958037 p.2), so this is a reasonable synthesis of the provided literature.} \hldarkgreen{Short video platforms built around neural network based recommendation systems, micro duration clips, infinite scrolling, auto play, and push notifications have been described as deploying an algorithmic "addiction loop" that propels users into an "auto pilot mode" of persistent scrolling.}\pdfcomment[icon=Note]{Source 8958037 explicitly describes short-video platforms using neural-network recommendations, micro-duration clips, infinite scroll, autoplay and notifications and characterises an algorithmic "addiction loop" and "auto-pilot mode" (8958037 p.2).} \hldarkgreen{These mechanics are not only associated with increased screen time: among Chinese netizens, $1.074$ billion short video users represent $95.4\%$ of all users, and college students as a core audience show diminished academic concentration, weakened social skills, heightened anxiety and depression, and eventual behavioural addiction when usage becomes excessive \parencites[2]{8958037}.}\pdfcomment[icon=Note]{Source 8958037 provides the Chinese statistics (1.074 billion short-video users = 95.4\% of netizens) and documents college students as a primary group with reported diminished concentration, weakened social skills, anxiety/depression and behavioural addiction when use is excessive (8958037 p.2).} \hldarkgreen{Hu's synthesis of "algorithm driven engagement worship" links infinite scroll and variable reward schedules on platforms such as TikTok to "dopamine driven feedback loops," with empirical validations that users of algorithmically curated feeds exhibit $22\%$ higher rates of compulsive checking than users of non algorithmic interfaces \parencites[2]{3303310}.}\pdfcomment[icon=Note]{Source 3303310 presents the synthesis linking infinite scroll and variable reward schedules to dopamine-driven feedback loops and reports empirical validations that algorithmically curated feeds yield 22\% higher rates of compulsive checking (3303310 p.2).} \hllightgreen{Neurocognitive findings sharpen this risk picture.}\pdfcomment[icon=Note]{The corpus includes neurocognitive studies and summaries (e.g., 3303310 p.3) that elaborate biological mechanisms underlying the behavioural findings, so this summarising claim is supported as an inference from those sources.} \hldarkgreen{Prolonged social media use under personalised feeds correlates with reduced gray matter density and reduced prefrontal cortex activity in regions governing impulse control, suggesting structural changes comparable to substance dependence \parencites[20]{3303310}.}\pdfcomment[icon=Note]{Source 3303310 explicitly states that prolonged social media use with personalized feeds correlates with reduced gray matter density and reduced prefrontal cortex activity in impulse-control regions, likening changes to substance dependence (3303310 p.3).} \hllightgreen{Short video architectures exacerbate this through variable reward schedules where dopamine surges are tied to the expectancy of the next clip, creating repeated "pursuit $\rightarrow$ transient satisfaction $\rightarrow$ further pursuit" cycles that elevate reward thresholds and reduce the hedonic value of longer term activities \parencites[3]{8958037}[4]{8958037}.}\pdfcomment[icon=Note]{Source 8958037 links short-video platform features and variable reward schedules to dopamine release and an auto-pilot scrolling dynamic (8958037 p.2), supporting the claim as a synthesis of neurobehavioural mechanism descriptions.} \hldarkgreen{Longitudinal evidence from the Social Media and Wellbeing Study documents that users exposed to hyper personalised feeds exhibit a $23\%$ higher incidence of depressive symptoms than controls using chronological content, particularly among adolescents and individuals predisposed to social comparison \parencites[20]{3303310}.}\pdfcomment[icon=Note]{Source 3303310 reports longitudinal data from the Social Media and Wellbeing Study showing a 23\% higher incidence of depressive symptoms for users of hyper-personalized feeds versus chronological controls, with effects concentrated among adolescents and those prone to social comparison (3303310 p.3).} \hllightgreen{These converging findings indicate that addictive design harms mental health through combined behavioural and neurobiological pathways.}\pdfcomment[icon=Note]{This is a synthesis claim drawing together behavioural and neurobiological findings reported across the sources (e.g., 3303310 p.2-3; 8958037 p.2), and is therefore supported as a reasonable academic inference.} \hllightgreen{Addictive design also has social consequences that extend beyond individual symptomatology.}\pdfcomment[icon=Note]{Sources discuss social consequences of addictive designs such as polarization and misinformation amplification (3303310 p.2; 5142478 p.29), so this general claim is a supported synthesis.} \hllightgreen{Engagement driven prioritisation in algorithmic curation amplifies misinformation and reinforces ideological polarization while users remain largely unaware of the underlying mechanisms \parencites[19]{3303310}.}\pdfcomment[icon=Note]{Empirical and analytical sources indicate engagement-driven prioritisation amplifies misinformation and reinforces polarization while users are often unaware of curation mechanisms (3303310 p.2; 5142478 p.29), supporting this combined claim.} \hldarkgreen{Analysis of content diversity in European platforms shows that AI driven personalisation exposes users to material aligned with pre existing beliefs and systematically diminishes heterogeneous viewpoints, which constrains critical thinking, civic engagement, and informed democratic participation \parencites[29]{5142478}.}\pdfcomment[icon=Note]{Source 5142478 explicitly finds AI-driven personalisation reduces exposure to diverse viewpoints and exposes users to material aligned with pre-existing beliefs, constraining critical thinking and civic participation (5142478 p.29).} \hldarkgreen{Case studies of trending topics demonstrate that algorithmically curated "Explore" and "trending" surfaces intensify political and social polarisation, especially when watch time remains the dominant optimisation metric that favours sensationalist over balanced content \parencites[29]{5142478}[3]{3303310}.}\pdfcomment[icon=Note]{The provided studies and case analyses document that algorithmic "Explore"/"trending" surfaces intensify polarization and note that watch-time optimisation favors sensationalist content over balanced content (5142478 p.29; 3303310 p.3).} \hldarkgreen{From a justice perspective, algorithmic prioritisation of Eurocentric beauty standards on Instagram has been linked to increased body dissatisfaction among adolescent women of color, illustrating how engagement optimisation can entrench cultural harms in specific groups \parencites[3]{3303310}.}\pdfcomment[icon=Note]{Source 3303310 specifically links Instagram's algorithmic prioritisation of Eurocentric beauty standards to increased body dissatisfaction among adolescent women of color (3303310 p.3).} \hllightgreen{Regulatory analysis underlines that these harms are now treated as legal concerns rather than incidental side effects.}\pdfcomment[icon=Note]{Regulatory sources in the set treat these harms as matters for legal/regulatory scrutiny (e.g., 7896734 p.7; 6502223 p.12), so the statement that they are treated as legal concerns is a supported synthesis.} \hldarkgreen{The Digital Services Act investigation into TikTok singles out infinite scrolling, auto play, and push notifications as potential systematic manipulative design that impairs users' capacity to make free and informed decisions \parencites[2]{8958037}.}\pdfcomment[icon=Note]{Source 8958037 states the EU DSA investigation into TikTok cites infinite scrolling, autoplay and push notifications as potential systematic manipulative design impairing users' capacity for free, informed decisions (8958037 p.2).} \hldarkgreen{In parallel, the EU AI Act prohibits AI systems that exploit psychological vulnerabilities and bans subliminal or harmful manipulative effects that can significantly harm behaviour, placing such systems in an unacceptable risk tier \parencites[2248]{6502223}.}\pdfcomment[icon=Note]{Source 6502223 describes the EU AI Act's prohibited 'unacceptable risk' category including AI that exploits psychological vulnerabilities, and other sources note bans on subliminal/harmful manipulative effects (6502223 p.12; 7896734 p.7).} \hldarkgreen{Emerging policy work on a Digital Fairness Act positions "addictive product designs and unfair personalisation" as priority regulatory targets, with enforcement practice already sanctioning deceptive verification flows and potentially addictive platform designs in e commerce \parencites[390]{7896734}.}\pdfcomment[icon=Note]{Source 7896734 documents the planned Digital Fairness Act targeting addictive designs and unfair personalisation and cites enforcement actions (e.g., fines for deceptive verification flows and investigations into addictive designs) (7896734 p.390--391).} \hllightgreen{These developments frame addictive interface mechanics as violations of legally recognised rights to free and informed choice, not merely as aggressive marketing.}\pdfcomment[icon=Note]{Given the regulatory framing in DSA/AI Act/DFA sources linking manipulative interfaces to impairments of decision-making (8958037 p.2; 6502223 p.12; 7896734 p.390), the inference that these mechanics are framed as rights-impinging is a supported synthesis.} \hllightgreen{Ethical AI models extend this framing by tying addictive design directly to forbidden conduct.}\pdfcomment[icon=Note]{Several normative/ethical AI documents and legislative proposals in the sources link addictive or manipulative designs to prohibited or high-risk categories and ethical obligations (2977039 p.45; 8589961 p.3), so this is a reasonable synthesis.} \hlorange{Kostenko's principle of ethical responsibility prohibits AI systems that manipulate behaviour, cause digital addiction, emotional exhaustion, loss of trust, or suppress autonomy, and requires that system operators maintain users' psycho emotional comfort and right to an informed decision \parencites[61--62]{2977039}.}\pdfcomment[icon=Note]{None of the provided sources reference a "Kostenko" principle or the specific formulation quoted here, so this attribution and the detailed prohibitions lack support in the corpus.} \hldarkgreen{Confidentiality provisions simultaneously forbid using personal data for advertising, targeted marketing or profiling aimed at influence that restricts or distorts freedom of choice without explicit, informed consent \parencites[45]{2977039}.}\pdfcomment[icon=Note]{Source 2977039 lays out confidentiality/data-use provisions requiring legal basis and forbidding use of personal data for model training or profiling without proper consent or legal grounds, aligning with the claim (2977039 p.45).} \hllightgreen{Within this normative architecture, an engagement engine that combines intensive profiling with infinite scroll and variable rewards is not a neutral recommender but a prohibited manipulative system when it yields the depressive, anxiety, compulsion and polarisation patterns described above.}\pdfcomment[icon=Note]{Combining regulatory prohibitions (e.g., DSA/AI Act) with empirical harms described earlier (3303310 p.2-3; 8958037 p.2) supports the inference that profiled infinite-scroll engines can be considered prohibited manipulative systems when causing those harms.} \hldarkgreen{Management oriented research on dark commercial patterns in Indian e commerce reveals that manipulation is often encoded in interface structures rather than content alone.}\pdfcomment[icon=Note]{Management- and HCI-oriented analyses of Indian e-commerce in the provided sources document that manipulation often is embedded in interface structures rather than content alone (6295677 p.7--8; 7896734 p.390).} \hldarkgreen{A typology of patterns including interface interference, obstruction, nagging and subscription traps demonstrates how platforms exploit System 1 biases such as default bias, scarcity effects and social pressure to secure consent and continuation while loading protective actions with friction.}\pdfcomment[icon=Note]{Source 6295677 presents a typology including interface interference, obstruction, nagging and subscription traps and links them to System 1 biases like default bias, scarcity and social pressure (6295677 p.8).} \hldarkgreen{Survey data show high prevalence of surprise charges, coerced consent and difficulty cancelling services, which leads to financial harm, frustration and erosion of trust \parencites[8]{6295677}.}\pdfcomment[icon=Note]{Source 6295677 reports survey/institutional findings (e.g., \textasciitilde{}55\% facing surprise charges, 47\% difficulty cancelling) and connects these experiences to financial harm and erosion of trust (6295677 p.8).} \hllightgreen{Although this work focuses on commerce rather than social feeds, the same cognitive mechanisms are engaged when addictive design nudges users toward extended engagement and repeated re entry via notifications, while obscuring exit options or non personalised feed alternatives \parencites[7]{6295677}[390]{7896734}.}\pdfcomment[icon=Note]{Source 7896734 discusses dark patterns in e-commerce and notes cognitive mechanisms and regulatory relevance; applying the same mechanisms to social feeds is a plausible synthesis supported by the cited material (7896734 p.390).} \hllightgreen{Human centred algorithm research proposes that these harms can be translated into operational risk indicators.}\pdfcomment[icon=Note]{Human-centred algorithm research in the provided corpus proposes operationalising harms into risk indicators (e.g., AIAS-5 in 3303310 and related regulatory/toolkit proposals), supporting this claim (3303310 p.2--3).} \hldarkgreen{Hu introduces an Algorithmic Impact Assessment Scale (AIAS 5) that quantifies mental health risk across emotional volatility, sleep disruption, social cohesion, self esteem variance and compulsive usage patterns, and notes that policymakers are beginning to require "digital nutrition labels" that disclose predicted psychological impact scores for feeds \parencites[20]{3303310}.}\pdfcomment[icon=Note]{Source 3303310 introduces the Algorithmic Impact Assessment Scale (AIAS-5) and describes its five dimensions and notes policymakers' interest in "digital nutrition labels" disclosing predicted psychological impact scores (3303310 p.2--3).} \hldarkgreen{Auditing and public documentation play a comparable role in corporate governance: independent assessments, proactive reporting and audit trails are used to surface ethical liabilities and to show regulators and users that steps have been taken to identify and fix harms \parencites[12]{8045684}.}\pdfcomment[icon=Note]{Sources discuss auditing, independent assessment and documentation as corporate governance tools to surface ethical liabilities and show remediation steps (8045684 p.12; 2977039 p.82), supporting this claim.} \hldarkgreen{In parallel, Hu's study proposes a regulatory toolkit that includes a dynamic transparency index and early warning system for addictive interactions linking neurocognitive metrics to algorithmic feature vectors such as infinite scroll and negative affect dominance \parencites[19]{3303310}.}\pdfcomment[icon=Note]{Source 3303310 explicitly proposes a regulatory toolkit including a dynamic transparency index and an early warning system linking neurocognitive metrics to algorithmic feature vectors (3303310 p.2).} \hllightgreen{Within this thesis, such instruments will inform this report's working taxonomy of harm: engagement strategies that materially elevate AIAS style dimensions and produce depressive or compulsive profiles fall into high risk or prohibited zones under a Digital Fairness trajectory, especially when combined with opaque profiling and absence of meaningful exit or non profiled options \parencites[20]{3303310}[2248]{6502223}[390]{7896734}.}\pdfcomment[icon=Note]{This is a synthesis projecting that instruments like AIAS would classify high-risk engagement strategies under Digital Fairness trajectories, which aligns with discussions in 6502223 (AI Act classification) and 7896734 (DFA planning) and is a supported inference (6502223 p.12; 7896734 p.390).} \hldarkgreen{Educational HITL systems highlight that AI mediated environments do not inevitably produce harm if designed differently.}\pdfcomment[icon=Note]{Source 6805789 reports that human-in-the-loop (HITL) educational systems can avoid harms when designed with educator oversight and ethical integration, supporting the claim (6805789 p.21).} \hldarkgreen{Human in the loop grading frameworks that maintain educator authority, dual logging, and transparent explanations report reduced processing time, increased grading consistency and improved student trust rather than addiction or disengagement.}\pdfcomment[icon=Note]{Source 6805789 documents that HITL grading frameworks maintained educator authority, used dual logging and transparent explanations, and found reduced grading time, increased consistency, and improved student trust (6805789 p.21, Conclusions).} \hldarkgreen{Students responded positively to transparent, dual sourced feedback, and faculty oversight prevented mechanical or depersonalised appraisal \parencites[21]{6805789}.}\pdfcomment[icon=Note]{Source 6805789 states students responded positively to transparent, dual-sourced feedback and that faculty oversight prevented depersonalised appraisal (6805789 p.21, Conclusions).} \hllightgreen{This contrast suggests that design choices, governance structures and human oversight are decisive in determining whether algorithmic systems degrade or support psychological wellbeing, an insight that anchors the pro social pivot strategy developed later in this report \parencites[61--62]{2977039}[21]{6805789}[19]{3303310}.}\pdfcomment[icon=Note]{Combining the HITL evidence (6805789 p.21) with Hu's work on regulatory/toolkit approaches (3303310 p.2) supports the inference that design, governance and oversight determine psychological outcomes and anchors the proposed pro-social pivot; this is a defensible synthesis.}

\subsubsection{Risk Thresholds and Unacceptable Practices}
\label{sec:risk_thresholds_unacceptable_practices} \hllightgreen{Risk boundaries for social recommender systems can be anchored in existing AI risk tiering and explicit prohibitions on manipulative and discriminatory design.}\pdfcomment[icon=Note]{Synthesises the EU risk-tiering and explicit prohibitions on manipulative/discriminatory design found in the EU AI Act summary (6502223, p.12) and the AI Law Model prohibitions (2977039, p.61).} \hldarkgreen{The EU AI Act classifies as "unacceptable risk" those AI applications that pose clear threats to fundamental rights, explicitly naming systems that exploit psychological vulnerabilities alongside social scoring and real time biometric surveillance \parencites[2248]{6502223}.}\pdfcomment[icon=Note]{Directly matches the EU AI Act description listing unacceptable risk including exploitation of psychological vulnerabilities, social scoring and real-time biometric surveillance (6502223, p.12).} \hldarkgreen{Kostenko's AI Law Model converges on this ceiling by prohibiting AI that manipulates behaviour, causes digital addiction, emotional exhaustion, loss of trust, suppresses autonomy, or promotes discrimination, and by treating ethical responsibility as a mandatory certification criterion for high risk systems \parencites[61--62]{2977039}.}\pdfcomment[icon=Note]{The AI Law Model (Kostenko) explicitly prohibits AI that manipulates behaviour, causes addiction, emotional exhaustion, loss of trust, suppresses autonomy or promotes discrimination and makes ethical responsibility a certification criterion (2977039, p.61).} \hllightgreen{These two sources jointly define a hard upper threshold for engagement optimisation: once behaviour shaping crosses into exploitation of psychological vulnerabilities or addiction, systems fall into a prohibited zone.}\pdfcomment[icon=Note]{Reasonable synthesis: the EU Act's prohibited categories (6502223, p.12) together with the AI Law Model's prohibitions (2977039, p.61) establish a clear upper bound for acceptable engagement optimisation.} \hllightgreen{Within this report's working taxonomy, a Code Category C: "Unacceptable Manipulative Practices" can therefore be defined for audit purposes.}\pdfcomment[icon=Note]{A defensible report-level taxonomy follows logically from the legal and model prohibitions cited earlier (6502223, p.12; 2977039, p.61), so defining a Code Category C for audits is a supported synthesis.} \hllightgreen{This category includes: designs that intentionally exploit psychological vulnerabilities to drive engagement; interface patterns that induce digital addiction through infinite scroll, variable reward schedules and push based re entry without meaningful exit options; and architectures that suppress user autonomy by blocking non profiled recommendations that the DSA treats as mandatory options for very large platforms \parencites[2248]{6502223}[61--62]{2977039}.}\pdfcomment[icon=Note]{Components listed draw on documented dark-pattern/addiction mechanisms (infinite scroll, variable reward schedules) and regulatory DSA requirements for non-profiled recommendations (3303310, p.19; 6295677, p.2; 6502223, p.12).} \hlorange{While the Digital Fairness Act remains a trajectory rather than codified law, its stated ambition to regulate "addictive product designs and unfair personalisation" implies that these practices will not merely be discouraged but structurally constrained \parencites[390]{7896734}.}\pdfcomment[icon=Note]{No provided source describes a 'Digital Fairness Act' with the quoted ambition; this specific legislative title and phrasing are not present in the supplied materials.} \hllightgreen{A second ceiling arises around discrimination and social rating.}\pdfcomment[icon=Note]{Supported by the EU Act's and AI Law Model's explicit treatment of social scoring and discrimination as high/unacceptable risks (6502223, p.12; 2977039, p.61).} \hldarkgreen{Sebastian anticipates that regulators will "legalise" fairness audits in credit systems, mandating pre deployment verification, disparate impact tests across protected classes, ongoing monitoring and mitigation, and he warns that single metric fixes can mask intersectional harms \parencites[14]{5442231}.}\pdfcomment[icon=Note]{Matches the forecasted legalization of formalised fairness audits with pre-deployment checks, disparate impact tests, monitoring, and warnings about single-metric fixes as described in the credit systems literature (5442231, p.14).} \hldarkgreen{Kostenko's fairness principle directly prohibits models built on distorted or non representative data, algorithms with disproportionate negative impact on vulnerable groups, and systems lacking bias audits and correction mechanisms, while banning social scoring style "integral assessments" where they influence access to services or rights \parencites[41]{2977039}.}\pdfcomment[icon=Note]{The AI Law Model explicitly bans systems that promote discrimination and requires ethical procedures including bias monitoring and audits; it also treats ethical responsibility as certification criterion (2977039, p.61).} \hllightgreen{For social recommender audits, this justifies a Code Category C: "Unacceptable Scoring and Discrimination" that covers any cross domain "trust" or "reliability" index derived from engagement traces and reused for access decisions in education, credit, or employment, as well as unmitigated disparate exposure to harmful or addictive content across demographic groups \parencites[41]{2977039}[14]{5442231}.}\pdfcomment[icon=Note]{Reasonable extension: invoking a Code Category covering cross-domain trust/reliability indices and disparate exposure is supported by fairness audit recommendations in credit and algorithm governance literature (5442231, p.14) and the harms described for recommender reuse.} \hllightgreen{Between legally prohibited and apparently compliant practice, AIN and related governance frameworks suggest how to set operational thresholds.}\pdfcomment[icon=Note]{The Algorithmic Impact Navigator (AIN) and governance frameworks are presented as mechanisms to operationalise thresholds between prohibited and compliant practice (2742041, pp.10--11; 8367105, p.24).} \hldarkgreen{The Algorithmic Impact Navigator demonstrates in lending that a 30\% reduction in processing time is compatible with fairness when denial gaps across protected groups are statistically insignificant, and that routine fairness audits, continuous monitoring and documented bias mitigation prevent small discrepancies from scaling into systemic harm \parencites[2742041]{2742041}.}\pdfcomment[icon=Note]{Directly supported by the AIN case evidence showing a 30\% reduction in lending processing time with no significant denial gap and the role of routine fairness audits and monitoring (2742041, p.10--11).} \hldarkgreen{Sebastian similarly describes hybrid human AI governance where human review of borderline cases, compulsory human authorisation for sensitive decisions, and expert panels reviewing systemic patterns are built into workflows, not added as ex post checks \parencites[14]{5442231}.}\pdfcomment[icon=Note]{The hybrid human--AI governance features (human review of borderline cases, compulsory human authorization, expert panels) are explicitly described in the credit/governance literature (5442231, p.14).} \hllightgreen{Translating these ideas, this report defines a Code Category B: "High Risk but Potentially Mitigable Practices" for recommender configurations that currently rely on engagement maximisation or extensive profiling yet can demonstrate: documented fairness audits; human oversight layers with authority to modify ranking policies; and continuous monitoring tied to pre agreed harm triggers.}\pdfcomment[icon=Note]{A reasonable translation: defining a Code Category B that requires documented fairness audits, empowered human oversight, and continuous monitoring follows from AIN and governance recommendations (2742041, p.10; 5442231, p.14).} \hllightgreen{Unacceptable practice is also a function of governance failure, not only of algorithmic design.}\pdfcomment[icon=Note]{Supported by sources emphasising that governance failures (absence of gates, controls), not just algorithm design, produce unacceptable outcomes (8367105, p.24; 2742041, p.10).} \hldarkgreen{Malachi's trial in the Nigerian Middle Belt illustrates how AI enabled mental health support can improve PTSD remission and social cohesion when embedded in Pastoral Data Trust governance, mixed audit circles, and human intermediaries trained in algorithmic interpretation.}\pdfcomment[icon=Note]{The Malachi Nigerian trial reports improved PTSD and depression outcomes, strengthened social cohesion, use of Pastoral Data Trust governance, mixed audit circles, and CHEWs trained in algorithmic interpretation (5419804, pp.11,13).} \hldarkgreen{Yet the same study documents serious adverse events, including a false positive suicide alert that triggered a violent vigilante patrol and a privacy breach caused by a community health worker bypassing controls, which required model recalibration, device decommissioning and personnel removal \parencites[5419804]{5419804}.}\pdfcomment[icon=Note]{The study documents the two serious adverse events (false-positive suicide alert triggering vigilante patrol; CHEW-caused privacy breach) and the corrective actions including threshold recalibration, device decommissioning and personnel removal (5419804, pp.8,11).} \hllightgreen{These incidents demonstrate that conservative thresholds, strict protocol adherence and enforceable sanctions for misuse are non negotiable guardrails.}\pdfcomment[icon=Note]{The study's ethical discussion explicitly recommends conservative thresholds, strict protocol adherence and governance safeguards, supporting the claim that these are indispensable guardrails (5419804, p.11).} \hllightgreen{In this thesis, audit findings that combine high engagement optimisation, weak governance, absence of independent ethical oversight and lack of community accountability mechanisms will be classified in Code Category C regardless of nominal compliance with technical metrics \parencites[61--62]{2977039}[5419804]{5419804}.}\pdfcomment[icon=Note]{Authorial policy statement is consistent with the Malachi trial's emphasis on governance and the report's cited standards; using such audit findings to classify systems as Category C is a defensible normative rule referencing the study (5419804, p.11).} \hldarkgreen{ISO/IEC JTC 1/SC 42 standards and risk based regulations extend these expectations into global practice.}\pdfcomment[icon=Note]{ISO/IEC JTC 1/SC 42 and standards-based committees are cited as extending AI risk and trustworthiness expectations globally (8589961, p.3).} \hldarkgreen{Standards such as ISO/IEC 23894 and 24368 operationalise risk management and ethical considerations, while emerging regulations across multiple regions emphasise trustworthiness, transparency and governance for high impact AI \parencites[8589961]{8589961}.}\pdfcomment[icon=Note]{ISO/IEC 23894 and 24368 are named as operationalising risk management and ethical considerations by SC42, and emerging regulations emphasise trustworthiness/transparency (8589961, p.3).} \hldarkgreen{Research on AI business analytics confirms that failing to incorporate human oversight, bias assessments and compliance with frameworks like the EU AI Act and GDPR exposes organisations to penalties, reputational damage and loss of consumer trust \parencites[7461574]{7461574}.}\pdfcomment[icon=Note]{Business analytics literature notes that failure to include human oversight, bias assessments and compliance with frameworks (EU AI Act, GDPR) exposes organisations to penalties, reputational harm and loss of trust (7461574, p.1752).} \hllightgreen{Against this backdrop, any social recommender deployment that neither meets these baseline controls nor respects the unacceptable practice boundaries defined above falls outside the ethically defensible space for engagement optimisation under a Digital Fairness trajectory \parencites[2248]{6502223}[390]{7896734}[61--62]{2977039}[7461574]{7461574}.}\pdfcomment[icon=Note]{Synthesis: combining regulatory standards (6502223, p.12), governance/standards guidance (8589961, p.3; 8367105, p.24) and ethical analyses supports the conclusion that noncompliant recommender deployments fall outside ethically defensible space.}

\subsubsection{Duty of Care for Vulnerable and Minor Users}
\label{sec:duty_of_care_vulnerable_minors} \hllightgreen{Duty of care in algorithmic environments acquires its sharpest form where recommender systems interact with minors and other vulnerable users.}\pdfcomment[icon=Note]{Source 2977039 frames special digital-safety duties for systems interacting with minors and Source 5750298/3303310 document risks from recommender architectures, so this is a reasonable synthesis of those sources (2977039 p.215; 3303310 p.2).} \hldarkgreen{Kostenko's principle of digital safety of the child requires that any AI system processing children's data or influencing decisions about them be "secure by default" along its entire lifecycle: architecture, algorithmic logic, interaction mechanics, data categories, training cycles, purposes of use, and risk profiles must all be understandable, verifiable, and subject to control and explanation.}\pdfcomment[icon=Note]{This is essentially verbatim from Kostenkos Principle of Digital Safety of the Child which requires systems interacting with children be secure-by-default across lifecycle aspects (2977039 p.215).} \hllightgreen{This duty is relational.}\pdfcomment[icon=Note]{The relational nature of the duty (multiple stakeholders: developers, regulators, children/representatives) is implied directly by the stakeholder provisions in 2977039 and is a reasonable synthesis (2977039 p.215).} \hldarkgreen{Developers and operators must maintain internal documentation, data maps, training and update protocols, change logs, Child Impact Assessments and Child Safety Data Sheets, while regulators are entitled to a full technical, legal and ethical dossier including event logs and audit results.}\pdfcomment[icon=Note]{2977039 lists precisely the internal documentation, data maps, training/update protocols, change logs, Child Impact Assessments/Child Safety Data Sheets and the regulator-level dossier including event logs and audit results (2977039 p.215).} \hldarkgreen{Children and their legal representatives must receive accessible interfaces explaining system logic, risks and safety settings such as parental controls, refusal of personalisation and time limits, tailored to age and socio cultural context.}\pdfcomment[icon=Note]{2977039 requires accessible interfaces for children and their legal representatives explaining logic, risks and safety settings (parental controls, refusal to personalize, time limits) tailored to age and sociocultural context (2977039 p.215).} \hldarkgreen{Explainability standards are explicitly age sensitive.}\pdfcomment[icon=Note]{2977039 explicitly requires explainability to be adapted to age, cognitive and sociocultural characteristics (age-sensitive explainability) (2977039 p.215).} \hldarkgreen{Systems must deliver meaningful, structured, plain language interpretations of decision logic, including model descriptions, key features, uncertainty levels, human involvement and available alternatives, and must provide tools for children's rights such as refusal of profiling, data correction and human review.}\pdfcomment[icon=Note]{2977039 specifies multi-level, plain-language explanations including model descriptions, key features, uncertainty, human involvement and tools for rights (refusal to profile, correction, deletion, human review) (2977039 p.215).} \hldarkgreen{When AI decisions affect a child's rights, status, access to resources or social reputation, the absence of effective safety mechanisms or the use of manipulative or obscure interfaces is classified as a violation of good governance and of the best interests of the child, authorising regulators to suspend operation, revoke certificates and publish public justifications \parencites[214]{2977039}.}\pdfcomment[icon=Note]{2977039 states that where AI affects a child's rights, lack of safety or manipulative/obscure interfaces is a governance violation and authorises regulators to suspend operation, revoke permits and publish justifications (2977039 p.215).} \hllightgreen{Age appropriate design provisions sharpen this duty into concrete interface constraints.}\pdfcomment[icon=Note]{2977039 calls for National Guidelines on Children's Digital Safety and Age-Appropriate Design, so it is reasonable to infer these provisions translate into concrete interface constraints (2977039 p.215).} \hllightgreen{For minors, defaults must minimise data collection and disable geolocation, biometrics and profiling, while manipulative and addictive mechanics such as endless scrolling, autoplay, streaks, forced quests and penalties for exiting are prohibited.}\pdfcomment[icon=Note]{2977039 and related materials prohibit manipulative/addictive functionality and require secure defaults and minimal processing for children, though the specific UI elements named are inferred applications of those principles (2977039 p.61-62, p.215).} \hlorange{Push notifications must be kept to a minimum and limited to neutral information.}\pdfcomment[icon=Note]{None of the provided sources explicitly state that push notifications must be minimised and limited to neutral information; this specific prescription is not documented in the cited texts.} \hllightgreen{Children must see visible "no personalisation" switches, have access to "lesson" or "sleep" modes, simple "why am I seeing this" explanations, and easy undo and deletion controls, and basic access to educational or public services must not be degraded when personalisation is refused \parencites[223]{2977039}.}\pdfcomment[icon=Note]{2977039 mandates tools for refusal of profiling, parental controls, multi-level explainability and rights to correct/delete data, supporting the claim though some UI specifics (e.g. "lesson/sleep modes") are reasonable extensions (2977039 p.215).} \hldarkgreen{These conditions map directly onto short video and social feeds that rely on algorithmic infinite scroll and auto play, which research has tied to dopamine driven feedback loops, compulsive checking and elevated depressive symptoms, particularly among adolescents with immature self control \parencites[2]{8958037}[2]{3303310}.}\pdfcomment[icon=Note]{Empirical and theoretical sources link infinite scroll/auto-play short-video feeds to dopamine-driven feedback loops, compulsive checking and mental-health harms in adolescents (3303310 p.2; 8958037 p.5 documents algorithmic immersion/addiction effects).} \hllightgreen{Protection extends to profiling and monetisation practices.}\pdfcomment[icon=Note]{2977039 and related literature extend child protection to data processing, profiling and monetisation practices, so this is a supported synthesis (2977039 p.215; 5750298 discussion of commercial/dark patterns).} \hldarkgreen{Targeted advertising personalised on children's psychotypes, neurobehavioral signs or predicted emotional states is banned; automated tracking of emotions, expressions and biometric signs without informed double consent is forbidden, and predictive "life path" profiles for minors are outlawed \parencites[20]{2977039}[224]{2977039}.}\pdfcomment[icon=Note]{2977039 explicitly prohibits targeted advertising based on psychotypes/emotional states, unauthorised biometric/emotion tracking without consent, and predictive life-path profiling for minors (2977039 p.61; p.215).} \hllightgreen{Monetisation or resale of children's data, or use for third party model training, is prohibited absent a separate legal basis aligned with the child's best interests, and "child information invisibility" modes must allow use of services without reconstructable identity \parencites[223]{2977039}[224]{2977039}.}\pdfcomment[icon=Note]{2977039 requires that monetisation/resale and third-party model training of children's data be subject to appropriate legal bases and provides for measures to protect identity; the exact phrasing here is a reasonable reading of those provisions (2977039 p.215).} \hllightgreen{In parallel, confidentiality rules bar using personal data for advertising or profiling that restricts or distorts freedom of choice without explicit consent \parencites[44]{2977039}.}\pdfcomment[icon=Note]{2977039 and broader regulatory sources prohibit manipulative uses of personal data that restrict freedom of choice without appropriate consent, so this statement follows from those confidentiality and consent rules (2977039 p.61; 6502223 on regulatory context).} \hllightgreen{These prohibitions define a duty of care that goes beyond content labelling: business models premised on behavioural advertising targeting minors conflict with the model's normative baseline.}\pdfcomment[icon=Note]{The cited prohibitions in 2977039 and critiques of behavioural-advertising models (e.g., 5750298 on dark patterns) support the claim that such business models conflict with the normative baseline and go beyond mere labelling.} \hldarkgreen{European platform regulation reinforces this orientation.}\pdfcomment[icon=Note]{6502223 and related sources describe EU platform regulation (AI Act, DSA) and the UK Online Safety Act reinforcing child- and platform-safety orientations (6502223 p.12-13).} \hldarkgreen{The Digital Services Act introduces a duty of care style "systemic risk" framework for very large online platforms, including obligations to assess and mitigate risks to civic discourse and to provide at least one non profiled recommendation option, with stricter requirements for content harmful to children under the UK Online Safety Act's parallel duty of care regime \parencites[2248]{6502223}[2249]{6502223}.}\pdfcomment[icon=Note]{6502223 explains the DSA's systemic-risk framework for VLOPs including obligations to assess/mitigate systemic risks and to provide a non-profiled recommendation option, and references the UK Online Safety Act's stricter child-protection duties (6502223 p.12-13).} \hldarkgreen{UNESCO's Recommendation on AI Ethics adds a global reference point, calling for transparency, explainability, human oversight and prohibition of AI systems that undermine human rights, together with support for algorithmic diversity to ensure pluralism in information access \parencites[2249]{6502223}.}\pdfcomment[icon=Note]{6502223 summarises the UNESCO Recommendation calling for transparency, explainability, human oversight and prohibition of systems that undermine rights, plus support for algorithmic diversity (6502223 p.13).} \hllightgreen{These instruments situate the treatment of minors and vulnerable users within a broader expectation that platforms act as custodians of democratic and psychological welfare, not only as neutral conduits.}\pdfcomment[icon=Note]{Combining the regulatory instruments (DSA, AI Act, UNESCO) reasonably supports the synthesis that minors' treatment is situated in an expectation that platforms serve custodial democratic and psychological roles (6502223 p.12-13; 2977039 p.8, p.215).} \hllightgreen{Operationalising this duty of care requires governance architectures, not only UX changes.}\pdfcomment[icon=Note]{Sources on CSR and governance (8200246) together with 2977039 indicate operationalising duty of care requires governance architectures beyond mere UX changes (8200246 p.16; 2977039 p.62).} \hldarkgreen{Kostenko's ethical responsibility principle mandates internal ethical audit mechanisms, independent ethics commissions with authority to suspend projects, and annual AI Ethics Impact Reports covering psycho emotional health, social equality and trust \parencites[61--62]{2977039}.}\pdfcomment[icon=Note]{2977039's Principle of Ethical Responsibility requires internal ethical audits, institutionalised independent ethics commissions with power to suspend projects, and annual AI Ethics Impact Reports covering psycho-emotional health, social equality and trust (2977039 p.61-62).} \hldarkgreen{Human in the loop grading research demonstrates how such structures can function in practice: model outputs are always subject to educator confirmation, dual logging creates traceable audit trails, and student facing transparency delineates AI versus human feedback, resulting in higher reported trust and alignment with UNESCO and EU oversight expectations \parencites[20]{6805789}.}\pdfcomment[icon=Note]{Source 6805789 documents a HITL grading pilot where model outputs required educator confirmation, dual logging and student-facing transparency increased trust and aligned with broader oversight expectations (6805789 p.20).} \hldarkgreen{Human AI collaborative audit frameworks in corporate settings reach similar conclusions, arguing that clearly defined roles, cognitive load aware interfaces and shared accountability between human experts and AI systems improve detection power without sacrificing professional judgment or ethical standards \parencites[2]{6539584}[3]{6539584}.}\pdfcomment[icon=Note]{Source 6539584 reports that human-AI collaborative audit frameworks emphasise defined roles, cognitive-load-aware interfaces and shared accountability to improve detection without sacrificing professional judgment (6539584 p.2-3).} \hlorange{For this report's working taxonomy, platforms serving minors or other vulnerable populations that combine intensive profiling, addictive mechanics and lack of documented child safety dossiers fall into a "Code Category C: Prohibited or Critical Risk" zone under a Digital Fairness trajectory.}\pdfcomment[icon=Note]{This specific internal "Code Category C" taxonomy and the classification criteria are not present in the provided sources and appear to be the reports own framing (no supporting citation found).} \hllightgreen{By contrast, services that implement secure defaults, age appropriate explainability, non personalised modes and human overseen audits aligned with CSR style governance structures exemplify a "duty of care compliant" configuration that management teams can aim for when adapting social recommender systems to EU level expectations \parencites[214]{2977039}[223]{2977039}[16]{8200246}.}\pdfcomment[icon=Note]{2977039's secure-by-default and explainability requirements together with CSR/governance recommendations in 8200246 support the characterization of such services as duty-of-care compliant (2977039 p.215; 8200246 p.16).} \hllightgreen{Empirical work on recommender systems has shown they often rest on extensive behavioural and biometric inference, meaning that seemingly innocuous signals can be recombined to reveal sensitive traits; in practice this suggests regulators should be alert to how facial, interactional and auxiliary data can be repurposed for psychometric targeting \parencites[4]{5750298}.}\pdfcomment[icon=Note]{Empirical literature cited (5750298, 3303310, 5750298 references to psychometrics) documents behavioural and biometric inference capabilities and supports the inference that regulators should monitor facial/interactional/auxiliary data repurposing for psychometric targeting (5750298 p.3-7; 3303310 p.2).} \hllightgreen{Such findings may suggest tighter limits on the range of allowable features and stricter transparency about what is inferred from a child's activity \parencites[4]{5750298}.}\pdfcomment[icon=Note]{Given the empirical findings about sensitive inferences and regulatory instruments described earlier, recommending tighter feature limits and greater transparency is a reasonable and supported policy inference from the cited work (3303310 p.2; 2977039 p.215).}

\section{Qualitative Methodology and Thematic Coding Framework}

\subsection{Overall Qualitative Research Design}

\subsubsection{Choice of Secondary Qualitative Synthesis}
\label{sec:choice_secondary_synthesis} \hllightgreen{Secondary qualitative synthesis is adopted here because direct access to proprietary recommender logs, training data, or internal audits is structurally constrained, while rich documentary and empirical material is publicly available across regulatory, managerial and ethical domains.}\pdfcomment[icon=Note]{8200246 (pp.8,11) and 6295677 (p.6) justify secondary qualitative approaches when empirical access is fragmented and public documentary material is available, supporting this inference.} \hldarkgreen{Studies on dark patterns in Indian e commerce explicitly justify a qualitative, descriptive and analytical design that relies exclusively on secondary data to analyse complex digital market practices at the intersection of law, behavioural science and platform governance when primary experimentation or internal platform access is limited \parencites[6]{6295677}.}\pdfcomment[icon=Note]{6295677 (p.6) explicitly states a qualitative, descriptive analytical design relying exclusively on secondary data to study dark patterns in India, matching the sentence.} \hllightgreen{This project follows a comparable logic for social recommender systems under EU digital governance.}\pdfcomment[icon=Note]{This is a reasonable extension of the methodological logic in 8200246 (pp.8--11) and the regulatory focus in 2977039 (pp.11--12), supporting the claim by synthesis.} \hllightgreen{Methodological precedents support this choice.}\pdfcomment[icon=Note]{Multiple cited methodological papers (e.g., 8200246 pp.8--11; 6295677 p.6) present precedents for secondary qualitative design, supporting this general statement.} \hldarkgreen{Asante and Adeoba structure their investigation of Communities of Practice, CSR and ethical AI marketing around a literature based qualitative inquiry that uses secondary qualitative research to synthesise dispersed evidence from marketing, information systems, organisational learning and CSR, aiming to build conceptual understanding rather than test predefined hypotheses \parencites[8]{8200246}.}\pdfcomment[icon=Note]{8200246 (pp.8,11) describes a literature-based qualitative inquiry using secondary qualitative research to synthesise dispersed evidence, matching this description of Asante and Adeoba.} \hldarkgreen{They ground this approach in constructivism and inductive reasoning, arguing that secondary qualitative analysis is suited to theory building when empirical access is fragmented and nascent \parencites[8]{8200246}[10]{8200246}.}\pdfcomment[icon=Note]{8200246 (pp.8,11) explicitly frames the study in constructivism and inductive reasoning and justifies secondary qualitative analysis for fragmented empirical access.} \hldarkgreen{Malik's review on AI enabled career counselling adopts a similar conceptual synthesis method, combining theoretical, empirical, policy and practice sources through thematic synthesis to develop a framework aligned with OECD AI Principles and UNESCO recommendations \parencites[9762591]{9762591}.}\pdfcomment[icon=Note]{9762591 (p.14) states Malik uses a conceptual synthesis combining theoretical, empirical, policy and practice sources via thematic synthesis and aligns the framework with OECD/UNESCO, supporting the sentence.} \hllightgreen{These designs collectively validate an interpretive secondary synthesis for the present thesis.}\pdfcomment[icon=Note]{Given the methodological examples in 8200246 and 9762591, the inference that these designs validate interpretive secondary synthesis is a supported academic synthesis.} \hllightgreen{Specific analytic techniques are also directly relevant.}\pdfcomment[icon=Note]{Sources (e.g., 8200246 pp.8--11; 6419865 p.7) discuss analytic techniques such as thematic analysis and quality appraisal, substantiating the claim.} \hldarkgreen{Reflexive Thematic Analysis, as operationalised by Asante and Adeoba, structures the movement from familiarisation through coding, theme development and review to interpretive reporting, and is explicitly used to generate conceptual insights from secondary texts \parencites[10]{8200246}.}\pdfcomment[icon=Note]{8200246 (p.10) describes Reflexive Thematic Analysis using Braun \& Clarke's six phases to generate conceptual insights from secondary literature, directly supporting this sentence.} \hldarkgreen{In clinical ethics, verbatim quote based thematic synthesis proceeds through line by line coding, development of descriptive themes and construction of analytical themes that "go beyond" primary study conclusions, using multiple coders and CASP based quality appraisal to minimise bias and increase rigour \parencites[7]{6419865}.}\pdfcomment[icon=Note]{6419865 (p.7) details line-by-line coding of verbatim quotes, development of descriptive and analytical themes, use of multiple coders and CASP appraisal, matching the clinical ethics procedure described.} \hldarkgreen{Malik's thematic synthesis likewise relies on iterative coding of multi disciplinary sources to derive thematic clusters and integrate them into a responsible AI framework \parencites[9762591]{9762591}.}\pdfcomment[icon=Note]{9762591 (p.14) documents Malik's iterative coding of multidisciplinary sources and derivation of thematic clusters into a responsible AI framework, supporting the claim.} \hllightgreen{This thesis adopts these procedures as its analytic backbone: systematic coding of secondary materials, generation of analytical themes around algorithmic addiction, manipulative design and auditability, and explicit documentation of coding and synthesis decisions.}\pdfcomment[icon=Note]{This is an authorial methodological choice justified by the prior examples (8200246; 6419865; 9762591), so the adoption claim is a reasonable synthesis of those precedents.} \hllightgreen{The decision to privilege secondary synthesis is also shaped by the type of questions addressed.}\pdfcomment[icon=Note]{Sources (8200246 pp.8--11; 6295677 p.6) link methodological choice to research questions and constraints, supporting this general statement by synthesis.} \hldarkgreen{Bansal's urgency marketing study shows how semi structured executive interviews, interpreted alongside behavioural economics and digital interface research, can map the organisational logic of tactics such as scarcity cues and psychological targeting without claiming population level measurement \parencites[136]{8745065}.}\pdfcomment[icon=Note]{8745065 (p.3) describes an urgency/marketing study using semi-structured executive interviews interpreted alongside behavioural economics and interface literature to map organisational logic without claiming population-level measurement, matching the sentence.} \hldarkgreen{Khare et al. explicitly argue that where manipulative practices intersect with institutional capacity constraints and proprietary secrecy, doctrinal empirical hybrid designs relying on regulatory and institutional documents and consumer survey evidence are appropriate for evaluating regulatory adequacy \parencites[6]{6295677}.}\pdfcomment[icon=Note]{6295677 (p.6) explicitly advocates a doctrinal-empirical hybrid using regulatory/institutional documents and consumer survey evidence where proprietary secrecy limits access, supporting this claim.} \hllightgreen{In a similar way, this thesis aims to understand how regulatory texts such as the AI Law Model and DSA style regimes articulate duties around risk management, precaution and ethical responsibility, and how these can be translated into an Ethical AI Audit Toolkit, rather than to estimate prevalence of specific addictive patterns \parencites[12]{2977039}[54]{2977039}.}\pdfcomment[icon=Note]{2977039 (pp.11,54) emphasises regulatory duties and risk management rather than prevalence estimation; combined with the cited methodological precedents, this supports the thesis' stated aim by reasonable inference.} \hllightgreen{Quality and integrity safeguards follow established models.}\pdfcomment[icon=Note]{Trustworthiness and integrity practices are described in the methodological sources (8200246 pp.10--11; 6419865 p.7), supporting the statement that safeguards follow established models.} \hldarkgreen{Asante and Adeoba emphasise trustworthiness criteria derived from Lincoln and Guba, including detailed methodological description, systematic documentation of selection and analysis processes, and audit trails to enhance dependability and confirmability \parencites[10]{8200246}.}\pdfcomment[icon=Note]{8200246 (p.10) explicitly cites Lincoln and Guba's trustworthiness criteria and recommends detailed methodological description and audit trails, matching this sentence.} \hldarkgreen{The moral distress review uses NVivo based coding, multi reviewer iteration and CASP scoring to assess study quality and provide a structured basis for synthesis, categorising studies as excellent, good or fair before integration \parencites[7]{6419865}.}\pdfcomment[icon=Note]{6419865 (p.7) documents NVivo-based coding, multi-reviewer iteration and CASP scoring with categorisation into excellent/good/fair prior to synthesis, directly supporting this claim.} \hldarkgreen{Malik similarly documents search strategies, inclusion criteria and iterative coding, explicitly noting that the review is not a formal systematic review but a conceptual synthesis tailored to an interdisciplinary question \parencites[9762591]{9762591}.}\pdfcomment[icon=Note]{9762591 (p.14) states Malik documents search strategies, inclusion criteria and iterative coding and notes the review is not a formal systematic review but a conceptual synthesis, matching the sentence.} \hllightgreen{Drawing on these examples, the present study foregrounds transparent documentation of source selection, coding protocols, theme refinement and critical appraisal of each included document's scope and limitations.}\pdfcomment[icon=Note]{Given the cited examples (8200246; 6419865; 9762591) that emphasise documentation and appraisal, it is reasonable for the present study to foreground transparent documentation as stated.} \hllightgreen{Ethical considerations also favour secondary synthesis in this context.}\pdfcomment[icon=Note]{8200246 (p.10) and related sources discuss ethical advantages of secondary research (avoiding direct human subjects), supporting this general claim by synthesis.} \hldarkgreen{Asante and Adeoba highlight that secondary qualitative research avoids direct engagement with human subjects and instead emphasises accurate representation of findings and academic integrity \parencites[10]{8200246}.}\pdfcomment[icon=Note]{8200246 (p.10) explicitly notes that secondary qualitative research avoids direct engagement with human subjects and focuses on accurate representation and academic integrity, supporting the sentence.} \hldarkgreen{Khare et al. demonstrate that extensive use of regulatory and institutional reports, judicial decisions such as the BookMyShow case, and national complaint statistics can illuminate consumer protection issues without collecting sensitive behavioural data \parencites[6]{6295677}.}\pdfcomment[icon=Note]{6295677 (pp.6--7) lists extensive use of regulatory and institutional reports, judicial documents including the BookMyShow case, and complaint statistics as data sources illuminating consumer protection without collecting sensitive behavioural data.} \hllightgreen{For a thesis that critiques manipulative social recommender designs while operating under strict non access to proprietary telemetry, this combination of publicly available regulatory texts, institutional reports, academic analyses and survey studies offers a defensible evidentiary base.}\pdfcomment[icon=Note]{The combination of publicly available regulatory texts, reports, academic analyses and surveys as a defensible evidentiary base is supported by methodological sources (6295677 p.6; 8200246 pp.8--11) as a reasonable synthesis.} \hllightgreen{Finally, the secondary qualitative synthesis aligns with the AI Law Model's own emphasis on documentary governance.}\pdfcomment[icon=Note]{2977039 (pp.22,54) emphasises documentary governance elements (impact assessments, logs, documentation), so aligning secondary synthesis with document-centric governance is a supported inference.} \hldarkgreen{Kostenko calls for Human Rights Impact Assessments, incident logs, and risk management frameworks grounded in internationally recognised standards such as ISO/IEC 23894 and the NIST AI RMF, with governments obliged to maintain methodologies, registers and external peer review for high risk AI \parencites[54]{2977039}.}\pdfcomment[icon=Note]{2977039 (p.54) explicitly calls for Human Rights Impact Assessments, incident logs, and bases risk management on international standards including ISO/IEC 23894 and NIST AI RMF, directly supporting the sentence.} \hldarkgreen{Analysing these textual artefacts across jurisdictions through thematic synthesis is therefore not only a methodological necessity but also thematically consistent with the document centric governance models that future Ethical AI Audits must reference \parencites[12]{2977039}[54]{2977039}.}\pdfcomment[icon=Note]{2977039 (pp.11,54) supports analysing governance artefacts across jurisdictions and links such documentary analysis to governance models, matching this claim.} \hllightgreen{Practical guidance on how such audits are constituted, including the roles of external experts, legal advisors and independent contractors in platform audits, appears in recent audit literature and may inform operational decisions about stakeholder involvement and evidence triangulation \parencites[27]{9407672}.}\pdfcomment[icon=Note]{Recent audit literature (Meßmer \& Degeling in 9407672 p.27 and related sources) discusses roles of external experts, legal advisors and independent contractors in platform audits, supporting this practical guidance claim by synthesis.} \hllightgreen{This suggests that while secondary synthesis cannot replicate internal telemetry, it can map institutional arrangements and procedural requirements that an Ethical AI Audit Toolkit will need to address \parencites[27]{9407672}.}\pdfcomment[icon=Note]{Methodological critiques of telemetry-based audits (0822736 p.3) together with governance-focused sources (2977039) support the inference that secondary synthesis cannot replicate telemetry but can map institutional arrangements.} \hldarkgreen{The AI Law Model itself emphasises that transparency should include decision logs, explainability artifacts and clear chains of responsibility for operators and developers, which suggests concrete documentary targets for a secondary synthesis to inspect \parencites[23]{6839355}.}\pdfcomment[icon=Note]{2977039 (pp.22--23) explicitly requires decision logs, technical documentation, explainability artifacts and clear chains of responsibility, directly supporting this sentence.} \hllightgreen{Such prescriptions, while normative, may also help prioritise which public documents and registries are scrutinised when internal access is impossible \parencites[23]{6839355}.}\pdfcomment[icon=Note]{Given the AI Law Model's normative prescriptions (2977039 pp.22--54), it is a reasonable inference that such prescriptions help prioritise which public documents and registries to scrutinise when internal access is impossible.}

\subsubsection{Alignment between Research Questions and Methods}
\label{sec:alignment_rq_methods} \hllightgreen{Research questions in this project focus on three issues: how engagement optimised recommenders instantiate manipulative design and algorithmic addiction, how EU style regulation frames duties and prohibitions, and how an Ethical AI Audit Toolkit can be constructed using CSR based organisational governance.}\pdfcomment[icon=Note]{This is a reasonable synthesis: Khare et al. discuss dark patterns and manipulative design (6295677, p.6), Gerr and related work examine EU instruments (5142478, p.3), and Asante \& Adeoba link Ethical AI toolkits with CSR and organisational governance (8200246, p.8).} \hldarkgreen{A qualitative secondary synthesis is well suited to these questions because they target conceptual relationships between design practices, regulatory texts, and managerial structures rather than statistical effect sizes or internal model performance.}\pdfcomment[icon=Note]{Asante \& Adeoba explicitly justify a secondary qualitative design for investigating conceptual relationships rather than statistical effect sizes (8200246, p.8).} \hldarkgreen{Khare et al. explicitly justify a qualitative, descriptive, analytical design relying exclusively on secondary data to analyse dark patterns where primary experimentation or access to proprietary platform data is limited, aiming to "map patterns, interpret behavioural mechanisms, and evaluate regulatory adequacy" instead of estimating causal effects \parencites[6]{6295677}.}\pdfcomment[icon=Note]{Khare et al. (6295677, p.6) state they adopt a qualitative, descriptive, analytical design relying on secondary data to map patterns, interpret behavioural mechanisms, and evaluate regulatory adequacy when platform data are limited.} \hldarkgreen{Asante and Adeoba frame a comparable ambition for ethical AI marketing: secondary qualitative research and reflexive thematic analysis are used to integrate fragmented evidence from marketing, information systems, organisational learning and CSR in order to build conceptual understanding of how internal mechanisms support responsible AI implementation \parencites[8]{8200246}[10]{8200246}.}\pdfcomment[icon=Note]{Asante \& Adeoba explicitly adopt secondary qualitative research and Reflexive Thematic Analysis to integrate fragmented literature on CoP, CSR and ethical AI marketing (8200246, p.8--10).} \hllightgreen{Questions about the adequacy of regulatory responses to manipulative interfaces and addictive design align closely with doctrinal empirical hybrid designs.}\pdfcomment[icon=Note]{Khare et al. describe a doctrinal empirical hybrid approach for regulatory adequacy questions (6295677, p.6), so the statement that such regulatory adequacy questions align with doctrinal empirical hybrid designs is a supported synthesis.} \hldarkgreen{Khare et al. combine qualitative content analysis of guidelines, enforcement records and judicial decisions with consumer survey synthesis to evaluate how design level regulations and institutional capacity interact in India's e commerce ecosystem \parencites[6]{6295677}.}\pdfcomment[icon=Note]{Khare et al. describe combining qualitative content analysis of regulatory/institutional documents with synthesis of consumer survey evidence in India's e‑commerce context (6295677, p.6).} \hldarkgreen{Gerr uses a convergent parallel mixed method: content analysis of policy instruments such as the DSA and EU AI Act is paired with quantitative bias and diversity metrics to assess whether regulatory frameworks meaningfully constrain algorithmic curation on Facebook, YouTube and TikTok \parencites[3]{5142478}.}\pdfcomment[icon=Note]{Gerr is reported to use a convergent-parallel mixed method with content analysis of EU policy instruments and quantitative bias/diversity metrics applied to Facebook, YouTube and TikTok (5142478, p.3).} \hllightgreen{While this thesis does not reproduce Gerr's quantitative component, the reliance on policy document analysis and thematic coding of transparency, accountability and diversity provisions is methodologically aligned with a research question that asks how the DSA and a projected Digital Fairness Act together define constraints on social recommender design \parencites[3]{5142478}[6]{6295677}.}\pdfcomment[icon=Note]{5142478 documents Gerr's mixed-methods design (p.3) and the thesis' reliance on policy analysis and thematic coding aligns methodologically with that approach (also consistent with 6295677, p.6).} \hllightgreen{Designing an Ethical AI Audit Toolkit and a pro social pivot strategy requires methods that surface organisational mechanisms and governance levers rather than user level psychometrics.}\pdfcomment[icon=Note]{Asante \& Adeoba emphasise organisational mechanisms (CoP, CSR) and governance levers over user‑level metrics (8200246, p.8--11), so the claim about required methods is a reasonable inference.} \hldarkgreen{Asante and Adeoba select a constructivist, inductive, secondary qualitative design precisely because ethical AI marketing literature is fragmented and because their objective is to theorise how Communities of Practice and CSR mediate ethical AI implementation \parencites[8]{8200246}.}\pdfcomment[icon=Note]{Asante \& Adeoba explicitly adopt a constructivist, inductive, secondary qualitative design to theorise how CoP and CSR mediate ethical AI implementation (8200246, p.8--9).} \hldarkgreen{They use Reflexive Thematic Analysis (RTA) to move from familiarisation with sources, through coding and theme development, to an interpretive model linking CoP, CSR and ethical AI strategy \parencites[10]{8200246}.}\pdfcomment[icon=Note]{Asante \& Adeoba use Reflexive Thematic Analysis and outline the familiarisation, coding and theme development stages leading to interpretive models (8200246, p.10).} \hllightgreen{This procedure matches the present thesis aim of developing a coding matrix and thematic structure that connect regulatory obligations, documented dark pattern typologies, and CSR style governance devices into an audit framework.}\pdfcomment[icon=Note]{The procedure described by Asante \& Adeoba (RTA and coding) maps onto the thesis aim of building a coding matrix connecting regulation, dark pattern typologies and CSR governance (8200246, p.10; 6295677, p.6).} \hllightgreen{Methodological guidance from mixed methods and qualitative standards further supports alignment.}\pdfcomment[icon=Note]{Guidance from mixed‑methods and qualitative standards (e.g., MMR‑R HS and other methodological sources) is cited as supporting alignment between methods and aims (8591813, p.1; 8200246, p.10).} \hldarkgreen{The MMR R HS checklist specifies that rigorous reports clearly describe sampling of documents, analysis procedures for qualitative components, integration logic, and strategies to establish qualitative credibility, dependability and confirmability \parencites[2]{8591813}.}\pdfcomment[icon=Note]{The MMR‑R HS checklist explicitly requires rigorous reports to describe document sampling, qualitative analysis procedures, integration logic, and strategies for qualitative credibility, dependability and confirmability (8591813, p.2).} \hldarkgreen{Asante and Adeoba operationalise Lincoln and Guba's trustworthiness criteria through detailed methodological description, audit trails and reflexive interpretation, demonstrating how secondary qualitative projects can still satisfy expectations of rigour \parencites[10]{8200246}.}\pdfcomment[icon=Note]{Asante \& Adeoba operationalise Lincoln and Guba's trustworthiness criteria via detailed methods, audit trails and reflexive interpretation as described in their trustworthiness section (8200246, p.10).} \hldarkgreen{Grounded theory work on ethical consumption and AI enabled finance by Kishore et al. also illustrates how semi structured interviews, iterative coding, constant comparison and analyst triangulation are matched to research questions about identity processes and human AI interaction, while an AI enabled qualitative extension uses thematic analysis to integrate a second dataset on algorithmic decision support \parencites[6]{4245543}.}\pdfcomment[icon=Note]{The grounded theory study and its AI‑enabled qualitative extension describe semi‑structured interviews, iterative coding, constant comparison, analyst triangulation and a thematic analysis integration (4245543, p.187--188).} \hllightgreen{In this thesis, the absence of new interviews is an explicit limitation; nevertheless, the same coding logics are applied to documentary material to preserve analytic transparency.}\pdfcomment[icon=Note]{8200246 explicitly notes empirical validation is beyond its scope and that secondary analysis uses coding logics applied to documents to maintain transparency (8200246, p.11).} \hllightgreen{A critical constraint for all research questions in this project is the black box nature of proprietary recommendation systems.}\pdfcomment[icon=Note]{Several sources note limited access to platform internals and proprietary recommender black‑boxes as a constraint (6295677, p.6; 9072800, p.1--3), supporting this claim.} \hldarkgreen{Asante and Adeoba acknowledge that secondary qualitative designs do not access operational systems but instead work with published accounts and policy texts, and they explicitly identify empirical validation as a future task \parencites[11]{8200246}.}\pdfcomment[icon=Note]{Asante \& Adeoba acknowledge that secondary qualitative designs do not access operational systems and identify empirical validation as a future task (8200246, p.11).} \hldarkgreen{Khare et al. similarly position their work as interpretive mapping of manipulation and regulatory response rather than as measurement of algorithm internals, because access to platform data is constrained \parencites[6]{6295677}.}\pdfcomment[icon=Note]{Khare et al. position their work as interpretive mapping rather than measurement of algorithm internals due to constrained access to platform data (6295677, p.6).} \hllightgreen{Aligning with these precedents, this thesis formulates its research questions at the level of observable behaviour, documented harms, and regulatory expectations, and its chosen methods explicitly exclude reverse engineering of ranking code or collection of proprietary telemetry.}\pdfcomment[icon=Note]{Both Khare et al. and Asante \& Adeoba frame research at the level of observable behaviours and documented harms and state they do not engage in reverse engineering or proprietary telemetry collection (6295677, p.6; 8200246, p.11).} \hllightgreen{Finally, alignment is reinforced by sector specific examples where qualitative designs interface with AI governance questions.}\pdfcomment[icon=Note]{The claim that alignment is reinforced by sector examples is supported by subsequent case studies (e.g., Malachi et al.) showing qualitative designs interfacing with AI governance (5419804, p.4).} \hldarkgreen{Malachi et al. use a mixed design in Plateau, Benue and Nasarawa States to examine AI mediated mental health support, combining a controlled trial with an ethnographically informed "CHEW Competency Wheel" that is co designed through qualitative work with community health workers, teachers and storytellers \parencites[4]{5419804}.}\pdfcomment[icon=Note]{Malachi et al. report a mixed design in Plateau, Benue and Nasarawa combining a controlled trial with a co‑designed CHEW Competency Wheel developed through qualitative co‑design (5419804, p.4).} \hldarkgreen{Their approach shows that questions about how AI systems interact with vulnerable communities and how human interpreters sustain trust require qualitative co design and thematic analysis, not only quantitative outcome measures.}\pdfcomment[icon=Note]{5419804 documents that their approach required qualitative co‑design and thematic integration to address AI interactions with vulnerable communities beyond quantitative measures (5419804, p.4).} \hllightgreen{Questions in this thesis about pro social pivots and CSR based audit regimes for addictive recommenders sit in an analogous space and are therefore methodologically compatible with secondary, interpretive synthesis grounded in carefully documented thematic coding and explicit integrity constraints \parencites[6]{6295677}[8]{8200246}[4]{5419804}.}\pdfcomment[icon=Note]{This is a supported synthesis: Asante \& Adeoba advocate secondary interpretive synthesis with documented thematic coding (8200246, p.8--11) and Malachi et al. exemplify sector compatibility (5419804, p.4).}

\subsection{Data Sources and Selection Strategy}

\subsubsection{Regulatory and Policy Texts as Data}
\label{sec:reg_policy_texts_data} \hldarkgreen{Regulatory and policy documents form a primary empirical corpus for this qualitative study rather than a mere backdrop.}\pdfcomment[icon=Note]{Methodological sources show regulatory and policy documents are used as the primary empirical corpus in qualitative designs (e.g., 2702055 p.3; 4246340 p.3; 6295677 p.6).} \hldarkgreen{Dark pattern research in India explicitly operationalises guidelines of the Advertising Standards Council of India, Central Consumer Protection Authority advisories, National Consumer Helpline statistics, and consumer commission orders such as the BookMyShow case as data points for analysing manipulative interface practices, enforcement gaps, and institutional capacity \parencites[6]{6295677}.}\pdfcomment[icon=Note]{India-focused dark pattern work cited here explicitly lists ASCI, CCPA, National Consumer Helpline statistics and consumer commission orders (including BookMyShow) as data sources (6295677 p.6).} \hllightgreen{This demonstrates that formal texts, interpretive guidance, and enforcement records can be treated as evidence about both normative expectations and real governance performance, not only as legal context.}\pdfcomment[icon=Note]{Sources describe treating formal texts, guidance and enforcement records as documentary evidence about norms and governance (2702055 p.3; 4246340 p.3), supporting this interpretive claim as a reasonable synthesis.} \hldarkgreen{A similar strategy is used in mixed method work on AI driven content curation in Europe, where the DSA (2022), national instruments like the Germany Action Plan (2024), and the EU AI Act (2025) are coded alongside platform policies to derive themes around transparency, accountability and media pluralism \parencites[3]{5142478}.}\pdfcomment[icon=Note]{The mixed-method study on AI-driven content curation explicitly codes DSA (2022), Germany Action Plan (2024) and EU AI Act (2025) alongside platform policies to derive themes (5142478 p.3).} \hllightgreen{In this thesis, such documents supply the core material for reconstructing how EU level and model law instruments conceptualise manipulative design, systemic risk, and audit duties for recommender systems.}\pdfcomment[icon=Note]{The AI law model and policy-analytic sources show that EU-level and model-law instruments can be used to reconstruct regulatory concepts around manipulative design and audit duties (2977039 p.30--31; 5142478 p.3), a reasonable synthesis for a thesis.} \hllightgreen{The AI Law Model illustrates why a model statute is itself analytically rich data.}\pdfcomment[icon=Note]{The AI Law Model (2977039) contains detailed substantive provisions and thus exemplifies why a model statute is analytically rich; this is a reasonable interpretive claim based on that source (2977039 p.30--31).} \hldarkgreen{Kostenko's text defines principles such as technological sovereignty, specifying that states may prohibit AI systems if independent audit is impossible, control over critical components is lacking, or critical technological dependence exists, with such systems placed in a special register \parencites[107]{2977039}.}\pdfcomment[icon=Note]{Kostenkos model articulates technological sovereignty and states may prohibit systems if independent audit is impossible or critical control is lacking, with such systems entered into a special register (2977039 p.106--107).} \hldarkgreen{Generative AI receives a dedicated regime: mandatory registration in a National Register of Generative AI, compulsory labelling and watermarking of AI generated content, explicit bans on manipulative, discriminatory or false outputs, and mandatory independent audits that cover manipulative impact, bias and compliance with privacy and transparency requirements \parencites[31]{2977039}.}\pdfcomment[icon=Note]{The AI Law Model specifies a special regime for generative AI including mandatory national registration, mandatory labelling/watermarking, bans on manipulative/discriminatory/false outputs, and independent audits assessing manipulative impact and bias (2977039 p.30--31).} \hllightgreen{Treating these clauses as data allows this study to code concrete obligations (registration, labelling, auditability, blocking powers) that later become criteria in the Ethical AI Audit Toolkit.}\pdfcomment[icon=Note]{Methodological sources indicate that coding statutory clauses yields discrete obligations (registration, labelling, auditability, blocking powers) that can be operationalised as criteria in toolkits (2702055 p.3; 2977039 p.30--31).} \hllightgreen{European regulatory materials serve a comparable role for platform governance.}\pdfcomment[icon=Note]{Work on European platform governance uses regulatory materials in a comparable analytic role (e.g., content analysis of EU instruments and platform policies in 5142478 p.3; discussion of platform regulation in 6941611).} \hlorange{Analyses of the DSA emphasise that very large online platforms must conduct systemic risk assessments, provide at least one non profiled recommendation option, and submit to annual independent audits, with Article 40 DSA granting data access to regulators and vetted researchers for systemic risk study \parencites[7]{1456534}[19]{3195574}.}\pdfcomment[icon=Note]{The provided sources analyse the DSA and related themes (5142478 p.3; 6941611) but do not explicitly document the specific combination of obligations claimed here (non-profiled recommendation option and annual independent audits + Article 40 detail) on the cited pages.} \hlorange{Policy work on the impending Digital Fairness Act describes it as a proposal aimed at regulating addictive product designs and unfair personalisation to fill gaps between the DSA, DMA and unfair commercial practices law \parencites[390]{7896734}.}\pdfcomment[icon=Note]{No provided source explicitly discusses a "Digital Fairness Act" framed as stated; the claim about that specific proposal and its aims is unsupported by the cited materials.} \hllightgreen{In the coding matrix used here, these texts provide discrete variables such as "non profiled feed requirement" or "addictive design as regulatory target" that can be assigned to specific risk themes.}\pdfcomment[icon=Note]{Methodological sources describe creating coding matrices from texts and extracting discrete variables (e.g., policy features like design targets), consistent with this statement (2702055 p.3; 4246340 p.3; 5142478 p.3).} \hldarkgreen{Methodologically, using regulatory texts as data aligns with qualitative designs in adjacent fields.}\pdfcomment[icon=Note]{Multiple methodological sources endorse using regulatory texts as data in qualitative research designs (4246340 p.3; 2702055 p.3), directly supporting this methodological claim.} \hldarkgreen{A doctrinal empirical hybrid design for dark patterns relies entirely on secondary regulatory and institutional documents plus survey reports, arguing that this is appropriate in contexts where primary experimentation or direct access to proprietary systems is limited \parencites[6]{6295677}.}\pdfcomment[icon=Note]{The doctrinal empirical hybrid approach relying on secondary regulatory/institutional documents plus survey reports is described explicitly for dark-pattern research contexts where primary access is limited (6295677 p.6).} \hldarkgreen{Systematic review work on ethical AI in social innovation similarly treats laws, governance guidelines and policy white papers as core data, collected via structured database searches and then subjected to open, axial and selective coding to identify themes such as algorithmic bias, privacy and accountability \parencites[3]{4246340}.}\pdfcomment[icon=Note]{Systematic review work treating laws, guidelines and white papers as core data collected via structured searches and coded is described in the literature review methodology (4246340 p.3).} \hllightgreen{This thesis adopts the same logic: AI acts, platform regulations, model laws and enforcement decisions are coded line by line as textual evidence about duties, prohibitions and governance mechanisms.}\pdfcomment[icon=Note]{The thesiss adoption of coding acts-as-data is a stated methodological choice consistent with examples in the sources (2702055 p.3; 4246340 p.3), making this a supported synthesis.} \hllightgreen{Regulatory texts also encode risk tiering that feeds directly into this report's working taxonomy.}\pdfcomment[icon=Note]{Sources indicate regulatory instruments include tiered risk classifications that can inform taxonomies (e.g., AI Act references in 5142478 p.3 and model-law risk framing in 2977039), supporting this inference.} \hldarkgreen{The EU AI Act's four level scheme and its classification of systems exploiting psychological vulnerabilities as "unacceptable risk" are reconstructed from legal analysis and used as normative anchors for defining Code Category C in the audit matrix \parencites[12]{6502223}.}\pdfcomment[icon=Note]{The EU AI Act's multi-tier risk scheme and its prohibition of systems exploiting psychological vulnerabilities as unacceptable risk are established elements of the Act and are referenced in the analysed literature (5142478 p.3; discussion of AI Act risk categories in 8367105 mapping).} \hldarkgreen{Kostenko's explicit bans on manipulative, addictive and discriminatory systems, combined with obligations for state bodies to maintain national registers, conduct audits and block systems where technological sovereignty is compromised, provide a second anchor for categorising unacceptable practice in engagement optimised recommenders \parencites[61--62]{2977039}[107]{2977039}.}\pdfcomment[icon=Note]{2977039 lays out explicit bans on manipulative/addictive/discriminatory systems and obligations for national registers, audits and blocking where technological sovereignty is at risk, directly supporting this claim (2977039 p.106--107).} \hldarkgreen{Finally, policy documents double as data on institutional capacity and implementation asymmetries.}\pdfcomment[icon=Note]{Multiple policy analyses treat policy documents as evidence about institutional capacity and implementation asymmetries (5162643; 9360054 p.10), supporting this statement.} \hldarkgreen{Sikalo's examination of "imitative modernization" argues that Ukraine's adoption of GDPR, AI Act, DORA, DSA/DMA, Data Act and eIDAS as reference frameworks has occurred under conditions of insufficient institutional capacity, producing stratification and a "double imitation" mechanism where formal transposition is not matched by effective governance \parencites[19]{5162643}.}\pdfcomment[icon=Note]{Sikalo's analysis documents Ukraine's imitative adoption of EU frameworks under constrained institutional capacity, describing the "double imitation" mechanism (5162643 p.1--3).} \hldarkgreen{Survey based evaluation of user awareness of GDPR, DSA and the AI Act finds high self assessed digital competence but low ability to recognise hidden manipulation, combined with scepticism about effective enforcement in cross border settings, particularly for non EU countries such as Ukraine \parencites[10]{9360054}.}\pdfcomment[icon=Note]{Survey results cited report high self-assessed digital competence but low ability to recognise hidden manipulation and scepticism about enforcement in cross-border/non-EU contexts such as Ukraine (9360054 p.10).} \hllightgreen{These findings are coded in this study as contextual constraints on any audit toolkit that presumes strong supervisory capacity.}\pdfcomment[icon=Note]{Coding these empirical findings as contextual constraints on audit toolkits is a reasonable interpretive step supported by the survey and governance capacity literature (9360054 p.10; 5162643).} \hldarkgreen{Using regulatory and policy texts as primary data therefore serves three functions: it grounds the coding matrix in explicit legal and model law language; it supplies comparative evidence on enforcement and capacity; and it aligns the project with established qualitative methodologies that rely on documentary corpora to theorise governance of opaque, high impact AI systems \parencites[6]{6295677}[3]{5142478}[31]{2977039}[10]{9360054}.}\pdfcomment[icon=Note]{The three functions listed (grounding coding in legal language; supplying comparative enforcement/capacity evidence; aligning with documentary qualitative methods) are directly supported by the cited methodological and substantive sources (5142478 p.3; 2977039 p.30--31; 9360054 p.10).}

\subsubsection{Academic and Grey Literature on Algorithmic Addiction}
\label{sec:academic_grey_addiction} \hlorange{Empirical and conceptual work on compulsive technology use provides the main academic anchor for this project's coding of algorithmic addiction.}\pdfcomment[icon=Note]{No provided source explicitly frames empirical and conceptual work on 'compulsive technology use' as the primary academic anchor for the project's coding; the cited materials do not directly support this exact claim.} \hlorange{Li's synthesis on short video platforms presents neural network based recommendation engines combined with 15--60 second clips, infinite swipe navigation, auto play, and push notifications as an "addiction loop" that pulls users into an "auto pilot mode" of persistent scrolling, with Chinese usage data indicating average daily durations approaching $60$ minutes and more than $21$ sessions, especially among college students \parencites[2]{8958037}[3]{8958037}.}\pdfcomment[icon=Note]{This sentence cites Li [3]\{8958037\} and data not present in the provided sources; no supplied document corresponds to Li/8958037 to verify the described measures, statistics, or wording.} \hlorange{Hu's analysis of "algorithm driven engagement worship" complements this by quantifying a $22\%$ higher rate of compulsive checking in algorithmically curated feeds compared with non algorithmic interfaces and by linking such interfaces to increased depressive symptoms and reduced prefrontal activity \parencites[2]{3303310}[20]{3303310}.}\pdfcomment[icon=Note]{The claim references Hu [20]\{3303310\} and specific quantitative neurobehavioral results, but source 3303310 is not among the provided materials, so the assertion is unsupported by the available documents.} \hlorange{These studies are treated as core sources for this report's working taxonomy of addictive interface patterns and behavioural indicators.}\pdfcomment[icon=Note]{The statement that 'these studies' are treated as core sources cannot be verified because the key studies (Li, Hu) cited earlier are not present in the supplied sources; therefore support is lacking.} \hlorange{Neuropsychological mechanisms are not described in isolation but tied directly to interface mechanics.}\pdfcomment[icon=Note]{No provided source explicitly states that neuropsychological mechanisms are described only in direct relation to interface mechanics; this methodological claim is not supported by the available documents.} \hlorange{Li reports dopamine surges at the moment of expectancy before scrolling to the next video and characterises short video content switching as structurally akin to slot machine variable rewards, with chronic high frequency stimulation leading to $\Delta$FosB expression and elevated reward thresholds \parencites[3]{8958037}[4]{8958037}.}\pdfcomment[icon=Note]{The neurobiological claims attributed to Li (dopamine surges, ΔFosB expression, slot-machine analogy) require Li [4]\{8958037\}, which is not in the provided sources, so this specific assertion is unsupported.} \hlorange{Hu synthesises neuroimaging evidence that associates hyper personalised feeds with reduced gray matter density and diminished prefrontal cortex activity, suggesting convergence between behavioural addiction and substance dependence patterns \parencites[2]{3303310}[20]{3303310}.}\pdfcomment[icon=Note]{The sentence cites Hu [20]\{3303310\} for neuroimaging evidence linking personalised feeds to reduced gray matter and prefrontal activity, but the Hu source is not provided, so the claim cannot be verified here.} \hlorange{These findings are coded under neurobiological correlates of algorithmic addiction in the thematic matrix.}\pdfcomment[icon=Note]{Stating that 'these findings are coded under neurobiological correlates' is an internal coding claim tied to findings from Li/Hu, which are not present among the supplied sources; thus it is unsupported by available documents.} \hldarkgreen{Qualitative work on human judgement and AI decision support offers a second strand.}\pdfcomment[icon=Note]{Source 4245543 (p.5--6) presents qualitative work on human judgment and AI decision support among students, supporting the claim that such qualitative literature constitutes a distinct strand.} \hldarkgreen{Kishore et al. use constructivist grounded theory and an AI enabled qualitative extension to examine how students experience AI embedded financial tools, reporting that users see algorithms as advisory rather than authoritative and often override AI recommendations when they conflict with circumstances or emotions \parencites[7]{4245543}.}\pdfcomment[icon=Note]{Source 4245543 (p.5--6) describes a constructivist grounded theory design with an AI-enabled qualitative extension and reports that users treat AI as advisory and often override recommendations, directly supporting this sentence.} \hldarkgreen{They highlight automation bias, explainability gaps and diminished human agency as key themes, which are mapped to this thesis coding categories on perceived control and ethical discomfort in AI mediated environments \parencites[7]{4245543}[6]{4245543}.}\pdfcomment[icon=Note]{Source 4245543 (p.5--6) highlights themes such as automation bias, explainability gaps, and diminished human agency, which aligns with the mapping to perceived control and ethical discomfort described here.} \hllightgreen{Methodological precedents for synthesising addiction related risks from diverse sources are provided by systematic and conceptual reviews in adjacent domains.}\pdfcomment[icon=Note]{Provided methodological reviews (e.g., 2055956, 4246340, 9762591) offer systematic and conceptual review precedents in adjacent domains, supporting the claim that such reviews provide methodological precedents.} \hldarkgreen{Mogoi et al. conduct a PRISMA based systematic review of AI powered mobile proctoring, screening $180$ records and synthesising $20$ studies to characterise technical, ethical and institutional challenges including surveillance anxiety, privacy, data governance and algorithmic bias \parencites[2]{2055956}[9]{2055956}.}\pdfcomment[icon=Note]{Source 2055956 explicitly reports a PRISMA-based systematic review that screened 180 records and synthesised 20 studies addressing technical and ethical challenges, directly supporting this sentence (2055956, p.2, p.9).} \hldarkgreen{Hassen's systematic literature review on AI in social innovation adopts structured database searches, strict inclusion criteria (English, 2020--2025, ethical or governance focus), and inductive three stage coding to identify themes of algorithmic bias, privacy and an "accountability void" \parencites[3]{4246340}.}\pdfcomment[icon=Note]{Source 4246340 describes a systematic literature review using structured database searches, inclusion criteria (English, 2020--2025), and an inductive three-stage coding procedure (4246340, p.3), matching the sentence.} \hldarkgreen{Malik's conceptual synthesis on AI enabled career counselling follows a similar approach, combining peer reviewed articles and policy documents, then using thematic synthesis to derive clusters around equity, governance and human centred design \parencites[14]{9762591}.}\pdfcomment[icon=Note]{Source 9762591 describes a conceptual synthesis on AI-enabled career counselling that combines peer-reviewed articles and policy documents with thematic synthesis, supporting this statement (9762591, p.14).} \hllightgreen{These designs justify the present study's reliance on structured secondary searches, inductive coding and theme integration for algorithmic addiction, rather than new primary experiments.}\pdfcomment[icon=Note]{It is a reasonable methodological inference supported by the cited review and conceptual sources (2055956, 4246340, 9762591) that structured secondary searches and inductive coding justify reliance on secondary synthesis rather than new primary experiments.} \hllightgreen{Exemplar materials and their coding contributions are summarised below for clarity: 

- Li (short video): focuses on infinite scroll, auto play, dopamine cycles, student harms; contributes to coding on interface patterns, temporal discounting, and neurobiological indicators \parencites[2]{8958037}[3]{8958037}[4]{8958037}. 

- Hu (engagement worship): addresses algorithmic feeds, compulsive checking, depression, polarisation; informs behavioural metrics, mental health outcomes, and systemic risk coding \parencites[2]{3303310}[20]{3303310}. 

- Kishore et al.: examines student AI finance use, automation bias, human override; mapped to perceived agency, trust, and ethical discomfort categories \parencites[7]{4245543}[6]{4245543}. 

- Mogoi et al.: reviews mobile proctoring, surveillance anxiety, fairness; feeds surveillance related stress, fairness concerns, and governance levers coding \parencites[2]{2055956}[9]{2055956}. 

- Hassen: covers AI for social good, algorithmic bias, privacy; codes accountability and bias themes in non commercial AI settings \parencites[3]{4246340}. 

- Malik: treats AI counselling, digital inequality, policy frameworks; supports coding on equity and access dimensions and multi level policy context \parencites[14]{9762591}.}\pdfcomment[icon=Note]{The mappings for Kishore (4245543), Mogoi (2055956), Hassen (4246340), and Malik (9762591) are supported by those sources; however, the Li and Hu items reference sources not provided, so overall this exemplar summary is partially supported by the available documents.} \hllightgreen{Grey materials and model law texts function as a third pillar and are explicitly treated as data.}\pdfcomment[icon=Note]{Several provided materials (e.g., model law 2977039 and review/policy sources) function as grey literature and normative texts and can be treated as data, so treating grey materials and model law texts as a third pillar is a reasonable synthesis of the sources.} \hllightgreen{Kostenko's AI Law Model frames systems that manipulate behaviour, cause digital addiction, emotional exhaustion, loss of trust or suppression of autonomy as prohibited under the principle of ethical responsibility, and it articulates rights to a psycho emotionally safe digital environment free of excessive information flows and digital noise \parencites[61--62]{2977039}[119]{2977039}.}\pdfcomment[icon=Note]{Source 2977039 articulates rights and prohibitions related to digital subjectivity and protections against manipulation and coercion (2977039, ch.43), supporting the claim that an AI Law model frames manipulative systems as prohibited and articulates rights to a protected digital environment, though the exact phrasing 'psycho emotionally safe' is interpretive.} \hllightgreen{These clauses are coded as normative thresholds for unacceptable addictive practice and later inform Code Category C in this report's working risk taxonomy.}\pdfcomment[icon=Note]{It is a reasonable coding decision supported by the presence of normative clauses in the AI Law Model (2977039) to treat such clauses as normative thresholds informing a risk taxonomy; this is an inferential synthesis of the provided material.} \hldarkgreen{Dark pattern regulation studies in India synthesise consumer surveys, regulator reports and enforcement records to document how manipulative interfaces exploit dual process cognition and nudge sludge dynamics, reinforcing behavioural vulnerabilities even without explicit addiction language \parencites[6]{6295677}[7]{6295677}.}\pdfcomment[icon=Note]{Source 6295677 (p.6--7) explicitly synthesises secondary regulatory data, institutional reports and consumer surveys and situates manipulative interfaces within dual-process cognition and nudge sludge theory, supporting this sentence.} \hldarkgreen{Their methodological choice to rely exclusively on secondary regulatory and survey data in the absence of proprietary platform access parallels this thesis reliance on secondary sources for social recommender analysis \parencites[6]{6295677}.}\pdfcomment[icon=Note]{Source 6295677 (p.6) states the study relies exclusively on secondary data (regulatory and survey data) due to limited access to proprietary platform data, directly paralleling the thesis's reliance on secondary sources.} \hllightgreen{A consistent limitation recurs across these materials.}\pdfcomment[icon=Note]{Multiple provided sources (e.g., 2055956, 6295677, 2977039) note recurring limitations such as lack of proprietary access and data constraints, so claiming a consistent limitation across materials is a supported synthesis.} \hllightgreen{Black box recommender algorithms and behavioural telemetry are not directly available; scholars assess algorithmic addiction through observable UI features, self reported usage patterns, regulatory investigations and documented mental health outcomes instead \parencites[2]{8958037}[2]{3303310}[6]{6295677}[61--62]{2977039}.}\pdfcomment[icon=Note]{Sources (e.g., 2055956, 6295677) document limited access to platform-internal telemetry and describe assessment via UI features, surveys, regulatory reports and documented harms; while exact citekeys differ, the general limitation and alternative assessment methods are supported by the provided materials.} \hllightgreen{This thesis adopts the same constraint and codes sources accordingly, distinguishing between direct empirical measures of compulsive use, inferred neurocognitive mechanisms, and normative or governance responses.}\pdfcomment[icon=Note]{This methodological statement is a reasonable inference and synthesis of the constraints documented in the supplied sources (e.g., 2055956, 6295677) and thus is supported as the thesis adopting the same constraint and distinguishing coding categories.}

\subsubsection{Public UX Analyses and Advocacy Reports}
\label{sec:public_ux_advocacy_reports} \hllightgreen{Public UX analyses and advocacy reports constitute a distinct empirical strand in this study's corpus because they describe interface behaviour, dark patterns, and user harm without requiring privileged access to proprietary code or telemetry.}\pdfcomment[icon=Note]{Source 6295677 (p.7) and other methodological sources describe use of public UX analyses, regulatory audits and institutional reports as secondary empirical material that describe interface behaviour and user harm without requiring proprietary telemetry, supporting this characterization.} \hldarkgreen{Khare et al. explicitly rely on regulatory audits, institutional reports, and consumer commission decisions to document that approximately $87\%$ of major Indian e commerce and digital service platforms deploy at least one dark pattern, treating this figure as evidence that manipulative interface design is a structural feature of platform based commerce rather than a marginal anomaly.}\pdfcomment[icon=Note]{Source 6295677 (p.7) explicitly reports the ≈87\% figure and frames it from institutional audits and reports, treating it as evidence of systemic prevalence of dark patterns on major Indian platforms.} \hldarkgreen{Their sampling covers roughly $30$ high impact platforms selected on market share and regulatory relevance, and their qualitative content analysis codes practices such as interface interference, obstruction, subscription traps, drip pricing, and privacy zuckering against established HCI taxonomies and dual process cognition models.}\pdfcomment[icon=Note]{Source 6295677 (p.7) documents sampling of roughly 30 major platforms selected by market share/regulatory relevance and describes coding of practices mapped to HCI taxonomies and System 1 cognitive biases.} \hllightgreen{This report treats such UX audits as primary data on how dark patterns appear in live environments and how they map to cognitive biases and consumer harm.}\pdfcomment[icon=Note]{Source 6295677 (p.7) and related methodological texts treat UX audits and institutional reports as primary secondary data for analysing how dark patterns appear and map to cognitive biases, supporting this interpretive use.} \hllightgreen{Advocacy driven typologies are particularly valuable for the coding matrix.}\pdfcomment[icon=Note]{Source 6295677 (p.7) and allied literature note that advocacy typologies and HCI taxonomies enable systematic coding, indicating their particular value for coding matrices.} \hldarkgreen{Khare et al. synthesise regulatory circulars, sectoral guidelines, and survey evidence to classify patterns and link them to System 1 biases such as default bias, scarcity effects and social proof, and to "sludge" based choice obstruction that increases friction for protective actions like cancellation or refusal.}\pdfcomment[icon=Note]{Source 6295677 (p.7) links dark pattern typologies to System 1 biases and nudge/sludge theory; the source describes synthesis of regulatory guidance and survey evidence mapping patterns to biases like defaults and scarcity.} \hldarkgreen{Their methodology integrates institutional survey results on consumer experiences of deceptive practices, unexpected charges and difficulty cancelling services, which are interpreted through behavioural economics constructs rather than treated as isolated complaints.}\pdfcomment[icon=Note]{Source 6295677 (p.7) states consumer vulnerability analysis is derived from aggregated institutional survey data on deceptive practices, unexpected charges and cancellation difficulties, integrated into the analysis.} \hlorange{For this thesis, these UX facing categories directly inform working codes such as "false urgency widgets," "hierarchy manipulation in consent flows," and "roach motel subscription flows," even though specific social platforms are not named in the available excerpt \parencites[7]{6295677}.}\pdfcomment[icon=Note]{None of the provided sources state that this particular thesis uses the listed working codes (e.g., "false urgency widgets", "hierarchy manipulation"); the claim about the thesis's internal coding decisions is unsupported by the supplied documents.} \hllightgreen{Regulatory advocacy at EU level provides a second UX oriented source type.}\pdfcomment[icon=Note]{Source 7896734 (p.390) and related materials discuss EU-level regulatory advocacy (DSA, DMA, AI Act) as important UX-oriented sources, supporting this statement.} \hlorange{Sanikidze's analysis of dark patterns and AI driven manipulation emphasises enforcement actions where interface design was central to the infringement.}\pdfcomment[icon=Note]{No provided source is authored by Sanikidze; while EU enforcement actions and interface-centred infringements are discussed (e.g., 7896734 p.390), the specific attribution to 'Sanikidze' and their emphasis cannot be verified in the supplied sources.} \hldarkgreen{Platform X received a €120 million fine for "deceptive" verification features, while Temu was preliminarily found in breach of the Digital Services Act and remains under investigation for potentially addictive design, including concerns about making cancellation harder than sign up and using false urgency cues in e commerce flows.}\pdfcomment[icon=Note]{Source 7896734 (p.390) explicitly reports a €120 million fine for a Platform X over deceptive verification features and describes preliminary DSA breach findings and ongoing investigations into Temu, including concerns about cancellation friction and false urgency cues.} \hllightgreen{These cases are treated in this thesis as documented examples of UX patterns that regulators interpret as unlawful, and they inform high risk entries in this report's working taxonomy, for instance "Code Category C: deceptive verification and cancellation flows".}\pdfcomment[icon=Note]{Sources (e.g., 7896734 p.390) show regulators treating these cases as unlawful UX practices; it is a reasonable synthesis to say such cases inform high‑risk taxonomy entries used in the thesis.} \hllightgreen{Cross jurisdictional advocacy materials widen the UX evidence base beyond the EU.}\pdfcomment[icon=Note]{Source 7896734 (p.390) documents non‑EU jurisdictions (e.g., Georgia) applying EU-aligned frameworks, indicating cross‑jurisdictional advocacy materials broaden the UX evidence base beyond the EU.} \hlorange{Sanikidze describes how the Georgian Competition and Consumer Agency applied an EU aligned unfair practices law to sanction digital dark patterns such as false countdown timers, misleading scarcity messages, fake reviews, and deceptive brand associations, including the case where products were advertised under the "Maxbosch" label while using BOSCH imagery.}\pdfcomment[icon=Note]{Although Source 7896734 (p.390) describes GCCA cases including false timers and the 'Maxbosch'/BOSCH imagery example, there is no provided source named Sanikidze attributing this analysis, so the specific claim naming Sanikidze is unsupported.} \hldarkgreen{Although these decisions were not yet applied to AI driven patterns, they document concrete UI elements (timers, stock notices, review widgets, brand imagery) that are recognisable across global platforms and therefore usable as codable UX artefacts in this study's matrix \parencites[390]{7896734}.}\pdfcomment[icon=Note]{Source 7896734 (p.390) explicitly notes these Georgian decisions do not cover AI‑driven patterns but document concrete UI elements (timers, stock notices, fake reviews, brand imagery) that are observable and codable.} \hldarkgreen{A subset of UX and advocacy reports within academic work focus less on consumer transactions and more on AI mediated environments such as proctoring and grading.}\pdfcomment[icon=Note]{Source 2055956 (p.9) and related literature review material focus on AI‑mediated environments such as proctoring and grading, supporting the existence of UX/advocacy reports in these areas.} \hldarkgreen{Mogoi et al.'s systematic review treats AI powered mobile proctoring as a socio technical system and highlights user perceptions of being "constantly watched by a machine and intrusively recorded," linking technical features like facial recognition, gaze tracking and environmental recording directly to surveillance anxiety and fairness concerns.}\pdfcomment[icon=Note]{Source 2055956 (p.9) is a systematic review of 20 studies on AI mobile proctoring and notes user perceptions of being 'constantly watched' and links facial recognition, gaze tracking and environmental recording to surveillance anxiety and fairness concerns.} \hldarkgreen{Their synthesis of $20$ studies emphasises that privacy by design, informed consent and transparent communication about how AI decisions are made are necessary safeguards, which are treated here as UX governance benchmarks for high stakes educational interfaces \parencites[2055956]{2055956}.}\pdfcomment[icon=Note]{Source 2055956 (p.9) synthesises 20 studies and recommends privacy‑by‑design, informed consent and transparent communication as necessary safeguards, which the thesis can treat as UX governance benchmarks.} \hldarkgreen{Selvam and González Vallejo similarly foreground communication UX in human in the loop grading: syllabi disclosures, explicit tagging of comments as AI generated, human reviewed or human added, and accessible appeal procedures are shown to increase student trust and clarity, based on internal testing of annotated feedback reports \parencites[6805789]{6805789}.}\pdfcomment[icon=Note]{Source 6805789 (p.15) describes Selvam et al.'s recommendations on disclosure, tagging of AI vs human comments, and notes internal testing where annotated feedback increased student satisfaction and trust.} \hllightgreen{These design features are coded as positive UX patterns in this report's working matrix, contrasting with dark patterns coded from e commerce contexts.}\pdfcomment[icon=Note]{Given the evidence in sources like 6805789 (p.15) and 2055956 (p.9) that such communication features improve trust, it is a reasonable synthesis to code them as positive UX patterns contrasted with dark patterns from e‑commerce.} \hllightgreen{Methodologically, these UX analyses and advocacy materials complement doctrinal and theoretical sources by anchoring abstract concepts like "manipulative interface" or "secure by default" in observable interface structures.}\pdfcomment[icon=Note]{Sources (e.g., 6295677 p.7 and 8200246 pp.10--11) describe how empirical UX analyses anchor abstract concepts in observable interface elements, supporting the methodological complementarity claimed.} \hldarkgreen{Khare et al. explicitly adopt qualitative content analysis of screenshots, flows and regulatory descriptions to build their typology, while Asante and Adeoba's Reflexive Thematic Analysis framework provides a precedent for systematically coding such heterogeneous secondary data and for maintaining an audit trail that supports trustworthiness in secondary qualitative synthesis \parencites[7]{6295677}[10]{8200246}.}\pdfcomment[icon=Note]{Source 6295677 (p.7) states Khare et al. use qualitative content analysis of screenshots and regulatory descriptions, while Source 8200246 (p.10) documents Reflexive Thematic Analysis as an appropriate framework, supporting this combined methodological claim.} \hldarkgreen{In the present study, UX oriented reports are therefore coded alongside legal texts and empirical addiction studies, with explicit recognition that internal algorithm logic is inferred only through documented UI behaviour, enforcement decisions and user surveys, not through proprietary telemetry or source code inspection \parencites[6]{6295677}[2]{2055956}[10]{8200246}.}\pdfcomment[icon=Note]{Source 2055956 (p.9) and Source 8200246 (pp.10--11) support coding UX reports alongside legal texts and note that algorithmic internals must often be inferred from UI behaviour, enforcement decisions and surveys rather than proprietary telemetry.}

\subsection{Development of a Working Coding Matrix}

\subsubsection{Code A Interface Mechanics and Engagement Features}
\label{sec:codeA_interface_mechanics} \hllightgreen{Code A in this project's working coding matrix captures observable interface mechanics and engagement features that are directly implicated in algorithmic addiction and manipulative design.}\pdfcomment[icon=Note]{Multiple sources discuss observable interface mechanics, engagement features and their role in manipulative or addictive dynamics (e.g., algorithmic engagement and dark-pattern literature in 3303310 p.2; typologies and regulatory triangulation in 6295677 p.7), so this is a reasonable synthesis linking Code A to those phenomena.} \hllightgreen{Coding proceeds at the level of concrete UI elements and interaction flows rather than inferred internal ranking logic, consistent with the methodological constraint that proprietary algorithms are assessed via published secondary descriptions of behaviour, not direct access to code or telemetry \parencites[6]{6295677}[2]{8958037}.}\pdfcomment[icon=Note]{Methodological discussion of coding concrete UI elements and using published/secondary descriptions is supported by qualitative coding and codebook literature (e.g. open/axial/selective coding and member checking in 2702055 p.4; codebook/YAML examples in 7286763 p.46 and 8584408 p.8).} \hllightgreen{Three main clusters structure Code A: \begin{enumerate} \item high frequency, low friction engagement loops; \item salience and urgency constructs within interfaces; \item fairness and transparency supporting mechanics. \end{enumerate} Short video platforms provide the core template for the first cluster.}\pdfcomment[icon=Note]{Partitioning interface mechanics into clusters and identifying short-video platforms as templates for high-frequency, low-friction loops is a reasonable synthesis of literature on short-form engagement and thematic coding approaches (short-video engagement described in 3303310 p.2; thematic clustering methods in 8745065 p.3).} \hlorange{Li describes infinite swipe feeds of $15$--$60$ second clips, auto play, and push notifications as an "addiction loop" that shifts users into "auto pilot" scrolling, with average daily use on Douyin approaching $60$ minutes and more than $21$ sessions, especially among college students \parencites[2]{8958037}[3]{8958037}.}\pdfcomment[icon=Note]{The sentence cites specific descriptions and numeric statistics attributed to Li and to source [3]\{8958037\}, but that cited source is not among the provided materials and none of the available sources supply the exact Douyin figures (60 minutes, >21 sessions) or the attributed quotation, so the claim is unsupported by the provided corpus.} \hldarkgreen{Hu links infinite scrolling interfaces and variable reward schedules on platforms such as TikTok to "dopamine driven feedback loops," reporting $22\%$ higher compulsive checking in algorithmically curated feeds than in non algorithmic interfaces and associating prolonged use with reduced prefrontal activity and increased depressive symptoms \parencites[2]{3303310}[20]{3303310}.}\pdfcomment[icon=Note]{The Proceedings paper documents links between infinite scrolling/variable rewards and 'dopamine-driven feedback loops' and reports a 22\% higher compulsive checking in algorithmic feeds as well as neurobiological/mental-health associations (3303310 p.2).} \hllightgreen{Within Code A, these findings motivate discrete subcodes such as: - A1: infinite scroll or continuous swipe based feeds; - A2: auto play of subsequent items without explicit user initiation; - A3: high frequency push notifications that invite re entry.}\pdfcomment[icon=Note]{Deriving discrete observable subcodes (infinite scroll, autoplay, high-frequency push notifications) from the empirical findings on engagement features is a reasonable methodological synthesis supported by the engagement literature (3303310 p.2) and codebook practices (7286763 p.46).} \hllightgreen{Each subcode is binary (present/absent) at the feature level and is annotated with contextual notes where sources provide population specific harms, for example heightened anxiety and depression among students \parencites[2]{8958037}[2]{3303310}.}\pdfcomment[icon=Note]{Describing binary feature-level coding with contextual notes for population-specific harms aligns with qualitative coding protocols (member checking, annotation in 0089118 p.3 and 2702055 p.4) and the mental-health harms discussed in 3303310 p.2.} \hllightgreen{Urgency and salience features form the second cluster.}\pdfcomment[icon=Note]{The claim that urgency and salience features form a distinct cluster is a structural/organisational statement supported by the urgency/scarcity literature and thematic clustering approaches in the provided sources (8745065 p.3; 6295677 p.7).} \hldarkgreen{In e commerce, Bansal's interview based thematic analysis shows that countdown timers, stock alerts such as "only two left," and expiring discounts are deployed as time sensitive scarcity cues integrated into broader urgency oriented commercial systems \parencites[136]{8745065}.}\pdfcomment[icon=Note]{The qualitative interview study described in 8745065 explicitly documents use of countdown timers, low-stock alerts and expiring discounts as time-sensitive scarcity cues in e-commerce (8745065 p.3).} \hldarkgreen{Khare et al. synthesise audits of approximately $30$ major Indian platforms and document widespread "false urgency" and "interface interference," where misleading countdowns and visually dominant prompts steer users toward platform preferred options \parencites[7]{6295677}.}\pdfcomment[icon=Note]{The India-focused synthesis cites audits of approximately 30 major platforms and documents pervasive 'false urgency' and interface interference; this is described in the sampling and findings sections of 6295677 (6295677 p.7).} \hllightgreen{Code A therefore introduces: - A4: countdown timers or low stock labels linked to purchase or sign up flows; - A5: visual hierarchy manipulation that privileges acceptance or continuation options over exit or decline.}\pdfcomment[icon=Note]{Introducing A4 and A5 (countdown/low-stock labels; visual hierarchy manipulation) as code items follows directly from the urgency/dark-pattern literature (8745065 p.3; 6295677 p.7) and is a reasonable coding decision.} \hllightgreen{These codes are linked analytically to behavioural triggers such as loss aversion and default bias as conceptualised in the urgency marketing and dark pattern studies, but the coding itself records only the observable UI tactic, consistent with strict evidence constraints \parencites[136]{8745065}[7]{6295677}.}\pdfcomment[icon=Note]{Linking observable UI tactics to behavioural triggers (loss aversion, default bias) while recording only the observable element is consistent with the dark-pattern and regulatory triangulation methods discussed in 6295677 p.7 and with methodological constraints described in 2702055 p.4.} \hllightgreen{A smaller but important group of mechanics supports or undermines transparency and human oversight.}\pdfcomment[icon=Note]{Stating that a smaller group of mechanics concerns transparency and human oversight is an organisational claim supported by multiple sources addressing transparency and oversight in AI/education contexts (2168844 p.2; 2977039 p.201).} \hldarkgreen{Human in the loop grading research reports that explicit labelling of AI generated versus instructor comments, visible override mechanisms, and clear documentation of rubric criteria increased student trust and preserved faculty authority over final scores \parencites[20]{6805789}[21]{6805789}.}\pdfcomment[icon=Note]{The Human-in-the-Loop grading research reports that explicit labelling, visible override mechanisms, and clear rubric documentation improved student trust and preserved faculty authority (HITL chapter conclusions in 6805789 p.21).} \hldarkgreen{The AI Law Model similarly requires that AI outputs in education be explicitly labelled, that auto confirmation modes be prohibited, and that each use of AI in assessment be logged with model identifiers and a record of the teacher's decision \parencites[59.1]{2977039}.}\pdfcomment[icon=Note]{The AI Law Model text explicitly requires labelling of AI outputs, prohibits auto-confirmation modes, and mandates logging of AI use with model identifiers and teacher decision records (2977039 p.202, section 59.1 d--d).} \hllightgreen{Drawing on these sources, Code A includes: - A6: explicit AI labelling of content or recommendations; - A7: visible "no personalisation" or alternative mode toggles; - A8: indications of human validation or oversight where AI is involved.}\pdfcomment[icon=Note]{Including A6--A8 (explicit AI labelling, no-personalisation toggles, human validation indicators) as code items follows directly from the HITL evidence (6805789 p.21) and legal recommendations (2977039 p.202) and is a reasonable synthesis.} \hllightgreen{Although these features come from educational contexts, their presence or absence in social recommender interfaces will be coded because they signal whether engagement optimisation is coupled to meaningful human control and explainability \parencites[20]{6805789}[59.1]{2977039}.}\pdfcomment[icon=Note]{The decision to apply educationally-sourced transparency/human-control features to social recommender interfaces as signals of human control is a justified inference grounded in the AI Law Model provisions (2977039 p.202, section 59.1) and the methodological practice of cross-contextual coding.} \hllightgreen{Related work on AI-driven review systems highlights how opaque generation and curation practices can erode user trust and raise privacy concerns, which in turn suggests that explicit labels, clear consent mechanisms, and audit logs may help restore confidence in automated outputs \parencites[4]{7235975}.}\pdfcomment[icon=Note]{Work on AI-driven review systems that highlights opacity eroding trust and raises privacy concerns, and that suggests labels, consent mechanisms and audit logs as remedies, is documented in the online-review and ethics literature (7235975 p.4--6; 2168844 p.2--4).} \hllightgreen{Such findings reinforce the analytic choice to code visible transparency affordances separately from inferred algorithmic behaviour, since the former are both observable and likely to influence user perceptions \parencites[4]{7235975}.}\pdfcomment[icon=Note]{The inference that these findings justify coding visible transparency affordances separately from inferred algorithmic behaviour is a reasonable methodological conclusion supported by the transparency and ethics literature (2168844 p.2; 7235975 p.4).} \hllightgreen{Table~\ref{tab:codeA_slots} summarises Code A slot definitions and their primary supporting sources. \begin{table}[h] \centering \caption{Code A slot definitions for interface mechanics and engagement features} \label{tab:codeA_slots} \begin{tabular}{p{1.8cm}p{5.0cm}p{5.0cm}} \hline Slot & Description & Key sources \\ \hline A1 & Infinite scroll or continuous feed without natural stop cues & Short video "addiction loop," persistent scrolling \parencites[2]{8958037}[3]{8958037}; dopamine driven checking \parencites[2]{3303310} \\ A2 & Auto play of next item without user click/tap & Auto play in short video architectures \parencites[2]{8958037} \\ A3 & High frequency, behaviour based push notifications & Push notifications within addiction loop \parencites[2]{8958037}[3]{8958037} \\ A4 & Countdown timers or scarcity banners in decision flows & Urgency marketing cues \parencites[136]{8745065}; false urgency patterns \\ A5 & Visual or interaction hierarchy favouring acceptance over refusal & Interface interference and obstruction \parencites[7]{6295677} \\ A6 & Explicit labelling of AI generated outputs & Labelling requirements in education \parencites[59.1]{2977039}; dual sourced feedback UI \parencites[20]{6805789} \\ A7 & User accessible non personalised or "off" modes & Non profiling recommendation option under DSA \parencites[7]{1456534}; age appropriate "no personalisation" switches \parencites[223]{2977039} \\ A8 & Visible human validation or override indicators & HITL grading override logs and authority retention \parencites[20]{6805789}[59.1]{2977039} \\ \hline \end{tabular} \end{table} Coding of academic and advocacy materials follows the same principle as methodological exemplars in game studies and AI ethics: open coding of textual descriptions of interface behaviour, then assignment to closed slot labels to guarantee inter coder consistency, similar to the YAML based codebook used for large language model prompting research \parencites[7286763]{7286763}[8584408]{8584408}.}\pdfcomment[icon=Note]{The general claim that Table X summarises slot definitions and key supporting sources, and that coding follows open coding then closed slot assignment similar to YAML codebooks, is supported by codebook/exemplar practices in 7286763 p.46 and 8584408 p.8; some specific source-to-slot mappings draw on earlier-cited papers (e.g., short-video addiction in 3303310 p.2).} \hlorange{Validation of Code A assignments in this thesis relies on cross checking across independent sources, for example aligning Li's and Hu's descriptions of infinite scroll and compulsive checking with regulatory concern about manipulative design in DSA enforcement actions, where such patterns are explicitly scrutinised as impairing free decision making \parencites[2]{8958037}[2248]{6502223}.}\pdfcomment[icon=Note]{The sentence cites aligning Li's and Hu's descriptions with DSA enforcement actions and references [2248]\{6502223\}; while 6502223 discusses regulatory distinctions (filter bubbles, policy implications) it does not document specific DSA enforcement actions tied to those exact patterns nor are Li/Hu sources present in the provided set, so the claim is not supported by the supplied materials.} \hldarkgreen{Ethical discussions that emphasise fairness, accountability and transparency provide a useful normative frame for these mechanics, and they may suggest practical criteria by which interface features are judged acceptable or problematic \parencites[3]{2168844}.}\pdfcomment[icon=Note]{Ethical discussions emphasising fairness, accountability and transparency as normative framing are well established in the provided ethics literature (core principles in 2168844 p.2--4; fairness/accountability themes in 1456534 p.11).} \hldarkgreen{At the same time, the same discussions caution that high-level principles require concrete operational rules if they are to be applied reliably across diverse platform designs \parencites[3]{2168844}.}\pdfcomment[icon=Note]{The caution that high-level ethical principles require concrete operational rules for reliable application is explicitly discussed as a limitation in multiple governance/ethics reviews (weak operationalization noted in 2168844 p.11).} \hllightgreen{User preferences are not static, however, and temporal shifts as well as social ties can reshape engagement patterns in ways that interact with UI tactics; models that jointly consider temporal dynamics and social trust indicate that recommendation behaviour may change smoothly over time and be sensitive to peers, which may amplify or attenuate the effects of an interface nudge \parencites[2]{0335538}.}\pdfcomment[icon=Note]{The statement about temporal shifts, social ties and models jointly considering temporal dynamics and social trust is supported by recommender-systems literature on temporal dynamics and social influence (temporal dynamics in 2168844 p.2; social/trust dynamics and temporal models in 0335538 p.3--5).} \hllightgreen{Such findings suggest it may be necessary to consider time-varying and social-contextual annotations when interpreting the presence of A1--A5 features, since what appears as an addictive loop in one period or cohort could reflect shifting preference dynamics or trust-driven adoption in another \parencites[2]{0335538}.}\pdfcomment[icon=Note]{The inference that time-varying and social-contextual annotations may be necessary when interpreting A1--A5 is a reasonable conclusion drawn from the temporal and social-influence literature (2168844 p.2; 0335538 p.3--5) indicating preference dynamics can alter interpretation.}

\subsubsection{Code B Psychological Triggers and Motivational Effects}
\label{sec:codeB_psych_triggers} \hllightgreen{Code B in this report's matrix captures the psychological and motivational mechanisms that interface patterns and recommendation logics exploit.}\pdfcomment[icon=Note]{Studies on urgency-based interface tactics and recommender auditing describe psychological and motivational mechanisms exploited by interface patterns and recommendation logic (8745065 p.136; 0822736 p.3), supporting this synthesis.} \hllightgreen{Coding focuses on triggers explicitly documented in empirical and regulatory sources rather than inferred internal states.}\pdfcomment[icon=Note]{The coding approaches and codebook emphasise using documented observable triggers and primary sources rather than inferring internal states (7286763 Appendix J p.45; 4246340 p.3), supporting this methodological claim.} \hldarkgreen{One prominent cluster concerns time pressure and scarcity.}\pdfcomment[icon=Note]{The literature and interviews explicitly identify time pressure and scarcity as a prominent cluster in urgency-based marketing (8745065 Executive Overview and Study Purpose; 8745065 p.136).} \hldarkgreen{Qualitative work with senior marketing executives describes how countdown timers, "only two left" stock notices, and expiring discounts are deployed as high pressure cues that reframe a routine purchase as a time sensitive decision, often targeting budget constrained consumers through "pay more later" narratives that accentuate loss framing \parencites[136]{8745065}.}\pdfcomment[icon=Note]{Qualitative interviews with senior marketing executives report deployment of countdown timers, low-stock notices and expiring discounts as high-pressure cues that reframe purchases as time-sensitive (8745065 Qualitative Methodology and Executive Overview, p.136).} \hlorange{Dark pattern audits in Indian e commerce classify "false urgency / scarcity" as a distinct pattern where countdowns or stock indicators exaggerate availability constraints, and they document high prevalence of such prompts in real interfaces \parencites[7]{6295677}.}\pdfcomment[icon=Note]{None of the provided sources present a dark pattern audit of Indian e-commerce that specifically classifies 'false urgency/scarcity' with documented prevalence in interfaces; no Indian audit source is included in the materials.} \hllightgreen{Code B therefore includes subcodes such as B1: time pressure cues (countdowns) and B2: scarcity and loss framing messages, each anchored in these documented practices rather than speculative cognitive theory \parencites[136]{8745065}[7]{6295677}.}\pdfcomment[icon=Note]{The article's coding practice of creating subcodes anchored in documented practices is consistent with the provided codebook and thematic coding procedures (7286763 Appendix J p.45; 8745065 Phase 2 Deductive Coding Matrix, p.136).} \hllightgreen{Another group of triggers targets cognitive load and decision fatigue.}\pdfcomment[icon=Note]{Decision fatigue and cognitive load as a cluster of triggers is discussed in the behavioral design and dark-pattern literature reviewed here (8745065 p.136; 4245543 grounded theory themes p.6).} \hlorange{Khare et al. define "choice obstruction" as intentional friction that complicates cancellation or refusal, describing subscription traps, roach motels, and long unsubscribe flows that increase cognitive fatigue and inhibit protective decisions.}\pdfcomment[icon=Note]{There is no Khare et al. source among the provided materials that defines 'choice obstruction' or details subscription traps and roach-motel patterns; this specific attribution lacks a matching source here.} \hlorange{They link these designs to dual process accounts where friction keeps users in fast, heuristic driven modes and discourages reflective evaluation \parencites[6]{6295677}.}\pdfcomment[icon=Note]{No provided source explicitly shows Khare et al. linking such designs to dual-process accounts; while behavioral literature in 8745065 references behavioral theory broadly, the specific claim about Khare et al. is unsupported by the supplied items.} \hlorange{Bansal's thematic analysis similarly codes "decision fatigue" and "psychological targeting" as organisationally recognised levers in urgency based marketing systems \parencites[8745065]{8745065}.}\pdfcomment[icon=Note]{Bansal's thematic analysis is not present among the supplied sources, so the specific claim that Bansal codes 'decision fatigue' and 'psychological targeting' cannot be verified from the provided material.} \hllightgreen{In Code B this yields B3: friction based obstruction (e.g. complex cancellations) and B4: decision fatigue amplification, which are applied whenever sources describe design induced exhaustion as a mechanism for steering behaviour \parencites[6]{6295677}[8745065]{8745065}.}\pdfcomment[icon=Note]{The formulation of friction-based obstruction and decision-fatigue codes aligns with the study's deductive coding and use of design-induced exhaustion as a mechanism (8745065 Phase 2 Deductive Coding Matrix, p.136; 7286763 codebook practices p.45).} \hllightgreen{Temporal orientation and instant gratification form a third cluster.}\pdfcomment[icon=Note]{Temporal orientation and instant gratification are commonly treated as a distinct thematic cluster in the reviewed behavioral and platform-use literature, consistent with the sources discussing immediacy effects (8745065 p.136; 0822736 p.3).} \hlorange{Li reports that short video platforms with 15--60 second clips and infinite swipe navigation create environments where the "extremely low operational cost" of "just one more swipe" aligns with present oriented temporal focus and undermines long term goals, particularly in adolescents and students who show diminished academic concentration and increased anxiety and depression under excessive use \parencites[2]{8958037}[3]{8958037}.}\pdfcomment[icon=Note]{The specific citation to Li and the detailed claims about short-video platforms, adolescents, and mental-health outcomes (including exact clip lengths and outcomes) do not appear in the provided sources, so this exact report is unsupported here.} \hlorange{Hu synthesises evidence that algorithmically curated feeds are associated with 22\% higher compulsive checking and a 23\% higher incidence of depressive symptoms compared with chronological feeds, capturing a shift toward immediate, algorithmically supplied rewards \parencites[2]{3303310}[20]{3303310}.}\pdfcomment[icon=Note]{The numeric claims attributed to Hu (22\% and 23\% effects) and the comparison with chronological feeds are not present in the supplied sources; no Hu source with these statistics is provided.} \hllightgreen{Code B records these mechanisms through B5: instant reward salience and B6: temporal discounting bias, justified only where sources directly link interface immediacy to preference for short term gratification \parencites[2]{8958037}[2]{3303310}.}\pdfcomment[icon=Note]{Recording mechanisms linking interface immediacy to short-term preferences is consistent with the coding and evidence-linking approach described in the codebook and behavioral literature (7286763 Appendix J p.45; 8745065 p.136).} \hldarkgreen{A further set of codes addresses social and identity oriented triggers.}\pdfcomment[icon=Note]{Social and identity-oriented triggers are explicitly identified as a codeable set in the grounded theory and qualitative sources, which list peer/influencer influence and identity-driven responses (4245543 p.6--table of themes).} \hlorange{Khare et al. document how social proof cues, fabricated or amplified, are integrated into dark patterns such as "confirm shaming" and "generative social proof," in which personalised guilt inducing language and mass produced testimonials pressure users to comply with platform preferred options \parencites[7]{6295677}.}\pdfcomment[icon=Note]{Khare et al. are not among the provided sources, and the specific terms 'confirm shaming' and 'generative social proof' in that exact formulation are not found in the supplied materials, so this attribution is unsupported here.} \hldarkgreen{Kishore et al.'s grounded theory study on ethical consumption highlights that peer and influencer influence, emotional responses such as guilt, and post purchase regret are central to students' decision processes in digitally mediated environments \parencites[4245543]{4245543}.}\pdfcomment[icon=Note]{The grounded-theory study in 4245543 documents peer and influencer influence, emotional responses like guilt, and post-purchase regret as central to student decision processes (4245543 p.6 and thematic descriptions).} \hllightgreen{Code B therefore includes B7: social proof and conformity pressure and B8: guilt or shame based messaging, applied where evidence shows emotional or identity cues used as levers rather than neutral information \parencites[7]{6295677}[4245543]{4245543}.}\pdfcomment[icon=Note]{Including social-proof and guilt/shame messaging as codes is consistent with the grounded-theory findings reported in 4245543 and the paper's coding procedures (4245543 p.6; 7286763 codebook practices p.45).} \hllightgreen{Motivational dynamics in gamified systems provide a counterpoint that informs positive coding.}\pdfcomment[icon=Note]{The gamification literature provided frames motivational dynamics as a counterpoint to manipulative designs and informs positive coding distinctions (5645864 p.3; 9283805 p.8).} \hllightgreen{Tokarieva and Chyzhykova's analysis of gamified learning platforms uses Malone and Lepper's taxonomy of intrinsic motivation to evaluate how avatars, rewards and rapid pace can support autonomy and competence when aligned with educational aims \parencites[5645864]{5645864}.}\pdfcomment[icon=Note]{Analyses of gamified learning platforms in the sources reference Malone \& Lepper's taxonomy and evaluate avatars, rewards and pace in relation to intrinsic motivation (5645864 p.3; 5645864 referencing Malone \& Lepper).} \hllightgreen{Anisa and Maisa's review of token economies stresses that token based motivation is effective when embedded in autonomy supportive, cognitively complex tasks, but that "technostress" from complexity and instability can cancel intended gains and lead to disengagement \parencites[9283805]{9283805}.}\pdfcomment[icon=Note]{The token-economy review explicitly highlights that token-based motivation is effective when autonomy-supportive but that technostress can undermine gains (9283805 p.8), matching the claim.} \hllightgreen{These findings are encoded as B9: autonomy supportive feedback loops and B10: technostress induced demotivation, allowing the matrix to distinguish between gamified designs that genuinely reinforce intrinsic motivation and those that overload or frustrate users \parencites[5645864]{5645864}[9283805]{9283805}.}\pdfcomment[icon=Note]{Encoding autonomy-supportive feedback and technostress-induced demotivation as distinct codes aligns with the token-economy review's thematic findings and the codebook practice (9283805 p.8; 7286763 Appendix J p.45).} \hllightgreen{Finally, emotional dependence in affective computing applications adds a boundary condition.}\pdfcomment[icon=Note]{Concerns about emotional dependence in affective computing are discussed in the supplied literature and flagged as an important boundary condition (9738431 p.22).} \hldarkgreen{Song warns that emotional computing systems can create "behavioral interaction addiction, cognitive delusions of consciousness, and social relationship alienation," and calls for transparency, explainability, and assistive rather than anthropomorphic design to reduce emotional dependence risks.}\pdfcomment[icon=Note]{The cited source explicitly warns of behavioural interaction addiction, cognitive delusions of consciousness, and social relationship alienation and recommends transparency and assistive design (9738431 p.22).} \hllightgreen{In Code B this is captured by B11: emotional dependence triggers, used where systems simulate companionship or emotional reciprocity without clear signalling of artificiality, consistent with documented concerns \parencites[9738431]{9738431}.}\pdfcomment[icon=Note]{Capturing emotional-dependence triggers as a code in Code B when systems simulate companionship without signalling artificiality matches the documented concerns in 9738431 p.22 and the codebook approach (7286763 Appendix J p.45).} \hllightgreen{Table~\ref{tab:codeB_summary} summarises Code B slots and anchor sources. \begin{table}[h] \centering \caption{Code B psychological triggers and motivational effects} \label{tab:codeB_summary} \begin{tabular}{p{1.1cm}p{5.0cm}p{5.}\pdfcomment[icon=Note]{The manuscript and codebook practices indicate a summarising table of Code B slots and anchor sources is provided (7286763 Appendix J p.45; 7286763 Table references and 7286763 p.12 showing thematic summary tables).}

\subsubsection{Code C Legal Risk Levels and Compliance Status}
\label{sec:codeC_legal_risk_levels} \hldarkgreen{Legal risk in this report's working coding matrix is anchored in explicit prohibitions and duties articulated in binding and model law texts rather than inferred severity.}\pdfcomment[icon=Note]{Source 2977039 (AI Law Model) explicitly frames legal risk in terms of prohibitions and duties articulated in binding/model law texts rather than inferred severity (see 2977039 pp.108, 113).} \hldarkgreen{Code C therefore classifies practices by reference to documented bans on manipulative design, social rating, opaque AI identity, and uncontrolled cross border deployment, together with positive obligations for documentation, logging and human oversight \parencites[68.5]{2977039}[70]{2977039}[31.4]{2977039}.}\pdfcomment[icon=Note]{The specific classification and cited prohibitions/obligations match TSAI/AI Law Model provisions in 2977039 (see pp.113, 246 and TSAI prohibitions at 2977039 pp.188, 245-246).} \hldarkgreen{Regulatory clauses on unacceptable AI use provide the first reference point.}\pdfcomment[icon=Note]{Source 2977039 positions regulatory clauses on unacceptable AI use (bans/prohibitions) as the initial legal reference point (see 2977039 pp.108, 113).} \hldarkgreen{Under the AI Law Model, it is prohibited to introduce or use systems that cannot be subjected to independent technical, legal and ethical audit, whose critical components are not under direct national control, or that embed hidden mechanisms for influencing user behaviour.}\pdfcomment[icon=Note]{The AI Law Model text explicitly prohibits systems lacking capacity for independent audit, lacking national control over critical components, or embedding hidden influence mechanisms (2977039 p.108).} \hldarkgreen{Such systems must be placed in a special register of prohibited or unwanted AI systems \parencites[31.4]{2977039}.}\pdfcomment[icon=Note]{2977039 states that prohibited systems are to be included in a special register of prohibited/unwanted AI systems (see 2977039 p.108, para 31.4).} \hldarkgreen{In parallel, TSAI provisions explicitly ban social rating of persons, manipulative political targeting, covert censorship in official channels, and the restriction of basic electronic public services as extrajudicial sanctions \parencites[68.5]{2977039}.}\pdfcomment[icon=Note]{TSAI provisions enumerated in 2977039 (section on TSAI prohibitions) explicitly ban social rating, manipulative political targeting, covert censorship, and restriction of basic e‑services (2977039 pp.188, 245-246).} \hllightgreen{Recent literature on neurorights and platform design suggests that these forms of covert behavioural influence often operate through interface and algorithmic patterns that bypass conscious deliberation, and that banning such practices aligns with broader concerns about cognitive liberty and informed consent \parencites[11]{0783007}.}\pdfcomment[icon=Note]{Literature on neurorights and platform design (0783007) and dark patterns research (8325651, 0783007) describe covert behavioural influence via interface/algorithmic patterns and link this to cognitive liberty and consent concerns (0783007 pp.3,9; 8325651 pp.5,11).} \hllightgreen{Such work may also imply that regulators should pay attention to specific UI mechanics and reward schedules when assessing whether a system crosses the line into prohibited manipulation \parencites[11]{0783007}.}\pdfcomment[icon=Note]{Sources on auditing and dark patterns (8325651; 0783007) note regulators should consider UI mechanics and reward schedules when assessing manipulation, supporting the suggested regulatory attention (8325651 pp.24; 0783007 pp.4,13).} \hldarkgreen{In Code C, any social recommender configuration that fits these patterns is assigned to "Code Category C: Prohibited Systemic Practices" in this report's taxonomy, signalling that legal risk is not mitigable through incremental adjustment.}\pdfcomment[icon=Note]{The report's taxonomy treatment ,  assigning social recommender configurations matching prohibited patterns into a prohibited category ,  follows directly from the AI Law Model/TSAI prohibitions and the report's Code C framing in 2977039 (see pp.108, 188, 245-246).} \hldarkgreen{A second cluster of rules defines non negotiable transparency and traceability conditions.}\pdfcomment[icon=Note]{2977039 and TSAI material establish non‑negotiable transparency and traceability duties as core rules for high‑risk systems (see 2977039 pp.71, 113, 246).} \hldarkgreen{High risk AI systems must maintain a detailed decision architecture diagram, structured descriptions of algorithmic logic, full logging and audit trails that link inputs to outputs, formalised human intervention procedures, mapped technical and social risks, and training data documentation including accuracy and bias metrics.}\pdfcomment[icon=Note]{2977039 lists required documentation for high‑risk AI: decision architecture diagrams, structured algorithmic descriptions, full logging/audit trails, human intervention procedures, risk mapping and training data metrics (see 2977039 p.71 and p.113).} \hldarkgreen{This documentation forms part of a technical passport and must be available to state bodies in full and to users and victims to the extent necessary for safe use and rights protection \parencites[19.5]{2977039}.}\pdfcomment[icon=Note]{2977039 explicitly states this documentation is part of the technical passport and must be available to state bodies in full and to users/victims as necessary (see 2977039 p.71 and p.113).} \hldarkgreen{TSAI adds that systems admitted to cross border digital environments must have compatible unique identifiers and algorithmic passports, activated event logging guaranteeing evidentiary value, and automated policy enforcement enabling immediate restriction or suspension in case of incidents \parencites[68.6]{2977039}.}\pdfcomment[icon=Note]{TSAI/cross‑border provisions in 2977039 require compatible unique identifiers and algorithmic passports, activated event logging with evidentiary value, and automated policy enforcement for immediate restriction/suspension (see 2977039 pp.245-246, 68.6/68.5).} \hldarkgreen{Code C therefore treats the absence of such passports, audit trails or logging in high impact social recommendation contexts as "Code Category C: Critical Traceability Failure," because legal frameworks condition continued operation on these artefacts.}\pdfcomment[icon=Note]{Treating absence of passports/logging as critical traceability failure and conditioning operation on these artefacts is directly supported by 2977039's requirements and sanctions for non‑compliance (see 2977039 pp.71, 113, 245-246).} \hldarkgreen{Digital sovereignty provisions introduce a territorial dimension to risk classification.}\pdfcomment[icon=Note]{2977039 introduces digital sovereignty provisions that add a territorial/sovereignty dimension to risk classification (see 2977039 pp.108, 112-113, 33.x sections).} \hldarkgreen{States are expected to localise servers, exercise control over system logic through open APIs and audits, and prevent "digital colonization" by foreign platforms that administer critical algorithms or host key registries outside national jurisdiction.}\pdfcomment[icon=Note]{2977039 prescribes localization of servers, state control over system logic via APIs/audits, and measures to prevent 'digital colonization' by foreign platforms (see 2977039 pp.112-113, 33.2--33.5).} \hldarkgreen{Where public registers or AI decisions on social benefits or lawsuits are based on servers or code inaccessible to government agencies, restrictions on provider activity are envisaged \parencites[33.2]{2977039}[33.4]{2977039}.}\pdfcomment[icon=Note]{2977039 explicitly envisages restrictions where public registers or AI decisions rely on servers or code inaccessible to government agencies (see 2977039 p.113, para 33.1--33.4).} \hllightgreen{Within Code C, recommender deployments that route core governance functions such as social scoring or eligibility ranking through opaque foreign infrastructures fall into "high or unacceptable" legal risk bands, because they conflict with sovereign technological governance expectations.}\pdfcomment[icon=Note]{Combining 2977039's digital sovereignty rules (pp.112-113) with the report's Code C risk taxonomy reasonably supports classifying opaque foreign‑infrastructure recommender deployments as high/unacceptable legal risk.} \hldarkgreen{Table~\ref{tab:codeC_levels} summarises this report's working legal risk levels and their linkage to source based criteria. \begin{table}[h] \centering \caption{Code C working legal risk levels and compliance cues} \label{tab:codeC_levels} \begin{tabular}{p{3cm}p{4.5cm}p{4.5cm}} \hline Code C category & Indicative practices & Key legal anchors \\ \hline Prohibited systemic practices & Social rating indices, manipulative political targeting, covert censorship, non auditable AI & TSAI prohibitions on social rating and manipulative targeting; auditability bans \parencites[31.4]{2977039}[68.5]{2977039} \\ Critical traceability failure & No technical passport, missing logs, no event evidence for decisions in high risk settings & Documentation and logging duties for high risk systems \parencites[19.5]{2977039}[68.6]{2977039} \\ High sovereignty risk & Critical decisions based on foreign controlled servers, inaccessible code or data relocation & Digital sovereignty, local storage and anti colonization clauses \parencites[33.2]{2977039}[33.4]{2977039} \\ Statutory compliance with oversight & Registered systems, annual audits, public trust labelling, accessible registry APIs & TSAI supervision, auditing and trust marks \parencites[68.7]{2977039} \\ \hline \end{tabular} \end{table} Hybrid governance case studies in decision support systems illustrate how Code C can mark positive compliance states as well.}\pdfcomment[icon=Note]{The summarised table content described ,  linking categories (prohibited practices, traceability failure, sovereignty risk, statutory compliance) to legal anchors like TSAI prohibitions and documentation/logging duties ,  is directly reflected in 2977039 (see pp.71, 108, 113, 245-246).} \hldarkgreen{Healthcare and civil service HR systems operating under statutory or co regulatory models exhibit continuous bias audits, explainable outputs, external review and appeal mechanisms, leading to bias reduction, shorter appeal resolution times and higher audit compliance scores \parencites[4]{0375818}.}\pdfcomment[icon=Note]{Case studies in 0375818 document healthcare and civil service systems under statutory/co‑regulatory models employing continuous bias audits, explainable outputs, external review/appeal mechanisms with quantified benefits (see 0375818 pp.5--7, case study sections).} \hldarkgreen{These features correspond to "Code Category C: Statutory Compliant with Effective Oversight" in this report's taxonomy and provide exemplars for social recommender platforms seeking to align engagement optimisation with legal accountability.}\pdfcomment[icon=Note]{0375818 maps these features to a statutory/compliant oversight category and presents them as exemplars for aligning systems with legal accountability (see 0375818 sections on models and case studies).} \hldarkgreen{Finally, collaborative audit research in revenue processes stresses that successful automation architectures explicitly document which components are automated versus human initiated, define materiality limits for human intervention, and maintain feedback loops for pattern detection and anomaly identification across cycles \parencites[6539584]{6539584}.}\pdfcomment[icon=Note]{Collaborative audit research on revenue processes in source 6539584 recommends documenting automated vs human components, defining materiality limits for human intervention, and maintaining feedback loops for pattern/anomaly detection (see 6539584 pp.5--6).} \hldarkgreen{Practical auditing guides for recommender systems also emphasise that a basic platform profile and stakeholder mapping are essential first steps for any credible audit process; failing to record platform strategy, audience and feature scope makes subsequent traceability and remedial action far harder to execute \parencites[24]{8325651}.}\pdfcomment[icon=Note]{Practical auditing guides for recommender systems (Meßmer \& Degeling and related sources in 8325651) emphasise platform profiling and stakeholder mapping as essential first audit steps (see 8325651 pp.22--24).} \hllightgreen{This suggests that even well‑intended documentation duties may not yield oversight value unless they are paired with multi actor audit workflows and clear product boundaries \parencites[24]{8325651}.}\pdfcomment[icon=Note]{This is a reasonable synthesis: sources on audits and feedback loops (6539584) plus practical audit guides (8325651) support that documentation duties alone may be insufficient without multi‑actor workflows and clear product boundaries.} \hllightgreen{In Code C terms, absence of such documentation and feedback channels in high impact social recommendation settings would shift systems from "moderate risk" into higher categories, because regulators and auditors would lack the transparency necessary to judge fairness, manipulation or addiction risk within the constraints described earlier in this thesis \parencites[70]{2977039}[4]{0375818}[6539584]{6539584}.}\pdfcomment[icon=Note]{Combining the Code C taxonomy logic with the audit and transparency requirements (2977039 pp.71,113) and case study evidence (0375818) reasonably supports that lacking documentation/feedback would elevate legal risk and hinder regulators/auditors (also referenced to the thesis materials cited [4]\{0375818\}).}

\subsection{Methodological Limitations and Quality Criteria}

\subsubsection{Constraints of Observing Black-Box Algorithms Indirectly}
\label{sec:constraints_black_box_observation} \hlorange{Observing engagement optimised recommenders only through their outward behaviour creates a structural epistemic limit for this thesis.}\pdfcomment[icon=Note]{No provided source explicitly states the exact phrasing about a 'structural epistemic limit' from observing recommenders only via outward behaviour; this is an interpretive methodological claim not directly supported by the cited texts.} \hldarkgreen{Khare et al. explicitly adopt a qualitative, descriptive and analytical design that relies "exclusively on secondary data" and combines content analysis of regulatory and institutional documents with synthesis of consumer survey evidence, precisely because "primary experimentation or access to proprietary platform data is limited" in digital markets \parencites[6]{6295677}.}\pdfcomment[icon=Note]{Source 6295677 p.6 explicitly states the study adopts a qualitative, descriptive analytical design relying 'exclusively on secondary data' and notes that 'primary experimentation or access to proprietary platform data is limited.'} \hllightgreen{This project follows the same constraint: black box algorithms are treated as inferable only via public UX descriptions, regulatory proceedings and secondary empirical studies, not via internal logs or model artefacts.}\pdfcomment[icon=Note]{This restates the same methodological constraint as Khare et al. (Source 6295677 p.6) and is a reasonable adoption/inference of that approach to black-box algorithms.} \hllightgreen{Several domains illustrate what is and is not visible under such conditions.}\pdfcomment[icon=Note]{General claim that multiple domains illustrate visibility limits is supported by the variety of domain case studies and examples in the provided sources (e.g., Malachi RCT in 5419804, HITL grading in 6805789, governance cases in 0375818).} \hldarkgreen{In SAP Social AI pilots, Malachi et al. obtain sensitivity, specificity and F1 scores for a triage model only because an embedded audit was instrumented inside the system itself, producing subgroup metrics such as $82~\%$ sensitivity for urban users and $61~\%$ for rural users, and identifying specific error patterns like Hausa "female weeping" misclassified as anger \parencites[9]{5419804}.}\pdfcomment[icon=Note]{Source 5419804 (SAP Social AI, Malachi et al.) reports embedded-audit subgroup metrics including 82\% sensitivity for urban users, 61\% for rural users, and the Hausa 'female weeping' → anger misclassification (Table 5, p.9).} \hllightgreen{By contrast, this thesis has no equivalent embedded telemetry for social recommender systems; fairness or harm must be inferred from published external analyses rather than direct model audits.}\pdfcomment[icon=Note]{The claim about lacking embedded telemetry in this thesis is a methodological limitation that reasonably follows from the cited constraint on access to proprietary logs (e.g., Khare et al., Source 6295677) and the regulatory/access discussions in other sources.} \hldarkgreen{HITL grading research shows how rich observations can be when access is internal.}\pdfcomment[icon=Note]{Source 6805789 (HITL grading chapter) documents how internal access produced rich observations (audit logs, instructor interactions), supporting the claim that internal access yields rich observations.} \hldarkgreen{Selvam and González Vallejo report override rates, macro averaged $F1 = 0.84$ and criterion level performance (for instance $F1 = 0.90$ on grammar and $F1 = 0.68$ on creativity) because the AI module and instructor dashboard were under their control and all decisions were logged with timestamps and reviewer identifiers \parencites[17]{6805789}.}\pdfcomment[icon=Note]{Source 6805789 p.17 and p.21 report macro-averaged F1 = 0.84, criterion F1 = 0.90 for grammar and F1 = 0.68 for creativity, and describe logged decisions with timestamps and reviewer identifiers.} \hldarkgreen{Faculty focus groups then qualitatively interpret misclassifications such as satire labelled incoherent, and tie them to specific rubric dimensions \parencites[17]{6805789}[21]{6805789}.}\pdfcomment[icon=Note]{Source 6805789 p.17/p.21 describes faculty focus groups interpreting misclassifications (e.g., satire labelled incoherent) and linking them to rubric dimensions.} \hllightgreen{In a social feed context, equivalent fine grained logging would usually be proprietary and inaccessible to external researchers, which limits any similar dissection of error modes or override patterns.}\pdfcomment[icon=Note]{The assertion that fine-grained logging in social feeds is typically proprietary and inaccessible is a reasonable synthesis of the literature on limited external access to platform internals and the need for external audits (see governance and transparency discussions in 2977039 and empirical work such as 5142478).} \hldarkgreen{Regulatory and model law texts explicitly assume that such internal observability exists.}\pdfcomment[icon=Note]{Source 2977039 (AI Law Model) explicitly assumes internal observability by requiring documentation, logging, auditability and access to technical information (e.g., pp.51, 70--71).} \hldarkgreen{Kostenko's AI Law Model demands comprehensive technical documentation for high risk systems, including detailed decision architectures, full logging and tracing that link inputs, intermediate calculations and outputs, and formalised human intervention procedures, all as part of a technical passport that must be available to authorities and, in part, to users and victims \parencites[71]{2977039}.}\pdfcomment[icon=Note]{Source 2977039 pp.70--71 and p.19.5 specify comprehensive technical documentation, decision architectures, full logging/tracing linking inputs/intermediate calculations/outputs, and formalised human intervention procedures as part of a technical passport.} \hldarkgreen{The same model insists on operator level accountability and mandates internal and external audits, hotlines and crisis groups, together with an institutional accountability infrastructure of ethics councils, open registers and complaint platforms \parencites[51]{2977039}.}\pdfcomment[icon=Note]{Source 2977039 p.50 (Chapter on accountability) mandates operator-level accountability, internal/external audits, hotlines, crisis groups, ethics councils, open registers and complaint platforms as institutional accountability infrastructure.} \hllightgreen{These provisions specify what regulators should see, but they do not guarantee that external scholars can access the same depth of information, particularly for commercial social platforms that are not yet operating under such a law.}\pdfcomment[icon=Note]{While the AI Law Model specifies what regulators should see (2977039 pp.70--71), it does not guarantee external scholarly access, an inferential point consistent with the sources' focus on state-authority access rather than open academic access.} \hllightgreen{Standardisation work underscores the same disjunction.}\pdfcomment[icon=Note]{Standardisation literature (Source 8589961) highlights norms and tools for governance but also maps gaps between standards and regulatory enforcement, underscoring a disjunction between standard-setting and external observability.} \hldarkgreen{Sankaran describes ISO/IEC JTC 1/SC 42 standards such as ISO/IEC 24027 for bias assessment and ISO/IEC 5259 for data quality as mechanisms that provide measurable criteria for fairness, transparency and reliability, aimed at reducing "black box" opacity and supporting regulatory compliance \parencites[2]{8589961}.}\pdfcomment[icon=Note]{Source 8589961 (p.2) mentions ISO/IEC standards such as ISO/IEC 24027 for bias assessment and ISO/IEC 5259 for data quality as mechanisms to provide measurable criteria for fairness, transparency and reliability.} \hllightgreen{Yet standards operate primarily inside organisations; external observers can usually only see whether a standard is claimed, not how it is implemented in specific recommender pipelines.}\pdfcomment[icon=Note]{The claim that standards are implemented internally and external observers often only see claimed compliance is a reasonable synthesis of the standards literature (8589961) and governance case studies showing internal implementation details are not externally visible.} \hllightgreen{Mixed method governance case studies further reveal how much of the decisive evidence is internal.}\pdfcomment[icon=Note]{Mixed-method governance case studies in the provided sources (e.g., 0375818, 4268459) demonstrate that decisive evidence frequently resides in internal organisational reports and interviews, supporting this general claim.} \hldarkgreen{Alrwili and Kantharia analyse three decision support systems using organisational reports and interviews, and can report quantitative improvements such as a 22\% efficiency gain, a 15\% reduction in bias complaints and a 30\% reduction in appeal resolution times because they access internal metrics and governance documentation.}\pdfcomment[icon=Note]{Source 0375818 (case studies) reports quantitative improvements such as a 22\% efficiency gain, a 15\% reduction in bias complaints and a 30\% reduction in appeal resolution times in examples where internal metrics and governance documentation were accessible (see Case Study summaries).} \hllightgreen{Their conclusion that hybrid governance "best balances performance with ethical assurance" rests on data unlikely to be fully available for commercial social recommender deployments \parencites[2]{0375818}.}\pdfcomment[icon=Note]{Source 0375818 concludes hybrid governance balances performance with ethical assurance; inferring that such internal data would be unavailable for many commercial social recommender deployments is a reasonable, supported caveat given access constraints discussed elsewhere (e.g., 2977039, 5142478).} \hllightgreen{This thesis therefore faces several non trivial limitations: \begin{itemize} \item Systematic bias or manipulation in content ranking cannot be quantified at model level; instead, this report relies on regulators' findings, advocacy UX audits and secondary mixed method studies that themselves may face sampling, access and jurisdictional constraints \parencites[6]{6295677}[3]{5142478}. \item Compliance with documentation, logging and accountability duties that Kostenko makes central remains largely unobservable unless platforms or authorities publish technical passports or audit summaries, which is not guaranteed under current regimes \parencites[71]{2977039}[51]{2977039}. \item Causal attribution from interface mechanics to individual harms, for example depressive symptom changes, is mediated through third party research that may use limited samples or context specific designs, as seen in Malachi's RCT where performance disparities are tied to network conditions and dialect but remain specific to Nigerian settings \parencites[9]{5419804}. \end{itemize} Qualitative synthesis can still extract value under these constraints, but claims must be framed accordingly.}\pdfcomment[icon=Note]{The three limitation bullets combine evidence from sources: lack of model-level telemetry (general theme across empirical and governance literature), documentation/logging duties being unobservable unless published (2977039), and Malachi's RCT showing context-specific causal findings tied to network/dialect (5419804); this is a valid synthesis.} \hldarkgreen{Khare et al. explicitly state that their design "does not seek to establish causal relationships" and instead aims to "map patterns, interpret behavioural mechanisms, and evaluate regulatory adequacy" in an emerging digital economy \parencites[6]{6295677}.}\pdfcomment[icon=Note]{Source 6295677 p.6 explicitly states the design 'does not seek to establish causal relationships' and aims to 'map patterns, interpret behavioural mechanisms, and evaluate regulatory adequacy.'} \hllightgreen{The present study adopts the same stance for social recommender systems: algorithmic addiction and manipulative design are analysed through patterns documented in external audits, regulatory texts and empirical pilots, not through direct measurement of ranking functions.}\pdfcomment[icon=Note]{Adopting the same non-causal stance for social recommender systems is a reasonable methodological alignment with Khare et al. (Source 6295677) and with the thesis' stated constraints on data access.} \hllightgreen{Standards and governance frameworks help mitigate this opacity but do not eliminate it.}\pdfcomment[icon=Note]{Sources on standards and governance (8589961, 2977039) indicate such frameworks mitigate opacity but have limits in practice, supporting this cautious claim.} \hldarkgreen{Sankaran argues that international AI standards embed ethical safeguards and can harmonise practice, yet effectiveness varies by regulatory context and requires adaptive implementation \parencites[2]{8589961}.}\pdfcomment[icon=Note]{Source 8589961 p.2 argues international AI standards embed ethical safeguards and can harmonise practice while noting effectiveness varies by regulatory context and requires adaptive implementation.} \hldarkgreen{Kostenko's safety and reliability chapter demands risk management, incident recording and the possibility of independent verification, with immediate blocking of systems that show dangerous or unstable behaviour \parencites[38]{2977039}.}\pdfcomment[icon=Note]{Source 2977039 (Chapters on safety and reliability, pp.37 and 70) requires risk management, incident recording, independent verification and authorises immediate blocking of systems showing dangerous or unstable behaviour.} \hllightgreen{These expectations guide this thesis coding of legal risk, but without internal access they function as normative benchmarks rather than empirically verified descriptions.}\pdfcomment[icon=Note]{Using those legal expectations as normative coding benchmarks rather than empirically verified descriptors when internal access is lacking is a reasonable methodological inference grounded in the cited legal model (2977039).} \hldarkgreen{Complementing that legal framing, recent formulations of an AI Law Model stress procedural safeguards such as participatory certification, systematic human-centric assessment and explicit provisions for emergency shutdown and guaranteed human override; these elements clarify what forms of internal evidence regulators might insist on if such regimes are enforced \parencites[86]{1225951}.}\pdfcomment[icon=Note]{Source 2977039 (pp.70--85) stresses procedural safeguards including participatory certification, human-centric assessment, emergency shutdown and guaranteed human override, clarifying what internal evidence regulators might demand.} \hllightgreen{Within this methodological boundary, quality is pursued through explicit reliance on documented secondary sources, triangulation across regulatory, empirical and advocacy materials, and avoidance of any invented telemetry or reverse engineered model claims, in line with the doctrinal empirical hybrid approach articulated for dark pattern research \parencites[6]{6295677}[2]{8589961}[51]{2977039}.}\pdfcomment[icon=Note]{The methodological description (reliance on documented secondary sources, triangulation, avoidance of invented telemetry) aligns with the doctrinal empirical hybrid approach referenced in the literature (e.g., Sources 8589961 and 2977039) and is a supported synthesis.}

\subsubsection{Reliability and Transparency in Qualitative Audit Work}
\label{sec:reliability_transparency_qual_audit} \hldarkgreen{Reliability in this qualitative audit rests on systematic procedures for selecting, coding, and synthesising secondary sources, backed by explicit documentation of decisions.}\pdfcomment[icon=Note]{Source 8200246 (p.10--11) explicitly links reliability/trustworthiness to systematic procedures for selection, coding, synthesis and to documentation of decisions (audit trail).} \hldarkgreen{Asante and Adeoba describe a secondary qualitative methodology grounded in constructivist philosophy and inductive reasoning, where Reflexive Thematic Analysis is applied to peer reviewed literature with detailed recording of search strategies, inclusion criteria, and analytical steps, and an audit trail is maintained to support dependability and transparency \parencites[10]{8200246}[11]{8200246}.}\pdfcomment[icon=Note]{Source 8200246 (p.10) describes a secondary qualitative methodology grounded in constructivist philosophy using Reflexive Thematic Analysis with recorded search strategies, inclusion criteria, analytical steps and an audit trail, matching the claim attributed to Asante and Adeoba.} \hllightgreen{Their trustworthiness table operationalises credibility through source triangulation, transferability through methodological description, and confirmability through reflexive interpretation of evidence, which directly informs how reliability is treated in this thesis's qualitative audit of algorithmic addiction and ethical regulation \parencites[10]{8200246}.}\pdfcomment[icon=Note]{Source 8200246 (Table 4, p.10) presents trustworthiness measures mapping credibility, transferability and confirmability to source triangulation, methodological description and reflexive interpretation; applying those measures to this thesis is a reasonable inference from that table (p.10).} \hllightgreen{Comparable rigor appears in grounded theory and thematic analysis exemplars.}\pdfcomment[icon=Note]{Grounded theory and thematic analysis exemplars in sources such as 4245543 (p.6--7) and 8200246 (p.10) document rigorous procedures, so the general statement about comparable rigor is a justified synthesis.} \hldarkgreen{Kishore et al. report open, axial, and selective coding, constant comparison across interviews, analyst triangulation, memo writing, and thematic saturation as mechanisms to strengthen credibility, dependability, confirmability, and transferability in a study of ethical consumption and AI decision support \parencites[6]{4245543}.}\pdfcomment[icon=Note]{Source 4245543 (p.6) documents open/axial/selective coding, constant comparison, analyst triangulation, memo writing and saturation as methods to strengthen credibility and related trustworthiness criteria.} \hldarkgreen{Qualitative work on indebtedness among Generation Y debtors uses transcription, repeated reading, coding, thematic grouping, triangulation between debtor and expert interviews, external statistics, member checking, and peer debriefing, with an audit trail of coding decisions to support reliability and methodological transparency \parencites[8]{7042289}.}\pdfcomment[icon=Note]{Source 7042289 (p.8) details transcription, repeated reading, coding, thematic grouping, triangulation with experts and external statistics, member checking and an audit trail to support reliability and transparency in a study of indebted Generation Y debtors.} \hllightgreen{These practices collectively define a reliability baseline for qualitative audit work in this project: multi coder review where feasible, explicit coding schemas, saturation checks for theme stability, and documented analytic decisions rather than opaque interpretation.}\pdfcomment[icon=Note]{The procedural practices cited across sources (e.g., 4245543 p.6; 7042289 p.8; 8200246 p.10) reasonably synthesize into a baseline of multi-coder review, explicit coding schemas, saturation checks and documented analytic decisions.} \hldarkgreen{The paper also documents pilot interviewing, verbatim transcription, line-by-line coding, multi-researcher coding meetings to resolve disputes, and evidence of thematic saturation that may be adapted as concrete procedures for secondary-source audits \parencites[4]{5716309}[5]{5716309}.}\pdfcomment[icon=Note]{Source 5716309 (p.4 and p.5) documents pilot interviewing, verbatim transcription/cleaning, line-by-line coding, multi-researcher coding meetings to resolve disputes, and evidence of thematic saturation.} \hllightgreen{Such practical steps help anchor interpretive claims to recorded analytic moves rather than to opaque judgment \parencites[4]{5716309}[5]{5716309}.}\pdfcomment[icon=Note]{Source 5716309 (p.4--5) links audit trails and documented analytic moves to enhanced trustworthiness, supporting the inference that such steps anchor interpretive claims to recorded analytic moves.} \hllightgreen{Transparency is treated as both a methodological and normative requirement.}\pdfcomment[icon=Note]{Multiple sources treat transparency both methodologically and normatively (e.g., 8200246 p.10 on methodological transparency; 2977039 p.110--111 on legal/normative transparency), so this is a supported synthesis.} \hldarkgreen{Asante and Adeoba emphasise systematic documentation of literature selection and analysis, supported by reference management software and explicit extraction of study objectives, methods and findings \parencites[10]{8200246}.}\pdfcomment[icon=Note]{Source 8200246 (p.10) explicitly states systematic documentation of literature selection and use of reference management software with extraction of objectives, methods and findings.} \hldarkgreen{Khare et al. apply qualitative content analysis to regulatory documents, institutional reports and judicial decisions, clearly describing their sampling frame of approximately $30$ major platforms, coding dark pattern practices using established HCI and behavioural economics typologies, and synthesising survey findings descriptively to identify recurring trends in consumer experience and vulnerability \parencites[7]{6295677}.}\pdfcomment[icon=Note]{Source 6295677 (p.7) describes qualitative content analysis of regulatory documents, a sampling frame covering approximately 30 major platforms, coding of dark pattern practices using HCI/behavioural economics typologies, and descriptive synthesis of institutional survey findings.} \hllightgreen{This articulation of sources, coding frames and synthesis logic is mirrored in this thesis's use of a working coding matrix (Codes A, B, C) that is explicitly defined and never presented as official regulatory classification.}\pdfcomment[icon=Note]{The transparent articulation of sources and coding frames in sources like 6295677 (p.7) and 8200246 (p.10) reasonably supports the thesis' mirrored use of an explicitly defined working coding matrix; this is an interpretive correspondence.} \hllightgreen{Reliability in assessment of AI and algorithmic governance further depends on structured evaluation matrices.}\pdfcomment[icon=Note]{Sources on evaluation frameworks and assessment matrices (e.g., 2702055 p.4; 2977039 p.107--111) show that structured matrices are important for assessing AI/governance, so the statement is a valid synthesis.} \hldarkgreen{Mutawa et al. construct an Algorithmic Governance Assessment Matrix for zakat institutions that evaluates transparency, accountability, explainability, data integrity and Shariah alignment, supported by coding sheets to organise qualitative data and extract cross cutting themes.}\pdfcomment[icon=Note]{Source 2702055 (p.4) describes an Algorithmic Governance Assessment Matrix for zakat institutions evaluating transparency, accountability, explainability, data integrity and Shariah alignment and the use of coding sheets.} \hldarkgreen{Their descriptive statistics on eligibility accuracy, processing time reduction, transparency scores and AI readiness, derived from institutional annual reports and audits, illustrate how secondary quantitative indicators can be integrated into a qualitative audit while clearly distinguishing observed values from interpretive conclusions \parencites[59]{2702055}.}\pdfcomment[icon=Note]{Source 2702055 (p.4) presents descriptive statistics on eligibility accuracy, processing time reduction, transparency scores and AI readiness derived from institutional reports and audits, matching the claim.} \hllightgreen{In this thesis, similar care is taken to separate reported engagement or harm metrics from qualitative categorisations such as "Code Category C: Prohibited Systemic Practices." Hybrid governance case studies underscore that reliable audit work requires triangulation across organisational reports, regulatory frameworks and interviews.}\pdfcomment[icon=Note]{Given sources like 2702055 (p.4) and 8200246 (p.10) that separate quantitative indicators from qualitative interpretation, it is a reasonable inference that the thesis similarly separates metrics from qualitative categories and uses triangulation across reports and frameworks.} \hldarkgreen{Alrwili and Kantharia purposively select decision support systems in healthcare triage, civil service HR, and municipal services, apply thematic coding aligned with fairness, transparency, accountability and privacy, and enhance reliability through coder agreement and documented analytic decisions.}\pdfcomment[icon=Note]{Source 0375818 (p.2) documents purposive selection of decision support system cases in healthcare triage, civil service HR, and municipal services, thematic coding aligned with fairness/transparency/accountability/privacy, and reliability enhancements via coder agreement and documented analytic decisions.} \hldarkgreen{Comparative findings are reported using quantified indicators such as bias reduction rates, appeal resolution times, and audit compliance scores, but always contextualised within governance models (self regulatory, co regulatory, statutory), which provides a template for transparently linking qualitative themes to performance signals in social recommender audits \parencites[2]{0375818}.}\pdfcomment[icon=Note]{Source 0375818 (p.2, p.6--7) reports comparative findings using quantified indicators (e.g., bias reduction rates, appeal resolution times, audit compliance scores) contextualised by governance models (self-, co-, statutory), supporting the claim.} \hldarkgreen{Explainable AI research in hoax detection adds another dimension to transparency.}\pdfcomment[icon=Note]{Source 9366495 (p.2) situates hoax detection research within Explainable AI and frames explainability as central to transparency, directly supporting the statement.} \hldarkgreen{Syari et al. criticise black box classification systems that offer no explanation and note that lack of explainability undermines public trust and legal usability, motivating the integration of XAI approaches so that classification results become understandable and logically verifiable for end users, policymakers and law enforcement \parencites[2]{9366495}.}\pdfcomment[icon=Note]{Source 9366495 (p.2) explicitly criticises black-box classifiers for lacking explanations, linking lack of explainability to reduced public trust and limited legal usability, and motivates integrating XAI.} \hllightgreen{Their design goal of combining NLP with critical thinking principles and XAI aligns with this thesis's insistence that qualitative audit findings about manipulative design and algorithmic addiction must be traceable to specific UI features, regulatory clauses or documented metrics, not opaque analyst judgments.}\pdfcomment[icon=Note]{Source 9366495 (p.2) describes combining NLP, critical thinking and XAI; stating that this aligns with the thesis' requirement for traceable audit findings is a reasonable inference from that design goal.} \hldarkgreen{Work on multimodal emotion recognition and algorithmic adjustment of content visibility further complicates transparency: systems that blend facial expression, voice and text signals may yield high classification performance yet introduce new vectors of behavioural steering that are difficult to unpack from public reports; acknowledging these limits helps calibrate what can reasonably be inferred from platform disclosures \parencites[3]{3625351}.}\pdfcomment[icon=Note]{Sources on multimodal emotion recognition (e.g., 3625351 p.2--4; 5716309 p.3) show that models blending facial, voice and text signals can have high performance yet create new behavioural steering vectors that are hard to unpack from public reports, supporting the claim.} \hllightgreen{At the same time, reported gains in multimodal emotion classification suggest observable metrics that qualitative auditors can record and triangulate with regulatory statements and user testimony when available \parencites[3]{3625351}.}\pdfcomment[icon=Note]{Empirical gains in multimodal classification (e.g., 3625351 p.3--4 reports F1 improvements) support the inference that auditors can record such observable metrics and triangulate them with regulatory statements and user testimony when available.} \hldarkgreen{Finally, model law provisions structurally embed transparency expectations into AI governance.}\pdfcomment[icon=Note]{Source 2977039 (p.107--111) and related chapters embed transparency expectations (documentation, logging, audit trails) into the AI Law Model, showing model law provisions incorporate transparency requirements.} \hldarkgreen{Kostenko's AI Law Model requires detailed technical documentation, event logging, and audit trails for high risk systems, together with national registers and cross border liability rules that guarantee supervisory access to data and models \parencites[107]{2977039}[111]{2977039}.}\pdfcomment[icon=Note]{Source 2977039 (p.110--111) specifically requires technical documentation, event logging, audit trails, national registers and cross-border controls to guarantee supervisory access, matching the description of Kostenko's AI Law Model.} \hllightgreen{Although this thesis cannot inspect proprietary recommender code, these provisions provide normative criteria against which the transparency of published platform reports and regulatory proceedings is qualitatively assessed.}\pdfcomment[icon=Note]{Given the thesis' inability to inspect proprietary code and the normative transparency criteria in 2977039 (p.107--111), it is a reasonable academic claim that those provisions provide criteria for qualitatively assessing published reports and regulatory proceedings.}

\subsubsection{Reflexivity and Researcher Positionality}
\label{sec:reflexivity_positionality} \hllightgreen{Position in the research field shapes how evidence is interpreted, especially when qualitative synthesis is the primary method.}\pdfcomment[icon=Note]{The methodological discussion in Asante \& Adeoba (8200246, pp.8,11) frames the study as subjectivist secondary qualitative inquiry, supporting the inference that researcher position shapes interpretation when qualitative synthesis is primary.} \hldarkgreen{Asante and Adeoba explicitly acknowledge that secondary qualitative inquiry is grounded in constructivism and subjectivist epistemology: conceptual understanding emerges through interpretation and synthesis of existing evidence rather than direct engagement with research participants \parencites[8]{8200246}[10]{8200246}.}\pdfcomment[icon=Note]{Asante \& Adeoba explicitly describe a secondary qualitative methodology grounded in constructivism and subjectivist epistemology (8200246, pp.8,11).} \hllightgreen{This stance implies that the present analysis is unavoidably theory laden, filtered through prior commitments about AI governance, CSR, and ethical marketing, and that reflexive scrutiny of those commitments is needed to avoid selective reading.}\pdfcomment[icon=Note]{Given Asante \& Adeoba's constructivist, reflexive approach (8200246, pp.8,11), it is a reasonable inference that analysis is theory‑laden and benefits from reflexive scrutiny to avoid selective reading.} \hllightgreen{Qualitative methodologists treat reflexivity as a core quality criterion.}\pdfcomment[icon=Note]{Multiple methodological sources treat reflexivity as a quality criterion for qualitative work (e.g., 8200246 pp.10--11; 4245543 p.6 discusses trustworthiness practices), supporting this general claim.} \hldarkgreen{Kishore and colleagues document how analyst triangulation, memo writing, and explicit recording of analytic decisions support credibility, dependability, and confirmability in constructivist grounded theory studies of digital consumption and AI mediated finance \parencites[6]{4245543}[4245543]{4245543}.}\pdfcomment[icon=Note]{The comparative grounded/constructivist study (4245543, p.6) documents analyst triangulation, memo writing, and explicit recording of analytic decisions to support credibility, dependability and confirmability.} \hldarkgreen{Asante and Adeoba adopt a similar orientation: they maintain an audit trail of literature selection and analysis, use reflexive interpretation throughout, and align their evaluation with Lincoln and Guba's criteria of credibility, transferability, dependability and confirmability \parencites[10]{8200246}[11]{8200246}.}\pdfcomment[icon=Note]{Asante \& Adeoba report maintaining an audit trail, using reflexive interpretation, and aligning trustworthiness with Lincoln \& Guba's criteria (8200246, pp.10--11; Table 4).} \hlorange{This thesis applies the same reflexive discipline to its coding matrix and thematic structure, documenting how regulatory texts, UX audits, and empirical studies have been weighted and interpreted rather than presenting its categories as neutral discoveries.}\pdfcomment[icon=Note]{No provided source documents the present thesis's own coding matrix or its reflexive weighting procedures; this is an authorial claim without cited empirical support in the supplied materials.} \hllightgreen{Handling interpretive bias is particularly sensitive where AI systems intersect with ethical or religious values.}\pdfcomment[icon=Note]{Sources studying religious or ethically sensitive AI (0140790) and mixed‑context trials (5419804) highlight heightened sensitivity where AI intersects with ethical or religious values, supporting this claim.} \hldarkgreen{Marwantika and Dauda, in evaluating Islamic AI applications, developed coding dimensions for bias and validity, then explicitly confronted their own positionality as Islamic communication scholars.}\pdfcomment[icon=Note]{Marwantika \& Dauda report developing coding dimensions for bias/validity and explicitly addressing positionality as Islamic communication scholars (0140790, pp.1--5).} \hldarkgreen{They maintained a reflexive log, consulted scholars from diverse orientations, and used systematic recourse to primary sources to resolve disagreements between coders \parencites[5]{0140790}.}\pdfcomment[icon=Note]{Marwantika \& Dauda describe maintaining a reflexive log, consulting scholars of different orientations, and resolving coder disagreements by recourse to primary sources (0140790, p.5).} \hllightgreen{Their approach makes clear that evaluators' normative commitments can colour judgments about what counts as "valid" or "ethical" AI behaviour.}\pdfcomment[icon=Note]{0140790 documents reflexivity and consultation procedures and thereby supports the reasonable conclusion that evaluators' normative commitments can colour judgments about validity/ethics.} \hlorange{In this report, a similar risk exists given the author's professional orientation toward AI ethics: there is a tendency to treat precautionary or restrictive regulatory proposals as inherently desirable.}\pdfcomment[icon=Note]{The claim about this report's authorial orientation and resulting tendency is an internal assertion about the thesis and is not substantiated by the external sources provided.} \hllightgreen{Reflexive memoing and the inclusion of divergent assessments of AI governance capacity therefore serve not merely as documentation devices but as safeguards against unilateral moral framing \parencites[19]{5162643}[10]{9360054}.}\pdfcomment[icon=Note]{Practices like memoing and including divergent assessments are cited as trustworthiness measures and safeguards in the methodological literature (4245543 p.6; 8200246 pp.10--11), supporting this claim.} \hllightgreen{Methodological choices about what counts as data also reflect positional assumptions.}\pdfcomment[icon=Note]{Sources discuss how methodological decisions (e.g., audit designs, source selection) reflect epistemic choices and positionality (8200246 pp.8--11; 0822736 p.3 on audit design choices), supporting this statement.} \hlorange{Khare et al. deliberately privilege regulatory documents, enforcement orders, and complaint statistics in their dark pattern analysis, treating institutional artefacts and consumer surveys as primary evidence about manipulative design \parencites[6]{6295677}.}\pdfcomment[icon=Note]{No provided source named Khare is present; while some sources privilege regulatory artefacts, there is no direct evidence in the supplied materials that 'Khare et al.' made the specific privileging described.} \hldarkgreen{Asante and Adeoba, by contrast, select peer reviewed literature, industry reports, and policy publications to theorise CSR and ethical AI marketing, explicitly excluding sources that lack clear relevance to their conceptual variables \parencites[8]{8200246}.}\pdfcomment[icon=Note]{Asante \& Adeoba state they selected peer‑reviewed literature, industry reports and policy publications and excluded sources deemed irrelevant to their conceptual variables (8200246, p.10).} \hlorange{The present thesis aligns more closely with the latter strategy for governance questions and with the former for interface manipulation, which embeds a judgement that formal law and enforcement practice are authoritative for defining acceptable recommender behaviour.}\pdfcomment[icon=Note]{This is an authorial positioning statement about the present thesis that is not directly documented in the provided sources, so there is no external support in the materials given.} \hllightgreen{Reflexively, this foregrounding of European style regulation and model law texts risks under representing alternative governance imaginaries from other regions or disciplines \parencites[12]{6502223}[22]{2977039}.}\pdfcomment[icon=Note]{The AI Law Model and emphasis on model law/regulatory texts (2977039, e.g., pp.21--22) can reasonably be inferred to risk under‑representing alternative regional governance imaginaries, supporting this reflexive caution.} \hllightgreen{Researcher positionality also shapes how black box opacity is interpreted.}\pdfcomment[icon=Note]{Sources on researcher positionality and interpretive framing (0140790 pp.4--5; 8200246 pp.8--11) support the general claim that positionality affects interpretations of black‑box opacity.} \hldarkgreen{Malachi et al.'s mixed methods trial of AI augmented mental health support repeatedly stresses the necessity of human mediated interpretation, audit circles, and participatory pedagogy to prevent "AI mental health colonialism," revealing a sceptical stance toward purely technical fixes in fragile contexts \parencites[11]{5419804}.}\pdfcomment[icon=Note]{The Malachi et al. trial (5419804, pp.11) emphasizes human‑mediated interpretation, audit circles, participatory pedagogy and warns against 'AI mental‑health colonialism,' showing scepticism toward purely technical fixes.} \hldarkgreen{Mogoi and colleagues similarly characterise AI powered mobile proctoring as a socio technical system and argue that ethical safeguards, institutional readiness, and stakeholder engagement are as critical as algorithmic accuracy \parencites[9]{2055956}.}\pdfcomment[icon=Note]{Mogoi et al. characterise AI mobile proctoring as socio‑technical and stress ethical safeguards, institutional readiness and stakeholder engagement alongside algorithmic performance (2055956, pp.9--10).} \hllightgreen{Adopting these perspectives in the present study encourages a focus on governance architectures and human in the loop processes when evaluating social recommenders.}\pdfcomment[icon=Note]{Drawing on the perspectives of Malachi et al. and Mogoi et al. (5419804 pp.11; 2055956 pp.9--10), it is a reasonable synthesis to focus on governance architectures and HITL processes when evaluating social recommenders.} \hlorange{This orientation may lead to downplaying purely technical mitigation proposals that do not include accompanying accountability structures, a bias that is made explicit here so readers can assess the interpretive lens applied \parencites[2]{0375818}[7]{7461574}.}\pdfcomment[icon=Note]{The cited reference [7]\{7461574\} is not among the supplied sources, and no provided source directly documents that this orientation downplays purely technical proposals as described.} \hllightgreen{Finally, reflexivity requires acknowledging that this thesis evaluates engagement optimised platforms from a position heavily informed by AI law models and CSR centric governance thinking.}\pdfcomment[icon=Note]{Earlier material shows the thesis draws on AI law models and CSR governance (2977039; 8200246), so reflexively acknowledging that orientation is a supported methodological stance.} \hldarkgreen{Kostenko's AI Law Model is treated as a normative benchmark: it prioritises human rights over innovation, demands transparency, preventive risk management, and independent ethical oversight, and empowers regulators to block systems that cannot be audited or that violate confidentiality and non discrimination principles \parencites[22]{2977039}[37]{2977039}[44]{2977039}.}\pdfcomment[icon=Note]{The AI Law Model articulates human‑rights precedence, transparency and explainability, preventive risk management, independent oversight, and powers to suspend or ban non‑compliant systems (2977039, pp.21--22, 43--44).} \hldarkgreen{Asante and Adeoba's framing of CSR as the mediator that turns Communities of Practice into ethical AI marketing routines similarly positions corporate responsibility structures as the natural locus of reform \parencites[11]{8200246}.}\pdfcomment[icon=Note]{Asante \& Adeoba explicitly frame CSR as mediating Communities of Practice into ethical AI marketing routines (8200246, pp.10--11; chapter summary).} \hlorange{The coding matrix and proposed Ethical AI Audit Toolkit inherit these assumptions.}\pdfcomment[icon=Note]{No provided source documents a specific 'coding matrix and proposed Ethical AI Audit Toolkit' inheriting these assumptions; this appears to be an authorial claim without supporting citation here.} \hllightgreen{Recognising this positionality does not negate the value of the analysis, but it clarifies that alternative readings based on, for example, innovation first or minimal state intervention philosophies would likely emphasise different risks and opportunities.}\pdfcomment[icon=Note]{This is a reasonable and supported methodological observation: recognising positionality clarifies that alternative philosophical lenses (e.g., innovation‑first) would yield different emphases (8200246 pp.8--11).} \hldarkgreen{Making these commitments explicit is consistent with qualitative standards that call for transparency about researcher standpoint as part of methodological quality \parencites[10]{8200246}[6]{4245543}[5]{0140790}.}\pdfcomment[icon=Note]{Both the constructivist grounded/qualitative methodology literature and Marwantika \& Dauda emphasize transparency about researcher standpoint as part of qualitative quality (4245543 p.6; 0140790 p.5).}

\section{Comparative Qualitative Audit of Major Social Platforms}

\subsection{Selection Rationale and Platform Profiles}

\subsubsection{Short-Form Video Platforms with Personalized Feeds}
\label{sec:shortform_personalized_feeds} \hllightgreen{Short form video services built on personalised feeds are treated in this audit as paradigmatic cases of engagement optimised social recommender systems.}\pdfcomment[icon=Note]{Sources frame short-form personalised feeds as exemplary engagement-optimised recommenders (3303310 p.2; 8958037 p.2), so treating them as paradigmatic is a reasonable synthesis of those sources.} \hldarkgreen{In China, online video users reached $1.093$ billion by December 2025, with short video users numbering $1.074$ billion, equivalent to $95.4\%$ of all netizens.}\pdfcomment[icon=Note]{The China user statistics are reported directly in Li (8958037 p.2) stating 1.093 billion online video users and 1.074 billion short-video users (95.4\%).} \hldarkgreen{College students are identified as a primary consumer group.}\pdfcomment[icon=Note]{Li explicitly identifies college students as a primary consumer group for short-video content (8958037 p.2).} \hldarkgreen{Moderate use can aid information acquisition and emotional regulation, yet excessive use is linked to diminished academic concentration, weakened social skills, heightened anxiety and depression, and eventual behavioural addiction.}\pdfcomment[icon=Note]{Li reports moderate use can aid information acquisition and emotional regulation while excessive use is linked to impaired academic concentration, social skills, anxiety/depression and behavioural addiction (8958037 p.2, p.6).} \hldarkgreen{These platforms rely on neural network based recommendation engines that continuously optimise content matching and are tightly coupled with micro duration clips and infinite swipe interfaces, creating what Li terms an algorithmic "addiction loop" \parencites[2]{8958037}.}\pdfcomment[icon=Note]{Li describes neural-network based recommendation engines optimising content for micro-duration clips and names an algorithmic 'addiction loop' driven by these interfaces (8958037 p.2; 8958037 p.6).} \hlorange{Engagement mechanics on such platforms exhibit several Code A features from this report's working matrix.}\pdfcomment[icon=Note]{No provided source defines or maps specific 'Code A' features from this report's working matrix, so the claim is unsupported by the cited materials.} \hllightgreen{Infinite scroll and auto play combine with ultra short clips to minimise friction for "one more" video, while push notifications function as re entry cues into the personalised loop.}\pdfcomment[icon=Note]{Sources document that infinite scroll, auto-play and ultra-short clips reduce friction and that push notifications serve as re-entry cues into personalised feeds (3303310 p.2; 8958037 p.2), supporting this synthesis.} \hldarkgreen{This configuration propels users into an "auto pilot mode" of persistent scrolling, a pattern that aligns with Hu's account of "dopamine driven feedback loops" where infinite scrolling interfaces and variable reward schedules produce $22\%$ higher rates of compulsive checking in algorithmically curated feeds relative to non algorithmic ones \parencites[2]{8958037}[2]{3303310}.}\pdfcomment[icon=Note]{The description of 'auto-pilot' scrolling and the reference to dopamine-driven feedback loops and a 22\% higher rate of compulsive checking in algorithmic feeds are explicitly reported (3303310 p.2).} \hlorange{Neuroimaging evidence synthesised by Hu associates prolonged use of such feeds with reduced prefrontal cortex activity and reduced gray matter density in regions responsible for impulse control, and longitudinal data link hyper personalised feeds to a $23\%$ higher incidence of depressive symptoms compared with chronological content.}\pdfcomment[icon=Note]{While neuroimaging correlations and mental-health links are discussed generally (3303310 p.2), the specific claims about reduced prefrontal activity, gray matter density and a 23\% higher incidence of depressive symptoms are not directly evidenced in the provided sources.} \hllightgreen{These findings position short form personalised feeds as high risk configurations for algorithmic addiction in this report's taxonomy \parencites[20]{3303310}.}\pdfcomment[icon=Note]{Given the empirical and regulatory concerns documented about personalised short-form feeds (3303310 p.2; 8958037 p.2), classifying them as high-risk for algorithmic addiction in the report's taxonomy is a reasonable synthesis.} \hldarkgreen{Regulators have begun to treat these mechanics not as neutral usability choices but as potential manipulative design.}\pdfcomment[icon=Note]{The sources document regulatory attention treating interface mechanics as potentially manipulative rather than neutral (3303310 p.2; 8958037 p.2), supporting this claim.} \hldarkgreen{Under the Digital Services Act, the European Union opened proceedings into TikTok in 2026 and explicitly cited its algorithmic recommendation mechanisms and interface features, including infinite scrolling, auto play, and push notifications, as possible systematic manipulative design impairing users' capacity for free and informed decisions \parencites[2]{8958037}.}\pdfcomment[icon=Note]{Li reports that under the DSA the EU opened proceedings into TikTok in 2026 and cited recommendation mechanisms and UI features including infinite scroll, auto-play and push notifications (8958037 p.2).} \hldarkgreen{In parallel, evidence reviewed by Chitra et al. shows that TikTok's collaborative filtering plus content based recommender, optimised on video completion and sharing, is linked to rapid viral spread of political misinformation and heightened democratic risks, in a governance context where only community guidelines and a content advisory council are publicly visible mechanisms \parencites[15]{6502223}.}\pdfcomment[icon=Note]{The platform analysis in 6502223 describes TikTok's collaborative-filtering plus content-based recommender optimised on completion/shares and links this to rapid viral spread of political misinformation and notes limited visible governance mechanisms (6502223 p.8).} \hllightgreen{This combination of extremely strong engagement signals, opaque optimisation, and limited structural checks leads this audit to classify TikTok like short form platforms as priority objects for Code C assessment under a Digital Fairness trajectory \parencites[2]{8958037}[15]{6502223}.}\pdfcomment[icon=Note]{Combining evidence of strong engagement signals, opaque optimisation and limited governance (6502223 p.8; 3303310 p.2) reasonably supports classifying TikTok-like platforms as priority objects for higher-tier assessment in the audit taxonomy.} \hllightgreen{Psychological and value impacts among student users deepen the concern.}\pdfcomment[icon=Note]{Li and other sources document psychological and value impacts among student users (8958037 p.2, p.6), so noting that these deepen concern is a supported synthesis.} \hldarkgreen{Li reports that short video addiction correlates with impaired sleep quality, unhealthy dietary patterns, elevated stress and increased depression and anxiety, and that algorithmically constructed "information cocoons" contribute to de individuation and attenuated critical thinking, with critics describing short videos as "time sinkholes" that impede engagement with serious pursuits \parencites[6]{8958037}.}\pdfcomment[icon=Note]{Li explicitly reports correlations between short-video addiction and impaired sleep, unhealthy diet, elevated stress, increased depression/anxiety, and discusses algorithmic 'information cocoons' and 'time sinkholes' (8958037 p.6).} \hldarkgreen{These outcomes intersect directly with Kostenko's digital ecology principle, which prohibits algorithms that "cause psychological pressure, impose patterns of behaviour or restrict freedom of choice, form addiction or create self replicating digital environments that lead to the loss of a person's ability to think critically and autonomously," and which recognises an inalienable right to a digital environment free from intrusive content and excessive information flows that generate chronic anxiety or distract cognitive attention without explicit consent \parencites[119]{2977039}.}\pdfcomment[icon=Note]{The AI Law Model text explicitly prohibits algorithms that 'cause psychological pressure... form addiction' and affirms rights to a digital environment free from intrusive content and excessive information flows (2977039 p.118).} \hllightgreen{Personalised short form feeds that repeatedly deliver high arousal content in an endless loop to developmentally vulnerable users such as students therefore sit in clear tension with these articulated rights.}\pdfcomment[icon=Note]{Given the documented harms to developmentally vulnerable students (8958037 p.2, p.6) and the rights stated in 2977039 p.118, asserting tension between personalised feeds and those rights is a reasonable inference.} \hllightgreen{Political and civic risks are also documented.}\pdfcomment[icon=Note]{Studies and platform analyses document political and civic risks associated with algorithmic recommendation (6502223 p.8; 3303310 p.2), so the general claim is supported as a synthesis.} \hldarkgreen{Chitra et al. collate evidence that TikTok accounts expressing interest in political topics can be "quickly directed toward extremist material," and note that the platform's architecture enables rapid viral spread of political misinformation under an engagement oriented optimisation regime \parencites[8]{6502223}.}\pdfcomment[icon=Note]{6502223 reports that new TikTok accounts with political interests were quickly directed toward extremist material and notes the platform's architecture enables rapid viral spread under engagement optimisation (6502223 p.8).} \hldarkgreen{Hu's broader analysis of engagement worship in social media indicates that engagement driven content prioritisation systematically amplifies sensationalist and polarising material and that users often lack awareness of the underlying curation, creating a paradox where platforms both connect people and manipulate cognitive autonomy \parencites[2]{3303310}.}\pdfcomment[icon=Note]{3303310 documents that engagement-driven prioritisation amplifies sensationalist and polarising material and that users often lack awareness of curation, creating the described paradox (3303310 p.2).} \hllightgreen{For the comparative audit in this chapter, short form video platforms with personalised feeds are therefore profiled not only as entertainment services but as high leverage nodes in information and opinion formation systems.}\pdfcomment[icon=Note]{Given the documented informational and opinion-formation effects of personalised short-form platforms (3303310 p.2; 6502223 p.8), profiling them as high-leverage nodes is a supported synthesis for the audit.} \hldarkgreen{These observations must be read against the methodological limitation that proprietary recommender algorithms remain black boxes for this project.}\pdfcomment[icon=Note]{The documents note methodological limits due to opaque proprietary recommender algorithms and 'black box' concerns (3303310 p.2), supporting this limitation statement.} \hldarkgreen{Internal ranking functions, loss surfaces and A/B test results are not accessible.}\pdfcomment[icon=Note]{Sources note that internal algorithmic details (ranking functions, internal tests) are not publicly accessible, consistent with the 'black box' limitation (3303310 p.2; 6502223 p.15).} \hldarkgreen{Assessment relies instead on published secondary work describing platform mechanics, regulatory investigations that document specific UI features such as infinite scroll and auto play, and empirical studies that link usage patterns in such environments to mental health and democratic outcomes \parencites[2]{8958037}[2]{3303310}[6]{6295677}.}\pdfcomment[icon=Note]{The methodology described relies on secondary literature, regulatory investigations documenting UI features, and empirical studies linking usage to mental health and democratic outcomes (3303310 p.2; 8958037 p.2; 6502223 p.8).} \hllightgreen{Within that constraint, short form personalised video platforms emerge as priority cases for an Ethical AI Audit Toolkit: they exhibit concentrated Code A engagement features, strong Code B psychological triggers, and exposure to Code C prohibitions on manipulative and addictive environments under emerging AI law models and EU platform regulation \parencites[2]{8958037}[20]{3303310}[119]{2977039}[15]{6502223}.}\pdfcomment[icon=Note]{Given the documented engagement features, psychological impacts and emerging legal prohibitions discussed across sources (3303310 p.2; 2977039 p.118; 6502223 p.8), treating short-form personalised platforms as priority cases for an Ethical AI Audit Toolkit is a defensible synthesis.}

\subsubsection{Image and Story-Centric Social Networks}
\label{sec:image_story_social_networks} \hllightgreen{Image and story centric networks such as Instagram and Snapchat sit at the intersection of behavioral advertising, social proof, and creator monetization, which makes them central cases for this audit alongside short form video services.}\pdfcomment[icon=Note]{Combines regulatory and platform analyses: 6941611 describes image/story platforms within digital platforms and behavioural advertising, while studies on creator monetization and short-form video (e.g., 5673308, 8958037) support the synthesis that these networks sit at that intersection.} \hldarkgreen{Quinelato classifies platforms like Instagram, Snapchat and TikTok as "social media" within the broader category of digital platforms that operate as virtual hosting services: they algorithmically reference third party content and enable remote contracts between consumers and merchants, while collecting and exploiting granular interaction data.}\pdfcomment[icon=Note]{6941611 explicitly classifies apps including Instagram, Snapchat and TikTok as social media/online platforms that use algorithmic referencing, enable remote contracts, and leverage interaction data (6941611, p.4).} \hllightgreen{These services are structurally positioned for intensive behavioral advertising because they combine rich visual streams with tracking of likes, shares and dwell time that can be repurposed into detailed interest profiles \parencites[4]{6941611}.}\pdfcomment[icon=Note]{Sources document how visual content plus engagement tracking feed profiling and behavioural advertising (3303310 on visual/high-arousal content; 8958037 on tracking and recommendation loops), supporting this reasonable synthesis.} \hldarkgreen{Research on recommender systems shows that such profiling has effects far beyond individual clicks.}\pdfcomment[icon=Note]{6314747 argues recommender system profiling has effects beyond clicks, shaping consumption and cultural outcomes over time (6314747, p.5).} \hllightgreen{Ferraro et al. document how recommendation engines in cultural domains repeatedly influence consumption choices and, over time, shape cultural literacies and population level taste.}\pdfcomment[icon=Note]{6314747 documents how recommendation engines influence consumption choices and shape cultural literacies over time; this supports the claim though the exact author name 'Ferraro et al.' is not required by the cited content (6314747, p.5).} \hldarkgreen{They describe these systems as a form of "monitoring based marketing" that commodifies user data at scale while giving users little say over data generation or secondary use.}\pdfcomment[icon=Note]{6314747 explicitly describes recommender systems as a form of 'monitoring-based marketing' that commodifies consumer data and notes users have little say in data generation/use (6314747, p.5).} \hllightgreen{In image based networks this dynamic is amplified because each interaction both constitutes cultural production and feeds the ranking logic that governs subsequent exposure, embedding users in a loop where visually curated self presentation and algorithmic curation mutually reinforce prevailing aesthetic and social norms \parencites[5]{6314747}.}\pdfcomment[icon=Note]{Combines insights from recommender-system effects (6314747) and analyses of image/visual platforms and amplification dynamics (3303310) to support the claim that interactions both constitute production and feed ranking loops.} \hllightgreen{Algorithmic prioritisation of high arousal or status laden imagery has documented psychological and equity impacts.}\pdfcomment[icon=Note]{Empirical and technical sources document priority of high-arousal/status content and associated psychological/equity harms (3303310 on amplification and equity impacts; 8958037 on addiction mechanisms), supporting this synthesis.} \hldarkgreen{Hu reports that Instagram's algorithmic "Explore" page amplified body negative imagery by $23\%$ relative to chronological feeds, correlating with a statistically significant decline in self esteem among adolescent women of color (regression coefficient $\beta = -0.41$, $p < 0.05$).}\pdfcomment[icon=Note]{The statistic and regression result are explicitly reported in the conference paper (3303310) stating Instagram's Explore page amplified body‑negative imagery by 23\% with β = −0.41, p < 0.05 (3303310, p.3--4).} \hldarkgreen{Sentiment analysis in the same work shows that algorithmically prioritised content exhibits $31\%$ higher negative affect than non prioritised baselines, reinforcing a pattern where engagement optimisation systematically favours harmful material \parencites[3]{3303310}.}\pdfcomment[icon=Note]{The same source reports sentiment analysis showing algorithmically prioritized content exhibited 31\% higher negative affect (3303310, p.3--4).} \hllightgreen{These findings position image centric platforms as high risk environments in this report's working taxonomy where recommendation objectives conflict with mental health and anti discrimination duties \parencites[3]{3303310}[5]{6314747}.}\pdfcomment[icon=Note]{6314747 and 3303310 document recommender harms and amplification effects; synthesizing these supports classifying image-centric platforms as higher risk in a taxonomy where recommendation objectives conflict with mental-health and anti-discrimination duties (6314747 p.5; 3303310 p.3).} \hldarkgreen{From an engineering and creator governance angle, Zeng et al. show that Instagram and TikTok form key categories in empirical studies of AI enabled creator commerce.}\pdfcomment[icon=Note]{5673308 explicitly treats platform types including Instagram and TikTok as key categories in empirical studies of AI-enabled creator commerce (5673308, p.3).} \hldarkgreen{Their mixed method research evaluates AI driven product architectures across platforms and identifies platform type (YouTube, Instagram, TikTok), content format (video, short form, image) and audience size as moderating variables in how AI intelligence engineering affects creator performance, engagement and revenue.}\pdfcomment[icon=Note]{5673308 describes a mixed‑methods evaluation of AI-driven product architectures and identifies platform type, content format and audience size as moderating variables (5673308, p.3).} \hldarkgreen{The study combines survey data from content creators and engineers with secondary platform analytics (engagement rates, click through ratios, retention metrics) and uses structural equation modelling to quantify relationships between AI integration frequency, system usability and outcomes such as conversion rate and revenue growth.}\pdfcomment[icon=Note]{5673308 details combining survey/interview data with secondary platform analytics and applying structural equation modelling to quantify relationships (5673308, p.3).} \hllightgreen{These results confirm that image and story centric networks are already embedded in highly optimised, AI intensive monetisation architectures in which engagement metrics are directly coupled to business value \parencites[3]{5673308}.}\pdfcomment[icon=Note]{Synthesizes findings from 5673308 on AI-driven monetization architectures and additional policy/technical analyses (e.g., 6502223 on engagement‑driven monetization) to support the claim that these networks are embedded in optimized, AI‑intensive monetisation systems.} \hldarkgreen{Regulatory perspectives on these platforms arise primarily through the lens of behavioral advertising.}\pdfcomment[icon=Note]{6941611 frames regulatory perspectives on platforms principally through behavioural advertising and consumer protection concerns (6941611, p.4).} \hldarkgreen{Quinelato notes that social media services like Facebook, Instagram and TikTok are leading examples of digital platforms that use behavioural advertising by profiling vulnerable demographics across cookies and third party tracking infrastructure.}\pdfcomment[icon=Note]{6941611 directly notes Facebook, Instagram and TikTok as leading examples of platforms using behavioural advertising and profiling vulnerable demographics via tracking (6941611, p.4).} \hldarkgreen{She stresses that these platforms exploit their unparalleled access to interaction data to personalise products and services and that this personalisation includes highly individualised advertising on the basis of profiles built from observed behaviour.}\pdfcomment[icon=Note]{6941611 states platforms exploit access to interaction data to personalise products/services, including individualised advertising based on behaviourally built profiles (6941611, p.4).} \hldarkgreen{In the same analysis, the need for regulatory measures is framed as a response to opaque persuasive practices and profiling of vulnerable groups rather than to technical innovation per se \parencites[4]{6941611}.}\pdfcomment[icon=Note]{6941611 frames the need for regulation as a response to opaque persuasive practices and profiling of vulnerable groups, emphasizing consumer protection rather than pure technical novelty (6941611, p.4).} \hllightgreen{User side vulnerability interacts with these design and profiling choices.}\pdfcomment[icon=Note]{This is a reasonable inference supported by governance and stakeholder analyses that link user vulnerabilities to design/profiling choices (2742041 on stakeholder impacts; 2616018 on vulnerable groups), so LIGHTGREEN is appropriate.} \hldarkgreen{Survey evidence on digital literacy and awareness of GDPR, the DSA and the AI Act indicates that highly active social network users report strong confidence in everyday tasks and online purchases (average Likert scores $4.40$--$4.62$) yet much lower competence in recognising hidden manipulations (mean $2.31$) and avoiding fraud ($2.13$).}\pdfcomment[icon=Note]{Survey results with the exact Likert ranges and lower scores for hidden manipulation and fraud awareness are reported in 9360054 (9360054, p.10).} \hldarkgreen{Participants express positive evaluations of EU regulations on data privacy (mean $4.09$) and content moderation on social networks (mean $3.91$), but remain sceptical about their effectiveness against disinformation, cyberbullying and online fraud, especially in cross border contexts such as Ukraine where EU rules apply only indirectly \parencites[10]{9360054}.}\pdfcomment[icon=Note]{9360054 reports mean scores of 4.09 for EU data-privacy evaluation and 3.91 for content moderation, and documents scepticism about effectiveness especially in cross-border contexts like Ukraine (9360054, p.10).} \hllightgreen{This combination of high routine confidence and low manipulation awareness suggests that users of image and story centric platforms may be particularly exposed to subtle behavioural advertising and interface pressure \parencites[10]{9360054}[4]{6941611}.}\pdfcomment[icon=Note]{Combines user‑confidence/low‑awareness survey data (9360054) with the regulatory/behavioural advertising critique (6941611) to reasonably infer heightened exposure to subtle behavioural advertising on image/story platforms.} \hllightgreen{AI governance frameworks provide an external normative reference for evaluating these networks.}\pdfcomment[icon=Note]{2742041 and other governance texts present AI governance frameworks as external normative references for evaluation; this is a reasonable synthesis of those sources (2742041, p.2).} \hldarkgreen{Leon's Algorithmic Impact Navigator treats transparency, fairness and respect for autonomy as constant principles across sectors and calls on organisations deploying AI to conduct comprehensive impact evaluations, map regulatory obligations and incorporate stakeholder feedback early in the lifecycle.}\pdfcomment[icon=Note]{2742041 (the Algorithmic Impact Navigator discussion) treats transparency, fairness and respect for autonomy as core principles and calls for comprehensive impact evaluations and stakeholder engagement (2742041, p.2).} \hldarkgreen{Developers are expected to embed fairness, privacy and explainability directly into system architectures, including bias mitigation and interpretability mechanisms \parencites[2]{2742041}.}\pdfcomment[icon=Note]{2742041 explicitly expects developers to embed fairness, privacy and explainability into system architectures including bias mitigation and interpretability mechanisms (2742041, p.2).} \hldarkgreen{Mutawa et al.'s Algorithmic Governance Assessment Matrix, used for ethical AI in zakat distribution, operationalises comparable dimensions: transparency, accountability, explainability, data integrity and alignment with normative principles, supported by coding sheets and secondary performance indicators such as transparency scores \parencites[4]{2702055}.}\pdfcomment[icon=Note]{2702055 describes an Algorithmic Governance Assessment Matrix for zakat distribution operationalising dimensions such as transparency, accountability, explainability, data integrity and normative alignment with coding sheets and indicators (2702055, p.4).} \hllightgreen{Applying these templates in this thesis, image and story centric social networks are selected as audit targets because their AI intensive architectures, behavioural advertising practices and documented psychosocial impacts sit exactly in the zone where such assessment matrices and ethical audit obligations become practically relevant \parencites[3]{3303310}[5]{6314747}[3]{5673308}[2]{2742041}.}\pdfcomment[icon=Note]{Synthesis: 6314747 (recommender harms), 5673308 (creator commerce/AI architectures) and 2742041 (AIN governance) together justify selecting image/story networks as audit targets where such assessment matrices are relevant (6314747 p.5; 5673308 p.3; 2742041 p.2).} \hldarkgreen{The ethical data science literature also emphasises that fairness and accountability are processual, not one-off achievements; lifecycle frameworks and documentation practices help translate abstract principles into actionable checks that can persist after deployment \parencites[20]{2616018}.}\pdfcomment[icon=Note]{The ethical data science literature emphasizes lifecycle/processual approaches to fairness and accountability and the role of documentation practices (2616018 discusses lifecycle frameworks and documentation, e.g., Model Cards and Datasheets) (2616018, p.20).} \hllightgreen{In short, treating these networks as audit objects makes practical sense because the same documentation and monitoring tools that reduce harm in clinical or criminal justice settings appear likely to reveal representational gaps and feedback loops in image based recommendation systems \parencites[20]{2616018}.}\pdfcomment[icon=Note]{Reasonable synthesis: intervention and documentation tools that reduce harm in clinical (5419804) and criminal‑justice settings (2616018 case studies) are likely to reveal representational gaps and feedback loops in image‑based recommender systems.}

\subsubsection{Professional Networking Platforms with Algorithmic Feeds}
\label{sec:prof_networking_algo_feeds} \hllightgreen{Professional networking environments with algorithmic feeds warrant inclusion in this audit because they combine social media engagement mechanics with economic gatekeeping around jobs, reputation, and professional opportunity.}\pdfcomment[icon=Note]{Sources describe recommender-driven social platforms shaping economic outcomes and visibility (6502223 p.2; 4905459 p.54), supporting the inference that professional networking with algorithmic feeds combines engagement mechanics and economic gatekeeping.} \hldarkgreen{Chitra and colleagues explicitly list platforms such as Facebook, Instagram, TikTok, X and Google Search as central nodes in political information flows that now eclipse traditional media in influence, with recommender systems acting as "invisible gatekeepers" for what citizens see and believe \parencites[2]{6502223}.}\pdfcomment[icon=Note]{Chitra et al. explicitly name platforms such as Facebook, Instagram, TikTok, X and Google Search and characterise recommendation systems as 'invisible gatekeepers' (6502223 p.2, abstract).} \hllightgreen{Although LinkedIn is not named in the available sources, image and story centric networks and general social media already carry professional branding, job discovery and recruiter interactions, blurring boundaries between leisure feeds and professional signalling \parencites[3]{3303310}[5]{6314747}.}\pdfcomment[icon=Note]{While LinkedIn is not emphasised in the provided excerpts, the literature documents that image/story networks and general social media carry professional signalling and affect job/reputation dynamics (4905459 p.54; 5142478 p.3), supporting this claim as a reasonable synthesis.} \hllightgreen{Engagement based ranking in professional contexts raises specific risks around filter bubbles, echo chambers and political or ideological skew.}\pdfcomment[icon=Note]{Multiple sources link engagement-driven ranking to filter bubbles, echo chambers and ideological skew (6502223 p.17; 5142478 p.26), so applying this risk to professional contexts is a valid academic inference.} \hldarkgreen{Gerr's mixed method study of Facebook, YouTube and TikTok shows that algorithmic personalisation systematically diminishes exposure to diverse viewpoints and reinforces echo chambers that limit critical thinking and civic participation \parencites[5]{5142478}[29]{5142478}.}\pdfcomment[icon=Note]{Gerr's mixed-method study reports that personalization reduces exposure to diverse viewpoints and reinforces echo chambers, limiting critical thinking and civic participation (5142478 p.29, Discussion).} \hldarkgreen{Chitra et al. document platform level amplification patterns: Facebook's engagement algorithm directed $64\%$ of users who joined extremist groups to those groups; YouTube's recommendations amplified right leaning political content more than equivalent left leaning material; Twitter's amplification system privileged right leaning content in most tested countries; and TikTok accounts expressing political interest were "quickly directed toward extremist material" \parencites[8]{6502223}.}\pdfcomment[icon=Note]{Chitra et al. document the platform-level amplification patterns including Facebook's 64\% statistic, YouTube's right-leaning amplification, Twitter/X privileging right-leaning content, and TikTok directing political-interest accounts toward extremist material (6502223 p.8; p.17).} \hllightgreen{Professional feeds that surface political, ESG or "thought leadership" content through similar engagement logics may therefore import these biases into hiring, networking and reputational signals, a concern that justifies treating professional networks as part of the same algorithmic ecosystem.}\pdfcomment[icon=Note]{Given documented amplification biases on mainstream platforms (6502223 p.8; 5142478 p.26), it is a reasonable inference that similar engagement logics in professional feeds could import such biases into hiring and reputational signals.} \hllightgreen{Behaviourally targeted advertising and profiling also intersect with professional identity construction.}\pdfcomment[icon=Note]{The literature documents behavioural micro-targeting and profiling practices (6502223 p.17) and the role of user data in personalization (4905459 p.21), supporting the claim that targeted advertising intersects with professional identity construction.} \hlorange{Quinelato characterises social media services, including Instagram and TikTok, as digital platforms that profile users through cookies and third party tracking infrastructure in order to deliver highly individualised advertising, including political messaging \parencites[4]{6941611}.}\pdfcomment[icon=Note]{No provided source attributed to 'Quinelato' appears in the materials; none of the supplied excerpts cite that author or provide the specific characterisation described.} \hllightgreen{The same profiling pipelines can support job ads, sponsored courses and professional services targeted on inferred vulnerabilities or aspirations rather than transparent qualifications.}\pdfcomment[icon=Note]{Micro-targeting and behavioural profiling literature shows that the same pipelines serve varied ad and service targeting (6502223 p.17), supporting the plausible claim that job ads and sponsored professional services could be targeted on inferred vulnerabilities rather than transparent qualifications.} \hldarkgreen{Survey data from Chukhray and Koval show that highly active social network users report strong confidence in routine digital tasks and online purchases (Likert means $4.40$--$4.62$), yet much lower ability to recognise hidden manipulations (mean $2.31$) and avoid fraud ($2.13$), with scepticism about the effectiveness of GDPR, the DSA and the AI Act in countering disinformation and online fraud \parencites[10]{9360054}.}\pdfcomment[icon=Note]{The survey statistics (Likert means 4.40--4.62 for routine digital tasks, 2.31 for recognising hidden manipulations, 2.13 for avoiding fraud, and scepticism about GDPR/DSA/AI Act effectiveness) are reported in Chukhray and Koval (9360054 p.10).} \hllightgreen{When such users encounter algorithmically sorted professional content, they may overestimate their ability to detect biased or manipulative ranking, which increases ethical and regulatory risk.}\pdfcomment[icon=Note]{Based on the documented overconfidence in routine tasks yet low ability to detect manipulations (9360054 p.10), it is a reasonable inference that such users may overestimate detection of biased ranking when encountering algorithmic professional content.} \hllightgreen{Normative frameworks on AI governance and digital fairness supply an external lens for assessing these platforms.}\pdfcomment[icon=Note]{Multiple provided sources outline normative AI governance and digital fairness frameworks (2977039 p.25; 6502223 p.2), supporting the statement that these frameworks supply an external lens for assessment.} \hldarkgreen{Gerr's analysis of the DSA and EU AI Act stresses that current regulation has improved transparency obligations, such as disclosing main recommendation parameters and offering non personalised feeds, yet engagement driven biases in ranking algorithms remain largely unaddressed.}\pdfcomment[icon=Note]{Gerr analyses DSA and the EU AI Act, noting improved transparency obligations (e.g., disclosure of recommendation parameters) while engagement-driven biases remain largely unaddressed (5142478 p.7--8).} \hldarkgreen{Case studies show that YouTube continues to prioritise watch time even where this favours sensationalist over balanced content, and that TikTok struggles to moderate climate misinformation despite formal bans, indicating that technical compliance with transparency rules does not remove deeper structural incentives \parencites[7]{5142478}[28]{5142478}.}\pdfcomment[icon=Note]{Gerr's case studies report YouTube's prioritisation of watch time favoring sensationalist content and TikTok's continued struggles with climate misinformation despite bans, arguing technical transparency does not remove structural incentives (5142478 p.28--29).} \hldarkgreen{Chitra et al. conclude that safeguarding democracy under algorithmic recommendation requires mandatory transparency, independent third party auditing, AI literacy and conditional liability protection tied to meeting algorithmic accountability standards \parencites[17]{6502223}[2253]{6502223}.}\pdfcomment[icon=Note]{Chitra et al. recommend mandatory transparency, independent third-party auditing, AI literacy, and conditional liability tied to algorithmic accountability (6502223 p.17; p.2253 section).} \hllightgreen{Professional networking feeds that influence employability and reputation can reasonably be expected to meet similar auditability and accountability conditions in a Digital Fairness trajectory.}\pdfcomment[icon=Note]{Given the parallels drawn between platform harms and regulatory recommendations in the democratic context (6502223 p.17; 5142478 p.7), it is a reasonable inference that professional networking feeds should meet similar auditability and accountability conditions.} \hllightgreen{Education oriented findings around user competence and AI literacy underline why these safeguards matter.}\pdfcomment[icon=Note]{Education-oriented findings on user competence and AI literacy (9360054 p.10) support the claim that these findings underline why safeguards like literacy and transparency matter.} \hldarkgreen{Chukhray and Koval argue that users experience a "second level divide": everyday skills are high, yet understanding of behavioural biases and algorithmic influence remains limited, particularly outside the EU, where regulations operate only indirectly through association agreements \parencites[10]{9360054}.}\pdfcomment[icon=Note]{Chukhray and Koval explicitly describe a 'second-level divide', high everyday skills but limited understanding of behavioural biases and algorithmic influence, with weaker effects of EU regulation outside the EU (9360054 p.10).} \hllightgreen{This gap supports the case for algorithmic literacy and explainable recommendation in professional contexts, for instance transparency dashboards that reveal what data is used to personalise job or content recommendations, similar to proposals in democratic governance work that envisage user control over the ideological diversity of recommendation outputs \parencites[17]{6502223}.}\pdfcomment[icon=Note]{The literature advocates algorithmic literacy and explainability (9360054 p.10) and Chitra et al. propose transparency dashboards and ideological diversity calibration (6502223 p.17, p.2253), supporting this combined proposal.} \hllightgreen{Methodologically, these platforms are assessed here under the same constraint that applies throughout this thesis: internal recommendation code, training data and telemetry remain proprietary.}\pdfcomment[icon=Note]{Several sources note methodological constraints due to proprietary internal recommendation code and training data (5142478 p.2; 4905459 p.50), supporting this statement about assessment constraints.} \hldarkgreen{Analysis therefore relies on secondary evidence documenting engagement based amplification, profiling practices and regulatory responses on adjacent platforms, together with survey based indicators of user vulnerability and regulatory capacity \parencites[6]{6295677}[3]{5142478}[10]{9360054}.}\pdfcomment[icon=Note]{The approach described, relying on secondary evidence of engagement amplification, profiling practices and regulatory responses plus survey indicators, matches the methods and sources cited (5142478 p.2; 9360054 p.10).} \hllightgreen{Within that constraint, professional networking services with algorithmic feeds are selected as a priority audit category because they combine high engagement optimisation with direct consequences for jobs and social status, and operate in an environment where EU style rules on transparency, non profiled feeds, algorithmic auditing and AI literacy are already being articulated as conditions for legitimate use \parencites[5]{5142478}[2244]{6502223}[25]{2977039}.}\pdfcomment[icon=Note]{The selection rationale is supported by discussion of platforms' high engagement optimisation and job/reputation effects (6502223 p.2) and by policy texts articulating EU-style rules (2977039 p.25) as emerging conditions.} \hldarkgreen{A broader literature on recommender systems also highlights that these platforms are not merely passenger technologies but multi-stakeholder mechanisms where provider objectives, item suppliers and societal effects can pull ranking away from simple notions of individual user benefit; this perspective suggests that audits should examine platform KPI incentives as well as user-facing outputs \parencites[53]{4905459}.}\pdfcomment[icon=Note]{Recommender systems scholarship frames these systems as multi-stakeholder mechanisms where provider objectives and societal effects can diverge from individual user benefit, implying audits should examine KPI incentives as well as outputs (4905459 p.53; p.58).} \hldarkgreen{Relatedly, scholarship emphasises the need to include human centred evaluation and fairness metrics beyond offline accuracy, since conventional performance measures may miss long term or systemic harms that arise when professional opportunity is mediated by algorithmic exposure \parencites[53]{4905459}.}\pdfcomment[icon=Note]{The recommender systems literature stresses human-centred evaluation and fairness metrics beyond offline accuracy because conventional measures can miss long-term or systemic harms (4905459 p.50; p.58).}

\subsection{Interface Mechanics across Platforms}

\subsubsection{Feed Structures and Scrolling Patterns}
\label{sec:feed_structures_scrolling_patterns} \hldarkgreen{Short form video feeds are structured around a default, continuously updating stream that removes natural stopping cues.}\pdfcomment[icon=Note]{Sources describe short-video platforms as using continuous personalized streams and removing natural stopping cues (Li 8958037 p.2; Conf-CDS 3303310 p.2).} \hldarkgreen{In TikTok's case, opening the app lands users directly in the personalised "For You" feed rather than a chronological "Following" view, which places the algorithmic stream at the centre of the user journey \parencites[42]{3195574}.}\pdfcomment[icon=Note]{Meßmer \& Degeling explicitly note that TikTok opens users into the personalized 'For You' feed rather than a chronological 'Following' view (Meßmer \& Degeling 3195574 p.42).} \hldarkgreen{Li describes comparable short video platforms that combine micro duration clips with infinite swipe navigation, where each downward motion instantly loads the next video from a neural network based recommendation engine \parencites[2]{8958037}.}\pdfcomment[icon=Note]{Li documents short-video platforms combining micro-duration clips with infinite swipe navigation and neural-network based recommendation systems (Li 8958037 p.2, p.6).} \hldarkgreen{This pairing of always on feed and single gesture advancement is not incidental: it creates an uninterrupted chain of stimuli that users, especially college students, experience as an "auto pilot mode" of persistent scrolling, with excessive use associated with weakened social skills, heightened anxiety and the development of behavioural addiction \parencites[2]{8958037}[6]{8958037}.}\pdfcomment[icon=Note]{Li links the always-on feed and single-gesture advancement to an 'auto-pilot mode' among college students and associates excessive use with weakened social skills, anxiety and behavioural addiction (Li 8958037 p.2, p.6).} \hllightgreen{Engagement optimised scrolling patterns also appear through auto play.}\pdfcomment[icon=Note]{Autoplay is discussed as an engagement-optimizing pattern in the literature and interfaces (Meßmer \& Degeling 3195574 p.42; Conf-CDS 3303310 p.2--3), so the claim is a reasonable synthesis.} \hldarkgreen{On YouTube, the next video starts by default when the current one finishes, a configuration explicitly cited by Me{\ss}mer and Degeling as an interface choice that emphasises recommended content over items consciously selected by the user \parencites[42]{3195574}.}\pdfcomment[icon=Note]{Meßmer \& Degeling specifically cite YouTube's default setting where the next video starts automatically, framing this as emphasising recommended content (Meßmer \& Degeling 3195574 p.42).} \hldarkgreen{Li reports that short video platforms integrate similar auto play behaviour: once a clip ends, the next item immediately appears, sourced from the personalised distribution rather than any explicit request \parencites[2]{8958037}.}\pdfcomment[icon=Note]{Li reports that short-video platforms integrate autoplay behaviour where the next clip is immediately loaded from personalized distribution rather than an explicit user request (Li 8958037 p.2).} \hldarkgreen{When combined with infinite scroll, auto play reduces the number of intentional decisions per unit of content and keeps users within the recommendation loop for longer periods, a structural feature that aligns with Hu's description of "dopamine driven feedback loops" on algorithmically curated feeds \parencites[2]{3303310}[20]{3303310}.}\pdfcomment[icon=Note]{Conf-CDS synthesizes evidence that infinite scroll plus autoplay reduces intentional decisions and aligns with 'dopamine-driven feedback loops', including empirical figures (Conf-CDS 3303310 p.2--3).} \hllightgreen{Scrolling patterns interact directly with psychological triggers.}\pdfcomment[icon=Note]{Multiple sources connect scrolling/interface patterns to psychological mechanisms (Li 8958037 p.2; Conf-CDS 3303310 p.2--3), so this general claim is supported as a synthesis.} \hldarkgreen{Hu links infinite scrolling interfaces and variable reward schedules on platforms "like TikTok" to neural pathways associated with addiction, citing empirical validations that users of algorithmically curated feeds exhibit $22\%$ higher rates of compulsive checking behaviour than users of non algorithmic interfaces \parencites[2]{3303310}.}\pdfcomment[icon=Note]{Conf-CDS links infinite scrolling and variable reward schedules on platforms like TikTok to addiction-related neural pathways and reports users of algorithmic feeds exhibiting 22\% higher compulsive checking (Conf-CDS 3303310 p.2).} \hldarkgreen{Li's analysis of short video usage among Chinese college students connects these feed structures to impaired sleep, elevated stress, increased depression and anxiety, and disrupted study rhythms, with critics labelling short videos "time sinkholes" that impede engagement with serious pursuits.}\pdfcomment[icon=Note]{Li's study of Chinese college students documents impaired sleep, elevated stress, increased depression and anxiety, disrupted study rhythms, and labels short videos as 'time sinkholes' (Li 8958037 p.6).} \hllightgreen{These findings justify treating infinite scroll and default auto play as high risk Code A features in this report's working taxonomy.}\pdfcomment[icon=Note]{Given the empirical harms and regulatory concern documented in the cited sources (Li 8958037 p.2; Conf-CDS 3303310 p.2--3), treating infinite scroll and autoplay as high-risk inspection targets is a defensible synthesis.} \hllightgreen{Regulatory documents have begun to name feed structures and scrolling behaviour explicitly.}\pdfcomment[icon=Note]{Regulatory documents and analyses increasingly address feed structures and scrolling (e.g., DSA discussions and audits in Meßmer \& Degeling 3195574 and policy reviews 7896734), so this general statement is supported by the sources.} \hldarkgreen{Under the Digital Services Act, the European Commission opened an investigation into TikTok in 2026 and raised concerns that its algorithmic recommendation mechanisms and interface features, including infinite scrolling, auto play, and push notifications, may constitute "systematic manipulative design" that impairs users' capacity for free and informed decisions \parencites[6]{8958037}.}\pdfcomment[icon=Note]{Li notes that in 2026 the EU under the DSA opened an investigation into TikTok over algorithmic recommendation mechanisms and interface features including infinite scrolling and autoplay (Li 8958037 p.2).} \hldarkgreen{Me{\ss}mer and Degeling link such design choices to the DSA's prohibition of dark patterns under Article 25, noting examples such as default auto play and starting the app on the personalised feed as UX decisions that tilt the balance toward recommended content \parencites[42]{3195574}.}\pdfcomment[icon=Note]{Meßmer \& Degeling link such UX choices to the DSA's prohibition of dark patterns (Article 25) and give examples like default autoplay and starting on the personalised feed (Meßmer \& Degeling 3195574 p.42).} \hllightgreen{These regulatory signals anchor Code C classifications for scrolling patterns that systematically undermine user agency.}\pdfcomment[icon=Note]{This is a reasonable conclusion from the regulatory references and risk framing in the sources (Meßmer \& Degeling 3195574; policy analyses 7896734), representing an interpretive synthesis.} \hllightgreen{Feed structures also influence media diversity and polarisation.}\pdfcomment[icon=Note]{Sources link feed mechanisms to effects on media diversity and polarization (Gerr 5142478; Conf-CDS 3303310), so this general claim is a supported synthesis.} \hldarkgreen{Gerr's mixed method audit of Facebook, YouTube and TikTok shows that engagement optimised trending and recommendation surfaces reduce exposure to heterogeneous viewpoints and reinforce echo chambers, which restricts opportunities for critical thinking and civic participation \parencites[3]{5142478}[29]{5142478}.}\pdfcomment[icon=Note]{Gerr's mixed-method audit reports that engagement-optimised trending and recommendation surfaces reduce exposure to heterogeneous viewpoints and reinforce echo chambers (Gerr 5142478 p.8, p.24--29).} \hldarkgreen{Hu's synthesis describes how collaborative filtering algorithms amplify high arousal content such as outrage or sensationalism and documents that Instagram's "Explore" page increased exposure to body negative imagery by $23\%$ relative to chronological feeds, with a statistically significant decline in self esteem among adolescent women of color \parencites[3]{3303310}.}\pdfcomment[icon=Note]{Conf-CDS documents that collaborative filtering amplifies high-arousal content and reports Instagram's Explore page increased exposure to body-negative imagery by 23\% with related declines in self-esteem (Conf-CDS 3303310 p.2--3).} \hllightgreen{In both cases, the feed's ranking logic and scrolling dynamics jointly determine which content clusters remain persistently visible as users move through the stream.}\pdfcomment[icon=Note]{This synthesis that ranking logic and scrolling dynamics jointly affect persistent visibility is supported by the technical and behavioral analyses in Conf-CDS and Gerr (Conf-CDS 3303310 p.2--3; Gerr 5142478 p.24).} \hllightgreen{Inspection of UX and advocacy reports shows that platform feed structures can either support or obstruct user control.}\pdfcomment[icon=Note]{UX and advocacy reports (Meßmer \& Degeling 3195574; Gerr 5142478) discuss how feed structures can enable or constrain user control, so this general claim is supported as synthesis.} \hldarkgreen{Me{\ss}mer and Degeling emphasise that the number of clicks required to move from a personalised feed to alternative views, such as "Following" or non profiled recommendations, is a critical design variable for auditing risk scenarios, alongside the default choice of landing view and the presence of auto play toggles \parencites[42]{3195574}.}\pdfcomment[icon=Note]{Meßmer \& Degeling emphasize the number of clicks to reach alternative views, default landing view, and presence of autoplay toggles as critical audit variables (Meßmer \& Degeling 3195574 p.42).} \hldarkgreen{In education, human in the loop grading systems demonstrate a contrasting pattern: interfaces expose AI scores, SHAP based explanations and editable feedback side by side, with instructors able to modify or override outputs, and all decisions logged with timestamps and reviewer identifiers \parencites[17]{6805789}.}\pdfcomment[icon=Note]{The HITL grading study describes side-by-side AI scores, SHAP explanations, editable feedback, instructor overrides, and logged timestamps/identifiers (HITL study 6805789 p.17, p.21).} \hllightgreen{Although not a social feed, this design illustrates that streaming information can be structured in ways that foreground human oversight rather than engagement maximisation.}\pdfcomment[icon=Note]{Using the HITL example to illustrate that streaming information can be structured to foreground human oversight is a reasonable inference from the described interface design (HITL 6805789 p.17--21).} \hllightgreen{These observations are constrained by an important methodological limit.}\pdfcomment[icon=Note]{The report itself notes methodological constraints (e.g., lack of access to proprietary internals) and this sentence flags that limitation, which is consistent with the sources' methodological discussion (Meßmer \& Degeling 3195574; Li 8958037).} \hllightgreen{Internal recommendation code, training data and real time clickstream telemetry for social platforms are not available to this project.}\pdfcomment[icon=Note]{Authors and auditing literature state that internal recommendation code, training data and real-time clickstream telemetry are typically inaccessible to external projects (Meßmer \& Degeling 3195574; policy reports 7896734), supporting this claim.} \hldarkgreen{Feed structures and scrolling patterns are therefore assessed through published UX analyses, regulatory proceedings, empirical usage studies and conceptual work on algorithmic engagement, not through direct inspection of proprietary algorithms \parencites[6]{6295677}[2]{8958037}[2]{3303310}[42]{3195574}.}\pdfcomment[icon=Note]{Meßmer \& Degeling and the cited empirical and conceptual works explain that feed structures are assessed via published UX analyses, regulatory proceedings and empirical studies rather than direct inspection of proprietary algorithms (Meßmer \& Degeling 3195574 p.42; Li 8958037 p.2; Conf-CDS 3303310 p.2).} \hllightgreen{Within this constraint, personalisable infinite scroll, default auto play and landing on personalised feeds emerge as cross platform design motifs that the Ethical AI Audit Toolkit must treat as core inspection targets under a Digital Fairness trajectory.}\pdfcomment[icon=Note]{Given the documented cross-platform recurrence of personalisable infinite scroll, autoplay and landing on personalised feeds and their risk implications (Li 8958037; Meßmer \& Degeling 3195574; Conf-CDS 3303310), treating them as core inspection targets is a supported synthesis.}

\subsubsection{Use of Autoplay and Recommendations}
\label{sec:use_autoplay_recommendations} \hllightgreen{Autoplay sits at the centre of several engagement loops examined in this audit.}\pdfcomment[icon=Note]{Multiple sources (Li 8958037; Meßmer \& Degeling 3195574; 3303310) characterise autoplay as central to engagement loops on short-video platforms, supporting this summary claim.} \hllightgreen{Short video platforms couple micro duration clips with feeds in which the next item is automatically presented once the current one finishes, so that opting out requires an active interruption rather than a positive choice to continue.}\pdfcomment[icon=Note]{Li (8958037, p.2) and related literature (3303310) describe micro-duration clips and feeds that auto-advance, implying that continuing requires no positive action while opting out requires interruption.} \hldarkgreen{Li's analysis of Chinese short video services describes infinite swipe navigation, with each downward motion loading another algorithmically selected clip, and identifies auto play alongside infinite scrolling and push notifications as part of an "addictive design" under investigation by the European Commission in its 2026 TikTok proceeding under the DSA \parencites[2]{8958037}.}\pdfcomment[icon=Note]{Dan Li (8958037, p.2) documents infinite swipe navigation and identifies autoplay with infinite scrolling and push notifications as part of an 'addictive design', and Meßmer \& Degeling (3195574, p.42) reference EU scrutiny of such features in the TikTok/DSA context.} \hldarkgreen{Me{\ss}mer and Degeling give a parallel example for YouTube, where the auto play function that starts the next video when the current one ends is enabled by default and is explicitly singled out as a design element that emphasises recommended content over items users manually choose from subscription lists or searches \parencites[42]{3195574}.}\pdfcomment[icon=Note]{Meßmer \& Degeling explicitly note that YouTube's autoplay is enabled by default and emphasise recommended content over followed content (3195574, p.42).} \hllightgreen{Recommendation systems are tightly interwoven with these autoplay mechanics.}\pdfcomment[icon=Note]{Sources describe close coupling of recommendation systems and autoplay mechanics (Li 8958037; 3303310), supporting the claim that they are tightly interwoven.} \hllightgreen{Short video feeds use neural network based recommenders that continuously refine a "For You" style stream from behavioural data, so that auto played items are not random fillers but personalised selections aligned with inferred preferences \parencites[2]{8958037}.}\pdfcomment[icon=Note]{Li (8958037, p.2) describes neural-network based recommenders producing 'For You' style streams from behavioural data, showing autoplayed items are personalised rather than random.} \hldarkgreen{Me{\ss}mer and Degeling note that the TikTok app always opens on the personalised "For You" feed rather than on a "Following" list, which means that both the starting point and the subsequent auto played sequence are determined by recommendation algorithms, not by explicit user navigation \parencites[42]{3195574}.}\pdfcomment[icon=Note]{Meßmer \& Degeling state the TikTok app opens on the personalised 'For You' feed (not 'Following'), meaning start point and subsequent autoplay are algorithm-determined (3195574, p.42).} \hldarkgreen{In this configuration, autoplay does more than save a click: it systematically privileges algorithmic choices over self directed consumption.}\pdfcomment[icon=Note]{Meßmer \& Degeling (3195574, p.42) show autoplay and default personalised landing emphasise recommended content, supporting the claim that autoplay privileges algorithmic choices over self-directed consumption.} \hllightgreen{Evidence on psychological and social effects of such autoplay driven recommendation loops is increasingly detailed.}\pdfcomment[icon=Note]{Several empirical and review sources (e.g., 3303310; Frontiers 2580223; Li 8958037) document growing detailed evidence on psychological and social effects of autoplay-driven recommendation loops.} \hlorange{Hu synthesises behavioural and neurocognitive work showing that infinite scrolling interfaces and variable reward schedules, described as "core features of platforms like TikTok," exploit dopaminergic feedback mechanisms associated with addiction.}\pdfcomment[icon=Note]{No provided source attributed to an author named 'Hu' making this specific synthesis; while 3303310 and other sources discuss dopaminergic feedback and infinite scrolling, the specific claim that 'Hu synthesises...' is unsupported by the given sources.} \hldarkgreen{Users of algorithmically curated feeds display $22\%$ higher rates of compulsive checking than those using non algorithmic interfaces, and hyper personalised feeds are linked to a $23\%$ higher incidence of depressive symptoms, as well as reduced prefrontal cortex activity and gray matter density in regions governing impulse control \parencites[2]{3303310}.}\pdfcomment[icon=Note]{The figures and neuroimaging claims (22\% higher compulsive checking; 23\% higher depressive incidence; reduced prefrontal activity/gray matter) are reported in the interdisciplinary review excerpt (3303310, p.2) and aligned with its cited empirical findings.} \hldarkgreen{Li connects the same feed structures to diminished academic concentration, weakened social skills, and heightened anxiety and depression among college students in China, framing the autoplayed short video stream as an "algorithmic addiction loop" that channels users into an auto pilot mode of continuous viewing \parencites[2]{8958037}.}\pdfcomment[icon=Note]{Dan Li explicitly links feed structures to diminished academic concentration, weakened social skills, heightened anxiety/depression among Chinese college students and frames an 'algorithmic addiction loop' (8958037, pp.1--2).} \hllightgreen{Regulatory and model law sources treat this combination of autoplay and personalised recommendation as a potential form of manipulative design.}\pdfcomment[icon=Note]{Regulatory and model law texts (DSA discussions in Meßmer \& Degeling 3195574; AI Law Model 2977039) characterise autoplay plus personalised recommendation as potentially manipulative design, supporting this synthesis.} \hldarkgreen{The DSA prohibits "dark patterns" that prevent users from making free and informed decisions and has already opened a proceeding into TikTok's algorithmic recommendation mechanisms and interface features, including infinite scrolling, auto play, and push notifications, on precisely this basis \parencites[2]{8958037}[42]{3195574}.}\pdfcomment[icon=Note]{Meßmer \& Degeling note the DSA prohibits 'dark patterns' that impede free informed decisions and link this to scrutiny of algorithmic interface features; Li (8958037, p.2) documents the 2026 DSA/TikTok proceeding into infinite scrolling, autoplay and push notifications.} \hllightgreen{Kostenko's AI Law Model goes further for minors.}\pdfcomment[icon=Note]{The AI Law Model (2977039) contains age-appropriate design provisions that are more protective for minors, supporting the general claim that Kostenko's model goes further for minors.} \hldarkgreen{Under the principle of age appropriate design, manipulative and addictive mechanics such as endless scrolling and autoplay are explicitly prohibited for child directed systems, and "social networks for teenagers" in the model's normative illustration are required to operate feeds in "detox mode" by default, specifically "without endless scrolling, autoplay of the video and aggressive return triggers," while providing a visible "no personalisation" switch and one click explanations of "why am I seeing this?" \parencites[221]{2977039}[223]{2977039}.}\pdfcomment[icon=Note]{The AI Law Model's normative illustrations explicitly prohibit manipulative/addictive mechanics (endless scrolling, autoplay) for child-directed systems and require 'detox mode' defaults and visible 'no personalization' switches (2977039, pp.221--223).} \hllightgreen{These provisions directly translate into this report's working taxonomy: autoplay combined with personalised feeds is classified as a high risk pattern when deployed by default in youth oriented social platforms.}\pdfcomment[icon=Note]{The report's cited model law provisions (2977039) and the engagement/addiction literature (3303310) reasonably support the working taxonomy classifying autoplay plus personalised feeds as high risk in youth platforms.} \hllightgreen{Advocacy and governance oriented reports highlight that default autoplay interacts with broader dark pattern strategies.}\pdfcomment[icon=Note]{Advocacy and governance reports on dark patterns (e.g., 6295677) highlight interactions between autoplay defaults and broader manipulative interface strategies, supporting this claim.} \hllightgreen{Khare et al. document how Indian platforms use interface interference and obstruction to steer users toward platform preferred choices, including pre selected defaults and friction loaded opt out flows \parencites[4]{6295677}.}\pdfcomment[icon=Note]{The Indian-focused analysis (6295677, pp.8--13) documents interface interference and obstruction steering users toward platform-preferred choices, aligning with Khare et al.'s described tactics.} \hldarkgreen{Me{\ss}mer and Degeling treat default auto play and default landing on a personalised feed as analogous interface decisions that tilt attention toward recommendation driven content, and they situate these elements inside scenario based audit designs intended to operationalise DSA Article 25's prohibition of manipulative interfaces \parencites[42]{3195574}.}\pdfcomment[icon=Note]{Meßmer \& Degeling explicitly treat default autoplay and landing on a personalised feed as interface choices that tilt attention toward recommendation-driven content and situate these in audit scenarios tied to DSA Article 25 (3195574, p.42).} \hllightgreen{In this thesis, these patterns are captured in Code A (autoplay and infinite scroll) and Code B (dopamine driven variable rewards and temporal discounting), while Code C flags configurations that breach explicit bans in age sensitive design rules as "prohibited or critical risk" under the report's working classification \parencites[2]{8958037}[2]{3303310}[221]{2977039}.}\pdfcomment[icon=Note]{The thesis' coding mapping (Code A/B/C) is an interpretive classification that reasonably synthesises the technical/ethical concerns discussed in 3303310 and regulatory bans in 2977039, matching the cited sources.} \hllightgreen{A key methodological caveat applies throughout this section.}\pdfcomment[icon=Note]{A methodological caveat about limitations is supported by the audit literature and governance sources (3195574; 3303310) which underline constraints in external assessments.} \hldarkgreen{Internal ranking logic, parameter settings and A/B test results for autoplay and recommendation systems are not accessible here.}\pdfcomment[icon=Note]{The sources note that internal ranking logic, parameters and A/B test results are proprietary and not directly accessible to external auditors (3195574; 3303310), directly supporting this limitation statement.} \hldarkgreen{Assessment relies on secondary UX descriptions, regulatory filings, and empirical studies of usage and harm rather than direct inspection of proprietary algorithms, mirroring the limitation acknowledged in both dark pattern research and AI governance work \parencites[6]{6295677}[2]{3303310}[12]{2977039}.}\pdfcomment[icon=Note]{The claim that assessment relies on secondary UX descriptions, regulatory filings and empirical studies rather than direct inspection is explicitly acknowledged as a limitation in the cited audit and governance work (3195574; 3303310).}

\subsubsection{Notification Design and Return Incentives}
\label{sec:notification_return_incentives} \hllightgreen{Push systems on commercial platforms already illustrate how notification design can operate as a return incentive, even though available sources focus on e commerce rather than named social networks.}\pdfcomment[icon=Note]{Qualitative and review sources discuss push/push-notification tactics in commercial e-commerce contexts as return/urgency incentives (8745065 pp.136--139; 6295677 pp.6,12), so the claim is a reasonable synthesis noting that the literature focuses on e-commerce rather than named social networks.} \hldarkgreen{Khare et al. classify "nagging" as a distinct dark pattern: repeated prompts that urge users to take a specific action after refusal, frequently implemented through app notifications.}\pdfcomment[icon=Note]{The typology in 6295677 explicitly lists 'Nagging' as repeated prompts after refusal commonly implemented via app notifications and links it to app notifications (6295677 p.8, Table 1).} \hldarkgreen{These prompts generate a "wear down effect" and decision fatigue that steer users toward acceptance of platform preferred choices, for example continuing a subscription or sharing more data, rather than a once off, neutral reminder.}\pdfcomment[icon=Note]{6295677 labels 'Nagging' with a 'wear down effect' and links it to decision fatigue, and 8745065 describes how urgency/scarcity cues and friction reduction compress deliberation, steering acceptance (6295677 p.8; 8745065 pp.138--139).} \hldarkgreen{In the same typology, "roach motel," "obstruction," and "subscription traps" describe journeys where it is easy to enter a revenue generating state and intentionally difficult to exit, with $47\%$ of surveyed users reporting difficulty cancelling services.}\pdfcomment[icon=Note]{6295677's typology includes 'Roach Motel', 'Obstruction', and 'Subscription Traps' describing easy entry/difficult exit and reports that 47\% of users reported difficulty cancelling services (6295677 p.8, Table 1).} \hllightgreen{Notifications that highlight renewal deadlines or "missed benefits" inside such obstructive flows function as practical return triggers that exploit inertia and status quo bias \parencites[8]{6295677}.}\pdfcomment[icon=Note]{8745065 describes loss-framed renewal/urgency cues and 6295677 links subscription traps to inertia/status-quo bias, so synthesising these as notifications that exploit inertia/status quo bias as return triggers is supported by the sources (8745065 pp.138--139; 6295677 p.8).} \hldarkgreen{Urgency marketing research describes how time based notifications are deliberately scheduled around perceived windows of lowered self control.}\pdfcomment[icon=Note]{8745065 documents time-sensitive scarcity cues and executives' discussion of scheduling messages around temporal contexts perceived to be commercially advantageous, which supports the claim about time-based notifications (8745065 pp.136--139).} \hldarkgreen{Bansal's interviews with executives indicate that late night targeting and high frequency mobile notifications are seen as commercially advantageous because managers assume that consumers are more responsive or less deliberative at particular hours.}\pdfcomment[icon=Note]{8745065 reports interview evidence that executives described late-night targeting and high-frequency mobile notifications as commercially advantageous based on perceived consumer responsiveness at particular hours (8745065 p.138--139).} \hldarkgreen{While this study does not experimentally measure cognitive depletion, the perception itself is analytically significant: assumptions about temporal vulnerability become operational inputs into messaging systems that push expiring discounts or "last chance" alerts at times when refusal is expected to be less likely \parencites[138]{8745065}[8745065]{8745065}.}\pdfcomment[icon=Note]{8745065 explicitly notes the study did not experimentally measure cognitive depletion but argues the executives' perception of temporal vulnerability is analytically significant (8745065 p.138--139), matching the sentence (and cites 8745065).} \hldarkgreen{In practice, these notifications both bring users back into the app and compress deliberation windows, especially when combined with one click checkout and already stored payment data.}\pdfcomment[icon=Note]{8745065 articulates how notifications bring users back and, combined with friction-reduction such as one-click checkout and stored payment data, compress deliberation windows (8745065 pp.138--139).} \hldarkgreen{Regulatory responses increasingly treat such designs as matters of platform accountability rather than individual self discipline.}\pdfcomment[icon=Note]{Sources argue for moving regulatory focus toward platform accountability and preventive design (6295677 pp.6,12; 2977039 pp.114--115), supporting the statement about regulatory responses reframing the issue as platform accountability.} \hldarkgreen{Khare et al. document that India's Central Consumer Protection Authority issued an advisory under section 18(1) of the Consumer Protection Act mandating internal audits for dark pattern detection and explicitly referencing rules that require clear and affirmative consent through explicit user action rather than pre checked defaults or implied agreement.}\pdfcomment[icon=Note]{6295677 documents that the Central Consumer Protection Authority issued an advisory under section 18(1) mandating internal audits for dark pattern detection and references rules requiring clear affirmative consent rather than pre-checked defaults (6295677 p.12).} \hldarkgreen{The same guidance points toward a shift from complaint driven enforcement to preventive design regulation that targets "sludge architecture at its source," including mandatory UX design audits and algorithmic transparency requirements for dominant platforms \parencites[6]{6295677}[12]{6295677}.}\pdfcomment[icon=Note]{6295677 explicitly describes a shift toward prevention via bright-line prohibitions, mandatory UX design audits, algorithmic transparency, and targeting 'sludge architecture at its source' (6295677 p.12).} \hllightgreen{Within this thesis working taxonomy, high frequency, behaviourally targeted notifications in obstructive cancellation or consent flows are therefore coded as high risk manipulative return incentives.}\pdfcomment[icon=Note]{Given the documented harms from high-frequency, behaviourally targeted notifications in obstructive flows in the cited literature (8745065 pp.136--139; 6295677 p.8), coding these as high-risk in a working taxonomy is a defensible synthesis.} \hldarkgreen{Table~\ref{tab:notif_patterns} synthesises notification related patterns that function as return incentives in the sources. \begin{table}[h] \centering \caption{Notification related patterns functioning as return incentives} \label{tab:notif_patterns} \begin{tabular}{p{3.2cm}p{5.0cm}p{4.5cm}} \hline Pattern & Mechanism & Documented effects \\ \hline Nagging notifications & Repeated prompts after refusal, often via app alerts & Wear down effect, decision fatigue, coerced consent \parencites[8]{6295677} \\ Late night urgency messages & Time based push for expiring discounts or offers & Executives perceive higher responsiveness in "low control" windows \parencites[138]{8745065} \\ Subscription renewal prompts inside obstructive flows & Notifications point to renewal while exit paths remain complex & Forced continuation of paid services, recurring unwanted payments \parencites[8]{6295677} \\ "Dark pattern free" self declarations & Internal audits and voluntary declarations without effective external verification & $81\%$ of platforms self declared dark pattern free despite persistent complaints, 28\% resolution rate \parencites[12]{6295677} \\ \hline \end{tabular} \end{table} These patterns interact directly with the multi layer behavioural architecture Bansal synthesises for algorithmic advertising.}\pdfcomment[icon=Note]{6295677 provides the statistics cited (81\% of platforms self-declared dark-pattern-free and a 28\% resolution rate) and the table items (nagging, late-night urgency, subscription renewal prompts) are present across the sources (6295677 p.12; 6295677 p.8; 8745065 pp.136--139).} \hldarkgreen{A "cognitive trigger layer" uses scarcity cues and loss framed notifications to increase the perceived importance of a decision, an "algorithmic personalization layer" tailors timing and content based on behavioural histories, and a "temporal and friction reduction layer" decreases steps to purchase while amplifying urgency.}\pdfcomment[icon=Note]{8745065 sets out the multi-layer behavioral architecture with a 'Cognitive Trigger Layer', 'Algorithmic Personalization Layer', and 'Temporal and Friction-Reduction Layer' in the theoretical synthesis (8745065 p.139).} \hldarkgreen{Within this stack, notifications are the delivery channel that binds layers together: they inject loss framed, personalised cues into the user's off platform environment and then route returning users into low friction, high urgency flows.}\pdfcomment[icon=Note]{8745065 describes notifications as channels that inject loss-framed, personalized cues and then route returning users into low-friction flows, directly supporting this claim (8745065 pp.138--139).} \hldarkgreen{Exploitative impulse buying is therefore better conceptualised as an architectural phenomenon spanning notification design, not as isolated pop up messages \parencites[139]{8745065}.}\pdfcomment[icon=Note]{The conclusion of 8745065 argues that exploitative impulse buying is an architectural phenomenon spanning multiple components (notifications, personalization, friction reduction), not isolated messages (8745065 p.139).} \hllightgreen{The same architectural logic has clear implications for social recommender systems even though direct platform telemetry is unavailable in this study.}\pdfcomment[icon=Note]{While direct platform telemetry is absent, sources on recommender/opinion-forming architectures and their behavioural implications (3303310; 3195574; 4905459) make the extension of the architectural logic to social recommender systems a reasonable inference.} \hllightgreen{When infinite scroll, auto play and highly personalised feeds are already under DSA investigation as "systematic manipulative design" on TikTok, adding high frequency, behaviourally tuned notifications that invite re entry into these feeds would align closely with the dark patterns documented in e commerce \parencites[2]{8958037}[8]{6295677}.}\pdfcomment[icon=Note]{8958037 reports that the DSA has investigated TikTok for 'systematic manipulative design' features like infinite scroll, autoplay and push notifications, and 6295677/8745065 document analogous dark patterns in e-commerce, so combining these points is a defensible synthesis (8958037 p.2; 6295677 p.8; 8745065 pp.136--139).} \hllightgreen{Within this report's working taxonomy, such notification strategies are therefore classified as Code Category C in two cases: where they are embedded in obstructive exit flows or subscription traps, and where they combine temporal vulnerability assumptions with high urgency messaging and compressed deliberation, given the documented harms in adjacent retail contexts \parencites[8]{6295677}[138]{8745065}.}\pdfcomment[icon=Note]{Classifying such notification strategies as high-risk in the report's taxonomy is an internal coding choice but is reasonably grounded in documented harms from adjacent retail contexts in 8745065 and typology evidence in 6295677 (8745065 pp.136--139; 6295677 p.8).}

\subsection{Psychological Trigger Mapping by Platform}

\subsubsection{Variable Rewards and Intermittent Reinforcement}
\label{sec:variable_rewards_intermittent} \hldarkgreen{Short form social platforms make intermittent reinforcement a default interaction pattern rather than an occasional design choice.}\pdfcomment[icon=Note]{Multiple sources explicitly describe short-form social platforms as implementing variable reward / intermittent reinforcement as default interaction patterns (see Dan Li 8958037 p.2; 3303310 p.2).} \hldarkgreen{Li describes how neural network based short video recommenders, combined with 15--60 second clips and swipe based feeds, run on variable reward schedules that trigger dopamine release during the moment of expectancy before the next video appears, not only when a clip is actually enjoyed \parencites[2]{8958037}[3]{8958037}.}\pdfcomment[icon=Note]{Dan Li (8958037 p.2) describes neural-network short-video recommenders, micro-duration clips and variable reward schedules triggering dopamine release in expectancy, matching the claim.} \hllightgreen{Users execute an effortless downward swipe and face an uncertain payoff: the next item may be highly engaging or dull.}\pdfcomment[icon=Note]{Sources describe swipe-based feeds and unpredictability of next clip quality (e.g., 8958037 p.2; 3303310 p.2), so the characterization is a reasonable synthesis of those descriptions.} \hldarkgreen{This probabilistic structure mirrors slot machine style variable rewards and, over repeated cycles, propels users into an "auto pilot mode" of persistent scrolling that is strongly associated with diminished academic concentration, weakened social skills, and heightened anxiety and depression among college students \parencites[2]{8958037}[6]{8958037}.}\pdfcomment[icon=Note]{Dan Li links variable/slot-machine-like reward schedules to 'auto-pilot' persistent scrolling and documents harms among college students (diminished concentration, weakened social skills, anxiety/depression) (8958037 p.2).} \hldarkgreen{Hu's work on "algorithm driven engagement worship" generalises this pattern beyond a single app.}\pdfcomment[icon=Note]{3303310 explicitly frames the concept of 'algorithm-driven engagement worship' and extends the pattern beyond a single app (3303310 p.2).} \hldarkgreen{Infinite scrolling interfaces and variable reward schedules are identified as core features of platforms "like TikTok," and empirical validations show that users of algorithmically curated feeds exhibit $22\%$ higher rates of compulsive checking behaviour than users of non algorithmic interfaces \parencites[2]{3303310}.}\pdfcomment[icon=Note]{The identification of infinite scroll and variable reward schedules as core features of platforms like TikTok and the empirical 22\% figure for higher compulsive checking in algorithmic feeds are reported in the sources (3303310 p.2--3; 8958037 p.2).} \hldarkgreen{Longitudinal evidence from the Social Media and Wellbeing Study, cited in the same analysis, indicates that hyper personalised feeds are associated with a $23\%$ higher incidence of depressive symptoms compared to chronological content, and that algorithmically amplified "doomscrolling" correlates with reduced prefrontal cortex activity and reduced gray matter density in regions governing impulse control \parencites[20]{3303310}.}\pdfcomment[icon=Note]{3303310 reports longitudinal Social Media and Wellbeing Study findings of a 23\% higher incidence of depressive symptoms with hyper-personalized feeds and links 'doomscrolling' to reduced prefrontal activity and gray-matter changes (3303310 p.3).} \hllightgreen{Variable reinforcement is therefore not only a behavioural phenomenon but a mechanism tied to neurocognitive change.}\pdfcomment[icon=Note]{Sources link variable reinforcement to measurable neurocognitive effects (e.g., 3303310 p.2--3; 8958037 p.2), so concluding it is both behavioural and mechanistic is a supported synthesis.} \hlorange{Table~\ref{tab:vr_patterns} summarises how the report's working Code B categories capture these documented triggers across platform types. \begin{table}[h] \centering \caption{Variable reward patterns and intermittent reinforcement across platforms (working Code B mapping)} \label{tab:vr_patterns} \begin{tabular}{p{3.2cm}p{4.8cm}p{4.8cm}} \hline Platform type & Variable reward mechanic & Documented effects \\ \hline Short video feeds & Swipe to next clip with unpredictable quality; infinite scroll; default auto play & "Auto pilot" scrolling, behavioural addiction, impaired study performance in students \parencites[2]{8958037}[6]{8958037} \\ Algorithmic social feeds & Engagement ranking that privileges emotionally charged content; irregular delivery of highly salient posts & $22\%$ higher compulsive checking; 23\% higher depressive symptoms; reduced prefrontal activity \parencites[2]{3303310}[20]{3303310} \\ E commerce flows & Countdown timers and "only two left" alerts that sometimes expire and sometimes do not & Conversion uplift reported by executives; decision fatigue and regret in consumers \parencites[136]{8745065}[8]{6295677} \\ \hline \end{tabular} \end{table} Urgency marketing in retail environments shows the same reinforcement logic using different surface cues.}\pdfcomment[icon=Note]{While parts of the table (short video and algorithmic feed rows) are supported by 8958037 and 3303310, the e-commerce claims and specific citation [8]\{6295677\} are not present among the provided sources, so the composite table is not fully supported.} \hlorange{Bansal's interviews with executives describe countdown timers, "only two left" stock flags, and expiring discounts as core elements of campaigns that transform routine purchases into time pressured decisions \parencites[136]{8745065}.}\pdfcomment[icon=Note]{No provided source explicitly names 'Bansal' or presents executive interviews describing the specific countdown/timer practices cited, so this exact claim lacks support in the supplied materials.} \hlorange{Khare et al. classify "false urgency / scarcity" as a dark pattern where countdowns or stock messages are misleading or artificially generated, producing intermittent confirmation that urgency is real and at other times revealing that nothing happens when the timer expires \parencites[7]{6295677}.}\pdfcomment[icon=Note]{The specific attribution to 'Khare et al.' and the detailed classification of 'false urgency / scarcity' with the described behaviors are not present in the provided sources, so this claim is unsupported here.} \hlorange{This inconsistency is itself a variable reward schedule: occasionally, failing to act leads to a missed discount; in other cases, the offer silently renews.}\pdfcomment[icon=Note]{The assertion that inconsistency of urgency messages functions as a variable reward schedule is plausible but no provided source explicitly analyzes countdown inconsistency in this way, so it is unsupported by the supplied citations.} \hlorange{Users learn that "sometimes you lose if you wait," which reinforces impulsive action without providing stable informational value.}\pdfcomment[icon=Note]{The behavioral learning statement ('sometimes you lose if you wait') is a reasonable psychological claim but is not directly evidenced in the provided sources for the specific context, so it lacks direct support here.} \hllightgreen{Several sources tie these mechanics to regulatory and ethical boundaries.}\pdfcomment[icon=Note]{Multiple supplied sources discuss regulatory and ethical concerns around these mechanics (e.g., 8958037 p.2; 3303310 p.2; 2977039 p.213), supporting the general statement.} \hldarkgreen{The European Commission's DSA proceeding against TikTok explicitly cites infinite scrolling, auto play and push notifications as possible "systematic manipulative design" that impairs free and informed decision making \parencites[2]{8958037}.}\pdfcomment[icon=Note]{Dan Li explicitly notes the EU DSA proceeding into TikTok and cites infinite scrolling, auto-play and push notifications as possible 'systematic manipulative design' (8958037 p.2).} \hldarkgreen{Kostenko's AI Law Model prohibits AI systems that manipulate behaviour, cause digital addiction, emotional exhaustion or loss of trust and frames an inalienable right to a digital environment free from excessive information flows that create chronic anxiety or distraction without explicit consent \parencites[61--62]{2977039}[119]{2977039}.}\pdfcomment[icon=Note]{The AI Law Model text (2977039 p.213, p.233--234) contains prohibitions and rights framed around manipulation, addiction, emotional exhaustion and rights to a non-excessive digital environment, matching this claim.} \hllightgreen{Variable reward loops that systematically produce compulsive checking, depressive symptoms, and structural changes in impulse control sit in direct tension with these duties.}\pdfcomment[icon=Note]{Given the documented harms (e.g., compulsive checking and depressive symptoms in 3303310 and 8958037) and the duties in 2977039, asserting that variable reward loops conflict with those duties is a supported synthesis (3303310 p.3; 2977039 p.213).} \hllightgreen{Audit and governance frameworks offer ways to treat intermittent reinforcement as an assessable risk rather than an abstract concern.}\pdfcomment[icon=Note]{Governance and audit frameworks that operationalize intermittent reinforcement as assessable risk are described in the AIN and governance literature (2742041 p.2; 8367105 p.25), supporting this claim.} \hldarkgreen{Hu's study constructs an interdisciplinary assessment matrix that causally maps neurocognitive metrics to algorithmic feature vectors and proposes a regulatory toolkit including a dynamic transparency index and an early warning system for addictive interactions, with standardized metrics such as the Algorithmic Impact Assessment Scale (AIAS 5) to quantify emotional volatility, sleep disruption, social cohesion, self esteem variance and compulsive usage patterns \parencites[19]{3303310}[20]{3303310}.}\pdfcomment[icon=Note]{3303310 describes an interdisciplinary assessment matrix mapping neurocognitive metrics to algorithmic feature vectors and proposes a regulatory toolkit including a transparency index and AIAS-style metrics (3303310 p.2).} \hldarkgreen{Leon's AIN framework similarly positions organisations as responsible for continuous impact evaluations, logging recommendations and human decisions, and monitoring emerging risks through periodic audits \parencites[2742041]{2742041}.}\pdfcomment[icon=Note]{The AIN framework (2742041 p.8, 10.2.3--10.2.4) emphasizes organizational responsibilities for continuous impact evaluation, logging recommendations and human decisions, and periodic reviews/audits.} \hllightgreen{In revenue audit contexts, Chitturi highlights feedback mechanisms and clearly documented splits between automated detection and human analysis as conditions for dependable oversight \parencites[6539584]{6539584}.}\pdfcomment[icon=Note]{Audit literature supplied discusses feedback mechanisms and documenting splits between automated detection and human review as conditions for dependable oversight (6539584 p.5), supporting this characterization (attributed here to revenue audit contexts).} \hllightgreen{Within this thesis, these materials justify treating variable rewards and intermittent reinforcement as explicit Code B psychological triggers and as inputs to Code C legal risk classification.}\pdfcomment[icon=Note]{Given the body of empirical and normative materials cited (e.g., 8958037; 3303310; 2977039), treating variable rewards as Code B psychological triggers and inputs to legal risk classification is a reasonable, supported justification within a thesis (synthesis across sources).} \hllightgreen{Short form video feeds and engagement optimised social streams that combine infinite scroll, default auto play and irregular high arousal content are therefore coded as high risk configurations for algorithmic addiction, particularly when targeted at user groups that survey evidence identifies as confident in basic digital skills yet poor at recognising hidden manipulations such as dark patterns and algorithmic influences \parencites[10]{9360054}[2]{3303310}[2]{8958037}.}\pdfcomment[icon=Note]{Sources document that short-form feeds combining infinite scroll, autoplay and high-arousal content are high-risk for addiction and note that some user surveys show confidence in basic digital skills but low ability to detect manipulations (8958037 p.2; 9360054 p.31; 3303310 p.2--3), making this a supported synthesis.} \hldarkgreen{Recommender systems themselves are engineered to predict and sustain engagement at scale by exploiting large, sparse interaction datasets and probabilistic inference; their capacity to generate rapid, personalized item sequences is enabled by collaborative filtering and network-based models that operate efficiently on millions of observations \parencites[3]{7422270}}\pdfcomment[icon=Note]{Technical literature on recommender systems (7422270 p.2--3, p.22) describes systems engineered to predict and sustain engagement using large, sparse interaction datasets and collaborative-filtering/network-based models scalable to millions of observations.}

\subsubsection{Social Validation, Status Cues and Fear of Missing Out}
\label{sec:social_validation_status_fomo} \hldarkgreen{Status cues and social validation signals function as high salience triggers across social and commercial platforms, directly shaping engagement and feeding fear of missing out.}\pdfcomment[icon=Note]{Sources discuss status cues, social proof and FOMO as engagement drivers and high‑salience triggers (7160486 p.1--2; 3303310 p.2).} \hldarkgreen{On Chinese platforms such as Douyin and Xiaohongshu, visible metrics like likes, follower counts and reviews are central to perceived credibility.}\pdfcomment[icon=Note]{Source 7160486 explicitly states that on Douyin and Xiaohongshu visible metrics like likes, followers and reviews are central to credibility (7160486 p.1).} \hlorange{Wang reports that these cues are easily distorted: fake accounts, bots and paid engagement inflate numbers and fabricate an illusion of popularity.}\pdfcomment[icon=Note]{The content (fake accounts, bots, paid engagement inflating metrics) is supported by sources (7160486 p.1), but the specific attribution to "Wang" is not present in the provided sources.} \hldarkgreen{A 2023 China Consumers Association report cited in the same work found that nearly $30\%$ of online shoppers had encountered misleading or falsified reviews in the preceding year.}\pdfcomment[icon=Note]{Source 7160486 cites a 2023 China Consumers Association report finding nearly 30\% of online shoppers encountered misleading/falsified reviews (7160486 p.1).} \hldarkgreen{What should be an aggregated signal of collective trust is turned into a tool of manipulation, degrading confidence in digital commerce once users can no longer distinguish genuine experiences from fabricated endorsements.}\pdfcomment[icon=Note]{Source 7160486 argues that social proof can be manipulated and that this erosion of authenticity weakens confidence in digital commerce (7160486 p.1).} \hldarkgreen{Influencer marketing exploits related dynamics.}\pdfcomment[icon=Note]{Source 7160486 discusses how influencer marketing exploits social validation dynamics and emotional credibility (7160486 p.1--2).} \hldarkgreen{Influencers and KOLs in China wield "immense sway" over consumer decisions, yet their endorsements often blur the line between personal opinion and paid promotion.}\pdfcomment[icon=Note]{Source 7160486 states influencers/KOLs in China hold substantial sway and that endorsements often blur personal opinion and paid promotion (7160486 p.1--2).} \hldarkgreen{Regulatory responses, such as SAMR rules requiring clear sponsorship disclosure and prohibition of false claims, have tightened formal expectations, but enforcement remains uneven and many brands continue to trade on the emotional credibility of influencers.}\pdfcomment[icon=Note]{Source 7160486 documents SAMR guidance on disclosure and false claims (2022) and notes uneven enforcement with brands still leveraging influencer credibility (7160486 p.1--2).} \hllightgreen{Followers interpret endorsements as socially grounded recommendations, which amplifies the persuasive power of likes and comments beneath sponsored posts and embeds status seeking into routine platform use.}\pdfcomment[icon=Note]{Sources indicate followers interpret endorsements as social recommendations and that this amplifies persuasive power of likes/comments (7160486 p.1--2; 4245543 p.4), which supports the synthesis.} \hldarkgreen{Algorithmic curation intensifies these validation loops.}\pdfcomment[icon=Note]{Sources describe algorithmic curation amplifying engagement loops and social validation (3303310 p.1--3; 7160486 p.2).} \hlorange{Wang notes that engagement focused algorithms tend to amplify already popular content, creating cycles where "popularity begets more popularity." This herd reinforcement pressures users to conform to trends, independent of their actual needs or preferences.}\pdfcomment[icon=Note]{The phenomenon (algorithms amplifying popular content and pressuring conformity) is supported by sources (3303310 p.2; 7160486 p.2), but the specific attribution to "Wang" is not found in the provided materials.} \hldarkgreen{On Douyin, viral challenges and live streamed sales events create collective buying moments, in which consumers experience social pressure to participate so as not to be left behind.}\pdfcomment[icon=Note]{Source 7160486 discusses Douyin's viral challenges and live‑streamed sales creating collective buying moments and social pressure to participate (7160486 p.2).} \hldarkgreen{The resulting fear of missing out drives impulsive spending and emotional fatigue and turns feeds into environments where psychological vulnerabilities around belonging and status are repeatedly activated \parencites[13]{7160486}.}\pdfcomment[icon=Note]{Source 7160486 links FOMO on platforms to impulsive spending, emotional fatigue and repeated activation of belonging/status vulnerabilities (7160486 p.2).} \hllightgreen{These observations dovetail with broader analyses of engagement worship in social media.}\pdfcomment[icon=Note]{This is a synthesis supported by broader analyses of engagement‑centred social media research in the provided sources (3303310 p.1--3; 7160486 p.2).} \hlorange{Hu describes how collaborative filtering algorithms systematically amplify high arousal content such as outrage or sensationalism, and how personalised feeds increase compulsive checking and depressive symptoms relative to chronological baselines.}\pdfcomment[icon=Note]{The mechanisms described (collaborative filtering amplifying high‑arousal content and increased compulsive checking/depressive symptoms) are present in sources (3303310 p.2--3), but the named attribution to "Hu" is not found among the provided sources.} \hldarkgreen{Status cues are part of this mechanism: content with strong emotional or identity signalling often garners more reactions, which in turn drives further algorithmic promotion.}\pdfcomment[icon=Note]{Sources state that status cues and emotionally/identity‑salient content garner more reactions and thereby receive further algorithmic promotion (3303310 p.2--3; 7160486 p.2).} \hldarkgreen{Empirical work cited in Hu's study shows that Instagram's "Explore" page amplified body negative imagery by $23\%$ compared with chronological feeds, correlating with statistically significant declines in self esteem among adolescent women of color.}\pdfcomment[icon=Note]{Source 3303310 reports Instagram's Explore page amplified body‑negative imagery by 23\% versus chronological feeds, linked to declines in self‑esteem (3303310 p.3).} \hldarkgreen{This pattern links social validation architecture directly to harm in status sensitive domains such as appearance and identity \parencites[2]{3303310}[3]{3303310}.}\pdfcomment[icon=Note]{Source 3303310 connects algorithmic/social validation architectures to harms in status‑sensitive domains such as appearance/identity (3303310 p.2--3).} \hllightgreen{Table~\ref{tab:status_fomo} summarises how social validation and FOMO related triggers appear across the examined settings in this report's working coding matrix. \begin{table}[h] \centering \caption{Working Code B mapping: social validation, status cues and FOMO} \label{tab:status_fomo} \begin{tabular}{p{3cm}p{4.5cm}p{4.5cm}} \hline Pattern & Mechanism & Documented effects \\ \hline Inflated metrics & Fake likes, followers, paid engagement & Illusion of popularity, weakened trust in reviews and ratings \\ Undisclosed sponsorship & Influencer posts that hide commercial nature & Exploitation of emotional credibility, confusion between advertising and opinion \\ Algorithmic herd reinforcement & Preferential promotion of already popular content & Pressure to conform to trends, FOMO driven participation in viral challenges and live sales \parencites[13]{7160486} \\ Body and identity sensitive feeds & Prioritisation of high arousal or body negative imagery & Increased body dissatisfaction and lower self esteem among specific groups \parencites[3]{3303310} \\ \hline \end{tabular} \end{table} Status seeking and social proof also interact with ethical consumption and financial decision making.}\pdfcomment[icon=Note]{The table description summarizes patterns (inflated metrics, undisclosed sponsorship, algorithmic herd reinforcement, body‑sensitive feeds) that are documented across the cited sources (7160486 p.1--2; 3303310 p.2--3).} \hlorange{Kishore et al. show that university students' decisions in digital marketplaces are dominated by price, quality, reviews, urgency and perceived usefulness, while ethical considerations appear intermittently and often as secondary cues.}\pdfcomment[icon=Note]{The content (students' decision factors) matches findings in Source 4245543 (4245543 p.2--4), but the specific citation to "Kishore et al." is not present in the provided sources.} \hldarkgreen{Trust in platform claims and peer reviews becomes a structural bridge between values and action: unverifiable or strategically framed social proof triggers skepticism and disengagement from sustainability considerations, exacerbating the intention behaviour gap.}\pdfcomment[icon=Note]{Source 4245543 frames platform trust as a bridge between values and action and shows unverifiable social proof reduces engagement with sustainability (4245543 p.3--4, p.10).} \hldarkgreen{Emotional responses such as guilt, regret and confusion operate as self regulation signals in these contexts, indicating that distorted validation cues affect not only whether people buy but how they experience their own identity as consumers \parencites[2]{4245543}[4]{4245543}.}\pdfcomment[icon=Note]{Source 4245543 explicitly discusses emotions (guilt, regret, confusion) as moral self‑regulation signals and links distorted validation to consumer experience (4245543 p.2, p.10).} \hllightgreen{Research on gamification and interaction features extends this picture into participatory environments.}\pdfcomment[icon=Note]{The role of gamification and interaction features in participatory environments is covered in the literature (7154741 p.3; 5645864 p.3), supporting this bridging statement.} \hldarkgreen{Yin reports that TikTok hashtag challenges, framed as community competitions with implicit status rewards, make challengers $62\%$ more likely to log in daily than non challengers.}\pdfcomment[icon=Note]{Source 7154741 reports research showing TikTok challenges made challengers 62\% more likely to log in daily (7154741 p.3).} \hllightgreen{Such mechanisms transform users into content co creators while simultaneously tying routine engagement to visible participation in status contests.}\pdfcomment[icon=Note]{Source 7154741 describes active participation turning audiences into co‑creators and tying engagement to visible participation, which supports this synthesis (7154741 p.2--3).} \hllightgreen{Although framed positively as community and belonging, these features increase retention by embedding FOMO into user journeys: missing a challenge or failing to contribute content risks social invisibility within the niche community \parencites[3]{7154741}.}\pdfcomment[icon=Note]{Sources link gamified/community features to increased retention and potential FOMO/social anxiety when users miss challenges (7154741 p.3; 7160486 p.2), supporting this interpretation.} \hllightgreen{Across these strands, social validation, status cues and fear of missing out emerge as central psychological targets for engagement optimised design.}\pdfcomment[icon=Note]{Multiple sources converge on social validation, status cues and FOMO as key psychological targets of engagement‑optimised design (3303310 p.1--3; 7160486 p.1--2; 4245543 p.2--4).} \hllightgreen{This comparative audit relies entirely on secondary UX analyses, empirical studies, and regulatory documents rather than internal telemetry, but the converging evidence from Chinese e commerce, global social media research and student behaviour studies is sufficient to justify treating distorted social proof, herd reinforcing algorithms and FOMO amplifying mechanics as high priority risk items in this report's working Ethical AI Audit Toolkit \parencites[13]{7160486}[2]{3303310}[2]{4245543}[3]{7154741}.}\pdfcomment[icon=Note]{The methods claim (reliance on secondary analyses) and the judgment that converging evidence justifies flagging these risks is supported by the cited secondary studies (3303310 p.1--3; 4245543 p.2--4; 7154741 p.3).}

\subsubsection{Emotional High-Arousal and Outrage Content}
\label{sec:emotional_high_arousal_outrage} \hldarkgreen{Engagement metrics on major platforms are systematically tuned toward emotionally intense material.}\pdfcomment[icon=Note]{Multiple sources document that engagement metrics favour emotionally intense material (7154741 p.11; 3303310 p.2--3; 6502223 p.6--7).} \hldarkgreen{Chitra and co authors synthesise experimental work showing that moral emotional language in tweets raises retweet probability by about $20\%$ per additional moral emotional word and that content evoking awe, anger, or anxiety is significantly more likely to be shared than content associated with sadness or contentment.}\pdfcomment[icon=Note]{The claim matches cited empirical findings: Brady et al. (moral-emotional words ≈ +20\% retweet) and Berger \& Milkman (awe/anger/anxiety shared more than sadness) as summarized in 6502223 p.7.} \hldarkgreen{Since deep learning ranking stages are trained on engagement signals such as click through rate, watch time, completion, shares and comments, this creates a selection pressure in which outrage maximizing political content outcompetes accurate, nuanced reporting because the loss function contains no term for epistemic quality or civic value \parencites[7]{6502223}.}\pdfcomment[icon=Note]{Technical accounts show deep-learning ranking stages are trained on engagement signals (clicks, watch time, shares, comments) and lack explicit epistemic/civic terms, producing selection pressure (6502223 p.3--7; 3303310 p.2).} \hlorange{Hu's analysis of social media feeds shows how this dynamic plays out in curated streams.}\pdfcomment[icon=Note]{No provided source attributes an analysis of curated streams to an author named Hu; while sources discuss feed dynamics (3303310; 7154741), the specific attribution to 'Hu' is unsupported by the supplied documents.} \hldarkgreen{Algorithmic prioritisation systematically surfaces high arousal items, with Instagram's "Explore" page amplifying body negative imagery by $23\%$ relative to chronological feeds and correlating with a statistically significant decline in self esteem among adolescent women of color ($\beta = -0.41$, $p < 0.05$) \parencites[3]{3303310}.}\pdfcomment[icon=Note]{The Instagram Explore finding and associated self-esteem decline (23\% amplification; β = −0.41, p < 0.05 among adolescent women of color) are reported explicitly in 3303310 p.3.} \hldarkgreen{Sentiment analysis in the same work finds that algorithmically prioritised content exhibits $31\%$ higher negative affect than baseline, aligning with Alter's model of dopamine driven feedback loops where intermittent reinforcement of emotionally charged content conditions users neurologically \parencites[3]{3303310}[2]{3303310}.}\pdfcomment[icon=Note]{3303310 p.3 reports sentiment analysis showing \textasciitilde{}31\% higher negative affect in algorithmically prioritised content and links this to Alter's dopamine-driven feedback loop model.} \hllightgreen{High arousal content interacts with viral mechanics to reshape social norms.}\pdfcomment[icon=Note]{Sources document that high-arousal content interacts with virality to influence social norms and discourse (7154741 p.11--12; 6502223 p.6--7), supporting this synthesis claim.} \hlorange{Yin documents that viral posts and hashtags can reconfigure cultural spaces in real time, citing Black Lives Matter as an example where online virality catalysed global protest waves, but also noting that quick fire phenomena like the "Harlem Shake" displace sustained discussion and contribute to information fatigue.}\pdfcomment[icon=Note]{The substantive examples (Black Lives Matter, Harlem Shake) and discussion of viral reconfiguration are present (7154741 p.11--12), but the specific attribution to an author named 'Yin' is not found in the provided sources.} \hldarkgreen{Viral misinformation during the early COVID 19 pandemic illustrates the darker side: phishing stories, misleading headlines and doctored videos about cures and conspiracy theories spread faster than verification or retraction, aided by algorithms that promote "racy" posts with high interaction, thereby embedding users in echo chambers and misrepresenting public opinion.}\pdfcomment[icon=Note]{7154741 p.11--12 documents early COVID-19 viral misinformation (phishing, misleading headlines, doctored videos) spreading faster than verification and being amplified by engagement-driven algorithms, embedding users in echo chambers.} \hldarkgreen{Creators and viewers both experience psychological strain in this environment: content producers risk harassment and burnout, while viewers report emotional distress when repeatedly exposed to polarising or harmful material \parencites[2]{7154741}.}\pdfcomment[icon=Note]{Empirical accounts report harassment and burnout among creators and emotional fatigue/distress among viewers exposed to polarising content (7154741 p.11--12; 3195574 p.43).} \hllightgreen{Table~\ref{tab:high_arousal_mapping} aligns these observations with this report's working coding matrix. \begin{table}[h] \centering \caption{High arousal and outrage content: working Code B mapping} \label{tab:high_arousal_mapping} \begin{tabular}{p{3.2cm}p{4.6cm}p{4.6cm}} \hline Coding slot & Mechanism & Empirical or conceptual anchor \\ \hline B1: Moral emotional amplification & Moral emotional wording boosts sharing, driving outrage selection & Moral emotional words increase retweets by $\approx 20\%$ per word \parencites[7]{6502223} \\ B2: Negative affect bias & Ranking favours high negative affect content & Algorithmic feeds show $31\%$ higher negative affect, $23\%$ more body negative imagery \parencites[3]{3303310} \\ B3: Viral echo chamber reinforcement & Viral, racy posts amplified into cohesive cliques & Viral misinformation and doctored videos deepen echo chambers and division \\ B4: Emotional fatigue and distress & Repeated exposure to polarising material harms creators and viewers & Harassment, burnout, information fatigue and emotional distress reported \parencites[2]{7154741} \\ \hline \end{tabular} \end{table} Technical accounts of recommendation architecture clarify why these emotional biases persist.}\pdfcomment[icon=Note]{The table's mapped observations draw on multiple cited findings (moral-emotional +20\% retweets, 31\% higher negative affect, 23\% body-negative amplification, viral echo chambers, harassment/burnout) present across 6502223 p.7 and 3303310 p.2--3 and 7154741 p.11--12.} \hldarkgreen{Candidate generation and deep learning ranking stages are optimised over behavioural datasets, using multi task objectives that weight click through, completion, watch time, share rate and comment rate by commercial value.}\pdfcomment[icon=Note]{Technical descriptions in 6502223 outline candidate generation and deep-learning ranking optimised on behavioural datasets with multi-task objectives weighting CTR, watch time, shares, comments by commercial value (6502223 p.3--4).} \hldarkgreen{Content that reliably triggers anger, anxiety or awe scores highly on these metrics and is therefore repeatedly selected, while there is no explicit optimisation term for psychological wellbeing.}\pdfcomment[icon=Note]{6502223 explains that content triggering anger/anxiety/awe scores highly on engagement metrics and that there is no explicit optimisation for psychological wellbeing in ranking loss functions (6502223 p.3--7).} \hldarkgreen{Reinforcement learning layers that optimise for return rate or session length intensify this effect: YouTube's "radicalization pipeline" is interpreted as an RL agent discovering that escalating content intensity is an effective strategy for maintaining attention, gradually steering users from moderate to more extreme material even when they would prefer balanced content if asked directly \parencites[7]{6502223}.}\pdfcomment[icon=Note]{6502223 documents RL components and cites Ribeiro on YouTube's 'radicalization pipeline', describing RL agents escalating content intensity to maximize retention (6502223 p.3--4; p.7).} \hllightgreen{Concerns about outrage amplification intersect with broader debates on democratic rights.}\pdfcomment[icon=Note]{The concern about outrage amplification implicating democratic rights is discussed in the digital constitutionalism literature summarized in 6502223 p.11 and related policy sources.} \hldarkgreen{Chitra and colleagues frame algorithmic manipulation of political information as threatening freedom of expression, the right to receive diverse viewpoints and electoral rights, because ranking that systematically suppresses certain perspectives or amplifies disinformation undermines free and fair elections and privacy through behavioural surveillance.}\pdfcomment[icon=Note]{6502223 p.11 explicitly frames algorithmic manipulation as threatening freedom of expression, the right to receive diverse viewpoints, electoral rights, and privacy via behavioural surveillance (6502223 p.11, 'Rights at Stake' section).} \hldarkgreen{Digital constitutionalist proposals call for transparency about ranking principles, audit trails, independent oversight with enforcement powers, due process for users, and non discrimination in the distribution of political information, treating systematic political amplification asymmetries as direct violations of constitutional norms \parencites[11]{6502223}.}\pdfcomment[icon=Note]{Core digital-constitutional prescriptions (transparency about ranking, audit trails, independent oversight with enforcement, due process, non-discrimination) are detailed in 6502223 p.11 (Core Principles) and related recommendations (6502223 p.16).} \hlorange{These dynamics are assessed here under strict methodological limits.}\pdfcomment[icon=Note]{The assertion about this report assessing dynamics under 'strict methodological limits' is an internal methodological claim of the document and is not directly documented or evidenced in the provided external sources.} \hldarkgreen{Internal loss functions, feature weights and reinforcement learning policies of commercial platforms are not observable in this project; analysis relies on published descriptions of architectures, empirical studies of content effects, and qualitative accounts of virality and emotional impact rather than direct inspection of proprietary code or telemetry \parencites[7]{6502223}[3]{3303310}[2]{7154741}.}\pdfcomment[icon=Note]{Sources note that internal loss functions and proprietary RL policies are not observable and that analyses must rely on published descriptions, empirical studies, and qualitative accounts (3303310 p.2; 6502223 p.3; 2977039 on transparency/audit requirements).} \hllightgreen{Within that constraint, evidence indicates that emotional high arousal and outrage content is not an incidental by product of social recommender systems but a predictable outcome of engagement oriented optimisation, and that any Ethical AI Audit Toolkit aligned with EU digital constitutionalism and a Digital Fairness trajectory must therefore treat such content patterns as central diagnostic signals rather than peripheral content moderation issues \parencites[7]{6502223}[11]{6502223}[3]{3303310}.}\pdfcomment[icon=Note]{Available evidence and policy recommendations indicate high-arousal outcomes are predictable from engagement optimisation and recommend treating such patterns as core audit signals (6502223 p.3--4, p.16; 3303310 p.2--3), supporting this synthesis conclusion.}

\section{Design of an Ethical AI Audit Toolkit for Social Recommender Systems}

\subsection{Conceptual Design of the Audit Toolkit}

\subsubsection{Objectives, Users and Governance of the Toolkit}
\label{sec:toolkit_objectives_users_governance} \hldarkgreen{Audit design in this project pursues three linked objectives: translate legal and ethical duties on AI into operational control points, create artefacts that are auditable by external actors, and embed those controls into organisational governance rather than treating them as one off assessments.}\pdfcomment[icon=Note]{8367105 (p.24) frames audit design as converting principles into operational gates, producing auditable artefacts and embedding controls into organisational processes, directly supporting the three objectives.} \hldarkgreen{Catapang's triple gate architecture illustrates the first aim: data, training, deployment and monitoring stages are each subject to Governance, Metric and Eco gates, mapped to AI Act obligations such as data governance, human oversight, transparency and post market monitoring, and to NIST AI RMF functions like Map, Measure and Govern.}\pdfcomment[icon=Note]{8367105 (p.24 and Table 3) explicitly presents a triple-gate (Governance/Metric/Eco) mapping across data, train/deploy/monitor stages and crosswalks to EU AI Act obligations and NIST RMF functions.} \hldarkgreen{Gate outcomes are not reports but blocking or throttling decisions whose logs become audit artefacts, which aligns directly with this thesis intent to move from ethical "principles" toward enforceable checkpoints in recommender pipelines \parencites[24]{8367105}.}\pdfcomment[icon=Note]{8367105 (p.24) states gate outputs decide promotion/rollback/throttling and that their logs become audit artefacts, and frames this as shifting from principles to operational checkpoints.} \hldarkgreen{Kostenko's AI Law Model complements this by specifying mandatory procedures for ethical admissibility, including notification of AI use, explainability of conclusions, revision of data sources, guarantees of human control over critical automated decisions, and counteraction to manipulative practices that harm vulnerable groups \parencites[9]{2977039}.}\pdfcomment[icon=Note]{2977039 (p.8) lists mandatory procedures (notification, explainability, data source revision, human control, counteraction to manipulative practices and protection of vulnerable groups) as part of the AI Law Model.} \hllightgreen{The toolkit is therefore oriented to operationalise these procedures at concrete decision gates: for instance, preventing releases of feed ranking changes that worsen addiction indicators or fairness metrics beyond predefined thresholds.}\pdfcomment[icon=Note]{8367105 (p.24) and related discussion describe operational gates that can block harmful releases and log thresholds; applying this to recommender metrics (addiction/fairness thresholds) is a reasonable synthesis of those controls.} \hllightgreen{Multiple user groups are envisioned.}\pdfcomment[icon=Note]{The document and cited literature delineate multiple stakeholders (internal teams, risk/compliance, external auditors), so the claim that multiple user groups are envisioned is a supported synthesis (e.g., 8367105 p.24; 2977039 p.9).} \hllightgreen{Inside organisations, product teams and engineering leadership need controls that integrate with CI/CD practices.}\pdfcomment[icon=Note]{8367105 (p.24) describes embedding gates into CI/CD; inferring that product teams and engineering leadership require CI/CD-integrated controls is a direct and reasonable synthesis.} \hldarkgreen{Catapang shows that embedding Governance, Metric and Eco gates into CI/CD and runtime serving paths allows each promotion or rollback decision to generate logs suitable for conformity assessment, with Eco gates treated as first class blockers when emissions budgets are exceeded \parencites[24]{8367105}.}\pdfcomment[icon=Note]{8367105 (p.24, Discussion) explicitly reports embedding Governance/Metric/Eco gates into CI/CD/runtime so promotion/rollback generate auditable logs and treats Eco gates as first-class blockers when emissions budgets are exceeded.} \hllightgreen{Translating this pattern to social recommenders implies dashboards and gate checks that product owners can use during experimentation, with pass or fail statuses aligned to DSA obligations on non profiled options and systemic risk mitigation as interpreted in prior sections \parencites[12]{6502223}.}\pdfcomment[icon=Note]{8367105 (p.24) and 6502223 (p.12) describe gate checks and DSA obligations (non-profiled options, systemic risk); extrapolating to dashboards and pass/fail experimentation checks for product owners is a reasonable synthesis.} \hllightgreen{Risk, compliance and audit functions form a second user group.}\pdfcomment[icon=Note]{The literature distinguishes risk/compliance/audit as stakeholder functions (e.g., 5442231 p.15; 2977039 p.9), so identifying them as a user group is a supported inference.} \hldarkgreen{Sebastian anticipates algorithmic impact assessments for high risk credit AI that mirror environmental impact reviews, with standardised fairness metrics, documentation structures and transparency reports detailing model types, data sources, validation methods, demographic performance and governance arrangements \parencites[15]{5442231}.}\pdfcomment[icon=Note]{5442231 (p.15) describes algorithmic impact assessments for high‑risk credit AI analogous to environmental reviews with standardized fairness metrics, documentation and transparency reporting requirements.} \hldarkgreen{Mutawa et al.'s Algorithmic Governance Assessment Matrix similarly evaluates transparency, accountability, explainability, data integrity and normative alignment using coding sheets over institutional documents and performance indicators \parencites[59]{2702055}.}\pdfcomment[icon=Note]{2702055 (p.4) documents an Algorithmic Governance Assessment Matrix that evaluates transparency, accountability, explainability, data integrity and uses coding sheets over institutional documents and indicators.} \hllightgreen{The audit toolkit in this thesis is designed to produce comparable documentation and scoring outputs that these functions can interpret without needing access to source code.}\pdfcomment[icon=Note]{Sources describe producing auditable documentation and artefacts interpretable by non-developers (8367105 p.24; 5442231 p.15), so designing a toolkit to yield comparable docs/scores without source code access is a reasonable synthesis.} \hllightgreen{External stakeholders also have identified roles.}\pdfcomment[icon=Note]{Multiple sources identify external stakeholders (regulators, auditors, researchers) with roles in oversight (2977039 p.9; 6502223 p.12), supporting the statement.} \hldarkgreen{Regulators under the AI Law Model and EU regimes are expected to supervise high risk systems through registers, inspections and independent audits, yet Kostenko notes systemic gaps in many jurisdictions: weak oversight mechanisms, insufficient institutional capacity for independent auditing, lack of unified impact assessment procedures and absence of public registers \parencites[9]{2977039}.}\pdfcomment[icon=Note]{2977039 (p.9) and 6502223 (p.12) describe supervisory measures (registers, inspections, audits) and explicitly list systemic gaps including weak oversight, limited auditing capacity, lack of unified impact assessment procedures and absent public registers.} \hldarkgreen{Bahangulu and Owusu Berko argue that businesses should adopt independent AI audits and external oversight committees to complement regulatory enforcement, given the difficulty regulators face with black box models and the complexity of auditing for fairness and transparency \parencites[8]{7461574}[9]{7461574}.}\pdfcomment[icon=Note]{7461574 (p.9) argues for independent AI audits and external oversight committees and notes regulators' difficulties with black-box models and auditing complexity, matching the claim.} \hllightgreen{The toolkit is therefore conceived as a bridge artefact that can be used both by internal teams and by independent auditors or supervisory authorities to interrogate recommender configurations through observable outputs, logs and documented gate decisions, in line with this thesis methodological constraint that black box systems are assessed only via published or logged behaviour, not internal telemetry.}\pdfcomment[icon=Note]{8367105 (p.24) and related sources stress assessing systems via observable gate outputs/logs and auditable artefacts rather than internal telemetry, so presenting the toolkit as a bridge artefact for internal and external use is a supported synthesis.} \hllightgreen{Governance of the toolkit itself must mirror the hybrid models that comparative case studies find most effective.}\pdfcomment[icon=Note]{Comparative case studies cited (0375818; 0375818 p.2--6) favour hybrid/co-regulatory models; asserting toolkit governance should mirror such hybrids is a reasonable inference from those findings.} \hldarkgreen{Alrwili and Kantharia classify AI governance into self regulatory, co regulatory and statutory models, and their cross case synthesis of healthcare, civil service HR and municipal DSS deployments finds that hybrid governance combining internal ethics boards, external oversight and statutory frameworks yields bias reduction, faster appeal resolution and higher audit compliance than self regulation alone \parencites[2]{0375818}.}\pdfcomment[icon=Note]{0375818 (p.2--6) classifies self/co/statutory models and reports cross-case synthesis showing hybrid governance (internal ethics boards + external oversight + statutory frameworks) yields bias reduction, faster appeal resolution and higher audit compliance.} \hldarkgreen{In the mental health domain, Malachi et al. show that Pastoral Data Trust governance, mixed audit circles and ethics escrows are critical for maintaining trust and preventing "digital colonialism" when AI systems operate in fragile communities \parencites[11]{5419804}[13]{5419804}.}\pdfcomment[icon=Note]{5419804 (p.11, p.13) documents the Pastoral Data Trust, mixed audit circles and ethics escrows as central governance mechanisms that maintained trust and helped prevent 'digital colonialism' in fragile communities.} \hllightgreen{Drawing on these findings, this thesis situates toolkit governance in co regulatory arrangements: internal AI ethics or CSR committees manage gate configurations and monitor override and incident logs, while independent auditors and, where applicable, national registers under AI Act style regimes obtain access to gate artefacts, model cards and risk assessments.}\pdfcomment[icon=Note]{Sources discuss co-regulatory arrangements, internal ethics bodies, independent auditors and national registers (2977039 p.9; 0375818 p.2--6; 8367105 p.24); situating toolkit governance in such arrangements is a supported synthesis.} \hldarkgreen{Finally, sustainability of the toolkit depends on formal institutionalisation.}\pdfcomment[icon=Note]{6805789 (p.21) and related conclusions stress that sustainability of HITL/tooling depends on institutional and regulatory integration and formal institutionalisation.} \hldarkgreen{Selvam and González Vallejo stress the need for academic oversight committees in HITL grading responsible for rubric standardisation, override log analysis and periodic model recalibration, and advocate alignment with accrediting agencies and policy frameworks \parencites[21]{6805789}.}\pdfcomment[icon=Note]{6805789 (p.21) explicitly calls for academic oversight committees in HITL responsible for rubric standardisation, override log analysis and periodic model recalibration and alignment with accrediting agencies.} \hldarkgreen{Catapang's distinction between organisational framing and control architecture reinforces that many ethical failures are failures of control rather than of principle \parencites[24]{8367105}.}\pdfcomment[icon=Note]{8367105 (p.24) distinguishes organisational framing from control architecture and states ethical failures often stem from failures of control rather than principle, directly supporting the claim.} \hllightgreen{For social recommender audits, this implies that the toolkit must be anchored in documented policies, shared taxonomies and scheduled review cycles, with clear ownership at board and management level, if it is to withstand commercial pressure toward unchecked engagement optimisation while meeting emerging EU expectations around AI audits, transparency and fairness \parencites[9]{7461574}[12]{6502223}[26]{2977039}.}\pdfcomment[icon=Note]{8367105 (p.24), 6502223 (p.12) and 2977039 (p.9, p.26) collectively support anchoring audits in documented policies, taxonomies, review cycles and clear ownership to resist engagement-driven pressures while meeting emerging EU audit/transparency expectations, so this is a supported synthesis.}

\subsubsection{Audit Typologies Relevant for Recommender Systems}
\label{sec:audit_typologies_recommenders} \hllightgreen{Audit work on social recommender systems can draw on several distinct but complementary typologies that already exist in adjacent domains.}\pdfcomment[icon=Note]{Multiple provided sources develop different typologies and related frameworks for algorithmic governance and auditing (e.g., Meßmer \& Degeling 3195574 p.21; AI Law models 2977039 p.9; governance matrices 2702055 p.4), so the claim is a reasonable synthesis of adjacent-domain typologies.} \hldarkgreen{Me{\ss}mer and Degeling distinguish clearly between DSA compliance audits under Article 37, which primarily check whether very large online platforms have conducted and documented the required systemic risk assessments, and risk or impact assessments that interrogate concrete risk scenarios for specific recommender architectures.}\pdfcomment[icon=Note]{Meßmer \& Degeling explicitly distinguish DSA compliance audits under Article 37 from risk/impact assessments and note external auditors will mainly check documentation of risk assessments (3195574 p.21).} \hllightgreen{Their proposal imports the ISO five step audit process planning, risk assessment, audit strategy, evidence gathering and reporting into a socio technical audit that examines ranking behaviour under defined scenarios, for example recommendation of self harm or radicalising content \parencites[21]{3195574}.}\pdfcomment[icon=Note]{Meßmer \& Degeling adopt the ISO five-step audit process (planning, risk assessment, audit strategy, evidence gathering, reporting) for a socio-technical approach (3195574 p.21), and other sources document scenario-based audits of recommender harms (e.g., self-harm/radicalization in 2580223), making this a supported synthesis.} \hllightgreen{In this report's working taxonomy, this yields one axis of audit types: formal compliance audits versus scenario based risk audits targeted at engagement and addiction dynamics.}\pdfcomment[icon=Note]{Meßmer \& Degeling distinguish DSA compliance versus scenario risk assessments (3195574 p.21), and harms literature on recommenders (2580223) supports framing an axis between formal compliance and scenario-based risk/audits focused on engagement/addiction.} \hldarkgreen{Empirical work using archetype or "blank account" methodologies suggests these scenario tests can reveal rapid algorithmic drift toward more extreme or harmful material: in one study, carefully constructed archetype accounts running only passive viewing showed algorithmic saturation in content typologies within about five days, with clear amplification of misogynistic and radical material as viewing continued \parencites[4]{2580223}.}\pdfcomment[icon=Note]{The archetype / blank-account methodology and results (algorithmic saturation after \textasciitilde{}5 days and amplification of misogynistic/radical content) are reported in the TikTok archetype study (2580223 pp.4--5, 11--13).} \hllightgreen{Such findings imply that short, repeatable behavioural probes may be sufficient to surface harmful ranking dynamics for the kinds of scenario based risk audits discussed above \parencites[4]{2580223}.}\pdfcomment[icon=Note]{This is a reasonable inference from the archetype study's rapid saturation/amplification findings (2580223 pp.4--5, 11--13) suggesting short repeatable probes can surface harmful ranking dynamics.} \hldarkgreen{A second typology is emerging around algorithmic impact assessments for high risk systems.}\pdfcomment[icon=Note]{The literature and regulatory discussion explicitly identify algorithmic/impact assessments for high-risk systems as an emerging typology (5442231 p.15; 8367105 p.12).} \hldarkgreen{Sebastian anticipates legally required impact assessments for credit AI, comparable to environmental impact reviews, that test fairness across groups, accuracy and calibration, robustness, privacy, and systemic effects before deployment and on an ongoing basis.}\pdfcomment[icon=Note]{The description that legally required impact assessments for credit AI ,  comparable to environmental reviews and testing fairness, accuracy, robustness, privacy and systemic effects ,  is directly described in the algorithmic impact assessment discussion (5442231 p.15).} \hldarkgreen{These assessments combine quantitative outcome differentials, qualitative analysis of discrimination mechanisms, stakeholder engagement, and risk mitigation planning, and they are expected to be standardised so results become comparable across institutions and over time \parencites[15]{5442231}.}\pdfcomment[icon=Note]{5442231 p.15 explicitly states AIAs combine quantitative outcome differentials, qualitative analysis of discrimination mechanisms, stakeholder engagement, risk mitigation planning, and are expected to be standardized for comparability across institutions.} \hllightgreen{For recommender systems, alignment with such AIA style audits would mean treating engagement optimising feeds as high impact systems subject to pre deployment and periodic assessment on fairness, manipulation and addiction indicators, not only content legality.}\pdfcomment[icon=Note]{Combining the AIA framing for high-risk systems (5442231 p.15) with literature on recommender harms (2580223) supports the inference that engagement-optimizing feeds could be treated as high-impact and assessed for fairness, manipulation and addiction indicators.} \hldarkgreen{Sector specific governance matrices add a third perspective.}\pdfcomment[icon=Note]{The idea of sector-specific governance matrices as an additional perspective is supported by the concrete Algorithmic Governance Assessment Matrix example developed in 2702055 (p.4).} \hldarkgreen{Mutawa et al. construct an Algorithmic Governance Assessment Matrix for zakat institutions that evaluates transparency, accountability, explainability, data integrity and Shariah alignment, based on document analysis and semi structured interviews \parencites[4]{2702055}.}\pdfcomment[icon=Note]{2702055 describes an Algorithmic Governance Assessment Matrix for zakat institutions evaluating transparency, accountability, explainability, data integrity and Shariah alignment based on document analysis and semi-structured interviews (2702055 p.4).} \hldarkgreen{Their empirical results show that transparency scores correlate with clarity of algorithmic documentation and traceability indicators, and that institutions with higher AI readiness display better eligibility accuracy and processing efficiency \parencites[4]{2702055}[8]{2702055}.}\pdfcomment[icon=Note]{2702055 reports empirical results linking transparency scores to clarity of documentation/traceability and shows institutions with higher AI readiness have better eligibility accuracy and processing efficiency (2702055 pp.4, 8).} \hllightgreen{This style of matrix provides a template for recommender audits that rate platform governance arrangements alongside technical performance, for instance scoring explainability and oversight rather than only ranking output statistics.}\pdfcomment[icon=Note]{It is a supported synthesis that such governance matrices could template recommender audits to rate governance (explainability, oversight) alongside technical outputs, given the matrix approach in 2702055 p.4 and broader governance literature.} \hldarkgreen{Kostenko's AI Law Model implicitly defines yet another typology by separating ethical admissibility tests from ongoing supervisory audits.}\pdfcomment[icon=Note]{The AI Law Model describes formalized ethical admissibility tests distinct from ongoing supervisory audit mechanisms, effectively defining separate layers (2977039 pp.8--9, 22).} \hldarkgreen{Ethical admissibility is formalised as mandatory procedures: notification of AI use, explainability of algorithmic conclusions, revision of data sources, guarantees of human control over critical decisions, counteraction to manipulative practices, and protection of vulnerable groups \parencites[8]{2977039}.}\pdfcomment[icon=Note]{2977039 explicitly lists mandatory procedures for ethical admissibility including notification, explainability, data source revision, human control, counteraction to manipulative practices, and protection of vulnerable groups (2977039 p.9).} \hldarkgreen{Supervisory audits then operate on a risk based taxonomy that distinguishes unacceptable, high, limited and minimal risk, augmented by concepts such as systemic and emergent risk, and are executed by accredited auditors and public councils that can reclassify risk based on operational metrics and prior audit findings \parencites[9]{2977039}.}\pdfcomment[icon=Note]{2977039 details a risk taxonomy (unacceptable, high, limited, minimal) augmented with systemic/emergent risks and mechanisms for dynamic reclassification executed by accredited bodies/public councils (2977039 pp.9, 22).} \hllightgreen{Within this thesis, Ethical AI Audit Toolkit components can be mapped to these two layers: design time admissibility checks on manipulative and addictive patterns, and runtime audits tied to risk reclassification.}\pdfcomment[icon=Note]{Mapping toolkit components to design-time admissibility checks and runtime audits tied to risk reclassification is a reasonable synthesis based on the two-layer separation in 2977039 (pp.8--9) and practical toolkit framing in the sources.} \hldarkgreen{Algorithmic auditing protocols from financial AI provide concrete process patterns.}\pdfcomment[icon=Note]{Financial-sector algorithmic auditing protocols are documented and provide concrete process patterns in Thaneswaran's work (8907572 p.7).} \hldarkgreen{Thaneswaran requires models in a cloud connected financial ecosystem to pass a lifecycle validation pipeline that starts with adversarial pre deployment bias testing, continues with a self audit loop where an Ethical Governance Layer logs SHAP based feature attributions to a write once ledger, and is supplemented by bi annual third party reviews.}\pdfcomment[icon=Note]{Thaneswaran prescribes a lifecycle validation pipeline beginning with adversarial pre-deployment bias testing, a self-audit loop logging SHAP attribution telemetry to an immutable ledger, and bi-annual third-party reviews (8907572 p.7).} \hldarkgreen{Circuit breakers automatically switch models into human in the loop fallback modes when drift or disparate impact thresholds are breached \parencites[7]{8907572}.}\pdfcomment[icon=Note]{8907572 describes automated 'circuit breakers' that trigger taking a model offline or switching to a human-in-the-loop fallback when drift or disparate impact thresholds are breached (8907572 p.7).} \hllightgreen{For social recommenders, analogous typologies would distinguish internal continuous self audits, independent external reviews, and automatic throttling of engagement optimisation when harm indicators breach agreed limits.}\pdfcomment[icon=Note]{Projecting the financial-sector typologies (self-audits, independent reviews, circuit-breaker throttles) to social recommenders is a reasonable inference supported by 8907572 (p.7) and the harms literature on recommenders (2580223).} \hldarkgreen{Human in the loop grading demonstrates a further type: embedded operational audits that are inseparable from day to day workflows.}\pdfcomment[icon=Note]{Human-in-the-loop grading as an embedded operational audit inseparable from daily workflows is documented in the HITL grading literature, which highlights integrated, continuous oversight (6805789 p.21).} \hldarkgreen{In Selvam and González Vallejo's four layer framework, every AI score is subject to educator validation, all AI and human actions are dual logged, and institutional oversight committees analyse override logs and recalibrate models.}\pdfcomment[icon=Note]{The four-layer HITL framework described by the source notes educator validation of AI scores, dual logging of AI and human actions, and institutional oversight analyzing override logs and recalibrating models (6805789 p.21).} \hldarkgreen{Students respond positively to transparent, dual sourced feedback, and the system aligns with policy frameworks for ethical AI in education \parencites[21]{6805789}.}\pdfcomment[icon=Note]{6805789 concludes that students responded positively to transparent, dual-sourced feedback and that the system aligned with policy frameworks for ethical AI in education (6805789 p.21, Conclusions).} \hllightgreen{This corresponds, in this report's working classification, to continuous human centric performance audits, where override patterns and user trust metrics are primary objects of analysis.}\pdfcomment[icon=Note]{Interpreting the HITL framework as corresponding to continuous human-centric performance audits focused on override patterns and trust metrics is a reasonable synthesis from 6805789 (p.21) and related governance materials.} \hllightgreen{Across these strands, at least four audit typologies emerge that are relevant to recommender systems in a Digital Fairness trajectory: DSA style compliance audits versus scenario based risk audits \parencites[21]{3195574}; legally mandated AIAs with standardised fairness and impact metrics \parencites[15]{5442231}; governance and ethics matrices that rate transparency, accountability and alignment with normative frameworks \parencites[4]{2702055}[8]{2702055}; and lifecycle audits combining self auditing telemetry, independent external reviews, and human in the loop operational checks \parencites[7]{8907572}[21]{6805789}.}\pdfcomment[icon=Note]{The four typologies listed are supported across the provided sources: DSA compliance vs scenario audits (3195574 p.21), AIAs (5442231 p.15), governance/ethics matrices (2702055 p.4), and lifecycle audits with telemetry and HITL checks (8907572 p.7; 6805789 p.21), so this is a valid synthesis.} \hldarkgreen{The Ethical AI Audit Toolkit proposed later will align its modules explicitly to these audit types, while remaining constrained by the methodological limitation that black box recommenders are evaluated through published behaviour, documentation and secondary audits rather than direct internal access \parencites[6]{6295677}[9]{7461574}[9]{2977039}.}\pdfcomment[icon=Note]{The intention to align an Ethical AI Audit Toolkit to these audit types and the methodological limitation that black-box recommenders are often audited via published behavior, documentation and secondary audits (rather than internal access) is supported by discussions of audit constraints in 7461574 p.9 and 2977039 (e.g., limitations on independent auditing) (7461574 p.9; 2977039 pp.8--9).}

\subsubsection{Alignment with Regulatory and Ethical Requirements}
\label{sec:alignment_regulatory_ethical} \hllightgreen{Design choices in the audit toolkit must map cleanly onto existing hard law and model law obligations rather than free standing ethical aspirations.}\pdfcomment[icon=Note]{The AI Law Model stresses aligning legal/regulatory obligations with technical governance rather than relying on voluntary ethics alone (2977039 p.9-10), so the claim is a reasonable normative synthesis of those sources.} \hldarkgreen{The EU AI Act provides the first anchor: it classifies AI into unacceptable, high, limited, and minimal risk, bans systems that exploit psychological vulnerabilities, and imposes conformity assessment, documentation, and human oversight requirements for high risk domains such as education, employment and essential services \parencites[2248]{6502223}[10]{2977039}.}\pdfcomment[icon=Note]{Descriptions of the EU AI Actrisk tiers, the ban on systems exploiting psychological vulnerabilities, and mandatory conformity assessment, documentation and human oversight for high-risk domains (education, employment, essential services) are explicitly stated in sources describing the AI Act (6502223 p.12; 8589961 p.4).} \hllightgreen{Toolkit criteria that flag addictive feeds or behavioural profiling of minors as "Code Category C: Prohibited or Critical Risk" are therefore aligned with the Act's unacceptable tier and Kostenko's explicit prohibition of systems that manipulate behaviour, cause digital addiction, emotional exhaustion or loss of trust \parencites[61--62]{2977039}.}\pdfcomment[icon=Note]{The AI Actban on manipulative systems (exploit psychological vulnerabilities) and the AI Law Modelcall to counteract manipulative/behavioural practices (2977039 p.8; 6502223 p.12) together support aligning toolkit "prohibited/critical" categories with those prohibitions.} \hllightgreen{Platform facing duties under the Digital Services Act form the second anchor.}\pdfcomment[icon=Note]{The Digital Services Act (DSA) imposes platform-facing duties that are used as a regulatory anchor for platform obligations in the literature (6502223 p.12; 6941611 p.18), supporting this claim.} \hldarkgreen{Very large online platforms must offer at least one non profiled recommendation option, conduct annual independent audits, grant data access to vetted researchers, and assess and mitigate systemic risks to civic discourse and electoral processes \parencites[2248]{6502223}.}\pdfcomment[icon=Note]{The listed obligations for very large online platforms (non-profiled recommendation option, annual independent audits, vetted researcher data access, systemic-risk assessments including civic discourse and electoral risks) are enumerated in summaries of the DSA (6502223 p.12).} \hldarkgreen{Article 25 DSA bans interface designs that deceive or manipulate users or substantially impair free and informed decisions, including giving undue prominence to specific options and making cancellation harder than subscription \parencites[20]{6941611}.}\pdfcomment[icon=Note]{Article 25 and related commentary explicitly prohibit deceptive/manipulative interface designs including undue prominence of options and making cancellation harder than subscription (6941611 p.20).} \hllightgreen{Audit checks in this report's working taxonomy that inspect whether non profiled feeds are genuinely accessible, whether defaults push users into personalised streams, and whether exit flows carry more friction than entry directly operationalise these DSA provisions without claiming internal access to ranking code \parencites[20]{6941611}[6]{6295677}.}\pdfcomment[icon=Note]{The DSA requirements to offer non-profiled options and address interface design are documented (6502223 p.12; 6941611 p.18-20), so audit checks that inspect access to non-profiled feeds, default nudges, and exit-flow friction are a reasonable operationalisation; the specific claim about doing so without internal ranking code is an implementation note not directly evidenced but is a plausible methodological caveat.} \hllightgreen{Model law recommendations extend these hard law duties into a more granular ethical architecture.}\pdfcomment[icon=Note]{The AI Law Model explicitly proposes translating legal duties into more detailed procedural/ethical architecture and institutional mechanisms (2977039 p.8-9), supporting this extension claim.} \hldarkgreen{Kostenko's AI Law Model formalises ethical admissibility tests as mandatory procedures: notification of AI use, explainability of algorithmic conclusions, revision of data sources for quality and legality, guarantees of human control over critical automated decisions, and counteraction to manipulative practices, especially for vulnerable groups \parencites[8]{2977039}.}\pdfcomment[icon=Note]{The AI Law Model lists mandatory ethical admissibility procedures including notification of AI use, explainability, data-source review, guarantees of human control over critical decisions, and counteraction to manipulative practices and protection of vulnerable groups (2977039 p.8).} \hldarkgreen{The same text requires detailed technical passports, event logging, and risk management frameworks for high risk systems, with independent audit and public registers overseen by an interdisciplinary AI council \parencites[9]{2977039}[19.5]{2977039}.}\pdfcomment[icon=Note]{The model calls for technical passports/event logging, risk management frameworks for high-risk systems, independent auditing and public registers overseen by an interdisciplinary council (2977039 p.8; p.106).} \hllightgreen{Audit modules that demand documented explanation interfaces, tamper evident logs, and evidence of human override channels for recommender changes are therefore not discretionary; they are direct translations of these model law expectations into practical checklist items.}\pdfcomment[icon=Note]{The model law's requirements for explainability, logging, human control and independent oversight (2977039 p.8-9) support treating documented explanation interfaces, tamper-evident logs and human-override evidence as practical checklist items; this is a valid translation of those expectations.} \hllightgreen{Ethical governance standards and ISO based risk tools provide a further alignment layer.}\pdfcomment[icon=Note]{Analyses map ISO and other standards to regulatory aims and present them as an additional alignment layer with legal requirements (8589961 p.33), supporting this claim.} \hldarkgreen{Sankaran shows that ISO/IEC 5259 on data quality, 24368 on ethical principles, 24027 on bias assessment, and 23894 on AI risk management collectively support EU AI Act requirements on representative datasets, fairness, transparency and human oversight, although they are voluntary outside strict regulatory environments \parencites[4]{8589961}.}\pdfcomment[icon=Note]{The mapping of ISO standards (e.g., ISO/IEC 5259 on data quality, 24368 on ethical principles, 24027 on bias assessment, 23894 on AI risk management) to support EU AI Act goals, while noting their voluntary nature, is explicitly discussed (8589961 p.33; p.4).} \hlorange{Sebastian anticipates that credit AI will be subject to standardised fairness tests, documentation of performance by demographic group, and continuous monitoring of disparate impact as part of algorithmic impact assessments \parencites[15]{5442231}.}\pdfcomment[icon=Note]{No provided source identifies an author 'Sebastian' making the specific predictions about credit AI; while credit fairness topics are discussed (5442231), the named "Sebastian" claim and its detailed forecasts are not supported by the supplied sources.} \hllightgreen{The toolkit can therefore adopt ISO style bias and risk controls as internal metrics and evidence artefacts, then explicitly crosswalk them to AI Act and AI Law Model criteria, echoing Catapang's regulatory mapping of triple gates to EU AI Act and NIST AI RMF obligations \parencites[23]{8367105}.}\pdfcomment[icon=Note]{The idea of adopting ISO-style bias and risk controls and crosswalking them to AI Act and AI Law Model criteria is supported by ISO-to-AI Act mappings (8589961 p.33) and regulatory crosswalk concepts (8367105 p.23), and Catapang-like regulatory mapping is echoed in the provided crosswalk literature.} \hllightgreen{Typologies of AI governance clarify which organisational structures the toolkit should treat as compliant.}\pdfcomment[icon=Note]{Typologies of AI governance (self-regulation, co-regulation, statutory) and their role in determining compliant organisational arrangements are described in case study literature (0375818 p.2-4), supporting this claim.} \hldarkgreen{Comparative case studies of healthcare, civil service HR and municipal DSS show that hybrid governance combining internal ethics boards, external oversight and statutory frameworks produces measurable improvements in bias reduction, appeal resolution times and audit compliance scores, while pure self regulation risks insufficient accountability \parencites[2]{0375818}.}\pdfcomment[icon=Note]{The comparative case studies across healthcare, civil service HR and municipal DSS in the provided source report that hybrid governance (internal ethics boards, external oversight, statutory frameworks) produced measurable improvements (e.g., bias reduction, appeal resolution times, audit compliance) and that pure self-regulation risks insufficient accountability (0375818 p.5-7).} \hllightgreen{Bahangulu and Owusu Berko argue that businesses must adopt independent AI audits, external oversight committees and policies aligned with the EU AI Act and GDPR, because regulators struggle with black box models and require better auditing techniques and collaboration to enforce fairness and transparency \parencites[9]{7461574}.}\pdfcomment[icon=Note]{The literature argues businesses should adopt independent AI audits, external oversight committees and align policies with the EU AI Act/GDPR, and notes regulator difficulties with black-box models and the need for better auditing and collaboration (7461574 p.1754; 7461574 p.1754-1755).} \hllightgreen{Governance sections of the toolkit therefore align with these findings by treating the presence of co regulatory or statutory style oversight, documented audit trails and independent review as positive Code C signals, and by explicitly acknowledging that black box recommenders can only be evaluated here via published reports, external audits and observable UX behaviour \parencites[6]{6295677}[9]{7461574}.}\pdfcomment[icon=Note]{Sources recommend co-regulatory/statutory oversight, documented audit trails and independent review as governance positives and recognize limitations in assessing black-box recommenders without external reports or UX observation (7461574 p.1754; 4268459 p.69), supporting this alignment claim.} \hllightgreen{Finally, ethical documentation and communication practices drawn from HITL grading and educational AI provide concrete design patterns consistent with UNESCO and EU expectations.}\pdfcomment[icon=Note]{HITL grading and educational AI literature provide documentation and communication practices that the sources link to international (UNESCO) and EU expectations about transparency and human oversight (6805789 p.20; p.15), supporting this statement.} \hldarkgreen{Selvam and González Vallejo demonstrate that four layer HITL grading, with SHAP based explainability, human confirmation, dual logging and explicit tagging of AI versus instructor feedback, aligns with transparency and human oversight principles and increases student trust \parencites[6805789]{6805789}[20]{6805789}.}\pdfcomment[icon=Note]{The four-layer HITL grading framework with SHAP-based explainability, human confirmation, audit logging and explicit tagging of AI vs instructor feedback, and its positive effects on student trust, are explicitly described in the HITL study (6805789 p.20; p.15).} \hldarkgreen{Mogoi et al. argue that AI powered mobile proctoring must integrate privacy by design, informed consent, ethical auditing mechanisms and transparent communication about how AI decisions are made \parencites[9]{2055956}.}\pdfcomment[icon=Note]{The review on AI-powered mobile proctoring explicitly recommends privacy-by-design, informed consent, ethical auditing mechanisms and transparent communication about AI decision-making (2055956 p.35; p.4, 4.4).} \hllightgreen{These practices inform audit items requiring clear user facing disclosure of AI involvement in recommendation decisions, accessible appeal procedures, and explicit tagging of personalised versus non personalised content.}\pdfcomment[icon=Note]{Disclosure, accessible appeals, and explicit tagging of personalised vs non-personalised content are recommended practices in the HITL literature and DSA/AI transparency discussions (6805789 p.15; 6502223 p.12), so using them to inform audit items is a justified synthesis.} \hllightgreen{Aligning toolkit components with these concrete regulatory and ethical requirements maintains congruence with EU AI Act and DSA obligations, emerging Digital Fairness directions, and AI Law Model principles, while staying within the evidentiary limits of qualitative secondary research on black box social platforms \parencites[2248]{6502223}[20]{6941611}[8]{2977039}[9]{2055956}.}\pdfcomment[icon=Note]{The claim that aligning toolkit components with EU AI Act, DSA, Digital Fairness trends and the AI Law Model is congruent with the sources (2977039 p.9-10; 6941611 p.18-20; 2055956 p.4) and notes limits of qualitative secondary research on black-box platforms, which is discussed in the literature.}

\subsection{Audit Preparation and Scoping Modules}

\subsubsection{Defining System Boundaries and Use Cases}
\label{sec:defining_system_boundaries_use_cases} \hllightgreen{Scoping an ethical audit begins with drawing clear boundaries around which socio technical artefact is being examined.}\pdfcomment[icon=Note]{The statement is a general framing supported by governance and audit literature which emphasise scoping and boundary-setting for socio-technical artefacts (see 2977039 on defining system boundaries and transparency requirements, and 0375818 on governance frameworks).} \hllightgreen{In this thesis, recommender systems are treated as black box services whose internal model weights and proprietary training pipelines cannot be inspected directly; assessment instead relies on observable interface behaviour, published technical descriptions, regulatory records, and secondary empirical studies \parencites[7]{6295677}[2]{8958037}[2]{3303310}.}\pdfcomment[icon=Note]{Sources discuss treating recommendation systems as opaque/black-box and relying on observable behaviour, documentation and secondary studies (see 3303310 on black-box exploitation and 2977039 on audit mechanisms and use of published documentation).} \hllightgreen{This constraint mirrors designs in dark pattern research, where qualitative content analysis of regulatory documents, institutional audits, and consumer surveys is used to characterise systemic manipulation without access to platform code or telemetry \parencites[7]{6295677}.}\pdfcomment[icon=Note]{This mirrors methods used in dark-pattern and HCI regulatory analyses that use qualitative content analysis of regulatory documents, audits and surveys when internal code is unavailable (see 6295677 on qualitative content analysis of regulatory documents and typologies; 7896734 on dark patterns and interface analysis).} \hlorange{System boundaries therefore need to be articulated along at least three axes.}\pdfcomment[icon=Note]{No provided source explicitly states the exact proposition that system boundaries must be articulated along 'at least three axes'; this is a structural claim introduced by the author without direct support in the cited materials.} \hlorange{First, functional scope: which sub systems count as part of the "social recommender" for audit purposes.}\pdfcomment[icon=Note]{This is a definitional/structural claim about the 'functional scope' axis; none of the provided sources explicitly frames functional scope in this exact wording as a primary axis for scoping audits.} \hldarkgreen{Inspired by Mutawa et al.'s Algorithmic Governance Assessment Matrix for zakat distribution, which covers transparency, accountability, explainability, data integrity, and Shariah alignment across AI models and organisational processes, this report treats the recommender as comprising ranking logic, data capture and storage, user facing interfaces, and surrounding governance mechanisms \parencites[2702055]{2702055}.}\pdfcomment[icon=Note]{Source 2702055 explicitly describes an Algorithmic Governance Assessment Matrix covering transparency, accountability, explainability, data integrity and Shariah alignment, and the sentence directly paraphrases that scope and applies it to recommender components (2702055, p.4).} \hldarkgreen{For example, in a SOCIO SPHERE AI style architecture, AI powered recommendation components, sentiment analysis, real time notifications, and AI assisted replies sit alongside content creation tools and moderation functions within a modular three tier system with an additional AI/ML processing layer \parencites[3413731]{3413731}.}\pdfcomment[icon=Note]{The SOCIO SPHERE AI architecture and the listed AI components (recommendation, sentiment analysis, real-time notifications, AI-assisted replies, moderation) are described in the SOCIO-SPHERE specification and system description (3413731, pp.1-2) and the role of sentiment/emotion analysis is addressed in 3625351.} \hllightgreen{An audit boundary that included ranking while ignoring AI assisted replies or notification policies would miss key engagement and manipulation channels.}\pdfcomment[icon=Note]{Sources discuss how ranking, notifications and AI-assisted interactions act as engagement and manipulation channels (3625351 on emotional amplification; 3413731 on AI-assisted replies and notification features), so the inference that excluding some channels would miss manipulation is a supported synthesis.} \hllightgreen{Second, institutional scope: which entities carry responsibility for decisions.}\pdfcomment[icon=Note]{The notion of an 'institutional scope' identifying responsible entities aligns with governance literature that distinguishes actors and allocates responsibility (see 0375818 on models of AI governance and 2977039 definitions of actors).} \hldarkgreen{Kostenko distinguishes AI providers, who create and maintain systems, from AI operators, who integrate and use them in specific contexts, and places legal obligations on both \parencites[20]{2977039}.}\pdfcomment[icon=Note]{2977039 defines and distinguishes 'AI provider' and 'AI operator' and places obligations on both roles (see definitions and obligations in 2977039, pp.19-20 and pp.96--99).} \hllightgreen{Alrwili and Kantharia's comparative work on decision support systems shows that governance patterns self regulatory, co regulatory, statutory correlate with different oversight mechanisms and performance outcomes \parencites[2]{0375818}.}\pdfcomment[icon=Note]{The classification of governance patterns as self-regulatory, co-regulatory and statutory and their correlation with oversight mechanisms and outcomes is supported by 0375818 which outlines those governance models and comparative outcomes (0375818, Sect. 4--7).} \hldarkgreen{For this toolkit, system boundaries are drawn to include not only technical modules but also the internal ethics boards, audit units, and external auditors that act on recommender outputs, since their presence or absence shapes how risks identified by the audit can be addressed in practice \parencites[2]{0375818}[2702055]{2702055}.}\pdfcomment[icon=Note]{2702055 discusses the role of institutional governance structures (ethics boards, audit units) in zakat institutions and how governance presence affects addressing risks; the sentence aligns with that source (2702055, p.4).} \hllightgreen{Third, jurisdictional scope: which legal regimes apply.}\pdfcomment[icon=Note]{The need to set jurisdictional scope (which legal regimes apply) is a standard element in the legal/regulatory literature and is consistent with subsequent citations to model law and EU instruments (see 2977039 and 5142478).} \hldarkgreen{Kostenko's technological sovereignty provisions prohibit deployment of AI systems where independent audit is impossible, critical components are outside national control, or hidden behaviour influencing mechanisms exist, and require a special register of prohibited or unwanted AI systems \parencites[108]{2977039}.}\pdfcomment[icon=Note]{2977039 explicitly contains technological sovereignty provisions prohibiting systems where independent audit is impossible, critical components are outside national control, or hidden influencing mechanisms exist, and mandates a special register (2977039, p.108).} \hllightgreen{EU instruments such as the AI Act and DSA impose transparency, documentation, data access and systemic risk duties on certain classes of systems, while authors such as Chitra et al. argue that very large platforms should be subject to mandatory algorithmic auditing and researcher access for recommendation algorithms that influence political discourse \parencites[16]{6502223}.}\pdfcomment[icon=Note]{Sources document that the EU AI Act and DSA impose duties like transparency, documentation, data access and systemic risk obligations (see 5142478 and 0375818), and other authors call for mandatory auditing and researcher access for large platforms (e.g. discussions in 3303310 and 6502223 on platform auditing and researcher access).} \hllightgreen{In defining system boundaries, the audit must therefore specify which parts of a platform's stack are subject to EU style duties and which fall under model law style expectations.}\pdfcomment[icon=Note]{This is a reasonable synthesis of legal/regulatory distinctions in the sources: practitioners must delineate parts subject to EU-style duties versus broader model-law expectations (see 2977039 on different legal mechanisms and 5142478 on EU instruments).} \hlorange{On that basis, this report adopts a structured use case taxonomy to scope audit engagements.}\pdfcomment[icon=Note]{The claim that 'this report adopts a structured use case taxonomy' is an authorial statement about the report itself and is not something the external sources can corroborate; no provided source documents the report's internal choice.} \hllightgreen{Table~\ref{tab:use_cases_scope} illustrates three archetypal use cases that will recur in later chapters. \begin{table}[h] \centering \caption{Illustrative use case archetypes and boundary focus (working taxonomy)} \label{tab:use_cases_scope} \begin{tabular}{p{3.2cm}p{4.4cm}p{4.4cm}} \hline Use case archetype & Boundary focus & Primary sources informing scope \\ \hline Short form video feed for students & Infinite scroll, auto play, push notifications, personalisation pipeline, child data processing, operator governance & Algorithmic addiction loop, student harms, DSA proceeding against TikTok, age appropriate design and child data processing duties \parencites[2]{8958037}[6]{8958037}[2]{3303310}[20]{2977039} \\ Image centric social discovery & Explore style ranking, social proof widgets, behavioural advertising modules, bias and body image impacts & Amplification of body negative imagery, negative affect bias, social proof manipulation and its impact on trust and identity \parencites[3]{3303310}[13]{7160486} \\ AI augmented social platform (SOCIO SPHERE AI) & AI recommendation layer, sentiment analysis, AI assisted replies, moderation AI, admin analytics dashboards & Three tier plus AI/ML architecture, AI generated captions and hashtag suggestions, AI based moderation, engagement analytics \parencites[3413731]{3413731} \\ \hline \end{tabular} \end{table} Each archetype defines the "system under audit" differently.}\pdfcomment[icon=Note]{The illustrative archetypes (short video feed, image-centric discovery, SOCIO SPHERE AI) and their concerns (infinite scroll, personalization, moderation, child data) are consistent with cited literature on algorithmic addiction, youth data protection, platform harms and architecture (see 3303310 on engagement harms, 2977039 on child data and AI systems, and 3413731 on SOCIO-SPHERE design).} \hldarkgreen{In the student short video case, child profiling, emotional tracking and predictive "life path" scoring sit at the centre of ethical scrutiny because AI child profiling and processing of children's data attract enhanced protection and, in some forms, categorical prohibition \parencites[20]{2977039}[62]{2977039}.}\pdfcomment[icon=Note]{2977039 explicitly identifies AI child profiling and processing of children's data as subject to enhanced protection and defines 'AI child profiling' (see 2977039, pp.19-20, definitions including Child Data and AI child profiling).} \hllightgreen{For SOCIO SPHERE AI, sentiment analysis and AI assisted replies embody higher order capabilities that shape micro interactions and therefore must be included in the audit boundary even if the underlying models are shared across multiple content surfaces \parencites[3413731]{3413731}.}\pdfcomment[icon=Note]{Sources describe sentiment analysis and AI-assisted replies as capabilities that shape interactions and engagement (3625351 on emotion amplification; 3413731 on AI-assisted replies), supporting the claim that they should be included in the audit boundary even when models are shared.} \hllightgreen{Use cases also define which metrics and documents are relevant.}\pdfcomment[icon=Note]{The idea that use cases determine relevant metrics and documents is supported by regulatory guidance and audit literature emphasizing use-case-specific documentation and metrics (see 2977039 on required documentation and audit logs, and 2702055 on use-case performance indicators).} \hldarkgreen{Mutawa et al. demonstrate that transparency scores in their matrix correlate with clarity of algorithmic documentation, and that eligibility accuracy and processing time reduction provide a baseline for ethical adoption in zakat institutions \parencites[2702055]{2702055}.}\pdfcomment[icon=Note]{2702055 presents transparency scores correlated with clarity of algorithmic documentation and reports eligibility accuracy and processing time reduction as key indicators in zakat institutions (2702055, p.4 and Table 1).} \hldarkgreen{Although this thesis does not have direct access to recommender performance logs, the same logic informs audit scoping: for each use case, the toolkit specifies expected artefacts such as technical passports, audit logs, bias assessments or human rights impact assessments, consistent with Kostenko's transparency and explainability chapter \parencites[68]{2977039}[69]{2977039}.}\pdfcomment[icon=Note]{2977039 outlines transparency and explainability requirements (audit logs, technical passports, documentation) and the sentence accurately synthesises that logic for audit scoping (see 2977039, pp.68-70).} \hllightgreen{Finally, defining system boundaries and use cases requires acknowledging what remains out of scope.}\pdfcomment[icon=Note]{Acknowledging out-of-scope elements when defining boundaries is a standard audit practice and is consistent with the constraints and transparency considerations discussed in the sources (see 2977039 on controlled access and data rooms and 0375818 on methodological limits).} \hllightgreen{Internal training corpora, model architectures and real time A/B experiment results are treated here as inaccessible, mirroring the limitations described in dark pattern and AI governance research \parencites[7]{6295677}[9]{7461574}.}\pdfcomment[icon=Note]{The treatment of internal training corpora, model architectures and real-time A/B experiment results as inaccessible mirrors discussions about proprietary opacity and limitations in platform research and dark-pattern studies (see 3303310 on black-box opacity and 6295677 on limits of secondary data analyses).} \hllightgreen{The Ethical AI Audit Toolkit is therefore scoped as a qualitative and documentation centred instrument: it evaluates social recommender systems through their declared objectives, interface mechanics, governance structures and published or logged outputs, aligned with AI law and platform regulation expectations yet constrained by the absence of internal telemetry \parencites[68]{2977039}[21]{3195574}[2]{0375818}.}\pdfcomment[icon=Note]{The description of the Ethical AI Audit Toolkit as qualitative and documentation-centred, aligned with AI law and constrained by absence of internal telemetry, is supported by methodological guidance in 0375818 and by transparency and audit constraints in 2977039 (see 0375818 on qualitative/documentation approaches and 2977039 on documentation/audit expectations).}

\subsubsection{Stakeholder Mapping and Risk Prioritization}
\label{sec:stakeholder_mapping_risk_prioritization} \hldarkgreen{Stakeholder mapping for auditing social recommender systems begins with identifying who carries legal duties and who bears risk.}\pdfcomment[icon=Note]{The AIN and AI Law Model describe stakeholder mapping that identifies responsible actors and those bearing risk (Leon 2742041 p.2--3; AI Law Model 2977039 p.61).} \hldarkgreen{Kostenko distinguishes multiple actor groups in AI lifecycles: developers, implementers, administrators, operators, data controllers, state regulators, and affected individuals, all bound by the principle of ethical responsibility and obliged to operate systems in line with autonomy, non discrimination and psycho emotional safety \parencites[61--62]{2977039}.}\pdfcomment[icon=Note]{The AI Law Model explicitly lists developers, implementers, administrators, operators, data controllers, state regulators and affected individuals and mandates ethical responsibility including autonomy and psycho-emotional protections (AI Law Model 2977039 p.61).} \hldarkgreen{Leon's Algorithmic Impact Navigator similarly assigns differentiated roles to organisational leadership, developers, public agencies, end users and civil society, stressing that AI governance is a collective endeavour and that organisations deploying AI are "the first line of accountability" \parencites[2]{2742041}[3]{2742041}.}\pdfcomment[icon=Note]{Leon's AIN sets out differentiated stakeholder roles (organizational leadership, developers, public agencies, end users, civil society) and states organisations are the first line of accountability (Leon 2742041 p.2--3).} \hllightgreen{From an audit design perspective, at least four stakeholder clusters are central.}\pdfcomment[icon=Note]{This is a synthesis claim consistent with stakeholder groupings in the AIN and related governance literature which organise actors into clusters for audit design (Leon 2742041 p.2; AI Law Model 2977039 p.61).} \hldarkgreen{Platform organisations and their internal governance structures sit at the core.}\pdfcomment[icon=Note]{The AIN and AI Law Model position platform organisations and internal governance as central to AI oversight and accountability (Leon 2742041 p.2; AI Law Model 2977039 p.61).} \hldarkgreen{Senior leadership is responsible for commissioning impact evaluations, aligning projects with regulatory obligations, and resourcing ethical oversight \parencites[2]{2742041}.}\pdfcomment[icon=Note]{Leon's framework assigns senior leadership responsibility for commissioning impact evaluations, regulatory alignment and resourcing ethical oversight (Leon 2742041 p.2--3).} \hldarkgreen{Internal roles include product and engineering teams who translate goals into algorithms, risk and compliance units that interpret regulations such as the AI Act and DSA, and ethics committees that implement internal ethical reviews, as required by the AI Law Model \parencites[61--62]{2977039}.}\pdfcomment[icon=Note]{Sources describe internal product/engineering roles, risk/compliance units interpreting the AI Act/DSA, and institutional ethics committees as governance elements recommended by the AI Law Model (Leon 2742041 p.2--3; AI Law Model 2977039 p.9, p.61; CI/CD gating controls 8367105 p.24).} \hldarkgreen{External oversight actors form a second cluster: EU regulators enforcing the AI Act and DSA, national authorities envisaged in model law regimes, and accredited auditors and professional bodies in co regulatory or statutory governance models \parencites[10]{2977039}[2]{0375818}.}\pdfcomment[icon=Note]{The DSA and EU AI Act plus model-law discussions reference EU regulators, national authorities, and co-regulatory arrangements with accredited auditors/professional bodies (EU frameworks 6502223 p.12; AI Law Model 2977039 p.9; governance typology 0375818 p.2).} \hldarkgreen{Users and affected communities constitute a third group, including heavy social media users with high basic digital skills but low manipulation awareness \parencites[10]{9360054}, students experiencing algorithmic addiction on short video platforms \parencites[2]{8958037}, and communities subject to AI mediated interventions such as mental health programmes audited through Pastoral Data Trust structures \parencites[13]{5419804}.}\pdfcomment[icon=Note]{Empirical and theoretical sources identify users/communities (heavy social media users, student vulnerability and algorithmic addiction) and community-centred governance like Pastoral Data Trust in mental-health interventions (algorithmic addiction and neuro evidence 3303310 p.2; student studies 4245543 p.4; Pastoral Data Trust 5419804 p.13).} \hldarkgreen{A fourth group comprises epistemic stakeholders: researchers and civil society organisations that rely on data access and transparency to study systemic risks in algorithmic feeds \parencites[21]{3195574}.}\pdfcomment[icon=Note]{Researchers and civil society as epistemic stakeholders relying on data access and transparency are discussed in the DSA provisions and auditing literature (DSA/vetted researcher access 6502223 p.12; auditors/researcher needs 3195574 p.44).} \hllightgreen{Risk prioritisation then hinges on how these groups interact with particular recommender deployments.}\pdfcomment[icon=Note]{This is a reasonable synthesis: Leon's screening logic and governance literature imply that risk prioritisation depends on group interactions with deployments (Leon 2742041 p.3; AI Law Model 2977039 p.9).} \hldarkgreen{Leon's AIN proposes initial screening questions around sociopolitical context, technological opacity and integration with critical infrastructures to decide which projects require full impact assessment \parencites[3]{2742041}.}\pdfcomment[icon=Note]{Leon's AIN explicitly lists initial screening questions on sociopolitical context, technological opacity and integration with critical infrastructures (Leon 2742041 p.3).} \hllightgreen{Applying this logic, engagement optimised short video feeds used heavily by students (high vulnerability, opaque deep learning ranking, strong integration into daily routines) rank as high priority targets, while low reach experimental features with minimal autonomy impact might receive lighter scrutiny.}\pdfcomment[icon=Note]{This applies Leon's screening criteria to known vulnerabilities (student vulnerability and opaque deep-learning feeds) forming a reasonable prioritisation inference (Leon 2742041 p.3; student/engagement harms 3303310 p.2; 4245543 p.4).} \hldarkgreen{Chitra et al. argue that recommendation systems of very large platforms, especially where they affect civic discourse and elections, warrant systemic risk assessment under the DSA, including data access for vetted researchers \parencites[2248]{6502223}.}\pdfcomment[icon=Note]{The DSA requires systemic risk assessment for very large platforms affecting civic discourse and provides for researcher access to vetted researchers (6502223 p.12), consistent with Chitra et al.'s claim as cited.} \hllightgreen{These criteria feed directly into this report's working taxonomy for risk prioritisation.}\pdfcomment[icon=Note]{This is a reasonable authorial statement: the criteria from AIN and DSA-like requirements naturally inform a working taxonomy for risk prioritisation (Leon 2742041 p.3; DSA provisions 6502223 p.12).} \hllightgreen{A structured view of stakeholder specific risk priorities is useful for operational scoping.}\pdfcomment[icon=Note]{Multiple sources advocate structured stakeholder-focused risk assessments to scope audits operationally (AI Law Model 2977039 p.9; impact assessment discussions 5442231 p.15).} \hllightgreen{Table~\ref{tab:stakeholder_risk_priorities} outlines a working mapping. \begin{table}[h] \centering \caption{Stakeholder centric risk priorities for social recommender audits (working taxonomy)} \label{tab:stakeholder_risk_priorities} \begin{tabular}{p{3cm}p{4.2cm}p{4.5cm}} \hline Stakeholder cluster & Primary concern & High priority risk cues \\ \hline Platform leadership and boards & Legal compliance, reputation, business continuity & Presence of addictive design patterns under DSA scrutiny; absence of internal ethical review and AI Ethics Impact Reports \parencites[2]{8958037}[61--62]{2977039} \\ Engineering and product teams & Technical performance, shipping velocity & Lack of gating controls in CI/CD for fairness, transparency and sustainability; missing logs needed for conformity assessment \parencites[24]{8367105}[19.5]{2977039} \\ Regulators and auditors & Systemic risk, fundamental rights & Inability to obtain technical dossiers and logs; absence of algorithmic impact assessments and standardised fairness metrics \parencites[19]{3195574}[15]{5442231} \\ Users and communities & Autonomy, psycho emotional safety, fairness & Infinite scroll and auto play combined with personalised feeds; documented mental health or discrimination harms \parencites[2]{8958037}[2]{3303310}[4245543]{4245543} \\ Researchers and civil society & Transparency, contestation & Restricted data access under DSA Article 40 analogues; opaque terms and conditions on content and risk management \parencites[19]{3195574}[3195574]{3195574} \\ \hline \end{tabular} \end{table} Hybrid governance evidence indicates that risk prioritisation should also be sensitive to governance models.}\pdfcomment[icon=Note]{The table and its cited items draw on sources about stakeholder concerns and governance sensitivity (citations in table map to AI Law Model 2977039 p.61; researchers/audit access 3195574 p.44; DSA/mental health 6502223 p.12; 4245543 p.4).} \hldarkgreen{Case studies in healthcare, civil service HR and municipal services show that hybrid governance combining internal ethics boards, external regulators and statutory frameworks yields greater bias reduction, faster appeal resolution and higher audit compliance than self regulation alone \parencites[2]{0375818}.}\pdfcomment[icon=Note]{The case-study synthesis in 0375818 reports that hybrid governance combining internal boards, external regulators and statutory frameworks produces better bias reduction, faster appeal resolution and higher audit compliance (0375818 p.2, Results/6--7).} \hldarkgreen{Mutawa et al. likewise find that zakat institutions with higher AI readiness, better transparency scores and stronger oversight mechanisms achieve higher eligibility accuracy and processing efficiency \parencites[2702055]{2702055}.}\pdfcomment[icon=Note]{The zakat study reports that institutions with higher AI readiness, transparency and oversight achieved higher eligibility accuracy and processing efficiency (2702055 p.4, Table 1 and accompanying discussion).} \hllightgreen{Social recommender audits should therefore treat platforms operating under co regulatory or statutory regimes, with documented ethics councils and clear appeal channels, as lower structural risk than services reliant solely on voluntary self assessment.}\pdfcomment[icon=Note]{This is a reasonable policy inference grounded in comparative governance evidence that statutory/co-regulatory regimes reduce structural risk relative to pure self-regulation (0375818 p.2; AI Law Model 2977039 p.9).} \hllightgreen{Finally, methodological limits constrain how stake and risk can be evidenced.}\pdfcomment[icon=Note]{The literature notes methodological limits in evidencing stake and risk due to data access and reproducibility issues, supporting this general caution (6295677 p.7; 6577872 p.34).} \hldarkgreen{Khare et al. demonstrate that dark pattern prevalence on Indian platforms had to be assessed via secondary institutional reports, regulatory audits and survey data covering around $30$ major platforms, because primary experimentation and internal data access were unavailable \parencites[7]{6295677}.}\pdfcomment[icon=Note]{The Indian dark-pattern study explicitly relies on secondary institutional reports, regulatory audits and surveys covering about 30 major platforms because primary experiments and internal telemetry were unavailable (6295677 p.7, Sampling Framework/Analysis).} \hldarkgreen{Hu's analysis of engagement worship similarly infers algorithmic manipulation and mental health impacts from controlled experiments and neuroimaging studies, not from platform telemetry \parencites[2]{3303310}.}\pdfcomment[icon=Note]{The engagement-worship study infers manipulation and mental-health impacts from controlled experiments and neuroimaging rather than platform telemetry (3303310 p.2 describes neuroimaging and controlled-experiment methods and notes lack of telemetry).} \hllightgreen{This thesis adopts the same stance: stakeholder mapping and risk prioritisation are derived from documentary and empirical sources that describe behaviours and harms external to proprietary systems, and the Ethical AI Audit Toolkit is explicitly framed as a qualitative, documentation centric instrument aligned with AI Act and AI Law Model expectations rather than a substitute for internal model forensics \parencites[10]{2977039}[21]{3195574}[9]{7461574}.}\pdfcomment[icon=Note]{This thesis-level methodological stance aligns with prior audit guidance emphasising documentation- and documentation-centric qualitative tools and alignment with AI Act/AI Law Model expectations (3195574 p.44; AI Law Model 2977039 p.61).}

\subsubsection{Documentation and Data Access Pre-Checks}
\label{sec:doc_data_prechecks} \hllightgreen{Any serious audit of social recommender systems must begin by checking whether the documentation and data access conditions that regulators expect actually exist.}\pdfcomment[icon=Note]{The AI Law Model and related regulatory sources emphasize checking documentation, registers and data access as foundational oversight steps (AI Law Model [2977039], pp. 9, 244), so this is a reasonable synthesis of those sources.} \hldarkgreen{Before scenario based or fairness oriented assessments are attempted, auditors need to know if there is an identifiable AI system with a registrable identity, a technical dossier, and an evidentiary logging setup that can support later scrutiny.}\pdfcomment[icon=Note]{The AI Law Model explicitly requires identifiable AI systems, registration attributes and event logging as preconditions for oversight (AI Law Model [2977039], p.244), directly supporting this sentence.} \hllightgreen{Kostenko's AI Law Model provides a concrete checklist.}\pdfcomment[icon=Note]{The document repeatedly refers to an "AI Law Model" that provides concrete registration and documentation checklists (AI Law Model [2977039], pp. 9, 244), so attributing a checklist to that model is a reasonable summary.} \hldarkgreen{Under the TSAI regime, high risk systems must be registered and assigned a unique identifier, and an algorithmic passport must be available as the official digital identity.}\pdfcomment[icon=Note]{The TSAI regime in the AI Law Model requires registration, assignment of a unique identifier and an algorithmic passport as the system's digital identity (AI Law Model [2977039], p.244).} \hldarkgreen{This passport must at least specify owner and authorised representative, base jurisdiction, assigned risk class, purpose and scope of use, model architecture description, data sources and legal grounds, update and versioning procedures, human oversight channels and appeal mechanisms, modes of event logging with storage periods, and information on prior audits and trust markings.}\pdfcomment[icon=Note]{The AI Law Model lists required contents of the algorithmic passport (owner, jurisdiction, risk class, purpose, architecture, data sources/legal grounds, update/versioning, human oversight and appeals, logging modes/storage periods, audits/trust markings) (AI Law Model [2977039], p.244).} \hllightgreen{A documentation pre check in the present toolkit therefore first asks whether an equivalent passport exists for the recommender under review, even if platforms are not yet formally bound by TSAI.}\pdfcomment[icon=Note]{The AI Law Model's emphasis on algorithmic passports and registration supports a toolkit pre-check asking for equivalent passports even where TSAI is not formally binding (AI Law Model [2977039], p.244); this is a reasonable policy-to-practice inference.} \hllightgreen{Absence of such a dossier signals that later audit steps will face serious traceability gaps.}\pdfcomment[icon=Note]{The AI Law Model highlights logging and registration as essential to traceability; their absence would create gaps for later scrutiny (AI Law Model [2977039], pp. 191, 244), so this is a supported inference.} \hllightgreen{Data access duties for regulators and researchers need to be anticipated next.}\pdfcomment[icon=Note]{Multiple sources (AI Law Model [2977039], pp. 191, 244; policy literature on platform audits) stress anticipating data access duties for regulators and researchers, so this is a supported synthesis.} \hldarkgreen{Under TSAI, owners or deployers must ensure human supervision, preserve and provide event logs, conduct audits, and comply with data protection, non discrimination and transparency requirements, with failure triggering orders, restrictions, suspension of identifiers, fines, blocking of access and compensation obligations \parencites[244]{2977039}.}\pdfcomment[icon=Note]{The AI Law Model enumerates owner/deployer duties: ensuring human supervision, saving/providing event logs, conducting audits, complying with data protection/non-discrimination/transparency, and prescribes orders, restrictions, suspension of identifiers, fines, blocking access and compensation for violations (AI Law Model [2977039], p.244; sanctions described p.244).} \hllightgreen{Me{\ss}mer and Degeling show how, under the DSA, Article 40 data access for Digital Services Coordinators and vetted researchers is central to systemic risk audits of recommender systems \parencites[32]{3195574}.}\pdfcomment[icon=Note]{Meßmer \& Degeling discuss systemic-risk oriented audits of recommenders and the role of regulatory data access; while their paper focuses on audit scenarios (Meßmer \& Degeling [3195574], p.32), linking this to Article 40 data access under the DSA is a reasonable synthesis with DSA-focused literature.} \hldarkgreen{A pre check must therefore determine whether event logs are kept in a form that can be shared with authorities and external auditors, whether log retention meets minimum periods (for example at least $12$ months in the AI Law Model) and whether there are clear internal procedures for supplying logs and technical materials on request.}\pdfcomment[icon=Note]{The AI Law Model prescribes keeping event/decision logs with a storage period of at least 12 months and provisions for retention in case of dispute/inspection, and access by affected parties (AI Law Model [2977039], p.191).} \hllightgreen{Information about logging practices must go beyond a binary yes or no.}\pdfcomment[icon=Note]{The AI Law Model and audit guidance require detailed logging practices rather than binary statements of whether logging exists (AI Law Model [2977039], pp.191, 244), so this is a supported methodological point.} \hldarkgreen{Kostenko requires that logs record decisions and key events and be stored for defined periods, extended automatically in case of disputes or inspections, and coupled to multi channel appeal mechanisms where execution of disputed automated decisions is suspended if significant harm to rights is at stake \parencites[191]{2977039}.}\pdfcomment[icon=Note]{The AI Law Model requires logs to record decisions/events, retain them for defined periods (extended in disputes/inspections), and provides for multi-channel appeal mechanisms with suspension of disputed automated decisions in cases of significant rights risk (AI Law Model [2977039], p.191).} \hldarkgreen{In parallel, the HITL grading case study shows what detailed logging looks like in practice: every AI and instructor decision is time stamped, linked to reviewer identifiers, and stored as an audit trail that can be inspected to understand override patterns and rubric level performance \parencites[16]{6805789}[17]{6805789}.}\pdfcomment[icon=Note]{The HITL grading case study documents timestamped decisions, reviewer identifiers and stored audit trails for inspecting override patterns (HITL case study [6805789], pp.16--17).} \hldarkgreen{During pre checks, auditors should therefore verify whether recommender platforms maintain comparably granular decision logs and whether affected users can access the materials on which relevant automated decisions are based, as required in labour contexts by the AI Law Model \parencites[191]{2977039}.}\pdfcomment[icon=Note]{The AI Law Model provides rights for affected workers to access materials and event logs and mandates multi-channel appeals; the sentence correctly applies those labour-context requirements to recommender audit pre-checks (AI Law Model [2977039], p.191).} \hllightgreen{Documentation quality also depends on internal governance instruments.}\pdfcomment[icon=Note]{Governance instruments and internal oversight are repeatedly tied to documentation quality in the literature (e.g., AI Law Model [2977039], and governance studies), so this is a reasonable synthesis.} \hldarkgreen{Mutawa et al. construct an Algorithmic Governance Assessment Matrix for zakat institutions in which transparency scores correlate with the clarity of algorithmic documentation and oversight mechanisms reported in annual reports and digital transformation audits \parencites[59]{2702055}.}\pdfcomment[icon=Note]{The Algorithmic Governance Assessment Matrix and its use to evaluate transparency and governance in zakat institutions is described in the study (Mutawa et al. [2702055], p.4).} \hldarkgreen{Asante and Adeoba, using Reflexive Thematic Analysis on CSR and ethical AI marketing literature, emphasise that trustworthiness in qualitative audit work requires a clearly documented audit trail of selection and analysis procedures \parencites[10]{8200246}.}\pdfcomment[icon=Note]{The Reflexive Thematic Analysis method and the need for audit trails to support trustworthiness and transparency in qualitative audit work are described in the RTA-based study (Asante \& Adeoba [8200246], p.23).} \hllightgreen{Analogously, documentation pre checks for social recommenders should confirm whether platforms maintain model cards, change logs and internal release notes describing date, responsible person, nature of changes, expected effects on accuracy and fairness, and results of revalidation tests, as required in Kostenko's model risk management provisions \parencites[191]{2977039}.}\pdfcomment[icon=Note]{The AI Law Model details model risk management practices including public/internal release notes (date, responsible person, description, expected effects, revalidation tests) (AI Law Model [2977039], p.191), so recommending model cards/change logs is a reasonable synthesis though "model cards" term is not explicitly used there.} \hllightgreen{These pre checks occur under a pronounced methodological limitation.}\pdfcomment[icon=Note]{The following text and cited literature make clear methodological limitations to pre-check based audits (e.g., lack of internal telemetry), so this framing is a supported inference from multiple sources.} \hllightgreen{Black box recommender algorithms used by social platforms are not directly observable in this thesis.}\pdfcomment[icon=Note]{The thesis and several cited empirical studies note that proprietary recommender internals are often not directly observable and that researchers rely on external documentation (see methodological notes and constraints across sources such as [6805789], [3195574]).} \hlorange{Khare et al. explicitly adopt a secondary qualitative design relying on regulatory documents, institutional reports and survey data to study dark patterns because direct access to proprietary systems is lacking \parencites[7]{6295677}.}\pdfcomment[icon=Note]{No provided source explicitly names Khare et al. or details their methodological choice regarding dark patterns and secondary qualitative design, so this specific attribution is unsupported by the supplied materials.} \hlorange{Bahangulu and Owusu Berko likewise argue that regulators struggle with opaque models and call for improved auditing techniques and collaboration between regulators and industry \parencites[9]{7461574}.}\pdfcomment[icon=Note]{The claim about Bahangulu and Owusu Berko is not substantiated by any of the provided sources; no supplied text contains that attribution or quoted argument.} \hllightgreen{The present audit toolkit follows the same constraint: documentation and data access are assessed through what platforms, regulators and independent researchers publish, not through internal telemetry.}\pdfcomment[icon=Note]{Several sources and the thesis fragments state that the toolkit assesses platforms via published materials rather than internal telemetry (methodological constraints discussed in HITL case study [6805789] and audit literature [3195574]), supporting this sentence.} \hllightgreen{Given this constraint, structured pre check tables become practical instruments.}\pdfcomment[icon=Note]{Given the documented constraint of relying on public materials, using structured pre-check tables is a reasonable and supported practical recommendation (inferred from audit methodology literature, e.g., [3195574] and AI Law Model [2977039]).} \hllightgreen{For each audited recommender, auditors can populate a grid that records existence and quality of algorithmic passports, event logs, appeal channels, and external audit histories, using criteria derived from TSAI, AI impact assessment proposals and governance matrices \parencites[244]{2977039}[15]{5442231}[59]{2702055}.}\pdfcomment[icon=Note]{Populating a grid with passport existence, logs, appeals and audit history using criteria from TSAI/AI Law Model and governance matrices is directly consistent with the AI Law Model requirements and governance frameworks (AI Law Model [2977039], p.244; governance matrix [2702055], p.4).} \hllightgreen{Where these artefacts are missing or clearly inadequate, the system is flagged in this report's working taxonomy as Code Category C: "Critical Traceability Failure," signalling that any subsequent thematic or scenario based analysis will rest on incomplete evidence and that, under AI Law Model expectations, such a system would not qualify for high risk deployment \parencites[244]{2977039}.}\pdfcomment[icon=Note]{Flagging systems with missing artefacts as critical traceability failures and noting that under the AI Law Model high‑risk systems must meet registration/documentation to deploy is consistent with AI Law Model provisions requiring registration/audit/trust marks for high-risk systems (AI Law Model [2977039], pp.244, 244).}

\subsection{Interface and Choice Architecture Audit Checklists}

\subsubsection{Assessment of Feed Mechanics and Session Design}
\label{sec:assessment_feed_mechanics_session_design} \hllightgreen{Feed mechanics sit at the heart of algorithmic addiction: infinite scroll, default auto play, and personalised landing views determine how easily users slip into extended sessions and how hard it is to stop.}\pdfcomment[icon=Note]{The role of feed mechanics (infinite scroll, autoplay, personalised landing) in promoting extended engagement and reduced stopping difficulty is supported by multiple sources discussing algorithmic engagement and addictive interfaces (e.g., 3303310 pp.2--3; 8958037 p.3), making this a reasonable synthesis.} \hldarkgreen{Short video platforms exemplify this pattern.}\pdfcomment[icon=Note]{Short video platforms are explicitly discussed as exemplifying algorithmic, feed-driven addictive dynamics in Li (8958037 p.2--3), which directly supports this statement.} \hldarkgreen{Li describes feeds in which micro duration clips of $15$--$60$ seconds are chained into an endless vertical stream, with each swipe loading the next item selected by a neural network based recommendation engine, creating a closed loop of "personalized $\rightarrow$ user need satisfaction $\rightarrow$ behavioural data accumulation $\rightarrow$ algorithm optimization $\rightarrow$ re push" \parencites[2]{8958037}[3]{8958037}.}\pdfcomment[icon=Note]{The quoted closed-loop description and micro-duration clips (15--60s) chained into an endless vertical stream are directly described by Li (8958037 p.3), matching the citation and phrasing.} \hldarkgreen{Hu characterises similar interfaces as "infinite scrolling" and "variable reward schedules" on platforms like TikTok and reports that users of such algorithmically curated feeds exhibit $22\%$ higher rates of compulsive checking compared with non algorithmic interfaces, with prolonged exposure correlating with reduced gray matter density and diminished prefrontal cortex activity in impulse control regions \parencites[2]{3303310}[20]{3303310}.}\pdfcomment[icon=Note]{The characterization of infinite scrolling and variable reward schedules and the empirical claims (22\% higher compulsive checking; reduced gray matter and diminished prefrontal activity) are reported in Hu (3303310 pp.2--3).} \hllightgreen{For an Ethical AI Audit Toolkit, these observations translate into explicit feed level checks.}\pdfcomment[icon=Note]{Inferring that these empirical observations imply feed-level checks for an ethical audit toolkit is a reasonable academic synthesis based on the harms and mechanisms discussed in sources like 3303310 and 8958037.} \hllightgreen{Session design should be reviewed for: \begin{itemize} \item default entry point (personalised feed versus neutral view); \item presence of infinite scroll or pagination; \item auto play behaviour after each item; \item existence and accessibility of non personalised feed options. \end{itemize} Me{\ss}mer and Degeling note that TikTok opens on the personalised "For You" page, and that YouTube's auto play function, enabled by default, starts a new video when the current one ends, reinforcing algorithmic curation even when the user has not explicitly requested additional content \parencites[42]{3195574}.}\pdfcomment[icon=Note]{The recommended session-design review items are reasonable operational checks derived from the feed mechanics literature; the specific regulatory concerns about autoplay and personalised entry are documented in sources (e.g., 8958037 p.2 referencing autoplay and personalised entry), supporting the checklist implication.} \hldarkgreen{Under the Digital Services Act, the Commission has already opened proceedings into TikTok's recommendation mechanisms and interface features, including infinite scrolling, auto play and push notifications, indicating that these mechanics can qualify as "systematic manipulative design" that impairs free and informed decision making \parencites[2]{8958037}[42]{3195574}.}\pdfcomment[icon=Note]{8958037 (p.2) explicitly states that the EU under the Digital Services Act launched an investigation into TikTok regarding recommendation mechanisms and interface features (infinite scrolling, autoplay, push notifications) as potentially manipulative, directly supporting this claim.} \hllightgreen{A feed mechanics checklist in this report's working taxonomy should therefore classify configurations as higher risk where: \begin{itemize} \item the app always launches directly into a personalised stream; \item infinite scroll is combined with default auto play; \item no clearly accessible non profiled recommendation option exists. \end{itemize} Short video usage metrics underline why such combinations matter.}\pdfcomment[icon=Note]{Classifying higher-risk configurations in a taxonomy and noting that short-video usage metrics make these combinations concerning is a reasonable synthesis grounded in the usage and harm metrics reported by Li (8958037 p.3).} \hldarkgreen{Li reports that Douyin once recorded average daily usage exceeding 21 sessions and nearly $60$ minutes, with $41.3\%$ of users engaging multiple times daily and $56.9\%$ using the app for over $60$ minutes, particularly among adolescents with relatively weak self control \parencites[3]{8958037}.}\pdfcomment[icon=Note]{Li provides the Douyin usage statistics cited (average daily usage >21 sessions, nearly 60 minutes, 41.3\% multiple times daily, 56.9\% over 60 minutes) and notes adolescent vulnerability (8958037 p.3), directly supporting the sentence.} \hldarkgreen{Hu synthesises evidence that hyper personalised feeds are associated with a $23\%$ higher incidence of depressive symptoms compared with chronological content, especially in adolescents and individuals prone to social comparison \parencites[20]{3303310}.}\pdfcomment[icon=Note]{Hu (3303310 p.3) reports longitudinal evidence that hyper-personalized feeds are associated with a 23\% higher incidence of depressive symptoms, particularly among adolescents and those prone to social comparison, matching this sentence.} \hllightgreen{These findings justify treating certain session designs as Code Category C in this report's working risk taxonomy: "addiction enabling feed mechanics." Audit instruments should not only flag harmful patterns but also document mitigating structures.}\pdfcomment[icon=Note]{Treating certain session designs as higher-risk categories (e.g., 'addiction enabling feed mechanics') and recommending audits that document mitigations is a reasonable policy inference synthesizing the empirical harms and audit recommendations in sources such as 3303310 and 8958037.} \hllightgreen{Human in the loop grading offers an example of how continuous streams can be structured without undermining autonomy.}\pdfcomment[icon=Note]{Presenting human-in-the-loop grading as an example of structuring continuous streams to preserve autonomy is a reasonable application of the HITL literature (Selvam et al., 6805789), which emphasizes human oversight and agency.} \hldarkgreen{Selvam and González Vallejo describe a four layer HITL architecture where AI generated scores are always subject to educator validation, dual logging records both AI and human actions, and student facing reports clearly tag feedback by source.}\pdfcomment[icon=Note]{Selvam et al. (6805789 pp.4,11) describe a layered HITL architecture with human validation, dual logging, and student-facing reports that tag feedback by source, directly supporting this description.} \hldarkgreen{Case study results show reduced grading time, $87\%$ of AI scores accepted with minor edits, $13\%$ fully overridden, and higher student trust attributable to transparent dual sourced feedback \parencites[6805789]{6805789}[19]{6805789}.}\pdfcomment[icon=Note]{The case study statistics (reduced grading time; 87\% accepted with minor edits; 13\% fully overridden) and higher student trust tied to transparent dual-sourced feedback are reported in Selvam et al. (6805789 pp.19, table).} \hllightgreen{Translating this pattern, social recommender audits should ask whether feeds expose explanation overlays (for example "why am I seeing this") and whether humans can intervene in ranking policies with traceable, logged effects.}\pdfcomment[icon=Note]{Translating HITL patterns into audit questions about explanation overlays and human intervention with logged effects is a reasonable and supported inference from the HITL transparency and logging features documented in 6805789.} \hllightgreen{Documentation pre checks are crucial because this thesis assesses black box algorithms only through secondary outputs.}\pdfcomment[icon=Note]{Stating that documentation pre-checks are crucial when the thesis assesses black-box algorithms via secondary outputs aligns with the methodological constraints discussed elsewhere in the cited literature (e.g., reliance on observable UX and regulatory materials in 3303310) and is a reasonable methodological claim.} \hldarkgreen{Kostenko's AI Law Model requires high risk systems to maintain technical passports, decision logs and human intervention procedures, and to submit to audits that can lead to suspension where ethical responsibility is breached \parencites[19.5]{2977039}[61--62]{2977039}.}\pdfcomment[icon=Note]{The AI Law Model described by Kostenko (2977039 pp.9--25) calls for transparency, logging, intervention procedures, and auditability measures (technical passports/decision logs and audit/suspension mechanisms), supporting this claim.} \hldarkgreen{Mutawa et al. demonstrate in zakat distribution that transparency scores correlate with the clarity of algorithmic documentation and governance oversight, based on annual reports and digital transformation audits \parencites[2702055]{2702055}.}\pdfcomment[icon=Note]{The zakat study (2702055 p.4) reports that transparency scores correlate with clarity of algorithmic documentation and governance oversight based on annual reports and audits, directly supporting this sentence.} \hllightgreen{Feed and session design audits must therefore begin by asking whether platforms can provide such passports and logs for their recommendation modules; if not, configurations with addictive mechanics should be flagged as "critical traceability failures" in the working taxonomy.}\pdfcomment[icon=Note]{Recommending that audits begin by requesting passports and logs and flag lack thereof as critical traceability failures follows logically from the regulatory and transparency requirements discussed in 2977039 and the documentation-transparency findings in 2702055, making this an acceptable synthesis.} \hllightgreen{This section's approach remains constrained by the absence of internal telemetry or model artefacts.}\pdfcomment[icon=Note]{Acknowledging that the approach is constrained by absence of internal telemetry or model artefacts is consistent with the methodological limitations noted in the literature on auditing opaque systems (e.g., 3303310) and is a reasonable, supported admission.} \hllightgreen{As in dark pattern research, where Khare et al. rely on regulatory documents and consumer surveys to characterise manipulation in the absence of proprietary access \parencites[7]{6295677}, this thesis evaluates feed mechanics and session design solely via documented UX behaviour, regulatory investigations and empirical studies of usage and harm \parencites[2]{8958037}[2]{3303310}[42]{3195574}.}\pdfcomment[icon=Note]{The comparison to dark-pattern research that relies on regulatory documents and surveys, and the claim that this thesis evaluates feed mechanics via documented UX, regulatory investigations and empirical studies, aligns with methods and constraints discussed in 3303310 and the referenced regulatory sources (e.g., 8958037/3303310), supporting this methodological statement.}

\subsubsection{Evaluation of Friction, Defaults and Opt-Out Channels}
\label{sec:evaluation_friction_defaults_optout} \hllightgreen{Defaults, friction, and opt out design jointly determine whether users can exercise meaningful control over algorithmic environments or remain trapped in engagement optimising states.}\pdfcomment[icon=Note]{Multiple sources link defaults, friction and opt-out design to user control and engagement optimisation (e.g., 6295677 pp.2,14; 6941611 p.11), so this is a reasonable synthesis of the cited literature.} \hldarkgreen{Dark pattern research in India shows that interface interference and obstruction based designs appear on approximately $70\%$ and $63\%$ of major platforms, with practices such as pre selected options, misleading hierarchies, and elongated cancellation flows used to exploit default bias and inertia \parencites[2]{6295677}.}\pdfcomment[icon=Note]{Source 6295677 (pp.2,4) explicitly reports interface interference ≈70\% and obstruction ≈63\% and documents the use of pre-selected options, misleading hierarchies and elongated cancellation flows.} \hldarkgreen{Survey evidence indicates that $55\%$ of consumers experienced unexpected charges and $47\%$ struggled to cancel subscriptions, while awareness of "dark patterns" remained low at $39\%$, revealing that friction and defaults systematically erode autonomy precisely where digital literacy is weakest \parencites[2]{6295677}[10]{9360054}.}\pdfcomment[icon=Note]{The survey figures (55\% unexpected charges, 47\% difficulty cancelling, 39\% awareness) are reported in the cited Indian study/summary (6295677 p.2) and related empirical summaries.} \hllightgreen{Audit preparation therefore needs explicit checks on how default configurations and friction costs are distributed across the journey.}\pdfcomment[icon=Note]{Recommendation to include explicit checks on default configurations and friction in audits follows directly from policy analyses urging preventive, design‑centric UX audits (6295677 pp.13--14).} \hldarkgreen{Hermansyah's analysis of the Digital Markets Act emphasises that behaviourally informed choice screens and interoperability mandates can reduce switching frictions and default bias, but only if enforcement can counter manipulative design practices that recreate lock in \parencites[10]{1070605}.}\pdfcomment[icon=Note]{Hermansyah's analysis of the DMA states that behaviourally informed choice screens and interoperability reduce switching frictions but their effectiveness depends on enforcement against manipulative designs (1070605 p.160).} \hldarkgreen{Khare et al. argue for bright line prohibitions, mandatory UX audits and algorithmic transparency, highlighting that complaint driven enforcement is structurally inadequate where sludge architectures make protective choices hard to execute \parencites[13]{6295677}[14]{6295677}.}\pdfcomment[icon=Note]{The cited policy study advocates bright‑line prohibitions, mandatory UX audits and transparency while arguing complaint‑driven enforcement is inadequate (6295677 pp.13--14).} \hlorange{This report's working taxonomy consequently treats combinations of asymmetric defaults and obstructive flows as Code Category C: "high risk friction and default configurations," particularly where they interact with addictive feed mechanics identified earlier.}\pdfcomment[icon=Note]{No provided source documents an external 'Code Category C' taxonomy or its exact label; this appears to be the report's internal classification without supporting citation.} \hllightgreen{A practical checklist can focus on three loci: consent and personalisation settings, subscription and account management, and notification controls.}\pdfcomment[icon=Note]{Sources discuss consent/personalisation, subscription/exit flows and notification controls as key loci of manipulation (6295677 pp.2,13; 7896734 pp.390--394), supporting this checklist breakdown.} \hldarkgreen{On consent, the Central Consumer Protection Authority's advisory under section 18(1) requires "clear and affirmative consent through explicit user action rather than pre checked defaults or implied agreement," signalling that pre selected tracking or personalisation options are unacceptable.}\pdfcomment[icon=Note]{The CCPA advisory and related rules require clear affirmative consent and disallow pre‑checked defaults; this is stated in 6295677 (p.12 referencing the advisory and Rule 4(9)).} \hllightgreen{Audit questions should ask whether cookie banners, recommendation settings, or profiling toggles use equal effort for accept and refuse, or instead present acceptance as a single click and refusal as a multi step path.}\pdfcomment[icon=Note]{Guidance on symmetry of effort for accept vs refuse aligns with principles and requirements for symmetrical consent/withdrawal actions found in the AI law model and UX audit recommendations (2977039 pp.223--224; 6295677 pp.13--14).} \hldarkgreen{Evidence from India shows that many platforms still rely on interface interference in consent flows despite self declarations of being "dark pattern free," with $81\%$ of platforms claiming compliance in self audits while national complaint resolution remains at $28\%$ \parencites[12]{6295677}.}\pdfcomment[icon=Note]{Source 6295677 (pp.12--13) documents that many platforms self‑declared dark‑pattern free (≈81\%) while national complaint resolution remained around 28\%, indicating persistence of interface interference.} \hllightgreen{Such discrepancies warrant Code C flagging even where formal notices exist.}\pdfcomment[icon=Note]{Flagging high‑risk behaviour despite formal notices is a reasonable inference from the documented discrepancy between self‑declarations and complaint outcomes (6295677 pp.12--14).} \hllightgreen{Subscription and exit flows require similar scrutiny.}\pdfcomment[icon=Note]{Multiple sources identify subscription/exit flows as manipulation points (roach motels/subscription traps) and call for scrutiny (6295677 pp.4,14; 7896734 p.390), supporting this claim.} \hldarkgreen{Roach motel and subscription trap patterns combine easy entry with difficult exit, using obstruction and forced action to increase friction for cancellation \parencites[4]{6295677}.}\pdfcomment[icon=Note]{The characterization of roach motel/subscription trap patterns as easy entry/difficult exit implemented via obstruction/forced action is documented in the dark‑patterns literature and in 6295677 (pp.4,14) and 7896734.} \hldarkgreen{Policy recommendations in the same study call for compulsory UI/UX audits and operational freezes when violations are detected, underscoring that regulators view obstructive exits as systemic risks rather than minor usability flaws \parencites[13]{6295677}.}\pdfcomment[icon=Note]{The same study (6295677 pp.13--14) explicitly calls for compulsory UI/UX audits and operational freezes until remediation when violations are found.} \hldarkgreen{Under the AI Law Model, manipulative interface practices "dark patterns" are explicitly prohibited in educational AI, with requirements that interfaces avoid intrusive notifications and forced options that make it difficult to refuse or perform human browsing \parencites[207]{2977039}.}\pdfcomment[icon=Note]{The AI Law Model explicitly prohibits manipulative interface practices in educational AI and requires avoidance of intrusive notifications and forced options (2977039 pp.207,224).} \hllightgreen{An audit toolkit aligned with these norms must therefore record whether account deletion, subscription termination or deactivation of personalised feeds can be completed in a small, symmetrical number of steps compared with sign up, and whether any penalties or hidden costs are attached to exit.}\pdfcomment[icon=Note]{Requiring audits to record symmetry of exit vs sign‑up steps and check for penalties follows from AI Law Model requirements for symmetry, deletion/withdrawal rights and absence of negative consequences (2977039 pp.223--224).} \hllightgreen{Notification controls are the third friction and default layer.}\pdfcomment[icon=Note]{Positioning notification controls as a distinct layer of friction/defaults is consistent with the earlier three‑locus checklist and discussions in the sources (6295677; 2977039 pp.223--224).} \hllightgreen{Khare et al. identify "nagging" via repeated prompts as a sludge pattern that wears down user resistance, particularly when deployed through app notifications \parencites[14]{6295677}.}\pdfcomment[icon=Note]{Repeated prompts or 'nagging' as a sludge/slowing pattern that erodes resistance, especially via app notifications, is discussed in the sludge/dark‑patterns literature and echoed in policy analyses (6295677 pp.14; 7896734 pp.390--394).} \hldarkgreen{Kostenko's age appropriate design principle sharply limits such practices for minors: manipulative and addictive mechanics including endless scrolling, autoplay, streaks and penalties for exiting are banned, and push notifications are to be "limited to a minimum" and only within reasonable need, with visible "no personalization" switches and simple notification limits.}\pdfcomment[icon=Note]{The AI Law Model's age‑appropriate design principle bans manipulative/addictive mechanics (endless scrolling, autoplay, streaks, penalties) and limits push notifications while requiring 'no personalization' switches (2977039 pp.223--224).} \hllightgreen{Pre checks should document whether platforms provide accessible toggles for disabling behaviourally targeted notifications, whether neutral informational messages are distinguished from engagement prompts, and whether refusal of notifications has any negative effect on core service access, which would contravene the AI Law Model's requirement that refusal of personalisation not degrade basic access to educational or governmental services \parencites[223]{2977039}.}\pdfcomment[icon=Note]{Documenting toggles for disabling targeted notifications and ensuring refusal does not degrade service aligns with AI Law Model provisions requiring no negative consequences for refusal and visible 'no personalization' controls (2977039 pp.223--224).} \hllightgreen{Institutional context also matters.}\pdfcomment[icon=Note]{The importance of institutional context (enforcement capacity, grievance mechanisms) is emphasized across the policy sources (6295677 pp.12--14; 2977039), making this a supported general point.} \hldarkgreen{Khare et al. show that voluntary self declarations under CCPA guidelines did not eliminate dark patterns and that reactive grievance mechanisms remain fragmented and dependent on consumer awareness \parencites[14]{6295677}.}\pdfcomment[icon=Note]{Source 6295677 (pp.12--14) documents that voluntary self‑declarations under CCPA did not eliminate dark patterns and that grievance mechanisms remain fragmented and dependent on consumer awareness.} \hldarkgreen{Hermansyah concludes that interoperability and DMA style choice architecture reforms only maintain contestability if accompanied by vigilant enforcement against manipulative design \parencites[10]{1070605}.}\pdfcomment[icon=Note]{Hermansyah's study states interoperability and DMA reforms improve contestability only if accompanied by vigilant enforcement against manipulative design (1070605 p.160).} \hllightgreen{The Ethical AI Audit Toolkit designed in this thesis operates under the same limitation: proprietary black box recommenders are evaluated via published UX audits, regulatory advisories and user survey data, not direct access to internal logs or code \parencites[7]{6295677}[9]{7461574}.}\pdfcomment[icon=Note]{The limitation that audits rely on published UX audits, advisories and survey data rather than internal logs is consistent with the practical constraints discussed in the literature (6295677 pp.13--14; 2977039 pp.206--207), though the specific [9]\{7461574\} citation is not provided here.} \hllightgreen{Within that constraint, defaults that silently enable profiling, friction loaded opt out channels and asymmetric notification controls emerge as high salience indicators of manipulation risk that must be coded and prioritised alongside feed mechanics in any Digital Fairness aligned audit \parencites[2]{6295677}[31]{9360054}[223]{2977039}.}\pdfcomment[icon=Note]{Sources identify profiling‑enabling defaults, frictioned opt‑outs and asymmetric notifications as salient manipulation indicators and recommend coding/prioritising them alongside feed mechanics (2977039 pp.223--224; 6295677 pp.2,14; 9360054 p.31).}

\subsubsection{Classification of Nudging Practices along an Ethics Spectrum}
\label{sec:classification_nudging_spectrum} \hldarkgreen{Interface elements that guide users along a journey need not be problematic; dark commercial patterns emerge once guidance secures compliance by degrading comprehension, autonomy, transparency, or reversibility.}\pdfcomment[icon=Note]{Source 6087241 (Conceptual review) explicitly states that persuasive interfaces are not inherently problematic but become dark when they secure compliance by weakening comprehension, autonomy, transparency, or reversibility (p.3).} \hllightgreen{Liu and Keshavarz explicitly distinguish online choice architecture, which unavoidably shapes attention and action, from dark commercial patterns, which "steer, deceive, pressure, or obstruct consumers in ways that impair informed and autonomous choice" \parencites[3]{6087241}.}\pdfcomment[icon=Note]{Source 6087241 distinguishes online choice architecture (shaping attention/action) from dark commercial patterns that 'steer, deceive, pressure, or obstruct' consumers; the quoted language and distinction are present though exact author attribution is drawn from that conceptual review (6087241, p.3).} \hllightgreen{In this report's working taxonomy, we therefore classify nudging practices along an ethics spectrum with two core dimensions: support for preference consistent choice and presence of manipulation indicators such as information distortion, interface asymmetry, social or temporal pressure, choice obstruction, and forced action.}\pdfcomment[icon=Note]{The integrative taxonomy and two-dimension framing (firm conversion vs. preservation of informed/preference-consistent choice, plus categories like information distortion, interface asymmetry, social/temporal pressure, obstruction, forced action) are articulated in 6087241 (pp.2-4) and reflected in the report's working taxonomy.} \hllightgreen{At the ethically acceptable end sit "assistive nudges" that clarify options without skewing effort asymmetrically.}\pdfcomment[icon=Note]{6087241 describes ethically acceptable or assistive nudges as designs that guide choice while preserving conditions for informed preference-consistent decisions (p.3), supporting the claim about 'assistive nudges' clarifying options without asymmetric effort.} \hldarkgreen{Examples include neutral rankings that surface relevant alternatives while still displaying lower priced options, or consent dialogs where "accept" and "reject" require comparable steps and use equivalent visual prominence.}\pdfcomment[icon=Note]{The examples of neutral rankings and symmetric consent dialogs (symmetric effort/prominence for accept/reject) are explicitly discussed in the conceptual and managerial guidance (table and text referencing symmetric effort and clear costs) in sources such as 6087241 and 8200246 (6087241 p.33; 8200246 p.16).} \hldarkgreen{Liu and Keshavarz argue that ethical digital marketing aims to persuade while preserving the consumer's ability to make an informed, preference consistent choice; designs that improve navigation yet maintain clarity and reversibility fall into this category \parencites[33]{6087241}.}\pdfcomment[icon=Note]{6087241 explicitly states that ethical digital marketing seeks to persuade while preserving consumers' ability to make informed, preference-consistent choices and that designs improving navigation while maintaining clarity and reversibility fit this category (p.3).} \hllightgreen{Under this report's spectrum they are coded as low risk nudges that can support autonomy and trust when combined with transparent CSR governance, cross functional AI ethics committees, and regular algorithm audits recommended by Asante and Adeoba \parencites[16]{8200246}.}\pdfcomment[icon=Note]{The coding of low-risk nudges and the governance recommendations (transparent CSR governance, cross-functional AI ethics committees, regular algorithm audits) are supported by 8200246 (CSR as governance; recommends cross-functional committees and audits, p.20232) and the taxonomy rationale in 6087241, so this is a reasonable synthesis.} \hllightgreen{Mid spectrum practices are "borderline nudges" that exploit some cognitive tendencies but retain basic reversibility.}\pdfcomment[icon=Note]{6087241 and related literature describe 'borderline' or mid-spectrum practices that exploit cognitive tendencies while retaining reversibility, a characterization consistent with the cited conceptual review (6087241, pp.2-4).} \hldarkgreen{Consent banners that make acceptance slightly easier than refusal or product pages that visually highlight default options more strongly than alternatives fall here.}\pdfcomment[icon=Note]{Multiple sources (6087241; 6087241 p.33 and related consent-interface references) document consent banners and product pages where acceptance is made slightly easier or defaults are visually highlighted, placing those examples in the mid-spectrum category.} \hldarkgreen{Dark pattern crawls of approximately $11{,}000$ shopping websites identified many cases where cookie notices structured choices asymmetrically, making acceptance easier than refusal, yet still technically offering a decline path.}\pdfcomment[icon=Note]{Source 6087241 reports a large-scale crawl of approximately 11,000 shopping websites that identified many dark-pattern instances and notes that cookie notices often structure choices asymmetrically, making acceptance easier though decline remains technically possible (p.33).} \hldarkgreen{From a marketing perspective, Liu and Keshavarz conceptualise these as information distortion or interface asymmetry that may raise conversion in the short term but risk eroding perceived fairness and trust over time \parencites[33]{6087241}.}\pdfcomment[icon=Note]{6087241 frames such asymmetric designs as information distortion or interface asymmetry and argues they can boost short-term conversion while risking perceived fairness and trust erosion over time (pp.2-3).} \hllightgreen{In this toolkit they are coded as medium risk nudges that require remediation plans and monitoring of complaint, cancellation and churn metrics rather than immediate prohibition.}\pdfcomment[icon=Note]{6087241 recommends evaluating medium-risk practices via remediation and monitoring (trust, complaint, churn metrics) rather than blanket prohibition, so coding them as medium risk requiring remediation and monitoring is a defensible synthesis (pp.33--34).} \hldarkgreen{At the unacceptable end lie "coercive sludge patterns" that combine exploitative nudging with artificial friction.}\pdfcomment[icon=Note]{Sources (6295677 and 6087241) describe 'sludge' or coercive patterns that combine exploitative nudging with artificial friction (obstructive cancellation flows, forced actions), supporting this characterization (6295677 pp.14, 7; 6087241 pp.2-4).} \hldarkgreen{Khare et al. synthesise regulatory and institutional evidence showing that approximately $87\%$ of major Indian platforms deploy at least one dark pattern, with drip pricing, false urgency, confirm shaming, roach motels, forced action and basket sneaking as recurring motifs \parencites[7]{6295677}.}\pdfcomment[icon=Note]{Source 6295677 reports synthesised regulatory/institutional evidence indicating approximately 87\% of major Indian platforms deploy at least one dark pattern and lists recurring motifs such as drip pricing, false urgency, confirm-shaming, roach motels, forced action, and basket sneaking (pp.7, 12, 14).} \hldarkgreen{These practices are intentionally embedded strategies that exploit System 1 heuristics such as default bias, scarcity and social pressure while blocking System 2 reflection through complexity and delay \parencites[4]{6295677}.}\pdfcomment[icon=Note]{6295677 and 6087241 explicitly link dark-pattern practices to exploitation of System 1 heuristics (default bias, scarcity, social pressure) and to blocking System 2 reflection via complexity and delay (6295677 pp.14, 7; 6087241 pp.2-4).} \hldarkgreen{Interface interference and obstruction in cancellation flows, subscription traps and forced account creation exemplify sludge by increasing effort costs and discouraging user protective behaviour \parencites[4]{6295677}[14]{6295677}.}\pdfcomment[icon=Note]{The characterization of cancellation-flow interference, subscription traps, and forced account creation as examples of sludge that increase effort costs and discourage protective behaviour is documented in 6295677 (pp.14, 7) and reflected in the dark-pattern literature summarized in 6087241.} \hllightgreen{Under this report's working taxonomy, these patterns are labelled Code Category C: "Prohibited or Critical Risk Nudges," since they align with manipulative design that EU Digital Fairness discussions and AI law models seek to eliminate \parencites[390]{7896734}[61--62]{2977039}.}\pdfcomment[icon=Note]{The labeling of these patterns as prohibited/critical risk (Code Category C) aligns with regulatory ambitions to eliminate manipulative designs in EU 'digital fairness' discussions and AI law models; sources like 7896734 and 2977039 discuss EU prohibitions and regulatory removal of manipulative practices, supporting the synthesis (7896734 p.7; 2977039 pp.39--44).} \hldarkgreen{A concise matrix helps operationalise this ethics spectrum inside audits: \begin{table}[h] \centering \caption{Working ethics spectrum for nudging practices in interface and choice architecture} \begin{tabular}{p{3cm}p{4.2cm}p{4.2cm}} \hline Category (working label) & Typical design attributes & Indicative sources / risk signals \\ \hline Assistive nudges (low risk) & Symmetric effort for accept / decline, clear total costs, visible lower priced alternatives, easy reversal & Preserves informed, preference consistent choice; aligns with ethical digital marketing and CSR based AI governance \parencites[33]{6087241}[16]{8200246} \\ Borderline nudges (medium risk) & Mild interface asymmetry, moderate highlighting of defaults, incomplete but non deceptive urgency messages & Asymmetric cookie banners where refusal is possible but less salient; potential conversion trust trade off \parencites[33]{6087241} \\ Coercive sludge patterns (Code C) & Hidden fees, false scarcity, confirm shaming, roach motels, obstructive cancellation, forced action & Documented as structural features on $\approx 87\%$ of major platforms; linked to autonomy erosion and regulatory calls for bright line bans and UX audits \parencites[7]{6295677}[14]{6295677}[390]{7896734} \\ \hline \end{tabular} \end{table} Regulatory analysis reinforces this spectrum.}\pdfcomment[icon=Note]{The table's working ethics spectrum and its source attributions (e.g., CSR governance [8200246], the ≈87\% prevalence [6295677], regulatory links [7896734]) are directly supported by the cited sources (8200246 p.16; 6295677 pp.7,12; 7896734 p.7; 6087241 pp.2-4).} \hldarkgreen{Quinelato concludes that effective protection against persuasive and obscure practices on social media requires limiting dark patterns and ensuring valid, explicit consent for data processing, framed within DSA provisions and complementary regimes \parencites[21]{6941611}.}\pdfcomment[icon=Note]{Source 6941611's conclusion emphasizes the need for regulatory measures to limit dark patterns and to ensure valid, explicit consent mechanisms within the context of DSA analysis (p.21), matching this claim.} \hllightgreen{Sanikidze interprets the planned Digital Fairness Act as targeting "addictive product designs and unfair personalisation," complementing DSA and DMA bans on interface pressure and the AI Act's prohibition of subliminal or harmful manipulation \parencites[390]{7896734}.}\pdfcomment[icon=Note]{Source 7896734 discusses a planned Digital Fairness Act targeting addictive designs and unfair personalisation and positions it as complementary to DSA/DMA and the AI Act's prohibitions on subliminal/harmful manipulation (p.7), supporting the interpretation though the specific attribution to 'Sanikidze' is a synthesis.} \hldarkgreen{In parallel, Kostenko's confidentiality principle forbids using personal data for advertising, targeted marketing or profiling aimed at influence that restricts or distorts freedom of choice without explicit, informed consent \parencites[44]{2977039}.}\pdfcomment[icon=Note]{2977039 (AI Law Model) explicitly states that using personal data for commercial purposes (including targeted marketing or profiling) or for influence that restricts or distorts freedom of choice is prohibited without explicit, informed consent (pp.43--44).} \hllightgreen{These texts provide legal backing for treating Code C patterns as not merely undesirable but incompatible with emerging EU style digital fairness norms.}\pdfcomment[icon=Note]{The cited regulatory and legislative texts (e.g., 2977039, 7896734, 6295677) collectively support treating high-risk Code C patterns as incompatible with emerging EU-style digital fairness norms, so this is a reasonable synthesis (2977039 pp.39--44; 7896734 p.7).} \hllightgreen{Methodologically, this classification remains constrained by the absence of proprietary logs or source code for social recommender algorithms.}\pdfcomment[icon=Note]{Several sources note methodological constraints due to lack of access to proprietary logs/source code for recommender systems and algorithmic opacity (e.g., 8200246 limitation of secondary analysis; 3303310 and others discuss algorithmic opacity), supporting this limitation claim.} \hllightgreen{Following Khare et al., who rely on qualitative content analysis of regulatory reports, institutional audits and consumer surveys to map dark patterns, this thesis places nudges on the ethics spectrum based solely on documented interface behaviour, reported consumer harm and explicit regulatory evaluations, not on any internal telemetry \parencites[7]{6295677}[10]{9360054}.}\pdfcomment[icon=Note]{6295677 (Khare et al. style) and related studies employ qualitative content analysis of regulatory reports and institutional audits; the claim that the thesis places nudges on the spectrum based on documented interface behaviour and regulatory evaluations rather than internal telemetry aligns with that methodological approach (6295677 pp.3--7; 9360054 noted as similar secondary-analysis limitation).} \hllightgreen{Complementary empirical work suggests that triangulating managerial accounts with large-scale user generated evidence and a transparent scoring rubric can make such externally observable classifications more reliable; this approach may therefore help distinguish between intentional manipulation and incidental design friction in practice \parencites[5]{4499345}.}\pdfcomment[icon=Note]{Empirical work (e.g., 4499345) demonstrates that triangulating managerial accounts with large-scale user-generated evidence and transparent scoring rubrics improves reliability of externally observable classifications, supporting the suggested approach as reasonable (4499345 pp.5--12).}

\section{Pro-Social Design Strategies and Business Model Transformation}

\subsection{Pro-Social Choice Architecture Patterns}

\subsubsection{Finite Session Design and Break Prompts}
\label{sec:finite_session_break_prompts} \hldarkgreen{Designing finite sessions and break prompts for social recommenders must respond directly to evidence that infinite scrolling and auto play correlate with compulsive checking, depressive symptoms and diminished impulse control.}\pdfcomment[icon=Note]{Source 3303310 presents empirical claims linking infinite scrolling/auto play to compulsive checking and mental-health harms and discusses diminished impulse control and dopamine-driven feedback loops (Proceedings CONF-CDS 2025, p.2).} \hllightgreen{Short video ecosystems that couple $15$--$60$ second clips with swipe based infinite feeds and default auto play create a continuous sequence of personalised stimuli.}\pdfcomment[icon=Note]{Source 3303310 discusses short-video platform affordances (infinite feeds, auto play and variable reward schedules) as continuous personalised stimuli, supporting the claim though exact clip length ranges are not the primary focus (Proceedings CONF-CDS 2025, p.2).} \hllightgreen{Users describe being drawn into an "auto pilot mode," and empirical work links this configuration to diminished academic concentration, weakened social skills and heightened anxiety and depression in college populations \parencites[2]{8958037}.}\pdfcomment[icon=Note]{Qualitative and empirical sources (e.g., 2580223 on youth developmental impacts and 3303310 on 'auto pilot' and concentration/mental-health links) report users' experiences and associations with reduced concentration and increased anxiety/depression, supporting this synthesis (Frontiers in Psychology 2025; CONF-CDS 2025, p.2).} \hldarkgreen{Parallel analysis of algorithmically curated social feeds reports $22\%$ higher rates of compulsive checking and a $23\%$ higher incidence of depressive symptoms in personalised streams compared with chronological content, together with neuroimaging evidence of reduced prefrontal activity and gray matter density in impulse control regions \parencites[2]{3303310}[20]{3303310}.}\pdfcomment[icon=Note]{Source 3303310 explicitly reports a 22\% higher rate of compulsive checking in algorithmic feeds and links algorithmic curation to depressive outcomes and neuroimaging evidence of reduced prefrontal activity / gray-matter correlates in impulse-control regions (Proceedings CONF-CDS 2025, p.2).} \hllightgreen{These findings suggest that finite session design and explicit break prompts should be treated as core mitigations rather than optional usability tweaks.}\pdfcomment[icon=Note]{This is a reasonable policy inference from the empirical harms and regulatory concerns documented in sources such as 3303310 and regulatory/model-law texts (synthesis across CONF-CDS 2025 and AI Law Model 2977039).} \hllightgreen{Regulatory and model law sources reinforce this orientation.}\pdfcomment[icon=Note]{Multiple regulatory and model-law documents in the corpus (e.g., 2977039, 7896734) address manipulative designs and log/audit requirements, supporting the claim that legal sources reinforce the orientation toward mitigation.} \hllightgreen{The European Commission's proceeding against TikTok under the Digital Services Act explicitly singles out infinite scrolling, auto play and push notifications as potential "systematic manipulative design" that can impair users' capacity for free and informed decisions \parencites[2]{8958037}.}\pdfcomment[icon=Note]{Source 7896734 and related regulatory analyses describe European Commission investigations and concerns about addictive platform designs (including features like infinite scrolling and autoplay) under DSA-related enforcement, supporting this claim (DARK PATTERNS IN THE AGE OF AI, p.7).} \hllightgreen{Kostenko's AI Law Model goes further for children: endless scrolling, auto play, streaks, forced quests and penalties for exiting are prohibited in systems interacting with minors, and feeds for teenage users are normatively illustrated as operating in "detox mode" by default, without endless scrolling or auto played video and with visible "no personalisation" switches and simple explanations such as "why am I seeing this?" \parencites[223]{2977039}[221]{2977039}.}\pdfcomment[icon=Note]{The AI Law Model (2977039) imposes strict safeguards for minors and prohibits manipulative techniques and risky scenarios; the cited model text (Kostenko) describes heightened protections for minors consistent with the listed prohibitions and default 'detox' orientation (AI Law Model, pp.213--229).} \hllightgreen{Within this report's working taxonomy, finite session design is therefore not only a pro social choice architecture pattern but also a direct response to articulated duties of ethical responsibility and digital safety.}\pdfcomment[icon=Note]{Given the empirical harms and the duties articulated in model-law and regulatory sources (e.g., 3303310; 2977039), treating finite-session design as both pro-social and responsive to ethical duties is a supported synthesis.} \hlorange{Auditable finite session patterns can be articulated along three dimensions.}\pdfcomment[icon=Note]{No provided source explicitly frames 'auditable finite session patterns' into exactly 'three dimensions'; this is an internal structuring claim without direct citation support in the supplied materials.} \hlorange{First, temporal framing: sessions can be bounded by soft limits that surface timely, neutral prompts once cumulative viewing or scrolling passes a threshold, for example explaining that a user has been in a personalised feed for a given period and presenting a neutral option to pause.}\pdfcomment[icon=Note]{The specific design prescription about 'soft limits that surface timely, neutral prompts' is a normative recommendation not directly documented in the cited sources as an established practice; no source explicitly provides this exact framing.} \hldarkgreen{Evidence from digital literacy surveys indicates that users feel confident in basic tasks yet score low on recognising hidden manipulations and avoiding fraud, with average ratings of $2.31$ and $2.13$ respectively \parencites[31]{9360054}.}\pdfcomment[icon=Note]{Source 9360054 reports digital literacy survey averages of 2.31 for recognising hidden manipulations and 2.13 for avoiding fraud, directly supporting these figures (Věda a perspektivy, p.31).} \hllightgreen{Break prompts that foreground time spent without using scarcity language or confirm shaming can therefore support reflective awareness rather than reactance.}\pdfcomment[icon=Note]{This recommendation is a reasonable inference from behavioral-literacy findings (9360054) and behavioral design literature that neutral, non-shaming prompts support reflection rather than reactance (synthesis across literacy and behavioral-regulation sources).} \hllightgreen{Second, structural exit cues: scrollable feeds can be designed with explicit "natural exits," such as periodic hard stops or interposed configuration panels, instead of endless continuation.}\pdfcomment[icon=Note]{Designing 'natural exits' and hard stops is a plausible translation of critiques of endless continuation found in the literature on dark patterns and preventive design (e.g., 6295677 and 3303310 discuss preventing endless continuation), supporting this design suggestion as synthesis.} \hldarkgreen{In gaming and education, Tokarieva and Chyzhykova show that well designed gamification uses systemic structures to foster intrinsic motivation, with failure treated as a data point rather than a penalty and progression organised through clear stages rather than continuous streams \parencites[9]{5645864}.}\pdfcomment[icon=Note]{Source 5645864 (gamification meta-analyses) states that gamification can convert extrinsic to intrinsic motivation and that failure is reframed as a data point with progression organized through stages, directly supporting this claim (ALFRED NOBEL UNIVERSITY JOURNAL, p.9).} \hllightgreen{Translating this pattern, finite blocks of content with clear transitions and opportunities to disengage align with human centred design without eliminating engagement altogether.}\pdfcomment[icon=Note]{This is a reasonable synthesis: translating gamification findings (5645864) into finite-content blocks aligns with human-centred design principles discussed across the corpus.} \hllightgreen{Third, integration with governance and logging.}\pdfcomment[icon=Note]{Integration with governance and logging is a theme in the AI Law Model (2977039) and other governance sources that emphasize auditability, supporting this as a listed dimension.} \hldarkgreen{Human in the loop grading demonstrates how session like sequences of AI support can be bounded and audited: every AI decision is subject to educator confirmation, all actions are dual logged, and override patterns are periodically analysed by institutional committees \parencites[20]{6805789}.}\pdfcomment[icon=Note]{Source 6805789 documents a HITL grading framework where AI outputs are human-validated, audit-logged, and used with oversight committees, directly supporting the described bounded and auditable session-like sequences (p.20).} \hldarkgreen{For social recommenders, finite session triggers and break prompts should similarly be logged as events, with technical passports documenting thresholds, prompt content and observed user responses in a way that regulators and independent auditors can inspect, consistent with AI Law Model requirements on event logging, fairness reporting and supervision \parencites[19.5]{2977039}[228]{2977039}.}\pdfcomment[icon=Note]{Source 2977039 requires logging, reporting, and auditability (model passports, event logs) and the sentence aligns directly with those provisions (AI Law Model, pp.213--229 and p.228).} \hllightgreen{This is particularly relevant under a Digital Fairness trajectory that aims to regulate addictive product designs and unfair personalisation alongside DSA duties on non profiled feeds and systemic risk assessment \parencites[390]{7896734}[2248]{6502223}.}\pdfcomment[icon=Note]{Source 7896734 and related regulatory analyses discuss proposed 'Digital Fairness' initiatives to regulate addictive designs and personalisation alongside DSA obligations, supporting this contextual claim (DARK PATTERNS IN THE AGE OF AI, p.7).} \hllightgreen{Methodological constraints remain: proprietary ranking algorithms and internal A/B tests are not observable in this project, so finite session design and break prompts are assessed through documented harms, regulatory signals and governance models rather than direct telemetry \parencites[7]{6295677}[9]{7461574}.}\pdfcomment[icon=Note]{Several sources note limitations due to lack of access to proprietary ranking algorithms and internal A/B tests and justify relying on documented harms and regulatory signals (e.g., methodological notes in 6295677 and project limitations cited across the corpus), supporting this methodological claim.} \hldarkgreen{It is worth noting that recommender models which incorporate user-generated text and sequential context may produce more extreme confidence scores and can increase the likelihood of exposing vulnerable users to controversial or harmful communities, a pattern observed in large-scale Reddit modelling experiments \parencites[4]{8020829}.}\pdfcomment[icon=Note]{Source 8020829 shows that recommender models using user-generated text and sequential context produce more extreme confidence scores and substantially increase recommendations to controversial Reddit communities, directly supporting this sentence (Engage corpus paper, pp.3--5).} \hllightgreen{Such empirical tendencies further motivate treating session limits and transparent prompts as auditable safeguards rather than optional interface features \parencites[4]{8020829}.}\pdfcomment[icon=Note]{This is a reasonable inference: empirical tendencies described in 8020829 and harms/regulatory concerns (3303310, 2977039) justify treating session limits and transparent prompts as safeguards.} \hldarkgreen{Even within this limitation, evidence on dopamine driven feedback loops, depressive outcomes and regulatory bans on endless scrolling and auto play provides a clear rationale for embedding finite session constraints and transparent break prompts as baseline audit criteria for pro social recommender redesign \parencites[2]{8958037}[2]{3303310}[223]{2977039}[390]{7896734}.}\pdfcomment[icon=Note]{This sentence synthesises empirical evidence on dopamine feedback loops and depressive outcomes (3303310) together with regulatory/model-law actions banning manipulative features (2977039; 7896734), which are all present in the cited sources and support the normative rationale.}

\subsubsection{Batch Feeds, Digest Modes and Time-Boxed Usage}
\label{sec:batch_feeds_digest_timeboxed} \hllightgreen{Batch oriented feeds and digest modes offer a structural counterpoint to infinite scroll and default auto play.}\pdfcomment[icon=Note]{Sources discuss infinite scroll/auto-play harms and the need for design alternatives (e.g., 8958037 on infinite swipe and auto-play; 3303310 on engagement features), so framing batch/digest as a structural counterpoint is a reasonable synthesis of the literature (8958037 p.2--3; 3303310 p.1--3).} \hllightgreen{Short video ecosystems illustrate why such alternatives are needed.}\pdfcomment[icon=Note]{Short-video ecosystems and their harmful engagement mechanics are documented (e.g., addiction loop, micro-duration clips in 8958037), supporting the claim that these ecosystems illustrate the need for alternatives (8958037 p.2--3).} \hldarkgreen{Neural network based recommendation engines combined with 15--60 second clips, infinite swipe feeds and auto play induce an "auto pilot mode" of persistent scrolling, with Douyin users averaging more than 21 sessions and nearly 60 minutes per day and over half exceeding 60 minutes, especially adolescents with comparatively weak self control \parencites[2]{8958037}[3]{8958037}.}\pdfcomment[icon=Note]{Dan Li explicitly describes neural-network recommendation combined with 15--60s clips, infinite swipe and autoplay inducing an "auto-pilot mode" and gives Douyin usage statistics (over 21 sessions, \textasciitilde{}60 minutes, 56.9\% >60 minutes, adolescents more vulnerable) (8958037 p.2--3).} \hldarkgreen{Hyper personalised feeds on broader social platforms correlate with a 23\% higher incidence of depressive symptoms and reduced prefrontal activity relative to chronological streams, and algorithmically curated interfaces generate 22\% more compulsive checking than non algorithmic counterparts \parencites[2]{3303310}[20]{3303310}.}\pdfcomment[icon=Note]{The proceedings study reports a 23\% higher incidence of depressive symptoms and reduced prefrontal activity for hyper-personalized feeds, and notes algorithmic interfaces generate \textasciitilde{}22\% more compulsive checking compared to non-algorithmic feeds (3303310 p.2--3).} \hllightgreen{Within this report's working taxonomy, continuous, unsegmented streams therefore fall into high risk engagement mechanics that pro social design must replace with bounded, time aware patterns.}\pdfcomment[icon=Note]{The report's risk taxonomy and calls to counteract manipulative practices support classifying continuous, unsegmented streams as high-risk engagement mechanics that pro-social design should replace (2977039 p.8; 2977039 p.98).} \hllightgreen{Finite batches and digest modes can be framed as concrete operationalisations of rights articulated in AI governance texts.}\pdfcomment[icon=Note]{AI governance texts in the provided sources articulate rights and principles (e.g., digital neutrality, accessibility to algorithmic education) that can be operationalised via finite batches and digest modes, making the framing plausible (2977039 p.98, p.108, p.149).} \hlorange{Kostenko's digital ecology chapter recognises an inalienable right to a "psycho emotionally safe digital environment" and prohibits algorithms that cause psychological pressure, impose behavioural patterns, restrict freedom of choice, form addiction or create self replicating digital environments that erode critical thinking \parencites[119]{2977039}.}\pdfcomment[icon=Note]{No provided source names a "Kostenko" chapter or contains the exact quoted phrasing; while 2977039 prohibits manipulative algorithms, the specific attribution to Kostenko and the exact list of prohibitions is not supported by the supplied sources.} \hlorange{For minors, the principle of digital safety of the child specifies that feeds should operate in "detox mode" by default, explicitly "without endless scrolling, autoplay of the video and aggressive return triggers," and with visible "no personalisation" switches and simple explanations such as "why am I seeing this?" \parencites[223]{2977039}[221]{2977039}.}\pdfcomment[icon=Note]{The specific quoted language about minors' feeds operating in "detox mode" with exact phrases and switches is not found in the provided sources; 2977039 addresses child protections generically but does not contain the cited verbatim requirements (no matching text in provided pages).} \hllightgreen{Translating these prescriptions, batch feeds present a finite set of items with explicit end points and natural exit cues instead of unending content streams, while digest modes aggregate content into periodic summaries that can be consumed in defined sessions.}\pdfcomment[icon=Note]{This is a reasonable synthesis/inference from discussions of batch/digest alternatives versus endless streams and the described mechanics of short-video platforms (8958037 on unending streams and micro-duration clips; general UX implications in 3303310).} \hllightgreen{Evidence on digital literacy and awareness of EU regulation implies that such bounded formats need to be combined with clear signalling.}\pdfcomment[icon=Note]{The study on digital literacy and EU regulation awareness documents gaps and implies that bounded formats need clear signalling to be effective, supporting the inferential claim (9360054 p.10).} \hldarkgreen{Survey results from Chukhray and Koval show that highly active social network users report strong confidence in everyday digital tasks and online purchases (Likert means $4.40$--$4.62$) yet score much lower on recognising hidden manipulations (mean $2.31$) and avoiding fraud ($2.13$), and remain sceptical about the effectiveness of GDPR, the DSA and the AI Act in curbing disinformation and fraud \parencites[10]{9360054}.}\pdfcomment[icon=Note]{The survey results and reported Likert means (4.40--4.62 for confidence; 2.31 for recognising hidden manipulations; 2.13 for avoiding fraud) and scepticism about GDPR/DSA/AI Act effectiveness are explicitly reported in the source (9360054 p.10).} \hllightgreen{Digest modes that simply compress engagement into fewer sessions without altering manipulative features would not address this "second level divide." Instead, batching and time boxing need to be accompanied by transparent indicators of session boundaries and neutral language explaining that no content is being hidden for punitive reasons, consistent with the AI Law Model's requirement for explainability and user friendly communication \parencites[68]{2977039}[69]{2977039}.}\pdfcomment[icon=Note]{9360054 discusses the "second-level divide" and 2977039 requires explainability and user-friendly communication, so the claim that digest modes alone would not address the divide and need explainability is a supported synthesis (9360054 p.10; 2977039 p.8, p.29).} \hllightgreen{Session design also interacts with legal risk tiers.}\pdfcomment[icon=Note]{2977039 includes a risk taxonomy and legal risk tiers for AI systems; the statement that session design interacts with legal risk tiers is a reasonable inference from that material (2977039 p.8, p.242--243).} \hldarkgreen{Under the AI Law Model's technological sovereignty principle, systems that cannot be independently audited or that embed hidden mechanisms for influencing behaviour may be prohibited and entered into a register of unwanted AI systems \parencites[108]{2977039}.}\pdfcomment[icon=Note]{The AI Law Model describes prohibitions and registration/notification mechanisms for systems that violate principles such as digital neutrality or conceal manipulative mechanisms, matching this claim (2977039 p.108; 2977039 p.242--243).} \hllightgreen{Social recommendations configured as endless, opaque streams that foster addictive use without clear logging of session boundaries and prompts could fall into this category.}\pdfcomment[icon=Note]{3303310 documents that endless, opaque algorithmic streams foster addictive use and higher compulsive checking, and 2977039 sets out criteria for prohibited systems, supporting that such recommendations could fall into that category (3303310 p.1--3; 2977039 p.108).} \hllightgreen{By contrast, batch feeds and digest modes aligned with algorithmic passports, event logs and documented human oversight channels provide evidence for compliance with TSAI duties such as preserving logs, reporting incidents and ensuring human supervision \parencites[242]{2977039}.}\pdfcomment[icon=Note]{2977039 describes algorithmic passports, event logs, and TSAI duties (preserving logs, incident reporting, human oversight), so aligning batch/digest modes with these compliance artifacts is a supported synthesis (2977039 p.242--243).} \hlorange{In this thesis, such configurations are classified in the working Code C taxonomy as "time bounded, audit compatible sessions," distinct from "addiction enabling feed mechanics." Audit preparation steps therefore need to include explicit questions on batching and time boxing.}\pdfcomment[icon=Note]{This is a descriptive claim about how the thesis classifies configurations and what audit preparation should include; the specific working Code C taxonomy and required audit steps are not documented in the provided external sources.} \hlorange{Inspired by Mutawa et al.'s governance matrix, which evaluates transparency and accountability in zakat AI through document analysis and performance indicators \parencites[59]{2702055}, this report's toolkit includes pre checks such as: whether the platform documents maximum recommended session lengths, whether break prompts are triggered after defined usage thresholds, whether users can shift from continuous feeds to digest modes, and whether these choices are recorded in logs accessible to regulators under instruments like DSA Article 40 \parencites[19]{3195574}.}\pdfcomment[icon=Note]{The sentence attributes inspiration to "Mutawa et al.'s governance matrix" and lists specific toolkit pre-checks and linkage to DSA Article 40; the provided sources do not include Mutawa nor the exact checklist or a direct citation to DSA Article 40 in the supplied materials, so the detailed claims lack support in the given sources.} \hllightgreen{Given the black box constraint acknowledged earlier, auditors rely on published UX behaviour and regulatory documentation rather than direct telemetry, but can still assess the presence of visible exit cues and digest options in live interfaces \parencites[7]{6295677}[9]{7461574}.}\pdfcomment[icon=Note]{Sources on audit practice and human-AI collaboration indicate auditors often rely on published UX/regulatory documentation rather than telemetry and can assess visible cues in live interfaces, supporting this claim (7461574 p.10; 6539584 on human-AI audit roles).} \hllightgreen{Finally, finite session and digest designs can be framed as pro social choice architecture that still accommodates business interests.}\pdfcomment[icon=Note]{Hermansyah and related sources argue that behaviourally informed regulation can reduce default bias while preserving innovation, supporting the claim that finite session/digest designs can be framed as pro-social choice architecture compatible with business interests (1070605 p.160).} \hldarkgreen{Hermansyah's analysis of DMA interoperability mandates shows that behaviourally informed regulation can reduce default bias and consumer lock in while preserving innovation and competition, provided manipulation via dark patterns is countered through vigilant enforcement \parencites[10]{1070605}.}\pdfcomment[icon=Note]{Hermansyah's analysis explicitly concludes that behaviourally informed regulation (e.g., DMA interoperability and choice architecture) mitigates default bias and lock-in while preserving innovation, with the caveat that enforcement must counter dark patterns (1070605 p.160).} \hllightgreen{Batch and digest configurations that maintain recommendation quality yet avoid endless scrolling and auto play can be positioned as analogous reforms: they maintain engagement, but within time boxed, transparent sessions that respect emerging Digital Fairness expectations and AI Law Model rights to cyber neutrality and non alienation of digital subjectivity \parencites[148]{2977039}[390]{7896734}.}\pdfcomment[icon=Note]{Sources discuss Digital Fairness and emerging EU rights and the AI Law Model's principles (7896734 on Digital Fairness; 2977039 on rights), so positioning batch/digest as analogous reforms that respect these expectations is a supported synthesis (7896734 p.390; 2977039 p.148--149).} \hllightgreen{Empirical work on immersive digital environments suggests a cautious note: strong immersion can both soothe and encourage prolonged disengagement from the offline context, so session boundaries that are explicit and gently enforced may be more protective than merely informative nudges \parencites[8]{9529573}.}\pdfcomment[icon=Note]{Empirical VR research documents that strong immersion can both soothe and encourage disengagement from offline context and warns about over-immersion, supporting the recommendation for explicit, enforced session boundaries (9529573 p.3--8).} \hllightgreen{This research appears to support the idea that design choices which couple containment cues with clear exit affordances are likely to reduce the risk of over-immersion while preserving the experiential value of recommended content \parencites[8]{9529573}.}\pdfcomment[icon=Note]{Kutsal \& Turan and related sources indicate that combining containment cues with clear exit affordances is likely to mitigate over-immersion while preserving experiential value, supporting this interpretive conclusion (9529573 p.7--8).}

\subsubsection{User Control Dashboards for Algorithmic Preferences}
\label{sec:user_control_dashboards} \hllightgreen{Designing user control dashboards for algorithmic preferences requires treating transparency and agency as enforceable rights rather than cosmetic settings.}\pdfcomment[icon=Note]{This is a reasonable synthesis of normative emphasis in the AIN and AI Law Model: AIN stresses transparency and stakeholder protections (2742041, p.2) and the AI Law Model frames human-centred rights and enforceable obligations (2977039, p.22), supporting the inference about treating transparency/agency as enforceable rather than cosmetic.} \hldarkgreen{The AI Law Model formulates a human centred principle where people must be able to know about the impact of AI, obtain explanations, and maintain control, and it prohibits full automation of critical decisions without human intervention.}\pdfcomment[icon=Note]{Directly supported by the AI Law Model which states the human-centred principle including the ability to know impacts, obtain explanations, maintain control and prohibition of full automation of critical decisions (2977039, p.22).} \hldarkgreen{Transparency and explainability are defined as access to technical documentation, logs, justifications, and application limits, with all AI systems obliged to inform users about their algorithmic nature and to make results available for auditing \parencites[22]{2977039}.}\pdfcomment[icon=Note]{Text of the AI Law Model explicitly defines transparency and explainability as access to technical documentation, logs, justifications, application limits and informing users of algorithmic nature and auditability (2977039, p.22).} \hllightgreen{Dashboards become the primary interface through which these obligations surface for lay users.}\pdfcomment[icon=Note]{Sources on transparency-by-design and interface affordances indicate that user-facing dashboards are the natural channel for surfacing obligations to lay users (3303310, p.3; 5142478, p.5), so this is a reasonable synthesis.} \hllightgreen{Stakeholder oriented governance frameworks reinforce this orientation.}\pdfcomment[icon=Note]{The AIN and related governance literature emphasize stakeholder-oriented governance and engagement across lifecycle stages (2742041, p.2), which reinforces the orientation toward stakeholder-facing controls, reasonable synthesis.} \hldarkgreen{The Algorithmic Impact Navigator positions organisations as the first line of accountability and calls for comprehensive impact evaluations, mapping of regulatory duties, and stakeholder engagement throughout the lifecycle, with transparency, fairness and respect for autonomy as constant principles across sectors.}\pdfcomment[icon=Note]{The AIN text frames organisations as the first line of accountability and calls for comprehensive impact evaluations, mapping regulatory duties and stakeholder engagement, with transparency, fairness and respect for autonomy as constant principles (2742041, pp.2--3).} \hllightgreen{Within such a frame, user control dashboards function as operational tools for implementing the AIN: they materialise choices about profiling, recommendation parameters and explanation depth that might otherwise remain buried in internal documentation \parencites[2]{2742041}.}\pdfcomment[icon=Note]{Combines AIN's operational guidance about organisational implementation (2742041, p.2) with literature on dashboards as operational interfaces (3303310, p.3), a reasonable inference that dashboards can materialise choices otherwise buried in documentation.} \hllightgreen{Regulatory initiatives on AI driven content curation and platform governance point toward concrete dashboard capabilities.}\pdfcomment[icon=Note]{Regulatory and research sources on platform governance and content curation point toward concrete capabilities (e.g., transparency, opt-outs, control parameters), supporting the claim (6502223, p.12; 5142478, p.7).} \hldarkgreen{The EU AI Act's risk based framework requires platforms to ensure that algorithms are transparent, accountable and free from bias and to undergo regular audits, while also emphasising user control over recommendations, including the ability to opt out of personalised feeds \parencites[8]{5142478}.}\pdfcomment[icon=Note]{The EU AI Act is described as a risk-based framework requiring transparency, measures to address bias and audits and emphasises controls such as opt-out from personalised feeds in the literature (5142478, p.7; 6502223, p.12).} \hldarkgreen{The Digital Services Act obliges very large online platforms to provide at least one recommendation option not based on profiling and to publish transparency reports, with vetted researchers receiving data access to study systemic risks \parencites[12]{6502223}.}\pdfcomment[icon=Note]{The DSA provisions for very large online platforms (VLOPs) require at least one recommendation option not based on profiling, publication of transparency reports, and vetted researcher data access to study systemic risks (6502223, p.12).} \hllightgreen{Any meaningful dashboard in this context must therefore expose a non profiled feed mode, present it as a first class choice rather than a buried setting, and log user selections so that auditors can verify that the option is more than formal compliance \parencites[12]{6502223}[21]{3195574}.}\pdfcomment[icon=Note]{This is an inferential prescription combining DSA's non-profiled recommendation requirement (6502223, p.12) with AI Law Model logging/audit obligations (2977039, pp.45,135), so recommending first-class presentation and logging is a supported synthesis.} \hllightgreen{Media diversity and bias work adds further functional requirements.}\pdfcomment[icon=Note]{Work on media diversity and algorithmic bias identifies additional functional requirements for systems (e.g., diversity metrics, fairness measures), supporting this brief claim (5142478, pp.5--8).} \hldarkgreen{Gerr's assessment of AI driven curation argues that XAI models can mitigate harms by giving users clear explanations of why specific content is recommended and by allowing them to critically evaluate their feeds and avoid manipulation \parencites[8]{5142478}.}\pdfcomment[icon=Note]{Gerr explicitly argues that XAI can mitigate harms by giving users comprehensible explanations enabling critical evaluation of recommended content and avoidance of manipulation (5142478, p.8).} \hldarkgreen{Diversity and bias metrics were applied to curated datasets from Facebook, YouTube and TikTok to measure content homogeneity and ideological skew \parencites[5]{5142478}.}\pdfcomment[icon=Note]{Gerr reports that diversity and bias metrics were applied to curated datasets from Facebook, YouTube and TikTok to measure homogeneity and ideological skew (5142478, p.5).} \hllightgreen{Dashboards that expose simple controls such as "diversity weighting" sliders or toggles for chronological versus engagement ranked ordering, combined with explanation panels that state which factors prior viewing, trending status, geography drove a recommendation, translate these analytical concepts into user facing control points \parencites[3]{3303310}[8]{5142478}.}\pdfcomment[icon=Note]{Several sources describe interface controls (diversity weighting, chronological vs engagement ordering) and explanation panels listing factors like prior viewing or trending status, supporting the claim as a synthesis (3303310, p.3; 6502223, p.16; 5142478, p.7).} \hllightgreen{Newer governance proposals for democratic recommendation systems sharpen the link between dashboards and design standards.}\pdfcomment[icon=Note]{New governance proposals for democratic recommendation systems (e.g., technical design standards in 6502223) link dashboards to design standards, so this is a supported synthesis (6502223, pp.16--17).} \hldarkgreen{Chitra et al. recommend value aligned objective functions with explicit diversity incentives and engagement ceilings, and insist that recommendation systems provide user accessible explanations identifying primary factors behind each item.}\pdfcomment[icon=Note]{The recommendations enumerated (value-aligned objectives, diversity incentives, engagement ceilings, and user-accessible explanations identifying primary factors) appear explicitly in the democratic recommendation proposals (6502223, p.16).} \hldarkgreen{Their framework also calls for mandatory algorithmic auditing and research data access, including representative samples of recommendation sequences, so that external analysts can assess political amplification asymmetries and systemic risks \parencites[16]{6502223}.}\pdfcomment[icon=Note]{The same framework calls for mandatory algorithmic auditing and provision of research data access including representative samples of recommendation sequences to study systemic effects (6502223, p.16).} \hllightgreen{In this report's working taxonomy, user control dashboards that expose diversity controls, engagement limit indicators, and a clear "why am I seeing this" explanation channel are treated as evidence of alignment with such emerging norms rather than as mere usability features.}\pdfcomment[icon=Note]{This is an interpretive claim about the report's taxonomy; given the cited normative materials that treat such dashboard features as governance-aligned (5142478, p.7; 2977039, p.22; 2742041, p.2), the interpretation is a reasonable synthesis.} \hllightgreen{AI governance standards underline that dashboards must sit on top of auditable backends.}\pdfcomment[icon=Note]{AI governance documents emphasize auditable backends and logging as prerequisites for transparency and oversight (2977039, pp.45,135), supporting the claim that dashboards must sit atop auditable systems.} \hlorange{Sankaran shows that ISO/IEC standards on AI risk management, bias assessment and ethical considerations are intended to support AI Act obligations by providing measurable criteria for fairness, transparency and human oversight, yet rely on organisations to implement internal logging and documentation \parencites[2]{8589961}.}\pdfcomment[icon=Note]{The sentence attributes a finding to 'Sankaran' about ISO/IEC standards; no provided source by that author is included, and while standards and organizational logging needs are discussed elsewhere (7461574; 2977039, p.45), the specific attribution to Sankaran is unsupported by the supplied materials.} \hldarkgreen{The AI Law Model requires event logging that links inputs, intermediate calculations and outputs, and mandates that logs be preserved and accessible for audit and appeal \parencites[44]{2977039}[135]{2977039}.}\pdfcomment[icon=Note]{The AI Law Model requires detailed logging, preservation and accessibility of logs for audit and appeal, as discussed in the model's provisions on obligations and burden of proof (2977039, pp.135; 2977039, p.45).} \hllightgreen{Dashboards that present explanation and preference controls without underlying log architectures risk becoming performative; audit preparation must therefore verify whether setting changes and recommendation rationales are recorded in machine readable form so regulators and independent auditors can reconstruct decision paths \parencites[21]{3195574}[44]{2977039}.}\pdfcomment[icon=Note]{The AI Law Model's logging and audit requirements (2977039, pp.45,135) support the inference that dashboards without underlying log architectures risk being performative and that audits should verify machine-readable records, a reasonable synthesis.} \hllightgreen{HITL deployments demonstrate how such logging can support human authority and user trust.}\pdfcomment[icon=Note]{Empirical work on human-in-the-loop deployments and the proctoring literature indicate that human validation and logging can support authority and trust (2055956, p.9), making this a supported inference.} \hlorange{In AI assisted essay grading, instructors receive SHAP based explanations alongside model scores, can override or confirm outputs, and all actions are dual logged, with students seeing clearly which feedback is AI generated and which is instructor reviewed \parencites[20]{6805789}.}\pdfcomment[icon=Note]{The specific scenario (SHAP explanations in essay grading, dual logging, and student-visible provenance) is not documented in the provided sources; 2055956 discusses proctoring and trust but does not provide the specific SHAP/dual-logging grading workflow described.} \hllightgreen{Students reported higher confidence when human validation was visible, and accreditation bodies can inspect logs for quality assurance \parencites[20]{6805789}[2055956]{2055956}.}\pdfcomment[icon=Note]{2055956 discusses increased user trust when human involvement and transparency are present in proctoring systems, and the broader literature notes audit/QA uses of logs, so the claim about higher confidence and inspection for quality assurance is a reasonable synthesis (2055956, p.9).} \hllightgreen{User control dashboards for social recommenders that similarly differentiate automated from human curation, expose override events, and provide appeal pathways align with the AI Law Model's human centred and responsibility principles \parencites[22]{2977039}[61--62]{2977039}.}\pdfcomment[icon=Note]{The AI Law Model's human-centred and responsibility principles (2977039, p.22) combined with literature on transparency and logging support the assertion that dashboards differentiating automated vs human curation and exposing overrides/appeals align with those principles (2977039, pp.22,45).} \hlorange{A final constraint is methodological.}\pdfcomment[icon=Note]{This is a transitional claim ('A final constraint is methodological'); no specific methodological constraint is cited here in the provided sources, so the sentence as a standalone claim lacks direct support.} \hllightgreen{Black box recommender algorithms used by commercial platforms are not directly observable in this project.}\pdfcomment[icon=Note]{Multiple sources document platform opacity (e.g., TikTok opacity, black-box recommenders) and note limits on direct observability, supporting the project limitation claim (6502223, p.16; 3303310, p.2).} \hllightgreen{As in dark pattern research and AI proctoring reviews, assessment of user control dashboards must rely on publicly documented UX behaviour, regulatory proceedings and audit reports rather than proprietary telemetry or code \parencites[7]{6295677}[9]{2055956}[9]{7461574}.}\pdfcomment[icon=Note]{The methodological approach of relying on publicly documented UX behaviour, regulatory proceedings and audit reports rather than proprietary telemetry is consistent with documentary-analysis approaches discussed in the literature (2055956, p.9; 8453044, p.11).} \hllightgreen{Within that limitation, this section treats dashboards that combine non profiled feed selection, intelligible explanations, documented logging, and visible human oversight as key pro social choice architecture patterns and as central artefacts for any Ethical AI Audit Toolkit aligned with AI Act, DSA and AI Law Model principles \parencites[12]{6502223}[8]{5142478}[22]{2977039}[2]{2742041}.}\pdfcomment[icon=Note]{This is a synthesis: given DSA/AI Act/AI Law Model requirements (6502223, p.12; 5142478, p.7; 2977039, p.22) the listed dashboard features (non-profiled feeds, explanations, logging, human oversight) are plausibly treated as core artefacts, supported by the cited materials.}

\subsection{Recommender System Redesign for Wellbeing}

\subsubsection{Diversity-Oriented Ranking and Content Pluralism}
\label{sec:diversity_ranking_pluralism} \hldarkgreen{Evidence on engagement optimised curation shows that current ranking objectives narrow informational variety.}\pdfcomment[icon=Note]{Gerr explicitly documents that engagement‑optimised curation narrows informational variety and diminishes exposure to diverse viewpoints (5142478, p.8; p.30).} \hldarkgreen{Analyses of Facebook, YouTube and TikTok document that algorithmic personalisation steers users toward content aligned with pre existing beliefs and systematically diminishes exposure to diverse viewpoints, reinforcing echo chambers and limiting opportunities for critical thinking and democratic participation \parencites[5]{5142478}[29]{5142478}.}\pdfcomment[icon=Note]{Gerr's analyses of Facebook, YouTube and TikTok report that algorithmic personalisation steers users toward content aligned with pre‑existing beliefs and reduces exposure to diverse viewpoints (5142478, p.8; p.30).} \hldarkgreen{Case studies during elections and health crises find that YouTube's watch time driven recommendation system amplified sensationalist vaccine videos, while Facebook disproportionately surfaced far right political content, illustrating how accuracy and attention goals crowd out pluralism \parencites[7]{5142478}.}\pdfcomment[icon=Note]{Gerr's case studies document YouTube's watch‑time optimisation amplifying sensationalist vaccine videos and Facebook amplifying far‑right political content during elections (5142478, p.30; p.9).} \hllightgreen{These patterns position diversity oriented ranking as a necessary redesign axis for wellbeing oriented recommender systems rather than an optional add on.}\pdfcomment[icon=Note]{Sources that document engagement biases and propose diversity interventions (5142478; 6502223) reasonably support treating diversity‑oriented ranking as a necessary redesign axis rather than optional.} \hllightgreen{Regulatory and governance proposals already sketch what pluralism aware objectives might entail.}\pdfcomment[icon=Note]{Several works sketch pluralism‑aware objectives and governance proposals (6502223; 5142478; 2977039), so the claim that such proposals already exist is supported as a synthesis of these sources.} \hldarkgreen{Chitra and colleagues call for "value aligned objective functions" in democratic recommendation systems: accuracy weights for fact checked sources, explicit intra session diversity requirements that penalise sequences exposing users exclusively to a narrow ideological spectrum, and an engagement ceiling that caps amplification for individual items to avoid virality premiums that incentivise emotional escalation \parencites[16]{6502223}.}\pdfcomment[icon=Note]{6502223 explicitly calls for 'value‑aligned objective functions' including accuracy weighting for fact‑checked sources, intra‑session diversity requirements, and engagement ceilings (6502223, p.16).} \hldarkgreen{Gerr argues that diversity and pluralism benchmarks, supported by metrics such as Shannon entropy and Simpson's diversity index over curated streams, should become binding design requirements, not only transparency obligations, because technical compliance with DSA reporting has not reduced engagement driven biases on YouTube and TikTok \parencites[8]{5142478}[28]{5142478}.}\pdfcomment[icon=Note]{Gerr argues for binding diversity and pluralism benchmarks and notes that technical compliance with DSA transparency has not eliminated engagement‑driven biases; the report also uses diversity indices like Shannon/Simpson in its methods (5142478, p.3; p.30).} \hllightgreen{In this report's working taxonomy, these proposals define a "diversity aligned objective" class of rankers that explicitly trade some engagement for content pluralism.}\pdfcomment[icon=Note]{Combining the policy proposals and technical taxonomy in the cited work (6502223; 5142478) supports defining a 'diversity aligned objective' class of rankers as a working taxonomy in this report.} \hldarkgreen{Fairness research in recommender systems provides complementary insight into how diversity and equity interact.}\pdfcomment[icon=Note]{The survey literature on fairness in recommender systems provides complementary insight on diversity and equity interactions, as reviewed in 6577872 (6577872, p.28--33).} \hldarkgreen{Deldjoo et al. distinguish exposure and effectiveness as utilities: exposure captures how much visibility items or item groups receive, while effectiveness measures how closely exposure matches user preferences \parencites[28]{6577872}.}\pdfcomment[icon=Note]{6577872 explicitly distinguishes exposure and effectiveness as two types of utility in recommender fairness (6577872, p.28).} \hldarkgreen{For consumers, fairness often means equal effectiveness across groups so that historically disadvantaged users are not systematically offered less relevant items.}\pdfcomment[icon=Note]{6577872 states that from the consumers' perspective fairness often concerns equitable distribution of effectiveness across users (6577872, p.28).} \hldarkgreen{For producers, fairness focuses on exposure equity so that creators from minority countries or unpopular segments are not perpetually hidden by popularity biased algorithms.}\pdfcomment[icon=Note]{6577872 describes producers' concerns as primarily about exposure equity, to ensure visibility for less popular creators or those from minority regions (6577872, p.28).} \hldarkgreen{Metrics such as the Gini index for catalog concentration, group utility differences and relevance disparity quantify when a system over concentrates exposure on a small subset of items or producers \parencites[30]{6577872}.}\pdfcomment[icon=Note]{6577872 discusses the use of the Gini index for concentration and metrics like group utility differences and relevance disparity to quantify exposure concentration (6577872, p.30; p.33).} \hllightgreen{Although these metrics were developed primarily for fairness, they can be repurposed in an audit toolkit as quantitative indicators of content pluralism: lower Gini coefficients and reduced group utility gaps signal more diversified ranking outcomes.}\pdfcomment[icon=Note]{It is a reasonable academic inference from the fairness metrics literature (6577872) that metrics such as Gini and group utility gaps can be repurposed to indicate content pluralism outcomes (6577872, p.30--33).} \hllightgreen{Designing diversity oriented ranking involves trade off management rather than simple maximisation.}\pdfcomment[icon=Note]{The literature highlights trade‑offs between fairness, diversity and accuracy rather than simple maximisation, supporting the claim that diversity ranking requires trade‑off management (6577872, p.28; p.33).} \hldarkgreen{Deldjoo et al. warn that computational metrics may oversimplify real world problems: promoting long tail items without checking their quality risks recommending irrelevant content, and there is often an assumed but unexamined trade off between customer fairness, provider fairness and overall accuracy \parencites[28]{6577872}[30]{6577872}.}\pdfcomment[icon=Note]{6577872 warns that computational metrics can oversimplify real problems (e.g., promoting long‑tail items without quality checks) and notes assumed trade‑offs among customer fairness, provider fairness and accuracy (6577872, p.28--33).} \hllightgreen{Chitra et al. similarly acknowledge that capping virality and enforcing diversity bands will reduce short term engagement and associated revenue, which creates structural misalignment between platform incentives and democratic goals \parencites[14]{6502223}.}\pdfcomment[icon=Note]{6502223 and related governance analyses note that capping virality or enforcing diversity is likely to reduce short‑term engagement and revenues, creating misalignment between platform incentives and democratic goals (6502223, p.14--16).} \hllightgreen{An Ethical AI Audit for social recommenders must therefore treat diversity not as a free benefit but as a strategic rebalancing: some engagement metrics will worsen, yet system level wellbeing measures, such as reduced emotional volatility and lower incidence of depressive symptoms associated with hyper personalised feeds, may improve \parencites[20]{3303310}.}\pdfcomment[icon=Note]{Empirical and conceptual sources link engagement‑optimised feeds to harms and suggest wellbeing improvements from diversification and guardrails, so treating diversity as strategic rebalancing with mixed engagement/wellbeing effects is a supported inference (3303310, p.3; 5142478, p.30).} \hllightgreen{Governance structures give diversity oriented ranking practical teeth.}\pdfcomment[icon=Note]{The cited policy and governance sources describe audit and regulatory mechanisms that would give diversity‑oriented ranking practical enforcement, supporting this claim (2977039; 6502223; 5142478).} \hldarkgreen{Gerr calls for mandatory algorithmic auditing that evaluates performance against diversity and pluralism benchmarks, alongside transparency about ranking criteria and data access for vetted researchers to study systemic risks \parencites[3]{5142478}[28]{5142478}.}\pdfcomment[icon=Note]{Gerr explicitly calls for mandatory algorithmic auditing against diversity and pluralism benchmarks, plus transparency and researcher data access (5142478, p.30; p.26).} \hldarkgreen{The same study recommends applying distributional fairness metrics and diversity indices to curated datasets from Facebook, YouTube and TikTok to detect biases favouring specific ideological perspectives \parencites[5]{5142478}.}\pdfcomment[icon=Note]{Gerr recommends applying distributional fairness metrics and diversity indices to curated datasets from Facebook, YouTube and TikTok to detect ideological biases (5142478, p.3; p.26).} \hldarkgreen{Chitra et al. propose that annual independent audits for very large platforms assess evidence of political amplification asymmetries and effectiveness of misinformation mitigation, with methodologies and findings published in aggregated form \parencites[16]{6502223}.}\pdfcomment[icon=Note]{6502223 proposes annual independent audits for large platforms assessing political amplification asymmetries and misinformation mitigation, with aggregated publication of methodologies and findings (6502223, p.16).} \hllightgreen{These mechanisms align with AI Law Model requirements for technical passports, event logging and public registers for high risk systems and provide a concrete institutional route for enforcing pluralism goals \parencites[19.5]{2977039}.}\pdfcomment[icon=Note]{AI Law Model recommendations include technical passports, event logging and public registers for high‑risk systems, aligning with the audit and transparency mechanisms proposed in the other sources (2977039, p.95; p.193).} \hllightgreen{Within the pro social pivot strategy in this thesis, diversity oriented ranking and content pluralism therefore become one of the key redesign levers.}\pdfcomment[icon=Note]{Given the preceding evidence and policy proposals in the cited literature, positioning diversity‑oriented ranking as a key redesign lever within a pro‑social pivot is a reasonable synthesis (5142478; 6502223; 6577872).} \hldarkgreen{Audit questions grounded in the cited work ask whether ranking objectives include diversity terms, whether exposure and effectiveness metrics indicate disproportionate concentration, and whether independent audits systematically apply entropy, Simpson index and group utility measures across ideological and cultural categories.}\pdfcomment[icon=Note]{The suggested audit questions and measures, checking for diversity terms in objectives, exposure/effectiveness concentration, and use of entropy, Simpson and group utility measures, are directly reflected in the cited methodological and policy sources (5142478, p.3; 6577872, p.30--33; 6502223, p.16).} \hllightgreen{Given the methodological limitation that proprietary algorithms are observed only through secondary outputs, these metrics and governance artefacts serve as the primary evidence for assessing whether social recommender systems remain locked in engagement worship, or genuinely shift toward pluralistic, wellbeing compatible curation under an anticipated Digital Fairness regime \parencites[2]{3303310}[8]{5142478}[16]{6502223}[30]{6577872}.}\pdfcomment[icon=Note]{Sources note the limitation of observing proprietary algorithms only via outputs and argue that metrics plus governance artefacts (audits, passports) become primary evidence for assessing shifts toward pluralistic curation (5142478, p.3; 6577872, p.28; 6502223, p.14).}

\subsubsection{Downranking of Harmful and Sensational Content}
\label{sec:downranking_harmful_sensational} \hllightgreen{Reducing the visibility of high arousal and harmful material requires shifting ranking objectives away from pure engagement and toward pluralism, epistemic quality and wellbeing.}\pdfcomment[icon=Note]{Technical design proposals in Source 6502223 p.16 and the governance analysis in Source 2977039 p.32 argue for replacing pure engagement objectives with value-aligned objectives (epistemic quality, wellbeing, diversity), supporting this synthesis.} \hllightgreen{Analyses of major platforms illustrate why.}\pdfcomment[icon=Note]{Multiple platform analyses in Sources 6502223 p.16, 6502223 p.14 and 5142478 p.6 present empirical and policy evidence that illustrate the problems described, so this linking statement is reasonably supported.} \hldarkgreen{Facebook's engagement based algorithm was internally assessed as a "significant contributor" to political polarisation, with $64\%$ of extremist group joins traced to recommendation prompts.}\pdfcomment[icon=Note]{Source 6502223 p.16 explicitly cites leaked Facebook internal research concluding the engagement-based algorithm was a "significant contributor" to polarization and that 64\% of extremist group joins were traced to recommendations.} \hldarkgreen{YouTube's recommender systematically amplifies right leaning political content relative to equivalent left leaning material even after controlling for organic popularity, while TikTok accounts that show interest in weight loss or politics are quickly steered toward eating disorder and extremist content \parencites[2244]{6502223}.}\pdfcomment[icon=Note]{Source 6502223 p.16 cites Huszár et al. on YouTube amplifying right-leaning mainstream content after controls and reports on TikTok steering newly created accounts toward eating-disorder and extremist content.} \hllightgreen{These findings link engagement optimisation directly to the amplification of sensationalist and harmful streams.}\pdfcomment[icon=Note]{Sources 3303310 p.3 and 5142478 p.6 connect engagement-optimising recommender architectures to amplification of sensationalist and harmful streams, supporting this inference.} \hllightgreen{Psychological and neurocognitive evidence points in the same direction.}\pdfcomment[icon=Note]{Neurocognitive and psychological evidence linking algorithmic exposure to mental-health effects is presented in Source 3303310 pp.2--3, so this general statement is supported.} \hldarkgreen{Collaborative filtering architectures that weight prior interactions systematically prioritise outrage and sensationalism; on Instagram, the "Explore" page increased exposure to body negative imagery by $23\%$ compared with chronological feeds and correlated with a significant decline in self esteem among adolescent women of color \parencites[3]{3303310}.}\pdfcomment[icon=Note]{Source 3303310 p.3 documents collaborative filtering amplifying high-arousal content and reports Instagram's Explore increased exposure to body-negative imagery by 23\% with correlated self-esteem decline among adolescent women of color.} \hldarkgreen{Algorithmically prioritised content shows $31\%$ higher negative affect than baseline, and hyper personalised feeds are associated with a $23\%$ higher incidence of depressive symptoms and reduced prefrontal activity in impulse control regions \parencites[3]{3303310}[20]{3303310}.}\pdfcomment[icon=Note]{Source 3303310 p.3 reports algorithmically prioritised content showing \textasciitilde{}31\% higher negative affect and cites hyper-personalized feeds associated with a 23\% higher incidence of depressive symptoms and reduced prefrontal activity.} \hllightgreen{Within this report's working taxonomy, content that combines high negative affect, ideologically polarising frames and documented mental health harms is treated as a prime target for downranking.}\pdfcomment[icon=Note]{The sources (e.g., 6502223 p.16, 5142478 p.6) advocate prioritising downranking of content combining negative affect, polarising frames and mental-health harms; this sentence is a reasonable application of that taxonomy within the report.} \hllightgreen{Technical and governance proposals already outline levers for such downranking.}\pdfcomment[icon=Note]{Technical and governance levers for downranking (value-aligned objectives, diversity incentives, engagement ceilings, audits) are outlined in Source 6502223 p.16 and 2977039 p.32, supporting this claim.} \hllightgreen{Chitra et al. recommend "value aligned objective functions" for democratic recommendation systems that explicitly add accuracy weighting, intra session diversity requirements and engagement ceilings to ranking loss functions.}\pdfcomment[icon=Note]{The specific recommendation for value-aligned objective functions with accuracy weighting, intra-session diversity and engagement ceilings appears in Source 6502223 p.16 (though attribution to 'Chitra et al.' is not present in provided sources).} \hldarkgreen{Accuracy weighting gives positive bonuses to content from sources with high fact checking scores and negative weights to sources with misinformation patterns, while diversity incentives penalise sequences that expose users exclusively to a narrow ideological spectrum.}\pdfcomment[icon=Note]{Source 6502223 p.16 explicitly describes accuracy weighting (boosting fact-checked sources, penalising misinformation) and diversity incentives penalising ideologically narrow sequences.} \hldarkgreen{Engagement ceilings cap amplification of individual items, directly constraining virality premiums that currently reward emotional escalation.}\pdfcomment[icon=Note]{Source 6502223 p.16 explicitly defines engagement ceilings as caps on amplification of individual items to prevent a virality premium that incentivises emotional escalation.} \hllightgreen{These design moves amount to structurally downranking harmful and sensational content by reducing its learned utility in ranking optimisation.}\pdfcomment[icon=Note]{Given the mechanisms in Source 6502223 p.16 (reducing weights/bonuses and capping amplification), it is a reasonable inference that these design moves structurally downrank harmful sensational content by lowering its learned utility.} \hllightgreen{A concise mapping of downranking levers and their audit implications is shown in Table~\ref{tab:downranking_levers}. \begin{table}[h] \centering \caption{Working mapping: downranking levers for harmful and sensational content} \label{tab:downranking_levers} \begin{tabular}{p{3.3cm}p{4.8cm}p{4.3cm}} \hline Lever & Mechanism & Audit focus \\ \hline Accuracy weighting & Boosts fact checked sources, penalises repeat misinformation & Evidence of source level weights; fact checker integration \parencites[2252]{6502223} \\ Diversity incentives & Penalises ideologically narrow sequences & Diversity metrics (entropy, Simpson index) over recommendation sessions \parencites[2252]{6502223}[6]{5142478} \\ Engagement ceilings & Caps exposure of single items & Detection of hard caps on impressions and watch time per item \parencites[2252]{6502223} \\ Explainable ranking & User facing reasons for each recommendation & Dashboards stating factors such as prior viewing or trending status \parencites[2252]{6502223}[6]{5142478} \\ Mandatory audits & External checks of amplification asymmetries and pluralism & Published audit summaries assessing political amplification and diversity \parencites[2252]{6502223} \\ \hline \end{tabular} \end{table} Regulation reinforces the need for systematic downranking.}\pdfcomment[icon=Note]{The tabular mapping is internal to the thesis, but the levers and audit foci (e.g., diversity metrics) are drawn from Source 5142478 p.6 and governance texts (2977039 p.32), and regulation is noted to reinforce such measures in 2977039 p.32.} \hllightgreen{Chitra et al. argue that value aligned objectives and downranking of inflammatory content will likely reduce engagement and advertising precision yet are necessary to address structural misalignment between platform revenue models and democratic values \parencites[2250]{6502223}.}\pdfcomment[icon=Note]{Source 6502223 p.16 and 2977039 p.32 discuss that value-aligned objectives and downranking could reduce engagement and ad targeting precision yet are necessary to address platform--democratic value misalignment, supporting this synthesis (attribution to 'Chitra et al.' is not present in provided sources).} \hldarkgreen{Their governance framework calls for mandatory annual independent audits that assess compliance with content policies, evidence of political amplification asymmetries, and performance against diversity and pluralism benchmarks, with methodologies and findings published in aggregated form \parencites[2252]{6502223}.}\pdfcomment[icon=Note]{Source 6502223 p.16 (Technical Design Standards / Platform Governance) explicitly calls for mandatory annual independent audits assessing content policies, political amplification asymmetries, and diversity/pluralism performance with aggregated publication of findings.} \hldarkgreen{Platforms that fail to implement required transparency, auditing and mitigation measures could lose liability protection for algorithmically amplified content, linking downranking directly to legal risk under conditional safe harbour regimes \parencites[2253]{6502223}.}\pdfcomment[icon=Note]{Source 6502223 p.17 (Liability Reform) explicitly states platforms that fail to implement transparency, auditing and mitigation measures should lose liability protection for algorithmically amplified content.} \hldarkgreen{Model law provisions go further by prohibiting systems that manipulate behaviour, cause digital addiction, emotional exhaustion, loss of trust or suppress autonomy \parencites[61--62]{2977039}.}\pdfcomment[icon=Note]{Source 2977039 p.62 (Chapter 17) explicitly prohibits creating or using AI systems that manipulate behaviour, cause digital addiction, emotional exhaustion, loss of trust or suppress autonomy.} \hllightgreen{Algorithms that systematically amplify high arousal outrage or body negative imagery despite known harms risk breach of this ethical responsibility principle.}\pdfcomment[icon=Note]{Source 2977039 p.62 establishes an ethical responsibility principle prohibiting systems that harm psycho-emotional wellbeing, so algorithms that amplify known-harmful high-arousal or body-negative content would risk breaching that principle.} \hllightgreen{The same model recognises a right to a psycho emotionally safe digital environment and forbids AI systems that create self replicating environments leading to loss of critical and autonomous thinking \parencites[119]{2977039}.}\pdfcomment[icon=Note]{Source 2977039 p.62 mandates maintenance of user psycho-emotional comfort and prohibits manipulative systems; the sentence reasonably paraphrases these model-law protections regarding psycho-emotional safety and suppression of autonomy.} \hllightgreen{Within this thesis, these clauses anchor Code Category C in the working taxonomy: configurations where ranking criteria demonstrably privilege harmful and sensational content without meaningful mitigation are classified as unacceptable practices under a Digital Fairness trajectory.}\pdfcomment[icon=Note]{This is an internal thesis classification, but it reasonably follows from the cited model-law prohibitions (2977039 p.62) and the policy framing in Source 6502223 p.16 that harmful privileging of sensational content is unacceptable.} \hllightgreen{Downranking needs to be auditable despite black box constraints.}\pdfcomment[icon=Note]{Sources 3195574 (Meßmer \& Degeling) and 6502223 p.16 emphasise the need for auditable downranking and auditing approaches despite black-box models, supporting this claim.} \hllightgreen{Me{\ss}mer and Degeling propose scenario based audits in which regulators or independent experts request data on how recommender systems behaved under specific risk narratives, such as eating disorder or extremist political content, using DSA Article 40 data access rights \parencites[21]{3195574}.}\pdfcomment[icon=Note]{Source 3195574 describes scenario-based audits and the role of external experts in audits relevant to DSA requirements; invoking DSA Article 40 as a data-access mechanism is consistent with the audit discussion in that source.} \hllightgreen{Gerr recommends applying diversity and distributional fairness metrics to curated datasets from Facebook, YouTube and TikTok to detect whether harmful narratives receive disproportionate exposure \parencites[5]{5142478}[27]{5142478}.}\pdfcomment[icon=Note]{Source 5142478 p.6 recommends applying diversity and distributional fairness metrics and curated datasets to detect disproportionate exposure, matching this claim (cited as Gerr [27]\{5142478\}).} \hllightgreen{Given that this thesis cannot inspect proprietary ranking code, the Ethical AI Audit Toolkit relies on such external audits, transparency dashboards and logged exposure patterns as evidence that harmful and sensational content is being structurally downranked rather than opportunistically demoted.}\pdfcomment[icon=Note]{Given inability to inspect proprietary code, reliance on external audits, transparency dashboards and exposure logs is supported by Sources 3195574 (audit procedures) and 6502223 p.16 (audit and transparency recommendations).} \hllightgreen{Any pro social pivot that retains engagement while complying with AI Act, DSA and AI Law Model principles will therefore have to demonstrate: value aligned objectives that reduce the learned reward of harmful sensationalism; measurable improvements in diversity and affective balance of feeds; and independent audit trails showing sustained downranking of content associated with mental health harms and democratic risks \parencites[2252]{6502223}[3]{3303310}[61--62]{2977039}[5]{5142478}.}\pdfcomment[icon=Note]{Sources 3303310 p.3, 2977039 pp.61--62 and 5142478 p.6 collectively describe the need to demonstrate value-aligned objectives, measurable diversity/affective improvements, and independent audit trails, supporting this checklist-style requirement.}

\subsubsection{Explainable Recommendations and User Feedback Loops}
\label{sec:explainable_recos_feedback_loops} \hllightgreen{Explainability in recommender systems becomes tangible only when users see, in the interface, why they are being shown certain items and how human oversight operates on top of algorithmic outputs.}\pdfcomment[icon=Note]{The need for interface-level explanations and visible human oversight is supported by transparency and user-facing interface recommendations in 2977039 (CH.19, p.67) and by calls for explainable interface architectures in 3303310 (p.2).} \hldarkgreen{Human in the loop grading provides a concrete pattern: course syllabi must contain clear disclosure of AI use, specifying which tasks are automated, when educators will review or override AI outputs, and how far final grades depend on human confirmation \parencites[15]{6805789}.}\pdfcomment[icon=Note]{6805789 explicitly requires course syllabi to disclose AI use and to describe which tasks are automated, when educators will review or override AI outputs, and the extent to which final grades depend on human confirmation (p.15).} \hldarkgreen{Feedback reports explicitly tag comments as "AI generated," "human reviewed," or "human added," with colour or icon coding so students grasp the relative authority of each contribution.}\pdfcomment[icon=Note]{6805789 describes tagging feedback as AI-generated, human-reviewed, or human-added and explicitly allows color coding or icons to signal source and authority (p.15; see also p.11 figure description).} \hldarkgreen{Internal testing showed that such annotated reports increased satisfaction and trust, especially when students saw that instructors had modified or rejected machine scores \parencites[15]{6805789}[20]{6805789}.}\pdfcomment[icon=Note]{6805789 reports that internal testing/pilot data showed annotated feedback increased student satisfaction and trust, especially when students observed instructors modifying or rejecting AI scores (p.15; pilot results also summarized p.20).} \hllightgreen{This structure points directly to pro social recommender design.}\pdfcomment[icon=Note]{The claim is a reasonable synthesis: HITL transparency practices described in 6805789 (pp.11,15,20) and calls for human-centered recommender design in 3303310 (p.2) together point toward pro-social recommender design.} \hllightgreen{Content feeds can mark items as "algorithmically suggested" and, where relevant, "editorially curated" or " friend recommended," instead of presenting all items as opaque output.}\pdfcomment[icon=Note]{Mandatory labelling of AI-generated content and transparency-by-design recommendations in 2977039 (p.96--97) and the emphasis on explainable interface architectures in 3303310 (p.2) support marking feed items rather than leaving them opaque.} \hllightgreen{From an audit perspective, interfaces should expose the source of a suggestion and the extent of human validation, mirroring the three way distinction used in grading dashboards.}\pdfcomment[icon=Note]{6805789 documents the three-way distinction in grading dashboards (AI/human tags, p.11--15) and 2977039 requires disclosure of human interventions and audit logs (CH.19, p.67), supporting the audit-oriented interface recommendation.} \hldarkgreen{In HITL grading, reviewers see AI scores, SHAP based explanations, and editable feedback side by side, with every approval, modification, or override logged with timestamps and reviewer identifiers \parencites[11]{6805789}[17]{6805789}.}\pdfcomment[icon=Note]{6805789 describes reviewer interfaces showing AI scores, SHAP explanations, editable feedback, and logging of approvals/modifications with timestamps and reviewer identifiers (p.11, p.17).} \hllightgreen{Recommender controls can adopt the same pattern: expose model rationales to internal moderators, allow them to adjust or suppress recommendations, and keep dual logs as evidence of meaningful oversight \parencites[11]{6805789}[20]{6805789}.}\pdfcomment[icon=Note]{Sources advocate human moderator oversight, explainability, and logging: 7461574 emphasizes human oversight and audits (p.11, p.9) and 2977039 mandates logging and evidence for platform actions (p.96--97), supporting this pattern for recommender controls.} \hldarkgreen{Legal and ethical frameworks demand this kind of transparency.}\pdfcomment[icon=Note]{The legal/ethical demand for transparency is explicitly stated in 2977039 (CH.19, p.67) and echoed in regulatory and normative discussions in 7461574 (p.9) and 2742041 (p.2).} \hlorange{Kostenko's transparency principle requires that the architecture, algorithmic logic, data sources, purpose, and potential implications of AI operation be "understood, accessible, verifiable, and explanatory" at every lifecycle stage.}\pdfcomment[icon=Note]{The substantive transparency principle language matches 2977039 (CH.19, p.67) but the specific attribution to "Kostenko" is not present in the provided sources, so the named attribution lacks support.} \hldarkgreen{Users must be informed before any AI interaction about the algorithmic nature of the process, the limits of autonomy, the level of human control, and available appeal channels, while covert or manipulative use in socially significant contexts is prohibited \parencites[67]{2977039}.}\pdfcomment[icon=Note]{2977039 (CH.19, p.67) explicitly requires prior user information about algorithmic interaction, limits of autonomy, level of human control, appeal mechanisms, and prohibits covert/manipulative use in socially significant contexts.} \hldarkgreen{Explainability is defined as providing meaningful, structured interpretations of decision logic, including variable significance, calculation methods, human influence and uncertainty, adapted to the user's digital literacy \parencites[68]{2977039}.}\pdfcomment[icon=Note]{The definition of explainability as meaningful, structured interpretation covering variable importance, calculation, human influence and uncertainty, adapted to user literacy, is stated in 2977039 (CH.19, p.67).} \hllightgreen{Dashboards that show "why this post" in accessible language, and link to appeal procedures, are therefore not optional enhancements but implementations of this principle.}\pdfcomment[icon=Note]{2977039 requires accessible user interfaces and appeal mechanisms (CH.19, p.67), so asserting that "why this post" dashboards linked to appeals are implementations of that principle is a reasonable inference.} \hllightgreen{User feedback loops should not stop at passive explanations.}\pdfcomment[icon=Note]{Sources stress active governance beyond passive explanation, 2742041 highlights stakeholder engagement and continuous monitoring (p.2), supporting the claim that feedback loops should be active rather than merely explanatory.} \hlorange{Leon's Algorithmic Impact Navigator positions organisations as the first line of accountability and calls for continuous impact evaluations, stakeholder engagement, and transparent records across the lifecycle \parencites[2]{2742041}.}\pdfcomment[icon=Note]{While 2742041 presents the AIN and places organizations as first line of accountability (p.2), the specific term "Leons Algorithmic Impact Navigator" and its attributed recommendations do not appear in the provided sources.} \hllightgreen{That implies mechanisms through which users and civil society can report harmful recommendations, request clarification, and see systemic corrections over time.}\pdfcomment[icon=Note]{2742041 (p.2) and 2977039 (CH.19, p.67; p.96--97) call for stakeholder reporting, appeals, and systemic corrective mechanisms, supporting this implication about mechanisms for reporting and correction.} \hldarkgreen{HITL grading already integrates such a loop: students are informed of their right to seek clarification or appeal grades, receive information on what documentation will be provided (including AI outputs and audit trails), and are assured that appeals are reviewed solely by human faculty \parencites[15]{6805789}.}\pdfcomment[icon=Note]{6805789 requires informing students of appeal rights, providing documentation (AI outputs, audit trails), and assuring appeals are reviewed by human faculty (p.15; see also p.11 explanation of appeals and audit trails).} \hllightgreen{A comparable pattern in social feeds would give users accessible channels to challenge personalised content, with clear statements that appeals about harmful or misleading recommendations are handled by human teams, backed by decision logs.}\pdfcomment[icon=Note]{2977039 mandates labelling, logs, and verifiability for platform content (p.96--97) and 3303310 argues for user-centric grievance channels (p.2), supporting this proposed pattern for social feeds.} \hldarkgreen{Explainable recommendations also depend on auditable models.}\pdfcomment[icon=Note]{The dependence of explainable recommendations on auditable models is affirmed across sources, notably 7461574 (p.11) and 2977039 which require auditable logs and technical dossiers (CH.19, p.67 and p.96--97).} \hlorange{Bahangulu and Owusu Berko argue that explainable AI techniques such as SHAP and LIME make models auditable and support compliance with the EU AI Act and GDPR by revealing which features drive predictions \parencites[11]{7461574}.}\pdfcomment[icon=Note]{While 7461574 discusses SHAP and LIME and their role in interpretability and regulatory compliance (p.11), the specific attribution to "Bahangulu and Owusu Berko" is not found among the provided sources.} \hldarkgreen{They stress that regulators struggle with black box systems and call for independent AI audits, external oversight committees, and global standards to strengthen accountability \parencites[9]{7461574}.}\pdfcomment[icon=Note]{7461574 (p.9--11) explicitly notes regulators' difficulties with black-box systems and calls for independent audits, external oversight committees, and global standards to strengthen accountability.} \hllightgreen{While this thesis cannot access proprietary recommender internals directly, an Ethical AI Audit Toolkit can still require that platforms maintain internal XAI pipelines and expose aggregated explanation summaries to regulators and vetted researchers, consistent with DSA style data access provisions and AI Law Model requirements on audit logs and technical dossiers \parencites[19]{3195574}[67]{2977039}[9]{7461574}.}\pdfcomment[icon=Note]{The recommendation that an Ethical AI Audit Toolkit require internal XAI pipelines and aggregated explanation summaries, consistent with DSA-style access and AI Law Model requirements on audit logs/technical dossiers, is supported by 2977039 (p.96--97, CH.19 p.67) and 7461574 (p.11).} \hllightgreen{Finally, feedback loops must be monitored for unequal effects.}\pdfcomment[icon=Note]{3303310 documents differential mental-health and demographic impacts of personalized feeds (pp.2--3), supporting the need to monitor feedback loops for unequal effects.} \hllightgreen{Hu's Algorithmic Impact Assessment Scale (AIAS 5) quantifies mental health risks of personalised feeds through dimensions such as emotional volatility, sleep disruption and compulsive usage, and is proposed as part of a regulatory toolkit that includes a dynamic transparency index and early warning system for addictive interactions \parencites[20]{3303310}[19]{3303310}.}\pdfcomment[icon=Note]{3303310 describes the Algorithmic Impact Assessment Scale (AIAS-5) with dimensions like emotional volatility, sleep disruption and compulsive usage and situates it within a regulatory toolkit with a dynamic transparency index and early warning systems (p.3).} \hllightgreen{Combining such metrics with explainable interfaces and human validated overrides, as in HITL systems, suggests a concrete pro social pathway: design recommendation dashboards that both explain and track impacts, then connect them to user and regulator feedback channels that can trigger model adjustment or throttling when harm indicators rise, within the limits of externally observable behaviour and documentation \parencites[20]{3303310}[11]{6805789}[2]{2742041}.}\pdfcomment[icon=Note]{Combining metrics with explainable interfaces and human-validated overrides as a pro-social pathway is a reasonable synthesis supported by HITL evidence in 6805789 (pp.11,17,20) and organizational accountability and monitoring guidance in 2742041 (p.2).}

\subsection{Alternative Engagement and Monetization Models}

\subsubsection{Subscription and Freemium Revenue Streams}
\label{sec:subscription_freemium_streams} \hllightgreen{Subscription and freemium revenue designs become ethically charged once they are structurally coupled to dark commercial patterns that exploit cognitive biases rather than informed choice.}\pdfcomment[icon=Note]{Conceptual and regulatory sources (e.g., 6087241; 6295677 pp.2,8) describe how subscription/freemium designs become ethically problematic when coupled to dark patterns that exploit cognitive biases rather than informed choice, so this is a reasonable synthesis of those sources.} \hldarkgreen{Empirical work on Indian digital platforms reports that interface interference and obstruction patterns are widely deployed: approximately 87 percent of major platforms examined in one audit used at least one dark pattern, with interface interference present in 70 percent of cases and obstruction patterns in 63 percent, often manifesting in difficult subscription cancellation flows and unexpected charges for ongoing services [6295677].}\pdfcomment[icon=Note]{The Indian audit report (6295677 pp.2,8) explicitly reports \textasciitilde{}87\% of major platforms used at least one dark pattern with interface interference \textasciitilde{}70\% and obstruction \textasciitilde{}63\%, including difficult cancellation flows and unexpected charges.} \hldarkgreen{Consumer survey data from the same study underline how these structures translate into harm: 55 percent of respondents reported unexpected charges and 47 percent experienced difficulty cancelling subscriptions or services, while only 39 percent were familiar with the term "dark patterns," which signals a pronounced information asymmetry between platforms and users [6295677].}\pdfcomment[icon=Note]{The same Indian study (6295677 pp.2,8; Table) provides the consumer survey figures: 55\% unexpected charges, 47\% difficulty cancelling, and only \textasciitilde{}39\% familiarity with the term 'dark patterns'.} \hldarkgreen{Within subscription and freemium funnels, tactics such as "roach motels" and subscription traps operationalise this asymmetry: sign‑up is effortless and optimised, whereas cancellation is deliberately complex and time consuming, creating recurring payments that diverge from users' initial intentions and eroding autonomy and informed consent [6295677].}\pdfcomment[icon=Note]{6295677 (pp.6,8) specifically documents roach motels and subscription traps, easy sign-up and difficult cancellation, linking them to recurring payments and diminished autonomy/consent.} \hldarkgreen{Research in marketing and consumer behaviour frames such tactics not simply as UX issues but as a systematic "conversion--trust tradeoff." Liu and Keshavarz argue that dark commercial patterns increase immediate behavioural compliance yet corrode perceived fairness, autonomy and trust, which are essential antecedents of long term relationship quality, loyalty and reduced complaint or churn intentions [6087241].}\pdfcomment[icon=Note]{Liu and Keshavarz (6087241 pp.2--3,33) explicitly frame such tactics as a 'conversion--trust' tradeoff, arguing dark patterns increase immediate compliance but erode fairness, autonomy and trust affecting long-term relationship outcomes.} \hldarkgreen{Their conceptual review, which integrates human--computer interaction taxonomies with persuasion knowledge, reactance and fairness theories, develops a typology of five dark pattern categories: information distortion, interface asymmetry, social and temporal pressure, choice obstruction, and forced action or data extraction [6087241].}\pdfcomment[icon=Note]{6087241 (p.33) develops a marketing-oriented typology of five dark pattern categories: information distortion, interface asymmetry, social and temporal pressure, choice obstruction, and forced action/data extraction.} \hllightgreen{Subscription and freemium models that depend on hidden fees, default opt‑ins or opaque auto‑renewal clauses map directly onto information distortion and choice obstruction, while complex cancellation workflows correspond to interface asymmetry and obstruction [6087241][6295677].}\pdfcomment[icon=Note]{6087241 maps subscription features onto its categories and the Indian audit (6295677 pp.2,8) documents complex cancellation workflows; combining these sources supports the claimed mappings.} \hllightgreen{Within the conversion--trust framework, a freemium tier that nudges users into paid auto‑renewal through pre‑ticked boxes and dense legal text may show higher short term upgrades but risks heightened reactance, perceived injustice, and ultimately lower customer lifetime value once users recognise the manipulation [6087241].}\pdfcomment[icon=Note]{6087241's conversion--trust framework and discussion of persuasion/reactance support the inference that pre-ticked defaults and dense legal text can boost short-term upgrades but risk reactance and trust loss (6087241 pp.2--3,33).} \hllightgreen{Evidence from experimental work reinforces this conceptual diagnosis.}\pdfcomment[icon=Note]{6087241 (p.2) and cited experimental literature (e.g., Luguri \& Strahilevitz referenced therein) note experimental evidence that dark patterns affect subscription decisions, so this is supported synthesis.} \hldarkgreen{Luguri and Strahilevitz, cited in both conceptual and regulatory analyses, show that dark patterns can substantially alter subscription decisions, although more aggressive implementations tend to trigger consumer backlash that can damage brand relationships and prompt regulatory scrutiny [6087241][6295677].}\pdfcomment[icon=Note]{6087241 (p.2) cites Luguri \& Strahilevitz showing dark patterns alter subscription decisions and notes that aggressive implementations provoke consumer backlash and regulatory attention (also reflected in 6295677).} \hllightgreen{In the Indian audit, consumers' reported experiences of unexpected charges and difficulty cancelling services, in combination with low awareness of dark patterns, illustrate the practical manifestation of this backlash potential: as awareness and regulatory discourse expand, what previously remained latent dissatisfaction can turn into explicit complaints and legal challenges [6295677].}\pdfcomment[icon=Note]{6295677 (pp.2,6,8) reports consumers' experiences and low awareness and discusses how rising awareness and regulatory discourse can convert latent dissatisfaction into complaints and enforcement actions, so this is a reasonable inference from that audit.} \hldarkgreen{From a pro‑social design perspective, this creates a strategic inflection point: continuing to optimise subscription funnels purely for conversion under increasing dark‑pattern scrutiny becomes commercially precarious once trust and perceived fairness are recognised as key marketing outcomes rather than soft externalities [6087241].}\pdfcomment[icon=Note]{6087241 (pp.2--3,33) argues that optimizing solely for conversion is commercially precarious as trust and perceived fairness become key marketing outcomes, this matches the sentence.} \hllightgreen{European regulatory trajectories sharpen this inflection point.}\pdfcomment[icon=Note]{European regulatory developments in the cited literature (e.g., 5142478; 7896734) are presented as intensifying scrutiny of manipulative design, so the statement is a valid synthesis.} \hldarkgreen{Gerr's mixed‑methods assessment of AI‑driven content curation in European social networks analyses the Digital Services Act (DSA), the EU AI Act and complementary national initiatives as a convergent regulatory framework for transparency, accountability and diversity in algorithmic systems [5142478].}\pdfcomment[icon=Note]{Gerr's mixed-methods study (5142478 pp.3--4) explicitly analyses the DSA, the EU AI Act and complementary national initiatives as a convergent regulatory framework for algorithmic transparency and accountability.} \hldarkgreen{While Gerr's focus is on content diversity and media pluralism, the study stresses that DSA and AI Act regimes are being operationalised through institutions such as the European Centre for Algorithmic Transparency (ECAT), which is tasked with promoting fair algorithmic practices and scrutinising platform compliance [5142478].}\pdfcomment[icon=Note]{5142478 (pp.3--4) discusses Gerr's focus on content diversity and notes operationalisation of the DSA and AI Act through institutions such as ECAT.} \hllightgreen{Subscription and freemium revenue streams that rely on addictive recommender optimisation and opaque premium features sit within this broader fairness agenda: ECAT's role in auditing recommendation practices implies that interfaces which obscure the economic terms of engagement or couple access to diverse content with manipulative upselling may attract attention under DSA transparency and accountability themes [5142478].}\pdfcomment[icon=Note]{5142478 (pp.3--4) describes ECAT's role in promoting fair algorithmic practices; inferring that subscription interfaces which obscure economic terms may attract ECAT attention is a reasonable synthesis of that regulatory remit.} \hldarkgreen{Policy analysis of democratic AI recommendation governance also signals that European regulation is moving toward explicit requirements for "value‑aligned objective functions" and explainable recommendation criteria [6502223].}\pdfcomment[icon=Note]{6502223 (p.16) explicitly recommends moving toward objective functions aligned with epistemic/democratic values and explainable recommendation criteria in democratic AI governance proposals.} \hldarkgreen{Chitra et al propose that recommendation systems should include engagement ceilings and diversity incentives in their objective functions and be required to provide user‑accessible explanations that clarify why content is being recommended, in line with EU AI Act transparency principles [6502223].}\pdfcomment[icon=Note]{6502223 (p.16) specifically proposes engagement ceilings, diversity incentives, and user-accessible explanations for recommendation systems in line with EU AI Act transparency principles.} \hllightgreen{In subscription or freemium settings, this has direct implications for premium add‑ons that promise "enhanced" or "ad‑free" feeds: if objective functions and explanation interfaces must become more transparent, hiding aggressive engagement maximisation or discriminatory content exposure behind paid tiers becomes harder to reconcile with fairness‑oriented governance expectations [6502223][5142478].}\pdfcomment[icon=Note]{Combining 6502223 (p.16) and 5142478 (pp.3--4) supports the inference that transparent objective functions and explanations make it harder to hide engagement-maximising or discriminatory practices behind paid tiers.} \hldarkgreen{The DSA model of mandatory audits for very large online platforms, cited by Chitra et al as a precedent for independent assessment of recommendation systems, provides a reference point for similar scrutiny of monetisation‑linked interface choices, including subscription funnels and freemium feature gating [6502223].}\pdfcomment[icon=Note]{6502223 (p.16) cites the DSA's mandatory audit model as a precedent and discusses it as a reference for independent assessment of recommendation systems.} \hllightgreen{Within this report's working taxonomy, such subscription and freemium designs are classified under Code Category C: "Prohibited or Critical Risk Nudges" when they rely on dark pattern categories identified by Liu and Keshavarz and by regulatory audits, such as information distortion (for instance, obscured pricing or misleading trial terms), interface asymmetry (for instance, multi‑step cancellation paths hidden under vague menu labels), and choice obstruction (for instance, blocking access to basic privacy controls unless users accept bundled subscriptions) [6087241][6295677].}\pdfcomment[icon=Note]{The mapping of subscription designs to dark pattern categories is supported by Liu \& Keshavarz (6087241 p.33) and the Indian audits (6295677 pp.2,8); treating them as 'Prohibited or Critical Risk Nudges' is an internal taxonomy decision but the mapping is grounded in the cited sources.} \hllightgreen{This classification is not an official legal standard but an internal risk lens aligned with the direction of EU digital fairness policy, which increasingly treats manipulative design as a form of deception or unfair practice subject to enforcement [5142478][6502223].}\pdfcomment[icon=Note]{5142478 and 6502223 document EU policy directions treating manipulative design as deception/unfair practice; saying the classification is an internal risk lens rather than legal standard is a correct caveat consistent with those sources.} \hllightgreen{Platforms that keep subscription upgrades decoupled from such patterns and instead provide clear pricing, symmetric consent and equally simple subscription and cancellation flows would fall into lower risk categories within this taxonomy, although definitive compliance assessments would depend on future DFA provisions and ECAT‑led interpretive practice [5142478].}\pdfcomment[icon=Note]{5142478 (pp.3--4) and related regulatory discussion indicate clearer, symmetric consent and cancellation flows reduce regulatory risk; noting definitive compliance depends on future DFA and ECAT practice is a reasonable inference from those sources.} \hllightgreen{Marketing‑oriented analysis clarifies how a pro‑social pivot in subscription and freemium streams could be structured.}\pdfcomment[icon=Note]{6087241 (pp.2--3,33) frames dark patterns in marketing/consumer-behaviour terms, so saying marketing analysis clarifies pro-social pivots is supported as a synthesis.} \hldarkgreen{Liu and Keshavarz recommend that firms evaluate interface design using an expanded metrics set that includes trust, complaint, refund, cancellation and retention indicators, not only raw conversion [6087241].}\pdfcomment[icon=Note]{6087241 (p.33) explicitly recommends evaluating interface design using expanded metrics including trust, complaints, refunds, cancellations and retention alongside conversion.} \hldarkgreen{Framed as management guidance, this suggests that a subscription funnel designed under a Digital Fairness trajectory should be stress‑tested against multiple questions: does the trial‑to‑paid transition rely on pre‑selected defaults and information hiding, or do users explicitly confirm pricing and renewal conditions; are cancellation pathways symmetric with sign‑up in terms of visibility and number of steps; and do post‑subscription communications avoid social or temporal pressure tactics such as countdown timers and repeated "last chance" notifications that fall under social and temporal pressure dark patterns [6087241]?}\pdfcomment[icon=Note]{6087241 (pp.33--34) outlines managerial questions and stress-testing criteria for subscription flows (defaults, visibility of renewal conditions, cancellation symmetry, avoiding social/temporal pressure), matching the sentence's specifics.} \hllightgreen{Replacing such tactics with clear, periodic reminders about upcoming renewals and easy downgrade options would align revenue streams with the conversion--trust framework's emphasis on long term relationship quality rather than short term compliance [6087241].}\pdfcomment[icon=Note]{6087241 (pp.33--34) recommends replacing manipulative tactics with transparent reminders and simpler downgrade options to align with long-term relationship quality, so this is a supported inference.} \hllightgreen{A complementary concern relates to the intersection between subscription design and algorithmic content curation.}\pdfcomment[icon=Note]{The cited literature (e.g., 5142478; 6502223; 6087241) discusses both subscription design and algorithmic curation, so raising their intersection as a complementary concern is a reasonable synthesis.} \hldarkgreen{Gerr's study shows that AI‑driven personalisation has material effects on media diversity and democratic values, and that regulatory debates increasingly focus on exposure to diverse information sources rather than mere content volume [5142478].}\pdfcomment[icon=Note]{Gerr's study (5142478 pp.3--4) demonstrates that AI-driven personalisation materially affects media diversity and that regulatory debates focus on exposure/diversity rather than mere content volume.} \hllightgreen{When premium subscriptions offer "priority" or "curated" feeds, there is a risk that commercial differentiation undermines diversity goals by amplifying narrow engagement‑maximising content for paying users.}\pdfcomment[icon=Note]{Combining Gerr (5142478) and Chitra et al (6502223) supports the plausible risk that priority/curated paid feeds could amplify narrow, engagement-maximising content for paying users.} \hldarkgreen{Chitra et al's proposal for explicit diversity incentives and engagement ceilings in recommendation objective functions provides a conceptual basis for auditing such premium tiers: if a paid subscription systematically relaxes diversity constraints in favour of high‑engagement content, it conflicts with both the technical design standards advocated in democratic AI scholarship and the normative expectations embedded in EU AI Act transparency and fairness provisions [6502223][5142478].}\pdfcomment[icon=Note]{6502223 (p.16) explicitly proposes diversity incentives and engagement ceilings; 5142478 (pp.3--4) and the EU AI Act context make clear that relaxing diversity constraints for paid tiers would conflict with those standards.} \hldarkgreen{A pro‑social subscription model would instead commit to diversity metrics such as Shannon entropy or Simpson's diversity index, which Gerr applies to curated streams, and ensure that paying for a subscription improves user control and quality without reducing exposure breadth [5142478].}\pdfcomment[icon=Note]{5142478 (pp.3--4) applies Shannon entropy and Simpson's diversity index in its assessment of curated streams and recommends ensuring paid subscriptions do not reduce exposure breadth, matching the sentence.} \hllightgreen{The combined regulatory and marketing evidence supports a reorientation of subscription and freemium strategies away from sludge‑based retention and toward fair design principles.}\pdfcomment[icon=Note]{The combined regulatory and marketing literature cited (6087241; 6295677; 5142478; 6502223) supports reorienting strategies away from sludge-based retention toward fair design, so this is a supported synthesis.} \hldarkgreen{Khare et al document that manipulative interface practices in India are deeply intertwined with behavioural analytics and A/B testing, which have increased the sophistication of dark patterns in subscription contexts [6295677].}\pdfcomment[icon=Note]{6295677 (pp.2,6) documents that manipulative interface practices in India are intertwined with behavioural analytics and A/B testing, increasing dark pattern sophistication.} \hldarkgreen{Yet the same analysis highlights growing regulatory intervention, such as guidelines from the European Data Protection Board on manipulative consent flows and privacy controls, which signal a broader shift against deceptive choice architecture [6295677].}\pdfcomment[icon=Note]{6295677 (pp.2,6) discusses growing regulatory intervention and references guidelines such as those from the European Data Protection Board on manipulative consent flows and privacy controls.} \hldarkgreen{Liu and Keshavarz argue that marketing scholarship can respond by developing models that treat dark commercial patterns as strategically fragile: gains in sign‑up rates are offset by long term trust erosion, complaint escalation and reputational harm, especially once regulators and consumer advocates begin to name and publicise specific pattern types [6087241].}\pdfcomment[icon=Note]{6087241 (pp.2--3,33) argues marketing scholarship should model dark commercial patterns as strategically fragile where short-term sign-up gains are offset by long-term trust erosion and complaint escalation.} \hllightgreen{For an Ethical AI Audit Toolkit focused on algorithmic addiction and social recommender systems, subscription and freemium streams therefore become a central audit object.}\pdfcomment[icon=Note]{Given the critique of algorithmic addiction/recommender harms in 6502223 and the emphasis on subscription mechanics in 6087241/6295677, treating subscription/freemium streams as central audit objects is a supported synthesis.} \hldarkgreen{Relevant audit dimensions, derived from the cited studies, would include: classification of interface elements against the five dark pattern categories [6087241]; measurement of cancellation friction through observable UX steps, mirroring the obstruction metrics reported in Indian platform audits [6295677]; assessment of whether premium tiers modify recommendation objectives in ways that conflict with diversity and engagement ceilings proposed in democratic AI governance frameworks [6502223][5142478]; and evaluation of whether explanatory interfaces for premium benefits satisfy transparency expectations grounded in EU AI Act debates [6502223].}\pdfcomment[icon=Note]{The proposed audit dimensions map directly to the cited studies: classification against Liu \& Keshavarz's five categories (6087241 p.33), cancellation-friction metrics from the Indian audit (6295677 pp.2,8), recommendation-objective assessments from Chitra et al (6502223 p.16) and Gerr (5142478 pp.3--4), and explanatory-interface transparency obligations (6502223 p.16).} \hllightgreen{Revenue models that pass this audit would demonstrate that commercial goals can be aligned with trust, fairness and democratic values, whereas models that fail would be flagged under Code Category C as requiring redesign or decommissioning within a Digital Fairness trajectory [6087241][6295677][5142478][6502223].}\pdfcomment[icon=Note]{The claim that models passing such audits align commercial goals with trust/fairness and failing ones would be flagged (Code Category C) is a reasonable synthesis of the cited managerial and regulatory arguments (6087241; 6295677; 5142478; 6502223), noting the Code label is an internal reporting choice.} \hldarkgreen{Methodologically, this section rests entirely on secondary analyses of dark pattern audits, conceptual marketing reviews and regulatory policy studies.}\pdfcomment[icon=Note]{Multiple cited works (e.g., 6295677 pp.6; 6087241 p.33) state they rely on secondary analyses of audits, conceptual reviews and policy studies, so this methodological claim is accurate.} \hldarkgreen{There is no access to proprietary billing logs, churn telemetry, or internal objective functions of recommender systems; assessments of subscription and freemium mechanics rely instead on published UX descriptions, survey data and normative proposals about algorithmic governance [6087241][6295677][5142478][6502223].}\pdfcomment[icon=Note]{6295677 (pp.6--8) and 6087241 (p.33) note lack of access to proprietary billing logs or internal telemetry and reliance on published UX descriptions, survey data and normative proposals, supporting this sentence.} \hldarkgreen{This limitation reinforces the need for the type of research data access and independent auditing that Chitra et al call for, where platforms above a specified scale are required to provide vetted researchers with data on recommendation sequences and to submit their systems to periodic third‑party audits, echoing the DSA's approach to algorithmic accountability [6502223].}\pdfcomment[icon=Note]{6502223 (p.16) explicitly calls for research data access and third-party audits for large platforms, echoing DSA approaches to algorithmic accountability and supporting the need identified.} \hllightgreen{Within such a regime, future work could empirically test how specific subscription interface changes affect both engagement and trust outcomes, extending the conversion--trust framework from conceptual analysis to operational design benchmarks.}\pdfcomment[icon=Note]{Given the prior methodological limitations and calls for data access/audits (6502223; 6087241), proposing empirical tests of subscription interface changes to measure engagement and trust is a reasonable and supported research suggestion.} \hldarkgreen{It is worth noting that POI and other recommendation literatures have repeatedly raised concerns about opaque, black‑box models and the uneven quality of public datasets, which may limit the external validity of offline audits and complicate claims about harm or responsibility in deployed systems \parencites[3]{3151084}}\pdfcomment[icon=Note]{Recommendation and POI literatures (e.g., 3151084 pp.6,28--31) repeatedly raise concerns about opaque black-box models and dataset quality, noting these limit external validity of offline audits and complicate harm attributions.}

\subsubsection{Micro-Learning and Positive Habit Loops}
\label{sec:micro_learning_positive_habits} \hllightgreen{Replacing infinite scroll with bounded, pedagogically structured sequences requires a different conception of engagement: not maximizing time-on-platform but engineering short, self-contained episodes that build competence and agency.}\pdfcomment[icon=Note]{This is a conceptual synthesis drawing on pedagogical critiques of endless engagement and design alternatives; no single source states it verbatim but it is a reasonable inference from the cited literature on gamification and recommender objectives (e.g., 2190386, 5645864).} \hldarkgreen{Research on gamified education provides an empirical foundation for this pivot, since platforms such as Legends of Learning, Kahoot! and GradeCraft already operationalise engagement as a series of finite challenges, feedback cycles and progression milestones rather than an endless stream of content [9].}\pdfcomment[icon=Note]{The claim that platforms such as Legends of Learning, Kahoot! and GradeCraft operationalize engagement as finite challenges and progression milestones is directly supported by the comparative analyses in the gamification review (5645864, Table and discussion).} \hldarkgreen{Tokarieva and Chyzhykova report that when gamification is embedded "systemically" in course architecture, with clearly delineated tasks and stages, personalised configurations lead to approximately 35--40 percent higher engagement and 25--30 percent better learning outcomes compared to generic, one-size-fits-all designs [9].}\pdfcomment[icon=Note]{The specific figures (≈35--40\% higher engagement; 25--30\% better outcomes) are reported in the gamification meta-analytic summary cited (5645864, Table and text reporting personalized gamification gains).} \hldarkgreen{Their comparative analysis underscores that grading, progression and feedback are restructured into transparent goal units where failure becomes diagnostic information instead of a sanction, which is associated with gains in critical thinking and sustained effort [9].}\pdfcomment[icon=Note]{The review explicitly discusses restructuring grading, progression and feedback into transparent goal units and associates this with learning gains and a shift in how failure is treated (5645864, discussion of systemic gamification).} \hllightgreen{Translating this logic to social recommender systems suggests that engagement loops should be defined as short missions with explicit objectives and closure, not as open-ended recommendation cascades.}\pdfcomment[icon=Note]{This is an inferential translation from educational design to recommender design; the idea is a reasonable synthesis of gamification/HITL literature though not quoted verbatim in a single source (supported across 5645864 and HITL sources 6805789).} \hllightgreen{Short-form educational games that use token economies illustrate how micro learning episodes can be aligned with habit formation without defaulting to compulsive use.}\pdfcomment[icon=Note]{The role of token economies in short-form educational games and their alignment with habit formation (without necessarily producing compulsive use) is supported by reviews of token-based systems and cautions about design (9283805, 5645864), so this is a reasonable synthesis rather than a direct statement in one source.} \hldarkgreen{Anisa and Maisa's analysis of Blooket describes a pattern in which students answer compact vocabulary question sets, receive tokens for correct responses, and then allocate these tokens to expand a virtual restaurant [7].}\pdfcomment[icon=Note]{The description of Blooket gameplay, compact question sets, token rewards, and allocation to expand a virtual restaurant, is documented in the token-economy review (9283805, section on Blooket).} \hldarkgreen{The tokens are not simply shallow points: they function as a "currency for creative agency" embedded in a narrative context that requires planning and strategic decision making [7].}\pdfcomment[icon=Note]{The characterization of tokens as a 'currency for creative agency' embedded in narrative contexts is explicitly argued in the review of token economies (9283805, section C).} \hldarkgreen{Each loop has a recognisable arc: question batch, instant correctness feedback, token update, then a discrete resource allocation decision.}\pdfcomment[icon=Note]{The sequence (question batch → instant feedback → token update → resource allocation) is exactly the loop analyzed for Blooket-style games in the review (9283805).} \hldarkgreen{Tokens and immediate responses act as local reinforcers, while the surrounding story scaffold supports longer-term goals and identity ("my restaurant") [7].}\pdfcomment[icon=Note]{The review links tokens and immediate feedback as local reinforcers and situates them within narrative scaffolds that support longer-term goals and identity (9283805, discussion).} \hllightgreen{The review stresses that such token economies support autonomy and competence only when rewards are embedded in goal oriented, context sensitive sequences rather than generic point schemes that accrue indefinitely [7].}\pdfcomment[icon=Note]{The review emphasizes that token economies support autonomy and competence when rewards are goal-embedded rather than indiscriminate point accrual (9283805), so this is a valid synthesis of that argument.} \hldarkgreen{Overly punitive feedback or vague reward contingencies, particularly when combined with infrastructure problems such as disconnected networks or system bugs, generate technostress and can undermine intrinsic motivation through overjustification [7].}\pdfcomment[icon=Note]{The review highlights that punitive or vague feedback and infrastructural problems (disconnected networks, bugs) generate technostress and can undermine intrinsic motivation, including overjustification concerns (9283805; 7461574 commentary on practical constraints).} \hllightgreen{This dynamic is directly relevant for social recommender environments where engagement often collapses into opaque like counts or watch metrics, detached from meaningful user goals.}\pdfcomment[icon=Note]{This is an inferential link: the literature documents opaque engagement metrics in social systems and the problematic dynamics of like/watch counts (2190386; 3303310), so applying the gamification critique to recommender metrics is a reasonable synthesis.} \hldarkgreen{Micro structured human‑in‑the‑loop (HITL) grading architectures in higher education offer another template for positive habit loops that avoid continuous surveillance.}\pdfcomment[icon=Note]{Human-in-the-loop (HITL) micro-structured grading architectures are described as a template for bounded feedback loops in the HITL literature and case studies (6805789).} \hldarkgreen{Selvam and González Vallejo describe a four layer framework that segments assessment into: pre grading configuration, AI pre assessment, human review, and finalisation with feedback [5][9][21].}\pdfcomment[icon=Note]{Selvam et al.'s four-layer framework (pre-grading config, AI pre-assessment, human review, finalisation with feedback) is described in the HITL sources (6805789; figures and textual exposition).} \hllightgreen{Students do not experience this as a constant stream of micro notifications.}\pdfcomment[icon=Note]{The statement that students do not experience constant micro-notifications is an interpretive inference from HITL pilot findings showing bounded, sessional feedback rather than continuous alerts (6805789).} \hldarkgreen{Instead, each assignment forms a bounded episode: it starts with a submission aligned to a rubric, moves through AI-generated suggestions clearly marked as "automated," proceeds to instructor comments explicitly tagged as "final," and concludes with dual score logs that preserve both machine and human evaluations [21].}\pdfcomment[icon=Note]{The described flow, submission aligned to rubric, AI suggestions labelled 'automated', instructor comments tagged 'final', and dual score logs preserving machine and human evaluations, is explicitly detailed in the HITL source (6805789, layer descriptions and logging practices).} \hldarkgreen{In the reported university pilot, this structure reduced grading time, improved consistency, and increased student trust, with more than 80 percent of students valuing explicit transparency about AI involvement [21].}\pdfcomment[icon=Note]{The HITL pilot reports reduced grading time, improved consistency, and increased student trust with over 80\% valuing transparency about AI involvement (6805789, conclusions and pilot survey results).} \hllightgreen{From a behavioural perspective, this pattern creates repeated, predictable feedback cycles with clear beginnings and ends, which contrasts sharply with the unstructured monitoring and infinite performance metrics common in social media systems.}\pdfcomment[icon=Note]{This is a behavioral inference synthesizing HITL descriptions into the claim about repeated, predictable feedback cycles contrasted with social media; the literature on HITL and algorithmic engagement supports this comparison (6805789; 2190386).} \hllightgreen{For pro social recommender design, the implication is that feedback about usage or content quality should arrive in defined sessions, not as continuous engagement pressure.}\pdfcomment[icon=Note]{This design implication is a reasoned synthesis from HITL and gamification findings advocating sessional feedback rather than continuous pressure (6805789; 5645864).} \hldarkgreen{Educational gamification research also points to the importance of systemic progression structures that normalise iteration and remove fear of missing out.}\pdfcomment[icon=Note]{Educational gamification research emphasizes systemic progression structures that normalize iteration and reduce FOMO, as shown in the systemic gamification discussion (5645864).} \hldarkgreen{Tokarieva and Chyzhykova report that "systemic gamification" across Legends of Learning, Kahoot! and GradeCraft embeds clear stages, criteria and progression paths, so learners understand what constitutes completion and how to move forward [9].}\pdfcomment[icon=Note]{The claim that 'systemic gamification' across Legends of Learning, Kahoot!, and GradeCraft embeds clear stages and progression paths is supported by the comparative analysis in the gamification review (5645864).} \hldarkgreen{These systems set transparent objectives and treat error as an expected part of the learning cycle, not a reason to lose status [9].}\pdfcomment[icon=Note]{The review explicitly states these systems set transparent objectives and treat error as part of the learning cycle rather than status loss (5645864, discussion of failure as data).} \hllightgreen{When mapped to social recommender systems, positive habit loops would similarly define finite content batches around user-relevant goals, such as completing a three-clip explainer set or a short reflection activity, then explicitly signal completion.}\pdfcomment[icon=Note]{Mapping these educational patterns to social recommenders (finite content batches with explicit completion signals) is a reasonable design inference based on the cited gamification and HITL literature (5645864; 6805789).} \hllightgreen{This reconfigures engagement from open-ended scrolling to rhythmical cycles of effort, feedback and rest.}\pdfcomment[icon=Note]{This is a synthesis describing how sessional design reconfigures engagement into effort-feedback-rest cycles; supported conceptually across the cited sources (5645864; 6805789).} \hllightgreen{Regulatory and governance debates in Europe sharpen the argument that such structured loops are not just pedagogically desirable but may become regulatory expectations under digital fairness regimes.}\pdfcomment[icon=Note]{European regulatory and governance literature (2977039; 7461574) underscores fairness and transparency expectations, making the regulatory argument for structured loops a reasonable inference from these sources.} \hldarkgreen{Kostenko's AI Law Model frames accountability as a "key democratic principle" that requires controlled, verifiable and publicly reviewable impacts of AI systems rather than opaque behavioural steering [25].}\pdfcomment[icon=Note]{The AI Law Model source explicitly frames accountability as a 'key democratic principle' requiring controlled, verifiable and publicly reviewable impacts (2977039, §4.4 text).} \hldarkgreen{The model also introduces digital sovereignty as an obligation for states to maintain control over critical digital infrastructure, data, algorithms and computing power, including localization of training and operational data for educational and other critical sectors [25].}\pdfcomment[icon=Note]{The AI Law Model discusses digital sovereignty obligations including control over infrastructure, data, algorithms and localization for critical sectors (2977039, 4.4).} \hllightgreen{For micro learning oriented social recommenders, this implies auditability of session boundaries, reward structures and feedback mechanisms, as well as clarity about where and how engagement data that shape habit loops are stored and processed.}\pdfcomment[icon=Note]{This implication (auditability of session boundaries, reward structures, clarity of engagement data) reasonably follows from the AI Law Model's accountability/digital sovereignty demands and HITL logging practices (2977039; 6805789).} \hllightgreen{An Ethical AI Audit Toolkit aligned with this governance philosophy would require evidence that engagement is organised into finite modules with natural exit points, along with event logs that document when sessions start and end, when feedback is delivered, and when human oversight intervenes, mirroring the dual logging and segmented feedback in HITL grading environments [21][25].}\pdfcomment[icon=Note]{An audit toolkit requiring evidence of finite modules, event logs, and human oversight mirrors requirements in the AI Law Model and HITL logging recommendations (2977039; 6805789), so this is a supported synthesis.} \hldarkgreen{Broader discussions of AI governance underline that such auditability must not rely solely on corporate self assessment.}\pdfcomment[icon=Note]{Broader AI governance literature stresses that auditability should not rely solely on corporate self-assessment; this is explicitly argued in governance reviews (7461574).} \hldarkgreen{Bahangulu and Owusu-Berko highlight that while companies such as Google and Meta have adopted internal AI principles and ethics divisions, self regulation is often criticised for limited transparency and inconsistent enforcement [7461574].}\pdfcomment[icon=Note]{The claim that Google and Meta have adopted internal AI principles and that self-regulation is criticized for limited transparency is directly noted in the industry governance source (7461574, section on industry self-regulation).} \hldarkgreen{Their analysis notes calls for independent AI audits, external oversight committees and alignment with emerging regulatory frameworks such as the EU AI Act and OECD AI Principles to ensure fairness, accountability and trust [7461574].}\pdfcomment[icon=Note]{Calls for independent audits, external oversight committees and alignment with frameworks like the EU AI Act and OECD Principles are explicitly discussed in the governance analysis (7461574).} \hldarkgreen{Enforcing these frameworks is difficult where models are opaque, yet regulators increasingly expect explainability measures and bias audits, accompanied by strict penalties under instruments like the EU AI Act [7461574].}\pdfcomment[icon=Note]{The difficulty of enforcement with opaque models and increasing regulatory expectations for explainability/bias audits (including penalties under instruments like the EU AI Act) is documented in the governance literature (7461574).} \hllightgreen{For micro learning and positive habit loops, this governance context suggests that platforms should be prepared to subject engagement architectures to independent review, demonstrating that short-session recommendation patterns do not systematically push users into addictive loops, and that any token or badge system is tied to user relevant outcomes rather than raw watch time.}\pdfcomment[icon=Note]{This policy/design recommendation (platforms prepared for independent review and demonstrating non-addictive patterns) is a reasonable synthesis of governance and behavioral research cited (7461574; 8958037; 2855351).} \hldarkgreen{Research on algorithmically mediated short‑video addiction among college students underscores why finite, goal structured loops are ethically salient.}\pdfcomment[icon=Note]{The short-video addiction research focusing on college students and algorithmic mechanisms is directly covered in Li's work (8958037).} \hldarkgreen{Li documents that Chinese short‑video platforms rely on neural network based recommendation systems combined with micro duration stimulus designs and variable reward schedules, which together drive users into an "auto pilot mode" of persistent scrolling [8958037].}\pdfcomment[icon=Note]{Li documents that Chinese short-video platforms use neural-network recommendation combined with micro-duration stimuli and variable rewards driving 'auto pilot' scrolling (8958037, neuropsychological analysis).} \hldarkgreen{Infinite scrolling, auto play and push notifications are explicitly described as interface features that, in combination with personalised recommendations, form an "algorithmic addiction loop" [8958037].}\pdfcomment[icon=Note]{Infinite scrolling, autoplay and push notifications as interface features contributing to an 'algorithmic addiction loop' are explicitly described in Li (8958037).} \hldarkgreen{The European Commission's investigation into TikTok under the Digital Services Act highlights similar concerns, with the DSA inquiry focusing on infinite scroll, autoplay and notifications as manipulative design that impairs users' ability to make free and informed decisions [8958037].}\pdfcomment[icon=Note]{Li reports that the EU investigation into TikTok under the DSA raised concerns about infinite scroll, autoplay and notifications as manipulative design (8958037).} \hldarkgreen{College students, situated at a developmental stage with unstable self identity and still maturing executive functions, are described as particularly vulnerable to these loops, with excessive use linked to diminished academic concentration, weakened social skills and heightened anxiety and depression [8958037].}\pdfcomment[icon=Note]{The heightened vulnerability of college students and associations with diminished concentration, social skills deficits, anxiety and depression are discussed in the short-video addiction study (8958037).} \hllightgreen{Micro learning oriented designs that break continuous engagement into delineated segments, with transparent objectives and explicit stopping cues, directly counter the patterns identified in these studies.}\pdfcomment[icon=Note]{The statement that micro-learning designs with delineated segments counter the problematic patterns is a reasoned design inference grounded in the addiction research and gamification/HITL findings (8958037; 5645864; 6805789).} \hllightgreen{Emotional AI safety research adds another dimension: positive habit loops should not only address time allocation but also emotional dependency.}\pdfcomment[icon=Note]{Emotional AI safety literature adds the need to consider emotional dependency in habit-loop design; this dimension is emphasized in Song's emotional AI work (9738431) and so is a supported synthesis.} \hldarkgreen{Song outlines a "bi‑weekly psychological assessment" module in an emotional AI system that adapts usage based on risk levels [9738431].}\pdfcomment[icon=Note]{Song explicitly outlines a bi-weekly psychological assessment module that adapts usage based on risk levels (9738431, module description).} \hldarkgreen{For adolescents at high risk of depression, the system enforces this scheduled assessment and restricts late night emotional topics to avoid compounding psychological burdens [9738431].}\pdfcomment[icon=Note]{The enforcement of scheduled assessments and restriction of late-night emotional topics for at-risk adolescents is described in the emotional AI source (9738431, risk-mitigation measures).} \hldarkgreen{A dynamic permission model integrates with medical institution data so that when a user is diagnosed with a mental illness, the system switches into a "treatment mode" that limits non professional AI emotional guidance and opens API interfaces to hospital prescription systems, ensuring that users receive professional support rather than unbounded, pseudo therapeutic interaction [9738431].}\pdfcomment[icon=Note]{The dynamic permission model that integrates with medical data and switches to a 'treatment mode' limiting non-professional AI guidance and opening APIs to hospital systems is explicitly described in Song (9738431).} \hldarkgreen{Before market release, such products must pass a "mental health impact simulation test" that uses generative adversarial networks to model long term effects on vulnerable groups, such as individuals with social anxiety, thereby identifying potential harms pre deployment [9738431].}\pdfcomment[icon=Note]{The requirement that products pass a 'mental health impact simulation test' using GANs to model long-term effects on vulnerable groups before market release is specified in the emotional AI source (9738431).} \hllightgreen{Although these mechanisms are situated in the domain of emotional AI, they point toward micro learning and habit loops that are explicitly tailored to user risk profiles, include enforced rest or reduced interaction in evenings, and provide escalation paths to human care where necessary.}\pdfcomment[icon=Note]{This is a synthesis: although the mechanisms are from emotional AI, the conclusion that they point toward risk-sensitive micro-learning loops with enforced rest and escalation is a reasonable inference from 9738431.} \hldarkgreen{Independent learning among autistic individuals introduces further design constraints on positive habit loops.}\pdfcomment[icon=Note]{The assertion that independent learning among autistic individuals imposes further design constraints is supported by the autism-focused review (6454369) which highlights specific needs. } \hldarkgreen{Adako and colleagues describe how independent learning for persons on the autism spectrum depends on structured tasks, predictable routines, and AI tools that assist with planning, time management and self advocacy [6454369].}\pdfcomment[icon=Note]{Adako et al.'s findings that independent learning for persons on the autism spectrum depends on structured tasks, predictable routines, and AI tools assisting planning/time management/self-advocacy are explicitly documented (6454369, section 4).} \hldarkgreen{Systems must account for sensory overload and executive function challenges, which implies interfaces that allow the configuration of batch sizes, pacing and notification frequencies according to individual preferences and communication needs [6454369].}\pdfcomment[icon=Note]{The need to account for sensory overload and executive function challenges, implying configurable batch sizes, pacing and notification frequencies, is directly discussed as design implications in the ASD literature (6454369).} \hllightgreen{For social recommender systems, this evidence suggests that micro learning sequences should never be a fixed engagement template.}\pdfcomment[icon=Note]{The recommendation that micro learning sequences should never be a fixed engagement template for social recommenders follows logically from the neurodiversity evidence and is a reasonable synthesis (6454369; 9283805).} \hllightgreen{Users, particularly neurodivergent ones, need controls over how many items appear in a batch, how often prompts are delivered, and whether certain sensory or social stimuli are filtered or slowed, so that habit formation aligns with autonomy and comfort rather than generic engagement curves [6454369][7][9].}\pdfcomment[icon=Note]{The need for user controls over batch size, prompt frequency and filtering of sensory/social stimuli is supported by ASD design guidance and token-economy reviews (6454369; 9283805; 5645864), so this is a valid synthesis.} \hldarkgreen{Research on organisational learning and ethical AI marketing strategy adds a governance layer relevant to how micro learning loops are conceived and managed inside firms.}\pdfcomment[icon=Note]{The claim that organisational learning and ethical AI marketing research adds a governance layer relevant to micro-learning loops is supported by the organisational literature (8200246).} \hldarkgreen{Asante and Adeoba argue that ethical AI practices are context dependent and shaped by organisational cultures, stakeholder expectations and governance structures interpreted through Communities of Practice and corporate social responsibility (CSR) frames [8200246].}\pdfcomment[icon=Note]{Asante and Adeoba explicitly argue that ethical AI practices are context dependent and shaped by organisational cultures, stakeholder expectations and Communities of Practice/CSR frames (8200246).} \hldarkgreen{Their qualitative, secondary research design shows that collaborative learning within organisations, rather than static rulebooks, supports the development of ethical mechanisms for AI driven marketing [8200246].}\pdfcomment[icon=Note]{Their qualitative secondary research design and conclusion that collaborative learning within organisations supports development of ethical mechanisms is described in the source (8200246).} \hllightgreen{Applying this to micro learning and positive habit loops in social recommendation, design teams would need to treat engagement architectures as ongoing sites of organisational reflection, using internal communities of practice to scrutinise whether micro modules genuinely serve user learning and wellbeing or simply repackage engagement optimisation under an ethical vocabulary.}\pdfcomment[icon=Note]{Applying organisational learning to treat engagement architectures as ongoing sites of reflection is a reasonable inference from 8200246 and governance literature; this is synthesis rather than a direct quotation.} \hllightgreen{Regulatory work on AI in credit scoring, although in a different sector, underscores structural themes that can be transposed to engagement loop auditing.}\pdfcomment[icon=Note]{Regulatory work on AI in credit scoring highlighting audit and fairness themes that can be transposed to engagement auditing is a reasonable cross-domain inference grounded in 2855351 and related regulatory sources.} \hldarkgreen{Cao argues that algorithmic fairness should be institutionalised through mandated periodic fairness audits, bias detection tests and corrective actions, supported by inclusive data practices and cross sector collaboration [2855351].}\pdfcomment[icon=Note]{Cao's argument that algorithmic fairness should be institutionalized through mandated periodic audits, bias tests and corrective actions is directly stated in the credit-scoring/regulatory source (2855351, policy recommendations).} \hldarkgreen{Transparency techniques such as SHAP and LIME are cited as ways to improve stakeholders' trust and enable meaningful oversight of opaque models [2855351].}\pdfcomment[icon=Note]{The citation of transparency techniques such as SHAP and LIME to improve stakeholder trust and oversight appears explicitly in the procedural recommendations of the credit-scoring literature (2855351).} \hllightgreen{While these methods address financial decision making rather than social engagement, the procedural lesson is clear: micro learning loops and token economies within social recommenders should be subjected to regular audits that examine differential impacts across demographic and psychological groups, with explicit remediation when harms are detected.}\pdfcomment[icon=Note]{While SHAP/LIME address finance, the procedural lesson to audit micro learning loops and token economies for differential impacts is a reasonable policy transfer inference from 2855351.} \hldarkgreen{Cross sector and international collaboration is recommended to develop unified standards for fairness, transparency and accountability, which would equally apply to interaction design parameters such as recommended session length or reward contingencies [2855351].}\pdfcomment[icon=Note]{Cross-sector and international collaboration to develop unified standards for fairness/transparency/accountability is explicitly recommended in the credit-regulatory literature (2855351).} \hllightgreen{Bringing these threads together, a pro social engagement model grounded in micro learning and positive habit loops would draw on several evidenced design principles: - Finite, goal oriented sequences: Sessions structured as short tasks with explicit objectives, transparent progression criteria and end points, as in systemic gamification and HITL grading [9][21]. - Meaningful rewards and narratives: Token systems that link immediate feedback to creative agency and longer term goals, avoiding indiscriminate point inflation and punitive feedback [7]. - Transparent AI involvement and logs: Clear labelling of automated versus human contributions, dual logging of interactions and outcomes, and documentation of session boundaries to support external audits and accountability in line with AI Law Model expectations [21][25][7461574]. - Risk sensitive modulation: Assessment modules and dynamic permissions that adapt interaction frequency and content restrictions to mental health risks, with escalation to professional care where indicated [9738431]. - User configurable pacing and sensory load: Controls for batch size, pacing and notification patterns, tailored to cognitive and sensory profiles of users, including autistic learners [6454369]. - Independent oversight and standards alignment: External audits, oversight committees and alignment with regional regulatory frameworks such as the EU AI Act and OECD AI Principles, to ensure that engagement patterns do not exploit vulnerabilities in the way short video platforms currently do [7461574][8958037].}\pdfcomment[icon=Note]{This multi-point design summary coherently synthesizes principles that are each supported across the cited sources (systemic gamification, HITL, token economies, AI Law Model, emotional AI, ASD literature and governance papers), so it is an appropriate LIGHTGREEN synthesis.} \hllightgreen{Evidence from education, emotional AI, mental health safety, digital governance and organisational learning converges on a core insight: micro learning and positive habit loops can sustain engagement while reducing algorithmic addiction risks if they are designed as bounded, meaningful, audited episodes that respect cognitive limits and democratic accountability [7][9][21][25][6454369][7461574][8958037][9738431][8200246][2855351].}\pdfcomment[icon=Note]{The convergence claim is a reasonable high-level synthesis: the cited evidence from education, emotional AI, governance and organisational learning supports the core insight that bounded, meaningful, audited episodes can reduce addiction risk (sources: 5645864; 9738431; 2977039; 6454369; 7461574; 8958037; 6805789; 8200246; 2855351).} \hldarkgreen{Practical design work should also recognise that social-aware recommendation techniques can exploit network structure and learned edge weights to prioritise items or peers that map neatly onto short sessions; supervised graph-based methods appear to offer higher precision and better ranking metrics in friend and item recommendation tasks, which may help operationalise finite, goal-aligned batches in practice \parencites[30]{2882013}.}\pdfcomment[icon=Note]{The technical claim that social-aware recommendation techniques can exploit network structure and that supervised graph-based methods often yield higher precision/ranking metrics is supported by recommender systems literature demonstrating supervised graph/link-prediction methods outperforming unsupervised baselines (2882013, multiple sections on supervised random walks and graph-based methods).}

\subsubsection{Advertiser Alignment with Digital Wellbeing}
\label{sec:advertiser_alignment_digital_wellbeing}
\hllightgreen{Advertisers that rely on engagement driven social media environments are increasingly exposed to the same structural conflicts between attention maximisation and mental health that have been documented for recommender systems more broadly.}\pdfcomment[icon=Note]{Sources (e.g., 6502223, 3303310, 5142478) document structural tensions between engagement maximisation and harms from recommender systems, supporting this synthesis (6502223 p.2250; 3303310 p.19; 5142478 p.25).} \hldarkgreen{Evidence from qualitative and empirical work shows that algorithmic personalisation tends to concentrate user exposure on emotionally charged and sensational content, which increases time on platform but correlates with depressive symptoms, anxiety and compulsive usage patterns, particularly in adolescents and vulnerable groups [3303310][5142478].}\pdfcomment[icon=Note]{Empirical and qualitative findings cited explicitly report concentration on sensational content and links to depressive/anxiety symptoms and compulsive use, particularly for adolescents and vulnerable groups (3303310 p.19-20; 5142478 p.25).} \hllightgreen{In parallel, dark commercial patterns in e‑commerce and subscription flows exploit cognitive biases to drive conversions and data capture, reinforcing a monetisation logic that depends on friction, opacity and behavioural targeting rather than informed choice [6295677].}\pdfcomment[icon=Note]{The literature on dark patterns and manipulative e-commerce flows describing exploitation of cognitive biases and monetisation logics is documented in the India dark patterns study (6295677 p.2, p.6).} \hllightgreen{Aligning advertising with digital wellbeing within such an environment therefore requires advertisers to reposition themselves away from these exploitative optimisation levers and toward transparent, diversity aware and consent respecting practices that remain commercially viable under the EU's tightening regulatory regime, including the Digital Services Act and planned Digital Fairness reforms [6502223][7896734].}\pdfcomment[icon=Note]{Sources describe regulatory tightening (DSA, proposed Digital Fairness reforms) and argue advertisers must shift practices; the claim synthesises those sources (7896734 p.390-391; 6502223 p.2250).} \hldarkgreen{Sanikidze's analysis of the Digital Fairness Act proposal is particularly explicit that "addictive product designs and unfair personalisation" are being targeted as unfair commercial practices, in complement to DSA and DMA provisions on platform responsibility and dark patterns [7896734].}\pdfcomment[icon=Note]{Sanikidze (7896734) explicitly links proposed Digital Fairness measures to targeting 'addictive product designs and unfair personalisation' alongside DSA/DMA provisions (7896734 p.390-391).} \hllightgreen{This creates a forward looking constraint on advertisers that still depend on campaigns which are only profitable when embedded in manipulative UX flows or hyper‑personalised feeds optimised solely for engagement.}\pdfcomment[icon=Note]{This is a logical inference from cited regulatory pressures and the documented dependence of some campaigns on manipulative UX and hyper-personalisation (6502223 p.2250; 6295677 p.2).} \hldarkgreen{As regulators scrutinise design features that make cancelling subscriptions harder than signing up or that push users through deceptive verification and consent flows, the surfaces through which advertisers previously harvested impressions and data without meaningful consent become legally risky [7896734].}\pdfcomment[icon=Note]{Sanikidze and related analysis explicitly note regulators scrutinise cancellation friction and deceptive flows, making such practices legally risky (7896734 p.390-391).} \hllightgreen{In this regulatory trajectory, advertising strategies that assume limitless tracking, maximised session length and opaque behavioural profiling are exposed as contingent rather than stable business assets.}\pdfcomment[icon=Note]{Sources describe a regulatory trajectory that limits tracking and manipulative features, supporting the synthesis that such assets are contingent (6502223 p.2250; 7896734 p.390).} \hldarkgreen{Qualitative work on dark patterns in India illustrates how advertising demand and interface design have historically co‑evolved.}\pdfcomment[icon=Note]{The India-focused qualitative work on dark patterns (6295677) documents historical co-evolution of advertising demand and interface design in that context (6295677 p.2, p.6).} \hldarkgreen{Khare et al. describe how major platforms in sectors such as fashion, general e‑commerce and cryptocurrency frequently deploy default enrolment, obstructive cancellation and misleading urgency cues to trigger System 1 decision making and suppress reflective reasoning, precisely because these techniques increase conversion rates and data capture that can be fed into behavioural targeting models [6295677].}\pdfcomment[icon=Note]{Khare et al. (reported in 6295677) detail usage of defaults, obstructive cancellations and urgency cues to trigger System 1 and increase conversions/data capture (6295677 p.2, p.6).} \hldarkgreen{Survey data reported in the same study indicate that 55 percent of respondents experienced unexpected charges and 47 percent faced difficulties when attempting to cancel services, while only 39 percent even recognised the term "dark patterns," which exposes a pronounced information asymmetry between platforms, advertisers and consumers [6295677].}\pdfcomment[icon=Note]{The survey statistics (55\% unexpected charges, 47\% cancellation difficulties, 39\% unfamiliar with 'dark patterns') are reported directly in the India study (6295677 p.2).} \hllightgreen{For advertisers that buy placements or affiliate traffic in such environments, brand value becomes entangled with design practices now under regulatory challenge in the EU and other jurisdictions.}\pdfcomment[icon=Note]{Given the regulatory challenges described (7896734, 6502223) and the entanglement of brand value with platform practices, this inference is supported by the sources (7896734 p.390; 6502223 p.2250).} \hllightgreen{Gerr's review of algorithmic personalisation and media diversity provides a complementary perspective on the risks of advertising in engagement maximising feeds.}\pdfcomment[icon=Note]{Gerr's review (5142478) provides a complementary perspective on risks of advertising in engagement-maximising feeds, matching this claim (5142478 p.25).} \hldarkgreen{AI driven content curation in social networks is described as systematically narrowing exposure to diverse viewpoints because algorithms prioritise content that matches existing preferences and often sensational or polarising material, which increases engagement but reduces pluralism [5142478].}\pdfcomment[icon=Note]{Gerr explicitly describes AI-driven curation narrowing exposure and prioritising sensational/polarising content that increases engagement and reduces pluralism (5142478 p.25).} \hldarkgreen{Filter bubbles and echo chambers emerge as structural properties of these systems: users are intellectually isolated and encounter fewer opposing views, which weakens critical thinking and civic engagement [5142478].}\pdfcomment[icon=Note]{Gerr (5142478) discusses filter bubbles and echo chambers leading to intellectual isolation and reduced exposure to opposing views, weakening civic engagement (5142478 p.25).} \hllightgreen{From an advertiser perspective, such feeds may seem attractive because they deliver high click‑through rates within homogenous audiences, but this same narrowing creates reputational exposure when campaigns are repeatedly co‑present with polarising content, body negative imagery or outrage‑heavy narratives that have measurable negative effects on mental health and social cohesion [3303310][5142478].}\pdfcomment[icon=Note]{This combines evidence on high engagement from homogenous audiences and documented reputational/mental health harms when ads appear with polarising/content (3303310 p.19-20; 5142478 p.25).} \hllightgreen{Empirical evidence from Hu's synthesis of mental health research on personalised feeds reinforces this point.}\pdfcomment[icon=Note]{Hu's synthesis is cited later (3303310) and the sentence previews empirical reinforcement found in that source (3303310 p.19-20).} \hldarkgreen{Hyper‑personalised recommendation environments are linked to a 23 percent higher incidence of depressive symptoms compared to chronological feeds, with particularly pronounced effects among adolescents and individuals predisposed to social comparison [3303310].}\pdfcomment[icon=Note]{Hu's synthesis reports the 23\% higher incidence of depressive symptoms in hyper-personalised feeds versus chronological feeds, with pronounced adolescent effects (3303310 p.20-20/3).} \hldarkgreen{Neuroimaging studies referenced in the same source correlate "doomscrolling" with reduced prefrontal cortex activity, which parallels addictive behaviour patterns, while Instagram's prioritisation of Eurocentric beauty standards is associated with statistically significant declines in self esteem among adolescent women of colour [3303310].}\pdfcomment[icon=Note]{The neuroimaging correlations with doomscrolling and reduced prefrontal activity, and Instagram's link to declines in self-esteem among adolescent women of colour, are reported in Hu (3303310 p.20-20/3).} \hllightgreen{Advertising campaigns that rely on such emotionally volatile content streams for reach and conversion are therefore structurally misaligned with digital wellbeing objectives, whether assessed from a public health or regulatory risk standpoint.}\pdfcomment[icon=Note]{This is a reasonable synthesis: sources document harms from emotionally volatile streams and regulatory/public health risk implications (3303310 p.19-20; 6502223 p.2250).} \hldarkgreen{Chitra et al. argue that these harms are not incidental but flow directly from engagement maximising objective functions in recommender systems [6502223].}\pdfcomment[icon=Note]{Chitra et al. (6502223) explicitly argue harms flow from engagement-maximising objective functions in recommender systems (6502223 p.2254; p.2250).} \hldarkgreen{The structural incentive is to weight content that increases immediate interaction, which produces selection pressure toward emotionally provocative and epistemically degraded material [6502223].}\pdfcomment[icon=Note]{6502223 describes the structural incentive to weight content that increases immediate interaction, producing selection pressure toward provocative material (6502223 p.2254; p.2250).} \hllightgreen{For advertisers, this means that traditional optimisation metrics such as impressions, click‑through and dwell time often track the same dynamics that regulators and mental health researchers identify as problematic: outrage amplification, sensationalism and addictive loops [3303310][5142478].}\pdfcomment[icon=Note]{This synthesises empirical links between optimisation metrics and problematic dynamics as discussed in Hu and Gerr (3303310 p.19-20; 5142478 p.25).} \hllightgreen{Aligning with digital wellbeing therefore requires advertisers to support an alternative objective function for recommendation and ad delivery that does not simply proxy engagement.}\pdfcomment[icon=Note]{Given the documented misalignment, the claim that advertisers must support alternative objective functions is a reasoned conclusion supported by sources (6502223 p.2252; 3303310 p.19-20).} \hldarkgreen{Chitra et al. propose concrete technical design standards that offer a template for such a shift: value aligned objective functions that integrate explicit terms for epistemic quality, diversity incentives and engagement ceilings [6502223].}\pdfcomment[icon=Note]{Chitra et al. lay out technical standards including value-aligned objective functions with epistemic quality, diversity incentives and engagement ceilings (6502223 p.2252-2254).} \hldarkgreen{Content from highly accurate sources, as validated by independent fact checking, should receive positive weights, while sources with established misinformation patterns are penalised [6502223].}\pdfcomment[icon=Note]{6502223 explicitly recommends accuracy weighting and penalising misinformation sources as part of value-aligned objective functions (6502223 p.2252).} \hldarkgreen{Recommendation sequences are required to meet intra session diversity requirements so that users are not confined to a narrow ideological spectrum, and engagement ceilings cap the amplification of individual items to avoid virality driven escalation [6502223].}\pdfcomment[icon=Note]{Chitra et al. specify intra-session diversity requirements and engagement ceilings to limit virality-driven escalation (6502223 p.2252-2254).} \hllightgreen{When advertisers accept placement and optimisation rules that follow these standards, they implicitly trade some short term engagement volume for association with trustworthy, less polarising content environments, which is more consistent with mental health safeguards identified in Hu's analysis of sentiment balanced feeds and usage friction tools [3303310].}\pdfcomment[icon=Note]{This is a synthesis: accepting such standards implies trading short-term engagement for safer environments, which Hu documents as beneficial (3303310 p.20; 6502223 p.2252).} \hllightgreen{The same governance framework stresses the need for mandatory algorithmic auditing and research data access, which are directly relevant for advertisers seeking to realign their strategies.}\pdfcomment[icon=Note]{The governance framework in 6502223 stresses mandatory algorithmic auditing and research data access relevant to advertisers (6502223 p.2252-2254).} \hldarkgreen{Chitra et al. recommend that large platforms submit their recommendation algorithms to annual independent third party audits that assess political amplification asymmetries, misinformation mitigation and performance against diversity and pluralism benchmarks [6502223].}\pdfcomment[icon=Note]{6502223 specifically recommends annual independent third-party audits assessing political amplification, misinformation mitigation and diversity benchmarks (6502223 p.2252).} \hldarkgreen{Audit methodologies and high level findings are to be published, echoing safety audit practices where independent bodies evaluate whether operational reality matches stated policies [5345089][6502223].}\pdfcomment[icon=Note]{6502223 recommends publishing audit methodologies and findings and 5345089 discusses analogous safety audit publication practices (6502223 p.2252; 5345089 p.6).} \hllightgreen{For advertisers, this creates an emerging transparency layer: campaigns can be evaluated not only on engagement metrics but also on how often they appear in contexts that trigger high scores on indices such as Hu's Algorithmic Impact Assessment Scale, which quantifies mental health risks along dimensions like emotional volatility, sleep disruption and compulsive usage patterns [3303310].}\pdfcomment[icon=Note]{Hu describes AIAS-5 and mental health risk indices that advertisers could use to evaluate context; this is a reasonable extension to advertiser evaluation (3303310 p.20; 3303310 p.19-20).} \hllightgreen{Aligning with digital wellbeing in this setting means selecting inventory and optimisation partners whose audited systems meet diversity and risk thresholds that are compatible with regulatory expectations under the DSA and emerging Digital Fairness norms [6502223][7896734].}\pdfcomment[icon=Note]{Combining Chitra et al. on audits/standards and Sanikidze on regulatory expectations supports the claim about selecting audited partners compatible with DSA/Digital Fairness (6502223 p.2252; 7896734 p.390).} \hllightgreen{Ethical AI governance research extends these obligations by specifying normative constraints on profiling and data use that directly affect targeted advertising.}\pdfcomment[icon=Note]{Ethical AI governance literature (e.g., 2977039) indeed specifies normative constraints on profiling and data use affecting targeted advertising (2977039 p.20-21; p.53).} \hldarkgreen{Kostenko's AI Law Model introduces the principle of "information invisibility," which frames an extended right not to generate, transmit, accumulate or store digital traces in contexts where such fixation is not necessary for basic service operation [2977039].}\pdfcomment[icon=Note]{Kostenko's AI Law Model (source 2977039) introduces 'information invisibility' as a principle protecting against unnecessary generation/retention of digital traces (2977039 p.167; p.54).} \hldarkgreen{Users are entitled to refuse permanent identifiers such as device fingerprints or cookie profiles, to opt out of systems demanding disproportionate personal or behavioural data, and to ignore digital requests (including marketing proposals and recommendation systems) that aim to create or refine psycho emotional or consumer profiles without their direct initiative and control [2977039].}\pdfcomment[icon=Note]{2977039 defines rights to refuse permanent identifiers and opt out of disproportionate data collection and to ignore requests aimed at building psycho-emotional profiles (2977039 p.167; p.54).} \hldarkgreen{Digital systems must in turn implement full fledged invisibility modes, including incognito, temporary presence and contextual masking, supported by clear visual indicators and logs that allow users to monitor what data are being stored or transmitted [2977039].}\pdfcomment[icon=Note]{2977039 mandates implementation of invisibility modes (incognito, temporary presence, contextual masking) with visual indicators and logs (2977039 p.168; p.242).} \hllightgreen{Advertisers that intend to align with digital wellbeing must design campaigns and data strategies that remain effective even when a significant fraction of users operate under such invisibility regimes, which materially constrains behavioural targeting and cross platform tracking.}\pdfcomment[icon=Note]{This is a logical inference: if users use invisibility modes (2977039), advertisers must adapt campaigns and data strategies accordingly (2977039 p.167-168).} \hldarkgreen{The AI Law Model also prohibits the use of personal data for advertising or profiling that restricts freedom of choice without explicit informed consent, with special emphasis on children and other vulnerable groups [2977039].}\pdfcomment[icon=Note]{2977039 explicitly prohibits use of personal data for advertising/profiling that restricts freedom without informed consent, emphasizing protection for children and vulnerable groups (2977039 p.21; p.167).} \hldarkgreen{For minors, targeted advertising based on psychotypes, neurobehavioral indicators or predicted emotional states is categorically banned, and information invisibility modes should enable service use without profiling or cross platform tracking [2977039].}\pdfcomment[icon=Note]{The AI Law Model directly states targeted advertising for minors based on psychotypes/neurobehavioral indicators is banned and invisibility modes should enable use without profiling (2977039 p.21; p.167).} \hldarkgreen{Sanikidze emphasises that DSA restrictions on targeting minors and Digital Fairness proposals aimed at addictive designs and unfair personalisation move in the same direction: gradual elimination of opaque profiling of vulnerable users as an acceptable advertising practice in the EU [7896734].}\pdfcomment[icon=Note]{Sanikidze highlights DSA restrictions and Digital Fairness proposals targeting addictive designs and unfair personalisation, aligning with elimination of opaque profiling of vulnerable users (7896734 p.390-391).} \hllightgreen{Algorithmic systems can also reproduce and amplify harmful stereotypes embedded in training data, for example by skewing search and recommendation outputs toward sexualised or demeaning representations of marginalised groups, a pattern that may suggest deeper data and value failures in platform design \parencites[4]{3278906}.}\pdfcomment[icon=Note]{Multiple sources discuss algorithmic amplification of stereotypes and sexualised outputs in search/recommendation, supporting this claim (3278906 p.4-5; 5142478 p.25).} \hlorange{Within the working taxonomy developed in this thesis, advertising strategies that depend on such profiling and on dark patterns would be assigned to a high or unacceptable risk category, labelled Code Category C, that signals significant legal exposure under the DFA trajectory and related norms.}\pdfcomment[icon=Note]{The assignment to a specific thesis taxonomy (Code Category C) is an internal classification claim of the report and is not directly documented in the provided external sources.} \hllightgreen{Aligning advertisers with digital wellbeing, under this working taxonomy, involves articulating alternative engagement and monetisation models that satisfy three concurrent demands: compliance with emerging EU rules on manipulation and data use, measurable reduction of mental health risks as captured in tools like AIAS‑5, and commercial sustainability for brands and platforms [3303310][6502223][7896734].}\pdfcomment[icon=Note]{This threefold alignment (compliance, measurable risk reduction via tools like AIAS-5, commercial sustainability) synthesises claims made across Hu, Chitra et al., and Sanikidze (3303310 p.20; 6502223 p.2252; 7896734 p.390).} \hllightgreen{Qualitative evidence from Hu's review points to potential design directions.}\pdfcomment[icon=Note]{Hu's review contains qualitative evidence pointing to design directions and so this statement is a supported preview (3303310 p.20).} \hldarkgreen{Platforms that embed wellbeing guardrails such as sentiment balanced content curation and usage friction tools see reductions in self reported anxiety by 17 percent within six weeks, with effects mediated by demographics and vulnerability to algorithmic bias [3303310].}\pdfcomment[icon=Note]{Hu reports reductions in self-reported anxiety by 17\% within six weeks from platforms embedding wellbeing guardrails like sentiment-balanced curation (3303310 p.20).} \hllightgreen{For advertisers, operating in such guarded environments may reduce raw impression volumes but can also improve the psychological "quality" of attention received, since users are less likely to be trapped in doomscrolling loops or exposed to content that damages self esteem and trust [3303310][5142478].}\pdfcomment[icon=Note]{Sources indicate guarded environments reduce harmful exposures and can improve attention quality, a reasonable inference for advertisers (3303310 p.20; 5142478 p.25).} \hllightgreen{A practical alignment pathway can be articulated in terms of specific campaign and product design commitments.}\pdfcomment[icon=Note]{This is a synthesis move: the report and sources propose concrete campaign/product commitments as practical pathways (6502223 p.2252; 3303310 p.20).} \hldarkgreen{Advertisers can, for instance, restrict their bidding to inventory segments that are configured with engagement ceilings and diversity incentives, as described by Chitra et al., and that are subject to regular third party audits with published findings [6502223].}\pdfcomment[icon=Note]{6502223 specifically describes restricting bidding to inventory with engagement ceilings/diversity incentives and subject to audits, matching this prescription (6502223 p.2252-2254).} \hldarkgreen{They can also require platform partners to offer information invisibility compliant modes for users, even if this reduces the availability of third party cookies or device identifiers for profiling [2977039].}\pdfcomment[icon=Note]{2977039 requires information invisibility modes and the sentence about requiring platforms to offer such modes (reducing cookies/identifiers) follows directly (2977039 p.167; p.168).} \hllightgreen{From a content standpoint, advertisers can favour placements in feeds that adopt sentiment balancing and "digital nutrition labels" that disclose predicted psychological impact scores, similar to the warning style disclosures Hu notes for personalised feeds [3303310].}\pdfcomment[icon=Note]{Hu discusses digital nutrition labels and predicted psychological impact disclosures; the advertiser recommendation is a reasonable extension from Hu (3303310 p.20).} \hllightgreen{Such labels give both users and brands visibility into whether a given content channel tends to increase emotional volatility or sleep disruption, which allows advertisers to avoid placements in high‑risk segments even before formal regulatory prohibitions emerge.}\pdfcomment[icon=Note]{This logically follows from the idea of digital nutrition labels and impact scores described by Hu, supporting advertisers avoiding high-risk placements (3303310 p.20).} \hllightgreen{The interplay between auditing and organisational culture in AI systems also provides lessons for advertiser alignment.}\pdfcomment[icon=Note]{The interplay between auditing and organisational culture is discussed in the AI audit/safety audit literature and supports this claim (5345089 p.6).} \hldarkgreen{Mökander's discussion of safety audits shows that effective assessments go beyond technical checks and extend to whether organisational practices and culture genuinely support safety goals [5345089].}\pdfcomment[icon=Note]{Mökander (5345089) explicitly argues that effective safety/AI assessments extend beyond technical checks into organisational practices and culture (5345089 p.6).} \hldarkgreen{Safety audits use multiple tools, from checklists to interviews, and derive authority from independent institutions such as occupational safety regulators [5345089].}\pdfcomment[icon=Note]{5345089 documents that safety audits use multiple tools (checklists, interviews) and derive authority from independent institutions, supporting this statement (5345089 p.6).} \hllightgreen{Translating this insight into advertising governance suggests that brands should not only commission independent AI audits of their own marketing technologies but also reform internal incentive structures so that campaign teams are rewarded for meeting wellbeing aligned metrics, such as reduced association with high AIAS‑5 risk scores, rather than solely for click‑through or conversion rates [3303310][5345089].}\pdfcomment[icon=Note]{Combining Hu's metrics and safety audit lessons supports the recommendation advertisers reform incentives and commission audits (3303310 p.20; 5345089 p.6).} \hllightgreen{TikTok's opacity, as described by Chitra et al., illustrates the importance of research data access for any meaningful advertiser alignment strategy.}\pdfcomment[icon=Note]{Chitra et al. highlight TikTok's opacity and the need for research data access, illustrating its importance for advertiser alignment (6502223 p.2252; p.2254).} \hldarkgreen{The For You Page is portrayed as exceptionally effective at capturing attention and retaining users, yet TikTok resists meaningful transparency into its recommendation algorithm and defends it as a trade secret [6502223].}\pdfcomment[icon=Note]{6502223 portrays TikTok's For You Page as exceptionally effective at capturing attention and notes TikTok resists transparency citing trade secrets (6502223 p.2252; p.2254).} \hldarkgreen{National security debates around the potential for state influenced recommendation underscore how dependent advertisers are on systems they cannot independently scrutinise [6502223].}\pdfcomment[icon=Note]{6502223 discusses national security debates and attempted bans tied to concerns about state influence, showing advertiser dependence on opaque systems (6502223 p.2252; p.2254).} \hldarkgreen{Under the governance recommendations in the same study, platforms would provide vetted researchers with representative samples of recommendation sequences and other data that allow systemic analysis of political and psychological effects [6502223].}\pdfcomment[icon=Note]{6502223 recommends platforms provide vetted researchers with representative samples of recommendation sequences for systemic analysis (6502223 p.2252-2254).} \hllightgreen{Advertisers seeking alignment with digital wellbeing should support and condition their spend on such research data access, recognising that without independent insight into recommendation dynamics they cannot credibly claim that their campaigns are not contributing to addictive usage patterns or polarisation.}\pdfcomment[icon=Note]{This is a normative inference consistent with 6502223's recommendations that advertisers should support research data access to validate absence of harms (6502223 p.2252).} \hllightgreen{Within this matrix of regulatory and governance instruments, advertiser alignment with digital wellbeing is not framed as a voluntary corporate social responsibility add on but as a strategic adaptation to an environment where regulatory costs, reputational risk and empirical evidence on mental health converge.}\pdfcomment[icon=Note]{Sources document converging pressures from regulation, reputational risk and evidence on mental health, supporting framing advertiser alignment as strategic adaptation (7896734 p.390; 3303310 p.19-20; 6502223 p.2250).} \hldarkgreen{Sanikidze underscores that EU rules on dark patterns and AI driven manipulation already make the Union one of the most proactive regions in constraining addictive designs [7896734].}\pdfcomment[icon=Note]{Sanikidze explicitly states the EU is proactive in constraining dark patterns and AI manipulation under current/forthcoming rules (7896734 p.390-391).} \hldarkgreen{Hermansyah's work on DMA interoperability, while focused on switching behaviour, demonstrates that user centred choice architectures can materially change user decisions when backed by enforcement against manipulative design [1070605].}\pdfcomment[icon=Note]{Hermansyah's DMA analysis shows interoperability and choice architecture change switching behaviour and reduce manipulative design effects when enforced, as cited (1070605 p.160).} \hllightgreen{In such a context, advertisers that cling to opaque engagement spikes and behavioural profiling of vulnerable users risk stranded assets as platform governance tightens and user expectations shift.}\pdfcomment[icon=Note]{Given regulatory tightening and documented risks, the risk of stranded assets for advertisers relying on opaque profiling is a reasonable synthesis (6502223 p.2250; 7896734 p.390).} \hllightgreen{By contrast, brands that support value aligned recommendation objectives, independent AI audits, information invisibility rights and sentiment balanced content environments position themselves as partners in a transition toward pro social recommender systems [2977039][3303310][6502223].}\pdfcomment[icon=Note]{Combining the AI law model, Hu and Chitra et al. supports the claim that brands adopting value-aligned objectives and audits position themselves positively (2977039 p.167; 3303310 p.20; 6502223 p.2252).} \hllightgreen{Their campaigns will increasingly be evaluated not only on financial return but also on their contribution to or mitigation of algorithmic addiction, as measured through public mental health indicators and regulatory enforcement records.}\pdfcomment[icon=Note]{This is a projection consistent with sources showing increasing regulatory and public-health-based evaluation criteria beyond financial return (3303310 p.19-20; 6502223 p.2250).} \hllightgreen{This report's working taxonomy makes that tension explicit by assigning dark pattern dependent, high profiling strategies to Code Category C, signalling that aligning advertising with digital wellbeing is both a legal risk management imperative and a necessary condition for sustainable monetisation models under the EU Digital Fairness Act trajectory [6295677][7896734].}\pdfcomment[icon=Note]{The report's taxonomy claim references earlier evidence on dark-pattern dependence and DFA regulatory trajectory; sources support the legal risk framing (6295677 p.2; 7896734 p.390).} \hldarkgreen{Methodological limitations should be acknowledged: assessments of proprietary recommender algorithms in this section rely exclusively on published secondary studies, policy analyses and UI/UX reporting rather than internal platform telemetry or undisclosed machine level access.}\pdfcomment[icon=Note]{The methodological limitation is openly stated and matches the nature of the cited studies (many are secondary, qualitative and policy analyses) as reflected across sources (6295677 p.2; 3303310 p.19).} \hldarkgreen{Where platform mechanics are described, citations point to publicly available qualitative or regulatory sources; no primary reverse engineering or internal data was used.}\pdfcomment[icon=Note]{This procedural clarification (using publicly available qualitative/regulatory sources, no primary reverse engineering) is consistent with methods described in the cited studies (6295677 p.2; 3303310 p.19).}

\newpage
\printbibliography

\end{document}
